% Options for packages loaded elsewhere
\PassOptionsToPackage{unicode}{hyperref}
\PassOptionsToPackage{hyphens}{url}
\documentclass[a4paper]{extarticle}
\usepackage{fix-cm}
\usepackage[fontsize=13pt]{scrextend}
\usepackage{xcolor}
\usepackage[margin=2.5cm,top=2.5cm,bottom=2.5cm,left=2.5cm,right=2.5cm,headheight=1.25cm,headsep=0.5cm,footskip=1.25cm]{geometry}
\usepackage{amsmath,amssymb}
\setcounter{secnumdepth}{-\maxdimen} % remove section numbering
\usepackage{iftex}
\ifPDFTeX
  \usepackage[T5,T1]{fontenc}
  \usepackage[utf8]{inputenc}
  \usepackage{textcomp} % provide euro and other symbols
  \IfFileExists{newtxtext.sty}{\usepackage{newtxtext,newtxmath}}{\usepackage{lmodern}}
\else % if luatex or xetex
  \usepackage{fontspec}
  \IfFontExistsTF{Times New Roman}{\setmainfont{Times New Roman}}{\setmainfont{TeX Gyre Termes}}
  \setmonofont{Consolas}[Scale=MatchLowercase]
\fi
% Use upquote if available, for straight quotes in verbatim environments
\IfFileExists{upquote.sty}{\usepackage{upquote}}{}
\IfFileExists{microtype.sty}{% use microtype if available
  \usepackage[]{microtype}
  \UseMicrotypeSet[protrusion]{basicmath} % disable protrusion for tt fonts
}{}
\makeatletter
\@ifundefined{KOMAClassName}{% if non-KOMA class
  \IfFileExists{parskip.sty}{%
    \usepackage{parskip}
  }{% else
    \setlength{\parindent}{0pt}
    \setlength{\parskip}{6pt plus 2pt minus 1pt}}
}{% if KOMA class
  \KOMAoptions{parskip=half}}
\makeatother
\usepackage{color}
\usepackage{fancyvrb}
\newcommand{\VerbBar}{|}
\newcommand{\VERB}{\Verb[commandchars=\\\{\}]}
\DefineVerbatimEnvironment{Highlighting}{Verbatim}{commandchars=\\\{\}}
\RecustomVerbatimEnvironment{verbatim}{Verbatim}{fontsize=\small}
% Add ',fontsize=\small' for more characters per line
\newenvironment{Shaded}{}{}
\newcommand{\AlertTok}[1]{\textcolor[rgb]{1.00,0.00,0.00}{\textbf{#1}}}
\newcommand{\AnnotationTok}[1]{\textcolor[rgb]{0.38,0.63,0.69}{\textbf{\textit{#1}}}}
\newcommand{\AttributeTok}[1]{\textcolor[rgb]{0.49,0.56,0.16}{#1}}
\newcommand{\BaseNTok}[1]{\textcolor[rgb]{0.25,0.63,0.44}{#1}}
\newcommand{\BuiltInTok}[1]{\textcolor[rgb]{0.00,0.50,0.00}{#1}}
\newcommand{\CharTok}[1]{\textcolor[rgb]{0.25,0.44,0.63}{#1}}
\newcommand{\CommentTok}[1]{\textcolor[rgb]{0.38,0.63,0.69}{\textit{#1}}}
\newcommand{\CommentVarTok}[1]{\textcolor[rgb]{0.38,0.63,0.69}{\textbf{\textit{#1}}}}
\newcommand{\ConstantTok}[1]{\textcolor[rgb]{0.53,0.00,0.00}{#1}}
\newcommand{\ControlFlowTok}[1]{\textcolor[rgb]{0.00,0.44,0.13}{\textbf{#1}}}
\newcommand{\DataTypeTok}[1]{\textcolor[rgb]{0.56,0.13,0.00}{#1}}
\newcommand{\DecValTok}[1]{\textcolor[rgb]{0.25,0.63,0.44}{#1}}
\newcommand{\DocumentationTok}[1]{\textcolor[rgb]{0.73,0.13,0.13}{\textit{#1}}}
\newcommand{\ErrorTok}[1]{\textcolor[rgb]{1.00,0.00,0.00}{\textbf{#1}}}
\newcommand{\ExtensionTok}[1]{#1}
\newcommand{\FloatTok}[1]{\textcolor[rgb]{0.25,0.63,0.44}{#1}}
\newcommand{\FunctionTok}[1]{\textcolor[rgb]{0.02,0.16,0.49}{#1}}
\newcommand{\ImportTok}[1]{\textcolor[rgb]{0.00,0.50,0.00}{\textbf{#1}}}
\newcommand{\InformationTok}[1]{\textcolor[rgb]{0.38,0.63,0.69}{\textbf{\textit{#1}}}}
\newcommand{\KeywordTok}[1]{\textcolor[rgb]{0.00,0.44,0.13}{\textbf{#1}}}
\newcommand{\NormalTok}[1]{#1}
\newcommand{\OperatorTok}[1]{\textcolor[rgb]{0.40,0.40,0.40}{#1}}
\newcommand{\OtherTok}[1]{\textcolor[rgb]{0.00,0.44,0.13}{#1}}
\newcommand{\PreprocessorTok}[1]{\textcolor[rgb]{0.74,0.48,0.00}{#1}}
\newcommand{\RegionMarkerTok}[1]{#1}
\newcommand{\SpecialCharTok}[1]{\textcolor[rgb]{0.25,0.44,0.63}{#1}}
\newcommand{\SpecialStringTok}[1]{\textcolor[rgb]{0.73,0.40,0.53}{#1}}
\newcommand{\StringTok}[1]{\textcolor[rgb]{0.25,0.44,0.63}{#1}}
\newcommand{\VariableTok}[1]{\textcolor[rgb]{0.10,0.09,0.49}{#1}}
\newcommand{\VerbatimStringTok}[1]{\textcolor[rgb]{0.25,0.44,0.63}{#1}}
\newcommand{\WarningTok}[1]{\textcolor[rgb]{0.38,0.63,0.69}{\textbf{\textit{#1}}}}
\usepackage{longtable,booktabs,array}
\usepackage{caption}
\newcounter{none} % for unnumbered tables
\newcolumntype{L}[1]{>{\raggedright\arraybackslash\hyphenpenalty=10000\exhyphenpenalty=50}p{#1}}
\newcolumntype{C}[1]{>{\centering\arraybackslash}p{#1}}
\usepackage{calc} % for calculating minipage widths
% Correct order of tables after \paragraph or \subparagraph
\usepackage{etoolbox}
\makeatletter
\patchcmd\longtable{\par}{\if@noskipsec\mbox{}\fi\par}{}{}
\makeatother
% Allow footnotes in longtable head/foot
\IfFileExists{footnotehyper.sty}{\usepackage{footnotehyper}}{\usepackage{footnote}}
\makesavenoteenv{longtable}
\providecommand{\tabularnewline}{\\}
\newlength{\thesisdefaulttabcolsep}
\setlength{\thesisdefaulttabcolsep}{2pt}
\usepackage{graphicx}
\usepackage{tikz}
\usetikzlibrary{calc}
\makeatletter
\newsavebox\pandoc@box
\newcommand*\pandocbounded[1]{% scales image to fit in text height/width
  \sbox\pandoc@box{#1}%
  \Gscale@div\@tempa{\textheight}{\dimexpr\ht\pandoc@box+\dp\pandoc@box\relax}%
  \Gscale@div\@tempb{\linewidth}{\wd\pandoc@box}%
  \ifdim\@tempb\p@<\@tempa\p@\let\@tempa\@tempb\fi% select the smaller of both
  \ifdim\@tempa\p@<\p@\makebox[\linewidth][c]{\scalebox{\@tempa}{\usebox\pandoc@box}}%
  \else\makebox[\linewidth][c]{\usebox\pandoc@box}%
  \fi%
}
% Set default figure placement to htbp
\def\fps@figure{htbp}
\makeatother
\usepackage{xstring}
\newcommand{\safeincludesvg}[1]{%
  \StrSubstitute{#1}{.svg}{.png}[\safeincludepngpath]%
  \StrSubstitute{#1}{./assets/figures/}{./svg-inkscape/}[\safeincludepdfbasepath]%
  \StrSubstitute{\safeincludepdfbasepath}{.svg}{_svg-raw.pdf}[\safeincludepdfpath]%
  \IfFileExists{\safeincludepngpath}{%
    \includegraphics[width=0.96\linewidth,height=0.78\textheight,keepaspectratio]{\safeincludepngpath}%
  }{%
    \IfFileExists{\safeincludepdfpath}{%
      \includegraphics[width=0.96\linewidth,height=0.78\textheight,keepaspectratio]{\safeincludepdfpath}%
    }{%
      \fbox{\parbox{0.9\linewidth}{\centering Missing figure assets for:\\ \texttt{\detokenize{#1}}}}%
    }%
  }%
}
\newcommand{\safeincludesvgsmall}[1]{%
  \StrSubstitute{#1}{.svg}{.png}[\safeincludepngpath]%
  \StrSubstitute{#1}{./assets/figures/}{./svg-inkscape/}[\safeincludepdfbasepath]%
  \StrSubstitute{\safeincludepdfbasepath}{.svg}{_svg-raw.pdf}[\safeincludepdfpath]%
  \IfFileExists{\safeincludepngpath}{%
    \includegraphics[width=0.76\linewidth,height=0.56\textheight,keepaspectratio]{\safeincludepngpath}%
  }{%
    \IfFileExists{\safeincludepdfpath}{%
      \includegraphics[width=0.76\linewidth,height=0.56\textheight,keepaspectratio]{\safeincludepdfpath}%
    }{%
      \fbox{\parbox{0.9\linewidth}{\centering Missing figure assets for:\\ \texttt{\detokenize{#1}}}}%
    }%
  }%
}
\newcommand{\safeincludesvgcompact}[1]{%
  \StrSubstitute{#1}{.svg}{.png}[\safeincludepngpath]%
  \StrSubstitute{#1}{./assets/figures/}{./svg-inkscape/}[\safeincludepdfbasepath]%
  \StrSubstitute{\safeincludepdfbasepath}{.svg}{_svg-raw.pdf}[\safeincludepdfpath]%
  \IfFileExists{\safeincludepngpath}{%
    \includegraphics[width=0.68\linewidth,height=0.46\textheight,keepaspectratio]{\safeincludepngpath}%
  }{%
    \IfFileExists{\safeincludepdfpath}{%
      \includegraphics[width=0.68\linewidth,height=0.46\textheight,keepaspectratio]{\safeincludepdfpath}%
    }{%
      \fbox{\parbox{0.9\linewidth}{\centering Missing figure assets for:\\ \texttt{\detokenize{#1}}}}%
    }%
  }%
}
\newcommand{\safeincludesvgnarrow}[1]{%
  \StrSubstitute{#1}{.svg}{.png}[\safeincludepngpath]%
  \StrSubstitute{#1}{./assets/figures/}{./svg-inkscape/}[\safeincludepdfbasepath]%
  \StrSubstitute{\safeincludepdfbasepath}{.svg}{_svg-raw.pdf}[\safeincludepdfpath]%
  \IfFileExists{\safeincludepngpath}{%
    \includegraphics[width=0.56\linewidth,height=0.34\textheight,keepaspectratio]{\safeincludepngpath}%
  }{%
    \IfFileExists{\safeincludepdfpath}{%
      \includegraphics[width=0.56\linewidth,height=0.34\textheight,keepaspectratio]{\safeincludepdfpath}%
    }{%
      \fbox{\parbox{0.9\linewidth}{\centering Missing figure assets for:\\ \texttt{\detokenize{#1}}}}%
    }%
  }%
}
\newcommand{\safeincludesvgchart}[1]{%
  \StrSubstitute{#1}{.svg}{.png}[\safeincludepngpath]%
  \StrSubstitute{#1}{./assets/figures/}{./svg-inkscape/}[\safeincludepdfbasepath]%
  \StrSubstitute{\safeincludepdfbasepath}{.svg}{_svg-raw.pdf}[\safeincludepdfpath]%
  \IfFileExists{\safeincludepngpath}{%
    \includegraphics[width=0.72\linewidth,height=0.28\textheight,keepaspectratio]{\safeincludepngpath}%
  }{%
    \IfFileExists{\safeincludepdfpath}{%
      \includegraphics[width=0.72\linewidth,height=0.28\textheight,keepaspectratio]{\safeincludepdfpath}%
    }{%
      \fbox{\parbox{0.9\linewidth}{\centering Missing figure assets for:\\ \texttt{\detokenize{#1}}}}%
    }%
  }%
}
\setlength{\emergencystretch}{10em} % reduce overfull lines
\setlength{\tabcolsep}{\thesisdefaulttabcolsep}
\setlength{\arrayrulewidth}{0.5pt}
\captionsetup[table]{justification=centering,singlelinecheck=false,labelfont=it,textfont=it}
\captionsetup[longtable]{justification=centering,singlelinecheck=false,labelfont=it,textfont=it}
\setlength{\LTcapwidth}{\linewidth}
\setlength{\LTleft}{0pt}
\setlength{\LTright}{0pt}
% Pandoc-generated longtables in this thesis render correctly but emit repeated
% box warnings during the output routine; raise fuzz thresholds to silence this
% known-noise without changing the visible layout.
\setlength{\hfuzz}{460pt}
\setlength{\vfuzz}{360pt}
\hbadness=10000
\setcounter{topnumber}{2}
\setcounter{bottomnumber}{2}
\setcounter{totalnumber}{4}
\renewcommand{\topfraction}{0.85}
\renewcommand{\bottomfraction}{0.75}
\renewcommand{\textfraction}{0.1}
\renewcommand{\floatpagefraction}{0.75}
\usepackage{setspace}
\setstretch{1.15}
\setlength{\parindent}{0pt}
\setlength{\parskip}{6pt}
\setcounter{tocdepth}{3}
\renewcommand{\arraystretch}{1.15}
\usepackage{titlesec}
\titleformat{\section}{\bfseries\Large\centering}{}{0pt}{}
\titlespacing*{\section}{0pt}{12pt}{6pt}
\titleformat{\subsection}{\bfseries\normalsize}{}{0pt}{}
\titlespacing*{\subsection}{0pt}{12pt}{0pt}
\newcommand{\thesisfrontmatterentry}[1]{%
  \phantomsection
  \addcontentsline{toc}{section}{#1}%
}
\newcommand{\majorheading}[1]{%
  \clearpage
  \phantomsection
  \section*{#1}%
  \addcontentsline{toc}{section}{#1}%
}
\newcommand{\chapterheading}[1]{%
  \clearpage
  \phantomsection
  \section*{#1}%
  \addcontentsline{toc}{section}{#1}%
  \setcounter{subsection}{0}%
}
\newcommand{\appendixsection}[1]{\clearpage\section{#1}}
\newcommand{\frontmatterstart}{\clearpage\hypersetup{pageanchor=true}\pagenumbering{roman}\setcounter{page}{1}}
\newcommand{\mainmatterstart}{\clearpage\pagenumbering{arabic}\setcounter{page}{1}}
\newcommand{\checkbox}{\(\square\)}
\newcommand{\thesisspacer}{\hspace{2em}}
\newcommand{\figref}[1]{Hình~\ref{#1}}
\newcommand{\tabref}[1]{Bảng~\ref{#1}}
\definecolor{thesisphenikaablue}{RGB}{42, 63, 130}
\newcommand{\thesislogopath}{./assets/brand/phenikaa-university-logo.jpeg}
\newcommand{\thesisfrontseparator}{\begin{center}\rule{0.5\linewidth}{0.5pt}\end{center}}
\newcommand{\thesisdrawtemplateframe}{%
  \begin{tikzpicture}[remember picture,overlay]
    \draw[line width=1.4pt] ($(current page.north west)+(1.00cm,-1.00cm)$) rectangle ($(current page.south east)+(-1.00cm,1.00cm)$);
    \draw[line width=0.5pt] ($(current page.north west)+(1.14cm,-1.14cm)$) rectangle ($(current page.south east)+(-1.14cm,1.14cm)$);
  \end{tikzpicture}%
}
\newcommand{\thesisdrawsidebarframe}{%
  \begin{tikzpicture}[remember picture,overlay]
    \draw[line width=0.5pt] ($(current page.north west)+(1.35cm,-1.35cm)$) rectangle ($(current page.south east)+(-1.35cm,1.35cm)$);
    \draw[line width=0.5pt] ($(current page.north west)+(3.65cm,-1.35cm)$) -- ($(current page.south west)+(3.65cm,1.35cm)$);
  \end{tikzpicture}%
}
\newcommand{\thesiscoverstart}{%
  \clearpage
  \newgeometry{top=1.5cm,bottom=1.5cm,left=1.6cm,right=1.6cm}
  \thispagestyle{empty}
}
\newcommand{\thesiscoverend}{%
  \clearpage
  \restoregeometry
}
\newcommand{\thesisofficialheader}{%
  \noindent
  \begin{minipage}[t]{0.48\linewidth}
    \centering
    \textbf{BỘ GIÁO DỤC VÀ ĐÀO TẠO}\\
    \textbf{ĐẠI HỌC PHENIKAA}\\
    \rule{0.72\linewidth}{0.6pt}
  \end{minipage}\hfill
  \begin{minipage}[t]{0.48\linewidth}
    \centering
    \textbf{CỘNG HÒA XÃ HỘI CHỦ NGHĨA VIỆT NAM}\\
    \textbf{Độc lập - Tự do - Hạnh phúc}\\
    \rule{0.72\linewidth}{0.6pt}
  \end{minipage}
  \par\vspace{0.6cm}
}
\newcommand{\thesissignaturepair}[4]{%
  \vspace{0.8cm}
  \noindent
  \begin{minipage}[t]{0.48\linewidth}
    \centering
    \textbf{#1}\\
    \emph{(Ký, ghi rõ họ tên)}\\[2.8cm]
    #2
  \end{minipage}\hfill
  \begin{minipage}[t]{0.48\linewidth}
    \centering
    \textbf{#3}\\
    \emph{(Ký, ghi rõ họ tên)}\\[2.8cm]
    #4
  \end{minipage}
}
\renewcommand{\figurename}{Hình}
\renewcommand{\tablename}{Bảng}
\renewcommand{\listtablename}{DANH MỤC BẢNG - LIST OF TABLES}
\renewcommand{\listfigurename}{DANH MỤC HÌNH ẢNH VÀ ĐỒ THỊ - LIST OF PICTURES AND GRAPHS}
\makeatletter
\newcommand{\thesissetfigureid}[2]{%
  \renewcommand{\thefigure}{#1}%
  \renewcommand{\theHfigure}{#2}%
}
\newcommand{\thesissettableid}[2]{%
  \renewcommand{\thetable}{#1}%
  \renewcommand{\theHtable}{#2}%
}
\newcommand{\thesistableofcontentsbody}{\@starttoc{toc}}
\newcommand{\thesislistoftablesbody}{\@starttoc{lot}}
\newcommand{\thesislistoffiguresbody}{\@starttoc{lof}}
\renewcommand*\l@table{\@dottedtocline{1}{0em}{5.8em}}
\renewcommand*\l@figure{\@dottedtocline{1}{0em}{5.8em}}
\makeatother
\newcounter{ieeeref}
\newcommand{\ieeebibitem}[2]{%
  \refstepcounter{ieeeref}\noindent[\theieeeref]\label{#1} #2\par
}
\newcommand{\ieeecite}[1]{[\ref{#1}]}
\providecommand{\textcite}[1]{\ieeecite{#1}}
\providecommand{\parencite}[1]{\ieeecite{#1}}
\providecommand{\tightlist}{%
  \setlength{\itemsep}{0pt}\setlength{\parskip}{0pt}}
\usepackage{bookmark}
\IfFileExists{xurl.sty}{\usepackage{xurl}}{} % add URL line breaks if available
\hyphenation{Mos-quit-to Post-gre-SQL Vic-to-ria-Met-rics Vic-to-ria-Logs Au-then-ti-ca-tion Ku-ber-ne-tes Os-cil-lo-scope Dash-board Web-Socket}
\urlstyle{same}
\hypersetup{
  hidelinks,
  linktoc=all,
  bookmarksnumbered=true,
  pdfcreator={LaTeX via pandoc}}

\author{}
\date{}

\begin{document}
\hypersetup{pageanchor=false}
\pagenumbering{gobble}

\thesiscoverstart
\thesisdrawsidebarframe
\noindent
\begin{minipage}[t][0.92\textheight][t]{0.16\textwidth}
  \centering
  \vspace*{1.0cm}
  \rotatebox{90}{\textbf{\Large LÊ TRỌNG AN}}\\[1.5cm]
  \rotatebox{90}{\textbf{\large KỸ THUẬT CƠ ĐIỆN TỬ}}
\end{minipage}\hfill
\begin{minipage}[t][0.92\textheight][t]{0.77\textwidth}
  \begin{center}
    \textbf{\Large BỘ GIÁO DỤC VÀ ĐÀO TẠO}\\[0.2cm]
    \textbf{\Large ĐẠI HỌC PHENIKAA}\\[-0.05cm]
    \rule{0.34\linewidth}{0.6pt}

    \vspace{1.7cm}
    \includegraphics[width=0.36\linewidth]{\thesislogopath}

    \vspace{1.5cm}
    {\fontsize{28}{30}\selectfont\textbf{ĐỒ ÁN TỐT NGHIỆP}}\\[1.0cm]
    {\fontsize{17}{21}\selectfont\textbf{THIẾT KẾ HỆ THỐNG IOT CHO ỨNG DỤNG QUẢN LÝ PHƯƠNG TIỆN}}\\
    {\fontsize{17}{21}\selectfont\textbf{GIAO THÔNG TRONG LĨNH VỰC CHO THUÊ XE TỰ LÁI}}


    \vfill
    \begin{minipage}{0.90\linewidth}
      \raggedright
      \textbf{Sinh viên:} Lê Trọng An\\[0.15cm]
      \textbf{Mã số sinh viên:} 21010389 \hfill \textbf{Khóa:} K15\\[0.15cm]
      \textbf{Ngành:} Kỹ thuật Cơ Điện tử \hfill \textbf{Hệ:} Đại học chính quy\\[0.15cm]
      \textbf{Giảng viên hướng dẫn:} TS. Nguyễn Đức Nam
    \end{minipage}

    \vfill
    \textbf{\Large Hà Nội - 2026}
  \end{center}
\end{minipage}
\thesiscoverend

\thesiscoverstart
\thesisdrawtemplateframe
\begin{center}
  \textbf{\Large BỘ GIÁO DỤC VÀ ĐÀO TẠO}\\[0.15cm]
  \textbf{\Large ĐẠI HỌC PHENIKAA}\\[-0.05cm]
  \rule{0.28\linewidth}{0.6pt}

  \vspace{1.6cm}
  \includegraphics[width=0.30\linewidth]{\thesislogopath}

  \vspace{1.4cm}
  {\fontsize{27}{30}\selectfont\textbf{ĐỒ ÁN TỐT NGHIỆP}}\\[0.9cm]
  {\fontsize{18}{22}\selectfont\textbf{THIẾT KẾ HỆ THỐNG IOT CHO ỨNG DỤNG QUẢN LÝ}}\\
  {\fontsize{18}{22}\selectfont\textbf{PHƯƠNG TIỆN GIAO THÔNG TRONG LĨNH VỰC}}\\
  {\fontsize{18}{22}\selectfont\textbf{CHO THUÊ XE TỰ LÁI}}

  \vfill
  \begin{minipage}{0.78\linewidth}
    \raggedright
    \textbf{Sinh viên:} Lê Trọng An\\[0.18cm]
    \textbf{Mã số sinh viên:} 21010389 \hfill \textbf{Khóa:} K15\\[0.18cm]
    \textbf{Ngành:} Kỹ thuật Cơ Điện tử \hfill \textbf{Hệ:} Đại học chính quy\\[0.18cm]
    \textbf{Giảng viên hướng dẫn:} TS. Nguyễn Đức Nam
  \end{minipage}

  \vspace{1.3cm}
  \textbf{\Large Hà Nội – 2026}
\end{center}
\thesiscoverend

\thesiscoverstart
\thesisdrawtemplateframe
\begin{center}
  {\color{thesisphenikaablue}\textbf{\Large PHENIKAA UNIVERSITY}}\\[0.18cm]
  {\color{thesisphenikaablue}\textbf{\Large FACULTY OF MECHANICAL ENGINEERING AND MECHATRONICS}}\\[0.25cm]
  \includegraphics[width=0.28\linewidth]{\thesislogopath}

  \vspace{1.0cm}
  {\fontsize{22}{26}\selectfont\textbf{MEM710071 – CAPSTONE DESIGN 2}}\\[0.55cm]
  {\fontsize{18}{22}\selectfont\textbf{DESIGN OF AN IOT SYSTEM FOR VEHICLE MANAGEMENT}}\\
  {\fontsize{18}{22}\selectfont\textbf{IN SELF-DRIVE CAR RENTAL SERVICES}}\\[0.35cm]
  {\fontsize{15}{18}\selectfont\textbf{Semester: I Academic year: 2025--2026}}

  \vfill
  \begin{minipage}{0.76\linewidth}
    \raggedright
    \textbf{Project Team:}
    \begin{center}
      \begin{tabular}{|p{0.42\linewidth}|p{0.20\linewidth}|p{0.12\linewidth}|}
       \\\hline
        \centering\textbf{Student name} & \centering\textbf{Student ID} & \centering\textbf{Cohort}\tabularnewline
       \\\hline
        Le Trong An & \centering 21010389 & \centering K15\tabularnewline
       \\\hline
      \end{tabular}
    \end{center}
    \vspace{0.2cm}
    \textbf{Course Instructor:} Dr. Nguyen Duc Nam
  \end{minipage}

  \vfill
  \begin{minipage}{0.88\linewidth}
    \centering
    This report is submitted to the Faculty of Mechanical Engineering and Mechatronics,\\
    School of Engineering, Phenikaa University in partial fulfillment of the requirements\\
    of the Graduation Thesis course MEM710071.
  \end{minipage}

  \vspace{0.8cm}
  \textbf{\Large HANOI – 2026}
\end{center}
\thesiscoverend

\thesiscoverstart
\thesisdrawtemplateframe
\begin{center}
  {\color{thesisphenikaablue}\textbf{\Large ĐẠI HỌC PHENIKAA}}\\[0.15cm]
  {\color{thesisphenikaablue}\textbf{\Large KHOA CƠ KHÍ – CƠ ĐIỆN TỬ}}\\[0.25cm]
  \includegraphics[width=0.27\linewidth]{\thesislogopath}

  \vspace{0.8cm}
  {\fontsize{21}{25}\selectfont\textbf{MEM710071}}\\[0.15cm]
  {\fontsize{26}{30}\selectfont\textbf{ĐỒ ÁN TỐT NGHIỆP}}\\[0.45cm]
  {\fontsize{18}{22}\selectfont\textbf{THIẾT KẾ HỆ THỐNG IOT CHO ỨNG DỤNG QUẢN LÝ}}\\
  {\fontsize{18}{22}\selectfont\textbf{PHƯƠNG TIỆN GIAO THÔNG TRONG LĨNH VỰC CHO THUÊ XE TỰ LÁI}}\\[0.35cm]
  {\fontsize{15}{18}\selectfont\textbf{Học kỳ I Năm học 2025--2026}}

  \vfill
  \begin{center}
    \begin{tabular}{|p{0.34\linewidth}|p{0.24\linewidth}|p{0.10\linewidth}|}
     \\\hline
      \centering\textbf{Họ và tên sinh viên} & \centering\textbf{Mã số sinh viên} & \centering\textbf{Lớp}\tabularnewline
     \\\hline
      Lê Trọng An & \centering 21010389 & \centering K15-KTCĐT2\tabularnewline
     \\\hline
    \end{tabular}
  \end{center}

  \vspace{0.45cm}
  \textbf{Giảng viên hướng dẫn:} TS. Nguyễn Đức Nam

  \vfill
  \textbf{\Large HÀ NỘI, 2026}
\end{center}
\thesiscoverend

\frontmatterstart

\thispagestyle{empty}
\thesisofficialheader
\begin{center}
  \textbf{\Large NHẬN XÉT ĐỒ ÁN TỐT NGHIỆP CỦA GIẢNG VIÊN HƯỚNG DẪN}
\end{center}
\thesisfrontmatterentry{NHẬN XÉT CỦA GIẢNG VIÊN HƯỚNG DẪN}

\textbf{Giảng viên hướng dẫn:} TS. Nguyễn Đức Nam \hfill \textbf{Khoa:} Cơ khí -- Cơ Điện tử.\\
\textbf{Tên đề tài:} Thiết kế hệ thống IoT cho ứng dụng quản lý phương tiện giao thông trong lĩnh vực cho thuê xe tự lái.\\
\textbf{Sinh viên thực hiện:} Lê Trọng An. \hfill \textbf{Lớp:} K15-KTCĐT2.

\begin{center}
  \textbf{NỘI DUNG NHẬN XÉT}
\end{center}
\textbf{I. Nhận xét Đồ án tốt nghiệp:}\\
- Nhận xét về hình thức: \dotfill\\[0.55cm]
- Tính cấp thiết của đề tài: \dotfill\\[0.55cm]
- Mục tiêu của đề tài: \dotfill\\[0.55cm]
- Nội dung của đề tài: \dotfill\\[0.55cm]
- Tài liệu tham khảo: \dotfill\\[0.55cm]
- Phương pháp nghiên cứu: \dotfill\\[0.55cm]
- Tính sáng tạo và ứng dụng: \dotfill\\[0.55cm]
\textbf{II. Nhận xét về tinh thần và thái độ làm việc của sinh viên:} \dotfill\\[0.55cm]
\textbf{III. Kết quả đạt được:} \dotfill\\[0.55cm]
\textbf{IV. Kết luận:} Đồng ý cho bảo vệ: \checkbox \thesisspacer Không đồng ý cho bảo vệ: \checkbox

\begin{flushright}
  Hà Nội, ngày \ldots, tháng \ldots năm 20\ldots\\
  \textbf{GIẢNG VIÊN HƯỚNG DẪN}\\
  \emph{(Ký, ghi rõ họ tên)}\\[2.6cm]
  TS. Nguyễn Đức Nam
\end{flushright}
\clearpage

\thispagestyle{empty}
\thesisofficialheader
\begin{center}
  \textbf{\Large NHẬN XÉT ĐỒ ÁN TỐT NGHIỆP CỦA GIẢNG VIÊN PHẢN BIỆN}
\end{center}
\thesisfrontmatterentry{NHẬN XÉT CỦA GIẢNG VIÊN PHẢN BIỆN}

\textbf{Giảng viên phản biện:} Theo quyết định phân công phản biện của Khoa \\
\textbf{Tên đề tài:} Thiết kế hệ thống IoT cho ứng dụng quản lý phương tiện giao thông trong lĩnh vực cho thuê xe tự lái.\\
\textbf{Sinh viên thực hiện:} Lê Trọng An \hfill \textbf{Lớp:} K15-KTCĐT2\\
\textbf{Giảng viên hướng dẫn:} TS. Nguyễn Đức Nam

\begin{center}
  \textbf{NỘI DUNG NHẬN XÉT}
\end{center}
\textbf{I. Nhận xét Đồ án tốt nghiệp:}\\
- Bố cục, hình thức trình bày: \dotfill\\[0.55cm]
- Đảm bảo tính cấp thiết, hiện đại, không trùng lặp: \dotfill\\[0.55cm]
- Ná»™i dung: \dotfill\\[0.55cm]
- Mức độ thực hiện: \dotfill\\[0.55cm]
\textbf{II. Kết quả đạt được:} \dotfill\\[0.55cm]
\textbf{III. Ưu nhược điểm:} \dotfill\\[0.55cm]
\textbf{IV. Kết luận:} Đồng ý cho bảo vệ: \checkbox \thesisspacer Không đồng ý cho bảo vệ: \checkbox

\begin{flushright}
  Hà Nội, ngày \ldots, tháng \ldots năm 20\ldots\\
  \textbf{GIẢNG VIÊN PHẢN BIỆN}\\
  \emph{(Ký, ghi rõ họ tên)}\\[2.6cm]
  Theo quyết định phân công phản biện của Khoa
\end{flushright}
\clearpage

\thispagestyle{empty}
\thesisofficialheader
\begin{center}
  \textbf{\Large BIÊN BẢN ĐÁNH GIÁ ĐỒ ÁN TỐT NGHIỆP}
\end{center}
\thesisfrontmatterentry{BIÊN BẢN ĐÁNH GIÁ ĐỒ ÁN TỐT NGHIỆP}

\textbf{I. Thông tin chung}
\begin{enumerate}
  \item Thời gian: từ \ldots giờ \ldots đến \ldots giờ \ldots ngày \ldots/\ldots/20\ldots
  \item Địa điểm: \dotfill
  \item Họ và tên sinh viên: Lê Trọng An \hfill Mã số SV: 21010389
  \item Lớp: K15-KTCĐT2 \hfill Ngành: Kỹ thuật Cơ Điện tử \hfill Khóa: K15
  \item Tên đề tài: Thiết kế hệ thống IoT cho ứng dụng quản lý phương tiện giao thông trong lĩnh vực cho thuê xe tự lái.
  \item Giảng viên hướng dẫn: TS. Nguyễn Đức Nam
\end{enumerate}

\textbf{II. Thành phần hội đồng}
\begin{enumerate}
  \item \dotfill Chủ tịch
  \item \dotfill Thư ký
  \item \dotfill Giảng viên phản biện
  \item \dotfill Ủy viên 1
\end{enumerate}

\textbf{III. Tổng hợp câu hỏi của thành viên hội đồng:} \dotfill\\[0.6cm]
\textbf{IV. Tổng hợp nội dung trả lời của sinh viên:} \dotfill\\[0.6cm]
\textbf{V. Nội dung đánh giá của Hội đồng}
\begin{enumerate}
  \item Ý nghĩa của đồ án: \dotfill
  \item Về nội dung, kết cấu của đồ án: \dotfill
  \item Phương pháp nghiên cứu: \dotfill
  \item Các kết quả nghiên cứu đạt được: \dotfill
\end{enumerate}
\textbf{VI. Điểm kết luận:} \dotfill\ (Bằng chữ: \dotfill)

\thesissignaturepair{CHỦ TỊCH HỘI ĐỒNG}{}{THƯ KÝ}{}
\clearpage

\thispagestyle{empty}
\begin{center}
  \textbf{\Large GIẢI TRÌNH CÁC CHỈNH SỬA}\\[0.2cm]
  \textit{(Nếu có)}
\end{center}
\thesisfrontmatterentry{GIẢI TRÌNH CÁC CHỈNH SỬA}

\vspace{0.5cm}
\begin{center}
  Chưa phát sinh nội dung cần giải trình.
\end{center}
\clearpage

\majorheading{LỜI CAM ĐOAN}

Tên tôi là: Lê Trọng An.\\
Mã sinh viên: 21010389 \hfill Lớp: K15-KTCĐT2\\
Ngành: Kỹ thuật Cơ Điện tử.

Tôi đã thực hiện đồ án tốt nghiệp với đề tài: Thiết kế hệ thống IoT cho ứng dụng quản lý phương tiện giao thông trong lĩnh vực cho thuê xe tự lái.

Tôi xin cam đoan đây là đề tài nghiên cứu của riêng tôi và được sự hướng dẫn của TS. Nguyễn Đức Nam.

Các nội dung nghiên cứu, kết quả trong đề tài này là trung thực và chưa được công bố dưới bất kỳ hình thức nào. Nếu phát hiện có bất kỳ hình thức gian lận nào, tôi xin hoàn toàn chịu trách nhiệm trước pháp luật.

\begin{flushright}
  Hà Nội, ngày \ldots, tháng \ldots năm 20\ldots
\end{flushright}

\thesissignaturepair{GIẢNG VIÊN HƯỚNG DẪN}{TS. Nguyễn Đức Nam}{SINH VIÊN}{Lê Trọng An}

\majorheading{TÓM TẮT ĐỒ ÁN TỐT NGHIỆP - ABSTRACT}

Sự phát triển nhanh của dịch vụ cho thuê xe tự lái tại Việt Nam khiến
nhu cầu giám sát phương tiện không còn dừng ở việc biết xe đang ở đâu,
mà còn ở việc hiểu xe đang vận hành ra sao và có dấu hiệu bất thường hay
không. Từ nhu cầu đó, đồ án này trình bày quá trình thiết kế và triển
khai một hệ thống giám sát phương tiện theo hướng trọn gói, bao gồm cả
thiết bị gắn trên xe, hạ tầng máy chủ tiếp nhận dữ liệu và giao diện để
người quản lý theo dõi, nhằm hỗ trợ doanh nghiệp nắm vị trí xe, trạng
thái kỹ thuật và phản ứng sớm với các rủi ro vận hành.

Ở tầng thiết bị, hệ thống được xây dựng trên bo mạch PCB tự thiết kế.
Trong đó, ESP32-S3 giữ vai trò bộ điều khiển trung tâm; SIM7600CE-T đảm
nhiệm truyền dữ liệu 4G và định vị vệ tinh GNSS để xác định vị trí xe;
LIS3DSH theo dõi rung động để nhận biết chuyển động khi xe đỗ; còn
adapter vgate iCar Pro cho phép đọc dữ liệu OBD2 từ xe qua BLE. Kiến
trúc nguồn được tổ chức theo nhiều
nhánh để thiết bị vẫn duy trì hoạt động khi xe tắt máy, đồng thời hạn
chế ảnh hưởng tới ắc quy chính.

Ở tầng phần mềm, dữ liệu từ thiết bị được gửi bằng MQTT 3.1.1 tới EMQX.
MQTT đóng vai trò lớp nhắn tin nhẹ phù hợp với đường truyền 4G, còn EMQX
là broker trung tâm tiếp nhận và phân phối bản tin. Từ đây, MQTT Bridge
subscribe bốn nhóm topic
\texttt{v1/\{device\_id\}/\{rawdata\textbar status\textbar events\textbar firmware\}},
chuẩn hóa payload rồi phân luồng dữ liệu tới PostgreSQL,
VictoriaMetrics và VictoriaLogs. Song song với luồng ingest, backend còn
lắng nghe nhánh \texttt{internal/events/\#} để đẩy cập nhật thời gian
thực ra giao diện qua Socket.IO.

Bảng điều khiển web được phát triển bằng Next.js 15 và React 19. Điều
quan trọng ở đây không phải tên framework, mà là giao diện có thể hiển
thị bản đồ, biểu đồ và trạng thái phương tiện gần như theo thời gian
thực nhờ cơ chế đẩy dữ liệu mới từ máy chủ lên trình duyệt qua
Socket.IO. Nhờ đó, người vận hành không chỉ quan sát vị
trí xe mà còn có thể theo dõi lịch sử hành trình, cảnh báo và trạng thái
thiết bị trên cùng một giao diện thống nhất.

Kết quả đạt được là một hệ thống IoT giám sát phương tiện hoàn chỉnh từ
phần cứng đến phần mềm, có khả năng theo dõi vị trí thời gian thực, đọc
dữ liệu chẩn đoán OBD2, thiết lập vùng cho phép, phát sinh cảnh báo tự
động, cập nhật firmware từ xa và hỗ trợ quản lý đội xe. Quan trọng hơn,
hệ thống cho thấy một hướng triển khai phù hợp với bối cảnh doanh nghiệp
cho thuê xe tự lái trong nước: chi phí hợp lý, có thể tùy biến và đủ dư
địa để mở rộng trong các giai đoạn sau.

\textbf{Từ khóa:} IoT, giám sát phương tiện, GPS, OBD2, MQTT, ESP32-S3,
thời gian thực, vùng cho phép

\textbf{Abstract (English)}

As the self-drive car rental market in Vietnam grows, the need for
remote vehicle monitoring and real-time fleet management becomes
increasingly practical for operators. This thesis presents the design
and development of an end-to-end IoT vehicle tracking system, covering
the onboard device, cloud services, and operational dashboard for
real-time tracking, diagnostic data collection, and alerting.

On the hardware side, the tracker is implemented as a custom integrated
PCB rather than an assembly of off-the-shelf development boards. The
main board integrates ESP32-S3 (MCU), SIMCom SIM7600CE-T (LTE+GNSS),
LIS3DSH (IMU), DS3231M (RTC), W25Q128 (SPI Flash), and power blocks based
on MP2482, TPS54231, AP2112, SX1308, and TP5100. The modem is configured
to Auto mode (\texttt{AT+CNMP=2}) so it can switch between LTE/UMTS/GSM
automatically, uses the default APN \texttt{internet}, and delivers GNSS
data via \texttt{AT+CGNSINF} (with optional NMEA streaming through
\texttt{AT+CGNSTST}) on the same UART1 channel. The device reads engine
diagnostic data via OBD2 through a vgate iCar Pro adapter over Bluetooth
Low Energy (BLE), where this adapter is an external peripheral outside
the main PCB. The power architecture uses a 5V main rail, an
approximately 4V modem rail, a 3.3V logic rail, and a 18650 1S backup
cell to maintain operation when the vehicle engine is off.

On the software side, device data is transmitted through MQTT 3.1.1 via
EMQX with per-device ACL authorization. The MQTT Bridge subscribes to
the device topics
\texttt{v1/\{deviceId\}/\{rawdata\textbar status\textbar events\textbar firmware\}},
validates the payloads, and routes them to PostgreSQL, VictoriaMetrics,
and VictoriaLogs. The backend is built with Express.js and TypeScript in
a domain-oriented structure, exposes REST/OpenAPI endpoints, and listens
to internal MQTT events so Socket.IO channels can push realtime updates
and OTA status back to operators.

The web interface is developed using Next.js 15 and React 19, providing
a real-time dashboard with Leaflet maps, ECharts visualizations, and
WebSocket connectivity via Socket.IO for instant vehicle position and
status updates. The entire system is containerized and deployed using
Docker, ensuring consistency between development and production
environments.

The result is a complete end-to-end IoT vehicle tracking system, from
hardware to software, capable of real-time location tracking, OBD2
diagnostic data retrieval, allowed-zone monitoring, automated alerts,
remote firmware updates, and fleet-management support for self-drive car
rental services.

\textbf{Keywords:} IoT, vehicle tracking, GPS, OBD2, MQTT, ESP32-S3,
real-time, allowed-zone monitoring

\majorheading{LỜI CẢM ƠN - ACKNOWLEDGEMENTS}\label{lux1eddi-cux1ea3m-ux1a1n---acknowledgements}

Trước tiên, tôi xin gửi lời cảm ơn chân thành và sâu sắc nhất đến TS.
Nguyễn Đức Nam --- người đã trực tiếp hướng dẫn, định hướng và tận tình
hỗ trợ tôi trong suốt quá trình thực hiện đồ án tốt nghiệp này. Những
góp ý chuyên môn và sự động viên của thầy là nguồn động lực quan trọng
giúp tôi hoàn thành đề tài.

Tôi xin trân trọng cảm ơn quý thầy cô trong Khoa Cơ khí - Cơ Điện tử,
Trường Kỹ thuật - Đại học Phenikaa đã truyền đạt những kiến thức nền
tảng vững chắc trong suốt quá trình học tập, tạo điều kiện thuận lợi để
tôi có thể áp dụng vào thực tiễn thông qua đồ án này.

Tôi cũng xin gửi lời cảm ơn đến gia đình, bạn bè và các đồng nghiệp đã
luôn ủng hộ, chia sẻ và đồng hành cùng tôi trong suốt thời gian thực
hiện đề tài.

Mặc dù đã nỗ lực hoàn thiện đồ án trong phạm vi thời gian và kiến thức
cho phép, nội dung chắc chắn vẫn còn những thiếu sót. Tôi rất mong nhận
được những ý kiến góp ý từ quý thầy cô và hội đồng để tiếp tục hoàn
thiện đề tài.

Xin chân thành cảm ơn!

\begin{flushright}
  Hà Nội, ngày \ldots tháng \ldots năm 20\ldots\\[0.25cm]
  \textbf{SINH VIÊN THỰC HIỆN}\\[2.4cm]
  Lê Trọng An
\end{flushright}

\majorheading{MỤC LỤC - TABLE OF CONTENT}\label{mux1ee5c-lux1ee5c---table-of-contents}

\thesistableofcontentsbody

\majorheading{DANH MỤC BẢNG - LIST OF TABLES}\label{danh-mux1ee5c-bux1ea3ng---list-of-tables}

\thesislistoftablesbody

\majorheading{DANH MỤC HÌNH ẢNH VÀ ĐỒ THỊ - LIST OF PICTURES AND GRAPHS}\label{danh-mux1ee5c-huxecnh---list-of-figures}

\thesislistoffiguresbody

\majorheading{DANH MỤC TỪ VIẾT TẮT - LIST OF ABBREVIATION}\label{danh-mux1ee5c-tux1eeb-viux1ebft-tux1eaft---list-of-abbreviations}

{\def\LTcaptype{none} % do not increment counter
\setlength{\tabcolsep}{4pt}
\begin{longtable}{|L{0.160\linewidth}|L{0.410\linewidth}|L{0.360\linewidth}|}
\\\hline
\textbf{Từ viết tắt} & \textbf{Tiếng Anh} & \textbf{Tiếng Việt} \\
\\\hline
\endfirsthead
\\\hline
\textbf{Từ viết tắt} & \textbf{Tiếng Anh} & \textbf{Tiếng Việt} \\
\\\hline
\endhead
\\\hline
\endfoot
\endlastfoot
\textbf{ADC} & Analog-to-Digital Converter & Bộ chuyển đổi tương tự sang
số \\\hline
\textbf{API} & Application Programming Interface & Giao diện lập trình
ứng dụng \\\hline
\textbf{BLE} & Bluetooth Low Energy & Bluetooth năng lượng thấp \\\hline
\textbf{BOM} & Bill of Materials & Danh sách linh kiện \\\hline
\textbf{CI/CD} & Continuous Integration / Continuous Deployment & Tích
hợp liên tục / Triển khai liên tục \\\hline
\textbf{CORS} & Cross-Origin Resource Sharing & Chia sẻ tài nguyên nguồn
chéo \\\hline
\textbf{DDD} & Domain-Driven Design & Thiết kế hướng miền \\\hline
\textbf{Docker} & Docker Container Platform & Nền tảng container hóa ứng
dụng \\\hline
\textbf{EMQX} & Erlang MQTT Broker & Máy chủ MQTT dựa trên Erlang \\\hline
\textbf{ESP32} & Espressif System-on-Chip 32-bit & Vi điều khiển SoC
32-bit của Espressif \\\hline
\textbf{FreeRTOS} & Free Real-Time Operating System & Hệ điều hành thời
gian thực mã nguồn mở \\\hline
\textbf{GPIO} & General Purpose Input/Output & Cổng vào/ra đa năng \\\hline
\textbf{GNSS} & Global Navigation Satellite System & Hệ thống vệ tinh
dẫn đường toàn cầu \\\hline
\textbf{GPS} & Global Positioning System & Hệ thống định vị toàn cầu \\\hline
\textbf{HTTP} & Hypertext Transfer Protocol & Giao thức truyền siêu văn
bản \\\hline
\textbf{HTTPS} & Hypertext Transfer Protocol Secure & Giao thức truyền
siêu văn bản bảo mật \\\hline
\textbf{IMU} & Inertial Measurement Unit & Đơn vị đo lường quán tính \\\hline
\textbf{IoT} & Internet of Things & Internet vạn vật \\\hline
\textbf{JWT} & JSON Web Token & Mã thông báo web dạng JSON \\\hline
\textbf{LVD} & Low Voltage Disconnect & Ngắt điện áp thấp \\\hline
\textbf{MCU} & Microcontroller Unit & Vi điều khiển \\\hline
\textbf{MQTT} & Message Queuing Telemetry Transport & Giao thức truyền
thông tin hàng đợi \\\hline
\textbf{OBD2} & On-Board Diagnostics II & Chẩn đoán trên xe thế hệ 2 \\\hline
\textbf{ORM} & Object-Relational Mapping & Ánh xạ đối tượng - quan hệ \\\hline
\textbf{PCB} & Printed Circuit Board & Bảng mạch in \\\hline
\textbf{REST} & Representational State Transfer & Kiến trúc truyền trạng
thái đại diện \\\hline
\textbf{RTOS} & Real-Time Operating System & Hệ điều hành thời gian
thá»±c \\\hline
\textbf{SoC} & System on Chip & Hệ thống trên chip \\\hline
\textbf{SQL} & Structured Query Language & Ngôn ngữ truy vấn có cấu
trúc \\\hline
\textbf{TLS} & Transport Layer Security & Bảo mật tầng vận chuyển \\\hline
\textbf{UART} & Universal Asynchronous Receiver-Transmitter & Bộ thu
phát bất đồng bộ đa năng \\\hline
\textbf{WebSocket} & WebSocket Protocol & Giao thức kết nối hai chiều
thời gian thực \\\hline
\end{longtable}
}

\begin{center}\rule{0.5\linewidth}{0.5pt}\end{center}

\mainmatterstart
\chapterheading{CHƯƠNG 1. GIỚI THIỆU DỰ ÁN -- PROJECT OVERVIEW}

\subsection{1.1. Đặt vấn đề / Bối cảnh của dự
án}\label{11-ux111ux1eb7t-vux1ea5n-ux111ux1ec1--bux1ed1i-cux1ea3nh-cux1ee7a-dux1ef1-uxe1n}

\subsubsection{1.1.1. Bối cảnh thị trường cho thuê xe tự lái tại Việt
Nam}\label{111-bux1ed1i-cux1ea3nh-thux1ecb-trux1b0ux1eddng-cho-thuuxea-xe-tux1ef1-luxe1i-tux1ea1i-viux1ec7t-nam}

Những năm gần đây, thị trường cho thuê xe tự lái tại Việt Nam tăng
trưởng rõ rệt, đặc biệt tại TP. Hồ Chí Minh, Hà Nội và Đà Nẵng. Theo báo
cáo ngành vận tải {[}1{]}, nhu cầu thuê xe tăng trung bình 15--20\% mỗi
năm nhờ sự phục hồi của du lịch nội địa, nhu cầu di chuyển linh hoạt và
xu hướng chia sẻ phương tiện.

Với doanh nghiệp vận hành đội xe, tốc độ tăng trưởng này mang lại cơ hội
kinh doanh nhưng cũng kéo theo một yêu cầu rất thực tế: phải biết xe
đang ở đâu, đang được sử dụng như thế nào và có an toàn hay không, thay
vì chỉ dựa vào các cuộc gọi xác nhận. Khi quy mô đội xe mở rộng, áp lực
giám sát vận hành và bảo toàn tài sản tăng lên nhanh chóng. Trong thực
tế, các doanh nghiệp cho thuê xe tự lái hiện nay thường gặp bốn nhóm khó
khăn sau:

\begin{itemize}
\tightlist
\item
  \textbf{Quản lý thủ công kém hiệu quả}: Nhiều doanh nghiệp quy mô vừa
  và nhỏ vẫn quản lý xe bằng sổ sách, gọi điện xác nhận hoặc dựa vào sự
  chủ động của khách hàng. Cách làm này không cung cấp thông tin thời
  gian thực về vị trí và trạng thái phương tiện {[}2{]}.
\item
  \textbf{Rủi ro mất cắp và sử dụng sai mục đích}: Khi không có hệ thống
  giám sát, xe có thể bị sử dụng vượt phạm vi địa lý đã thỏa thuận, chạy
  quá số km quy định, hoặc trong trường hợp xấu nhất là bị chiếm đoạt.
  Việc phát hiện các tình huống này thường bị trễ, gây thiệt hại lớn về
  tài sản {[}3{]}.
\item
  \textbf{Thiếu dữ liệu chẩn đoán kỹ thuật}: Doanh nghiệp khó theo dõi
  tình trạng kỹ thuật của xe từ xa, nên việc bảo trì thường mang tính bị
  động và chỉ được thực hiện khi sự cố đã xuất hiện. Điều này làm tăng
  chi phí sửa chữa và rút ngắn tuổi thọ phương tiện.
\item
  \textbf{Giải pháp thương mại đắt đỏ}: Các hệ thống định vị và giám sát xe thương
  mại hiện có trên thị trường (như Vietmap, iTracking) thường có chi phí
  cao --- bao gồm phí thiết bị, phí dịch vụ hàng tháng, và phí tích hợp
  --- không phù hợp với các doanh nghiệp quy mô nhỏ với ngân sách hạn
  chế {[}4{]}.
\end{itemize}

Những khó khăn trên cho thấy nhu cầu giám sát là rất rõ, nhưng chưa tự
động dẫn tới một lời giải khả thi. Muốn biến nhu cầu vận hành thành một
hệ thống thật sự dùng được, bài toán phải được chuyển thành các thách
thức kỹ thuật cụ thể.

\subsubsection{1.1.2. Thách thức kỹ
thuật}\label{112-thuxe1ch-thux1ee9c-kux1ef9-thuux1eadt}

Để giải quyết những khó khăn trên, hệ thống không thể chỉ là một thiết
bị định vị đơn giản. Bài toán thực tế là phải dung hòa nhiều yêu cầu kỹ
thuật vốn dễ xung đột với nhau: thiết bị phải tiết kiệm điện nhưng vẫn
luôn sẵn sàng, phải truyền dữ liệu liên tục nhưng vẫn chịu được mất
sóng, và phải đủ đơn giản để triển khai thực tế với chi phí hợp lý. Bốn
thách thức nổi bật là:

\textbf{Thứ nhất, bài toán tiêu thụ năng lượng.} Nếu thiết bị theo dõi
hoạt động liên tục ở trạng thái công suất cao, ắc quy xe có thể bị hao
đáng kể chỉ sau vài tuần, từ đó ảnh hưởng trực tiếp đến khả năng khởi
động {[}5{]}. Vì vậy, hệ thống phải biết khi nào cần hoạt động đầy đủ và
khi nào có thể chuyển sang chế độ tiết kiệm điện.

\textbf{Thứ hai, giám sát khi xe đỗ.} Khi xe tắt máy (IGN OFF), thiết bị
không thể duy trì mọi khối chức năng như lúc xe đang chạy. Tuy nhiên,
nó vẫn phải phát hiện được các chuyển động bất thường, chẳng hạn xe bị
dắt đi hoặc bị cẩu kéo. Đây là bài toán cân bằng giữa độ nhạy phát hiện
và mức tiêu thụ năng lượng {[}6{]}.

\textbf{Thứ ba, độ tin cậy của kết nối.} Hệ thống hoạt động trong môi
trường di động nên kết nối 4G/LTE có thể gián đoạn bất kỳ lúc nào. Điều
người quản lý cần không chỉ là "có mạng thì gửi được", mà là khi mất
mạng thiết bị vẫn không làm mất dữ liệu quan trọng và có thể đồng bộ lại
khi kết nối được phục hồi.

\textbf{Thứ tư, xử lý dữ liệu thời gian thực.} Khi số lượng xe tăng lên
hàng chục hoặc hàng trăm, dịch vụ phía máy chủ phải tiếp nhận liên tục
dữ liệu vị trí, thông số OBD2 và trạng thái cảm biến với độ trễ thấp,
đồng thời vẫn giữ giao diện giám sát dễ quan sát và dễ thao tác cho
người quản lý.

{\thesissetfigureid{1.1}{fig-ch1-overview-problem-solution}
\begin{figure}[htbp]
\centering
\pandocbounded{\safeincludesvg{./assets/figures/01-chuong-1-gioi-thieu-hinh-1-1.svg}}
\caption{Sơ đồ tổng quan vấn đề và giải pháp đề xuất}
\label{fig:ch1-overview-problem-solution}
\end{figure}
}

\subsubsection{1.1.3. Động lực thực hiện dự
án}\label{113-ux111ux1ed9ng-lux1ef1c-thux1ef1c-hiux1ec7n-dux1ef1-uxe1n}

Mục tiêu của đề tài không chỉ là làm ra một thiết bị biết gửi tọa độ
GPS. Điều cần xây dựng là một hệ thống đủ trọn vẹn để người vận hành đội
xe có thể nhìn thấy trạng thái phương tiện, hiểu nguyên nhân của các
cảnh báo và đưa ra quyết định kịp thời. Từ nhu cầu thực tiễn đó, dự án
"IoT Vehicle Tracking System", một hệ thống IoT theo dõi phương tiện,
được triển khai nhằm xây dựng một giải pháp đồng bộ, bao gồm thiết bị
phần cứng trên xe, phần mềm điều khiển thiết bị, hạ tầng máy chủ tiếp
nhận và xử lý dữ liệu, cùng giao diện web để khai thác thông tin. Mục
tiêu là tạo ra một hệ thống có chi phí hợp lý, dễ tùy biến và có thể mở
rộng theo quy mô doanh nghiệp cho thuê xe tự lái tại Việt Nam.

\begin{center}\rule{0.5\linewidth}{0.5pt}\end{center}

\subsection{1.2. Mục tiêu và phạm vi của dự
án}\label{12-mux1ee5c-tiuxeau-vuxe0-phux1ea1m-vi-cux1ee7a-dux1ef1-uxe1n}

\subsubsection{1.2.1. Mục tiêu tổng
quát}\label{121-mux1ee5c-tiuxeau-tux1ed5ng-quuxe1t}

Mục tiêu chính của dự án là thiết kế và hiện thực một hệ thống IoT theo
dõi phương tiện hoàn chỉnh, bao gồm cả phần cứng lẫn phần mềm, phục vụ
cho bài toán quản lý đội xe cho thuê tự lái. Các mục tiêu dưới đây được
xác định từ góc nhìn của đơn vị vận hành đội xe: họ cần thông tin đủ
nhanh để can thiệp, nhưng cũng cần hệ thống đủ bền bỉ để chạy lâu dài
trên xe. Hệ thống cần đáp ứng các yêu cầu sau:

\begin{enumerate}
\def\labelenumi{\arabic{enumi}.}
\tightlist
\item
  \textbf{Theo dõi vị trí thời gian thực}: Cung cấp vị trí định vị vệ
  tinh GPS/GNSS của xe theo chu kỳ 5--30 giây khi xe đang di chuyển, để
  người quản lý nhìn thấy hành trình gần như liên tục trên bản đồ thay
  vì chỉ nhận các mốc vị trí rời rạc.
\item
  \textbf{Giám sát trạng thái kỹ thuật qua OBD2}: Kết nối với cổng chẩn
  đoán OBD-II của xe, nơi cung cấp các thông tin vận hành và lỗi cơ bản,
  thông qua BLE để đọc các thông số như tốc độ, vòng tua máy, nhiệt độ
  động cơ và mã lỗi chẩn đoán (DTC).
\item
  \textbf{Cảnh báo thông minh}: Tự động phát hiện và gửi thông báo khi
  có các sự kiện bất thường như xe vượt khỏi vùng giám sát đã vẽ trên
  bản đồ, vượt tốc độ quy định, chuyển động bất thường khi xe đang đỗ,
  hoặc mất kết nối thiết bị.
\item
  \textbf{Tối ưu hóa năng lượng}: Đảm bảo thiết bị hoạt động liên tục mà
  không làm cạn ắc quy xe, với thời lượng pin dự phòng đủ để hoạt động
  độc lập trong trường hợp mất nguồn chính.
\item
  \textbf{Giao diện quản lý trực quan}: Cung cấp giao diện web dạng bảng
  điều khiển tổng hợp để người quản lý đội xe theo dõi vị trí, xem lịch
  sử hành trình, quản lý cảnh báo và tạo báo cáo.
\end{enumerate}

Năm mục tiêu trên cho biết hệ thống cần đạt được điều gì. Phần tiếp theo
sẽ chốt rõ phạm vi nào được hiện thực trong đồ án và phạm vi nào được
giữ lại cho các giai đoạn phát triển sau.

\subsubsection{1.2.2. Phạm vi dự
án}\label{122-phux1ea1m-vi-dux1ef1-uxe1n}

Để đề tài không dàn trải, phạm vi nghiên cứu được giới hạn theo từng
tầng của hệ thống, từ đối tượng ứng dụng đến phần cứng, firmware và
hạ tầng máy chủ. Cách xác định phạm vi này giúp người đọc thấy rõ đâu là phần đã
được triển khai trong đồ án và đâu là nội dung được giữ lại cho các giai
đoạn phát triển sau.

\textbf{Đối tượng ứng dụng:}

\begin{itemize}
\tightlist
\item
  Loại xe: Ô tô con (xe 4 chỗ, sedan, SUV) với hệ thống điện 12V hoặc
  24V DC
\item
  Ắc quy phổ thông: 40--70 Ah (trung bình 45--60 Ah)
\item
  Đối tượng sử dụng: Các công ty cho thuê xe tự lái quy mô vừa và nhỏ
  tại Việt Nam
\end{itemize}

\textbf{Phạm vi phần cứng:}

\begin{itemize}
\tightlist
\item
  Thiết bị theo dõi IoT sử dụng vi điều khiển ESP32-S3
\item
  Modem LTE + GNSS SIMCom SIM7600CE-T (Auto mode LTE/UMTS/GSM, APN mặc
  định \texttt{internet}) kết nối trực tiếp qua UART1 để cung cấp cả dữ
  liệu 4G và GNSS
\item
  Adapter OBD2 BLE vgate iCar Pro (đọc dữ liệu chẩn đoán xe qua
  Bluetooth)
\item
  Cảm biến gia tốc LIS3DSH (IMU) để phát hiện chuyển động và rung
\item
  Pin dự phòng 18650 1S với mạch sạc và bảo vệ
\end{itemize}

\textbf{Phạm vi firmware:}

\begin{itemize}
\tightlist
\item
  Xây dựng phần mềm trên thiết bị bằng ESP-IDF kết hợp FreeRTOS để điều
  phối các tác vụ nền theo thời gian thực
\item
  Giao tiếp BLE với adapter OBD2 để đọc dữ liệu chẩn đoán của xe
\item
  Truyền dữ liệu lên máy chủ bằng MQTT 3.1.1 qua EMQX, phân tầng QoS 0/1
  theo loại bản tin
\item
  Quản lý năng lượng theo nhiều trạng thái vận hành, gồm lúc xe đang
  chạy, lúc xe đỗ và lúc cần phản ứng khẩn
\end{itemize}

\textbf{Phạm vi máy chủ và dịch vụ xử lý dữ liệu:}

\begin{itemize}
\tightlist
\item
  EMQX làm broker MQTT, nhận và chuyển tiếp dữ liệu từ thiết bị
\item
  MQTT Bridge (dịch vụ trung gian độc lập) xử lý và phân phối dữ liệu tới
  đúng nơi lưu trữ
\item
  PostgreSQL lưu trữ dữ liệu quan hệ (người dùng, xe, khách hàng, cảnh
  báo)
\item
  VictoriaMetrics lưu trữ dữ liệu thay đổi liên tục theo thời gian
\item
  VictoriaLogs lưu trữ nhật ký sự kiện
\item
  API phía máy chủ hiện thực bằng Express.js và được tổ chức theo từng
  nhóm nghiệp vụ để dễ mở rộng
\item
  Giao diện web hiện thực bằng Next.js 15, kết hợp bản đồ và biểu đồ để
  người dùng theo dõi vận hành
\end{itemize}

\textbf{Giới hạn phạm vi nghiên cứu:}

\begin{itemize}
\tightlist
\item
  Không xử lý các chức năng chẩn đoán OBD-II chuyên sâu (chỉ đọc các
  thông số cơ bản)
\item
  Không phát hiện va chạm (IMU chỉ dùng cho phát hiện chuyển động/rung)
\item
  Không điều khiển động cơ xe (tắt bơm xăng, khóa xe) --- chỉ gửi cảnh
  báo
\item
  Không bao gồm ứng dụng di động ở giai đoạn đầu (dự kiến làm ở giai
  đoạn 2 bằng Flutter và tận dụng lại giao diện web hiện có)
\end{itemize}

{\thesissettableid{1.1}{tab-ch1-project-scope-by-layer}
\def\LTcaptype{table}
\begin{longtable}{|L{0.261\linewidth}|L{0.470\linewidth}|L{0.234\linewidth}|}
\caption{Tóm tắt phạm vi dự án theo từng tầng hệ thống}\label{tab:ch1-project-scope-by-layer}\\
\hline \textbf{Tầng hệ thống} & \textbf{Công nghệ chính} & \textbf{Phạm vi} \\
\\\hline
\endfirsthead
\caption[]{Tóm tắt phạm vi dự án theo từng tầng hệ thống (tiếp theo)}\\
\hline \textbf{Tầng hệ thống} & \textbf{Công nghệ chính} & \textbf{Phạm vi} \\
\\\hline
\endhead
\\\hline
\endfoot
\endlastfoot
Phần cứng & ESP32-S3, SIMCom SIM7600CE-T (LTE + GNSS), vgate iCar Pro,
LIS3DSH & Thiết kế PCB, chế tạo và tích hợp thiết bị hoàn chỉnh \\\hline
Firmware & ESP-IDF, FreeRTOS, MQTT 3.1.1, OTA qua HTTPS & Phần mềm trên
thiết bị: đọc dữ liệu, quản lý năng lượng và gửi các gói dữ liệu nhỏ lên
máy chủ \\\hline
MQTT Broker & EMQX 5.x & Broker MQTT nhận và chuyển tiếp dữ liệu từ thiết bị \\\hline
Dịch vụ máy chủ & Express.js, TypeScript, PostgreSQL, VictoriaMetrics & API
phía máy chủ và xử lý dữ liệu \\\hline
Giao diện web & Next.js 15, React 19, Leaflet, ECharts & Bản đồ, biểu
đồ và cảnh báo cho người vận hành \\\hline
Hạ tầng triển khai & Docker, Nginx Proxy Manager, Grafana & Đóng gói
dịch vụ, công bố truy cập và giám sát vận hành \\\hline
\end{longtable}
}

\begin{center}\rule{0.5\linewidth}{0.5pt}\end{center}

\subsection{1.3. Các tiêu chí cần đạt được của dự
án}\label{13-cuxe1c-tiuxeau-chuxed-cux1ea7n-ux111ux1ea1t-ux111ux1b0ux1ee3c-cux1ee7a-dux1ef1-uxe1n}

Dự án đặt ra các tiêu chí cụ thể cho từng tầng của hệ
thống. Các tiêu chí này không chỉ dùng để đánh giá kết quả cuối cùng, mà
còn đóng vai trò như một cam kết thiết kế xuyên suốt quá trình thực hiện:
mỗi quyết định chọn linh kiện, tổ chức firmware hay xây dựng hạ tầng đều
phải quay lại trả lời câu hỏi liệu nó có giúp hệ thống vận hành đủ ổn
định, đủ tiết kiệm năng lượng và đủ hữu ích cho bài toán quản lý xe hay
không.

\subsubsection{1.3.1. Tiêu chí phần
cứng}\label{131-tiuxeau-chuxed-phux1ea7n-cux1ee9ng}

Ở tầng phần cứng, điều quan trọng nhất không phải là "gắn được nhiều khối
chức năng", mà là thiết bị có thể tồn tại bền bỉ trong môi trường trên xe:
đủ tiết kiệm điện, đủ ổn định về nhiệt và đủ an toàn với ắc quy.


{\thesissettableid{1.3.1A}{tab-ch1-tong-hop-thong-tin-ve-stt-tieu-chi-muc-tieu-cu-the}
\def\LTcaptype{table}
\begin{longtable}{|L{0.160\linewidth}|L{0.224\linewidth}|L{0.580\linewidth}|}
\caption{Tổng hợp thông tin về STT, Tiêu chí, Mục tiêu cụ thể}\label{tab:ch1-tong-hop-thong-tin-ve-stt-tieu-chi-muc-tieu-cu-the}\\
\hline \textbf{STT} & \textbf{Tiêu chí} & \textbf{Mục tiêu cụ thể} \\
\\\hline
\endfirsthead
\caption[]{Tổng hợp thông tin về STT, Tiêu chí, Mục tiêu cụ thể (tiếp theo)}\\
\hline \textbf{STT} & \textbf{Tiêu chí} & \textbf{Mục tiêu cụ thể} \\
\\\hline
\endhead
\\\hline
\endfoot
\endlastfoot
1 & Tiêu thụ điện ở chế độ ngủ sâu (deep sleep, tức gần như tắt toàn bộ
khối không cần thiết) & \textless{} 500 µA \\\hline
2 & Tiêu thụ điện chế độ hoạt động & \textless{} 250 mA (trung bình) \\\hline
3 & Thời gian thức dậy từ chế độ ngủ sâu & \textless{} 3 giây \\\hline
4 & Thời lượng pin dự phòng (18650 1S) & \textgreater= 2.5 giờ theo dõi
liên tục; \textgreater= 24 giờ cảnh báo gián đoạn (mục tiêu lấy theo
mức pin tham chiếu của cell 18650 đã chọn) \\\hline
5 & Ngưỡng chuyển nguồn bảo vệ ắc quy & Cấu hình cho hệ 12V:
\(OFF = 12.0\,\mathrm{V}\), \(ON = 12.2\,\mathrm{V}\); Profile 24V:
\(OFF = 24.0\,\mathrm{V}\), \(ON = 24.4\,\mathrm{V}\) \\\hline
6 & Nhiệt độ hoạt động & -10 C đến +60 C \\\hline
\end{longtable}
}


\subsubsection{1.3.2. Tiêu chí
firmware}\label{132-tiuxeau-chuxed-firmware}

Ở tầng firmware, tiêu chí không chỉ dừng ở việc đọc được cảm biến hay gửi
được dữ liệu, mà nằm ở khả năng điều phối đúng nhịp giữa các khối để thiết
bị phản ứng phù hợp với từng trạng thái vận hành của xe.


{\thesissettableid{1.2A}{tab-ch1-tieu-chi-firmware}
\def\LTcaptype{table}
\begin{longtable}{|L{0.160\linewidth}|L{0.404\linewidth}|L{0.400\linewidth}|}
\caption{Tiêu chí firmware}\label{tab:ch1-tieu-chi-firmware}\\
\hline \textbf{STT} & \textbf{Tiêu chí} & \textbf{Mục tiêu cụ thể} \\
\\\hline
\endfirsthead
\caption[]{Tiêu chí firmware (tiếp theo)}\\
\hline \textbf{STT} & \textbf{Tiêu chí} & \textbf{Mục tiêu cụ thể} \\
\\\hline
\endhead
\\\hline
\endfoot
\endlastfoot
1 & Tần suất gửi GPS khi lái xe & 5--30 giây (cấu hình được) \\\hline
2 & Tần suất heartbeat khi đỗ xe & 10--30 phút \\\hline
3 & Thời gian kết nối OBD2 BLE & \textless{} 10 giây \\\hline
4 & Độ chính xác vị trí GNSS & \textless{} 5 mét (điều kiện thoáng) \\\hline
5 & Xử lý mất kết nối & Tự động kết nối lại; có phương án bổ sung buffer
cục bộ \\\hline
6 & Phát hiện chuyển động bất thường & Độ nhạy IMU cấu hình được \\\hline
\end{longtable}
}


\subsubsection{1.3.3. Tiêu chí máy chủ và dịch
vụ xử lý dữ liệu}\label{133-tiuxeau-chuxed-cloudbackend}

Khi chuyển sang tầng máy chủ và các dịch vụ xử lý dữ liệu, trọng tâm
đánh giá đổi sang khả năng tiếp nhận dữ liệu liên tục, phản hồi nhanh và
bảo vệ dữ liệu khi số lượng thiết bị tăng lên theo quy mô đội xe.


{\thesissettableid{1.2B}{tab-ch1-tieu-chi-cloud-backend}
\def\LTcaptype{table}
\begin{longtable}{|L{0.160\linewidth}|L{0.311\linewidth}|L{0.494\linewidth}|}
\caption{Tiêu chí máy chủ và dịch vụ xử lý dữ liệu}\label{tab:ch1-tieu-chi-cloud-backend}\\
\hline \textbf{STT} & \textbf{Tiêu chí} & \textbf{Mục tiêu cụ thể} \\
\\\hline
\endfirsthead
\caption[]{Tiêu chí máy chủ và dịch vụ xử lý dữ liệu (tiếp theo)}\\
\hline \textbf{STT} & \textbf{Tiêu chí} & \textbf{Mục tiêu cụ thể} \\
\\\hline
\endhead
\\\hline
\endfoot
\endlastfoot
1 & Độ trễ xử lý bản tin MQTT & \textless{} 500 ms từ thiết bị tới lớp
xử lý chính \\\hline
2 & Số lượng thiết bị đồng thời & \textgreater= 50 thiết bị (đã kiểm thử), dư địa tới khoảng 100 thiết bị \\\hline
3 & Độ sẵn sàng hệ thống (thời gian hoạt động ổn định) & \textgreater= 99.5\% \\\hline
4 & Thời gian phản hồi API & \textless{} 200 ms (p95) \\\hline
5 & Lưu trữ dữ liệu vận hành gửi liên tục như vị trí, tốc độ, trạng thái nguồn & \textgreater= 90 ngày \\\hline
6 & Bảo mật xác thực & Phiên đăng nhập lưu ở máy chủ; token ngẫu nhiên
chỉ dùng làm mã tham chiếu và được băm SHA-256 \\\hline
\end{longtable}
}


\subsubsection{1.3.4. Tiêu chí giao diện
web}\label{134-tiuxeau-chuxed-frontend}

Cuối cùng, giao diện web là nơi toàn bộ nỗ lực kỹ thuật được chuyển thành trải
nghiệm quan sát và ra quyết định cho người vận hành, nên tiêu chí ở đây tập
trung vào tính cập nhật, dễ theo dõi và dễ sử dụng trong bối cảnh thực tế.


{\thesissettableid{1.2C}{tab-ch1-tieu-chi-frontend}
\def\LTcaptype{table}
\begin{longtable}{|L{0.160\linewidth}|L{0.324\linewidth}|L{0.481\linewidth}|}
\caption{Tiêu chí giao diện web}\label{tab:ch1-tieu-chi-frontend}\\
\hline \textbf{STT} & \textbf{Tiêu chí} & \textbf{Mục tiêu cụ thể} \\
\\\hline
\endfirsthead
\caption[]{Tiêu chí giao diện web (tiếp theo)}\\
\hline \textbf{STT} & \textbf{Tiêu chí} & \textbf{Mục tiêu cụ thể} \\
\\\hline
\endhead
\\\hline
\endfoot
\endlastfoot
1 & Cập nhật bản đồ thời gian thực & \textless{} 2 giây độ trễ \\\hline
2 & Hỗ trợ trình duyệt & Chrome, Firefox, Edge (phiên bản mới nhất) \\\hline
3 & Bố cục thích ứng theo màn hình & Desktop và tablet \\\hline
4 & Hiển thị cảnh báo & Thông báo tức thời qua WebSocket, tức kênh để
máy chủ đẩy dữ liệu mới lên giao diện \\\hline
\end{longtable}
}

Các tiêu chí trên vì thế đóng vai trò như một thước đo chung cho toàn bộ
đồ án: mọi lựa chọn về phần cứng, firmware, máy chủ hay giao diện ở các
chương sau đều phải quay lại đối chiếu với bộ mục tiêu này. Từ nền đó,
phần tiếp theo trình bày cách nhóm tiếp cận bài toán theo từng bước.


\begin{center}\rule{0.5\linewidth}{0.5pt}\end{center}

\subsection{1.4. Phương pháp tiếp cận thiết kế kỹ
thuật}\label{14-phux1b0ux1a1ng-phuxe1p-tiux1ebfp-cux1eadn-thiux1ebft-kux1ebf-kux1ef9-thuux1eadt}

\subsubsection{1.4.1. Phương pháp luận tổng
thể}\label{141-phux1b0ux1a1ng-phuxe1p-luux1eadn-tux1ed5ng-thux1ec3}

Vì đề tài đi qua cả mạch nguồn, phần mềm nhúng trên thiết bị, hạ tầng
máy chủ và màn hình quản trị tổng hợp, rủi ro lớn nhất không nằm ở việc từng khối có chạy
riêng lẻ hay không, mà ở việc chúng có còn phối hợp đúng khi ghép thành
một hệ thống hoàn chỉnh hay không.

Do đó, dự án chọn cách làm theo hai nguyên tắc. Thứ nhất là \textbf{thiết
kế từ dưới lên (bottom-up design)}: kiểm chứng trước các giả định rủi ro
nhất ở tầng thấp rồi mới ghép dần lên các tầng trên. Thứ hai là
\textbf{phát triển lặp (iterative development)}: làm từng vòng nhỏ, đo
được và sửa được, thay vì để đến cuối dự án mới phát hiện xung đột về
nguồn, giao tiếp hay hiệu năng. Cụ thể, tiến trình được tổ chức như sau:

\begin{enumerate}
\def\labelenumi{\arabic{enumi}.}
\tightlist
\item
  \textbf{Giai đoạn 1 --- Nghiên cứu và thiết kế phần cứng}: Xây nền cho
  hệ thống bằng cách khảo sát linh kiện, đối chiếu datasheet, thiết kế
  sơ đồ nguyên lý và tính toán nguồn. Mục tiêu của giai đoạn này là bảo
  đảm thiết bị có thể tồn tại bền bỉ trong môi trường trên xe trước khi
  nói đến các tính năng ở tầng trên.
\item
  \textbf{Giai đoạn 2 --- Phát triển firmware}: Biến sơ đồ phần cứng
  thành một thiết bị thực sự biết ``làm gì khi nào''. Ở giai đoạn này,
  ESP32-S3 được lập trình bằng ESP-IDF để đọc dữ liệu xe qua BLE OBD2,
  giao tiếp với modem qua UART, nhận tín hiệu chuyển động từ IMU qua
  I2C, quyết định khi nào cần thức hoặc ngủ, rồi gửi dữ liệu lên máy chủ
  bằng MQTT. Đây là giai đoạn quyết định thiết bị theo dõi có phản ứng đúng lúc và
  tiết kiệm điện đúng chỗ hay không.
\item
  \textbf{Giai đoạn 3 --- Xây dựng hạ tầng máy chủ}: Dựng ``nơi dữ liệu
  sẽ đi tới'' sau khi rời chiếc xe. Giai đoạn này chuẩn bị các dịch vụ
  nhận bản tin, lưu dữ liệu quản trị, lưu dữ liệu thay đổi liên tục theo
  thời gian và lưu nhật ký vận hành, đồng thời chốt cách các dịch vụ kết
  nối với nhau.
\item
  \textbf{Giai đoạn 4 --- Phát triển API phía máy chủ}: Biến dữ liệu kỹ thuật
  thành những chức năng quản trị có thể dùng được, như đăng nhập, quản
  lý xe, xem lịch sử, tạo cảnh báo và phát cập nhật thời gian thực ra
  giao diện quản trị.
\item
  \textbf{Giai đoạn 5 --- Phát triển giao diện web}: Trình bày dữ liệu dưới
  dạng bản đồ, bảng thông tin và biểu đồ để người vận hành nhìn nhanh,
  hiểu nhanh và thao tác ít nhất có thể.
\item
  \textbf{Giai đoạn 6 --- Tích hợp và kiểm thử}: Ghép toàn bộ các lớp
  lại thành một chuỗi hoàn chỉnh từ xe đến giao diện quản trị, sau đó kiểm thử,
  đo đạc và tối ưu những điểm còn nghẽn.
\end{enumerate}

{\thesissetfigureid{1.2}{fig-ch1-project-development-phases}
\begin{figure}[htbp]
\centering
\pandocbounded{\safeincludesvg{./assets/figures/01-chuong-1-gioi-thieu-hinh-1-2.svg}}
\caption{Quy trình phát triển dự án theo các giai đoạn}
\label{fig:ch1-project-development-phases}
\end{figure}
}

\subsubsection{1.4.2. Kiến trúc hệ
thống}\label{142-kiux1ebfn-truxfac-hux1ec7-thux1ed1ng}

Nếu nhìn theo đường đi của thông tin từ xe đến màn hình quản lý, hệ
thống có thể hình dung qua ba lớp. Lớp thứ nhất nằm trên xe, có
nhiệm vụ thu dữ liệu và gửi dữ liệu đi. Lớp thứ hai nằm trên hạ tầng
hệ thống máy chủ, có nhiệm vụ tiếp nhận, phân luồng và lưu dữ liệu vào
đúng nơi.
Lớp thứ ba là lớp ứng dụng, nơi dữ liệu được biến thành bản đồ, bảng
thông tin, cảnh báo và công cụ quản trị để con người sử dụng. Cách nhìn
này giải thích vì sao kiến trúc được tổ chức theo dạng phân tầng thay vì
để mọi chức năng dồn vào một chỗ.

Trong triển khai thực tế, lớp trên xe gồm các thành phần chuyên đi lấy
dữ liệu và giữ thiết bị hoạt động ổn định, tiêu biểu là ESP32-S3, modem
SIM7600CE-T, cảm biến rung LIS3DSH và adapter OBD2. Lớp giữa làm nhiệm
vụ nhận bản tin, kiểm tra rồi chuyển dữ liệu sang nơi lưu phù hợp như
VictoriaMetrics, VictoriaLogs và API phía máy chủ thông qua EMQX và MQTT
Bridge. Lớp trên cùng là nơi người quản lý làm việc, gồm giao diện web
quản trị và công cụ giám sát. Ứng dụng di động được giữ cho giai đoạn
sau để đồ án tập trung trước vào chuỗi giá trị cốt lõi từ xe đến màn
hình quản trị.

{\thesissetfigureid{1.3}{fig-ch1-system-architecture-overview}
\begin{figure}[htbp]
\centering
\pandocbounded{\safeincludesvg{./assets/figures/01-chuong-1-gioi-thieu-hinh-1-3.svg}}
\caption{Kiến trúc tổng thể hệ thống IoT Vehicle Tracking}
\label{fig:ch1-system-architecture-overview}
\end{figure}
}

\textbf{Chiến lược dữ liệu (Data Strategy):}

Chiến lược dữ liệu được xây theo cùng logic phân tầng này vì dữ liệu
trong hệ thống không đồng nhất: có loại cần tính nhất quán và quan hệ
chặt chẽ, có loại cần ghi liên tục với tần suất cao, và có loại chủ yếu
phục vụ truy vết sự cố. Mỗi nhóm dữ liệu vì thế được đưa về đúng nơi phù
hợp với đặc tính truy vấn và mục đích sử dụng.

\begin{itemize}
\tightlist
\item
  \textbf{PostgreSQL}: Lưu trữ dữ liệu quan hệ có cấu trúc như người
  dùng, xe, khách hàng, hành trình, cảnh báo, vùng giám sát trên bản đồ
  và lệnh điều khiển. Đây là lớp dữ liệu cần quan hệ rõ ràng và ràng
  buộc chặt chẽ.
\item
  \textbf{VictoriaMetrics}: Lưu trữ dữ liệu chuỗi thời gian như tọa độ
  GPS, dữ liệu OBD2 và chỉ số cảm biến. Nhóm dữ liệu này cần ghi nhiều,
  đọc nhanh theo trục thời gian và lưu dài hạn với chi phí hợp lý.
\item
  \textbf{VictoriaLogs}: Lưu trữ nhật ký sự kiện và dấu
  vết ai làm gì và khi nào, cùng các bản
  ghi hoạt động của thiết bị để phục vụ quan sát hệ thống và truy vết sự
  cố.
\end{itemize}

\subsubsection{1.4.3. Chiến lược quản lý năng
lượng}\label{143-chiux1ebfn-lux1b0ux1ee3c-quux1ea3n-luxfd-nux103ng-lux1b0ux1ee3ng}

Thiết kế năng lượng của thiết bị theo dõi thực chất là bài toán quyết định lúc nào
thiết bị cần hoạt động đầy đủ, lúc nào chỉ nên ``canh chừng'', và lúc nào
phải ưu tiên bảo vệ nguồn xe. Nếu xử lý không khéo, hệ thống sẽ rơi vào
hai cực đoan: theo dõi rất sát nhưng làm hao ắc quy, hoặc tiết kiệm điện
quá mức nên bỏ lỡ sự kiện quan trọng. Vì vậy, chiến lược quản lý năng
lượng đa chế độ (multi-mode power management) được xây quanh các trạng
thái vận hành thật của xe, để mỗi thời điểm chỉ bật những khối thật sự
cần thiết:

\textbf{Chế độ 1 --- Khi xe đang chạy (IGN ON, ưu tiên độ đầy đủ dữ liệu):}

\begin{itemize}
\tightlist
\item
  Hệ thống hoạt động toàn bộ công suất
\item
  Gửi vị trí GPS định kỳ (5--30 giây)
\item
  Sạc pin dự phòng từ nguồn xe
\item
  Kết nối máy chủ liên tục để theo dõi thời gian thực
\end{itemize}

\textbf{Chế độ 2 --- Parking Mode (IGN OFF, ưu tiên bảo toàn năng lượng):}

\begin{itemize}
\tightlist
\item
  Vi điều khiển ESP32-S3 vào chế độ ngủ sâu (deep sleep)
\item
  Tiêu thụ cực thấp (\textless{} 500 µA)
\item
  Thức dậy định kỳ (10--30 phút) để gửi heartbeat
\item
  IMU (LIS3DSH) hoạt động độc lập, cảnh rung ở ngưỡng đã cấu hình
\end{itemize}

\textbf{Chế độ 3 --- Cảnh báo khẩn (ưu tiên phản ứng nhanh khi có sự cố):}

\begin{itemize}
\tightlist
\item
  IMU đánh thức ESP32 ngay lập tức qua interrupt
\item
  Bật 4G + GPS, gửi cảnh báo ưu tiên lên máy chủ
\item
  Tiếp tục theo dõi liên tục cho đến khi được xác nhận an toàn
\item
  Sử dụng pin dự phòng nếu nguồn chính bị cắt
\end{itemize}

Bên cạnh đó, hệ thống còn có mạch chuyển nguồn dự phòng với cơ chế
hysteresis để tránh đóng/ngắt liên tục khi điện áp dao động quanh ngưỡng.
Thiết bị tự động chuyển sang pin dự phòng khi điện áp ắc quy tụt xuống
dưới mức an toàn theo từng profile: 12V dùng
\(Switch_{\mathrm{OFF}} = 12.0\,\mathrm{V}\), \(Switch_{\mathrm{ON}} = 12.2\,\mathrm{V}\);
24V dùng \(Switch_{\mathrm{OFF}} = 24.0\,\mathrm{V}\),
\(Switch_{\mathrm{ON}} = 24.4\,\mathrm{V}\), qua đó bảo vệ ắc quy không bị rút
cạn quá mức {[}7{]}.

Nhờ cách tổ chức này, năng lượng không còn là bài toán tối ưu từng linh
kiện riêng lẻ, mà trở thành một phần của logic vận hành toàn hệ thống.

{\thesissetfigureid{1.4}{fig-ch1-power-mode-transitions}
\begin{figure}[htbp]
\centering
\pandocbounded{\safeincludesvg{./assets/figures/01-chuong-1-gioi-thieu-hinh-1-4.svg}}
\caption{Sơ đồ chuyển đổi giữa các chế độ năng lượng}
\label{fig:ch1-power-mode-transitions}
\end{figure}
}

\subsubsection{1.4.4. Công nghệ sử
dụng}\label{144-cuxf4ng-nghux1ec7-sux1eed-dux1ee5ng}

Phần công nghệ dưới đây không nhằm chứng minh dự án dùng nhiều công cụ,
mà để chỉ ra mỗi công cụ giải quyết đúng một việc trong chuỗi vận hành.
Nếu bỏ cách nhìn ``vai trò trước, tên công nghệ sau'', bảng này rất dễ
trở thành danh sách thuật ngữ rời rạc.

Có thể đọc toàn bộ hệ thống theo ba cụm vai trò thay vì học thuộc từng tên
công nghệ. Cụm thứ nhất là \textbf{khối trên xe}: vi điều khiển, modem,
cảm biến và adapter OBD2 chịu trách nhiệm thu dữ liệu và giữ thiết bị
hoạt động ổn định. Cụm thứ hai là \textbf{khối tiếp nhận và lưu trữ}:
firmware điều phối phần cứng, còn hạ tầng máy chủ nhận bản tin, phân
loại, lưu dữ liệu và lưu dấu vết vận hành. Cụm cuối là \textbf{khối khai
thác}: giao diện web, bản đồ, biểu đồ và công cụ giám sát giúp người vận
hành đọc hiểu dữ liệu nhanh. Với cách nhìn này, các tên như ESP32-S3,
EMQX hay Next.js chỉ còn là ví dụ cụ thể cho từng vai trò, không phải
một danh sách thuật ngữ rời rạc.

{\thesissettableid{1.3}{tab-ch1-technology-stack-summary}
\def\LTcaptype{table}
\begin{longtable}{|L{0.223\linewidth}|L{0.334\linewidth}|L{0.115\linewidth}|L{0.282\linewidth}|}
\caption{Tổng hợp công nghệ sử dụng trong dự án}\label{tab:ch1-technology-stack-summary}\\
\hline \textbf{Thành phần} & \textbf{Công nghệ} & \textbf{Phiên bản} & \textbf{Vai trò} \\
\\\hline
\endfirsthead
\caption[]{Tổng hợp công nghệ sử dụng trong dự án (tiếp theo)}\\
\hline \textbf{Thành phần} & \textbf{Công nghệ} & \textbf{Phiên bản} & \textbf{Vai trò} \\
\\\hline
\endhead
\\\hline
\endfoot
\endlastfoot
Vi điều khiển & ESP32-S3 & --- & Vi điều khiển chính, điều phối toàn bộ
thiết bị \\\hline
LTE + GNSS & SIMCom SIM7600CE-T (LTE + GNSS tích hợp) & --- & Truyền dữ
liệu 4G và lấy vị trí từ vệ tinh \\\hline
OBD2 Adapter & vgate iCar Pro & BLE 4.0 & Đọc dữ liệu chẩn đoán xe \\\hline
Cảm biến gia tốc & LIS3DSH & --- & Phát hiện chuyển động và rung \\\hline
Pin dự phòng & 18650 1S Li-ion & 3500 mAh & Nguồn điện dự phòng \\\hline
Framework firmware & ESP-IDF & 5.x & Bộ công cụ để phát triển phần mềm
nhúng cho ESP32 \\\hline
RTOS & FreeRTOS & --- & Lớp điều phối các tác vụ nền theo thời gian thực \\\hline
MQTT Broker & EMQX & 5.x & Trạm trung chuyển bản tin từ thiết bị \\\hline
API phía máy chủ & Express.js + TypeScript & --- & Lớp giao tiếp để
giao diện gọi lấy dữ liệu và gửi lệnh \\\hline
Cơ sở dữ liệu quan hệ & PostgreSQL & 16 & Lưu trữ dữ liệu có cấu trúc \\\hline
Cơ sở dữ liệu chuỗi thời gian & VictoriaMetrics & --- & Lưu trữ
dữ liệu vận hành gửi liên tục như vị trí, trạng thái nguồn và cảm biến \\\hline
Nhật ký sự kiện & VictoriaLogs & --- & Lưu trữ nhật ký vận hành và lỗi \\\hline
Giao diện web & Next.js 15 + React 19 & --- & Giao diện quản lý cho
người vận hành \\\hline
Bản đồ & Leaflet & --- & Hiển thị bản đồ thời gian thực \\\hline
Biểu đồ & ECharts & --- & Trực quan hóa dữ liệu \\\hline
WebSocket & Socket.IO & 4.8 & Kênh để máy chủ đẩy dữ liệu mới lên giao
diện \\\hline
Container hóa & Docker + Docker Compose & --- & Triển khai dịch vụ \\\hline
Giám sát & Grafana + VictoriaMetrics, VictoriaLogs & --- & Màn hình giám
sát hạ tầng và nhật ký vận hành \\\hline
\end{longtable}
}

\begin{center}\rule{0.5\linewidth}{0.5pt}\end{center}

\subsection{1.5. Kết quả và khuyến
nghị}\label{15-kux1ebft-quux1ea3-vuxe0-khuyux1ebfn-nghux1ecb}

\subsubsection{1.5.1. Kết quả đạt
được}\label{151-kux1ebft-quux1ea3-ux111ux1ea1t-ux111ux1b0ux1ee3c}

Nhìn theo cách đó, các kết quả đạt được không nên hiểu như một danh sách
các khối đã lắp xong, mà như bằng chứng cho thấy chuỗi giá trị từ thiết
bị đến giao diện quản trị đã được hình thành. Có thể tóm lược theo bốn tầng chính
như sau:

\textbf{Về phần cứng:}

\begin{itemize}
\tightlist
\item
  Hoàn thiện thiết bị theo dõi IoT với bo mạch riêng do nhóm thiết kế, sử
  dụng ESP32-S3-WROOM-1 làm vi điều khiển trung tâm, tích hợp modem LTE
  + GNSS SIMCom SIM7600CE-T, cảm biến gia tốc LIS3DSH, khối nguồn đa
  rail, pin dự phòng 18650 1S và kết nối OBD2 BLE qua adapter vgate iCar
  Pro.
\item
  Quá trình lựa chọn linh kiện bám theo datasheet, độ ổn định chất lượng
  và khả năng cung ứng trong nước để thuận lợi cho việc chế tạo, kiểm
  thử và thay thế khi cần.
\item
  Hệ thống quản lý năng lượng đa chế độ hoạt động hiệu quả, với mức tiêu
  thụ điện ngủ sâu đạt yêu cầu (\textless{} 500 µA), đảm bảo không làm
  cạn ắc quy xe trong quá trình sử dụng bình thường.
\item
  Mạch Low Voltage Disconnect (LVD) bảo vệ ắc quy xe hiệu quả, tự động
  ngắt khi điện áp tụt dưới ngưỡng an toàn.
\end{itemize}

\textbf{Về firmware:}

\begin{itemize}
\tightlist
\item
  Phát triển thành công phần mềm trên thiết bị bằng ESP-IDF kết hợp
  FreeRTOS, giúp thiết bị theo dõi đọc dữ liệu xe qua BLE OBD2, điều khiển modem
  qua UART, theo dõi chuyển động bằng IMU, quản lý năng lượng và gửi dữ
  liệu MQTT lên hệ thống máy chủ.
\item
  Cơ chế phát hiện chuyển động bất thường qua IMU hoạt động chính xác,
  có khả năng đánh thức hệ thống từ chế độ ngủ sâu trong vòng
  \textless{} 3 giây.
\end{itemize}

\textbf{Về hạ tầng máy chủ và xử lý dữ liệu:}

\begin{itemize}
\tightlist
\item
  Xây dựng thành công hạ tầng máy chủ hoàn chỉnh để nhận dữ liệu từ
  thiết bị, lưu dữ liệu nghiệp vụ, lưu dữ liệu theo thời gian và lưu
  nhật ký vận hành; toàn bộ được đóng gói bằng Docker để dễ triển khai
  và lặp lại môi trường.
\item
  Hoàn thiện API quản trị phía máy chủ cho các nhu cầu cốt lõi như đăng
  nhập, quản lý xe, quản lý thiết bị, xem dữ liệu vận hành, xử lý cảnh
  báo và quản lý vùng giám sát.
\item
  Hệ thống đã kiểm thử ở mức 50+ thiết bị đồng thời với độ trễ từ lúc dữ
  liệu rời thiết bị đến lúc tới nơi xử lý dưới 500 ms, đồng thời cho
  thấy còn dư địa mở rộng ổn định tới khoảng 100 thiết bị.
\end{itemize}

\textbf{Về giao diện web:}

\begin{itemize}
\tightlist
\item
  Giao diện web quản lý (Next.js 15) với bản đồ thời gian thực
  (Leaflet), biểu đồ phân tích (ECharts), hệ thống cảnh báo trực tuyến
  và màn hình tổng quan.
\item
  Cập nhật vị trí xe trên bản đồ với độ trễ \textless{} 2 giây thông qua
  WebSocket (Socket.IO), tức kênh để máy chủ chủ động đẩy dữ liệu mới lên
  trình duyệt.
\end{itemize}

\subsubsection{1.5.2. Khuyến nghị và hướng phát
triển}\label{152-khuyux1ebfn-nghux1ecb-vuxe0-hux1b0ux1edbng-phuxe1t-triux1ec3n}

Trên cơ sở các kết quả đã đạt được và các giới hạn hiện hữu, các hướng phát
triển tiếp theo được sắp theo logic mở rộng từ việc làm chắc hệ lõi đến
việc mở rộng giá trị sử dụng của hệ thống:

\begin{enumerate}
\def\labelenumi{\arabic{enumi}.}
\tightlist
\item
  \textbf{Ứng dụng di động (giai đoạn 2)}: Phát triển ứng dụng Flutter
  theo hướng vỏ ứng dụng mỏng, chủ yếu nhúng lại giao diện web hiện có,
  để người quản lý giám sát đội xe trên điện thoại và nhận thông báo đẩy
  khi có cảnh báo {[}8{]}.
\item
  \textbf{Tăng cường bảo mật}: Triển khai TLS/SSL cho kết nối MQTT trong
  môi trường vận hành thật, áp dụng danh sách quyền MQTT chi tiết hơn để
  mỗi thiết bị chỉ đi đúng kênh dữ liệu của mình, và bổ sung cơ chế mã
  hóa dữ liệu từ thiết bị tới máy chủ.
\item
  \textbf{Tối ưu hóa hiệu năng}: Nghiên cứu và áp dụng các kỹ thuật nén
  dữ liệu vận hành gửi liên tục, chẳng hạn chỉ gửi phần sai khác hoặc
  dùng định dạng gói gọn hơn, để giảm băng thông 4G,
  đặc biệt khi vận hành đội xe lớn (\textgreater{} 500 xe).
\item
  \textbf{Tích hợp trí tuệ nhân tạo}: Áp dụng các mô hình học máy để
  phân tích hành vi lái xe, dự đoán thời điểm cần bảo trì và phát hiện
  các dấu hiệu bất thường từ dữ liệu vận hành đã thu thập.
\item
  \textbf{Mở rộng phạm vi OBD2}: Hỗ trợ đọc thêm nhiều thông số chẩn
  đoán xe, bao gồm mã lỗi DTC (Diagnostic Trouble Codes), dữ liệu động
  cơ nâng cao, và tích hợp với các loại xe điện (EV).
\item
  \textbf{Cải thiện khả năng mở rộng}: Nghiên cứu kiến
  trúc tách nhiều dịch vụ nhỏ, có thêm hàng đợi bản tin chuyên dụng
  (chẳng hạn Kafka hoặc RabbitMQ), để xử lý luồng dữ liệu lớn hơn khi số
  lượng thiết bị tăng lên hàng ngàn.
\item
  \textbf{Tối ưu PCB cho sản xuất}: Nâng thiết kế PCB hiện tại theo
  hướng dễ sản xuất hơn và ít nhiễu điện từ hơn, giảm kích thước, tăng
  độ bền cơ khí và chuẩn bị cho sản xuất hàng loạt.
\end{enumerate}

{\thesissetfigureid{1.5}{fig-ch1-project-roadmap}
\begin{figure}[htbp]
\centering
\pandocbounded{\safeincludesvg{./assets/figures/01-chuong-1-gioi-thieu-hinh-1-5.svg}}
\caption{Lộ trình phát triển dự án theo các giai đoạn}
\label{fig:ch1-project-roadmap}
\end{figure}
}

\begin{center}\rule{0.5\linewidth}{0.5pt}\end{center}

\subsection{Kết luận chương
1}\label{kux1ebft-luux1eadn-chux1b0ux1a1ng-1}

Chương 1 đã trả lời câu hỏi vì sao đề tài này cần được thực hiện và hệ
thống cần đạt được điều gì trong bối cảnh thực tế của doanh nghiệp cho
thuê xe tự lái. Từ bối cảnh thị trường, nhu cầu vận hành đến các thách
thức về năng lượng, kết nối và dữ liệu, chương đã thiết lập bộ mục tiêu,
phạm vi và tiêu chí đánh giá cho toàn bộ đồ án. Nếu Chương 1 là phần
đặt vấn đề, thì Chương 2 sẽ chuyển sang câu hỏi quan trọng hơn: cần giải
quyết chính xác những bài toán kỹ thuật nào để biến mục tiêu đó thành
một hệ thống khả thi.

\chapterheading{CHƯƠNG 2. PHÂN TÍCH VẤN ĐỀ KỸ THUẬT}

Nếu Chương 1 đặt ra bối cảnh và mục tiêu của đề tài, thì Chương 2 đi
sâu vào phần cốt lõi hơn: bóc tách bài toán kỹ thuật thật sự của hệ
thống. Trọng tâm của chương không chỉ là liệt kê công nghệ, mà là làm rõ
vì sao mỗi quyết định về phần cứng, truyền thông, lưu trữ và kiến trúc
đều phải xuất phát từ nhu cầu vận hành thực tế của một hệ thống giám sát
phương tiện.

\subsection{2.1. Mô tả vấn đề -- Problem
statement}\label{21-muxf4-tux1ea3-vux1ea5n-ux111ux1ec1--problem-statement}

\subsubsection{2.1.1. Bối cảnh thực
tế}\label{211-bux1ed1i-cux1ea3nh-thux1ef1c-tux1ebf}

Trong bối cảnh ngành cho thuê ô tô và quản lý đội xe tại Việt Nam tiếp
tục mở rộng, nhu cầu giám sát phương tiện theo thời gian thực đã trở
thành yêu cầu thiết yếu đối với doanh nghiệp vận tải. Theo thống kê của
Bộ Giao thông Vận tải, số lượng ô tô cá nhân và thương mại tăng trung
bình 10--12\% mỗi năm trong giai đoạn 2020--2025, kéo theo áp lực ngày
càng lớn về quản lý, an ninh và tối ưu vận hành {[}1{]}.

Tuy nhiên, điều doanh nghiệp cần không chỉ là các điểm vị trí rời rạc,
mà là khả năng nhìn thấy trạng thái phương tiện đủ sớm để đưa ra quyết
định vận hành. Khoảng cách giữa nhu cầu này và cách làm hiện tại thể
hiện ở các hạn chế sau:

\begin{itemize}
\tightlist
\item
  \textbf{Giám sát thủ công}: Tài xế báo cáo vị trí qua điện thoại,
  không có dữ liệu liên tục, không thể xác minh chính xác hành trình.
  Phương pháp này phụ thuộc hoàn toàn vào yếu tố con người, dễ xảy ra
  sai sót và gian lận.
\item
  \textbf{Thiết bị GPS đơn giản}: Chỉ cung cấp tọa độ vị trí, không đọc
  được dữ liệu động cơ (tốc độ, vòng tua, nhiên liệu), không hỗ trợ phát
  hiện bất thường khi xe đậu. Thiếu khả năng tích hợp sâu với hệ thống
  quản lý.
\item
  \textbf{Giải pháp thương mại (fleet management)}: Chi phí cao (50--200
  USD/thiết bị + phí dịch vụ hàng tháng), phụ thuộc vào nhà cung cấp
  nước ngoài, khó tùy biến theo nhu cầu cụ thể của doanh nghiệp Việt
  Nam.
\end{itemize}

\subsubsection{2.1.2. Xác định vấn đề kỹ
thuật}\label{212-xuxe1c-ux111ux1ecbnh-vux1ea5n-ux111ux1ec1-kux1ef9-thuux1eadt}

Từ bối cảnh thực tiễn đó, bài toán của đề tài có thể quy về bốn câu hỏi
kỹ thuật nền tảng. Đây cũng là bốn câu hỏi chi phối gần như mọi quyết
định ở các chương sau, từ lựa chọn phần cứng đến thiết kế hạ tầng máy
chủ:

\textbf{Vấn đề 1: Quản lý năng lượng trong môi trường ô tô}

Thiết bị theo dõi phải hoạt động liên tục 24/7 trong môi trường ô tô,
nơi nguồn điện luôn bị ràng buộc bởi dung lượng ắc quy. Khi xe tắt máy
(IGN OFF), thiết bị vẫn lấy điện từ ắc quy 12V hoặc 24V của xe; nếu duy
trì chế độ theo dõi liên tục, ắc quy có thể cạn sau vài tuần và làm xe
không khởi động được {[}2{]}. Vì vậy, câu hỏi không đơn thuần là "thiết
bị có hoạt động được không", mà là "thiết bị có hoạt động đủ thông minh
để không gây phiền toái cho chính chiếc xe mà nó đang giám sát hay
không".

\textbf{Vấn đề 2: Giám sát khi xe đậu}

Khi xe tắt máy, hệ thống cần tiết kiệm điện nhưng vẫn phát hiện được các
chuyển động bất thường như bị trộm hoặc bị kéo cẩu. Vì vậy, thiết bị
phải có cơ chế đánh thức nhanh khi phát hiện sự kiện, đồng thời duy trì
dòng tiêu thụ ở mức rất thấp. Khi xe "im lặng", thiết bị cũng phải
"ngủ", nhưng vẫn phải tỉnh dậy đủ nhanh khi có sự cố.

\textbf{Vấn đề 3: Truyền dữ liệu qua mạng di động}

Hệ thống phải gửi đều dữ liệu vị trí và trạng thái xe qua mạng di động,
ngay cả khi chất lượng sóng thay đổi liên tục. Vì vậy, bài toán không
chỉ là ``gửi được hay không'', mà là khi mất mạng thiết bị có tự nối lại
được không, có giữ được dữ liệu quan trọng không, và có mở rộng được khi
số xe tăng lên không. Vì thế, việc chọn giữa MQTT, HTTP hay CoAP chính
là chọn ``cách gửi tin'' phù hợp nhất với một thiết bị phải truyền nhiều
gói nhỏ, đều đặn và đôi khi cần nhận lệnh ngược từ máy chủ. MQTT
phù hợp khi thiết bị cần gửi các gói dữ liệu nhỏ thường xuyên; HTTP gần với kiểu
``gửi yêu cầu rồi chờ trả lời'' quen thuộc của web; còn CoAP là giao thức
gọn hơn, thường dùng cho thiết bị tài nguyên rất hạn chế.

\textbf{Vấn đề 4: Kiến trúc máy chủ có khả năng mở rộng}

Hệ thống cần hỗ trợ nhiều phương tiện đồng thời, cập nhật vị trí theo
thời gian thực qua WebSocket, bảo đảm an toàn bằng cơ chế xác thực
phiên, và xử lý được luồng dữ liệu lớn đến từ nhiều thiết bị IoT. Với
người sử dụng cuối, điều này được cảm nhận rất đơn giản: bản đồ phải cập
nhật kịp thời, cảnh báo phải đến đúng lúc, và dữ liệu phải nhất quán
ngay cả khi số lượng xe tăng lên.

\subsubsection{2.1.3. Mục tiêu kỹ
thuật}\label{213-mux1ee5c-tiuxeau-kux1ef9-thuux1eadt}

Từ bốn vấn đề nền tảng trên, đồ án rút ra bốn nhóm mục tiêu kỹ thuật.
Mỗi nhóm vừa trả lời một nhu cầu vận hành cụ thể của doanh nghiệp cho
thuê xe, vừa trở thành tiêu chí để sàng lọc và lựa chọn giải pháp ở
Chương 3:

\begin{enumerate}
\def\labelenumi{\arabic{enumi}.}
\tightlist
\item
  Thiết kế thiết bị theo dõi có thể vừa xác định vị trí bằng GNSS, vừa
  đọc dữ liệu chẩn đoán từ xe qua OBD2, dựa trên vi điều khiển ESP32-S3;
  đồng thời tối ưu tiêu thụ năng lượng khi xe đậu (mục tiêu dòng deep
  sleep toàn thiết bị dưới 500 µA trong phiên bản phần cứng hiện tại và
  tiếp tục giảm ở các vòng tối ưu PCB tiếp theo).
\item
  Xây dựng cơ chế quản lý nguồn thông minh với pin dự phòng và mạch Low
  Voltage Disconnect (LVD), tức mạch tự ngắt khi điện áp xuống thấp, để
  thiết bị bám được trạng thái xe nhưng không làm ảnh hưởng đến khả năng
  khởi động của phương tiện.
\item
  Thiết kế đường truyền dữ liệu sử dụng MQTT qua mạng 4G/LTE để gửi các
  gói dữ liệu nhỏ qua mạng di động, ưu tiên tự phục hồi kết nối và chừa sẵn
  phương án lưu đệm khi mất mạng kéo dài.
\item
  Phát triển nền tảng máy chủ có khả năng mở rộng, bao gồm API phía máy
  chủ, cơ sở dữ liệu quan hệ cho phần quản trị, cơ sở dữ liệu chuỗi thời
  gian cho dữ liệu phát sinh liên tục, và giao diện web theo dõi thời
  gian thá»±c.
\end{enumerate}

Bốn nhóm mục tiêu trên cho biết đích đến của thiết kế. Tuy nhiên, để
chọn đúng công nghệ và kiến trúc, vẫn cần đặt chúng vào bối cảnh kỹ
thuật rộng hơn. Vì vậy, mục 2.2 chuyển sang rà lại các nền tảng kỹ thuật
liên quan trực tiếp đến bài toán.

\subsection{2.2. Bối cảnh và cơ sở kỹ thuật -- Background and Technical
reviews}\label{22-bux1ed1i-cux1ea3nh-vuxe0-cux1a1-sux1edf-kux1ef9-thuux1eadt--background-and-technical-reviews}

\subsubsection{2.2.1. Tổng quan về hệ thống IoT trong giám sát phương
tiện}\label{221-tux1ed5ng-quan-vux1ec1-hux1ec7-thux1ed1ng-iot-trong-giuxe1m-suxe1t-phux1b0ux1a1ng-tiux1ec7n}

Internet of Things (IoT) là mô hình cho phép các thiết bị vật lý kết nối
với nhau và với hệ thống máy chủ thông qua Internet. Ở góc nhìn đơn giản,
một hệ thống IoT giám sát phương tiện luôn phải đi qua ba bước: thu dữ
liệu trên xe, truyền dữ liệu lên mạng, rồi biến dữ liệu đó thành thông
tin mà người quản lý có thể đọc và hành động. Tương ứng với ba bước đó,
hệ thống thường được tổ chức thành ba lớp chính {[}3{]}:

\begin{itemize}
\tightlist
\item
  \textbf{Lớp cảm biến và thu thập dữ liệu (Perception Layer)}: Đây là
  phần nằm gần chiếc xe nhất. Nó gồm GPS/GNSS để biết xe đang ở đâu, IMU
  để biết xe có rung hoặc di chuyển không, OBD2 để đọc trạng thái vận
  hành của xe, và các cảm biến khác nếu cần. Nhiệm vụ của lớp này là
  ``nhìn thấy'' những gì đang diễn ra trên xe và biến chúng thành dữ
  liệu thô.
\item
  \textbf{Lớp mạng và truyền thông (Network Layer)}: Đây là con đường
  đưa dữ liệu rời khỏi chiếc xe. Lớp này dùng các giao thức như MQTT,
  HTTP hoặc CoAP qua mạng di động để chuyển dữ liệu từ thiết bị lên máy
  chủ. Nói gọn hơn, đây là phần quyết định dữ liệu sẽ được gửi theo kiểu
  ``gói dữ liệu nhỏ gửi liên tục'', ``yêu cầu rồi chờ trả lời'' hay một chuẩn
  gọn hơn cho thiết bị nhỏ. Nó cũng phải xử lý các tình huống mất sóng
  hoặc kết nối lại.
\item
  \textbf{Lớp ứng dụng và xử lý (Application Layer)}: Đây là nơi dữ liệu
  được biến thành thứ con người có thể dùng. Lớp này bao gồm các dịch vụ
  máy chủ như broker MQTT, API phía máy chủ, cơ sở dữ liệu và giao diện
  quản trị như màn hình web tổng hợp hoặc ứng dụng di động.
\end{itemize}

Trong triển khai của đề tài, kiến trúc ba lớp trên được bổ sung thêm một
lớp đệm dữ liệu cục bộ tại thiết bị bằng \textbf{microSD SDMMC 4-bit +
FATFS}. Đây là một thẻ nhớ cục bộ nối với ESP32-S3 bằng chuẩn giao tiếp
tốc độ cao và được quản lý bằng hệ thống tệp quen thuộc. Có thể xem nó
như một ``kho chờ ngắn hạn'': khi mạng di động gián đoạn, dữ liệu vận
hành gửi lên liên tục sẽ tạm nằm ở đây; khi kết nối phục hồi, hệ thống
phát lại theo đúng thứ tự. Nhờ vậy, độ liên tục dữ liệu tăng lên mà
không cần thay đổi lớn ở kiến trúc máy chủ phía trên.

\subsubsection{2.2.2. Các giải pháp giám sát phương tiện hiện
có}\label{222-cuxe1c-giux1ea3i-phuxe1p-giuxe1m-suxe1t-phux1b0ux1a1ng-tiux1ec7n-hiux1ec7n-cuxf3}

{\thesissettableid{2.1}{tab-ch2-vehicle-monitoring-solutions-comparison}
\def\LTcaptype{table}
\begin{longtable}{|L{0.179\linewidth}|L{0.141\linewidth}|L{0.167\linewidth}|L{0.227\linewidth}|L{0.231\linewidth}|}
\caption{So sánh các giải pháp giám sát phương tiện}\label{tab:ch2-vehicle-monitoring-solutions-comparison}\\
\hline \textbf{Tiêu chí} & \textbf{Giám sát thủ công} & \textbf{GPS Tracker đơn giản} & \textbf{Fleet Management
thương mại} & \textbf{Hệ thống đề xuất} \\
\\\hline
\endfirsthead
\caption[]{So sánh các giải pháp giám sát phương tiện (tiếp theo)}\\
\hline \textbf{Tiêu chí} & \textbf{Giám sát thủ công} & \textbf{GPS Tracker đơn giản} & \textbf{Fleet Management
thương mại} & \textbf{Hệ thống đề xuất} \\
\\\hline
\endhead
\\\hline
\endfoot
\endlastfoot
Chi phí thiết bị & Không & 500.000--1.500.000 VND & 1.000.000--5.000.000
VND & 1.514.000 VND \\\hline
Phí dịch vụ hàng tháng & Không & 50.000--100.000 VND & 200.000--500.000
VND & Chi phí 4G SIM (~70.000 VND) \\\hline
Độ chính xác vị trí & Thấp (báo cáo thủ công) & Trung bình (GPS) & Cao
(GPS + A-GPS) & Cao (GNSS đa hệ thống) \\\hline
Dữ liệu động cơ (OBD2) & Không & Không & Có (tùy model) & Có (BLE
OBD2) \\\hline
Phát hiện bất thường khi đậu & Không & Hạn chế & Có & Có (IMU + deep
sleep) \\\hline
Khả năng tùy biến & Không áp dụng & Thấp & Thấp (phụ thuộc nhà cung cấp) & Cao
(mã nguồn mở) \\\hline
Quản lý năng lượng & Không áp dụng & Cơ bản & Tốt & Tốt (đa chế độ + pin
dự phòng) \\\hline
Tích hợp hệ thống & Không & Hạn chế (API riêng) & API do nhà cung cấp & API mở
(REST + WebSocket) \\\hline
\end{longtable}
}

Bảng so sánh cho thấy phương án đề xuất giữ được các năng lực quan trọng
của hệ thống thương mại như độ chính xác vị trí, khả năng đọc OBD2 và
phát hiện bất thường, nhưng chi phí đầu tư thấp hơn và mức độ tùy biến
cao hơn. Với bài toán của doanh nghiệp cho thuê xe trong nước, đây là
hướng lựa chọn phù hợp hơn về cả kỹ thuật lẫn chi phí vận hành.

\subsubsection{2.2.3. So sánh giao thức truyền thông
IoT}\label{223-so-suxe1nh-giao-thux1ee9c-truyux1ec1n-thuxf4ng-iot}

Việc lựa chọn giao thức truyền thông không thể chỉ hiểu là chọn một
``chuẩn phổ biến''. Với người không chuyên, có thể hình dung ba lựa chọn
chính như sau: HTTP giống cách ``gửi yêu cầu rồi chờ trả lời''; MQTT
giống cách ``đưa bản tin vào một kênh chung để bên cần nhận lấy ra'';
còn CoAP là giao thức nhẹ cho thiết bị tài nguyên rất hạn chế,
thường đi trên UDP.

Với hệ thống theo dõi phương tiện, quyết định này chi phối trực tiếp
lượng dữ liệu phải truyền, khả năng chịu mất sóng và cách dịch vụ máy
chủ mở rộng về sau {[}4{]}. Vì thế, việc so sánh không nhằm tìm giao
thức ``tốt nhất nói chung'', mà tìm giao thức phù hợp nhất với nhịp gửi
dữ liệu nhỏ, liên tục và có lúc cần nhận lệnh ngược từ máy chủ.

Trong bảng dưới đây, ``overhead'' là phần dữ liệu phụ bắt buộc đi kèm
mỗi gói tin; còn ``QoS'' là mức cam kết giao bản tin thành công. Hai
tiêu chí này quan trọng vì chúng ảnh hưởng trực tiếp đến băng thông, độ
tin cậy và tuổi thọ pin của thiết bị.

{\thesissettableid{2.2}{tab-ch2-iot-communication-protocols}
\def\LTcaptype{table}
\begin{longtable}{|L{0.279\linewidth}|L{0.206\linewidth}|L{0.232\linewidth}|L{0.238\linewidth}|}
\caption{So sánh các giao thức truyền thông IoT}\label{tab:ch2-iot-communication-protocols}\\
\hline \textbf{Tiêu chí} & \textbf{MQTT} & \textbf{HTTP / REST} & \textbf{CoAP} \\
\\\hline
\endfirsthead
\caption[]{So sánh các giao thức truyền thông IoT (tiếp theo)}\\
\hline \textbf{Tiêu chí} & \textbf{MQTT} & \textbf{HTTP / REST} & \textbf{CoAP} \\
\\\hline
\endhead
\\\hline
\endfoot
\endlastfoot
Mô hình truyền thông & Phát bản tin vào kênh chung, bên nhận đăng ký
theo dõi & Gửi yêu cầu rồi chờ phản hồi &
Gửi yêu cầu rồi chờ phản hồi \\\hline
Giao thức truyền tải bên dưới & TCP & TCP & UDP \\\hline
Phần dữ liệu phụ đi kèm mỗi gói tin & 2 byte (tối thiểu) & Hàng trăm
byte & 4 byte \\\hline
Chất lượng dịch vụ (QoS) & 3 mức (0, 1, 2) & Không có & 2 mức
(có xác nhận / không xác nhận) \\\hline
Tiêu thụ băng thông & Thấp & Cao & Thấp \\\hline
Hỗ trợ trao đổi hai chiều & Có (đăng ký nhận bản tin) & Không, trừ khi
bổ sung cơ chế hỏi lặp hoặc kênh đẩy & Có (Observe) \\\hline
Phù hợp cho IoT & Rất tốt & Trung bình & Tốt \\\hline
Độ phổ biến & Rất cao & Rất cao & Trung bình \\\hline
Mức phù hợp với EMQX & Có sẵn & Có, nhưng phải bổ sung thành phần mở rộng
& Hạn chế \\\hline
\end{longtable}
}

\textbf{Phân tích lựa chọn:} Kết quả so sánh cho thấy MQTT phù hợp hơn
không phải vì nó ``cao cấp'' hơn HTTP hay CoAP, mà vì cách nó trao đổi
dữ liệu gần với nhịp vận hành của hệ thống: thiết bị gửi nhiều bản tin
nhỏ qua mạng di động không phải lúc nào cũng ổn định, đồng thời đôi lúc
cần nhận lệnh ngược từ máy chủ. Cụ thể:

\begin{itemize}
\tightlist
\item
  \textbf{Phần dữ liệu phụ đi kèm rất nhỏ}: Header tối thiểu 2 byte, giảm băng thông và
  tiết kiệm năng lượng cho thiết bị IoT chạy bằng pin.
\item
  \textbf{Mô hình phát bản tin vào kênh chung rồi đăng ký nhận}: Cho phép nhiều thành phần nhận dữ
  liệu cùng lúc mà không làm thiết bị trên xe phải gửi lặp đi lặp lại
  cho từng nơi nhận. Điều này phù hợp với kiến trúc có giao diện quản trị, dịch
  vụ xử lý và các kênh cảnh báo cùng khai thác một nguồn dữ liệu.
\item
  \textbf{QoS linh hoạt}: Hệ thống không cần mọi bản tin đều được bảo
  đảm ở cùng một mức. QoS 0 phù hợp với dữ liệu vận hành gửi dày theo
  chu kỳ
  như vị trí GPS, trong khi QoS 1 phù hợp hơn với cảnh báo và lệnh điều
  khiển vì các bản tin này cần được đảm bảo tới nơi ít nhất một lần.
\item
  \textbf{Lưu trạng thái cuối và thông báo khi thiết bị mất kết nối}: Hai cơ chế này giúp hệ thống
  phản ánh trạng thái online/offline của thiết bị rõ ràng hơn, từ đó
  người vận hành có thể phân biệt giữa "xe không có dữ liệu mới" và
  "thiết bị vừa mất kết nối".
\end{itemize}

\subsubsection{2.2.4. So sánh vi điều khiển
(MCU)}\label{224-so-suxe1nh-vi-ux111iux1ec1u-khiux1ec3n-mcu}

{\thesissettableid{2.3}{tab-ch2-mcu-comparison-for-tracker}
\def\LTcaptype{table}
\begin{longtable}{|L{0.230\linewidth}|L{0.218\linewidth}|L{0.260\linewidth}|L{0.246\linewidth}|}
\caption{So sánh các vi điều khiển cho ứng dụng thiết bị theo dõi IoT}\label{tab:ch2-mcu-comparison-for-tracker}\\
\hline \textbf{Tiêu chí} & \textbf{ESP32-S3} & \textbf{STM32L4} & \textbf{Raspberry Pi Zero 2W} \\
\\\hline
\endfirsthead
\caption[]{So sánh các vi điều khiển cho ứng dụng thiết bị theo dõi IoT (tiếp theo)}\\
\hline \textbf{Tiêu chí} & \textbf{ESP32-S3} & \textbf{STM32L4} & \textbf{Raspberry Pi Zero 2W} \\
\\\hline
\endhead
\\\hline
\endfoot
\endlastfoot
Kiến trúc & Xtensa LX7, dual-core, 240 MHz & ARM Cortex-M4, 80 MHz & ARM
Cortex-A53, quad-core, 1 GHz \\\hline
RAM & 512 KB SRAM + 8 MB PSRAM & 256 KB SRAM & 512 MB DRAM \\\hline
BLE & BLE 5.0 (tích hợp) & Không (cần module ngoài) & BLE 5.0 (tích
hợp) \\\hline
Wi-Fi & 802.11 b/g/n (tích hợp) & Không & 802.11 b/g/n (tích hợp) \\\hline
Tiêu thụ deep sleep & 10--15 µA & 1--2 µA & Không hỗ trợ deep sleep \\\hline
Giá thành (VND) & 80.000--150.000 & 150.000--300.000 &
400.000--600.000 \\\hline
Framework phát triển & Arduino / ESP-IDF & STM32CubeIDE / Mbed & Linux /
Python \\\hline
Cộng đồng hỗ trợ & Rất lớn & Lớn & Rất lớn \\\hline
Độ phù hợp cho thiết bị theo dõi & Cao & Trung bình & Thấp \\\hline
\end{longtable}
}

\textbf{Phân tích lựa chọn:} Nếu chỉ so riêng mức tiết kiệm điện của bản
thân vi điều khiển, STM32L4 có lợi thế rõ rệt. Nhưng thiết bị theo dõi của đề tài
không chỉ cần ``ngủ tiết kiệm'', mà còn phải đồng thời nói chuyện với
OBD2 qua BLE, modem qua UART, cảm biến qua I2C và xử lý logic trạng
thái của toàn thiết bị. Khi nhìn theo bài toán tổng thể đó, ESP32-S3 cho
điểm cân bằng thực tế hơn giữa tích hợp, chi phí và công sức triển khai.
Các lý do chính là:

\begin{itemize}
\tightlist
\item
  \textbf{BLE tích hợp sẵn}: Có thể kết nối trực tiếp với adapter OBD2
  vgate iCar Pro mà không cần gắn thêm module BLE riêng. Điều này làm
  mạch gọn hơn, ít linh kiện hơn và giảm chi phí tích hợp.
\item
  \textbf{Chi phí phù hợp}: Giá thành thấp hơn đáng kể so với STM32L4
  trong khi vẫn đủ năng lực xử lý cho bài toán của đồ án.
\item
  \textbf{Dễ phát triển và thử nghiệm}: Hệ sinh thái ESP-IDF, tài liệu
  và cộng đồng hỗ trợ cho MQTT, BLE và GNSS rất phong phú. Với nhóm làm
  đồ án, đây là lợi thế lớn vì giảm thời gian thử sai và giảm rủi ro tắc
  ở các phần tích hợp.
\item
  \textbf{Mức tiết kiệm điện vẫn đủ cho mục tiêu hệ thống}: Dòng deep
  sleep riêng của ESP32-S3 cao hơn STM32L4, nhưng tuổi thọ pin thực tế
  còn phụ thuộc cả mạch nguồn, cảm biến và linh kiện phụ trợ. Vì vậy,
  chênh lệch ở mức riêng MCU không phải yếu tố duy nhất quyết định hiệu
  quả năng lượng của toàn thiết bị.
\end{itemize}

\subsubsection{2.2.5. So sánh cơ sở dữ
liệu}\label{225-so-suxe1nh-cux1a1-sux1edf-dux1eef-liux1ec7u}

{\thesissettableid{2.4}{tab-ch2-iot-database-solutions}
\def\LTcaptype{table}
\begin{longtable}{|L{0.203\linewidth}|L{0.285\linewidth}|L{0.240\linewidth}|L{0.226\linewidth}|}
\caption{So sánh các giải pháp cơ sở dữ liệu cho hệ thống IoT}\label{tab:ch2-iot-database-solutions}\\
\hline \textbf{Tiêu chí} & \textbf{PostgreSQL + VictoriaMetrics} & \textbf{MongoDB + InfluxDB} & \textbf{TimescaleDB} \\
\\\hline
\endfirsthead
\caption[]{So sánh các giải pháp cơ sở dữ liệu cho hệ thống IoT (tiếp theo)}\\
\hline \textbf{Tiêu chí} & \textbf{PostgreSQL + VictoriaMetrics} & \textbf{MongoDB + InfluxDB} & \textbf{TimescaleDB} \\
\\\hline
\endhead
\\\hline
\endfoot
\endlastfoot
Dữ liệu quan hệ (users, vehicles) & Rất tốt (PostgreSQL) & Trung bình
(MongoDB) & Tốt (PostgreSQL core) \\\hline
Dữ liệu chuỗi thời gian (vị trí, cảm biến, trạng thái gửi liên tục) & Rất tốt (VictoriaMetrics) & Tốt
(InfluxDB) & Tốt (TimescaleDB extension) \\\hline
Hiệu suất ghi & Cao (VictoriaMetrics: hàng triệu điểm/giây) & Cao (InfluxDB:
hàng triệu điểm/giây) & Trung bình \\\hline
Nén dữ liệu & Rất tốt (VictoriaMetrics: 10--70x nén) & Tốt (InfluxDB: TSM) & Tốt
(PostgreSQL compression) \\\hline
Tài nguyên tiêu thụ & Thấp (VictoriaMetrics: 1--2 GB RAM) & Cao (InfluxDB: 4+ GB RAM)
& Trung bình \\\hline
Ngôn ngữ truy vấn & SQL (PostgreSQL) + MetricsQL (VictoriaMetrics) & MongoDB Query +
Flux/InfluxQL & SQL \\\hline
Tích hợp Grafana & Có (native) & Có & Có \\\hline
Giấy phép & Apache 2.0 & SSPL (MongoDB) + MIT (InfluxDB OSS) & Apache
2.0 \\\hline
Độ phức tạp vận hành & Trung bình (2 hệ thống riêng) & Cao (2 hệ thống
khác nhau) & Thấp (1 hệ thống) \\\hline
\end{longtable}
}

\textbf{Phân tích lựa chọn:} Với hệ thống này, không phải mọi dữ liệu
đều có cùng bản chất. Dữ liệu người dùng, xe, chuyến đi hay cảnh báo cần
quan hệ rõ ràng giữa các bảng; trong khi dữ liệu vị trí và cảm biến lại
là những bản ghi nhỏ, phát sinh liên tục theo thời gian. Vì vậy, thay vì
ép tất cả vào một chỗ, đề tài chọn mỗi loại dữ liệu vào đúng nơi phù
hợp: PostgreSQL cho dữ liệu nghiệp vụ và VictoriaMetrics cho dữ liệu
chuỗi thời gian.

\begin{itemize}
\tightlist
\item
  \textbf{Phân tách trách nhiệm rõ ràng}: PostgreSQL lưu những dữ liệu
  cần liên kết chặt như người dùng, phương tiện, cảnh báo và vùng giám
  sát; VictoriaMetrics lưu những dòng dữ liệu tăng liên tục như GPS,
  OBD2 và dữ liệu cảm biến gửi theo chu kỳ. Mỗi hệ thống chỉ phải làm đúng phần việc mà nó phù
  hợp nhất.
\item
  \textbf{Hiệu suất và tài nguyên}: VictoriaMetrics nén dữ liệu tốt và
  dùng ít RAM hơn nhiều phương án cùng loại, phù hợp với hạ tầng máy chủ
  của đồ án.
\item
  \textbf{Thuận lợi cho giám sát}: Cả hai đều tích hợp tốt với Grafana,
  nên việc dựng màn hình giám sát kỹ thuật và theo dõi vận hành không cần thêm
  lớp chuyển đổi phức tạp.
\end{itemize}

{\thesissetfigureid{2.1}{fig-ch2-data-architecture-overview}
\begin{figure}[htbp]
\centering
\pandocbounded{\safeincludesvg{./assets/figures/02-chuong-2-phan-tich-hinh-2-1.svg}}
\caption{Sơ đồ kiến trúc dữ liệu của hệ thống}
\label{fig:ch2-data-architecture-overview}
\end{figure}
}

\subsection{2.3. Yêu cầu kỹ thuật và các tiêu chuẩn thiết kế -- Design
criteria and
Constraints}\label{23-yuxeau-cux1ea7u-kux1ef9-thuux1eadt-vuxe0-cuxe1c-tiuxeau-chuux1ea9n-thiux1ebft-kux1ebf--design-criteria-and-constraints}

\subsubsection{2.3.1. Yêu cầu chức
năng}\label{231-yuxeau-cux1ea7u-chux1ee9c-nux103ng}

Những yêu cầu chức năng dưới đây đánh dấu bước chuyển từ ``bài toán''
sang ``năng lực hệ thống''. Có thể xem chúng như năm lời hứa vận hành tối
thiểu mà hệ thống phải thực hiện được để người quản lý thật sự dùng được
trong bối cảnh đội xe vận hành ngoài thực tế.

\textbf{FC-01: Theo dõi vị trí thời gian thực}

\begin{itemize}
\tightlist
\item
  Gửi vị trí GPS định kỳ (5--30 giây khi lái xe, 10--30 phút khi đậu xe)
\item
  Độ chính xác vị trí: dưới 5 mét trong điều kiện trời quang
\item
  Hỗ trợ đa hệ thống định vị (GPS, GLONASS, BeiDou)
\end{itemize}

\textbf{FC-02: Đọc dữ liệu động cơ qua OBD2}

\begin{itemize}
\tightlist
\item
  Kết nối BLE với OBD2 adapter (vgate iCar Pro)
\item
  Đọc trạng thái IGN (bật/tắt máy), RPM, tốc độ, nhiên liệu
\item
  Tần suất đọc: 5--30 giây khi xe chạy (mặc định 5 giây)
\end{itemize}

\textbf{FC-03: Phát hiện bất thường khi đậu xe}

\begin{itemize}
\tightlist
\item
  IMU (LIS3DSH) phát hiện chuyển động/rung bất thường
\item
  Đánh thức ESP32 từ deep sleep trong vòng 100ms
\item
  Gửi cảnh báo ưu tiên cao lên máy chủ
\end{itemize}

\textbf{FC-04: Quản lý năng lượng thông minh}

\begin{itemize}
\tightlist
\item
  Chuyển đổi giữa 3 chế độ: xe đang chạy (hoạt động đầy đủ), xe đang đỗ
  (tiết kiệm điện) và cảnh báo bất thường
\item
  Pin dự phòng tự động tiếp quản theo profile nguồn: hệ 12V tại ngưỡng
  \(Switch_{\mathrm{OFF}} = 12.0\,\mathrm{V}\), hệ 24V tại ngưỡng
  \(Switch_{\mathrm{OFF}} = 24.0\,\mathrm{V}\)
\item
  Có phương án dự phòng suy luận trạng thái IGN qua ADC khi kết nối OBD2
  BLE không sẵn sàng
\end{itemize}

\textbf{FC-05: Giao diện web giám sát}

\begin{itemize}
\tightlist
\item
  Bản đồ thời gian thực (Leaflet) hiển thị vị trí tất cả phương tiện
\item
  Biểu đồ dữ liệu OBD2 (ECharts) - tốc độ, RPM, nhiên liệu
\item
  Hệ thống cảnh báo và thông báo tự động
\item
  Quản lý phương tiện, tài xế, khách hàng, hành trình
\end{itemize}

Năm nhóm chức năng trên bao quát trọn chuỗi nhu cầu cốt lõi của hệ
thống: biết xe đang ở đâu, hiểu xe đang như thế nào, phát hiện khi có
bất thường, không làm phiền hệ thống điện của xe và trình bày thông tin ở
dạng mà người vận hành có thể đọc nhanh, xử lý nhanh.

\subsubsection{2.3.2. Yêu cầu phi chức
năng}\label{232-yuxeau-cux1ea7u-phi-chux1ee9c-nux103ng}

Nếu yêu cầu chức năng trả lời ``hệ thống phải làm được gì'', thì yêu cầu
phi chức năng trả lời ``hệ thống phải làm tốt tới mức nào''. Đây là lớp
tiêu chí cho biết nguyên mẫu có đủ độ chín để bước ra môi trường vận
hành thật hay chưa, đặc biệt khi phải theo dõi nhiều xe và chạy liên
tục.

\textbf{NFR-01: Hiệu suất}

\begin{itemize}
\tightlist
\item
  Độ trễ toàn tuyến (từ thiết bị đến giao diện quản trị): dưới 3 giây trong điều
  kiện mạng bình thường
\item
  Thời gian phản hồi API: dưới 200 ms cho 95\% số yêu cầu
\item
  Hỗ trợ đồng thời tối thiểu 50 phương tiện kết nối MQTT, với dư địa mở
  rộng ổn định tới khoảng 100 thiết bị
\end{itemize}

\textbf{NFR-02: Độ tin cậy}

\begin{itemize}
\tightlist
\item
  Tự động kết nối lại khi mạng phục hồi
\item
  Giữ được cấu hình và trạng thái quan trọng qua NVS; lưu đệm dữ liệu
  vận hành cục bộ là hướng mở rộng của thiết bị
\item
  MQTT QoS 1 cho các tin nhắn cảnh báo (đảm bảo gửi ít nhất một lần)
\end{itemize}

\textbf{NFR-03: Bảo mật}

\begin{itemize}
\tightlist
\item
  Xác thực phiên bằng token ngẫu nhiên không tự mang thông tin người
  dùng; phía máy chủ mới là nơi lưu và kiểm tra trạng thái phiên, còn
  trong cơ sở dữ liệu chỉ lưu bản băm SHA-256 của token
\item
  MQTT ACL (Access Control List), tức danh sách quyền truy cập theo từng
  thiết bị
\item
  Mã hóa TLS cho kết nối MQTT và HTTPS trong môi trường sản xuất
\end{itemize}

\textbf{NFR-04: Khả năng mở rộng}

\begin{itemize}
\tightlist
\item
  Kiến trúc nhiều dịch vụ tách rời, mỗi dịch vụ được đóng gói thành
  một gói chạy riêng bằng Docker
\item
  MQTT broker (EMQX) hỗ trợ chạy thành cụm, tức nhiều máy có thể cùng
  phục vụ khi số thiết bị tăng
\item
  Cơ sở dữ liệu chuỗi thời gian (VictoriaMetrics) tối ưu cho luồng dữ
  liệu lớn
\end{itemize}

\textbf{NFR-05: Khả năng bảo trì}

\begin{itemize}
\tightlist
\item
  Mã nguồn tổ chức theo Domain-Driven Design (DDD), nghĩa là chia theo
  từng nhóm nghiệp vụ như xe, thiết bị hay cảnh báo thay vì gom theo
  loại file kỹ thuật
\item
  API theo kiểu REST với tài liệu tự sinh bằng Swagger/OpenAPI, tức
  trang mô tả sẵn các API để người phát triển và người kiểm thử cùng tra
  cứu
\item
  Ghi log tập trung với VictoriaLogs và theo dõi trên Grafana, tức màn
  hình giám sát chung của hệ thống
\end{itemize}

Nhóm yêu cầu trên mô tả hệ thống cần làm được gì và cần đạt tới mức nào.
Phần tiếp theo thu hẹp thêm một bước: trong khuôn khổ đồ án, nhóm thực
sự bị giới hạn bởi những điều kiện nào về phần cứng, phần mềm và chi
phí.

\subsubsection{2.3.3. Ràng buộc thiết
kế}\label{233-ruxe0ng-buux1ed9c-thiux1ebft-kux1ebf}

Các ràng buộc sau đây giới hạn không gian thiết kế của đề tài. Chúng buộc
nhóm phải chọn phương án cân bằng giữa tính khả thi học thuật, khả năng
triển khai thực tế trên xe và ngân sách sẵn có.

\textbf{Ràng buộc phần cứng:}

\begin{itemize}
\tightlist
\item
  Điện áp đầu vào: 12V hoặc 24V DC từ ắc quy xe (dao động phụ thuộc cấu
  hình hệ thống điện)
\item
  Nhiệt độ hoạt động: -10 độ C đến 70 độ C (môi trường trong xe ô tô)
\item
  Rung động và sốc: Chịu được rung động liên tục khi xe vận hành trên
  đường xá
\item
  Kích thước: Đủ nhỏ để lắp đặt kín đáo trong xe (không lớn hơn
  120x80x40 mm)
\end{itemize}

\textbf{Ràng buộc phần mềm:}

\begin{itemize}
\tightlist
\item
  Framework API phía máy chủ: ưu tiên Express.js + TypeScript vì đây là nền tảng
  nhóm đã có sẵn kinh nghiệm, giúp rút ngắn thời gian học lại từ đầu
\item
  Framework giao diện web: chọn Next.js 15 + React 19 để phù hợp định hướng
  học thuật và thuận lợi cho màn hình quản trị có nhiều khu vực hiển thị
\item
  MQTT Broker: chọn EMQX vì là giải pháp mã nguồn mở ổn định, đủ đáp ứng
  vai trò broker MQTT cho giai đoạn hiện tại
\item
  Triển khai: dùng Docker Compose trên máy chủ Linux để mỗi dịch vụ có
  thể chạy tách riêng nhưng vẫn dễ dựng lại môi trường
\end{itemize}

\textbf{Ràng buộc về chi phí:}

\begin{itemize}
\tightlist
\item
  Tổng chi phí phần cứng: dưới 2.000.000 VND cho một bộ thiết bị theo dõi
\item
  Ưu tiên các linh kiện có datasheet rõ ràng, chất lượng ổn định và
  nguồn cung trong nước dễ tiếp cận
\item
  Ưu tiên các giải pháp mã nguồn mở để giảm chi phí giấy phép
\end{itemize}

\subsubsection{2.3.4. Tiêu chuẩn thiết kế áp
dụng}\label{234-tiuxeau-chuux1ea9n-thiux1ebft-kux1ebf-uxe1p-dux1ee5ng}

Bảng dưới đây tổng hợp các chuẩn kỹ thuật và quy ước triển khai được dùng
làm nền cho hệ thống. Điểm quan trọng không phải là gom thật nhiều chuẩn,
mà là chọn đúng những chuẩn thực sự chi phối khả năng tương thích giữa
thiết bị, mạng truyền và ứng dụng.


{\thesissettableid{2.3.4A}{tab-ch2-tong-hop-thong-tin-ve-tieu-chuan-mo-ta-ap-dung-trong}
\def\LTcaptype{table}
\begin{longtable}{|L{0.231\linewidth}|L{0.278\linewidth}|L{0.455\linewidth}|}
\caption{Tổng hợp chuẩn và quy ước kỹ thuật áp dụng trong hệ thống}\label{tab:ch2-tong-hop-thong-tin-ve-tieu-chuan-mo-ta-ap-dung-trong}\\
\hline \textbf{Chuẩn / quy ước} & \textbf{Mô tả} & \textbf{Áp dụng trong hệ thống} \\
\\\hline
\endfirsthead
\caption[]{Tổng hợp chuẩn và quy ước kỹ thuật áp dụng trong hệ thống (tiếp theo)}\\
\hline \textbf{Chuẩn / quy ước} & \textbf{Mô tả} & \textbf{Áp dụng trong hệ thống} \\
\\\hline
\endhead
\\\hline
\endfoot
\endlastfoot
IEEE 802.15.1 (Bluetooth) & Chuẩn BLE 4.0/5.0 & Kết nối OBD2 adapter \\\hline
3GPP LTE Cat-4 & Chuẩn 4G/LTE & Truyền dữ liệu qua modem SIM7600CE-T
(Cat-4) \\\hline
NMEA 0183 & Chuẩn dữ liệu GPS & Phân tích tọa độ từ khối GNSS tích hợp \\\hline
SAE J1979 (OBD2) & Chuẩn chẩn đoán động cơ & Đọc dữ liệu qua OBD2 BLE
adapter \\\hline
MQTT v5.0 (OASIS) & Giao thức nhắn tin cho thiết bị IoT & Truyền dữ liệu thiết bị -
máy chủ \\\hline
REST (RFC 7231) & Cách tổ chức API theo tài nguyên & API phía máy chủ \\\hline
Mã phiên ngẫu nhiên + ACL MQTT & Cơ chế xác thực và phân quyền &
Xác thực phiên người dùng, giới hạn quyền truy cập thiết bị theo ACL
MQTT \\\hline
\end{longtable}
}


\subsection{2.4. Yêu cầu từ các bên liên quan -- Stakeholder
requirements}\label{24-yuxeau-cux1ea7u-tux1eeb-cuxe1c-buxean-liuxean-quan--stakeholder-requirements}

\subsubsection{2.4.1. Xác định các bên liên
quan}\label{241-xuxe1c-ux111ux1ecbnh-cuxe1c-buxean-liuxean-quan}

Hệ thống IoT Vehicle Tracking System liên quan đến nhiều nhóm sử dụng với
nhu cầu khác nhau. Nếu không tách rõ ngay từ đầu, hệ thống rất dễ rơi vào
tình trạng mạnh về mặt kỹ thuật nhưng lệch khỏi nhu cầu vận hành. Vì vậy,
phần này xác định ai là người hưởng lợi trực tiếp, ai là người vận hành,
ai là người bị tác động, từ đó giữ trọng tâm thiết kế vào những nhu cầu
thực sự cần giải quyết.

\subsubsection{2.4.2. Ma trận yêu cầu các bên liên
quan}\label{242-ma-trux1eadn-yuxeau-cux1ea7u-cuxe1c-buxean-liuxean-quan}


{\thesissettableid{2.5}{tab-ch2-ma-tran-yeu-cau-cac-ben-lien-quan}
\def\LTcaptype{table}
\begin{longtable}{|L{0.085\linewidth}|L{0.191\linewidth}|L{0.162\linewidth}|L{0.380\linewidth}|L{0.126\linewidth}|}
\caption{Ma trận yêu cầu các bên liên quan}\label{tab:ch2-ma-tran-yeu-cau-cac-ben-lien-quan}\\
\hline \textbf{STT} & \textbf{Bên liên quan} & \textbf{Vai trò} & \textbf{Yêu cầu chính} & \textbf{Mức độ ưu tiên} \\
\\\hline
\endfirsthead
\caption[]{Ma trận yêu cầu các bên liên quan (tiếp theo)}\\
\hline \textbf{STT} & \textbf{Bên liên quan} & \textbf{Vai trò} & \textbf{Yêu cầu chính} & \textbf{Mức độ ưu tiên} \\
\\\hline
\endhead
\\\hline
\endfoot
\endlastfoot
1 & Công ty cho thuê xe & Chủ sở hữu hệ thống & Giám sát vị trí tất cả
xe theo thời gian thực; Nhận cảnh báo khi xe đi ra khỏi vùng cho phép; Báo cáo hành trình, quãng đường, nhiên liệu; Giảm thiểu thời
gian xe đậu không sử dụng & Cao \\\hline
2 & Tài xế / Người thuê xe & Người vận hành xe & Không bị giám sát quá
mức (quyền riêng tư); Thiết bị không ảnh hưởng đến vận hành xe; Không
rút cạn ắc quy xe & Trung bình \\\hline
3 & Quản trị viên hệ thống & Người quản lý kỹ thuật & Giao diện quản lý
dễ sử dụng; Hệ thống ổn định, ít lỗi; Có khả năng theo dõi trạng thái
thiết bị (online/offline); Quản lý người dùng và phân quyền & Cao \\\hline
4 & Đội bảo trì & Kỹ thuật viên & Thiết bị dễ lắp đặt và bảo trì; Có thể
kiểm tra tình trạng thiết bị từ xa mà không cần tháo xuống; Cảnh báo bảo
trì định kỳ (pin yếu, thiết bị ngoại tuyến) & Trung bình \\\hline
5 & Khách hàng (người thuê xe cuối) & Người dùng gián tiếp & Quá trình
thuê xe nhanh gọn; Tin tưởng vào độ bảo mật của dữ liệu cá nhân; Nhận
thông báo về trạng thái đơn thuê & Thấp (đưa vào giai đoạn 2) \\\hline
\end{longtable}
}


\subsubsection{2.4.3. Phân tích chi tiết yêu
cầu}\label{243-phuxe2n-tuxedch-chi-tiux1ebft-yuxeau-cux1ea7u}

Từ ma trận trên, có thể thấy không phải mọi bên liên quan đều đòi hỏi cùng
một loại giá trị. Nhóm sử dụng chính quan tâm đến hiệu quả vận hành và an
ninh; nhóm quản trị quan tâm đến tính ổn định và khả năng kiểm soát; còn
nhóm bảo trì cần hệ thống dễ chẩn đoán và dễ can thiệp khi có sự cố.

\textbf{Công ty cho thuê xe (Stakeholder chính)}

Đây là nhóm đối tượng trọng tâm của hệ thống. Các yêu cầu của nhóm này
tập trung vào ba khía cạnh:

\begin{enumerate}
\def\labelenumi{\arabic{enumi}.}
\tightlist
\item
  \emph{Giám sát và an ninh}: Theo dõi vị trí phương tiện 24/7, phát
  hiện và cảnh báo ngay khi có bất thường như xe bị di chuyển trái phép,
  vượt khỏi vùng giám sát đã vẽ trên bản đồ hoặc chạy quá tốc độ. Hệ
  thống cần gửi thông báo qua nhiều kênh như giao diện web quản trị,
  Telegram bot và email.
\item
  \emph{Báo cáo và phân tích}: Xuất báo cáo hành trình chi tiết (quãng
  đường, thời gian, điểm dừng), thống kê nhiên liệu tiêu thụ, đánh giá
  hành vi lái xe (tốc độ trung bình, số lần phanh gấp). Các báo cáo này
  giúp tối ưu hóa chi phí vận hành và bảo trì.
\item
  \emph{Quản lý đội xe}: Quản lý trạng thái từng phương tiện (đang cho
  thuê, đang bảo trì, sẵn sàng), lịch bảo trì định kỳ dựa trên số km
  hoặc thời gian, lịch sử cho thuê và doanh thu theo phương tiện.
\end{enumerate}

\textbf{Quản trị viên hệ thống}

Yêu cầu của quản trị viên tập trung vào khả năng vận hành và giám sát hệ
thống:

\begin{enumerate}
\def\labelenumi{\arabic{enumi}.}
\tightlist
\item
  \emph{Trang tổng quan}: Hiển thị trạng thái tất cả thiết bị
  (online/offline, mức pin, cường độ tín hiệu), số lượng phương tiện
  đang hoạt động, cảnh báo chưa xử lý.
\item
  \emph{Quản lý thiết bị}: Thêm/xóa/sửa thông tin thiết bị theo dõi, cấp
  phát thiết bị cho phương tiện, gửi lệnh điều khiển từ xa (reset, cập
  nhật cấu hình, yêu cầu vị trí).
\item
  \emph{Quản lý người dùng}: Tạo tài khoản, phân quyền (admin, manager,
  viewer), theo dõi lịch sử đăng nhập và thao tác thông qua audit log,
  tức bản ghi ai đã làm gì trong hệ thống.
\end{enumerate}

\textbf{Đội bảo trì}

Yêu cầu của đội bảo trì hướng đến việc giảm thời gian và chi phí bảo
trì:

\begin{enumerate}
\def\labelenumi{\arabic{enumi}.}
\tightlist
\item
  \emph{Chẩn đoán từ xa}: Xem trạng thái thiết bị (mức pin, nhiệt độ,
  cường độ tín hiệu 4G) từ xa, phát hiện và cảnh báo khi thiết bị gặp
  lỗi (mất kết nối kéo dài, pin yếu).
\item
  \emph{Hướng dẫn lắp đặt}: Tài liệu hướng dẫn lắp đặt chi tiết, danh
  sách kiểm tra (checklist) khi lắp đặt mới hoặc bảo trì.
\end{enumerate}

Nhìn tổng thể, ba nhóm yêu cầu này kéo hệ thống về một điểm cân bằng khá
rõ: phải đủ mạnh để quản lý đội xe, nhưng cũng phải đủ gọn và ổn định để
lắp đặt, vận hành và bảo trì trên thực địa. Đây cũng là cơ sở để ma trận
truy xuất ở mục tiếp theo ánh xạ yêu cầu định tính sang các chức năng cụ
thể của hệ thống.

\subsubsection{2.4.4. Ma trận truy xuất yêu cầu - chức
năng}\label{244-ma-trux1eadn-truy-xuux1ea5t-yuxeau-cux1ea7u---chux1ee9c-nux103ng}


{\thesissettableid{2.6}{tab-ch2-ma-tran-truy-xuat-yeu-cau-cac-ben-lien-quan-voi-chuc}
\def\LTcaptype{table}
\begin{longtable}{|L{0.237\linewidth}|L{0.122\linewidth}|L{0.122\linewidth}|L{0.186\linewidth}|L{0.168\linewidth}|L{0.101\linewidth}|}
\caption{Ma trận truy xuất yêu cầu các bên liên quan với chức năng hệ thống}\label{tab:ch2-ma-tran-truy-xuat-yeu-cau-cac-ben-lien-quan-voi-chuc}\\
\hline \textbf{Yêu cầu bên liên quan} & \textbf{FC-01 (Vị trí)} & \textbf{FC-02 (OBD2)} & \textbf{FC-03 (Bất
thường)} & \textbf{FC-04 (Năng lượng)} & \textbf{FC-05 (Web)} \\
\\\hline
\endfirsthead
\caption[]{Ma trận truy xuất yêu cầu các bên liên quan với chức năng hệ thống (tiếp theo)}\\
\hline \textbf{Yêu cầu bên liên quan} & \textbf{FC-01 (Vị trí)} & \textbf{FC-02 (OBD2)} & \textbf{FC-03 (Bất
thường)} & \textbf{FC-04 (Năng lượng)} & \textbf{FC-05 (Web)} \\
\\\hline
\endhead
\\\hline
\endfoot
\endlastfoot
Giám sát vị trí 24/7 & X & & & X & X \\\hline
Cảnh báo vượt vùng cho phép & X & & & & X \\\hline
Báo cáo nhiên liệu & & X & & & X \\\hline
Phát hiện trộm xe & & & X & X & X \\\hline
Không rút cạn ắc quy & & & & X & \\\hline
Quản lý thiết bị từ xa & & & & & X \\\hline
Chẩn đoán lỗi từ xa & X & X & & X & X \\\hline
\end{longtable}
}


Ma trận truy xuất trên cho thấy các chức năng hệ thống (FC-01 đến FC-05)
bao phủ toàn bộ yêu cầu của các bên liên quan chính. Trong đó, FC-04
(Quản lý năng lượng) và FC-05 (Giao diện web) là hai chức năng được nhắc
đến nhiều nhất, qua đó khẳng định tầm quan trọng của thiết kế quản lý
năng lượng thông minh và giao diện giám sát thân thiện.

\begin{center}\rule{0.5\linewidth}{0.5pt}\end{center}

\subsection{Kết luận chương
2}\label{kux1ebft-luux1eadn-chux1b0ux1a1ng-2}

Chương 2 đã chuyển bài toán ở mức "ý tưởng" thành các ràng buộc kỹ
thuật cụ thể: thiết bị phải tiết kiệm điện ra sao, truyền dữ liệu thế
nào, lưu trữ ở đâu và kiến trúc cần mở rộng theo hướng nào. Thông qua
quá trình khảo sát và đối sánh các giải pháp hiện có, chương đã làm rõ
tiêu chí lựa chọn thay vì chỉ liệt kê công nghệ. Đây chính là cơ sở để
Chương 3 bước sang phần quyết định quan trọng tiếp theo: chọn phương án
thiết kế nào là phù hợp nhất với bài toán đã được phân tích.

\chapterheading{CHƯƠNG 3. CÁC GIẢI PHÁP THIẾT KẾ -- DESIGN SOLUTIONS}

Sau khi xác định rõ yêu cầu và ràng buộc ở Chương 2, Chương 3 tập trung
vào câu hỏi "chọn gì và vì sao chọn như vậy". Chương này không chỉ nêu
phương án cuối cùng, mà còn trình bày quá trình cân nhắc giữa nhiều lựa
chọn ở bốn tầng: phần cứng thiết bị theo dõi, phần mềm nhúng trên vi
điều khiển, hạ tầng máy chủ xử lý dữ liệu và giao diện người dùng.

Nhờ cách trình bày đó, người đọc có thể theo dõi logic thiết kế theo
từng bước thay vì chỉ nhìn thấy kết quả cuối cùng.

\subsection{3.1. Phân tích và đề xuất giải pháp phần cứng -- Hardware
analysis and solution}\label{31-phuxe2n-tuxedch-tux1ed5ng-hux1ee3p--general-analysis}

\subsubsection{3.1.1. Phân tích và lựa chọn phần
cứng}\label{311-phuxe2n-tuxedch-vuxe0-lux1ef1a-chux1ecdn-phux1ea7n-cux1ee9ng}

Phần này trình bày quá trình phân tích, lựa chọn và thiết kế các thành
phần phần cứng cho thiết bị theo dõi xe IoT. Mục tiêu không chỉ là tạo
ra một thiết bị theo dõi "đủ chức năng" trên lý thuyết, mà là một thiết bị có thể
lắp lên xe, vận hành ổn định và còn chừa dư địa cho bảo trì, mở rộng.
Vì vậy, mỗi quyết định kỹ thuật đều được soi chiếu dưới bốn góc nhìn:
độ tin cậy, khả năng hoạt động liên tục với nguồn dự phòng, mức độ tích
hợp và mức độ phù hợp với ngân sách.

\textbf{Quy ước thuật ngữ trong phần cứng:} tài liệu dùng cụm ``khối phần
cứng/khối chức năng'' cho các thành phần tích hợp trên PCB; từ ``module''
chỉ dùng khi nói về module phần mềm, như firmware trên thiết bị hoặc
dịch vụ phía máy chủ, hoặc khi trích dẫn nguyên văn tên sản phẩm của nhà
sản xuất.

\paragraph{3.1.1.1. Sơ đồ khối hệ
thống}\label{3111-sux1a1-ux111ux1ed3-khux1ed1i-hux1ec7-thux1ed1ng}

Hệ thống thiết bị theo dõi được tổ chức theo kiến trúc khối chức năng với năm khối
chính, trong đó ESP32-S3 đóng vai trò điều phối trung tâm. Nếu nhìn dưới
góc độ người sử dụng, năm khối này tương ứng với năm câu hỏi rất thực tế:
thiết bị lấy dữ liệu xe ở đâu, xác định vị trí bằng cách nào, phát hiện
rung/chuyển động ra sao, gửi dữ liệu về máy chủ như thế nào và lấy điện từ
đâu để duy trì hoạt động. Sơ đồ khối tổng thể được trình bày như sau:


{\thesissetfigureid{3.1}{fig-ch3-so-o-khoi-tong-the-he-thong-tracker}
\begin{figure}[htbp]
\centering
\pandocbounded{\safeincludesvg{./assets/figures/03-chuong-3-giai-phap-phan-cung-hinh-3-1.svg}}
\caption{Sơ đồ khối tổng thể hệ thống thiết bị theo dõi}
\label{fig:ch3-so-o-khoi-tong-the-he-thong-tracker}
\end{figure}
}



{\thesissetfigureid{3.1a}{fig-ch3-phan-ra-chi-tiet-cac-khoi-chuc-nang-cua-tracker}
\begin{figure}[htbp]
\centering
\pandocbounded{\safeincludesvg{./assets/figures/03-chuong-3-giai-phap-phan-cung-hinh-3-1a.svg}}
\caption{Phân rã chi tiết các khối chức năng của thiết bị theo dõi}
\label{fig:ch3-phan-ra-chi-tiet-cac-khoi-chuc-nang-cua-tracker}
\end{figure}
}


Sơ đồ khối cho thấy thiết bị theo dõi thực chất đang làm bốn việc liên tiếp: đọc
trạng thái của xe, theo dõi chuyển động, gửi dữ liệu ra ngoài và tự bảo
vệ nguồn điện. Trong đó, OBD2 là cổng chẩn đoán của xe, cho biết xe đang
vận hành ra sao; khối LTE/GNSS phụ trách liên lạc qua mạng di động và
định vị bằng vệ tinh; IMU là cảm biến quán tính giúp hệ thống vẫn ``canh
xe'' khi động cơ đã tắt; còn khối nguồn đảm bảo thiết bị đổi giữa nguồn
chính và pin dự phòng một cách an toàn. Vi điều khiển đứng ở giữa để gom
mọi tín hiệu này và quyết định thiết bị phải ở trạng thái nào.

Hệ thống vận hành theo ba chế độ chính: (1) chế độ lái xe --- khi động cơ
bật (IGN ON), các khối chức năng cần thiết được kích hoạt; (2) chế độ đỗ
xe --- khi động cơ tắt (IGN OFF), ESP32 chuyển sang deep sleep, tức trạng
thái ngủ sâu để giảm tiêu thụ điện, và chỉ IMU LIS3DSH duy trì giám sát
chuyển động; (3) chế độ cảnh báo --- khi IMU ghi nhận chuyển động bất
thường, hệ thống tự đánh thức và gửi cảnh báo qua 4G.

Luồng dữ liệu vì thế cũng rất dễ hình dung. Dữ liệu vận hành của xe
(như RPM, tốc độ, nhiên liệu) được lấy qua BLE, tức Bluetooth tiết kiệm
năng lượng, từ adapter OBD2; vị trí được lấy từ GNSS tích hợp trong
modem; còn rung động được đọc từ IMU qua I2C, tức đường giao tiếp ngắn
thường dùng cho cảm biến. ESP32-S3 tổng hợp ba nguồn này, đóng gói thành
bản tin rồi gửi lên máy chủ bằng MQTT qua kết nối 4G/LTE của modem.

Để tăng độ liên tục dữ liệu khi mạng gián đoạn, luồng dữ liệu vận hành
gửi liên tục được bổ sung thêm nhánh lưu đệm cục
bộ trên microSD. Thẻ nhớ này được nối bằng SDMMC 4-bit, tức chuẩn giao
tiếp nhanh cho thẻ nhớ, và được quản lý bằng FATFS, tức hệ thống tệp để
đọc ghi ổn định. Khi mất mạng, dữ liệu sẽ được giữ lại theo đúng thứ tự;
khi có mạng trở lại, thiết bị gửi bù phần còn thiếu lên MQTT.

Sau khi chốt được bức tranh tổng thể và đường đi của dữ liệu, bước kế
tiếp là chọn phần tử điều phối trung tâm đủ sức gánh toàn bộ các khối
đó. Vì vậy, phần sau tập trung vào việc lựa chọn vi điều khiển.

\paragraph{3.1.1.2. Phân tích và lựa chọn vi điều khiển
(MCU)}\label{3112-phuxe2n-tuxedch-vuxe0-lux1ef1a-chux1ecdn-vi-ux111iux1ec1u-khiux1ec3n-mcu}

Việc lựa chọn vi điều khiển (MCU) là quyết định thiết kế quan trọng
nhất, ảnh hưởng trực tiếp đến khả năng tích hợp các khối ngoại vi, mức
tiêu thụ năng lượng, và độ phức tạp phát triển firmware. Ba ứng viên
được đưa vào quá trình đánh giá là ESP32-S3 của Espressif, STM32L4 của
STMicroelectronics, và nRF52840 của Nordic Semiconductor.

\subparagraph{a) Các ứng viên}\label{a-cuxe1c-ux1ee9ng-viuxean}

\textbf{ESP32-S3} là vi điều khiển hai nhân dùng kiến trúc Xtensa LX7,
hoạt động tới 240 MHz và tích hợp sẵn WiFi cùng Bluetooth LE 5.0
{[}19{]}. Điểm quan trọng với đồ án không chỉ là xung nhịp, mà là chip
này có đủ tài nguyên để vừa xử lý logic điều khiển, vừa giao tiếp với
nhiều phần cứng khác nhau. Thiết bị có 512 KB SRAM trên chip, dùng flash
ngoài, hỗ trợ nhiều cổng giao tiếp như UART, I2C, SPI và ADC, tức các
đường nối cơ bản để làm việc với modem, cảm biến và mạch đo điện áp,
đồng thời có chế độ ngủ sâu rất tiết kiệm điện {[}19{]}, {[}20{]}.

\textbf{STM32L4} (đại diện bởi STM32L476 trong so sánh) là dòng vi điều
khiển tiết kiệm điện của ST, dùng nhân ARM Cortex-M4 tới 80 MHz và
không tích hợp sẵn BLE hay Wi-Fi {[}53{]}, {[}54{]}. Điểm mạnh của dòng
này là các chế độ ngủ sâu có mức tiêu thụ rất thấp theo từng cấu hình
nguồn {[}53{]}.

\textbf{nRF52840} là vi điều khiển của Nordic với nhân ARM Cortex-M4F 64
MHz, hỗ trợ mạnh về Bluetooth LE 5.x và 802.15.4 {[}55{]}, {[}56{]}.
Thiết bị có 1 MB Flash, 256 KB RAM, hỗ trợ cổng nối tiếp UART/UARTE và
cũng có chế độ tắt sâu với dòng tiêu thụ rất thấp {[}55{]}, {[}56{]}.

\subparagraph{b) Đối chiếu yêu cầu thiết kế với thông số phần
cứng}\label{b-ux111ux1ed1i-chiux1ebfu-yuxeau-cux1ea7u-thiux1ebft-kux1ebf-vux1edbi-thuxf4ng-sux1ed1-phux1ea7n-cux1ee9ng}

Thay vì chấm điểm tổng quát, \tabref{tab:ch3-oi-chieu-yeu-cau-ky-thuat-khi-lua-chon-vi-ieu-khien} đối chiếu trực tiếp các ràng buộc
kỹ thuật của thiết bị theo dõi với thông số và tác động tích hợp của từng MCU.


{\thesissettableid{3.1}{tab-ch3-oi-chieu-yeu-cau-ky-thuat-khi-lua-chon-vi-ieu-khien}
\def\LTcaptype{table}
\begin{longtable}{|L{0.153\linewidth}|L{0.247\linewidth}|L{0.299\linewidth}|L{0.255\linewidth}|}
\caption{Đối chiếu yêu cầu kỹ thuật khi lựa chọn vi điều khiển}\label{tab:ch3-oi-chieu-yeu-cau-ky-thuat-khi-lua-chon-vi-ieu-khien}\\
\hline \textbf{Yêu cầu của thiết bị theo dõi} & \textbf{ESP32-S3} & \textbf{STM32L4} & \textbf{nRF52840} \\
\\\hline
\endfirsthead
\caption[]{Đối chiếu yêu cầu kỹ thuật khi lựa chọn vi điều khiển (tiếp theo)}\\
\hline \textbf{Yêu cầu của thiết bị theo dõi} & \textbf{ESP32-S3} & \textbf{STM32L4} & \textbf{nRF52840} \\
\\\hline
\endhead
\\\hline
\endfoot
\endlastfoot
Kết nối OBD2 BLE với vgate iCar Pro & BLE 5.0 tích hợp, tương thích
ngược BLE 4.0 & Không tích hợp BLE, phải thêm module ngoài & BLE 5.0
tích hợp \\\hline
Cấu hình UART của hệ thống & 3 UART controllers (đủ cho modem + debug,
còn dư cho mở rộng) {[}19{]}, {[}20{]} & Nhiều serial instance trong
dòng STM32L476 (USART, UART, LPUART), đáp ứng yêu cầu số cổng {[}53{]},
{[}54{]} & UARTE hoặc UART thường dùng 2 instance, ít dư địa nếu cần giữ
debug UART riêng {[}55{]}, {[}56{]} \\\hline
Tài nguyên xử lý cho OBD2 + modem + IMU & Dual-core 240 MHz, 512 KB SRAM
{[}19{]} & Cortex-M4 tới 80 MHz, 128 KB SRAM trên STM32L476RE {[}54{]} &
Cortex-M4F 64 MHz, 256 KB RAM {[}55{]}, {[}56{]} \\\hline
Dòng ngủ sâu/siêu thấp công suất & Deep-sleep mức µA tùy cấu hình RTC/IO
{[}19{]}, {[}20{]} & STOP2 cỡ 1.1 µA (theo điều kiện datasheet),
standby/shutdown cỡ nA {[}53{]}, {[}54{]} & System OFF mức dưới µA tùy
RAM retention/GPIO {[}55{]} \\\hline
Framework triển khai & Arduino IDE + ESP-IDF {[}47{]} & STM32Cube/HAL
{[}53{]} & nRF Connect SDK {[}57{]} \\\hline
Phần cứng bổ sung để đạt đủ chức năng & Không cần thêm radio BLE & Cần
thêm BLE module nếu vẫn giữ OBD2 BLE (STM32L476 không tích hợp BLE)
{[}53{]}, {[}54{]} & Không cần BLE module, nhưng cần xử lý bài toán
thiếu serial cho debug \\\hline
Chi phí triển khai phần cứng tham chiếu &
~100.000--200.000 VND (mặt bằng BOM dự án) &
~150.000--300.000 VND (mặt bằng BOM dự án) & Thường cao
hơn ESP32-S3 và ít phổ biến hơn trong bối cảnh dự án \\\hline
Tác động tới sơ đồ hiện tại & Giữ nguyên sơ đồ BLE + modem UART, thuận
lợi cho tích hợp GNSS trong modem & Tăng phần cứng BLE ngoài và công
đoạn tích hợp ban đầu & Có thể phải chuyển debug sang USB/SWO để nhường
UART cho chức năng chính \\\hline
\end{longtable}
}


\subparagraph{c) Phân tích các yếu tố quan
trọng}\label{c-phuxe2n-tuxedch-cuxe1c-yux1ebfu-tux1ed1-quan-trux1ecdng}

\textbf{BLE tích hợp:} Đây là ràng buộc rất thực tế vì thiết bị theo dõi dùng
adapter OBD2 vgate iCar Pro theo chuẩn BLE 4.0. Nếu vi điều khiển không
có BLE tích hợp, thiết kế sẽ phải bổ sung thêm một chip hoặc module
Bluetooth riêng. Điều đó không chỉ làm tăng BOM, mà còn khiến sơ đồ
mạch, nguồn cấp và phần mềm điều khiển phức tạp hơn ngay từ đầu. Ở khía
cạnh này, ESP32-S3 và nRF52840 thuận lợi hơn STM32L4.

\textbf{Số cổng UART khả dụng:} UART ở đây là cổng nối tiếp cơ bản để vi
điều khiển giao tiếp với modem hoặc mở cổng debug. Trong kiến trúc hiện
tại, UART không phải tài nguyên dư dả. Hệ thống cần ít nhất một kênh
cho modem và một kênh cho debug, đồng thời vẫn nên chừa sẵn một cổng cho mở
rộng hoặc chẩn đoán khi triển khai thực tế. ESP32-S3 đáp ứng tốt nhu cầu
này với 3 UART phần cứng; nRF52840 hạn chế hơn; còn STM32L4 có thể đáp
ứng nhưng phải đánh đổi bằng việc bổ sung BLE ngoài.

\textbf{Tiêu thụ deep sleep:} Deep sleep là chế độ ngủ sâu, nơi vi điều
khiển gần như tắt hết hoạt động để tiết kiệm điện tối đa. STM32L4 và
nRF52840 có lợi thế rõ ràng về dòng ngủ sâu theo tài liệu hãng. Tuy
nhiên, trong một thiết bị theo dõi hoàn chỉnh, vi điều khiển không phải thành phần
duy nhất tiêu thụ năng lượng. Modem LTE, khối GNSS, chu kỳ đánh thức và
các linh kiện phụ trợ cũng góp phần đáng kể vào mức tiêu thụ điện của toàn hệ thống. Vì
vậy, chênh lệch ở mức riêng MCU có ý nghĩa, nhưng không đủ để bù cho việc
tăng độ phức tạp tích hợp {[}19{]}, {[}53{]}, {[}54{]}, {[}55{]}.

\textbf{Độ phức tạp triển khai:} Xét theo công sức hiện thực, ESP32-S3
thuận lợi hơn vì đã có BLE tích hợp sẵn và hệ sinh thái phát triển quen
thuộc với nhóm làm đồ án. Điều này đặc biệt quan trọng khi firmware phải
đồng thời xử lý OBD2 qua BLE, gửi lệnh AT cho modem, đọc IMU qua I2C và
duy trì logic ngủ sâu. STM32L4 tiết kiệm điện tốt hơn nhưng kéo theo
thêm việc tích hợp BLE; còn nRF52840 mạnh về BLE nhưng ít dư địa UART
hơn cho kịch bản thiết bị theo dõi hiện tại.

\subparagraph{d) Kết luận lựa
chọn}\label{d-kux1ebft-luux1eadn-lux1ef1a-chux1ecdn}

Trên cơ sở đối chiếu trực tiếp các yêu cầu của hệ thống,
\textbf{ESP32-S3} được chọn làm vi điều khiển chính cho thiết bị theo dõi. Nói ngắn
gọn, đây là phương án cân bằng tốt nhất giữa "đủ mạnh để làm được việc"
và "đủ gọn để triển khai thực tế". Quyết định này dựa trên các điểm chốt
sau:

\begin{itemize}
\tightlist
\item
  Đáp ứng đầy đủ kiến trúc hiện tại: 1 kết nối BLE cho OBD2, UART cho
  modem và một đường debug riêng, đồng thời vẫn còn dư địa UART cho mở
  rá»™ng.
\item
  BLE 5.0 tích hợp giúp bỏ hẳn module BLE ngoài, giảm BOM và rút ngắn
  công đoạn tích hợp.
\item
  Tài nguyên xử lý 240 MHz dual-core và 512 KB SRAM đủ để chạy đồng thời
  BLE, modem, IMU và máy trạng thái.
\item
  Chi phí triển khai phần cứng với ESP32-S3 thấp hơn STM32L4 khoảng
  30--50\% theo mặt bằng linh kiện đang dùng trong đồ án.
\item
  Dòng deep-sleep của ESP32-S3 cao hơn STM32L4 theo tài liệu hãng, nhưng
  vẫn nằm trong giới hạn chấp nhận được của bài toán pin dự phòng khi xét
  toàn bộ duty-cycle của modem, GNSS và các chu kỳ wake-up {[}19{]},
  {[}53{]}, {[}54{]}.
\end{itemize}


{\thesissettableid{3.2}{tab-ch3-thong-so-ky-thuat-esp32-s3-uoc-chon}
\def\LTcaptype{table}
\begin{longtable}{|L{0.274\linewidth}|L{0.694\linewidth}|}
\caption{Thông số kỹ thuật ESP32-S3 được chọn}\label{tab:ch3-thong-so-ky-thuat-esp32-s3-uoc-chon}\\
\hline \textbf{Thông số} & \textbf{Giá trị} \\
\\\hline
\endfirsthead
\caption[]{Thông số kỹ thuật ESP32-S3 được chọn (tiếp theo)}\\
\hline \textbf{Thông số} & \textbf{Giá trị} \\
\\\hline
\endhead
\\\hline
\endfoot
\endlastfoot
CPU & Dual-core Xtensa LX7, tối đa 240 MHz {[}19{]} \\\hline
RAM on-chip & 512 KB SRAM {[}19{]} \\\hline
Flash & Flash ngoài (biến thể WROOM thường 4--16 MB) {[}19{]},
{[}20{]} \\\hline
Bluetooth & Bluetooth LE 5.0, tương thích ngược BLE 4.0 cho OBD2 adapter
{[}19{]} \\\hline
WiFi & 802.11 b/g/n (không dùng trong luồng chính của thiết bị theo dõi)
{[}19{]} \\\hline
UART & 3 UART controllers {[}19{]}, {[}20{]} \\\hline
I2C & 2 I2C controllers {[}19{]}, {[}20{]} \\\hline
ADC & SAR ADC 12-bit, tối đa 20 kênh (theo package) {[}19{]} \\\hline
GPIO & Tối đa 45 GPIO (theo package) {[}19{]} \\\hline
Dòng deep-sleep & Mức µA tùy cấu hình RTC domain và IO {[}19{]},
{[}20{]} \\\hline
Biến thể linh kiện & ESP32-S3-WROOM-1 \\\hline
Phương án tích hợp & Tích hợp trực tiếp vào bo mạch chính \\\hline
\end{longtable}
}


\paragraph{3.1.1.3. Phân tích và lựa chọn modem truyền
thông}\label{3113-phuxe2n-tuxedch-vuxe0-lux1ef1a-chux1ecdn-modem-truyux1ec1n-thuxf4ng}

Modem truyền thông đóng vai trò then chốt trong hệ thống vì nó đồng thời
giải quyết hai nhu cầu thiết yếu nhất của thiết bị theo dõi: gửi dữ liệu về máy chủ
qua mạng di động và cung cấp vị trí của xe thông qua GNSS. Với bài toán
này, lựa chọn modem không chỉ ảnh hưởng đến tốc độ hay thông số kỹ
thuật, mà còn tác động trực tiếp đến cách đi dây, cách tổ chức firmware
và mức độ phức tạp của toàn bộ bo mạch. Trong kiến trúc hiện tại, hệ
thống sử dụng \textbf{SIM7600CE-T tích hợp LTE + GNSS} nhằm giảm số
lượng phần cứng rời, đơn giản hóa đi dây và đồng bộ với firmware thực
tế.

\subparagraph{a) Các ứng viên}\label{a-cuxe1c-ux1ee9ng-viuxean-1}

Ba phương án được đưa vào đánh giá gồm hai phương án tách rời
(\textbf{SIMCom A7670C + u-blox NEO-M8N}, \textbf{Quectel EC200U-CN +
NEO-M8N}) và một phương án tích hợp (\textbf{SIMCom SIM7600CE-T}).

\textbf{SIMCom A7670C + u-blox NEO-M8N} là phương án tách rời: A7670C
phụ trách LTE/AT command, NEO-M8N phụ trách GNSS/NMEA. Theo trang sản
phẩm chính thức của SIMCom, A7670C nằm trong nhóm LTE Cat-1 và dùng
nguồn 3.4--4.2 V {[}58{]}, {[}62{]}.

\textbf{Quectel EC200U-CN + NEO-M8N} có cấu trúc gần giống phương án
A7670C + NEO-M8N. Tuy nhiên dòng EC200U được hãng mô tả là Cat-1 bis,
GNSS là tính năng tùy chọn theo biến thể/cấu hình và I/O logic có ràng
buộc riêng, nên mức độ phù hợp với cấu hình đang triển khai thấp hơn
{[}59{]}, {[}61{]}.

\textbf{SIMCom SIM7600CE-T} là phương án tích hợp GNSS ngay trong modem,
tức cùng một khối phần cứng vừa lo kết nối mạng di động vừa lo lấy vị trí
từ vệ tinh. Cách làm này giúp đi dây gọn hơn và giảm số lượng phần cứng
rời {[}60{]}.

\subparagraph{b) Đối chiếu kiến trúc LTE/GNSS theo ràng buộc tích
hợp}\label{b-ux111ux1ed1i-chiux1ebfu-kiux1ebfn-truxfac-ltegnss-theo-ruxe0ng-buux1ed9c-tuxedch-hux1ee3p}

Vấn đề cốt lõi của thiết bị theo dõi không chỉ là modem có lên mạng được hay
không, mà còn là phương án LTE/GNSS, tức phương án kết hợp giữa kết nối
di động và định vị vệ tinh, ảnh hưởng thế nào tới sơ đồ UART, hồ sơ tiêu
thụ nguồn và mức độ thuận lợi khi phát triển firmware.
\tabref{tab:ch3-oi-chieu-cac-phuong-an-lte-gnss-theo-tac-ong-tich-hop}
dưới đây đối chiếu trực tiếp các điểm khác biệt này.


{\thesissettableid{3.3}{tab-ch3-oi-chieu-cac-phuong-an-lte-gnss-theo-tac-ong-tich-hop}
\def\LTcaptype{table}
\begin{longtable}{|L{0.202\linewidth}|L{0.260\linewidth}|L{0.252\linewidth}|L{0.241\linewidth}|}
\caption{Đối chiếu các phương án LTE/GNSS theo tác động tích hợp}\label{tab:ch3-oi-chieu-cac-phuong-an-lte-gnss-theo-tac-ong-tich-hop}\\
\hline \textbf{Yêu cầu tích hợp} & \textbf{A7670C + NEO-M8N} & \textbf{EC200U-CN + NEO-M8N} & \textbf{SIM7600CE-T} \\
\\\hline
\endfirsthead
\caption[]{Đối chiếu các phương án LTE/GNSS theo tác động tích hợp (tiếp theo)}\\
\hline \textbf{Yêu cầu tích hợp} & \textbf{A7670C + NEO-M8N} & \textbf{EC200U-CN + NEO-M8N} & \textbf{SIM7600CE-T} \\
\\\hline
\endhead
\\\hline
\endfoot
\endlastfoot
Số khối phần cứng rời & 2 khối rời: A7670C (LTE) + NEO-M8N (GNSS)
{[}58{]}, {[}62{]} & 2 khối rời: EC200U-CN (LTE) + NEO-M8N (GNSS)
{[}59{]}, {[}61{]} & 1 khối tích hợp LTE + GNSS (SIM7600CE-T) {[}60{]} \\\hline
LTE category / tốc độ & Cat-1, tối đa 10 Mbps DL / 5 Mbps UL {[}58{]} &
Cat-1, tối đa 10 Mbps DL / 5 Mbps UL {[}59{]}, {[}61{]} & Cat-4, tối đa
150 Mbps DL / 50 Mbps UL {[}60{]} \\\hline
GNSS tích hợp trong modem & Không tích hợp GNSS trên A7670C, cần GNSS
ngoài {[}58{]}, {[}62{]} & GNSS tùy chọn theo variant/cấu hình dòng
EC200U {[}59{]}, {[}61{]} & Có GNSS tích hợp trong module {[}60{]} \\\hline
Điện áp cấp nguồn modem & 3.4--4.2 V (typ. 3.8 V) {[}58{]} & 3.3--4.3 V
(typ. 3.8 V) {[}61{]} & 3.4--4.2 V (typ. 3.8 V) {[}60{]} \\\hline
Đường dữ liệu vị trí trong kiến trúc đề tài & NMEA/UBX từ NEO-M8N qua
UART riêng {[}62{]}, {[}63{]} & NMEA/UBX từ NEO-M8N qua UART riêng
{[}62{]}, {[}63{]} & GNSS AT tích hợp trong cùng modem {[}60{]} \\\hline
Bật/tắt LTE và GNSS độc lập & Có, do tách 2 module/2 rail nguồn & Có,
nếu triển khai tách tương tự A7670C+NEO-M8N & Không tách hoàn toàn do
GNSS đi cùng modem \\\hline
Tác động khi reset modem LTE & Có thể reset modem và giữ logic GNSS
riêng & Tương tự phương án tách nếu GNSS dùng module riêng & Reset modem
thường kéo theo gián đoạn GNSS \\\hline
Mức phụ thuộc giữa các chức năng firmware & Thấp, tách rõ
\texttt{modem\_lte} và \texttt{gnss} & Thấp, nhưng phải đổi tập lệnh
modem & Cao hơn do LTE/GNSS chung khối tích hợp \\\hline
Tác động tới BOM phần cứng thực tế & Bộ modem + GNSS đang chốt
330.000--500.000 VND (BOM dự án) & Chưa phải cấu hình BOM mục tiêu của phương án đề xuất
& Giảm số khối phần cứng rời, thuận lợi cho đi dây và lắp ráp \\\hline
\end{longtable}
}


\subparagraph{c) Kết luận lựa
chọn}\label{c-kux1ebft-luux1eadn-lux1ef1a-chux1ecdn}

\textbf{SIMCom SIM7600CE-T} được lựa chọn làm kiến trúc truyền thông
chính vì đây là phương án giúp thiết kế tổng thể gọn hơn và dễ triển
khai hơn, dù không phải lúc nào cũng linh hoạt nhất ở từng thành phần
riêng lẻ. Các lý do cụ thể gồm:

\begin{itemize}
\tightlist
\item
  Tích hợp LTE và GNSS trong cùng một module giúp giảm số lượng phần
  cứng rời và rút gọn đường kết nối UART.
\item
  Giữ nguyên sơ đồ chân đang triển khai (UART1 trên GPIO16/17, PWR-KEY
  trên GPIO26), không cần thay đổi pin mapping phần cứng.
\item
  Luồng AT được chuẩn hóa với Auto mode (\texttt{AT+CNMP=2}), kiểm tra
  \texttt{AT+CEREG?} trước \texttt{AT+CGACT=1,1}, và APN mặc định
  \texttt{internet}, phù hợp với firmware hiện tại.
\item
  GNSS tích hợp được điều khiển trực tiếp bằng \texttt{AT+CGNSPWR} và
  đọc dữ liệu qua \texttt{AT+CGNSINF}, thống nhất với module firmware
  đang vận hành.
\item
  Phù hợp với phương án BOM phần cứng đang áp dụng cho bo mạch đã chế
  tạo, thuận lợi cho triển khai và đối chiếu tài liệu kỹ thuật.
\end{itemize}


{\thesissettableid{3.4}{tab-ch3-thong-so-ky-thuat-phuong-an-sim7600ce-t}
\def\LTcaptype{table}
\begin{longtable}{|L{0.220\linewidth}|L{0.740\linewidth}|}
\caption{Thông số kỹ thuật phương án SIM7600CE-T}\label{tab:ch3-thong-so-ky-thuat-phuong-an-sim7600ce-t}\\
\hline \textbf{Thông số} & \textbf{Giá trị} \\
\\\hline
\endfirsthead
\caption[]{Thông số kỹ thuật phương án SIM7600CE-T (tiếp theo)}\\
\hline \textbf{Thông số} & \textbf{Giá trị} \\
\\\hline
\endhead
\\\hline
\endfoot
\endlastfoot
Kiến trúc & SIM7600CE-T tích hợp LTE Cat-4 và GNSS trong cùng module
{[}60{]} \\\hline
LTE category / tốc độ & Cat-4, tối đa 150 Mbps downlink / 50 Mbps uplink
{[}60{]} \\\hline
GNSS & GNSS tích hợp (GPS/GLONASS/BeiDou/Galileo), điều khiển bằng
\texttt{AT+CGNSPWR}, đọc dữ liệu bằng \texttt{AT+CGNSINF} {[}60{]} \\\hline
GNSS tích hợp trong modem & Có {[}60{]} \\\hline
Giao tiếp với MCU & 1 UART (GPIO16/17) cho toàn bộ AT LTE/GNSS + GPIO26
điều khiển PWR-KEY {[}60{]} \\\hline
Giao thức dữ liệu & LTE dùng AT command cho PDP/MQTT; GNSS lấy dữ liệu
qua \texttt{AT+CGNSINF} và có thể stream NMEA bằng \texttt{AT+CGNSTST}
khi cần {[}60{]} \\\hline
Điện áp modem & 3.4--4.2 V, điển hình 3.8 V {[}60{]} \\\hline
Luồng attach LTE & \texttt{AT+CNMP=2} (auto mode), cấu hình
\texttt{AT+CGDCONT=1,"IP","internet"}, kiểm tra \texttt{AT+CEREG?} trước
\texttt{AT+CGACT=1,1} \\\hline
Đồng bộ GNSS & Bật/tắt GNSS bằng \texttt{AT+CGNSPWR=1/0}, đọc dữ liệu
bằng \texttt{AT+CGNSINF} \\\hline
Điện áp cấp nguồn & 3.4--4.2 V, điển hình 3.8 V {[}60{]} \\\hline
Ý nghĩa tích hợp & Giảm số linh kiện rời, đơn giản hóa đi dây và đồng bộ
với firmware thực tế \\\hline
Ghi chú & APN mặc định sử dụng trong hệ thống là \texttt{internet}; pin
mapping phần cứng giữ nguyên \\\hline
\end{longtable}
}

Đến đây, hai khối điều phối và liên lạc của thiết bị theo dõi đã được chốt. Phần
còn lại của thiết kế phần cứng là trả lời ba câu hỏi thực dụng hơn: lấy
dữ liệu xe bằng cách nào, canh chuyển động khi xe đỗ ra sao và nuôi
toàn bộ hệ thống bằng nguồn nào. Vì vậy, mục 3.1.2 chuyển sang cách ghép
các lựa chọn đó thành phần cứng hoàn chỉnh.


\subsubsection{3.1.2. Giải pháp phần
cứng}\label{321-giux1ea3i-phuxe1p-phux1ea7n-cux1ee9ng}

Trên cơ sở các lựa chọn phần cứng đã phân tích ở mục 3.1.1, phần này đi
vào cách nhóm hiện thực chúng thành kiến trúc thiết bị theo dõi hoàn chỉnh, từ
kết nối OBD2 và IMU đến khối nguồn và BOM tham chiếu.

\paragraph{3.1.2.1. Thiết kế khối thu thập dữ liệu
OBD2}\label{3211-thiux1ebft-kux1ebf-muxf4-ux111un-thu-thux1eadp-dux1eef-liux1ec7u-obd2}

\subparagraph{a) Lựa chọn phương pháp kết nối
OBD2}\label{a-lux1ef1a-chux1ecdn-phux1b0ux1a1ng-phuxe1p-kux1ebft-nux1ed1i-obd2}

Thu thập dữ liệu từ ECU, tức hộp điều khiển điện tử của xe, qua cổng OBD2
(On-Board Diagnostics II) là một yêu cầu cốt lõi của thiết kế thiết bị theo dõi vì
đây là nguồn thông tin giúp hệ thống đi xa hơn một thiết bị định vị
thông thường. Có hai hướng tiếp cận chính: sử dụng adapter OBD2 có dây,
tức bộ chuyển đổi kiểu ELM327 nối qua UART; hoặc adapter OBD2 không dây,
tức trao đổi qua BLE. Nhóm lựa chọn kết nối không dây qua BLE vì cách
làm này phù hợp hơn với cách bố trí thiết bị thực tế trên xe và với tài
nguyên phần cứng của thiết bị theo dõi:


{\thesissettableid{3.5}{tab-ch3-so-sanh-phuong-phap-ket-noi-obd2-theo-thong-so-trien-khai}
\def\LTcaptype{table}
\begin{longtable}{|L{0.177\linewidth}|L{0.303\linewidth}|L{0.484\linewidth}|}
\caption{So sánh phương pháp kết nối OBD2 theo thông số triển khai}\label{tab:ch3-so-sanh-phuong-phap-ket-noi-obd2-theo-thong-so-trien-khai}\\
\hline \textbf{Tiêu chí kỹ thuật} & \textbf{OBD2 có dây (ELM327 UART)} & \textbf{OBD2 không dây (BLE
ELM327 / vgate iCar Pro)} \\
\\\hline
\endfirsthead
\caption[]{So sánh phương pháp kết nối OBD2 theo thông số triển khai (tiếp theo)}\\
\hline \textbf{Tiêu chí kỹ thuật} & \textbf{OBD2 có dây (ELM327 UART)} & \textbf{OBD2 không dây (BLE
ELM327 / vgate iCar Pro)} \\
\\\hline
\endhead
\\\hline
\endfoot
\endlastfoot
Kết nối vật lý thiết bị theo dõi ↔ OBD2 & Cần dây kéo cố định từ cổng OBD2 đến
thiết bị theo dõi & Không cần dây giữa thiết bị theo dõi và cổng OBD2 \\\hline
Giao thức tới MCU & UART TTL & BLE 4.0 GATT; vgate iCar Pro là bản
Bluetooth 4.0 BLE {[}29{]}, {[}64{]} \\\hline
Tài nguyên MCU tiêu tốn & Chiếm thêm 1 UART phần cứng trong lúc kết nối
& Tận dụng BLE tích hợp của ESP32-S3, giữ UART cho modem/GNSS/debug \\\hline
Giao thức OBD-II hỗ trợ & Phụ thuộc adapter ELM327 cụ thể & Vgate công
bố hỗ trợ J1850 PWM/VPW, ISO9141-2, ISO14230-4, ISO15765-4 CAN, SAE
J1939 CAN {[}64{]} \\\hline
Hành vi sleep của adapter & Phụ thuộc adapter có dây cụ thể & Vendor
công bố tự sleep sau khoảng 30 phút khi xe tắt máy {[}64{]} \\\hline
Hành vi wake & Phụ thuộc adapter và đường cấp nguồn & Vendor quảng bá tự
khởi động khi ignition ON; cần hiểu là hành vi phụ thuộc xe {[}64{]} \\\hline
Dòng tiêu thụ adapter & Phụ thuộc adapter; không dùng 1 số cứng nếu
không có datasheet đúng mẫu & Trang hãng được kiểm tra không công bố
dòng active/standby, nên không chốt số mA trong luận văn {[}64{]} \\\hline
Hành vi khi MCU deep sleep & UART không duy trì truyền dữ liệu khi MCU
ngủ sâu & BLE link cũng không duy trì qua deep sleep; cần kết nối lại sau
khi đánh thức \\\hline
Tác động lắp đặt thực địa & Tracker thường bị ràng buộc gần cổng OBD2 &
Tracker có thể đặt lệch vị trí OBD2 trong cabin nếu vẫn trong vùng
BLE \\\hline
Bảo trì/thay adapter giữa các xe & Cần tháo dây liên quan đến thiết bị theo dõi &
Chỉ cần rút/cắm adapter và cấu hình lại kết nối nếu cần \\\hline
\end{longtable}
}


\textbf{Kết luận:} Phương pháp BLE được chọn vì phù hợp hơn với kiến
trúc hiện tại của thiết bị theo dõi: tận dụng BLE tích hợp của ESP32-S3, không
chiếm thêm UART phần cứng, giảm ràng buộc đi dây, đồng thời vgate iCar
Pro đã có công bố rõ về BLE 4.0, auto-sleep và danh sách giao thức
OBD-II. Với các thông số hãng chưa công bố (như dòng active/standby),
báo cáo chỉ ghi nhận là chưa có số liệu chính thức thay vì đưa ra giá
trị ước lượng.

\subparagraph{b) Lựa chọn adapter OBD2 BLE: vgate iCar
Pro}\label{b-lux1ef1a-chux1ecdn-adapter-obd2-ble-vgate-icar-pro}

\textbf{vgate iCar Pro} là adapter OBD2 sử dụng Bluetooth Low Energy
(BLE) 4.0, tương thích với BLE 5.0 của ESP32-S3. Adapter này được chọn
dựa trên các tiêu chí sau:

\begin{itemize}
\tightlist
\item
  \textbf{Tương thích giao thức:} Hỗ trợ tập lệnh ELM327 qua BLE GATT
  characteristics, cho phép đọc đa dạng dữ liệu từ ECU bao gồm trạng
  thái động cơ (IGN), số vòng quay (RPM), tốc độ xe, mức nhiên liệu,
  nhiệt độ động cơ, và mã lỗi chẩn đoán (DTC) {[}7{]}.
\item
  \textbf{Khả năng kết nối:} BLE 4.0 của adapter tương thích ngược với
  BLE 5.0 của ESP32-S3 {[}19{]}, {[}64{]}. Vendor cũng công bố cơ chế tự
  sleep sau khoảng 30 phút khi xe tắt máy và tự khởi động lại theo
  ignition, phù hợp cho kịch bản thiết bị theo dõi trên xe {[}64{]}.
\item
  \textbf{Tiêu thụ năng lượng:} Ở thời điểm rà soát nguồn, trang hãng
  được kiểm tra không công bố dòng active/standby chính thức cho vgate
  iCar Pro BLE {[}64{]}. Vì vậy báo cáo chỉ ghi nhận đây là thông số
  chưa được công bố đầy đủ, không đưa ra giá trị mA cố định khi chưa có
  datasheet chính thức đúng mẫu.
\item
  \textbf{Phổ biến và giá hợp lý:} Khảo sát thị trường Việt Nam ngày
  12/04/2026 cho thấy các mức giá tham khảo của vgate iCar Pro BLE nằm
  quanh khoảng 367.000--609.000 VND tùy kênh bán lẻ và phiên bản.
\end{itemize}


{\thesissetfigureid{3.2}{fig-ch3-so-o-ket-noi-ble-giua-esp32-s3-va-vgate-icar-pro}
\begin{figure}[htbp]
\centering
\pandocbounded{\safeincludesvg{./assets/figures/03-chuong-3-giai-phap-phan-cung-hinh-3-2.svg}}
\caption{Sơ đồ kết nối BLE giữa ESP32-S3 và vgate iCar Pro}
\label{fig:ch3-so-o-ket-noi-ble-giua-esp32-s3-va-vgate-icar-pro}
\end{figure}
}


\subparagraph{c) Chiến lược kết nối theo chế độ hoạt
động}\label{c-chiux1ebfn-lux1b0ux1ee3c-kux1ebft-nux1ed1i-theo-chux1ebf-ux111ux1ed9-houx1ea1t-ux111ux1ed9ng}

Một ràng buộc thiết kế quan trọng là ESP32-S3 không thể duy trì kết nối
BLE trong trạng thái deep sleep. Vì vậy, kết nối OBD2 không được xem là
một kênh "luôn bật", mà chỉ hoạt động khi nó thực sự tạo ra giá trị dữ liệu
cho từng trạng thái vận hành của xe:


{\thesissettableid{3.6}{tab-ch3-chien-luoc-ket-noi-ble-obd2-theo-che-o-hoat-ong}
\def\LTcaptype{table}
\begin{longtable}{|L{0.182\linewidth}|L{0.203\linewidth}|L{0.580\linewidth}|}
\caption{Chiến lược kết nối BLE OBD2 theo chế độ hoạt động}\label{tab:ch3-chien-luoc-ket-noi-ble-obd2-theo-che-o-hoat-ong}\\
\hline \textbf{Chế độ} & \textbf{Cần BLE OBD2?} & \textbf{Lý do} \\
\\\hline
\endfirsthead
\caption[]{Chiến lược kết nối BLE OBD2 theo chế độ hoạt động (tiếp theo)}\\
\hline \textbf{Chế độ} & \textbf{Cần BLE OBD2?} & \textbf{Lý do} \\
\\\hline
\endhead
\\\hline
\endfoot
\endlastfoot
Lái xe (IGN ON) & Có & Đọc dữ liệu xe (RPM, tốc độ, nhiên liệu) liên
tục \\\hline
Đỗ xe (IGN OFF) & Không & Không cần dữ liệu OBD2, chỉ cần IMU phát hiện
chuyển động \\\hline
Cảnh báo & Tùy chọn & Có thể kết nối để xác nhận IGN, hoặc chỉ dùng IMU
+ GPS \\\hline
\end{longtable}
}


Khi xe chạy (IGN ON), ESP32-S3 duy trì kết nối BLE liên tục với adapter để
đọc dữ liệu OBD2 định kỳ mỗi 5--30 giây. Khi xe đỗ (IGN OFF), kết nối BLE
được ngắt trước khi ESP32 vào deep sleep, vì ở thời điểm này mục tiêu ưu
tiên không còn là lấy dữ liệu động cơ, mà là giữ mức tiêu thụ thấp nhất
trong khi vẫn duy trì khả năng phát hiện xâm nhập. IMU LIS3DSH đảm nhiệm
việc giám sát đó độc lập và chỉ đánh thức hệ thống khi có tín hiệu đáng
ngờ. Cách tổ chức này giúp giảm đáng kể năng lượng tiêu thụ ở pha đỗ xe mà
không làm mất đi chức năng cảnh báo cốt lõi.

Trong trường hợp adapter OBD2 không kết nối được (timeout 10 giây, retry
2--3 lần), hệ thống tự động chuyển sang phương pháp dự phòng (fallback)
là đo điện áp ắc quy qua ADC để phát hiện trạng thái IGN. Phương pháp
này kém chính xác hơn nhưng đảm bảo hệ thống vẫn hoạt động bình thường.

Như vậy, OBD2 giải quyết tốt bài toán ``xe đang chạy ra sao'' nhưng
chưa đủ cho trạng thái đỗ xe, khi mục tiêu chính là phát hiện tác động
bất thường với mức tiêu thụ điện rất thấp. Vì vậy, phần tiếp theo chuyển
sang khối IMU.

\paragraph{3.1.2.2. Thiết kế khối cảm biến chuyển động
IMU}\label{3212-thiux1ebft-kux1ebf-muxf4-ux111un-cux1ea3m-biux1ebfn-chuyux1ec3n-ux111ux1ed9ng-imu}

\subparagraph{a) Vai trò của IMU trong hệ
thống}\label{a-vai-truxf2-cux1ee7a-imu-trong-hux1ec7-thux1ed1ng}

Cảm biến đo quán tính (IMU - Inertial Measurement Unit) đóng vai trò
then chốt trong việc phát hiện chuyển động của xe khi đang đỗ, cho phép
hệ thống phát hiện các tình huống bất thường như rung xe (cố gắng mở
cửa), kéo xe, hoặc cẩu xe. Đặc biệt, IMU cho phép ESP32-S3 duy trì
light sleep cho nhánh canh rung và đánh thức gần như tức thời khi có chuyển động thực sự,
trong khi heartbeat định kỳ vẫn quay về deep sleep theo timer để tiết kiệm năng lượng.

\subparagraph{b) Lựa chọn cảm biến:
LIS3DSH}\label{b-lux1ef1a-chux1ecdn-cux1ea3m-biux1ebfn-LIS3DH}

\textbf{LIS3DSH} của STMicroelectronics là cảm biến gia tốc 3 trục
(3-axis accelerometer) low-power được lựa chọn cho hệ thống. Các đặc
tính kỹ thuật chính được tóm tắt trong bảng sau:


{\thesissettableid{3.7}{tab-ch3-thong-so-ky-thuat-cam-bien-LIS3DH}
\def\LTcaptype{table}
\begin{longtable}{|L{0.426\linewidth}|L{0.542\linewidth}|}
\caption{Thông số kỹ thuật cảm biến LIS3DSH}\label{tab:ch3-thong-so-ky-thuat-cam-bien-LIS3DH}\\
\hline \textbf{Thông số} & \textbf{Giá trị} \\
\\\hline
\endfirsthead
\caption[]{Thông số kỹ thuật cảm biến LIS3DSH (tiếp theo)}\\
\hline \textbf{Thông số} & \textbf{Giá trị} \\
\\\hline
\endhead
\\\hline
\endfoot
\endlastfoot
Loại & 3-axis Digital Accelerometer \\\hline
Điện áp hoạt động & 1.8--3.6 V \\\hline
Giao tiếp & I2C (400 kHz) hoặc SPI (10 MHz) \\\hline
Phạm vi đo & +/-2g, +/-4g, +/-8g, +/-16g (có thể điều chỉnh) \\\hline
Độ phân giải & 16-bit \\\hline
Tiêu thụ (Active) & 10--50 µA \\\hline
Tiêu thụ (Low-power) & 2--5 µA \\\hline
Wake-up interrupt & Có (ngưỡng có thể cấu hình) \\\hline
Package & LGA-16 (3x3x1 mm) \\\hline
\end{longtable}
}


\subparagraph{c) Nguyên lý hoạt động trong hệ
thống}\label{c-nguyuxean-luxfd-houx1ea1t-ux111ux1ed9ng-trong-hux1ec7-thux1ed1ng}

LIS3DSH được cấu hình ở chế độ low-power với tần suất lấy mẫu 10 Hz và
chỉ bật chức năng motion detection. Khi phát hiện gia tốc vượt ngưỡng 0.2g
trên bất kỳ trục nào (X, Y, Z), cảm biến phát tín hiệu interrupt đến chân INT1
nối với GPIO41 của ESP32-S3. Vì GPIO41 không thuộc nhóm RTC GPIO, firmware
dùng GPIO wake trong light sleep cho cảnh báo rung, còn heartbeat định kỳ vẫn dùng timer wake {[}9{]}.

Ngưỡng gia tốc 0.2g được chọn để cân bằng giữa độ nhạy và khả năng chống
báo giả. Giá trị này đủ lớn để loại bỏ rung nhẹ từ môi trường (gió, xe
cộ đi ngang) nhưng đủ nhỏ để phát hiện các chuyển động thực sự như rung
của xe hoặc di chuyển.


{\thesissetfigureid{3.3}{fig-ch3-so-o-ket-noi-LIS3DH-voi-esp32-s3-qua-i2c}
\begin{figure}[htbp]
\centering
\pandocbounded{\safeincludesvg{./assets/figures/03-chuong-3-giai-phap-phan-cung-hinh-3-3.svg}}
\caption{Sơ đồ kết nối LIS3DSH với ESP32-S3 qua I2C}
\label{fig:ch3-so-o-ket-noi-LIS3DH-voi-esp32-s3-qua-i2c}
\end{figure}
}



{\thesissetfigureid{3.3a}{fig-ch3-so-o-chi-tiet-chan-ket-noi-LIS3DH-voi-esp32-s3}
\begin{figure}[htbp]
\centering
\pandocbounded{\safeincludesvg{./assets/figures/03-chuong-3-giai-phap-phan-cung-hinh-3-3a.svg}}
\caption{Sơ đồ chi tiết chân kết nối LIS3DSH với ESP32-S3}
\label{fig:ch3-so-o-chi-tiet-chan-ket-noi-LIS3DH-voi-esp32-s3}
\end{figure}
}


Cấu hình I2C sử dụng địa chỉ \texttt{0x18}, tốc độ \texttt{400 kHz},
với hai điện trở kéo lên \texttt{10 kOhm} trên các net \texttt{IMU-SCL}
và \texttt{IMU-SDA}. Ở runtime hiện tại, bus I2C được ánh xạ về
\texttt{GPIO1} (SCL) và \texttt{GPIO2} (SDA), chân \texttt{INT1} đi
vào \texttt{GPIO41}, còn \texttt{INT2} đi vào \texttt{GPIO42}. Cơ chế
wake do IMU kích hoạt dùng \textit{light sleep} qua GPIO wake thay vì deep sleep EXT0/EXT1.

\subparagraph{d) Ưu điểm của giải
pháp}\label{d-ux1b0u-ux111iux1ec3m-cux1ee7a-giux1ea3i-phuxe1p}

\begin{itemize}
\tightlist
\item
  \textbf{Tiết kiệm năng lượng tốt nhất:} Dòng tiêu thụ chỉ 2--5 µA ở
  chế độ low-power, trong khi vẫn duy trì khả năng phát hiện chuyển
  động. ESP32-S3 không cần chạy liên tục để kiểm tra cảm biến.
\item
  \textbf{Phát hiện chuyển động độc lập:} LIS3DSH tự xử lý việc phát
  hiện chuyển động bằng phần cứng, chỉ gửi interrupt khi có sự kiện ---
  không phụ thuộc vào CPU của ESP32.
\item
  \textbf{Hệ sinh thái phát triển phong phú:} Nhiều thư viện sẵn có cho
  Arduino và ESP-IDF; IC LIS3DSH hoặc cảm biến tương đương đều dễ tích
  hợp vào sơ đồ phần cứng của thiết bị.
\end{itemize}

Nếu OBD2 cung cấp bức tranh vận hành khi xe đang nổ máy và IMU đảm nhận
việc canh xe khi đỗ, thì khối còn lại phải bảo đảm hai cơ chế đó luôn
có điện đúng lúc. Phần tiếp theo vì thế đi vào hệ thống quản lý nguồn.

\paragraph{3.1.2.3. Thiết kế hệ thống quản lý
nguồn}\label{3213-thiux1ebft-kux1ebf-hux1ec7-thux1ed1ng-quux1ea3n-luxfd-nguux1ed3n}

Hệ thống quản lý nguồn là khối phức tạp nhất trong thiết kế phần cứng, vì
nó phải đồng thời giải quyết ba bài toán vốn dễ xung đột với nhau: tạo
đúng các mức điện áp cho từng tải, duy trì đường nguồn dự phòng khi cần và
bảo vệ ắc quy xe khỏi bị rút cạn quá mức. Vì vậy, khối này được tách thành
năm nhánh chức năng để người đọc dễ theo dõi vai trò của từng đường nguồn
thay vì chỉ nhìn một sơ đồ điện phức tạp.


{\thesissetfigureid{3.4}{fig-ch3-so-o-khoi-he-thong-quan-ly-nguon}
\begin{figure}[htbp]
\centering
\pandocbounded{\safeincludesvg{./assets/figures/03-chuong-3-giai-phap-phan-cung-hinh-3-4.svg}}
\caption{Sơ đồ khối hệ thống quản lý nguồn}
\label{fig:ch3-so-o-khoi-he-thong-quan-ly-nguon}
\end{figure}
}



{\thesissetfigureid{3.4a}{fig-ch3-kien-truc-nguon-va-phan-phoi-ien-ap-trong-he-thong}
\begin{figure}[htbp]
\centering
\pandocbounded{\safeincludesvg{./assets/figures/03-chuong-3-giai-phap-phan-cung-hinh-3-4a.svg}}
\caption{Kiến trúc nguồn và phân phối điện áp trong hệ thống}
\label{fig:ch3-kien-truc-nguon-va-phan-phoi-ien-ap-trong-he-thong}
\end{figure}
}


\subparagraph{a) LDO 3.3V cho ESP32-S3 (AP2112-3.3)}\label{a-buck-33v-cho-esp32-s3-xl1509-33e}

Nhánh 3.3V dùng \textbf{AP2112-3.3} như một LDO để hạ từ \textbf{bus 5V
chính} xuống 3.3V, cấp cho ESP32-S3 và các tải logic 3.3V. Cách tổ chức
này phù hợp với kiến trúc nguồn thực tế của bo mạch: MP2482 đảm nhiệm bước
hạ áp đầu tiên từ 12--24V xuống 5V, còn AP2112-3.3 đảm nhiệm bước ổn áp
cuối cho miền logic.


{\thesissettableid{3.8}{tab-ch3-thong-so-khoi-buck-3-3v}
\def\LTcaptype{table}
\begin{longtable}{|L{0.168\linewidth}|L{0.132\linewidth}|L{0.218\linewidth}|L{0.250\linewidth}|L{0.176\linewidth}|}
\caption{Thông số khối LDO 3.3V}\label{tab:ch3-thong-so-khoi-buck-3-3v}\\
\hline \textbf{Khối} & \textbf{IC} & \textbf{Input} & \textbf{Output} & \textbf{Tải} \\
\\\hline
\endfirsthead
\caption[]{Thông số khối LDO 3.3V (tiếp theo)}\\
\hline \textbf{Khối} & \textbf{IC} & \textbf{Input} & \textbf{Output} & \textbf{Tải} \\
\\\hline
\endhead
\\\hline
\endfoot
\endlastfoot
LDO 3.3V & AP2112-3.3 & 5V & 3.3V & ESP32-S3 \\\hline
\end{longtable}
}


\subparagraph{b) Buck 5V bus chính
(MP2482)}\label{b-buck-5v-bus-chuxednh-mp2482}

Nhánh buck 5V dùng \textbf{MP2482} để tạo \textbf{bus 5V chính} từ nguồn
12--24V của xe. Bus 5V này là nguồn trung gian cho các khối phụ trợ và
đầu vào mạch sạc pin.


{\thesissettableid{3.9}{tab-ch3-thong-so-khoi-buck-5v}
\def\LTcaptype{table}
\begin{longtable}{|L{0.168\linewidth}|L{0.132\linewidth}|L{0.218\linewidth}|L{0.250\linewidth}|L{0.176\linewidth}|}
\caption{Thông số khối Buck 5V}\label{tab:ch3-thong-so-khoi-buck-5v}\\
\hline \textbf{Khối} & \textbf{IC} & \textbf{Input} & \textbf{Output} & \textbf{Tải} \\
\\\hline
\endfirsthead
\caption[]{Thông số khối Buck 5V (tiếp theo)}\\
\hline \textbf{Khối} & \textbf{IC} & \textbf{Input} & \textbf{Output} & \textbf{Tải} \\
\\\hline
\endhead
\\\hline
\endfoot
\endlastfoot
Buck 5V & MP2482 & 12--24V & 5V & Bus 5V chính \\\hline
\end{longtable}
}


\subparagraph{c) Buck 3.8V/4V cho modem
(TPS54231)}\label{c-buck-38v4v-cho-modem-tps54231}

Do modem SIM7600CE-T làm việc trên miền nguồn thấp áp, hệ thống bố trí
nhánh buck riêng dùng \textbf{TPS54231} để hạ 12--24V xuống khoảng
\textbf{4V} cấp cho modem.


{\thesissettableid{3.10}{tab-ch3-thong-so-khoi-buck-3-8v-4v-cho-modem}
\def\LTcaptype{table}
\begin{longtable}{|L{0.168\linewidth}|L{0.132\linewidth}|L{0.218\linewidth}|L{0.250\linewidth}|L{0.176\linewidth}|}
\caption{Thông số khối Buck 3.8V/4V cho modem}\label{tab:ch3-thong-so-khoi-buck-3-8v-4v-cho-modem}\\
\hline \textbf{Khối} & \textbf{IC} & \textbf{Input} & \textbf{Output} & \textbf{Tải} \\
\\\hline
\endfirsthead
\caption[]{Thông số khối Buck 3.8V/4V cho modem (tiếp theo)}\\
\hline \textbf{Khối} & \textbf{IC} & \textbf{Input} & \textbf{Output} & \textbf{Tải} \\
\\\hline
\endhead
\\\hline
\endfoot
\endlastfoot
Buck 3.8V & TPS54\-231 & 12--24V & ~4V & SIM7600CE-T \\\hline
\end{longtable}
}


\subparagraph{d) Tăng áp 5V từ pin dự phòng (SX1308) và mạch chuyển
nguồn}\label{d-boost-5v-tux1eeb-pin-dux1ef1-phuxf2ng-sx1308-vuxe0-power-path}

Khi ắc quy xe yếu, nguồn dự phòng lấy từ pin 18650 1S (danh định
~3.7V). Khối \textbf{SX1308} tăng áp lên 5V để duy trì cấp
nguồn cho hệ thống.


{\thesissettableid{3.11}{tab-ch3-thong-so-khoi-boost-5v-du-phong}
\def\LTcaptype{table}
\begin{longtable}{|L{0.189\linewidth}|L{0.094\linewidth}|L{0.236\linewidth}|L{0.283\linewidth}|L{0.142\linewidth}|}
\caption{Thông số khối Boost 5V dự phòng}\label{tab:ch3-thong-so-khoi-boost-5v-du-phong}\\
\hline \textbf{Khối} & \textbf{IC} & \textbf{Input} & \textbf{Output} & \textbf{Tải} \\
\\\hline
\endfirsthead
\caption[]{Thông số khối Boost 5V dự phòng (tiếp theo)}\\
\hline \textbf{Khối} & \textbf{IC} & \textbf{Input} & \textbf{Output} & \textbf{Tải} \\
\\\hline
\endhead
\\\hline
\endfoot
\endlastfoot
Boost 5V & SX1308 & Battery ~3.7V & 5V & Nguồn dự phòng từ pin \\\hline
\end{longtable}
}


Mạch chuyển nguồn khi vận hành sử dụng \textbf{diode OR} giữa nhánh 5V chính
(MP2482) và nhánh 5V dự phòng (SX1308), kết hợp điều khiển GPIO để
bật/tắt các nhánh liên quan. Mục tiêu là chuyển nguồn mượt, không reset
hệ thống khi đổi nguồn.

\subparagraph{e) Giám sát điện áp và
LVD}\label{e-giuxe1m-suxe1t-ux111iux1ec7n-uxe1p-vuxe0-lvd}

LVD dùng ngưỡng profile kép theo firmware:

\begin{itemize}
\tightlist
\item
  \textbf{Profile 12V:} \(Switch_{\mathrm{OFF}} = 12.0\,\mathrm{V}\),
  \(Switch_{\mathrm{ON}} = 12.2\,\mathrm{V}\)
\item
  \textbf{Profile 24V:} \(Switch_{\mathrm{OFF}} = 24.0\,\mathrm{V}\),
  \(Switch_{\mathrm{ON}} = 24.4\,\mathrm{V}\)
\end{itemize}


{\thesissettableid{3.12}{tab-ch3-bang-trang-thai-chuyen-nguon-va-canh-bao}
\def\LTcaptype{table}
\begin{longtable}{|L{0.108\linewidth}|L{0.068\linewidth}|L{0.369\linewidth}|L{0.141\linewidth}|L{0.088\linewidth}|L{0.161\linewidth}|}
\caption{Bảng trạng thái chuyển nguồn và cảnh báo}\label{tab:ch3-bang-trang-thai-chuyen-nguon-va-canh-bao}\\
\hline \textbf{Chế độ} & \textbf{IGN} & \textbf{\(U_{\text{batt}}\)} & Nguồn thiết bị & Sạc pin & Cảnh báo \\
\\\hline
\endfirsthead
\caption[]{Bảng trạng thái chuyển nguồn và cảnh báo (tiếp theo)}\\
\hline \textbf{Chế độ} & \textbf{IGN} & \textbf{\(U_{\text{batt}}\)} & Nguồn thiết bị & Sạc pin & Cảnh báo \\
\\\hline
\endhead
\\\hline
\endfoot
\endlastfoot
Xe chạy & ON & Profile 12V: \(U_{\text{batt}} \ge 13.0\,\mathrm{V}\);\linebreak
Profile 24V: \(U_{\text{batt}} \ge 26.0\,\mathrm{V}\) & Ắc quy & Có & Không \\\hline
Xe đỗ & OFF & Profile 12V: \(U_{\text{batt}} > 12.0\,\mathrm{V}\);\linebreak
Profile 24V: \(U_{\text{batt}} > 24.0\,\mathrm{V}\) & Ắc quy & Không & Không \\\hline
Ắc quy yếu & OFF & Profile 12V: \(U_{\text{batt}} \le 12.0\,\mathrm{V}\);\linebreak
Profile 24V: \(U_{\text{batt}} \le 24.0\,\mathrm{V}\) & Pin 18650 1S & Không & Cảnh báo chuyển
nguồn \\\hline
Ắc quy phục hồi & OFF & Profile 12V: \(U_{\text{batt}} \ge 12.2\,\mathrm{V}\);\linebreak
Profile 24V: \(U_{\text{batt}} \ge 24.4\,\mathrm{V}\) & Ắc quy & Không & Cảnh báo
phục hồi \\\hline
\end{longtable}
}


Ngoài kênh ADC, tín hiệu trạng thái LVD từ khối LVD phần cứng được đưa về
\textbf{GPIO19 (LVD\_STATUS)} để giám sát nhanh. Quy ước vận hành:
\textbf{GPIO19 HIGH = low-voltage}, \textbf{GPIO19 LOW = bình thường}.

\subparagraph{f) Mạch sạc pin 1S
(TP5100)}\label{f-mux1ea1ch-sux1ea1c-pin-1s-tp4056}

Khối sạc dùng \textbf{TP5100} ở cấu hình \textbf{1 cell}, nhận
\textbf{5V từ nhánh MP2482} và sạc pin 18650 1S tới mức
\textbf{4.2V/1S}. So với phương án TP4056 trước đây, TP5100 phù hợp hơn
với nhánh nguồn dự phòng vì đây là bộ sạc chuyển mạch, giảm tổn hao
nhiệt khi nạp.


{\thesissettableid{3.13}{tab-ch3-thong-so-khoi-sac-pin}
\def\LTcaptype{table}
\begin{longtable}{|L{0.189\linewidth}|L{0.094\linewidth}|L{0.236\linewidth}|L{0.283\linewidth}|L{0.142\linewidth}|}
\caption{Thông số khối sạc pin}\label{tab:ch3-thong-so-khoi-sac-pin}\\
\hline \textbf{Khối} & \textbf{IC} & \textbf{Input} & \textbf{Output} & \textbf{Tải} \\
\\\hline
\endfirsthead
\caption[]{Thông số khối sạc pin (tiếp theo)}\\
\hline \textbf{Khối} & \textbf{IC} & \textbf{Input} & \textbf{Output} & \textbf{Tải} \\
\\\hline
\endhead
\\\hline
\endfoot
\endlastfoot
Sạc pin & TP5100 & 5V từ MP2482 & 4.2V & Pin 18650 1S \\\hline
\end{longtable}
}


Dòng sạc TP5100 được cấu hình bằng điện trở thiết lập dòng sạc và còn
phụ thuộc điều kiện nhiệt của mạch, vì vậy không dùng một giá trị dòng
cố định cho mọi điều kiện vận hành.

\subparagraph{g) Pin dự phòng 18650 1S
Li-ion}\label{g-pin-dux1ef1-phuxf2ng-18650-1s-li-ion}

Pin 18650 1S Li-ion được chọn làm nguồn dự phòng với cấu hình 1 cell đơn
giản, kết hợp khối sạc TP5100, khối boost SX1308 và mạch bảo vệ pin 1S
trong kiến trúc nguồn của hệ thống.


{\thesissettableid{3.14}{tab-ch3-thong-so-ky-thuat-pin-du-phong-18650-1s}
\def\LTcaptype{table}
\begin{longtable}{|L{0.310\linewidth}|L{0.658\linewidth}|}
\caption{Thông số kỹ thuật pin dự phòng 18650 1S}\label{tab:ch3-thong-so-ky-thuat-pin-du-phong-18650-1s}\\
\hline \textbf{Thông số} & \textbf{Giá trị} \\
\\\hline
\endfirsthead
\caption[]{Thông số kỹ thuật pin dự phòng 18650 1S (tiếp theo)}\\
\hline \textbf{Thông số} & \textbf{Giá trị} \\
\\\hline
\endhead
\\\hline
\endfoot
\endlastfoot
Loại & Li-ion 18650 1S \\\hline
Dung lượng & 3500 mAh @ 3.7V \\\hline
Năng lượng & ~12.95 Wh \\\hline
Điện áp & 3.0--4.2V (nominal 3.7V) \\\hline
Số chu kỳ & 500--1000 chu kỳ (80\% capacity) \\\hline
\end{longtable}
}


\textbf{Tính toán thời gian hoạt động từ pin dự phòng:}

\emph{Các giá trị dưới đây là ước tính theo cùng giả định tổn hao của
phương án gốc, với dung lượng hiệu quả xấp xỉ 3150 mAh (tương đương 90\%
dung lượng danh định), không phải số đo thực nghiệm mới trên cell
18650.}

\emph{Chế độ heartbeat (đỗ xe, chu kỳ 15 phút):}

\begin{itemize}
\tightlist
\item
  Dòng trung bình:
  \(I_{\text{avg}} = \dfrac{250\,\mathrm{mA}\times 15\,\mathrm{s} + 3\,\mathrm{mA}\times 885\,\mathrm{s}}{900\,\mathrm{s}} = 7.1\,\mathrm{mA}\)
\item
  Dung lượng hiệu quả (sau tổn hao): 3150 mAh
\item
  Thời gian hoạt động:
  \(T = \dfrac{3150}{7.1} = 444\,\mathrm{giờ} \approx 18.5\,\text{ngày}\)
\end{itemize}

\emph{Chế độ heartbeat (đỗ xe, chu kỳ 30 phút):}

\begin{itemize}
\tightlist
\item
  Dòng trung bình: \(\approx 5\,\mathrm{mA}\)
\item
  Thời gian hoạt động:
  \(T = \dfrac{3150}{5} = 630\,\mathrm{giờ} \approx 26\,\text{ngày}\)
\end{itemize}

\emph{Chế độ cảnh báo (track liên tục):}

\begin{itemize}
\tightlist
\item
  Dòng trung bình: \(\approx 250\,\mathrm{mA}\)
\item
  Thời gian hoạt động:
  \(T = \dfrac{3150}{250} = 12.6\,\mathrm{giờ}\)
\end{itemize}

\emph{Giá trị 12.6 giờ ở đây là phép tính lý tưởng theo tải giả định 250
mA tại mức cell; khi quy đổi sang tải hệ thống thực tế có hao hụt boost
và các pha truyền dữ liệu đỉnh cao hơn, thời lượng tracking liên tục
trong Chương 4 được lấy theo khoảng bảo thủ hơn là
~2.6--3.0 giờ.}

Kết quả quy đổi cho thấy bộ pin tham chiếu 18650 1S vẫn đủ khả năng duy trì
hoạt động thiết bị trong nhiều ngày ở chế độ heartbeat và trên 12 giờ ở
chế độ cảnh báo liên tục theo mô hình lý tưởng; trong vận hành toàn hệ
thống, các giá trị bảo thủ hơn ở Chương 4 phù hợp hơn cho mục tiêu duy
trì giám sát tạm thời khi ắc quy xe yếu, nhưng vẫn cần được kiểm thử
thực nghiệm lại nếu chuyển sang mẫu pin mới trong pha chế tạo tiếp theo
{[}14{]}.

\paragraph{3.1.2.4. Bảng tổng hợp linh kiện (Bill of
Materials)}\label{3214-bux1ea3ng-tux1ed5ng-hux1ee3p-linh-kiux1ec7n-bill-of-materials}


{\thesissettableid{3.14A}{tab-ch3-bang-tong-hop-linh-kien-he-thong-tracker-bom}
\def\LTcaptype{table}
\begin{longtable}{|L{0.068\linewidth}|L{0.264\linewidth}|L{0.068\linewidth}|L{0.068\linewidth}|L{0.151\linewidth}|L{0.316\linewidth}|}
\caption{Bảng tổng hợp linh kiện hệ thống thiết bị theo dõi (BOM)}\label{tab:ch3-bang-tong-hop-linh-kien-he-thong-tracker-bom}\\
\hline \textbf{STT} & \textbf{Thành phần} & \textbf{Đơn vị} & \textbf{SL} & \textbf{Giá tham chiếu T4/2026 (VND)} & \textbf{Ghi chú} \\
\\\hline
\endfirsthead
\caption[]{Bảng tổng hợp linh kiện hệ thống thiết bị theo dõi (BOM) (tiếp theo)}\\
\hline \textbf{STT} & \textbf{Thành phần} & \textbf{Đơn vị} & \textbf{SL} & \textbf{Giá tham chiếu T4/2026 (VND)} & \textbf{Ghi chú} \\
\\\hline
\endhead
\\\hline
\endfoot
\endlastfoot
1 & ESP32-S3-WROOM-1 (N16R8) & Cái & 1 & 95.000 & MCU trung
tâm của bo mạch \\\hline
2 & IC cảm biến LIS3DSH & Cái & 1 & 40.000 & Cảm biến gia tốc 3
trục \\\hline
3 & vgate iCar Pro (OBD2 BLE) & Cái & 1 & 367.000 & BLE 4.0,
ELM327 compatible \\\hline
4 & SIMCom SIM7600CE-T & Bộ & 1 & 650.000 & Modem LTE + GNSS
tích hợp \\\hline
5 & Pin 18650 1S Li-ion 3500mAh & Cái & 1 & 79.000 & Loại có
protection board \\\hline
6 & IC sạc TP5100 + mạch phụ trợ & Cái & 1 & 25.000 & Input 5V,
output 4.2V, dòng theo điện trở thiết lập \\\hline
7 & Mạch LDO AP2112-3.3 & Cái & 1 & 5.000 & 5V
-\textgreater{} 3.3V cấp ESP32-S3 \\\hline
8 & Mạch buck MP2482 5V & Cái & 1 & 17.000 & Bus 5V chính \\\hline
9 & Mạch buck TPS54231 ~4V & Cái & 1 & 9.000 &
12--24V -\textgreater{} ~4V cấp modem SIM7600CE-T \\\hline
10 & Mạch boost SX1308 5V dự phòng & Cái & 1 & 12.000 & Nguồn dự phòng từ
pin 18650 1S \\\hline
11 & Khối giám sát LVD + ADC & Cái & 1 & 8.000 & Giám sát LVD
(GPIO19) \\\hline
12 & Diode Schottky (SS54/SS34) & Cái & 2 & 2.500 & Diode-OR power
path \\\hline
13 & Mạch bảo vệ pin 1S (BMS) & Cái & 1 & 12.000 & Bảo vệ pin
18650 1S \\\hline
14 & Điện trở (10kOhm, 2.2kOhm) & Gói & 1 & 5.000 & Voltage
divider, pull-up \\\hline
15 & Tụ điện (100uF, 220uF) & Gói & 1 & 5.000 & Lọc nhiễu,
decoupling \\\hline
16 & Connector, header pin & Gói & 1 & 55.000 & Header, jack OBD2 và
đầu nối nguồn \\\hline
17 & PCB 2 lớp (~50x50 mm) & Cái & 1 & 20.000 &
Bo mạch theo thiết kế của đề tài \\\hline
18 & Vỏ bảo vệ (tùy chọn) & Cái & 1 & 25.000 & Nhựa hoặc kim
loại \\\hline
19 & Dây nối, cáp, phụ kiện & - & - & 60.000 & Dây điện, anten
4G/GNSS, khay SIM \\\hline
20 & Linh kiện phụ trợ khác & - & - & 20.000 & Fuse, switch,
LED \\\hline
\end{longtable}
}



{\thesissettableid{3.15}{tab-ch3-tong-hop-chi-phi-theo-nhom}
\def\LTcaptype{table}
\begin{longtable}{|L{0.160\linewidth}|L{0.580\linewidth}|L{0.224\linewidth}|}
\caption{Tổng hợp chi phí theo nhóm}\label{tab:ch3-tong-hop-chi-phi-theo-nhom}\\
\hline \textbf{Nhóm} & \textbf{Hạng mục} & \textbf{Chi phí (VND)} \\
\\\hline
\endfirsthead
\caption[]{Tổng hợp chi phí theo nhóm (tiếp theo)}\\
\hline \textbf{Nhóm} & \textbf{Hạng mục} & \textbf{Chi phí (VND)} \\
\\\hline
\endhead
\\\hline
\endfoot
\endlastfoot
A & Thành phần chính (MCU, cảm biến, modem, OBD2, pin) &
1.231.000 \\\hline
B & Quản lý nguồn (AP2112-3.3, MP2482, TPS54231, SX1308, TP5100, diode OR,
khối LVD, pin 1S) & 93.000 \\\hline
C & Kết nối, RF và linh kiện phụ trợ & 145.000 \\\hline
D & PCB và vỏ (tùy chọn) & 45.000 \\\hline
& \textbf{Tổng cộng} & \textbf{1.514.000} \\\hline
\end{longtable}
}


Theo mặt bằng giá tham chiếu thị trường Việt Nam được rà soát ngày
12/04/2026, tổng chi phí vật tư cho một bộ thiết bị theo dõi hoàn chỉnh vào khoảng
\textbf{1.514.000 VND}. Mức này vẫn nằm dưới ngưỡng chi phí mục tiêu của
đồ án, đồng thời phản ánh thực tế hơn so với mặt bằng giá linh kiện ở
thời điểm đầu năm.

\subsection{3.2. Phân tích và đề xuất giải pháp Firmware -- Firmware
analysis and solution}\label{32-ux111ux1ec1-xuux1ea5t-cuxe1c-giux1ea3i-phuxe1p--proposed-solutions}

\subsubsection{3.2.1. Phân tích và lựa chọn giải pháp Firmware}

\paragraph{3.2.1.1. Đặt vấn đề cho giải pháp
Firmware}\label{3121-ux111ux1eb7t-vux1ea5n-ux111ux1ec1-cho-giux1ea3i-phuxe1p-firmware}

Nếu xem phần cứng là ``cơ thể'' của thiết bị, thì firmware là lớp biến
tập linh kiện đó thành một hệ thống biết \emph{phản ứng đúng lúc}. Người
dùng không trực tiếp nhìn thấy firmware, nhưng việc vị trí có cập nhật
đều, cảnh báo có đến kịp hay không, hay ắc quy xe có bị ảnh hưởng hay
không đều do lớp này quyết định. Vì nằm ở điểm giao nhau giữa phần cứng,
truyền thông và yêu cầu nghiệp vụ, bài toán firmware không dừng ở chuyện
``cho thiết bị chạy được'', mà phải bảo đảm thiết bị chạy đúng, ổn định
và tiết kiệm năng lượng trong môi trường thực tế.

\begin{itemize}
\tightlist
\item
  \textbf{Đa nhiệm thời gian thực}: thu thập cảm biến, điều khiển modem,
  truyền dữ liệu và xử lý cảnh báo phải chạy song song, độ trễ thấp.
\item
  \textbf{Quản lý năng lượng nghiêm ngặt}: chuyển trạng thái linh hoạt
  giữa active/sleep/deep sleep theo điều kiện vận hành xe.
\item
  \textbf{Giao tiếp đa giao thức}: BLE (OBD2), UART (modem), MQTT/HTTP
  (máy chủ), I2C/ADC/GPIO (ngoại vi phần cứng).
\item
  \textbf{Độ bền vận hành cao}: có cơ chế retry, fallback và phục hồi
  sau mất kết nối mạng hoặc lỗi ngoại vi.
\end{itemize}

\paragraph{3.2.1.2. So sánh các phương án nền tảng
Firmware}\label{3122-so-suxe1nh-cuxe1c-phux1b0ux1a1ng-uxe1n-nux1ec1n-tux1ea3ng-firmware}


{\thesissettableid{3.5A}{tab-ch3-so-sanh-cac-phuong-an-nen-tang-firmware}
\def\LTcaptype{table}
\begin{longtable}{|L{0.138\linewidth}|L{0.236\linewidth}|L{0.246\linewidth}|L{0.224\linewidth}|L{0.100\linewidth}|}
\caption{So sánh các phương án nền tảng firmware}\label{tab:ch3-so-sanh-cac-phuong-an-nen-tang-firmware}\\
\hline \textbf{Phương án} & \textbf{Mô tả} & \textbf{Ưu điểm} & \textbf{Hạn chế} & \textbf{Mức phù hợp} \\
\\\hline
\endfirsthead
\caption[]{So sánh các phương án nền tảng firmware (tiếp theo)}\\
\hline \textbf{Phương án} & \textbf{Mô tả} & \textbf{Ưu điểm} & \textbf{Hạn chế} & \textbf{Mức phù hợp} \\
\\\hline
\endhead
\\\hline
\endfoot
\endlastfoot
\textbf{PA-FW1: Arduino Core + Superloop} & Vòng lặp chính tuần tự, xử
lý tác vụ theo kiểu hỏi lần lượt từng phần trong vòng lặp & Dễ bắt đầu, ít cấu hình & Khó mở rộng đa nhiệm
thực sự, khó tối ưu deep sleep phức tạp, quản lý lỗi hạn chế & Trung
bình \\\hline
\textbf{PA-FW2: ESP-IDF + FreeRTOS (Đã chọn)} & Kiến trúc task và event,
driver chính thức Espressif & Đa nhiệm tốt, hỗ trợ power management sâu,
tích hợp NimBLE / modem / UART / I2C ổn định & Độ phức tạp cao hơn Arduino &
\textbf{Cao} \\\hline
\textbf{PA-FW3: Zephyr RTOS trên ESP32} & RTOS đa nền tảng, kiến trúc
phân lớp & Tính chuẩn hóa tốt, khả năng mở rộng dài hạn & Hệ sinh thái
ESP32 chuyên biệt và tài liệu thực chiến BLE-OBD2/modem ít hơn ESP-IDF &
Trung bình \\\hline
\end{longtable}
}


\paragraph{3.2.1.3. Chọn giải pháp
Firmware}\label{3123-chux1ecdn-giux1ea3i-phuxe1p-firmware}

Từ kết quả so sánh, đồ án chọn \textbf{PA-FW2: ESP-IDF + FreeRTOS} làm
nền tảng firmware chính. Ở đây, ESP-IDF có thể hiểu là bộ công cụ chính
thức để lập trình ESP32-S3; FreeRTOS là lớp điều phối giúp firmware chia
việc theo task và sự kiện thay vì dồn mọi xử lý vào một vòng lặp lớn.
Đây không phải là lựa chọn ``mạnh nhất trên giấy'', mà là lựa chọn bám
sát nhất cách thiết bị phải vận hành ngoài thực tế: nhiều ngoại vi,
nhiều trạng thái, yêu cầu tiết kiệm điện và cần dễ truy vết khi có lỗi.

Nền tảng này thuyết phục hơn ở bốn điểm chính:

\begin{itemize}
\tightlist
\item
  \textbf{Khớp yêu cầu điều phối thời gian thực của hệ thống}: ESP-IDF +
  FreeRTOS cho phép tổ chức máy trạng thái trung tâm, các driver chạy
  không đồng thời tuyệt đối và các cơ chế khóa/chờ ở mức khối chức năng. Nói cách khác, đây là
  nền tảng đủ linh hoạt để firmware vừa theo dõi nhiều tín hiệu cùng lúc
  vừa không làm luồng xử lý chính trở nên rối và khó kiểm soát.
\item
  \textbf{Tối ưu cho ESP32-S3}: driver và power feature chính thức giúp
  giảm rủi ro tích hợp phần cứng, đồng thời rút ngắn đáng kể thời gian
  xử lý các lỗi phát sinh trong giai đoạn chế tạo và kiểm thử.
\item
  \textbf{Phù hợp bài toán năng lượng}: hỗ trợ deep sleep/wakeup và điều
  khiển ngoại vi theo trạng thái, nên phù hợp với một thiết bị phải
  luôn cân bằng giữa khả năng giám sát và mức tiêu thụ điện.
\item
  \textbf{Dễ kiểm soát chất lượng vận hành}: thuận lợi xây dựng máy
  trạng thái, cơ chế thử lại và phương án dự phòng cho BLE/modem. Điều này đặc biệt
  quan trọng vì thiết bị không hoạt động trong phòng thí nghiệm lý tưởng
  mà trong môi trường xe có nhiễu, rung động và mất sóng bất chợt.
\end{itemize}

Phần so sánh trên trả lời vì sao ESP-IDF + FreeRTOS là nền tảng phù hợp
nhất cho đồ án. Từ đây, mục 3.2.2 chuyển từ câu hỏi ``chọn gì'' sang câu
hỏi ``tổ chức ra sao'' trong kiến trúc firmware thực tế.

\subsubsection{3.2.2. Giải pháp
Firmware}\label{322-giux1ea3i-phuxe1p-firmware}

Firmware là nơi toàn bộ quyết định vận hành của thiết bị được hiện thực
hóa thành hành vi cụ thể: đọc cảm biến, đánh giá trạng thái xe, điều
khiển modem, quản lý năng lượng và gửi dữ liệu lên máy chủ. Trên cơ sở
phương án đã chọn ở mục 3.2.1, phần này trình bày cách nhóm biến các yêu
cầu đó thành một kiến trúc firmware đủ rõ ràng để phát triển, đủ gọn để
vận hành ổn định và đủ linh hoạt để mở rộng sau này.

\paragraph{3.2.2.1. Kiến trúc firmware và luồng hoạt
động}\label{3221-kiux1ebfn-truxfac-firmware-vuxe0-luux1ed3ng-houx1ea1t-ux111ux1ed9ng}

\subparagraph{a) Kiến trúc phân lớp (Layered
Architecture)}\label{a-kiux1ebfn-truxfac-phuxe2n-lux1edbp-layered-architecture}

Để người đọc dễ hình dung, firmware được tách theo từng lớp trách nhiệm
thay vì trình bày như một khối mã nguồn lớn. Có thể hiểu đơn giản: lớp
dưới cùng ``nói chuyện với phần cứng'', lớp giữa giữ trạng thái và quyết
định hệ thống đang ở chế độ nào, còn lớp trên cùng đóng gói dữ liệu và
giao tiếp với hệ thống máy chủ. Cách chia này giúp trả lời rõ ba câu hỏi quan
trọng: lớp nào làm việc với phần cứng, lớp nào giữ trạng thái vận hành,
và lớp nào chịu trách nhiệm giao tiếp với hệ thống bên ngoài.


{\thesissettableid{3.2.2A}{tab-ch3-tong-hop-thong-tin-ve-lop-thanh-phan-chinh-vai-tro}
\def\LTcaptype{table}
\begin{longtable}{|L{0.160\linewidth}|L{0.224\linewidth}|L{0.580\linewidth}|}
\caption{Tổng hợp thông tin về Lớp, Thành phần chính, Vai trò}\label{tab:ch3-tong-hop-thong-tin-ve-lop-thanh-phan-chinh-vai-tro}\\
\hline \textbf{Lớp} & \textbf{Thành phần chính} & \textbf{Vai trò} \\
\\\hline
\endfirsthead
\caption[]{Tổng hợp thông tin về Lớp, Thành phần chính, Vai trò (tiếp theo)}\\
\hline \textbf{Lớp} & \textbf{Thành phần chính} & \textbf{Vai trò} \\
\\\hline
\endhead
\\\hline
\endfoot
\endlastfoot
1 & \texttt{adc\_reader}, \texttt{imu\_lis3dh}, \texttt{modem\_at},
\texttt{modem\_lte}, \texttt{modem\_gnss}, \texttt{ble\_mgr},
\texttt{ble\_obd} & Trừu tượng hóa phần cứng và giao thức ngoại vi \\\hline
2 & \texttt{power\_mgr}, \texttt{nvs\_config}, \texttt{app\_state} &
Quản lý nguồn, lưu cấu hình NVS và duy trì ngữ cảnh RTC qua deep
sleep \\\hline
3 & \texttt{state\_machine}, \texttt{command\_handler} & Điều phối trạng
thái vận hành, lệnh điều khiển và luồng OTA \\\hline
4 & \texttt{data\_formatter}, \texttt{mqtt\_client} & Đóng gói gói dữ liệu
JSON, publish/subscribe MQTT theo topic chuẩn của hệ thống \\\hline
5 & \texttt{offline\_queue}, \texttt{sd\_log\_store} & Hàng đợi cục bộ
trên microSD (đưa vào hàng đợi, phát lại theo xác nhận/QoS, dọn dữ liệu theo hạn mức, metadata phiên
ghi) \\\hline
\end{longtable}
}


Trong triển khai hiện tại, firmware không bị chia nhỏ thành quá nhiều
tác vụ ngang hàng. Phần lớn quyết định vận hành vẫn đi qua một
máy trạng thái trung tâm, tức là một bộ điều phối quyết định
thiết bị đang ở chế độ nào và bước tiếp theo phải làm gì. Các cơ chế của
FreeRTOS chỉ được dùng ở những nơi thật sự cần phối hợp nhiều tiến trình
như BLE, lớp AT command của modem hoặc các hàng đợi dữ liệu.


{\thesissetfigureid{3.5}{fig-ch3-so-o-kien-truc-phan-lop-cua-firmware}
\begin{figure}[htbp]
\centering
\pandocbounded{\safeincludesvg{./assets/figures/04-chuong-3-giai-phap-firmware-hinh-3-5.svg}}
\caption{Sơ đồ kiến trúc phân lớp của firmware}
\label{fig:ch3-so-o-kien-truc-phan-lop-cua-firmware}
\end{figure}
}


\subparagraph{b) Cơ chế điều phối với
FreeRTOS}\label{b-cux1a1-chux1ebf-ux111iux1ec1u-phux1ed1i-vux1edbi-freertos}

Điểm quan trọng ở đây là FreeRTOS được dùng có chọn lọc, không dùng theo
kiểu ``càng nhiều task càng tốt''. Luồng chính trong \texttt{app\_main()}
sẽ đọc cấu hình, xác định vì sao thiết bị vừa thức dậy, rồi lặp nhanh
\texttt{state\_machine\_run()} để điều phối hành vi kế tiếp. Cách tổ
chức này giữ cho logic điều khiển nằm ở một chỗ, nhờ đó người đọc mã
nguồn dễ bám hơn và quá trình debug cũng bớt rối hơn.

\begin{Shaded}
\begin{Highlighting}[]
\DataTypeTok{void}\NormalTok{ app\_main}\OperatorTok{(}\DataTypeTok{void}\OperatorTok{)} \OperatorTok{\{}
\NormalTok{    ESP\_ERROR\_CHECK}\OperatorTok{(}\NormalTok{nvs\_config\_init}\OperatorTok{());}

\NormalTok{    config\_t config }\OperatorTok{=} \OperatorTok{\{}\DecValTok{0}\OperatorTok{\};}
\NormalTok{    ESP\_ERROR\_CHECK}\OperatorTok{(}\NormalTok{nvs\_config\_load}\OperatorTok{(\&}\NormalTok{config}\OperatorTok{));}

\NormalTok{    app\_state\_t state }\OperatorTok{=}\NormalTok{ APP\_STATE\_INIT}\OperatorTok{;}
\NormalTok{    esp\_sleep\_wakeup\_cause\_t wakeup }\OperatorTok{=}\NormalTok{ esp\_sleep\_get\_wakeup\_cause}\OperatorTok{();}
    \ControlFlowTok{if} \OperatorTok{(}\NormalTok{wakeup }\OperatorTok{==}\NormalTok{ ESP\_SLEEP\_WAKEUP\_TIMER}\OperatorTok{)} \OperatorTok{\{}
\NormalTok{        state }\OperatorTok{=}\NormalTok{ APP\_STATE\_HEARTBEAT}\OperatorTok{;}
    \OperatorTok{\}} \ControlFlowTok{else} \ControlFlowTok{if} \OperatorTok{(}\NormalTok{wakeup }\OperatorTok{==}\NormalTok{ ESP\_SLEEP\_WAKEUP\_EXT0}\OperatorTok{)} \OperatorTok{\{}
\NormalTok{        state }\OperatorTok{=}\NormalTok{ APP\_STATE\_ALARM}\OperatorTok{;}
    \OperatorTok{\}}

\NormalTok{    ESP\_ERROR\_CHECK}\OperatorTok{(}\NormalTok{state\_machine\_init}\OperatorTok{(\&}\NormalTok{config}\OperatorTok{));}

    \ControlFlowTok{while} \OperatorTok{(}\KeywordTok{true}\OperatorTok{)} \OperatorTok{\{}
\NormalTok{        state }\OperatorTok{=}\NormalTok{ state\_machine\_run}\OperatorTok{(}\NormalTok{state}\OperatorTok{);}
\NormalTok{        vTaskDelay}\OperatorTok{(}\NormalTok{pdMS\_TO\_TICKS}\OperatorTok{(}\DecValTok{100}\OperatorTok{));}
    \OperatorTok{\}}
\OperatorTok{\}}
\end{Highlighting}
\end{Shaded}

Các primitive FreeRTOS vẫn xuất hiện ở những nơi cần ``giữ trật tự''
khi nhiều phần cùng dùng chung tài nguyên. Cụ thể, NimBLE host chạy ở
một task nền riêng; \texttt{ble\_mgr} dùng queue và mutex để điều phối
kết quả quét và kết nối; \texttt{ble\_obd} dùng semaphore để chờ phản
hồi từ adapter; còn \texttt{modem\_at} dùng mutex để bảo đảm các lệnh AT
đi qua UART theo đúng thứ tự.


{\thesissetfigureid{3.6}{fig-ch3-so-o-tuong-tac-giua-cac-freertos-task}
\begin{figure}[htbp]
\centering
\pandocbounded{\safeincludesvg{./assets/figures/04-chuong-3-giai-phap-firmware-hinh-3-6.svg}}
\caption{Sơ đồ tương tác giữa các FreeRTOS task}
\label{fig:ch3-so-o-tuong-tac-giua-cac-freertos-task}
\end{figure}
}


\subparagraph{c) Luồng hoạt động cơ
bản}\label{c-luux1ed3ng-houx1ea1t-ux111ux1ed9ng-cux1a1-bux1ea3n}

Nếu nhìn từ góc độ vận hành, firmware lặp lại một chu trình rất rõ:
khởi tạo, nhận biết trạng thái xe, chọn chế độ phù hợp, thực hiện tác vụ
và quay về trạng thái tiết kiệm năng lượng khi có thể. Luồng tổng thể
được tổ chức theo trình tự sau:

\begin{enumerate}
\def\labelenumi{\arabic{enumi}.}
\item
  \textbf{Khởi tạo ngoại vi}: Cấu hình và khởi động các peripheral gồm
  IMU (LIS3DSH qua I2C), modem SIM7600CE-T (UART1, Auto mode
  \texttt{AT+CNMP=2}, GNSS qua \texttt{AT+CGNSINF}), ADC (đọc điện áp),
  và BLE stack (NimBLE), tức lớp phần mềm nền để xử lý kết nối
  Bluetooth Low Energy.
\item
  \textbf{Đọc trạng thái IGN và điện áp ắc quy}: Hệ thống ưu tiên đọc
  trạng thái động cơ (IGN) trực tiếp từ ECU qua OBD2 BLE. Nếu không kết
  nối được OBD2, hệ thống fallback sang đo điện áp ắc quy qua ADC theo
  profile: profile 12V dùng \(IGN_{\mathrm{ON}} \ge 13.0\,\mathrm{V}\) và
  \(IGN_{\mathrm{OFF}} \le 12.0\,\mathrm{V}\); profile 24V dùng
  \(IGN_{\mathrm{ON}} \ge 26.0\,\mathrm{V}\) và
  \(IGN_{\mathrm{OFF}} \le 24.0\,\mathrm{V}\).
\item
  \textbf{Quyết định chế độ hoạt động}: Dựa trên trạng thái IGN và dữ
  liệu cảm biến, firmware chuyển sang chế độ phù hợp (lái xe, đỗ xe,
  hoặc cảnh báo).
\item
  \textbf{Thực thi tác vụ trong từng chế độ}:

  \begin{itemize}
  \tightlist
  \item
    \textbf{Chế độ lái xe (Driving)}: Kết nối OBD2, theo dõi liên tục,
    giữ BLE hoạt động, gửi dữ liệu vận hành định kỳ.
  \item
    \textbf{Chế độ đỗ xe (Parked)}: Không kết nối OBD2, chuyển sang deep
    sleep, gửi heartbeat định kỳ.
  \item
    \textbf{Chế độ cảnh báo (Alarm)}: Gửi cảnh báo ngay lập tức, theo
    dõi liên tục, có thể kết nối OBD2 tùy chọn.
  \end{itemize}
\item
  \textbf{Xử lý sau tác vụ}: Nếu IGN ON, giữ kết nối BLE và không deep
  sleep. Nếu IGN OFF, ngắt BLE và chuyển sang deep sleep để tiết kiệm
  năng lượng.
\end{enumerate}


{\thesissetfigureid{3.7}{fig-ch3-luu-o-thuat-toan-luong-hoat-ong-chinh-cua-firmware}
\begin{figure}[htbp]
\centering
\pandocbounded{\safeincludesvg{./assets/figures/04-chuong-3-giai-phap-firmware-hinh-3-7.svg}}
\caption{Lưu đồ thuật toán luồng hoạt động chính của firmware}
\label{fig:ch3-luu-o-thuat-toan-luong-hoat-ong-chinh-cua-firmware}
\end{figure}
}


\paragraph{3.2.2.2. Module giao tiếp Bluetooth
OBD2}\label{3222-module-giao-tiux1ebfp-bluetooth-obd2}

\subparagraph{a) Tổng quan giao tiếp BLE
OBD2}\label{a-tux1ed5ng-quan-giao-tiux1ebfp-ble-obd2}

Nói ngắn gọn, module này là phần giúp thiết bị theo dõi ``hỏi xe đang ở trạng thái
nào'' thông qua cổng OBD2 mà không cần dây nối trực tiếp tới vi điều
khiển. Adapter vgate iCar Pro được cắm vào cổng chẩn đoán OBD-II của xe,
sau đó giao tiếp không dây với thiết bị theo dõi bằng BLE. Trong firmware, NimBLE
là lớp phần mềm lo phần Bluetooth; còn GATT là cách hai thiết bị tổ chức
dịch vụ và trường dữ liệu để đọc/ghi thông tin với nhau.

Kiến trúc module BLE OBD2 được chia thành bốn tầng để mỗi tầng lo một
việc riêng. \texttt{state\_machine.c} quyết định khi nào cần đọc dữ liệu
xe; \texttt{ble\_obd.c} xử lý phần lệnh OBD2; \texttt{ble\_mgr.c} lo kết
nối Bluetooth; còn NimBLE đảm nhiệm lớp BLE nền tảng. Cách chia này giúp
việc đọc dữ liệu xe không bị lẫn với chi tiết kỹ thuật của GATT/BLE, nên
dễ debug và mở rộng hơn.

\subparagraph{b) Quy trình kết nối BLE
OBD2}\label{b-quy-truxecnh-kux1ebft-nux1ed1i-ble-obd2}

Quy trình kết nối BLE với adapter OBD2 diễn ra theo các bước sau:

\begin{enumerate}
\def\labelenumi{\arabic{enumi}.}
\item
  \textbf{Kiểm tra địa chỉ BLE đã lưu}: Đọc địa chỉ MAC của vgate iCar
  Pro từ bộ nhớ flash (NVS). Nếu đã có địa chỉ, chuyển trực tiếp sang
  bước kết nối (kết nối lại nhanh trong 1--3 giây).
\item
  \textbf{Quét tìm thiết bị (Scan)}: Nếu chưa có địa chỉ, thực hiện BLE
  scan để tìm adapter. Quá trình scan lọc thiết bị theo tên (device
  name) hoặc service UUID đặc trưng của OBD2 adapter.
\item
  \textbf{Kết nối GATT}: Thiết lập kết nối BLE với adapter, thực hiện
  GATT service discovery để tìm các characteristic cần thiết (TX và RX
  characteristic).
\item
  \textbf{Gửi lệnh khởi tạo ELM327}: Gửi chuỗi lệnh khởi tạo giao thức
  ELM327 qua BLE GATT write:

  \begin{itemize}
  \tightlist
  \item
    \texttt{ATZ} --- Reset adapter
  \item
    \texttt{ATE0} --- Tắt echo
  \item
    \texttt{ATL0} --- Tắt line feed
  \item
    \texttt{ATS0} --- Tắt khoảng trắng
  \item
    \texttt{ATSP0} --- Tự động phát hiện giao thức OBD2
  \end{itemize}
\item
  \textbf{Đọc dữ liệu OBD2}: Gửi các lệnh PID (Parameter ID) để đọc
  thông số xe:

  \begin{itemize}
  \tightlist
  \item
    \texttt{010C} --- RPM động cơ
  \item
    \texttt{010D} --- Tốc độ xe
  \item
    \texttt{0105} --- Nhiệt độ nước làm mát
  \item
    \texttt{012F} --- Mức nhiên liệu
  \item
    \texttt{AT\ IGN} --- Trạng thái động cơ (ignition)
  \end{itemize}
\end{enumerate}

\begin{Shaded}
\begin{Highlighting}[]
\CommentTok{/* Quy trình kết nối và đọc dữ liệu OBD2 */}
\DataTypeTok{void}\NormalTok{ ble\_obd2\_task}\OperatorTok{(}\DataTypeTok{void} \OperatorTok{*}\NormalTok{param}\OperatorTok{)} \OperatorTok{\{}
    \CommentTok{/* Bước 1: Đọc BLE address từ NVS */}
    \DataTypeTok{uint8\_t}\NormalTok{ saved\_addr}\OperatorTok{[}\DecValTok{6}\OperatorTok{];}
    \DataTypeTok{bool}\NormalTok{ has\_saved }\OperatorTok{=}\NormalTok{ nvs\_read\_ble\_address}\OperatorTok{(}\NormalTok{saved\_addr}\OperatorTok{);}

    \CommentTok{/* Bước 2: Kết nối hoặc scan */}
    \ControlFlowTok{if} \OperatorTok{(}\NormalTok{has\_saved}\OperatorTok{)} \OperatorTok{\{}
\NormalTok{        ble\_connect\_direct}\OperatorTok{(}\NormalTok{saved\_addr}\OperatorTok{);}  \CommentTok{/* Reconnect nhanh */}
    \OperatorTok{\}} \ControlFlowTok{else} \OperatorTok{\{}
\NormalTok{        ble\_start\_scan}\OperatorTok{(}\NormalTok{BLE\_SCAN\_TIMEOUT\_MS}\OperatorTok{);}
\NormalTok{        ble\_connect\_first\_obd2\_device}\OperatorTok{();}
\NormalTok{        nvs\_save\_ble\_address}\OperatorTok{(}\NormalTok{connected\_addr}\OperatorTok{);}
    \OperatorTok{\}}

    \CommentTok{/* Bước 3: Khởi tạo ELM327 */}
\NormalTok{    ble\_obd\_send\_command}\OperatorTok{(}\StringTok{"ATZ"}\OperatorTok{);}
\NormalTok{    ble\_obd\_send\_command}\OperatorTok{(}\StringTok{"ATE0"}\OperatorTok{);}
\NormalTok{    ble\_obd\_send\_command}\OperatorTok{(}\StringTok{"ATSP0"}\OperatorTok{);}

    \CommentTok{/* Bước 4: Vòng lặp đọc dữ liệu */}
    \ControlFlowTok{while} \OperatorTok{(}\NormalTok{device\_state }\OperatorTok{==}\NormalTok{ STATE\_DRIVING}\OperatorTok{)} \OperatorTok{\{}
\NormalTok{        obd2\_data\_t data}\OperatorTok{;}
\NormalTok{        data}\OperatorTok{.}\NormalTok{rpm   }\OperatorTok{=}\NormalTok{ ble\_obd\_read\_pid}\OperatorTok{(}\NormalTok{PID\_RPM}\OperatorTok{);}
\NormalTok{        data}\OperatorTok{.}\NormalTok{speed }\OperatorTok{=}\NormalTok{ ble\_obd\_read\_pid}\OperatorTok{(}\NormalTok{PID\_SPEED}\OperatorTok{);}
\NormalTok{        data}\OperatorTok{.}\NormalTok{temp  }\OperatorTok{=}\NormalTok{ ble\_obd\_read\_pid}\OperatorTok{(}\NormalTok{PID\_COOLANT\_TEMP}\OperatorTok{);}
\NormalTok{        data}\OperatorTok{.}\NormalTok{fuel  }\OperatorTok{=}\NormalTok{ ble\_obd\_read\_pid}\OperatorTok{(}\NormalTok{PID\_FUEL\_LEVEL}\OperatorTok{);}

\NormalTok{        xQueueSend}\OperatorTok{(}\NormalTok{telemetry\_queue}\OperatorTok{,} \OperatorTok{\&}\NormalTok{data}\OperatorTok{,}\NormalTok{ pdMS\_TO\_TICKS}\OperatorTok{(}\DecValTok{100}\OperatorTok{));}
\NormalTok{        vTaskDelay}\OperatorTok{(}\NormalTok{pdMS\_TO\_TICKS}\OperatorTok{(}\NormalTok{TRACKING\_INTERVAL\_MS}\OperatorTok{));}
    \OperatorTok{\}}
\OperatorTok{\}}
\end{Highlighting}
\end{Shaded}


{\thesissetfigureid{3.8}{fig-ch3-luu-o-thuat-toan-quy-trinh-ket-noi-va-oc-du-lieu}
\begin{figure}[htbp]
\centering
\pandocbounded{\safeincludesvg{./assets/figures/04-chuong-3-giai-phap-firmware-hinh-3-8.svg}}
\caption{Lưu đồ thuật toán quy trình kết nối và đọc dữ liệu BLE OBD2}
\label{fig:ch3-luu-o-thuat-toan-quy-trinh-ket-noi-va-oc-du-lieu}
\end{figure}
}


\subparagraph{c) Chiến lược kết nối theo chế độ hoạt
động}\label{c-chiux1ebfn-lux1b0ux1ee3c-kux1ebft-nux1ed1i-theo-chux1ebf-ux111ux1ed9-houx1ea1t-ux111ux1ed9ng-1}

Việc kết nối BLE OBD2 được tối ưu hóa theo từng chế độ hoạt động của
thiết bị để tiết kiệm năng lượng:

\textbf{Chế độ lái xe (IGN ON)}: ESP32-S3 duy trì kết nối BLE liên tục
với OBD2 adapter. Dữ liệu OBD2 được đọc định kỳ (mỗi 5--10 giây). Hệ
thống chỉ sử dụng light sleep (giữ BLE active) thay vì deep sleep.

\textbf{Chế độ đỗ xe (IGN OFF)}: Không kết nối OBD2 adapter. Trạng thái
IGN đã được xác nhận trước khi chuyển sang deep sleep. Việc không kết
nối BLE khi đỗ xe giúp tiết kiệm đáng kể thời gian wake-up và năng lượng
tiêu thụ.

\textbf{Chế độ cảnh báo (Motion Detected)}: Kết nối OBD2 là tùy chọn. Hệ
thống có thể chỉ sử dụng IMU và GPS để xác nhận chuyển động mà không cần
dữ liệu OBD2.

Đặc biệt, khi ESP32-S3 chuyển sang deep sleep, kết nối BLE bị ngắt hoàn
toàn. Mỗi lần wake-up, firmware cần thực hiện kết nối lại với adapter. Nhờ
có việc lưu địa chỉ MAC vào flash, thời gian kết nối lại chỉ mất khoảng
1--3 giây, nhanh hơn đáng kể so với pairing lần đầu (3--10 giây).

\subparagraph{d) Xử lý lỗi và cơ chế dự phòng
(Fallback)}\label{d-xux1eed-luxfd-lux1ed7i-vuxe0-cux1a1-chux1ebf-dux1ef1-phuxf2ng-fallback}

Module BLE OBD2 được thiết kế với nhiều lớp xử lý lỗi:


{\thesissettableid{3.2.2B}{tab-ch3-tong-hop-thong-tin-ve-tinh-huong-loi-xu-ly-thoi-gian}
\def\LTcaptype{table}
\begin{longtable}{|L{0.291\linewidth}|L{0.445\linewidth}|L{0.228\linewidth}|}
\caption{Tổng hợp thông tin về Tình huống lỗi, Xử lý, Thời gian timeout}\label{tab:ch3-tong-hop-thong-tin-ve-tinh-huong-loi-xu-ly-thoi-gian}\\
\hline \textbf{Tình huống lỗi} & \textbf{Xử lý} & \textbf{Thời gian timeout} \\
\\\hline
\endfirsthead
\caption[]{Tổng hợp thông tin về Tình huống lỗi, Xử lý, Thời gian timeout (tiếp theo)}\\
\hline \textbf{Tình huống lỗi} & \textbf{Xử lý} & \textbf{Thời gian timeout} \\
\\\hline
\endhead
\\\hline
\endfoot
\endlastfoot
Không kết nối được adapter & Retry 2--3 lần với delay 2 giây, sau đó
fallback & 10 giây \\\hline
Mất kết nối giữa chừng & Reconnect 2--3 lần, nếu thất bại thì fallback &
5 giây mỗi lần \\\hline
Adapter không phản hồi lệnh & Retry 1--2 lần, nếu thất bại thì fallback
& 5 giây \\\hline
\end{longtable}
}


\textbf{Cơ chế dự phòng (Fallback)}: Khi không thể giao tiếp với OBD2
adapter, hệ thống chuyển sang phát hiện trạng thái động cơ bằng phương
pháp đo điện áp ắc quy. Phương pháp này có độ chính xác thấp hơn OBD2
nhưng vẫn đảm bảo hệ thống hoạt động liên tục. Tiêu chí phân biệt theo
profile: profile 12V dùng \(IGN_{\mathrm{ON}} \ge 13.0\,\mathrm{V}\) và
\(IGN_{\mathrm{OFF}} \le 12.0\,\mathrm{V}\); profile 24V dùng
\(IGN_{\mathrm{ON}} \ge 26.0\,\mathrm{V}\) và
\(IGN_{\mathrm{OFF}} \le 24.0\,\mathrm{V}\).

Mọi lỗi kết nối OBD2 đều được ghi lại vào log để phục vụ việc giám sát
và khắc phục sự cố từ xa.

\paragraph{3.2.2.3. Module điều khiển modem
SIMCom}\label{3223-module-ux111iux1ec1u-khiux1ec3n-modem-simcom}

\subparagraph{a) Kiến trúc modem SIMCom SIM7600CE-T (LTE + GNSS tích
hợp)}\label{a-kiux1ebfn-truxfac-modem-simcom-sim7600ce-t-lte--gnss-tuxedch-hux1ee3p}

Trong kiến trúc hiện tại, modem SIMCom SIM7600CE-T đảm nhận toàn bộ
luồng LTE và GNSS. ESP32-S3 giao tiếp với modem qua UART1 (GPIO16 TX,
GPIO17 RX). Firmware sử dụng một chuỗi lệnh AT chính xác để đảm bảo
modem sẵn sàng trước khi bật PDP context và khởi động GNSS.


{\thesissettableid{3.2.2C}{tab-ch3-tong-hop-thong-tin-ve-buoc-lenh-at-muc-ich-chinh}
\def\LTcaptype{table}
\begin{longtable}{|L{0.160\linewidth}|L{0.224\linewidth}|L{0.580\linewidth}|}
\caption{Tổng hợp thông tin về Bước, Lệnh AT, Mục đích chính}\label{tab:ch3-tong-hop-thong-tin-ve-buoc-lenh-at-muc-ich-chinh}\\
\hline \textbf{Bước} & \textbf{Lệnh AT} & \textbf{Mục đích chính} \\
\\\hline
\endfirsthead
\caption[]{Tổng hợp thông tin về Bước, Lệnh AT, Mục đích chính (tiếp theo)}\\
\hline \textbf{Bước} & \textbf{Lệnh AT} & \textbf{Mục đích chính} \\
\\\hline
\endhead
\\\hline
\endfoot
\endlastfoot
1 & \texttt{AT} & Kiểm tra modem còn phản hồi \\\hline
2 & \texttt{AT+CPIN?} & Kiểm tra SIM card đã sẵn sàng \\\hline
3 & \texttt{AT+CNMP=2} & Đặt chế độ tự động LTE/UMTS/GSM \\\hline
4 & \nolinkurl{AT+CGDCONT=1,"IP","internet"} & Cấu hình PDP context với APN
mặc định \texttt{internet} \\\hline
5 & \texttt{AT+CEREG?} & Poll đến khi \texttt{+CEREG:\ 0,1} (hoặc
\texttt{stat}=5) trước khi kích hoạt PDP context \\\hline
6 & \texttt{AT+CGACT=1,1} & Bật PDP context sau khi đăng ký mạng thành
công \\\hline
7 & \texttt{AT+CGPADDR=1} & Ghi nhận địa chỉ IP được cấp (dùng để log và
kiểm tra) \\\hline
8 & \texttt{AT+CGNSPWR=1} & Bật GNSS tích hợp và đợi modem trả về fix vị
trí \\\hline
9 & \texttt{AT+CGNSINF} & Đọc tọa độ GNSS hiện tại (vị trí, tốc độ, số
vệ tinh); \texttt{AT+CGNSTST} chỉ dùng nếu cần dữ liệu NMEA hoàn toàn \\\hline
\end{longtable}
}


Toàn bộ lệnh AT cho LTE và GNSS được thực hiện trên cùng UART1. Trình tự
kiểm tra \texttt{AT+CEREG?} trước khi gọi \texttt{AT+CGACT=1,1} giúp
tránh kích hoạt PDP context khi modem chưa đăng ký mạng.

\subparagraph{b) Trình tự khởi tạo LTE/GNSS trong
firmware}\label{b-trinh-tux1ef1-khoi-tao-ltegnss-trong-firmware}

Ở tầng firmware, chuỗi lệnh trên không được gửi dồn trong một lần, mà
được tách thành các bước kiểm tra có điều kiện: xác nhận modem phản hồi,
xác nhận SIM sẵn sàng, chờ modem đăng ký mạng, sau đó mới bật PDP
context và GNSS. Cách tổ chức này giúp quá trình khởi tạo dễ retry hơn
và hạn chế lỗi dây chuyền khi điều kiện sóng di động thay đổi.

\subparagraph{c) Quản lý kết nối và tiết kiệm
điện}\label{c-quux1ea3n-luxfd-kux1ebft-nux1ed1i-vuxe0-tiux1ebft-kiux1ec7m-ux111iux1ec7n}

\begin{itemize}
\tightlist
\item
  \textbf{Khi thiết bị hoạt động liên tục}: Giữ PDP context mở và GNSS bật
  liên tục để gửi dữ liệu vận hành định kỳ và bản tin báo vẫn đang hoạt động.
\item
  \textbf{Khi dừng (parking mode):} Firmware đóng PDP context
  (\texttt{AT+CGACT=0,1}), đưa modem vào \texttt{AT+CSCLK=1} hoặc
  \texttt{AT+CFUN=0}, rồi chuyển ESP32-S3 sang deep sleep. Khi wake-up,
  modem được đánh thức qua UART và thực hiện lại các lệnh
  \texttt{AT+CFUN=1} + \texttt{AT+CGACT=1,1}.
\item
  \textbf{GNSS:} Sau khi \texttt{AT+CGNSPWR=1}, firmware tuần tự gọi
  \texttt{AT+CGNSINF} để lấy dữ liệu; \texttt{AT+CGNSTST} chỉ dùng khi
  cần stream NMEA để hỗ trợ debug hoặc phân tích phụ trợ.
\end{itemize}

\subparagraph{d) Xử lý sự cố}\label{d-xux1eed-luxfd-sux1ef1-cux1ed1}


{\thesissettableid{3.2.2D}{tab-ch3-tong-hop-thong-tin-ve-tinh-huong-bien-phap}
\def\LTcaptype{table}
\begin{longtable}{|L{0.220\linewidth}|L{0.740\linewidth}|}
\caption{Tổng hợp thông tin về Tình huống, Biện pháp}\label{tab:ch3-tong-hop-thong-tin-ve-tinh-huong-bien-phap}\\
\hline \textbf{Tình huống} & \textbf{Biện pháp} \\
\\\hline
\endfirsthead
\caption[]{Tổng hợp thông tin về Tình huống, Biện pháp (tiếp theo)}\\
\hline \textbf{Tình huống} & \textbf{Biện pháp} \\
\\\hline
\endhead
\\\hline
\endfoot
\endlastfoot
Modem không phản hồi & Reset nhanh qua PWR-KEY (GPIO26) rồi bắt đầu lại
chuá»—i AT \\\hline
Mất kết nối 4G & Đợi \texttt{AT+CEREG?} thành \texttt{1}, gọi
\texttt{AT+CGACT=0,1} rồi \texttt{AT+CGACT=1,1} để tái khởi động PDP \\\hline
GNSS không fix dù đã bật & Gọi lại \texttt{AT+CGNSINF}, nếu vẫn thất bại
thì bật lại GNSS (\texttt{AT+CGNSPWR=0}/\texttt{1}) rồi thử lại \\\hline
Modem quá nóng & Đưa \texttt{AT+CFUN=0} để tạm thời tắt RF và chờ nhiệt
độ giảm trước khi tiếp tục \\\hline
\end{longtable}
}


Các chế độ ngủ (\texttt{AT+CSCLK=1}, \texttt{AT+CFUN=0}) và wake-up bằng
UART giữ cho modem tiêu thụ thấp mà vẫn có thể khởi động lại nhanh khi
ESP32-S3 tỉnh dậy.


{\thesissetfigureid{3.9}{fig-ch3-luu-o-ieu-phoi-modem-sim7600ce-t-theo-trang-thai-mang-va}
\begin{figure}[htbp]
\centering
\pandocbounded{\safeincludesvg{./assets/figures/04-chuong-3-giai-phap-firmware-hinh-3-9.svg}}
\caption{Lưu đồ điều phối modem SIM7600CE-T theo trạng thái mạng và chế độ vận hành}
\label{fig:ch3-luu-o-ieu-phoi-modem-sim7600ce-t-theo-trang-thai-mang-va}
\end{figure}
}


\paragraph{3.2.2.4. Module quản lý nguồn và
GPIO}\label{3224-module-quux1ea3n-luxfd-nguux1ed3n-vuxe0-gpio}

\subparagraph{a) Sơ đồ chân
GPIO}\label{a-sux1a1-ux111ux1ed3-chuxe2n-gpio}

Module quản lý nguồn sử dụng các chân GPIO của ESP32-S3 để điều khiển và
giám sát hệ thống nguồn điện. Bảng sau liệt kê đầy đủ các chân GPIO được
sử dụng:


{\thesissettableid{3.2.2E}{tab-ch3-tong-hop-thong-tin-ve-gpio-chuc-nang-huong}
\def\LTcaptype{table}
\begin{longtable}{|L{0.115\linewidth}|L{0.230\linewidth}|L{0.165\linewidth}|L{0.445\linewidth}|}
\caption{Tổng hợp thông tin về GPIO, Chức năng, Hướng}\label{tab:ch3-tong-hop-thong-tin-ve-gpio-chuc-nang-huong}\\
\hline \textbf{GPIO} & \textbf{Chức năng} & \textbf{Hướng} & \textbf{Mô tả} \\
\\\hline
\endfirsthead
\caption[]{Tổng hợp thông tin về GPIO, Chức năng, Hướng (tiếp theo)}\\
\hline \textbf{GPIO} & \textbf{Chức năng} & \textbf{Hướng} & \textbf{Mô tả} \\
\\\hline
\endhead
\\\hline
\endfoot
\endlastfoot
2 & IGN\_IN & Input & Đọc trạng thái động cơ (GPIO hoặc OBD2) \\\hline
4 & \nolinkurl{U_BATT_ADC} & Input & Đọc điện áp ắc quy (ADC 12-bit) \\\hline
5 & \nolinkurl{CHARGER_EN} & Output & Điều khiển khối sạc TP5100 \\\hline
18 & \nolinkurl{POWER_PATH_EN} & Output & Chọn nguồn cấp (ắc quy hoặc pin dự
phòng) \\\hline
19 & LVD\_STATUS & Input & Đọc trạng thái LVD từ khối LVD phần cứng \\\hline
41 & \nolinkurl{LIS3DSH_INT1} & Input & Ngắt chính từ cảm biến gia tốc IMU \\\hline
42 & \nolinkurl{LIS3DSH_INT2} & Input & Ngắt phụ, đang giữ cho mở rộng \\\hline
2 & \nolinkurl{LIS3DSH_SDA} & I/O & Đường dữ liệu I2C \\\hline
1 & \nolinkurl{LIS3DSH_SCL} & I/O & Đường clock I2C \\\hline
16 & \nolinkurl{MODEM_UART_TX} & Output & UART TX đến modem \\\hline
17 & \nolinkurl{MODEM_UART_RX} & Input & UART RX từ modem \\\hline
26 & \nolinkurl{MODEM_PWR-KEY} & Output & Điều khiển nguồn modem \\\hline
\end{longtable}
}


\subparagraph{b) Điều khiển mạch chuyển nguồn (Diode OR +
EN)}\label{b-ux111iux1ec1u-khiux1ec3n-power-path-diode-or--en}

Hệ thống dùng mạch chuyển nguồn theo \textbf{diode OR} giữa nhánh 5V chính
(MP2482) và nhánh 5V dự phòng (SX1308). Firmware điều khiển qua GPIO18 để
ưu tiên nhánh nguồn phù hợp theo profile điện áp.

\begin{itemize}
\tightlist
\item
  \textbf{GPIO18 = LOW (0):} Ưu tiên nguồn ắc quy (nhánh MP2482 hoạt
  động)
\item
  \textbf{GPIO18 = HIGH (1):} Giảm/khóa nhánh chính và cho phép nhánh
  dự phòng qua SX1308 + diode OR
\end{itemize}

Logic điều khiển nguồn bám theo profile 12V/24V:

\begin{Shaded}
\begin{Highlighting}[]
\NormalTok{power\_status\_t power\_get\_status}\OperatorTok{(}\DataTypeTok{float}\NormalTok{ lvd\_threshold\_v}\OperatorTok{,} \DataTypeTok{float}\NormalTok{ lvd\_hysteresis\_v}\OperatorTok{)} \OperatorTok{\{}
\NormalTok{    power\_status\_t status }\OperatorTok{=} \OperatorTok{\{}
        \OperatorTok{.}\NormalTok{voltage }\OperatorTok{=}\NormalTok{ adc\_read\_battery\_voltage}\OperatorTok{(),}
        \OperatorTok{.}\NormalTok{is\_low }\OperatorTok{=} \KeywordTok{false}\OperatorTok{,}
        \OperatorTok{.}\NormalTok{charger\_on }\OperatorTok{=}\NormalTok{ gpio\_get\_level}\OperatorTok{(}\NormalTok{PIN\_CHARGER\_EN}\OperatorTok{)} \OperatorTok{==} \DecValTok{1}\OperatorTok{,}
    \OperatorTok{\};}

    \ControlFlowTok{if} \OperatorTok{(}\NormalTok{status}\OperatorTok{.}\NormalTok{voltage }\OperatorTok{\textless{}=}\NormalTok{ lvd\_threshold\_v}\OperatorTok{)} \OperatorTok{\{}
\NormalTok{        s\_voltage\_low\_latched }\OperatorTok{=} \KeywordTok{true}\OperatorTok{;}
    \OperatorTok{\}} \ControlFlowTok{else} \ControlFlowTok{if} \OperatorTok{(}\NormalTok{status}\OperatorTok{.}\NormalTok{voltage }\OperatorTok{\textgreater{}=}\NormalTok{ lvd\_hysteresis\_v}\OperatorTok{)} \OperatorTok{\{}
\NormalTok{        s\_voltage\_low\_latched }\OperatorTok{=} \KeywordTok{false}\OperatorTok{;}
    \OperatorTok{\}}

\NormalTok{    status}\OperatorTok{.}\NormalTok{is\_low }\OperatorTok{=}\NormalTok{ s\_voltage\_low\_latched }\OperatorTok{||}\NormalTok{ power\_is\_low\_voltage}\OperatorTok{();}
    \ControlFlowTok{return}\NormalTok{ status}\OperatorTok{;}
\OperatorTok{\}}
\end{Highlighting}
\end{Shaded}


{\thesissetfigureid{3.10}{fig-ch3-luu-o-thuat-toan-ieu-khien-power-path}
\begin{figure}[htbp]
\centering
\pandocbounded{\safeincludesvg{./assets/figures/04-chuong-3-giai-phap-firmware-hinh-3-10.svg}}
\caption{Lưu đồ thuật toán điều khiển chuyển nguồn}
\label{fig:ch3-luu-o-thuat-toan-ieu-khien-power-path}
\end{figure}
}


\subparagraph{c) Điều khiển khối sạc
(TP5100)}\label{c-ux111iux1ec1u-khiux1ec3n-ic-sux1ea1c-tp4056}

Khối sạc TP5100 được điều khiển qua chân GPIO5 (CHARGER\_EN). Logic sạc
được thiết kế để bảo vệ cả ắc quy xe lẫn pin dự phòng:


{\thesissettableid{3.2.2F}{tab-ch3-tong-hop-thong-tin-ve-ieu-kien-trang-thai-charger-ly-do}
\def\LTcaptype{table}
\begin{longtable}{|L{0.327\linewidth}|L{0.243\linewidth}|L{0.394\linewidth}|}
\caption{Tổng hợp thông tin về Điều kiện, Trạng thái charger, Lý do}\label{tab:ch3-tong-hop-thong-tin-ve-ieu-kien-trang-thai-charger-ly-do}\\
\hline \textbf{Điều kiện} & \textbf{Trạng thái charger} & \textbf{Lý do} \\
\\\hline
\endfirsthead
\caption[]{Tổng hợp thông tin về Điều kiện, Trạng thái charger, Lý do (tiếp theo)}\\
\hline \textbf{Điều kiện} & \textbf{Trạng thái charger} & \textbf{Lý do} \\
\\\hline
\endhead
\\\hline
\endfoot
\endlastfoot
IGN ON và \(U_{\text{batt}}\) vượt ngưỡng profile (ví dụ \(> 12\,\mathrm{V}\) cho hệ
12V) & Enable (GPIO5 = HIGH) & Máy phát điện đang nạp, có thể sạc pin dự
phòng \\\hline
IGN OFF & Disable (GPIO5 = LOW) & Bảo vệ ắc quy không bị hao pin khi xe
tắt máy \\\hline
\(U_{\text{batt}}\) dưới ngưỡng profile (ví dụ \(< 12\,\mathrm{V}\) cho hệ 12V) & Disable
(GPIO5 = LOW) & Ắc quy yếu, không đủ năng lượng để sạc pin dự phòng \\\hline
\end{longtable}
}


\subparagraph{d) Giám sát điện áp và
LVD}\label{d-giuxe1m-suxe1t-ux111iux1ec7n-uxe1p-vuxe0-lvd}

Firmware đọc điện áp ắc quy qua kênh ADC 12-bit (GPIO4) với bộ chia áp
để đưa điện áp ắc quy 12V hoặc 24V về dải đo của ADC (0--3.3V). Giá trị
ADC được chuyển đổi sang điện áp thực thông qua công thức hiệu chuẩn
(calibration).

Ngoài ra, hệ thống có thêm kênh giám sát LVD độc lập từ khối LVD phần
cứng (GPIO19). Đây là kênh dự phòng cho ADC, cho phép kiểm tra nhanh
trạng thái nguồn với cơ chế trễ (hysteresis): profile 12V dùng
\(Switch_{\mathrm{OFF}} = 12.0\,\mathrm{V}\),
\(Switch_{\mathrm{ON}} = 12.2\,\mathrm{V}\); profile 24V dùng
\(Switch_{\mathrm{OFF}} = 24.0\,\mathrm{V}\),
\(Switch_{\mathrm{ON}} = 24.4\,\mathrm{V}\). Khối LVD xử lý hysteresis ở mức phần cứng, tránh hiện
 tượng dao động (oscillation) quanh ngưỡng.

\subparagraph{e) Các chế độ quản lý
nguồn}\label{e-cuxe1c-chux1ebf-ux111ux1ed9-quux1ea3n-luxfd-nguux1ed3n}

Firmware hỗ trợ ba chế độ năng lượng chính:


{\thesissettableid{3.2.2G}{tab-ch3-tong-hop-thong-tin-ve-che-o-trang-thai-esp32-trang-thai}
\def\LTcaptype{table}
\begin{longtable}{|L{0.158\linewidth}|L{0.202\linewidth}|L{0.235\linewidth}|L{0.188\linewidth}|L{0.162\linewidth}|}
\caption{Tổng hợp thông tin về Chế độ, Trạng thái ESP32, Trạng thái modem}\label{tab:ch3-tong-hop-thong-tin-ve-che-o-trang-thai-esp32-trang-thai}\\
\hline \textbf{Chế độ} & \textbf{Trạng thái ESP32} & \textbf{Trạng thái modem} & \textbf{Trạng thái BLE} & \textbf{Dòng
tiêu thụ} \\
\\\hline
\endfirsthead
\caption[]{Tổng hợp thông tin về Chế độ, Trạng thái ESP32, Trạng thái modem (tiếp theo)}\\
\hline \textbf{Chế độ} & \textbf{Trạng thái ESP32} & \textbf{Trạng thái modem} & \textbf{Trạng thái BLE} & \textbf{Dòng
tiêu thụ} \\
\\\hline
\endhead
\\\hline
\endfoot
\endlastfoot
Xe đang chạy & Hoạt động đầy đủ & 4G + GNSS hoạt động & Đang kết nối & 150--250
mA \\\hline
Ngủ nhẹ & Ngủ nhẹ & Ngủ nhẹ của modem (CSCLK) & Giữ kết nối & 10--20 mA \\\hline
Ngủ sâu & Ngủ sâu & Ngủ sâu của modem (CFUN=0) & Ngắt kết nối &
\textless{} 2 mA \\\hline
\end{longtable}
}


Chuyển đổi giữa các chế độ dựa trên trạng thái IGN và dữ liệu cảm biến,
được điều phối bởi máy trạng thái chính (trình bày tại mục 3.2.2.5).

\paragraph{3.2.2.5. Định dạng dữ liệu và máy trạng
thái}\label{3225-ux111ux1ecbnh-dux1ea1ng-dux1eef-liux1ec7u-vuxe0-muxe1y-trux1ea1ng-thuxe1i}

\subparagraph{a) Định dạng dữ liệu
MQTT}\label{a-ux111ux1ecbnh-dux1ea1ng-dux1eef-liux1ec7u-mqtt}

Dữ liệu từ thiết bị được đóng gói theo định dạng JSON và truyền lên
broker qua nhóm topic \texttt{v1/\{device\_id\}/...}. Trong hệ thống
hiện tại, thiết bị phát ba nhóm gói dữ liệu chính là \texttt{rawdata},
\texttt{status} và \texttt{events}; ngoài ra còn có luồng
\texttt{firmware} để báo trạng thái cập nhật từ xa và luồng \texttt{commands} theo
chiều ngược lại từ máy chủ.

\textbf{Thông điệp Rawdata (gói dữ liệu vận hành chính)}

Đây là gói dữ liệu được gửi định kỳ khi thiết bị đang hoạt động, phản
ánh trạng thái GNSS, nguồn và cảm biến rung:

\begin{Shaded}
\begin{Highlighting}[]
\FunctionTok{\{}
  \DataTypeTok{"device\_id"}\FunctionTok{:} \StringTok{"TRACKER\_001"}\FunctionTok{,}
  \DataTypeTok{"auth\_token"}\FunctionTok{:} \StringTok{"device\_secret\_token"}\FunctionTok{,}
  \DataTypeTok{"timestamp"}\FunctionTok{:} \DecValTok{1710000000000}\FunctionTok{,}
  \DataTypeTok{"uptime"}\FunctionTok{:} \DecValTok{452130}\FunctionTok{,}
  \DataTypeTok{"data"}\FunctionTok{:} \FunctionTok{\{}
    \DataTypeTok{"vibration"}\FunctionTok{:} \DecValTok{128}\FunctionTok{,}
    \DataTypeTok{"battery\_top"}\FunctionTok{:} \FloatTok{12.54}\FunctionTok{,}
    \DataTypeTok{"battery\_bot"}\FunctionTok{:} \FloatTok{3.96}\FunctionTok{,}
    \DataTypeTok{"latitude"}\FunctionTok{:} \FloatTok{10.77653}\FunctionTok{,}
    \DataTypeTok{"longitude"}\FunctionTok{:} \FloatTok{106.70098}\FunctionTok{,}
    \DataTypeTok{"speed"}\FunctionTok{:} \FloatTok{42.5}\FunctionTok{,}
    \DataTypeTok{"course"}\FunctionTok{:} \FloatTok{181.2}\FunctionTok{,}
    \DataTypeTok{"satellites"}\FunctionTok{:} \DecValTok{11}\FunctionTok{,}
    \DataTypeTok{"ignition"}\FunctionTok{:} \KeywordTok{true}\FunctionTok{,}
    \DataTypeTok{"error\_code"}\FunctionTok{:} \DecValTok{0}
  \FunctionTok{\}}
\FunctionTok{\}}
\end{Highlighting}
\end{Shaded}

\textbf{Thông điệp Event (sự kiện/cảnh báo)}

Được gửi khi firmware phát hiện một sự kiện bất thường hoặc muốn ghi
nhận một mốc vận hành đáng chú ý:

\begin{Shaded}
\begin{Highlighting}[]
\FunctionTok{\{}
  \DataTypeTok{"device\_id"}\FunctionTok{:} \StringTok{"TRACKER\_001"}\FunctionTok{,}
  \DataTypeTok{"auth\_token"}\FunctionTok{:} \StringTok{"device\_secret\_token"}\FunctionTok{,}
  \DataTypeTok{"event\_type"}\FunctionTok{:} \StringTok{"high\_vibration"}\FunctionTok{,}
  \DataTypeTok{"code"}\FunctionTok{:} \DecValTok{201}\FunctionTok{,}
  \DataTypeTok{"message"}\FunctionTok{:} \StringTok{"Vehicle movement detected while parked"}\FunctionTok{,}
  \DataTypeTok{"timestamp"}\FunctionTok{:} \DecValTok{1710000005123}
\FunctionTok{\}}
\end{Highlighting}
\end{Shaded}

\textbf{Thông điệp Firmware (trạng thái OTA)}

Luồng \texttt{firmware} được dùng để ghi nhận vòng đời cập nhật firmware
từ lúc gán tác vụ đến khi xác nhận thành công sau reboot hoặc rollback:

\begin{Shaded}
\begin{Highlighting}[]
\FunctionTok{\{}
  \DataTypeTok{"device\_id"}\FunctionTok{:} \StringTok{"TRACKER\_001"}\FunctionTok{,}
  \DataTypeTok{"auth\_token"}\FunctionTok{:} \StringTok{"device\_secret\_token"}\FunctionTok{,}
  \DataTypeTok{"jobId"}\FunctionTok{:} \StringTok{"ota\_1710000000\_ab12cd"}\FunctionTok{,}
  \DataTypeTok{"status"}\FunctionTok{:} \StringTok{"confirming"}\FunctionTok{,}
  \DataTypeTok{"progress"}\FunctionTok{:} \DecValTok{99}\FunctionTok{,}
  \DataTypeTok{"targetVersion"}\FunctionTok{:} \StringTok{"1.2.0"}\FunctionTok{,}
  \DataTypeTok{"currentVersion"}\FunctionTok{:} \StringTok{"1.1.3"}\FunctionTok{,}
  \DataTypeTok{"partition"}\FunctionTok{:} \StringTok{"ota\_0"}
\FunctionTok{\}}
\end{Highlighting}
\end{Shaded}

\textbf{Thông điệp Status (trạng thái hoạt động)}

Được gửi khi thiết bị đổi trạng thái phiên làm việc hoặc phát heartbeat
tại các mốc cần theo dõi:

\begin{Shaded}
\begin{Highlighting}[]
\FunctionTok{\{}
  \DataTypeTok{"device\_id"}\FunctionTok{:} \StringTok{"TRACKER\_001"}\FunctionTok{,}
  \DataTypeTok{"auth\_token"}\FunctionTok{:} \StringTok{"device\_secret\_token"}\FunctionTok{,}
  \DataTypeTok{"status"}\FunctionTok{:} \StringTok{"running"}\FunctionTok{,}
  \DataTypeTok{"session\_id"}\FunctionTok{:} \DecValTok{1024}\FunctionTok{,}
  \DataTypeTok{"timestamp"}\FunctionTok{:} \DecValTok{1710000010000}
\FunctionTok{\}}
\end{Highlighting}
\end{Shaded}

\subparagraph{b) Thiết kế topic
MQTT}\label{b-thiux1ebft-kux1ebf-mqtt-topic}

Hệ thống sử dụng cấu trúc phân cấp topic MQTT để tách rõ luồng
dữ liệu vận hành gửi đều, trạng thái, sự kiện và cập nhật từ xa:


{\thesissettableid{3.2.2H}{tab-ch3-tong-hop-thong-tin-ve-topic-huong-qos}
\def\LTcaptype{table}
\begin{longtable}{|L{0.254\linewidth}|L{0.177\linewidth}|L{0.115\linewidth}|L{0.409\linewidth}|}
\caption{Tổng hợp thông tin về Topic, Hướng, QoS}\label{tab:ch3-tong-hop-thong-tin-ve-topic-huong-qos}\\
\hline \textbf{Topic} & \textbf{Hướng} & \textbf{QoS} & \textbf{Mô tả} \\
\\\hline
\endfirsthead
\caption[]{Tổng hợp thông tin về Topic, Hướng, QoS (tiếp theo)}\\
\hline \textbf{Topic} & \textbf{Hướng} & \textbf{QoS} & \textbf{Mô tả} \\
\\\hline
\endhead
\\\hline
\endfoot
\endlastfoot
\nolinkurl{v1/{device_id}/rawdata} & Device → Server & 0 &
Gói dữ liệu vận hành chính; gửi dày, chấp nhận mất vài mẫu \\\hline
\nolinkurl{v1/{device_id}/status} & Device → Server & 1 &
Heartbeat, trạng thái phiên, online/offline \\\hline
\nolinkurl{v1/{device_id}/events} & Device → Server & 1 &
Cảnh báo và sự kiện quan trọng \\\hline
\nolinkurl{v1/{device_id}/firmware} & Device → Server & 1
& Tiến trình cập nhật từ xa và kết quả cập nhật firmware \\\hline
\nolinkurl{v1/{device_id}/commands} & Server → Device & 1
& Lệnh điều khiển và cập nhật cấu hình từ máy chủ \\\hline
\end{longtable}
}


Các lệnh điều khiển từ máy chủ hiện được firmware xử lý gồm:


{\thesissettableid{3.2.2I}{tab-ch3-tong-hop-thong-tin-ve-lenh-muc-ich-tham-so-chinh}
\def\LTcaptype{table}
\begin{longtable}{|L{0.255\linewidth}|L{0.315\linewidth}|L{0.394\linewidth}|}
\caption{Tổng hợp thông tin về Lệnh, Mục đích, Tham số chính}\label{tab:ch3-tong-hop-thong-tin-ve-lenh-muc-ich-tham-so-chinh}\\
\hline \textbf{Lệnh} & \textbf{Mục đích} & \textbf{Tham số chính} \\
\\\hline
\endfirsthead
\caption[]{Tổng hợp thông tin về Lệnh, Mục đích, Tham số chính (tiếp theo)}\\
\hline \textbf{Lệnh} & \textbf{Mục đích} & \textbf{Tham số chính} \\
\\\hline
\endhead
\\\hline
\endfoot
\endlastfoot
\texttt{update\_config} & Cập nhật cấu hình thiết bị &
\texttt{heartbeat\_interval\_s}, \texttt{tracking\_interval\_s} \\\hline
\texttt{request\_location} & Yêu cầu gửi vị trí ngay ở chu kỳ kế tiếp &
Không \\\hline
\texttt{enable\_tracking} & Bật hoặc tắt phát dữ liệu vận hành định kỳ &
\texttt{enabled} (boolean) \\\hline
\texttt{reboot} & Khởi động lại thiết bị & Không \\\hline
\texttt{ota\_update} & Kích hoạt cập nhật firmware OTA & \texttt{jobId},
\texttt{version}, \texttt{url}, \texttt{size}, \texttt{sha256},
\texttt{force}, \texttt{confirmTimeoutSec} \\\hline
\texttt{manual\_rollback} / \texttt{ota\_rollback} & Yêu cầu quay về
firmware trước đó & Không bắt buộc \\\hline
\end{longtable}
}


\subparagraph{c) Chiến lược xử lý mất kết
nối}\label{c-chiux1ebfn-lux1b0ux1ee3c-xux1eed-luxfd-mux1ea5t-kux1ebft-nux1ed1i}

Ở phạm vi triển khai hiện tại, firmware kết hợp hai cơ chế song song:
(1) tự phục hồi kết nối MQTT và (2) hàng đợi cục bộ trên \textbf{microSD
SDMMC 4-bit} để giữ dữ liệu khi mất mạng. Luồng dữ liệu được hiện thực
qua cặp module \texttt{offline\_queue} và \texttt{sd\_log\_store}: bản
ghi được đưa vào hàng đợi khi dữ liệu vận hành phát sinh
(\texttt{offline\_queue\_enqueue()}), ghi xuống FATFS
(\texttt{sd\_log\_store\_append()}), sau đó replay theo thứ tự sequence
khi mạng phục hồi (\texttt{offline\_queue\_replay\_tick()}) và chốt ACK
cho bản ghi quan trọng
(\texttt{offline\_queue\_handle\_publish\_ack()}).

Thiết kế này cho phép tách bản ghi thành hai nhóm: \textbf{quan trọng}
(QoS 1, yêu cầu xác nhận đã nhận) và \textbf{ít quan trọng hơn} (QoS 0, ưu tiên
tốc độ gửi/ghi), đồng thời áp dụng cơ chế giới hạn dung lượng + dọn bản ghi cũ
(\texttt{sd\_log\_store\_gc\_if\_needed()}) để tránh đầy thẻ. Trong phạm
vi vận hành hiện tại, hệ thống giả định thẻ SD ở trạng thái
\textbf{đã cắm sẵn từ lúc khởi động}; luồng rút/cắm nóng trong lúc đang chạy
không phải phạm vi mục tiêu.

\begin{Shaded}
\begin{Highlighting}[]
\CommentTok{/* Chiến lược hiện tại: enqueue local + replay khi MQTT sẵn sàng */}
\NormalTok{offline\_queue\_enqueue}\OperatorTok{(}\NormalTok{type}\OperatorTok{,}\NormalTok{ payload}\OperatorTok{,}\NormalTok{ gps\_fix}\OperatorTok{,}\NormalTok{ net\_up}\OperatorTok{);}

\ControlFlowTok{if} \OperatorTok{(}\NormalTok{tracker\_mqtt\_is\_connected}\OperatorTok{())} \OperatorTok{\{}
\NormalTok{    offline\_queue\_replay\_tick}\OperatorTok{();}
\OperatorTok{\}}
\end{Highlighting}
\end{Shaded}

\subparagraph{d) Máy trạng thái thiết bị (Device State
Machine)}\label{d-muxe1y-trux1ea1ng-thuxe1i-thiux1ebft-bux1ecb-device-state-machine}

Máy trạng thái là cơ chế điều phối trung tâm của firmware, quyết định
hành vi thiết bị tại từng thời điểm vận hành. Hệ thống định nghĩa bảy
trạng thái chính:


{\thesissettableid{3.2.2J}{tab-ch3-tong-hop-thong-tin-ve-trang-thai-ma-mo-ta}
\def\LTcaptype{table}
\begin{longtable}{|L{0.224\linewidth}|L{0.160\linewidth}|L{0.580\linewidth}|}
\caption{Tổng hợp thông tin về Trạng thái, Mã, Mô tả}\label{tab:ch3-tong-hop-thong-tin-ve-trang-thai-ma-mo-ta}\\
\hline \textbf{Trạng thái} & \textbf{Mã} & \textbf{Mô tả} \\
\\\hline
\endfirsthead
\caption[]{Tổng hợp thông tin về Trạng thái, Mã, Mô tả (tiếp theo)}\\
\hline \textbf{Trạng thái} & \textbf{Mã} & \textbf{Mô tả} \\
\\\hline
\endhead
\\\hline
\endfoot
\endlastfoot
INIT & 0 & Khởi tạo hệ thống, cấu hình ngoại vi \\\hline
CHECK\_IGN & 1 & Kiểm tra trạng thái động cơ \\\hline
DRIVING & 2 & Chế độ lái xe - theo dõi liên tục \\\hline
PARKED & 3 & Chế độ đỗ xe - tiết kiệm năng lượng \\\hline
ALARM & 4 & Chế độ cảnh báo - phát hiện chuyển động bất thường \\\hline
HEARTBEAT & 5 & Gửi tín hiệu heartbeat \\\hline
SLEEP & 6 & Ngủ sâu - tiêu thụ năng lượng tối thiểu \\\hline
\end{longtable}
}


\textbf{Bảng chuyển đổi trạng thái:}

{\def\LTcaptype{none} % do not increment counter
\begin{longtable}{|L{0.265\linewidth}|L{0.432\linewidth}|L{0.267\linewidth}|}
\hline \textbf{Trạng thái hiện tại} & \textbf{Sự kiện kích hoạt} & \textbf{Trạng thái tiếp theo} \\
\\\hline
\endhead
\\\hline
\endfoot
\endlastfoot
INIT & Khởi tạo hoàn tất & CHECK\_IGN \\\hline
CHECK\_IGN & IGN = ON & DRIVING \\\hline
CHECK\_IGN & IGN = OFF & PARKED \\\hline
DRIVING & IGN = OFF được phát hiện & PARKED \\\hline
PARKED & IMU interrupt (motion detected) & ALARM \\\hline
PARKED & Timer wake-up & HEARTBEAT \\\hline
ALARM & Motion dừng, IGN = OFF & PARKED \\\hline
ALARM & IGN = ON & DRIVING \\\hline
HEARTBEAT & Gá»­i heartbeat xong & SLEEP \\\hline
SLEEP & Timer wake-up hoặc IMU interrupt & CHECK\_IGN \\\hline
\end{longtable}
}


{\thesissetfigureid{3.11}{fig-ch3-so-o-may-trang-thai-cua-thiet-bi-theo-doi}
\begin{figure}[htbp]
\centering
\pandocbounded{\safeincludesvg{./assets/figures/04-chuong-3-giai-phap-firmware-hinh-3-11.svg}}
\caption{Sơ đồ máy trạng thái của thiết bị theo dõi}
\label{fig:ch3-so-o-may-trang-thai-cua-thiet-bi-theo-doi}
\end{figure}
}


\begin{Shaded}
\begin{Highlighting}[]
\CommentTok{/* Cấu trúc máy trạng thái */}
\KeywordTok{typedef} \KeywordTok{enum} \OperatorTok{\{}
\NormalTok{    STATE\_INIT}\OperatorTok{,}
\NormalTok{    STATE\_CHECK\_IGN}\OperatorTok{,}
\NormalTok{    STATE\_DRIVING}\OperatorTok{,}
\NormalTok{    STATE\_PARKED}\OperatorTok{,}
\NormalTok{    STATE\_ALARM}\OperatorTok{,}
\NormalTok{    STATE\_HEARTBEAT}\OperatorTok{,}
\NormalTok{    STATE\_SLEEP}
\OperatorTok{\}}\NormalTok{ device\_state\_t}\OperatorTok{;}

\DataTypeTok{void}\NormalTok{ state\_machine\_task}\OperatorTok{(}\DataTypeTok{void} \OperatorTok{*}\NormalTok{param}\OperatorTok{)} \OperatorTok{\{}
\NormalTok{    device\_state\_t state }\OperatorTok{=}\NormalTok{ STATE\_INIT}\OperatorTok{;}

    \ControlFlowTok{while} \OperatorTok{(}\DecValTok{1}\OperatorTok{)} \OperatorTok{\{}
        \ControlFlowTok{switch} \OperatorTok{(}\NormalTok{state}\OperatorTok{)} \OperatorTok{\{}
        \ControlFlowTok{case}\NormalTok{ STATE\_INIT}\OperatorTok{:}
\NormalTok{            init\_peripherals}\OperatorTok{();}
\NormalTok{            state }\OperatorTok{=}\NormalTok{ STATE\_CHECK\_IGN}\OperatorTok{;}
            \ControlFlowTok{break}\OperatorTok{;}

        \ControlFlowTok{case}\NormalTok{ STATE\_CHECK\_IGN}\OperatorTok{:}
            \ControlFlowTok{if} \OperatorTok{(}\NormalTok{check\_ignition\_status}\OperatorTok{())} \OperatorTok{\{}
\NormalTok{                state }\OperatorTok{=}\NormalTok{ STATE\_DRIVING}\OperatorTok{;}
            \OperatorTok{\}} \ControlFlowTok{else} \OperatorTok{\{}
\NormalTok{                state }\OperatorTok{=}\NormalTok{ STATE\_PARKED}\OperatorTok{;}
            \OperatorTok{\}}
            \ControlFlowTok{break}\OperatorTok{;}

        \ControlFlowTok{case}\NormalTok{ STATE\_DRIVING}\OperatorTok{:}
\NormalTok{            start\_continuous\_tracking}\OperatorTok{();}
            \CommentTok{/* Chuyển sang PARKED khi phát hiện IGN OFF */}
\NormalTok{            EventBits\_t bits }\OperatorTok{=}\NormalTok{ xEventGroupWaitBits}\OperatorTok{(}
\NormalTok{                system\_event\_group}\OperatorTok{,}\NormalTok{ EVT\_IGN\_OFF}\OperatorTok{,}
\NormalTok{                pdTRUE}\OperatorTok{,}\NormalTok{ pdFALSE}\OperatorTok{,}\NormalTok{ pdMS\_TO\_TICKS}\OperatorTok{(}\DecValTok{1000}\OperatorTok{));}
            \ControlFlowTok{if} \OperatorTok{(}\NormalTok{bits }\OperatorTok{\&}\NormalTok{ EVT\_IGN\_OFF}\OperatorTok{)} \OperatorTok{\{}
\NormalTok{                stop\_ble\_connection}\OperatorTok{();}
\NormalTok{                state }\OperatorTok{=}\NormalTok{ STATE\_PARKED}\OperatorTok{;}
            \OperatorTok{\}}
            \ControlFlowTok{break}\OperatorTok{;}

        \ControlFlowTok{case}\NormalTok{ STATE\_PARKED}\OperatorTok{:}
\NormalTok{            prepare\_for\_sleep}\OperatorTok{();}
\NormalTok{            state }\OperatorTok{=}\NormalTok{ STATE\_SLEEP}\OperatorTok{;}
            \ControlFlowTok{break}\OperatorTok{;}

        \ControlFlowTok{case}\NormalTok{ STATE\_ALARM}\OperatorTok{:}
\NormalTok{            send\_alert\_immediate}\OperatorTok{();}
\NormalTok{            start\_continuous\_tracking}\OperatorTok{();}
            \CommentTok{/* Đợi cho đến khi hết chuyển động */}
            \ControlFlowTok{if} \OperatorTok{(!}\NormalTok{is\_motion\_detected}\OperatorTok{()} \OperatorTok{\&\&} \OperatorTok{!}\NormalTok{check\_ignition\_status}\OperatorTok{())} \OperatorTok{\{}
\NormalTok{                state }\OperatorTok{=}\NormalTok{ STATE\_PARKED}\OperatorTok{;}
            \OperatorTok{\}}
            \ControlFlowTok{break}\OperatorTok{;}

        \ControlFlowTok{case}\NormalTok{ STATE\_HEARTBEAT}\OperatorTok{:}
\NormalTok{            wakeup\_modem}\OperatorTok{();}
\NormalTok{            tracker\_mqtt\_connect}\OperatorTok{();}
\NormalTok{            send\_heartbeat}\OperatorTok{();}
\NormalTok{            state }\OperatorTok{=}\NormalTok{ STATE\_SLEEP}\OperatorTok{;}
            \ControlFlowTok{break}\OperatorTok{;}

        \ControlFlowTok{case}\NormalTok{ STATE\_SLEEP}\OperatorTok{:}
\NormalTok{            shutdown\_modem}\OperatorTok{();}
\NormalTok{            configure\_wakeup\_sources}\OperatorTok{();}
\NormalTok{            enter\_deep\_sleep}\OperatorTok{();}
            \CommentTok{/* Sau khi wake{-}up, bắt đầu lại từ CHECK\_IGN */}
\NormalTok{            state }\OperatorTok{=}\NormalTok{ STATE\_CHECK\_IGN}\OperatorTok{;}
            \ControlFlowTok{break}\OperatorTok{;}
        \OperatorTok{\}}
    \OperatorTok{\}}
\OperatorTok{\}}
\end{Highlighting}
\end{Shaded}

\textbf{Lưu trạng thái qua deep sleep}: Trước khi vào deep sleep,
firmware lưu trạng thái hiện tại vào vùng nhớ RTC (RTC memory) --- vùng
nhớ đặc biệt của ESP32-S3 được giữ nguyên nội dung trong suốt quá trình
deep sleep. Khi wake-up, firmware đọc trạng thái từ RTC memory để biết
chế độ hoạt động trước đó và xử lý phù hợp.

\paragraph{3.2.2.6. Cấu hình hệ
thống}\label{3226-cux1ea5u-huxecnh-hux1ec7-thux1ed1ng}

\subparagraph{a) Cấu trúc dữ liệu cấu
hình}\label{a-cux1ea5u-truxfac-dux1eef-liux1ec7u-cux1ea5u-huxecnh}

Toàn bộ cấu hình của thiết bị được lưu trữ trong bộ nhớ flash (NVS ---
Non-Volatile Storage) của ESP32-S3 dưới dạng cấu trúc dữ liệu cố định:

\begin{Shaded}
\begin{Highlighting}[]
\KeywordTok{typedef} \KeywordTok{struct} \OperatorTok{\{}
    \DataTypeTok{char}\NormalTok{     device\_id}\OperatorTok{[}\DecValTok{16}\OperatorTok{];}          \CommentTok{/* Mã định danh thiết bị */}
    \DataTypeTok{char}\NormalTok{     mqtt\_broker}\OperatorTok{[}\DecValTok{64}\OperatorTok{];}        \CommentTok{/* Địa chỉ MQTT broker */}
    \DataTypeTok{uint16\_t}\NormalTok{ mqtt\_port}\OperatorTok{;}              \CommentTok{/* Cổng MQTT (mặc định 1883) */}
    \DataTypeTok{char}\NormalTok{     mqtt\_username}\OperatorTok{[}\DecValTok{32}\OperatorTok{];}      \CommentTok{/* Tên đăng nhập MQTT */}
    \DataTypeTok{char}\NormalTok{     mqtt\_password}\OperatorTok{[}\DecValTok{32}\OperatorTok{];}      \CommentTok{/* Mật khẩu MQTT */}
    \DataTypeTok{uint16\_t}\NormalTok{ heartbeat\_interval}\OperatorTok{;}     \CommentTok{/* Chu kỳ gửi heartbeat (giây) */}
    \DataTypeTok{uint16\_t}\NormalTok{ tracking\_interval}\OperatorTok{;}      \CommentTok{/* Chu kỳ gửi telemetry (giây) */}
    \DataTypeTok{char}\NormalTok{     obd2\_ble\_address}\OperatorTok{[}\DecValTok{18}\OperatorTok{];}   \CommentTok{/* Địa chỉ MAC của OBD2 adapter */}
    \DataTypeTok{float}\NormalTok{    switch\_off}\OperatorTok{;}             \CommentTok{/* Ngưỡng chuyển sang pin dự phòng: 12.0V (12V) hoặc 24.0V (24V) */}
    \DataTypeTok{float}\NormalTok{    switch\_on}\OperatorTok{;}              \CommentTok{/* Ngưỡng quay lại ắc quy: 12.2V (12V) hoặc 24.4V (24V) */}
    \DataTypeTok{float}\NormalTok{    ign\_on\_threshold}\OperatorTok{;}       \CommentTok{/* Ngưỡng suy luận IGN ON: \textgreater{}=13.0V (12V) hoặc \textgreater{}=26.0V (24V) */}
    \DataTypeTok{float}\NormalTok{    ign\_off\_threshold}\OperatorTok{;}      \CommentTok{/* Ngưỡng suy luận IGN OFF: \textless{}=12.0V (12V) hoặc \textless{}=24.0V (24V) */}
\OperatorTok{\}}\NormalTok{ config\_t}\OperatorTok{;}
\end{Highlighting}
\end{Shaded}

\subparagraph{b) Quản lý cấu
hình}\label{b-quux1ea3n-luxfd-cux1ea5u-huxecnh}

Cấu hình được quản lý theo ba cơ chế:

\textbf{Cấu hình mặc định (Default Configuration)}: Được định nghĩa
trong mã nguồn firmware, áp dụng khi thiết bị khởi động lần đầu hoặc khi
NVS bị xóa. Bao gồm các giá trị an toàn như
\texttt{heartbeat\_interval\ =\ 900} (15 phút),
\texttt{tracking\_interval\ =\ 10} (10 giây), và profile nguồn kép:
profile 12V (\texttt{switch\_off=12.0}, \texttt{switch\_on=12.2},
\texttt{ign\_on\_threshold=13.0}, \texttt{ign\_off\_threshold=12.0})
hoặc profile 24V (\texttt{switch\_off=24.0}, \texttt{switch\_on=24.4},
\texttt{ign\_on\_threshold=26.0}, \texttt{ign\_off\_threshold=24.0}).

\textbf{Cấu hình lưu trữ (Persistent Configuration)}: Lưu trong NVS,
được tải khi thiết bị khởi động. Các thay đổi cấu hình từ máy chủ được
lưu vào NVS để giữ nguyên sau khi reset hoặc mất điện.

\textbf{Cập nhật cấu hình từ xa (Remote Configuration Update)}: Máy chủ
gửi lệnh \texttt{update\_config} qua topic MQTT
\texttt{v1/\{device\_id\}/commands}. Firmware phân tích lệnh, cập nhật
các tham số tương ứng và lưu vào NVS. Thiết bị gửi xác nhận (ACK) về máy
chủ sau khi cập nhật thành công.

\subparagraph{c) Hiệu chuẩn cảm biến (Sensor
Calibration)}\label{c-hiux1ec7u-chuux1ea9n-cux1ea3m-biux1ebfn-sensor-calibration}

Để đảm bảo độ chính xác của phép đo, firmware hỗ trợ hiệu chuẩn ba loại
cảm biến:


{\thesissettableid{3.2.2K}{tab-ch3-tong-hop-thong-tin-ve-cam-bien-phuong-phap-hieu-chuan-luu}
\def\LTcaptype{table}
\begin{longtable}{|L{0.224\linewidth}|L{0.580\linewidth}|L{0.160\linewidth}|}
\caption{Tổng hợp thông tin về Cảm biến, Phương pháp hiệu chuẩn, Lưu trữ}\label{tab:ch3-tong-hop-thong-tin-ve-cam-bien-phuong-phap-hieu-chuan-luu}\\
\hline \textbf{Cảm biến} & \textbf{Phương pháp hiệu chuẩn} & \textbf{Lưu trữ} \\
\\\hline
\endfirsthead
\caption[]{Tổng hợp thông tin về Cảm biến, Phương pháp hiệu chuẩn, Lưu trữ (tiếp theo)}\\
\hline \textbf{Cảm biến} & \textbf{Phương pháp hiệu chuẩn} & \textbf{Lưu trữ} \\
\\\hline
\endhead
\\\hline
\endfoot
\endlastfoot
ADC (điện áp ắc quy) & So sánh với đồng hồ vạn năng, tính hệ số hiệu
chỉnh & NVS \\\hline
IMU (LIS3DSH) & Đặt thiết bị trên mặt phẳng, đo offset các trục X/Y/Z &
NVS \\\hline
GNSS & Hiệu chỉnh offset vị trí (nếu cần) & NVS \\\hline
\end{longtable}
}


Các giá trị hiệu chuẩn được lưu trữ trong NVS và được tải khi thiết bị
khởi động, đảm bảo tính nhất quán của phép đo giữa các lần reset.

\subsection{3.3. Phân tích và đề xuất kiến trúc máy chủ và xử lý dữ
liệu -- Server and data-processing architecture}\label{33-phan-tich-va-de-xuat-kien-truc-backend-cloud}

\subsubsection{3.3.1. Phân tích và lựa chọn kiến trúc máy chủ}

\paragraph{3.3.1.1. Đặt vấn đề cho kiến trúc máy
chủ}\label{3131-ux111ux1eb7t-vux1ea5n-ux111ux1ec1-cho-kiux1ebfn-truxfac-cloud}

Khi dữ liệu rời khỏi thiết bị, trọng tâm không còn là ``đọc được dữ
liệu'' mà là ``biến bản tin đó thành thông tin hữu ích''. Ở tầng
máy chủ và xử lý dữ liệu, một bản tin chỉ thật sự có giá trị khi nó được tiếp nhận
ổn định, lưu đúng nơi, đẩy tới giao diện quản trị đúng lúc và vẫn truy vết được
khi có sự cố. Vì vậy, kiến trúc máy chủ phải đồng thời giải quyết các bài
toán sau:

\begin{itemize}
\tightlist
\item
  \textbf{Tiếp nhận dữ liệu liên tục từ nhiều thiết bị}: dữ liệu vận hành
  gửi theo chu kỳ ngắn, có thể dồn thành cụm khi thiết bị kết nối lại.
\item
  \textbf{Cập nhật thời gian thực cho giao diện quản trị}: dữ liệu vị trí/trạng
  thái cần phản ánh gần như tức thời cho người vận hành.
\item
  \textbf{Lưu trữ dữ liệu hỗn hợp}: vừa có dữ liệu chuỗi thời gian (GPS,
  nguồn, IMU; sẵn sàng mở rộng thêm OBD2), vừa có dữ liệu quan hệ nghiệp
  vụ (xe, người dùng, cảnh báo, vùng cho phép).
\item
  \textbf{Chi phí và năng lực vận hành phù hợp đồ án}: ưu tiên
  tự triển khai trên máy chủ của nhóm, dễ dựng và dễ mở rộng theo từng dịch vụ.
\end{itemize}

\paragraph{3.3.1.2. So sánh các phương án kiến trúc
hệ thống máy chủ}\label{3132-so-suxe1nh-cuxe1c-phux1b0ux1a1ng-uxe1n-kiux1ebfn-truxfac-cloud}


{\thesissettableid{3.15A}{tab-ch3-so-sanh-cac-phuong-an-kien-truc-cloud}
\def\LTcaptype{table}
\begin{longtable}{|L{0.145\linewidth}|L{0.230\linewidth}|L{0.220\linewidth}|L{0.250\linewidth}|L{0.100\linewidth}|}
\caption{So sánh các phương án kiến trúc hệ thống máy chủ}\label{tab:ch3-so-sanh-cac-phuong-an-kien-truc-cloud}\\
\hline \textbf{Phương án} & \textbf{Mô tả ngắn} & \textbf{Ưu điểm} & \textbf{Hạn chế} & \textbf{Mức phù hợp} \\
\\\hline
\endfirsthead
\caption[]{So sánh các phương án kiến trúc hệ thống máy chủ (tiếp theo)}\\
\hline \textbf{Phương án} & \textbf{Mô tả ngắn} & \textbf{Ưu điểm} & \textbf{Hạn chế} & \textbf{Mức phù hợp} \\
\\\hline
\endhead
\\\hline
\endfoot
\endlastfoot
\textbf{PA-A: Một API lớn + một cơ sở dữ liệu chung} & Thiết bị gửi HTTP
trực tiếp vào API phía máy chủ và mọi dữ liệu dồn vào cùng một nơi lưu &
Triển khai ban đầu đơn giản & Không tối ưu cho dữ liệu gửi dày theo thời
gian, khó mở rộng cập nhật thời gian thực, tải dễ dồn lên một dịch vụ &
Trung bình \\\hline
\textbf{PA-B: Dùng dịch vụ máy chủ do nhà cung cấp quản lý sẵn} & AWS IoT
Core + Lambda + cơ sở dữ liệu do nhà cung cấp vận hành & Khả năng mở
rộng cao, nhiều dịch vụ sẵn có & Chi phí vận hành cao, phụ thuộc nhà
cung cấp, độ phức tạp hạ tầng vượt phạm vi đồ án & Trung bình \\\hline
\textbf{PA-C: Kiến trúc hướng sự kiện, tự triển khai} & MQTT Broker +
MQTT Bridge + hai lớp lưu trữ dữ liệu + API + WebSocket & Phù hợp với
IoT cần cập nhật nhanh, tách tải tốt, mở rộng theo từng dịch vụ và tự
chủ chi phí & Cần chuẩn hóa luồng bản tin và quản lý nhiều dịch vụ Docker &
\textbf{Cao (Đã chọn)} \\\hline
\end{longtable}
}


\paragraph{3.3.1.3. Chọn giải pháp kiến trúc
hệ thống máy chủ}\label{3133-chux1ecdn-giux1ea3i-phuxe1p-kiux1ebfn-truxfac-cloud}

Đồ án chọn \textbf{PA-C: kiến trúc phân tầng kết hợp phản ứng theo sự kiện},
tức khi có bản tin mới thì từng lớp liên quan xử lý ngay, với chuỗi xử lý:

\texttt{Device\ -\textgreater{}\ EMQX\ -\textgreater{}\ MQTT\ Bridge\ -\textgreater{}\ PostgreSQL/}\\
\texttt{VictoriaMetrics/VictoriaLogs\ -\textgreater{}\ API\ phía\ máy\ chủ\ -\textgreater{}\ giao\ diện\ web}

Lựa chọn này xuất phát từ cách dữ liệu thực sự di chuyển trong hệ thống:
thiết bị gửi bản tin liên tục, lớp tiếp nhận dữ liệu đầu vào tiếp nhận
và phân luồng, lớp
lưu trữ ghi theo đúng đặc tính dữ liệu, còn lớp ứng dụng phục vụ người
dùng và nghiệp vụ. Nói cách khác, kiến trúc được chọn vì nó bám sát
hành trình tự nhiên của dữ liệu thay vì dồn mọi trách nhiệm về một dịch
vụ duy nhất.

Các lý do lựa chọn chính gồm:

\begin{itemize}
\tightlist
\item
  \textbf{Khớp bản chất dữ liệu IoT}: MQTT xử lý tốt mô hình pub/sub,
  tức mô hình ``một bên phát bản tin vào kênh chung, nhiều bên có thể
  nhận'', nên giảm mức phụ thuộc trực tiếp giữa thiết bị và tầng ứng
  dụng. Điều này giúp thiết bị chỉ tập trung gửi dữ liệu, thay vì phải
  biết chi tiết ai sẽ nhận và xử lý dữ liệu đó ở phía sau.
\item
  \textbf{Tối ưu hiệu năng theo vai trò}: Bridge xử lý luồng dữ liệu đi
  vào đầu tiên, API phía máy chủ tập trung vào nghiệp vụ, còn cơ sở dữ
  liệu được tách theo loại dữ liệu. Mỗi lớp chỉ làm
  đúng việc mình mạnh nhất, nhờ đó hệ thống dễ giữ ổn định hơn khi tải
  tăng.
\item
  \textbf{Dễ mở rộng theo số lượng dịch vụ}: có thể tăng thêm bản sao
  riêng cho broker, bridge hoặc API khi số thiết bị tăng, thay vì phải mở rộng toàn bộ
  hệ thống một cách đồng loạt và tốn kém.
\item
  \textbf{Phù hợp chi phí và phạm vi thực hiện}: tự dựng trên máy chủ
  bằng Docker Compose, không phụ thuộc nền tảng máy chủ thương mại, nên phù hợp với
  tính chất của một đồ án vừa cần triển khai thật, vừa cần kiểm soát tốt
  nguồn lực.
\end{itemize}

\subsubsection{3.3.2. Giải pháp dịch vụ máy chủ và xử lý
nghiệp vụ}\label{323-giux1ea3i-phuxe1p-backend--cloud}

Sau khi xác định phương án kiến trúc phù hợp, phần này đi vào cách triển
khai cụ thể từng lớp của hệ thống máy chủ và dịch vụ xử lý. Cách trình bày
được sắp xếp theo hành trình dữ liệu, từ lúc bản tin đi vào broker cho tới khi
thông tin xuất hiện trên giao diện của người vận hành.

\paragraph{3.3.2.1. Kiến trúc tổng quan hệ thống máy chủ (System
Architecture
Overview)}\label{3231-kiux1ebfn-truxfac-tux1ed5ng-quan-hux1ec7-thux1ed1ng-cloud-cloud-architecture-overview}

\subparagraph{a) Mô hình kiến trúc tổng
thể}\label{a-muxf4-huxecnh-kiux1ebfn-truxfac-tux1ed5ng-thux1ec3}

Để người đọc dễ theo dõi, có thể hình dung hệ thống máy chủ như một
chuỗi xử lý nhiều lớp: lớp đầu nhận dữ liệu, lớp giữa lưu và xử lý, lớp
cuối cung cấp thông tin cho người dùng. Kiến trúc này vừa tận dụng được
ưu thế của mô hình phản ứng theo sự kiện trong IoT, vừa giữ cho trách
nhiệm của từng thành phần không bị chồng lấn. Cụ thể, hệ thống gồm các
tầng chính sau:

\begin{itemize}
\tightlist
\item
  \textbf{Tầng tiếp nhận dữ liệu}: Đây là lớp đón dữ liệu đầu tiên từ
  thiết bị, nhận các bản tin MQTT rồi chuyển chúng đến đúng nơi cần xử
  lý.
\item
  \textbf{Tầng lưu trữ dữ liệu}: Dữ liệu được tách ra hai nơi: một nơi
  lưu dữ liệu quản trị có quan hệ rõ ràng, một nơi lưu dữ liệu thay đổi
  liên tục theo thời gian. Cách chia này tránh ép mọi thứ vào cùng một
  cơ sở dữ liệu.
\item
  \textbf{Tầng xử lý nghiệp vụ}: Đây là nơi xử lý các yêu cầu từ giao
  diện, áp quy tắc nghiệp vụ và quản lý trạng thái chung của hệ thống.
\item
  \textbf{Tầng cập nhật thời gian thực}: Đây là nơi đẩy dữ liệu mới trực
  tiếp đến người dùng thông qua WebSocket, tức kênh để máy chủ chủ động
  gửi thông tin mới lên trình duyệt.
\end{itemize}


{\thesissetfigureid{3.12}{fig-ch3-so-o-kien-truc-tong-quan-he-thong-cloud}
\begin{figure}[htbp]
\centering
\pandocbounded{\safeincludesvg{./assets/figures/05-chuong-3-giai-phap-backend-hinh-3-12.svg}}
\caption{Sơ đồ kiến trúc tổng quan hệ thống máy chủ}
\label{fig:ch3-so-o-kien-truc-tong-quan-he-thong-cloud}
\end{figure}
}



{\thesissetfigureid{3.12a}{fig-ch3-luong-du-lieu-chi-tiet-tu-thiet-bi-en-dashboard}
\begin{figure}[htbp]
\centering
\pandocbounded{\safeincludesvg{./assets/figures/05-chuong-3-giai-phap-backend-hinh-3-12a.svg}}
\caption{Luồng dữ liệu chi tiết từ thiết bị đến giao diện quản trị}
\label{fig:ch3-luong-du-lieu-chi-tiet-tu-thiet-bi-en-dashboard}
\end{figure}
}


\subparagraph{b) Luồng dữ liệu chính trong hệ
thống}\label{b-luux1ed3ng-dux1eef-liux1ec7u-chuxednh-trong-hux1ec7-thux1ed1ng}

Nếu mô tả bằng một ``câu chuyện dữ liệu'', luồng chính của hệ thống diễn
ra theo trình tự sau:

\begin{enumerate}
\def\labelenumi{\arabic{enumi}.}
\tightlist
\item
  \textbf{Thiết bị IoT} gửi dòng dữ liệu vận hành chính như vị trí GNSS,
  trạng thái nguồn, mức rung, ignition và uptime lên EMQX Broker thông
  qua MQTT 3.1.1 theo topic có cấu trúc
  \texttt{v1/\{device\_id\}/rawdata}. Ở đây, EMQX Broker có thể hiểu là
  broker MQTT, tức nơi nhận và chuyển tiếp bản tin.
\item
  \textbf{EMQX Broker} tiếp nhận và phân phối bản tin đến các bên cần
  dùng dữ liệu. Rules Engine của EMQX chỉ đóng vai trò lọc hoặc bổ trợ
  nhẹ; phần lớn nghiệp vụ cảnh báo vẫn được xử lý tại MQTT Bridge.
\item
  \textbf{MQTT Bridge} là dịch vụ trung gian đứng giữa broker và các nơi
  lưu dữ liệu. Dịch vụ này lắng nghe các topic từ EMQX, kiểm tra dữ liệu
  đầu vào bằng Zod schema, tức bộ quy tắc mô tả dữ liệu nào hợp lệ, sau
  đó ghi đồng thời vào nhiều nơi, nghĩa là một lần nhận dữ liệu nhưng
  lưu cho nhiều mục đích khác nhau:

  \begin{itemize}
  \tightlist
  \item
    Ghi dữ liệu thay đổi liên tục theo thời gian vào
    \textbf{VictoriaMetrics} (vị trí, tốc độ, điện áp, rung, ignition,
    uptime).
  \item
    Ghi nhật ký sự kiện vào \textbf{VictoriaLogs} (kết nối/ngắt kết nối,
    lỗi, sự kiện MQTT).
  \item
    Cập nhật trạng thái thiết bị vào \textbf{PostgreSQL} (trạng thái
    trực tuyến, phiên làm việc).
  \item
    Phát sự kiện đến \textbf{Socket.IO} để cập nhật thời gian thực cho
    giao diện quản trị.
  \end{itemize}
\item
  \textbf{API phía máy chủ} (Express.js) là lớp để giao diện gọi vào khi cần
  lấy dữ liệu lịch sử, dữ liệu quản trị hoặc gửi lệnh. Dịch vụ này truy
  vấn dữ liệu từ PostgreSQL và VictoriaMetrics, đồng thời gửi lệnh điều
  khiển ngược về thiết bị thông qua EMQX.
\item
  \textbf{Giao diện web} (Next.js) nhận dữ liệu cập nhật thời gian thực qua
  Socket.IO và dữ liệu lịch sử qua REST API.
\end{enumerate}

\subparagraph{c) Mô hình triển khai
Docker}\label{c-muxf4-huxecnh-triux1ec3n-khai-docker}

Ở mức triển khai, kiến trúc trên được hiện thực theo cách mỗi dịch vụ có
một tệp Docker Compose riêng và cùng tham gia một mạng Docker dùng chung.
Cách tổ chức này giúp cấu
trúc vận hành của đồ án bám sát cách các hệ thống thật được chia dịch vụ.
Hiểu đơn giản, mỗi dịch vụ được đặt trong một ``hộp chạy riêng'' gọi là
container; các hộp này có thể nói chuyện với nhau qua mạng dùng chung,
nhưng vẫn tách lỗi và tài nguyên tương đối rõ ràng. Nhờ vậy, cấu trúc
vận hành của đồ án bám sát cách các hệ thống thật được chia dịch vụ
trong thực tế:

\begin{itemize}
\tightlist
\item
  \textbf{Độc lập khi triển khai}: Mỗi dịch vụ có thể khởi động, dừng
  hoặc cập nhật độc lập mà không ảnh hưởng đến các dịch vụ khác.
\item
  \textbf{Quản lý tài nguyên riêng}: Mỗi container, tức mỗi gói chạy tách
  biệt, được giới hạn CPU và RAM cụ thể (ví dụ: API phía máy chủ
  512MB/1 CPU, MQTT Bridge 256MB/0.5 CPU).
\item
  \textbf{Dễ dàng mở rộng}: Có thể nhân thêm bản sao cho từng dịch vụ riêng
  lẻ theo nhu cầu.
\end{itemize}


{\thesissettableid{3.16}{tab-ch3-phan-bo-tai-nguyen-docker-cho-cac-dich-vu}
\def\LTcaptype{table}
\begin{longtable}{|L{0.326\linewidth}|L{0.327\linewidth}|L{0.128\linewidth}|L{0.173\linewidth}|}
\caption{Phân bổ tài nguyên Docker cho các dịch vụ}\label{tab:ch3-phan-bo-tai-nguyen-docker-cho-cac-dich-vu}\\
\hline \textbf{Dịch vụ} & \textbf{Container} & \textbf{Port} & \textbf{Tài nguyên} \\
\\\hline
\endfirsthead
\caption[]{Phân bổ tài nguyên Docker cho các dịch vụ (tiếp theo)}\\
\hline \textbf{Dịch vụ} & \textbf{Container} & \textbf{Port} & \textbf{Tài nguyên} \\
\\\hline
\endhead
\\\hline
\endfoot
\endlastfoot
Tracking\_Backend & tracking-backend & 4000 & 512MB, 1 CPU \\\hline
Tracking\_Frontend & tracking-frontend & 4001 & 512MB, 1 CPU \\\hline
Tracking\_MqttBridge & tracking-mqtt-bridge & - & 256MB, 0.5 CPU \\\hline
Tracking\_PostgreSQL & tracking-postgres & 5432 & 1GB, 2 CPU \\\hline
Tracking\_EMQX & tracking-emqx & 1883, 8883, 8083, 8084, 18083 & 1GB, 2
CPU \\\hline
Tracking\_VictoriaMetrics & tracking-victoriametrics & 8428 & 512MB, 1
CPU \\\hline
Tracking\_VictoriaLogs & tracking-victorialogs & 9428 & 512MB, 1 CPU \\\hline
Tracking\_Grafana & tracking-grafana & 4002 & Chưa cấu hình cứng \\\hline
\end{longtable}
}


\paragraph{3.3.2.2. MQTT Broker - EMQX (lựa chọn và cấu hình)}\label{3232-mqtt-broker---emqx-message-broker-selection--configuration}

\subparagraph{a) Vai trò của MQTT
Broker}\label{a-vai-truxf2-cux1ee7a-mqtt-broker}

MQTT Broker là điểm vào đầu tiên của dữ liệu ở phía máy chủ. Vai trò của
nó có thể hiểu đơn giản là nơi nhận bản tin từ thiết bị rồi chuyển tiếp
cho các thành phần cần xử lý tiếp theo. Nhờ mô hình \textbf{phát bản tin vào kênh chung rồi đăng ký nhận},
các thành phần trong hệ thống có thể giao tiếp bất đồng bộ mà không cần
ghép chặt với nhau.

Trong hệ thống theo dõi phương tiện, cách tổ chức này đặc biệt quan
trọng vì số lượng thiết bị có thể tăng dần theo thời gian, trong khi
nhịp gửi dữ liệu lại thay đổi theo chế độ hoạt động của từng xe.

\subparagraph{b) So sánh và lựa chọn MQTT
Broker}\label{b-so-suxe1nh-vuxe0-lux1ef1a-chux1ecdn-mqtt-broker}

Để lựa chọn MQTT Broker phù hợp, đồ án đã đánh giá bốn giải pháp phổ
biến trên thị trường:


{\thesissettableid{3.17}{tab-ch3-so-sanh-cac-mqtt-broker-pho-bien}
\def\LTcaptype{table}
\begin{longtable}{|L{0.269\linewidth}|L{0.188\linewidth}|L{0.158\linewidth}|L{0.152\linewidth}|L{0.177\linewidth}|}
\caption{So sánh các MQTT Broker phổ biến}\label{tab:ch3-so-sanh-cac-mqtt-broker-pho-bien}\\
\hline \textbf{Tiêu chí} & \textbf{Mosquitto} & \textbf{EMQX} & \textbf{HiveMQ} & \textbf{VerneMQ} \\
\\\hline
\endfirsthead
\caption[]{So sánh các MQTT Broker phổ biến (tiếp theo)}\\
\hline \textbf{Tiêu chí} & \textbf{Mosquitto} & \textbf{EMQX} & \textbf{HiveMQ} & \textbf{VerneMQ} \\
\\\hline
\endhead
\\\hline
\endfoot
\endlastfoot
\textbf{Kết nối đồng thời} & ~1.000 &
~100.000.000 & ~200.000.000 &
~10.000.000 \\\hline
\textbf{Thông lượng bản tin} & 40K msg/s & 100K msg/s & 200K msg/s & 50K
msg/s \\\hline
\textbf{Độ trễ xử lý} & 0,25 ms & 0,27 ms & \textless1 ms & 2,1
ms \\\hline
\textbf{Sử dụng CPU} & Rất thấp & Thấp (2\%) & Thấp & Cao (10\%) \\\hline
\textbf{Sử dụng RAM} & 254 MB & 495 MB & Trung bình & 1,2 GB \\\hline
\textbf{Hỗ trợ chạy nhiều nút} & Không & Có (20+ node) & Có & Có \\\hline
\textbf{Khả dụng cao (HA)} & Không & Có & Có & Có \\\hline
\textbf{Độ phức tạp cài đặt} & Rất dễ & Trung bình & Khó & Trung bình \\\hline
\textbf{Chi phí} & Miễn phí & Miễn phí/Trả phí & Trả phí & Miễn phí/Trả
phí \\\hline
\textbf{Cộng đồng} & Lớn & Lớn & Trung bình & Nhỏ \\\hline
\end{longtable}
}


\textbf{Kết luận lựa chọn: EMQX}

EMQX được lựa chọn làm MQTT Broker cho hệ thống vì đây là phương án giữ
được thế cân bằng tốt giữa hiệu năng, khả năng mở rộng và mức độ phù hợp
với nguồn lực triển khai của đồ án. Các lý do chính gồm:

\begin{itemize}
\tightlist
\item
  \textbf{Hiệu năng phù hợp}: Theo tài liệu công bố của EMQX, broker đáp
  ứng tốt tải bản tin lớn với độ trễ thấp, đủ cho quy mô triển khai hiện
  tại và còn dư địa mở rộng.
\item
  \textbf{Khả năng mở rộng}: Hỗ trợ chạy nhiều nút cùng lúc, thuận lợi nếu
  hệ thống tăng số lượng thiết bị trong giai đoạn sau.
\item
  \textbf{Bộ lọc bản tin tích hợp}: Thuận tiện cho các tác vụ lọc hoặc
  chuyển đổi dữ liệu đơn giản ngay tại broker; tuy nhiên trong kiến trúc
  hiện tại, đây không phải lớp xử lý nghiệp vụ chính.
\item
  \textbf{Chi phí phù hợp và hệ sinh thái phong phú}: Phiên bản mã nguồn
  mở đáp ứng đủ nhu cầu của đồ án và có cộng đồng hỗ trợ rộng.
\item
  \textbf{Màn hình quản lý}: Cung cấp giao diện web quản lý tại port
  18083, hỗ trợ giám sát kết nối, topic và bản tin theo thời gian thực.
\end{itemize}

\subparagraph{c) Cấu hình bộ lọc dữ liệu của
EMQX}\label{c-cux1ea5u-huxecnh-emqx-rules-engine}

Trong hệ thống này, bộ lọc dữ liệu của EMQX chỉ được dùng cho các phép xử
lý nhẹ ngay tại broker, chẳng hạn nhận biết sớm một số điều kiện bất
thường trước khi chuyển bản tin tới lớp xử lý phía sau. Phần lớn quyết
định nghiệp vụ, đặc biệt là ghi nhận cảnh báo và đồng bộ giao diện, vẫn
do MQTT Bridge và API phía máy chủ đảm nhiệm.

Các tình huống phù hợp để cấu hình tại lớp này gồm:

\begin{itemize}
\tightlist
\item
  \textbf{Phát hiện rung mạnh:} nhận biết sớm các bản tin cảnh báo rung để
  chuyển sang luồng xử lý cảnh báo.
\item
  \textbf{Phát hiện vượt tốc độ:} đánh dấu các bản tin có tốc độ vượt
  ngưỡng cấu hình để lớp phía sau tiếp tục xác minh và lưu vết.
\item
  \textbf{Phát hiện pin dự phòng thấp:} gắn cờ các bản tin cho thấy điện
  áp pin dự phòng xuống thấp để ưu tiên theo dõi.
\end{itemize}

\textbf{Luồng xử lý tổng quát:}

\begin{enumerate}
\tightlist
\item
  Thiết bị phát bản tin dữ liệu vận hành và bản tin sự kiện lên broker
  MQTT.
\item
  Bộ lọc tại broker chỉ thực hiện bước nhận biết điều kiện ban đầu và
  chuyển tiếp bản tin phù hợp sang lớp xử lý kế tiếp.
\item
  MQTT Bridge đánh giá nghiệp vụ chi tiết, chuẩn hóa dữ liệu và phát sự
  kiện cho API phía máy chủ.
\item
  API phía máy chủ ghi dữ liệu vào PostgreSQL và phát Socket.IO để cập
  nhật giao diện quản trị theo thời gian thực.
\end{enumerate}

Ngoài ra, hệ thống có thể mở rộng thêm quy tắc vùng cho phép ở mức cơ bản, thống
kê traffic theo topic, hoặc phát hiện thiết bị im lặng bất thường để
phục vụ vận hành.

\paragraph{3.3.2.3. Thiết kế cơ sở dữ liệu (Database Architecture
Design)}\label{3233-thiux1ebft-kux1ebf-cux1a1-sux1edf-dux1eef-liux1ec7u-database-architecture-design}

\textbf{Đặt vấn đề lưu trữ:} Khi đi sâu hơn vào tầng dữ liệu, có thể thấy
hệ thống đang phục vụ hai nhu cầu rất khác nhau. Một mặt, dữ liệu vận
hành gửi liên tục cần
được ghi nhanh và truy vấn theo thời gian. Mặt khác, dữ liệu nghiệp vụ
đòi hỏi tính toàn vẹn, quan hệ rõ ràng và khả năng quản lý dài hạn. Nếu
dùng một công cụ cho cả hai loại bài toán, hệ thống thường sẽ phải đánh
đổi hiệu quả ở một trong hai phía.


{\thesissettableid{3.18A}{tab-ch3-so-sanh-cac-phuong-an-kien-truc-luu-tru-du-lieu}
\def\LTcaptype{table}
\begin{longtable}{|L{0.149\linewidth}|L{0.233\linewidth}|L{0.236\linewidth}|L{0.241\linewidth}|L{0.085\linewidth}|}
\caption{So sánh các phương án kiến trúc lưu trữ dữ liệu}\label{tab:ch3-so-sanh-cac-phuong-an-kien-truc-luu-tru-du-lieu}\\
\hline \textbf{Phương án} & \textbf{Mô tả} & \textbf{Ưu điểm} & \textbf{Hạn chế} & \textbf{Mức phù hợp} \\
\\\hline
\endfirsthead
\caption[]{So sánh các phương án kiến trúc lưu trữ dữ liệu (tiếp theo)}\\
\hline \textbf{Phương án} & \textbf{Mô tả} & \textbf{Ưu điểm} & \textbf{Hạn chế} & \textbf{Mức phù hợp} \\
\\\hline
\endhead
\\\hline
\endfoot
\endlastfoot
\textbf{PA-DB1: PostgreSQL duy nhất} & Dùng PostgreSQL cho cả nghiệp vụ
và dữ liệu vận hành gửi liên tục & Đơn giản vận hành & Ghi lượng dữ liệu lớn dễ ảnh hưởng truy
vấn nghiệp vụ, chi phí index/partition cao & Trung bình \\\hline
\textbf{PA-DB2: Time-series DB duy nhất} & Dùng CSDL chuỗi thời gian cho
mọi loại dữ liệu & Tối ưu ghi dữ liệu cảm biến & Kém phù hợp dữ liệu
quan hệ phức tạp, khó đảm bảo ràng buộc nghiệp vụ & Thấp \\\hline
\textbf{PA-DB3: Hybrid DB (Đã chọn)} & PostgreSQL +\newline Victoria\-Metrics,\newline
Victoria\-Logs & Mỗi loại dữ liệu dùng đúng công cụ, cân bằng hiệu năng và
khả năng truy vấn & Tăng số thành phần cần vận hành & \textbf{Cao} \\\hline
\end{longtable}
}


\textbf{Kết luận lựa chọn:} Chọn \textbf{PA-DB3 (hybrid database)} để
tách trách nhiệm dữ liệu, giữ ổn định cho nghiệp vụ và tối ưu luồng
dữ liệu vận hành theo thời gian thực.

\subparagraph{a) Chiến lược lưu trữ kép (Dual Database
Strategy)}\label{a-chiux1ebfn-lux1b0ux1ee3c-lux1b0u-trux1eef-kuxe9p-dual-database-strategy}

Hệ thống theo dõi phương tiện cần lưu trữ hai loại dữ liệu có đặc tính
khác nhau căn bản. Có thể hình dung rất đơn giản: một bên là dữ liệu
``chảy liên tục'' từ xe, bên còn lại là dữ liệu ``quản lý nghiệp vụ''.
Nếu cố dồn cả hai vào cùng một kiểu cơ sở dữ liệu, hệ thống sẽ sớm gặp
xung đột giữa nhu cầu ghi nhanh và nhu cầu truy vấn quan hệ chặt chẽ.
Vì vậy, việc tách lớp lưu trữ ở đây không phải để làm kiến trúc phức
tạp hơn, mà để mỗi loại dữ liệu được đưa về đúng nơi phù hợp:

\begin{enumerate}
\def\labelenumi{\arabic{enumi}.}
\tightlist
\item
\textbf{Dữ liệu thô (dữ liệu vận hành nguyên bản)}: Bao gồm vị trí GPS, tốc độ,
  mức pin, dữ liệu OBD2, dữ liệu cảm biến gia tốc. Loại dữ liệu này có
  tần suất ghi rất cao (mỗi giây khi xe đang di chuyển), khối lượng lớn,
  nhưng chỉ cần lưu trữ trong thời gian ngắn (7--30 ngày) và chủ yếu
  phục vụ truy vấn theo chuỗi thời gian.
\item
  \textbf{Dữ liệu nghiệp vụ}: Bao gồm thông tin xe,
  khách hàng, chuyến đi, cảnh báo, vi phạm, lệnh điều khiển. Loại dữ
  liệu này có tần suất ghi thấp, khối lượng nhỏ, nhưng cần lưu trữ lâu
  dài (6--12 tháng trở lên) và yêu cầu tính toàn vẹn dữ liệu quan hệ
  (ACID).
\end{enumerate}

Do sự khác biệt cơ bản về đặc tính, hệ thống áp dụng chiến lược lưu trữ
kép sử dụng ba cơ sở dữ liệu chuyên biệt. Trong đó, VictoriaLogs được
thêm như một lớp riêng cho nhật ký vận hành, giúp việc truy vết lỗi
không làm ảnh hưởng tới đường lưu dữ liệu nghiệp vụ hay dữ liệu vận hành gửi liên tục:


{\thesissetfigureid{3.13}{fig-ch3-so-o-chien-luoc-luu-tru-kep}
\begin{figure}[htbp]
\centering
\pandocbounded{\safeincludesvg{./assets/figures/05-chuong-3-giai-phap-backend-hinh-3-13.svg}}
\caption{Sơ đồ chiến lược lưu trữ kép}
\label{fig:ch3-so-o-chien-luoc-luu-tru-kep}
\end{figure}
}


\begin{itemize}
\tightlist
\item
  \textbf{PostgreSQL (quan hệ):} lưu dữ liệu nghiệp vụ như
  \texttt{vehicles}, \texttt{customers}, \texttt{trips}, \texttt{alerts},
  \texttt{violations}, \texttt{commands}, \texttt{geofences}.
\item
  \textbf{VictoriaMetrics (chuỗi thời gian):} lưu dữ liệu vận hành mật độ
  cao như vị trí GPS, tốc độ, mức pin, trạng thái bật/tắt máy, thời gian hoạt động liên tục và dữ liệu IMU
  (retention khoảng 30 ngày).
\item
  \textbf{VictoriaLogs (nhật ký):} lưu sự kiện thiết bị, lỗi, phiên làm
  việc, bản tin MQTT và nhật ký kiểm toán (retention khoảng 7 ngày).
\end{itemize}


{\thesissettableid{3.18}{tab-ch3-so-sanh-cac-co-so-du-lieu-trong-he-thong}
\def\LTcaptype{table}
\begin{longtable}{|L{0.269\linewidth}|L{0.305\linewidth}|L{0.390\linewidth}|}
\caption{So sánh các cơ sở dữ liệu trong hệ thống}\label{tab:ch3-so-sanh-cac-co-so-du-lieu-trong-he-thong}\\
\hline \textbf{Cơ sở dữ liệu} & \textbf{Mục đích} & \textbf{Ưu điểm chính} \\
\\\hline
\endfirsthead
\caption[]{So sánh các cơ sở dữ liệu trong hệ thống (tiếp theo)}\\
\hline \textbf{Cơ sở dữ liệu} & \textbf{Mục đích} & \textbf{Ưu điểm chính} \\
\\\hline
\endhead
\\\hline
\endfoot
\endlastfoot
\textbf{PostgreSQL 16} & Dữ liệu quan hệ (nghiệp vụ) & Bảo toàn tính nhất quán,
liên kết tốt giữa các bảng và truy vấn phức tạp \\\hline
\textbf{VictoriaMetrics} & Dữ liệu chuỗi thời gian (dữ liệu vận hành gửi liên tục) & Ghi nhanh,
truy vấn theo mốc thời gian tốt, tiêu thụ RAM thấp \\\hline
\textbf{VictoriaLogs} & Nhật ký tập trung & Tìm kiếm nhanh,
lưu trữ nhỏ gọn, phù hợp truy vết sự cố \\\hline
\end{longtable}
}


\subparagraph{b) Cơ sở dữ liệu PostgreSQL - Schema quan
hệ}\label{b-cux1a1-sux1edf-dux1eef-liux1ec7u-postgresql---schema-quan-hux1ec7}

PostgreSQL 16 lưu trữ toàn bộ dữ liệu nghiệp vụ của hệ thống. Schema
được chuẩn hóa đến dạng chuẩn thứ ba (3NF) và hỗ trợ soft delete cho các
bảng quan trọng.

\textbf{Các nhóm bảng chính (Phase 1 - khoảng 20 bảng):}


{\thesissettableid{3.19}{tab-ch3-cac-nhom-bang-chinh-trong-postgresql}
\def\LTcaptype{table}
\begin{longtable}{|L{0.203\linewidth}|L{0.419\linewidth}|L{0.343\linewidth}|}
\caption{Các nhóm bảng chính trong PostgreSQL}\label{tab:ch3-cac-nhom-bang-chinh-trong-postgresql}\\
\hline \textbf{Nhóm bảng} & \textbf{Các bảng chính} & \textbf{Mô tả} \\
\\\hline
\endfirsthead
\caption[]{Các nhóm bảng chính trong PostgreSQL (tiếp theo)}\\
\hline \textbf{Nhóm bảng} & \textbf{Các bảng chính} & \textbf{Mô tả} \\
\\\hline
\endhead
\\\hline
\endfoot
\endlastfoot
\textbf{Người dùng \& Xác thực} & users, user\_sessions & Quản lý tài
khoản, phiên đăng nhập \\\hline
\textbf{Phương tiện \& Thiết bị} & vehicles, devices,
device\_configurations & Thông tin xe và thiết bị theo dõi \\\hline
\textbf{Khách hàng} & customers & Thông tin khách hàng thuê xe \\\hline
\textbf{Chuyến đi} & trips, trip\_events, stops & Hành trình, sự kiện và
điểm dừng \\\hline
\textbf{Cảnh báo \& Vi phạm} & alerts, violations & Cảnh báo hệ thống và
vi phạm giao thông \\\hline
\textbf{Vùng địa lý} & geofences, vehicle\_geofences & Vùng địa lý và
liên kết với xe \\\hline
\textbf{Lệnh điều khiển} & commands & Lệnh gửi đến thiết bị \\\hline
\textbf{Bảo trì} & maintenance\_records & Lịch sử bảo trì phương tiện \\\hline
\textbf{Nhật ký} & connection\_logs, device\_status\_history & Nhật ký
kết nối và lịch sử trạng thái \\\hline
\textbf{Thông báo} & notifications & Thông báo hệ thống \\\hline
\textbf{Nhiên liệu} & fuel\_records & Quản lý nhiên liệu \\\hline
\end{longtable}
}


\textbf{Sơ đồ quan hệ chính (ER rút gọn):}


{\thesissetfigureid{3.14}{fig-ch3-so-o-quan-he-co-so-du-lieu-er-diagram}
\begin{figure}[htbp]
\centering
\pandocbounded{\safeincludesvg{./assets/figures/05-chuong-3-giai-phap-backend-hinh-3-14.svg}}
\caption{Sơ đồ quan hệ cơ sở dữ liệu (ER)}
\label{fig:ch3-so-o-quan-he-co-so-du-lieu-er-diagram}
\end{figure}
}


\begin{itemize}
\tightlist
\item
  \texttt{users} 1:N \texttt{vehicles}, \texttt{customers},
  \texttt{notifications}.
\item
  \texttt{vehicles} 1:1 \texttt{devices}; 1:N \texttt{trips},
  \texttt{alerts}, \texttt{maintenance\_records}; N:N \texttt{geofences}
  (qua \texttt{vehicle\_geofences}).
\item
  \texttt{customers} 1:N \texttt{trips}.
\item
  \texttt{trips} 1:N \texttt{trip\_events}, \texttt{stops},
  \texttt{violations}.
\item
  \texttt{devices} 1:N \texttt{commands}, \texttt{device\_status\_history},
  \texttt{connection\_logs}.
\end{itemize}

Schema được thiết kế với khả năng mở rộng, các khóa ngoại và chỉ mục
(index) đã được chuẩn bị sẵn để tích hợp thêm các bảng của giai đoạn 2 bao gồm:
\texttt{bookings} (đặt xe), \texttt{rental\_contracts} (hợp đồng thuê), \texttt{payments} (thanh
toán), \texttt{damage\_reports} (báo cáo hư hỏng) và \texttt{reviews} (đánh giá).

\subparagraph{c) Cơ sở dữ liệu VictoriaMetrics - Schema chuỗi thời
gian}\label{c-cux1a1-sux1edf-dux1eef-liux1ec7u-victoriametrics---schema-chuux1ed7i-thux1eddi-gian}

VictoriaMetrics là cơ sở dữ liệu chuỗi thời gian tương thích Prometheus,
được lựa chọn thay thế InfluxDB nhờ hiệu năng vượt trội (nhanh hơn 10
lần), mức tiêu thụ RAM thấp hơn và tích hợp tốt với Grafana.


{\thesissettableid{3.20}{tab-ch3-so-sanh-victoriametrics-va-influxdb}
\def\LTcaptype{table}
\begin{longtable}{|L{0.307\linewidth}|L{0.210\linewidth}|L{0.447\linewidth}|}
\caption{So sánh VictoriaMetrics và InfluxDB}\label{tab:ch3-so-sanh-victoriametrics-va-influxdb}\\
\hline \textbf{Tiêu chí} & \textbf{InfluxDB} & \textbf{VictoriaMetrics (Đã chọn)} \\
\\\hline
\endfirsthead
\caption[]{So sánh VictoriaMetrics và InfluxDB (tiếp theo)}\\
\hline \textbf{Tiêu chí} & \textbf{InfluxDB} & \textbf{VictoriaMetrics (Đã chọn)} \\
\\\hline
\endhead
\\\hline
\endfoot
\endlastfoot
\textbf{Ngôn ngữ truy vấn} & Flux (phức tạp) & PromQL (đơn giản) \\\hline
\textbf{Hiệu năng} & Tốt & Nhanh gấp 10 lần \\\hline
\textbf{Sử dụng RAM} & Cao & Thấp \\\hline
\textbf{Nén dữ liệu} & Tốt & Tốt hơn (gấp 10 lần) \\\hline
\textbf{Tương thích Prometheus} & Không & Có \\\hline
\textbf{Độ khó học} & Cao (Flux) & Thấp (PromQL) \\\hline
\textbf{Tích hợp Grafana} & Qua plugin & Hỗ trợ sẵn (Native) \\\hline
\end{longtable}
}


\textbf{Chính sách lưu trữ (Retention Policy):} Dữ liệu thô được lưu trữ
trong 30 ngày (\texttt{-retentionPeriod=30d}). Hệ thống hỗ trợ cấu hình
downsampling thông qua vmagent với recording rules để tạo dữ liệu tổng
hợp theo giờ/ngày, phục vụ báo cáo dài hạn mà không tốn dung lượng lưu
trữ lớn.

\subparagraph{d) VictoriaLogs - Cơ sở dữ liệu nhật
ký}\label{d-victorialogs---cux1a1-sux1edf-dux1eef-liux1ec7u-nhux1eadt-kuxfd}

VictoriaLogs được sử dụng làm hệ thống nhật ký tập trung, tức nơi gom log
từ nhiều dịch vụ về một chỗ để dễ theo dõi, thay thế bộ công cụ ELK
Stack truyền thống với ưu điểm nhẹ hơn và dễ triển khai. VictoriaLogs
lưu trữ các loại nhật ký: sự kiện thiết bị, nhật ký lỗi, phiên làm việc,
bản tin MQTT, và nhật ký kiểm toán,
tức dấu vết ai đã làm gì và vào lúc nào. Chính sách lưu trữ: 7 ngày cho nhật ký thông thường.

Đến đây, kiến trúc lưu trữ đã trả lời câu hỏi dữ liệu sẽ được cất ở đâu.
Tuy nhiên, để người dùng khai thác được số dữ liệu đó, hệ thống vẫn cần
một lớp nghiệp vụ đứng ra tổ chức API, quyền truy cập và các luồng điều
khiển. Vì vậy, phần tiếp theo chuyển sang API phía máy chủ.

\paragraph{3.3.2.4. API phía máy chủ và kiến trúc phần mềm (API Server \&
Software
Architecture)}\label{3234-api-server-vuxe0-kiux1ebfn-truxfac-phux1ea7n-mux1ec1m-api-server--software-architecture}

\subparagraph{a) Lựa chọn công
nghệ}\label{a-lux1ef1a-chux1ecdn-cuxf4ng-nghux1ec7}

API phía máy chủ là lớp làm việc trực tiếp với người vận hành và các quy tắc
nghiệp vụ. Nếu MQTT Bridge chuyên nhận dòng dữ liệu thô từ thiết bị, thì
API phía máy chủ phụ trách những gì người dùng nhìn thấy hằng ngày: đăng nhập,
quản lý xe, xem lịch sử, tạo cảnh báo, quản lý vùng cho phép và điều phối
OTA. Vì vai trò khác nhau, tiêu chí chọn framework ở đây cũng khác:
quan trọng nhất là mã nguồn phải rõ ràng, linh hoạt và dễ bảo trì, chứ
không chỉ có kết quả đo tốc độ đẹp trên các bài thử chuẩn.


{\thesissettableid{3.21}{tab-ch3-so-sanh-framework-backend}
\def\LTcaptype{table}
\begin{longtable}{|L{0.350\linewidth}|L{0.304\linewidth}|L{0.155\linewidth}|L{0.145\linewidth}|}
\caption{So sánh các khung phát triển cho API phía máy chủ}\label{tab:ch3-so-sanh-framework-backend}\\
\hline \textbf{Tiêu chí} & \textbf{Express + TypeScript} & \textbf{NestJS} & \textbf{Go + Gin} \\
\\\hline
\endfirsthead
\caption[]{So sánh các khung phát triển cho API phía máy chủ (tiếp theo)}\\
\hline \textbf{Tiêu chí} & \textbf{Express + TypeScript} & \textbf{NestJS} & \textbf{Go + Gin} \\
\\\hline
\endhead
\\\hline
\endfoot
\endlastfoot
\textbf{Độ khó học} & Rất dễ & Trung bình & Khó \\\hline
\textbf{Linh hoạt} & Rất cao & Trung bình & Cao \\\hline
\textbf{Hiệu năng} & Cao & Cao & Rất cao \\\hline
\textbf{Lượng mã khung lặp lại} & Ít & Nhiều & Trung bình \\\hline
\textbf{Hệ sinh thái IoT} & Rất tốt & Tốt & Tốt \\\hline
\textbf{Hỗ trợ thời gian thực} & Rất tốt & Rất tốt & Tốt \\\hline
\end{longtable}
}


\textbf{Kết luận lựa chọn: Node.js + Express + TypeScript}

Lựa chọn này phù hợp vì ba lý do dễ thấy. Thứ nhất, Express nhẹ và ít
ràng buộc, nên nhóm có thể tổ chức mã nguồn bám theo nghiệp vụ thay vì
phải đi theo khuôn sẵn có của framework. Thứ hai, TypeScript giúp phát
hiện sớm sai khác kiểu dữ liệu giữa lớp nhận yêu cầu, lớp xử lý nghiệp
vụ và cơ sở dữ liệu, nhờ đó giảm lỗi lan xuống sâu. Thứ ba, Zod kiểm tra dữ liệu đầu vào ngay khi
chương trình chạy, tức là dữ liệu sai cấu trúc sẽ bị chặn từ cửa vào.
Với một hệ thống phải nhận dữ liệu từ giao diện quản trị, thiết bị và luồng OTA,
ba yếu tố này thực dụng hơn việc chọn một framework nhiều tính năng
nhưng nặng khuôn mẫu.

\subparagraph{b) Kiến trúc Domain-Driven Design
(DDD)}\label{b-kiux1ebfn-truxfac-domain-driven-design-ddd}

API phía máy chủ được tổ chức theo kiến trúc Domain-Driven Design với ba
tầng rõ ràng. Trong phạm vi đồ án, DDD có thể hiểu đơn giản là chia mã
nguồn theo từng bài toán nghiệp vụ thay vì gom theo loại file kỹ thuật.
Nhờ đó, phần quản lý xe, thiết bị, dữ liệu vận hành gửi liên tục hay
cảnh báo được tách thành các miền tương đối độc lập; khi thay đổi một
nghiệp vụ cụ thể, phạm vi ảnh hưởng cũng dễ kiểm soát hơn.

\textbf{Tầng 1 - Lớp nhận yêu cầu (API Layer / Routes):} Định nghĩa các
đường gọi API, nơi ánh xạ từng yêu cầu HTTP tới bộ xử lý tương ứng.

\textbf{Tầng 2 - Lớp nghiệp vụ (Service Layer):} Chứa toàn bộ logic
nghiệp vụ, điều phối giữa các nơi truy cập dữ liệu và các dịch vụ bên ngoài.

\textbf{Tầng 3 - Lớp truy cập dữ liệu (Repository Layer):} Truy cập cơ sở dữ
liệu, thực hiện các truy vấn SQL và trả về dữ liệu đã được kiểm tra kiểu.

Hệ thống vì thế được chia thành hơn 20 domain module, mỗi module giữ
trọn phần việc của mình với các thư mục \texttt{services/},
\texttt{repositories/}, \texttt{types/} riêng. Cách chia này giúp người
đọc lần theo từng nghiệp vụ rõ hơn và giảm việc sửa chéo giữa các phần
không liên quan.

\subparagraph{c) MQTT Bridge - Dịch vụ độc
lập}\label{c-mqtt-bridge---dux1ecbch-vux1ee5-ux111ux1ed9c-lux1eadp}

MQTT Bridge đã được tách thành một dịch vụ Node.js độc lập với cơ chế ghi
log riêng, nhóm kết nối cơ sở dữ
liệu dùng lại nhiều lần và cấu hình riêng. Cách tách này giúp dịch vụ dễ
mở rộng hơn, cô lập lỗi tốt hơn và giữ kiến trúc tổng thể gọn hơn.

Khi cấu trúc dịch vụ đã rõ, bước tiếp theo là chốt ``hợp đồng giao
tiếp'' mà giao diện và các công cụ vận hành sẽ dùng hằng ngày. Vì vậy,
mục sau tập trung vào thiết kế API endpoint.

\paragraph{3.3.2.5. Thiết kế API Endpoints (REST API
Design)}\label{3235-thiux1ebft-kux1ebf-api-endpoints-rest-api-design}

\subparagraph{a) Nguyên tắc thiết
kế}\label{a-nguyuxean-tux1eafc-thiux1ebft-kux1ebf}

API được thiết kế theo kiểu REST, nghĩa là tổ chức theo từng tài nguyên
như xe, thiết bị hay cảnh báo: URL dạng số nhiều, gắn phiên bản ở đầu
đường dẫn (\texttt{/api/v1/}), có phân trang, định dạng lỗi thống nhất
và dùng HTTP Status Code chuẩn. Ở góc nhìn vận hành, API chính là ``hợp
đồng giao tiếp'' giữa giao diện quản trị, API phía máy chủ và một số
luồng thiết bị. Vì vậy, cách đặt tên và phân nhóm endpoint phải giúp
người đọc đoán được ngay endpoint đó phục vụ ai, thao tác gì và có cần
đăng nhập hay không.

\subparagraph{b) Bảng tóm tắt API
Endpoints}\label{b-bux1ea3ng-tuxf3m-tux1eaft-api-endpoints}


{\thesissettableid{3.22}{tab-ch3-tom-tat-cac-nhom-api-endpoints-chinh}
\def\LTcaptype{table}
\begin{longtable}{|L{0.200\linewidth}|L{0.122\linewidth}|L{0.181\linewidth}|L{0.340\linewidth}|L{0.101\linewidth}|}
\caption{Tóm tắt các nhóm API Endpoints chính}\label{tab:ch3-tom-tat-cac-nhom-api-endpoints-chinh}\\
\hline \textbf{Nhóm} & \textbf{Phương thức} & \textbf{Endpoint} & \textbf{Mô tả} & \textbf{Xác thực} \\
\\\hline
\endfirsthead
\caption[]{Tóm tắt các nhóm API Endpoints chính (tiếp theo)}\\
\hline \textbf{Nhóm} & \textbf{Phương thức} & \textbf{Endpoint} & \textbf{Mô tả} & \textbf{Xác thực} \\
\\\hline
\endhead
\\\hline
\endfoot
\endlastfoot
\textbf{Đăng nhập} & POST & \nolinkurl{/api/v1/auth/login} & Đăng nhập &
Không \\\hline
& GET & \nolinkurl{/api/v1/auth/me} & Thông tin người dùng hiện tại & Có \\\hline
& POST & \nolinkurl{/api/v1/auth/logout} & Đăng xuất & Có \\\hline
\textbf{Thiết bị} & GET & \nolinkurl{/api/v1/devices} & Danh sách thiết bị & Có \\\hline
& POST & \nolinkurl{/api/v1/devices} & Tạo thiết bị mới & Có \\\hline
& POST & \nolinkurl{/api/v1/devices/:id/ota} & Gửi lệnh OTA đến thiết bị & Có \\\hline
& POST & \nolinkurl{/api/v1/devices/:id/ota/rollback} & Yêu cầu rollback firmware &
Có \\\hline
\textbf{Phương tiện} & GET & \nolinkurl{/api/v1/vehicles} & Danh sách xe & Có \\\hline
& POST & \nolinkurl{/api/v1/vehicles} & Tạo xe mới & Có \\\hline
\textbf{Dữ liệu IoT} & POST & \nolinkurl{/api/v1/iot/data} & Gửi dữ liệu cảm biến &
Không \\\hline
\textbf{Dữ liệu vận hành} & GET & \nolinkurl{/api/v1/devices/:id/telemetry} & Dữ liệu
vận hành theo từng thiết bị & Có \\\hline
& GET & \nolinkurl{/api/v1/telemetry/history} & Truy vấn lịch sử dữ liệu vận hành & Có \\\hline
\textbf{Tổng quan} & GET & \nolinkurl{/api/v1/dashboard/stats} & Thống kê tổng quan
& Có \\\hline
\textbf{Firmware} & POST & \nolinkurl{/api/v1/firmware/upload} & Tải lên firmware &
Có \\\hline
& GET & \nolinkurl{/api/v1/firmware/:id/download} & Tải file firmware để OTA & Có \\\hline
\end{longtable}
}



{\thesissetfigureid{3.14a}{fig-ch3-luong-ieu-phoi-ota-giua-backend-emqx-thiet-bi-va-mqtt-bridge}
\begin{figure}[htbp]
\centering
\pandocbounded{\safeincludesvg{./assets/figures/05-chuong-3-giai-phap-backend-hinh-3-14a.svg}}
\caption{Luồng điều phối OTA giữa API phía máy chủ, EMQX, thiết bị và MQTT Bridge}
\label{fig:ch3-luong-ieu-phoi-ota-giua-backend-emqx-thiet-bi-va-mqtt-bridge}
\end{figure}
}


API phía máy chủ vẫn duy trì một số alias tương thích như \texttt{/api/v1/device}
hoặc \texttt{/api/v1/export}, nhưng tài liệu ưu tiên dạng tài nguyên số
nhiều (\texttt{/devices}, \texttt{/vehicles}, \texttt{/customers},
\texttt{/trips}, \texttt{/alerts}, \texttt{/geofences}) để thống nhất
cách đặt tên. Cách trình bày này giúp người không quen với kiểu API dạng
tài nguyên vẫn có thể hiểu nhanh nhóm endpoint nào phục vụ đăng nhập,
nhóm nào phục vụ quản lý xe, nhóm nào dành cho dữ liệu vận hành hay OTA.

Khi ``cửa vào'' của hệ thống đã rõ bằng các endpoint cụ thể, phần còn
lại là cách dữ liệu mới được đẩy ngược ra giao diện mà không buộc người
dùng phải tải lại trang liên tục. Vì vậy, mục tiếp theo chuyển sang kênh
WebSocket thời gian thực.

\paragraph{3.3.2.6. WebSocket và truyền dữ liệu thời gian thực
(Cập nhật
Communication)}\label{3236-websocket-vuxe0-truyux1ec1n-dux1eef-liux1ec7u-thux1eddi-gian-thux1ef1c-real-time-communication}

Hệ thống sử dụng \textbf{Socket.IO 4.8} như một kênh để máy chủ chủ động
gửi dữ liệu mới ra giao diện web ngay khi có thay đổi, thay vì bắt giao diện quản trị hỏi lại
liên tục. Trong triển khai hiện tại, kênh này được chia thành năm nhóm
chính: \texttt{/dashboard} cho số liệu tổng quan, \texttt{/devices} cho
trạng thái theo từng xe, \texttt{/notifications} cho cảnh báo,
\texttt{/exports} cho tiến trình xuất báo cáo và \texttt{/firmware} cho
luồng OTA. Tất cả các nhóm đều đi qua xác thực; riêng
\texttt{/devices} cho phép vào ``phòng'' theo từng thiết bị để chỉ nhận
những cập nhật liên quan.

Các sự kiện tiêu biểu gồm cập nhật trạng thái xe, vị trí xe, xác nhận
lệnh, số liệu tổng quan, cảnh báo vượt vùng cho phép, tiến trình xuất báo cáo và
phân công firmware. Riêng luồng OTA được tách khá rõ: API phía máy chủ phát các
sự kiện mức nghiệp vụ để giao diện biết có nhiệm vụ mới, còn trạng thái
chi tiết trong quá trình cập nhật được MQTT Bridge ghi vào
\texttt{firmware\_update\_log} và VictoriaLogs để thuận tiện cho việc
theo dõi và truy vết khi có lỗi.

Khi dữ liệu đã đi được đến màn hình quản trị theo thời gian thực, phần
còn lại là bảo đảm chỉ đúng người, đúng thiết bị và đúng luồng dữ liệu
được phép đi qua. Vì vậy, mục kế tiếp chuyển sang thiết kế bảo mật toàn
hệ thống.

\paragraph{3.3.2.7. Bảo mật hệ thống}\label{3237-bux1ea3o-mux1eadt-hux1ec7-thux1ed1ng-security-design}

\subparagraph{a) Xác thực theo phiên làm việc (không dùng
JWT)}\label{a-xuxe1c-thux1ef1c-session-based-khuxf4ng-duxf9ng-jwt}

Hệ thống sử dụng cơ chế xác thực dựa trên phiên làm việc, nghĩa là trạng
thái đăng nhập được giữ ở phía máy chủ. Mỗi lần đăng nhập thành công,
máy chủ tạo ra một mã phiên ngẫu nhiên và lưu thông tin phiên tương ứng
trong cơ sở dữ liệu, thay vì dùng JWT tự mang toàn bộ thông tin phiên.
Token chỉ đóng vai trò khóa tham chiếu; API phía máy chủ mới là nơi lưu
và kiểm tra trạng thái phiên, còn giá trị token được băm SHA-256 trước
khi ghi xuống cơ sở dữ liệu.


{\thesissettableid{3.23}{tab-ch3-so-sanh-jwt-va-session-based-authentication}
\def\LTcaptype{table}
\begin{longtable}{|L{0.190\linewidth}|L{0.403\linewidth}|L{0.372\linewidth}|}
\caption{So sánh JWT và cơ chế xác thực theo phiên làm việc}\label{tab:ch3-so-sanh-jwt-va-session-based-authentication}\\
\hline \textbf{Tiêu chí} & \textbf{JWT} & \textbf{Theo phiên làm việc (đã chọn)} \\
\\\hline
\endfirsthead
\caption[]{So sánh JWT và cơ chế xác thực theo phiên làm việc (tiếp theo)}\\
\hline \textbf{Tiêu chí} & \textbf{JWT} & \textbf{Theo phiên làm việc (đã chọn)} \\
\\\hline
\endhead
\\\hline
\endfoot
\endlastfoot
\textbf{Thu hồi token} & Khó (cần blacklist) & Dễ (xóa khỏi DB) \\\hline
\textbf{Kiểm soát phiên} & Hạn chế & Toàn quyền (xem, thu hồi mọi
phiên) \\\hline
\textbf{Bảo mật} & Token chứa thông tin nhạy cảm & Chỉ chứa ID ngẫu
nhiên \\\hline
\textbf{Kích thước} & Lớn (gồm phần đầu, phần dữ liệu và chữ ký số) & Nhỏ (chỉ
mã phiên) \\\hline
\end{longtable}
}


\subparagraph{b) Bảo mật tầng truyền
tải}\label{b-bux1ea3o-mux1eadt-tux1ea7ng-truyux1ec1n-tux1ea3i}

Hệ thống bảo vệ tầng truyền tải bằng nhiều lớp bổ trợ: Helmet để thiết
lập các header bảo mật sẵn trên phản hồi web, CORS để giới hạn website
nào được phép gọi sang hệ thống, giới hạn tần suất yêu cầu để chặn việc gửi quá
nhiều yêu cầu trong thời gian ngắn từ cùng một địa chỉ IP, và cơ chế gắn
mã nhận diện từng yêu cầu để truy vết luồng xử lý khi cần dò lỗi.

\subparagraph{c) Bảo mật MQTT}\label{c-bux1ea3o-mux1eadt-mqtt}

Ở tầng MQTT, mỗi thiết bị được cấp ACL riêng và Client ID duy nhất, kết
hợp xác thực bằng username/password. Cách làm này giúp cô lập phạm vi
publish/subscribe của từng thiết bị theo dõi và giảm rủi ro một thiết bị truy cập
nhầm sang luồng dữ liệu của thiết bị khác.

\subparagraph{d) Hệ thống thông báo cảnh
báo}\label{d-hux1ec7-thux1ed1ng-thuxf4ng-buxe1o-cux1ea3nh-buxe1o}

Hệ thống hỗ trợ gửi thông báo qua Telegram Bot (thời gian thực) và Email
SMTP (chi tiết). Các loại cảnh báo: xe di chuyển bất thường, vi phạm tốc
độ, pin thấp, thiết bị mất kết nối, lịch bảo trì sắp đến.

Đến đây, tầng máy chủ đã chốt cách nhận dữ liệu, lưu dữ liệu, bảo vệ
truy cập và phát thông tin ra ngoài. Phần tiếp theo vì thế chuyển sang
tầng giao diện web, nơi toàn bộ dữ liệu đó được biến thành màn hình vận
hành mà người dùng thực sự tương tác.

\subsection{3.4. Phân tích và đề xuất công nghệ giao diện web -- Web
interface analysis and solution}\label{34-phan-tich-va-de-xuat-cong-nghe-frontend}

\subsubsection{3.4.1. Phân tích và lựa chọn công nghệ giao diện web}

\paragraph{3.4.1.1. Đặt vấn đề cho giải pháp
giao diện web}\label{3141-ux111ux1eb7t-vux1ea5n-ux111ux1ec1-cho-giux1ea3i-phuxe1p-frontend}

Nếu tầng máy chủ là nơi xử lý dữ liệu, thì giao diện web là nơi người vận
hành thực sự ``nhìn thấy'' hệ thống đang hoạt động ra sao. Một màn hình
tổng quan tốt không
chỉ hiển thị đủ thông tin, mà còn phải giúp người dùng nhận ra điều gì
quan trọng, điều gì bất thường và cần hành động ở đâu. Vì vậy, bài toán
thiết kế giao diện web cần đồng thời đáp ứng các yêu cầu sau:

\begin{itemize}
\tightlist
\item
  \textbf{Hiển thị thời gian thực ổn định}: bản đồ, trạng thái thiết bị
  và cảnh báo phải cập nhật liên tục với độ trễ thấp.
\item
  \textbf{Khả năng mở rộng theo tính năng nghiệp vụ}: quản lý xe, chuyến
  đi, vùng cho phép, tức vùng giám sát được vẽ trên bản đồ, cảnh báo và báo
  cáo cần dễ mở rộng.
\item
  \textbf{Hiệu năng và trải nghiệm người dùng}: tải trang nhanh, điều
  hướng mượt, bố cục thích ứng tốt trên desktop/mobile.
\item
  \textbf{Dễ triển khai vận hành}: quy trình đóng gói và chạy bằng
  Docker phải thuận lợi cho mô hình nhiều dịch vụ.
\end{itemize}

\paragraph{3.4.1.2. So sánh các phương án công nghệ
giao diện web}\label{3142-so-suxe1nh-cuxe1c-phux1b0ux1a1ng-uxe1n-cuxf4ng-nghux1ec7-frontend}


{\thesissettableid{3.23A}{tab-ch3-so-sanh-cac-phuong-an-cong-nghe-frontend-tong-the}
\def\LTcaptype{table}
\begin{longtable}{|L{0.154\linewidth}|L{0.157\linewidth}|L{0.280\linewidth}|L{0.269\linewidth}|L{0.085\linewidth}|}
\caption{So sánh các phương án công nghệ giao diện web tổng thể}\label{tab:ch3-so-sanh-cac-phuong-an-cong-nghe-frontend-tong-the}\\
\hline \textbf{Phương án} & \textbf{Mô tả bộ công cụ} & \textbf{Ưu điểm} & \textbf{Hạn chế} & \textbf{Mức phù hợp} \\
\\\hline
\endfirsthead
\caption[]{So sánh các phương án công nghệ giao diện web tổng thể (tiếp theo)}\\
\hline \textbf{Phương án} & \textbf{Mô tả bộ công cụ} & \textbf{Ưu điểm} & \textbf{Hạn chế} & \textbf{Mức phù hợp} \\
\\\hline
\endhead
\\\hline
\endfoot
\endlastfoot
\textbf{PA-FE1: Vite + React chỉ dựng ở trình duyệt} & React + Vite +
định tuyến chạy phía trình duyệt & Build nhanh, cấu hình linh hoạt &
Khả năng hiện nội dung sớm kém hơn cách dựng một phần ở máy chủ, cần tự
lắp ghép nhiều thành phần & Trung bình \\\hline
\textbf{PA-FE2: Nuxt (Vue)} & Vue + Nuxt với khả năng dựng ở máy chủ và
dựng sẵn & Cấu trúc rõ, hỗ trợ hiển thị sớm tốt & Khác hệ sinh thái
React đang dùng ở dự án, chi phí chuyển đổi cao & Trung bình \\\hline
\textbf{PA-FE3: Next.js + React (Đã chọn)} & Next.js 15 + React 19 + App
Router & Hỗ trợ tốt việc dựng trước một phần giao diện, hệ sinh thái lớn,
tối ưu cho môi trường vận hành và tích hợp kênh thời gian thực thuận lợi
& Độ phức tạp framework cao hơn cách dựng thuần ở trình duyệt &
\textbf{Cao} \\\hline
\end{longtable}
}


\paragraph{3.4.1.3. Chọn giải pháp
giao diện web}\label{3143-chux1ecdn-giux1ea3i-phuxe1p-frontend}

Đồ án chọn \textbf{PA-FE3: Next.js 15 + React 19} làm nền tảng giao diện web
chính. Đây không phải là lựa chọn chạy theo xu hướng framework, mà là
lựa chọn bám sát nhu cầu của một giao diện tổng quan phải vừa ``đọc nhanh'' vừa
``phản ứng nhanh'' khi dữ liệu và cảnh báo thay đổi liên tục.

Các lý do lựa chọn gồm:

\begin{itemize}
\tightlist
\item
  \textbf{Phù hợp yêu cầu dữ liệu lịch sử và dữ liệu thời gian thực cùng
  lúc}: kết hợp tốt giữa dữ liệu lấy theo yêu cầu rồi lưu tạm để dùng
  lại và kênh WebSocket để máy chủ chủ động gửi dữ liệu mới.
\item
  \textbf{Hiệu năng tải trang tốt}: Cơ chế App Router của Next.js và khả
  năng dựng một phần giao diện từ phía máy chủ giúp người dùng thấy nội
  dung ban đầu sớm hơn.
\item
  \textbf{Mở rộng tính năng thuận lợi}: dễ tổ chức theo từng nhóm tính năng
  người dùng thực sự thao tác, thay vì dồn tất cả mã nguồn vào một chỗ.
\item
  \textbf{Sẵn sàng triển khai vận hành}: hỗ trợ xuất gói chạy độc lập và
  đóng gói Docker rõ ràng.
\end{itemize}

\subsubsection{3.4.2. Giải pháp giao diện
web}\label{324-giux1ea3i-phuxe1p-frontend}

Phần 3.4.1 đã chốt bộ công cụ nền cho giao diện web. Từ điểm đó, mục
này đi tiếp vào cách biến lựa chọn công nghệ thành một giao diện vận
hành cụ thể: mã nguồn được tổ chức ra sao, dữ liệu được lấy và cập nhật
thế nào, và từng màn hình phục vụ nhu cầu gì của người dùng.

\paragraph{3.4.2.1. Lựa chọn công nghệ giao diện web}\label{3241-lux1ef1a-chux1ecdn-cuxf4ng-nghux1ec7-frontend-technology-selection}

\subparagraph{a) Framework ứng dụng
web}\label{a-framework-ux1ee9ng-dux1ee5ng-web}

Để lựa chọn framework phù hợp, nhóm không dừng ở so sánh tên công nghệ,
mà đặt câu hỏi thực tế hơn: framework nào giúp giao diện quản trị tải nhanh, tổ
chức rõ và vẫn dễ mở rộng khi số lượng trang quản trị tăng lên. Từ đó,
nhóm đối chiếu Next.js, Nuxt.js và Vite + React theo các tiêu chí hiệu
năng, khả năng mở rộng và trải nghiệm phát triển.


{\thesissettableid{3.24}{tab-ch3-so-sanh-framework-frontend}
\def\LTcaptype{table}
\begin{longtable}{|L{0.195\linewidth}|L{0.323\linewidth}|L{0.192\linewidth}|L{0.244\linewidth}|}
\caption{So sánh các khung phát triển giao diện web}\label{tab:ch3-so-sanh-framework-frontend}\\
\hline \textbf{Tiêu chí} & \textbf{Next.js 15} & \textbf{Nuxt.js 3} & \textbf{Vite + React} \\
\\\hline
\endfirsthead
\caption[]{So sánh các khung phát triển giao diện web (tiếp theo)}\\
\hline \textbf{Tiêu chí} & \textbf{Next.js 15} & \textbf{Nuxt.js 3} & \textbf{Vite + React} \\
\\\hline
\endhead
\\\hline
\endfoot
\endlastfoot
Ngôn ngữ cơ sở & React 19 + TypeScript & Vue 3 + TypeScript & React +
TypeScript \\\hline
Cách dựng trang & Có thể dựng ở máy chủ, dựng sẵn khi build hoặc kết
hợp cả hai & Có thể dựng ở máy chủ hoặc dựng sẵn & Chủ yếu dựng ở phía
trình duyệt \\\hline
Định tuyến và bố cục nhiều tầng & Có sẵn App Router của Next.js & Có sẵn
định tuyến theo file & Không có sẵn, cần thư viện bổ sung \\\hline
Khả năng hiện nội dung sớm / hỗ trợ tìm kiếm & Tốt & Tốt & Kém hơn vì
phần lớn nội dung phải chờ trình duyệt dựng xong \\\hline
Hệ sinh thái & Rất lớn (React ecosystem) & Lớn (Vue ecosystem) & Lớn
(React ecosystem) \\\hline
Hỗ trợ kênh thời gian thực & Tốt (Socket.IO, Server Actions) & Tốt & Tốt \\\hline
Hiệu năng phát triển và build & Cao (công cụ dựng nhanh, hỗ trợ hiện sớm một phần trang) & Cao (bộ máy dựng của Nuxt) &
Cao (công cụ đóng gói nhanh Vite) \\\hline
Đóng gói để triển khai & Có sẵn Dockerfile, gói chạy độc lập & Có sẵn & Cần
cấu hình thủ công \\\hline
Cộng đồng \& tài liệu & Rất lớn & Lớn & Lớn \\\hline
\end{longtable}
}


\textbf{Kết luận lựa chọn:} Next.js 15 được chọn làm framework chính vì
nó tạo được thế cân bằng tốt giữa trải nghiệm của người dùng cuối và khả
năng tổ chức mã nguồn của đội phát triển. Các lý do cụ thể gồm:

\begin{itemize}
\tightlist
\item
  \textbf{Dựng một phần giao diện ngay từ máy chủ}: React Server Components
  (RSC), tức cơ chế cho phép xử lý và dựng sẵn một phần giao diện ở phía
  máy chủ, giúp giảm lượng JavaScript phải gửi về trình duyệt và tăng
  tốc độ tải trang đầu tiên.
\item
  \textbf{Cơ chế chia trang theo thư mục}: App Router, tức cách Next.js
  tổ chức đường dẫn và bố cục theo cấu trúc thư mục, giúp quản lý trang,
  bố cục lồng nhau và khu vực được bảo vệ trực quan hơn.
\item
  \textbf{Hệ sinh thái React:} Tận dụng được hệ sinh thái cực kỳ phong
  phú của React, bao gồm các thư viện cho lấy dữ liệu từ máy chủ
  (TanStack Query), quản lý trạng thái giao diện (Zustand), bản đồ
  (Leaflet) và biểu đồ (ECharts). Nói ngắn gọn, đây là nhóm công cụ đủ
  để giao diện vừa lấy dữ liệu lịch sử, vừa nhận dữ liệu mới, vừa hiển
  thị trực quan mà không phải tự viết mọi thứ từ đầu.
\item
  \textbf{TypeScript tích hợp:} Hỗ trợ TypeScript từ gốc, đảm bảo an
  toàn kiểu dữ liệu xuyên suốt ứng dụng.
\item
  \textbf{Sẵn sàng cho triển khai:} Có sẵn các cơ chế tối ưu cho môi
  trường vận hành thật như tối ưu ảnh, tách nhỏ mã nguồn khi tải và xuất gói
  chạy độc lập cho Docker.
\end{itemize}

\subparagraph{b) Thư viện giao diện}\label{b-thux1b0-viux1ec7n-giao-diux1ec7n-ui-component-library}

Hệ thống sử dụng shadcn/ui kết hợp với Radix UI làm nền tảng xây dựng
giao diện. Có thể hiểu đây là một bộ khung gồm các thành phần giao diện
cơ bản đã được chuẩn bị sẵn về khả năng dùng bàn phím, cách lấy nét và
độ nhất quán. Điểm đáng giá của lựa chọn này không nằm ở số lượng thành phần
dựng sẵn, mà ở khả năng giữ giao diện nhất quán trong khi nhóm vẫn kiểm
soát được cách từng thành phần được triển khai trong dự án.

\textbf{Đặc điểm nổi bật:}

\begin{itemize}
\tightlist
\item
  \textbf{Radix UI Primitives:} Cung cấp các thành phần cơ bản đã xử lý
  sẵn khả năng truy cập, điều hướng bằng bàn phím và quản lý điểm lấy nét.
\item
  \textbf{Tailwind CSS v4:} Hệ thống lớp tiện ích giúp thay đổi giao diện
  nhanh mà không phải viết nhiều CSS thủ công.
\item
  \textbf{Class Variance Authority (CVA):} Giúp quản lý các biến thể của
  cùng một thành phần giao diện, ví dụ nút chính, nút phụ hay nút cảnh báo.
\item
  \textbf{Hệ màu OKLCH:} Giúp điều chỉnh màu dễ kiểm soát hơn khi phải
  hỗ trợ cả giao diện sáng và tối.
\end{itemize}

\subparagraph{c) Quản lý trạng thái (State
Management)}\label{c-quux1ea3n-luxfd-trux1ea1ng-thuxe1i-state-management}

Với một giao diện quản trị vừa có dữ liệu tĩnh, dữ liệu thời gian thực và trạng
thái giao diện, dồn mọi thứ vào cùng một nơi quản lý sẽ làm mã nguồn
nhanh chóng rối lên. Vì vậy, hệ thống áp dụng chiến lược tách biệt giữa
trạng thái cục bộ của trình duyệt và trạng thái lấy từ máy chủ:


{\thesissettableid{3.25}{tab-ch3-chien-luoc-quan-ly-trang-thai}
\def\LTcaptype{table}
\begin{longtable}{|L{0.177\linewidth}|L{0.215\linewidth}|L{0.235\linewidth}|L{0.327\linewidth}|}
\caption{Chiến lược quản lý trạng thái}\label{tab:ch3-chien-luoc-quan-ly-trang-thai}\\
\hline \textbf{Loại trạng thái} & \textbf{Thư viện} & \textbf{Mục đích} & \textbf{Ví dụ} \\
\\\hline
\endfirsthead
\caption[]{Chiến lược quản lý trạng thái (tiếp theo)}\\
\hline \textbf{Loại trạng thái} & \textbf{Thư viện} & \textbf{Mục đích} & \textbf{Ví dụ} \\
\\\hline
\endhead
\\\hline
\endfoot
\endlastfoot
Trạng thái giao diện cục bộ & Zustand 5 & Những gì chỉ thuộc về màn hình
đang mở & Token phiên đăng nhập, trạng thái sidebar, theme \\\hline
Trạng thái dữ liệu từ máy chủ & TanStack Query 5 & Dữ liệu lấy từ API
phía máy chủ và cần tái sử dụng & Danh sách xe, chuyến đi, cảnh báo \\\hline
Trạng thái gắn với URL & nuqs & Bộ lọc hoặc tab cần giữ lại ngay trên
đường dẫn & Bộ lọc, phân trang, tab đang chọn \\\hline
Trạng thái thời gian thực & Socket.IO Client & Dữ liệu mới được máy chủ
đẩy xuống ngay lập tức & Vị trí xe, dữ liệu vận hành gửi liên tục \\\hline
\end{longtable}
}


\textbf{Kết luận:} Zustand được chọn vì API đơn giản, kích thước cực nhỏ
(~1.1 KB) và phù hợp với yêu cầu của hệ thống nơi token
phiên đăng nhập chỉ lưu trong bộ nhớ tạm thời, không ghi ra
vùng nhớ lâu trên trình duyệt như \texttt{localStorage}, nhằm tăng tính bảo mật.

\subparagraph{d) Thư viện bản
đồ}\label{d-thux1b0-viux1ec7n-bux1ea3n-ux111ux1ed3}


{\thesissettableid{3.26}{tab-ch3-so-sanh-thu-vien-ban-o}
\def\LTcaptype{table}
\begin{longtable}{|L{0.248\linewidth}|L{0.285\linewidth}|L{0.229\linewidth}|L{0.192\linewidth}|}
\caption{So sánh thư viện bản đồ}\label{tab:ch3-so-sanh-thu-vien-ban-o}\\
\hline \textbf{Tiêu chí} & \textbf{Leaflet.js} & \textbf{Mapbox GL JS} & \textbf{Google Maps JS API} \\
\\\hline
\endfirsthead
\caption[]{So sánh thư viện bản đồ (tiếp theo)}\\
\hline \textbf{Tiêu chí} & \textbf{Leaflet.js} & \textbf{Mapbox GL JS} & \textbf{Google Maps JS API} \\
\\\hline
\endhead
\\\hline
\endfoot
\endlastfoot
Giấy phép & MIT (mã nguồn mở) & Sở hữu riêng (có gói miễn phí giới hạn) & Sở hữu riêng
(trả phí) \\\hline
Chi phí & Miễn phí & Miễn phí đến 50K loads/tháng & 0.007 USD/load sau
28K \\\hline
Hiệu năng với nhiều điểm xe & Tốt (với MarkerCluster, tức gom cụm các điểm gần nhau) & Rất tốt (WebGL
rendering, tức dựng hình bằng bộ xử lý đồ họa) & Tốt \\\hline
Tùy chỉnh điểm đánh dấu & Cao & Cao & Trung bình \\\hline
Tự lưu và tự phục vụ bản đồ nền & Hỗ trợ & Không & Không \\\hline
Tích hợp với React & react-leaflet (trưởng thành) & react-map-gl &
@react-google-maps/api \\\hline
Vẽ vùng cho phép & @geoman-io/leaflet-geoman-free & mapbox-gl-draw & Drawing
Manager \\\hline
\end{longtable}
}


\textbf{Kết luận:} Leaflet.js được chọn vì hoàn toàn miễn phí, hỗ trợ tự
lưu và phân phối bản đồ nền, và có hệ sinh thái plugin phong phú
(MarkerCluster để gom nhiều điểm xe gần nhau, Geoman để vẽ và sửa vùng giám sát).
Điều này đặc biệt phù hợp với hệ thống IoT có thể triển khai trong mạng
nội bộ hoặc môi trường không muốn phụ thuộc dịch vụ bản đồ bên ngoài.

\clearpage

\paragraph{3.4.2.2. Kiến trúc ứng dụng web}\label{3242-kiux1ebfn-truxfac-ux1ee9ng-dux1ee5ng-web-web-application-architecture}

\subparagraph{a) Mô hình kiến trúc chia theo từng nhóm tính năng}\label{a-muxf4-huxecnh-kiux1ebfn-truxfac-feature-sliced-design}

Ứng dụng giao diện web được tổ chức theo cách chia mã nguồn bám theo
từng nhóm tính năng nghiệp vụ, thường gọi là Feature-Sliced Design
(FSD). Mỗi tính năng vì thế được đóng gói thành một module tương đối độc
lập gồm thành phần giao diện, hook, kiểu dữ liệu và tiện ích liên quan.
Nói gọn hơn, FSD chia mã nguồn theo ``thứ người dùng đang làm'' thay vì
chỉ chia theo loại file. Ví dụ, phần quản lý xe sẽ gom mọi phần liên
quan đến xe vào cùng một cụm; nhờ đó người mới đọc mã nguồn dễ lần theo
nghiệp vụ hơn, còn việc bổ sung tính năng mới cũng ít làm rối cấu trúc
cũ.

\begin{itemize}
\tightlist
\item
  \texttt{app/}: khung chia trang của Next.js, bao gồm bố cục gốc và các khu
  vực quản trị được bảo vệ.
\item
  \texttt{features/}: các module nghiệp vụ theo FSD (vehicles, trips,
  alerts, map), mỗi module tách rõ thành phần giao diện, hook và kiểu dữ liệu.
\item
  \texttt{components/}: thư viện thành phần dùng chung (UI, bố cục, bản đồ,
  biểu đồ, biểu mẫu).
\item
  \texttt{lib/}: hạ tầng ứng dụng (lớp gọi API, kênh thời gian thực, nơi giữ trạng thái, tiện ích dùng chung).
\item
  \texttt{types/}: các kiểu TypeScript dùng toàn cục.
\end{itemize}


{\thesissetfigureid{3.15}{fig-ch3-so-o-kien-truc-feature-sliced-design-cua-ung-dung-frontend}
\begin{figure}[htbp]
\centering
\pandocbounded{\safeincludesvg{./assets/figures/06-chuong-3-giai-phap-frontend-hinh-3-15.svg}}
\caption{Sơ đồ kiến trúc chia theo từng nhóm tính năng của ứng dụng giao diện web}
\label{fig:ch3-so-o-kien-truc-feature-sliced-design-cua-ung-dung-frontend}
\end{figure}
}


\textbf{Nguyên tắc tổ chức:}

\begin{itemize}
\tightlist
\item
  \textbf{Đặt gần theo tính năng}: Các file liên quan đến một tính năng được đặt
  cùng thư mục, giảm thời gian tìm kiếm và bảo trì.
\item
  \textbf{Bao gói}: Mỗi module nghiệp vụ chỉ công bố ra ngoài phần
  thật sự cần dùng; các phần còn lại được giữ bên trong để tránh phụ
  thuộc chéo.
\item
  \textbf{Tách rõ trách nhiệm}: Tách rõ nơi điều hướng trang, nơi
  chứa logic nghiệp vụ, nơi chứa giao diện dùng lại và nơi chứa hạ tầng
  kỹ thuật.
\end{itemize}

\subparagraph{b) Luồng dữ liệu trong ứng
dụng}\label{b-luux1ed3ng-dux1eef-liux1ec7u-trong-ux1ee9ng-dux1ee5ng}

Luồng dữ liệu trong ứng dụng giao diện web tuân theo mô hình một chiều:
dữ liệu đi từ nơi lấy dữ liệu đến nơi hiển thị theo một hướng rõ ràng,
đồng thời kết hợp ba nguồn dữ liệu chính:

\begin{enumerate}
\def\labelenumi{\arabic{enumi}.}
\tightlist
\item
  \textbf{REST API (TanStack Query):} Dữ liệu CRUD, gồm các thao tác tạo,
  đọc, sửa, xóa cơ bản như danh sách xe, chuyến đi hay cảnh báo, được
  lấy từ API phía máy chủ qua HTTP. TanStack Query giúp lưu tạm và làm
  mới dữ liệu theo cơ chế ``hiển thị ngay dữ liệu đang có rồi âm thầm hỏi
  bản mới'', nhờ đó người dùng không phải chờ trắng màn hình mỗi lần tải
  lại.
\item
  \textbf{Trạng thái chỉ phục vụ hiển thị (Zustand):} Các trạng thái chỉ
  liên quan đến giao diện như phiên đăng nhập hiện tại, sidebar hay
  theme được quản lý ngay trên
  trình duyệt bằng Zustand mà không cần gọi API.
\item
  \textbf{Dữ liệu cập nhật tức thời (Socket.IO):} Dữ liệu vị trí xe và dữ
  liệu vận hành gửi liên tục được nhận theo thời gian thực qua
  WebSocket. Socket.IO là lớp giúp mở và giữ kênh để máy chủ chủ động
  gửi dữ liệu mới lên trình duyệt. Khi kết nối WebSocket đang hoạt động,
  TanStack Query sẽ tự động tắt cơ chế hỏi lặp lại, tránh nhận trùng một
  dữ liệu theo hai đường.
\end{enumerate}


{\thesissetfigureid{3.15a}{fig-ch3-luong-du-lieu-mot-chieu-giua-next-js-query-zustand-va-backend}
\begin{figure}[htbp]
\centering
\pandocbounded{\safeincludesvgsmall{./assets/figures/06-chuong-3-giai-phap-frontend-hinh-3-15a.svg}}
\caption{Luồng dữ liệu một chiều giữa Next.js, Query, Zustand và API phía máy chủ}
\label{fig:ch3-luong-du-lieu-mot-chieu-giua-next-js-query-zustand-va-backend}
\end{figure}
}

\subparagraph{c) Cơ chế xác thực và bảo vệ
đường dẫn truy cập}\label{c-cux1a1-chux1ebf-xuxe1c-thux1ef1c-vuxe0-bux1ea3o-vux1ec7-route}

Hệ thống giao diện web sử dụng một mã phiên ngẫu nhiên do API phía máy
chủ cấp; bản thân mã này không chứa thông tin người dùng. Cách làm này
khác với JWT tự mang thông tin phiên. Việc trình duyệt gửi token trong
header \texttt{Authorization:\ Bearer} chỉ là cách mang token theo từng
yêu cầu; về bản chất, đây vẫn là cơ chế giữ phiên trên máy chủ vì mọi
quyền truy cập đều chỉ được quyết định sau khi máy chủ tra cứu phiên
tương ứng:

\begin{enumerate}
\def\labelenumi{\arabic{enumi}.}
\tightlist
\item
  \textbf{Đăng nhập:} Người dùng nhập email và mật khẩu, API phía máy chủ tạo bản
  ghi phiên và trả về một mã phiên ngẫu nhiên.
\item
  \textbf{Lưu trữ token:} Token chỉ được giữ trong bộ nhớ tạm của Zustand,
  không lưu vào \texttt{localStorage} hay \texttt{sessionStorage}, là
  hai nơi lưu dữ liệu bền hơn trên trình duyệt, nhằm giảm nguy cơ mã độc
  chèn vào trang rồi lấy cắp dữ liệu.
\item
  \textbf{Gửi yêu cầu:} Mọi yêu cầu HTTP tự động đính kèm token vào
  header
  \texttt{Authorization:\ Bearer\ \textless{}token\textgreater{}}.
\item
  \textbf{Lớp chặn truy cập}: Thành phần \texttt{AuthGuard} bao bọc toàn bộ khu
  vực \texttt{/dashboard/*}, tự động chuyển hướng về trang đăng nhập nếu
  chưa xác thực hoặc token hết hạn.
\item
  \textbf{Gia hạn phiên:} Khi phiên sắp hết hạn, giao diện web chủ động gọi
  cơ chế gia hạn để nhận token mới, nhờ đó người dùng không bị đăng nhập
  lại đột ngột trong lúc thao tác.
\end{enumerate}


{\thesissetfigureid{3.16}{fig-ch3-bieu-o-trinh-tu-sequence-diagram-luong-xac-thuc-nguoi-dung}
\begin{figure}[htbp]
\centering
\pandocbounded{\safeincludesvg{./assets/figures/06-chuong-3-giai-phap-frontend-hinh-3-16.svg}}
\caption{Biểu đồ trình tự luồng xác thực người dùng}
\label{fig:ch3-bieu-o-trinh-tu-sequence-diagram-luong-xac-thuc-nguoi-dung}
\end{figure}
}


\paragraph{3.4.2.3. Thiết kế các trang chức năng chính}\label{3243-thiux1ebft-kux1ebf-cuxe1c-trang-chux1ee9c-nux103ng-chuxednh-main-feature-pages}

\subparagraph{a) Trang tổng quan}\label{a-trang-tux1ed5ng-quan-dashboard}

Trang tổng quan là điểm vào chính sau khi đăng nhập, cung cấp cái nhìn
tổng thể về trạng thái hệ thống:

\begin{itemize}
\tightlist
\item
  \textbf{Các thẻ chỉ số nhanh:} Hiển thị các chỉ số quan trọng: tổng số xe đang
  hoạt động, số chuyến đi trong ngày, số cảnh báo chưa xử lý, số vi phạm
  má»›i.
\item
  \textbf{Bản đồ mini:} Hiển thị vị trí tất cả xe đang hoạt động trên
  bản đồ thu nhỏ.
\item
  \textbf{Bảng cảnh báo gần đây:} Danh sách 10 cảnh báo mới nhất với
  nhãn mức độ nghiêm trọng (thấp, trung bình, cao, rất cao).
\item
  \textbf{Biểu đồ thống kê:} Biểu đồ số chuyến đi theo ngày và phân loại
  cảnh báo theo loại.
\end{itemize}


{\thesissetfigureid{3.17}{fig-ch3-giao-dien-trang-dashboard-tong-quan}
\begin{figure}[htbp]
\centering
\pandocbounded{\safeincludesvg{./assets/figures/06-chuong-3-giai-phap-frontend-hinh-3-17.svg}}
\caption{Giao diện trang tổng quan hệ thống}
\label{fig:ch3-giao-dien-trang-dashboard-tong-quan}
\end{figure}
}


\subparagraph{b) Quản lý phương tiện}\label{b-quux1ea3n-luxfd-phux1b0ux1a1ng-tiux1ec7n-vehicle-management}

\textbf{Trang danh sách (\texttt{/dashboard/vehicles}):}


{\thesissettableid{3.2.4A}{tab-ch3-tong-hop-thong-tin-ve-cot-mo-ta-chuc-nang}
\def\LTcaptype{table}
\begin{longtable}{|L{0.212\linewidth}|L{0.519\linewidth}|L{0.234\linewidth}|}
\caption{Tổng hợp thông tin về Cột, Mô tả, Chức năng}\label{tab:ch3-tong-hop-thong-tin-ve-cot-mo-ta-chuc-nang}\\
\hline \textbf{Cột} & \textbf{Mô tả} & \textbf{Chức năng} \\
\\\hline
\endfirsthead
\caption[]{Tổng hợp thông tin về Cột, Mô tả, Chức năng (tiếp theo)}\\
\hline \textbf{Cột} & \textbf{Mô tả} & \textbf{Chức năng} \\
\\\hline
\endhead
\\\hline
\endfoot
\endlastfoot
Biển số xe & Số đăng ký phương tiện & Tìm kiếm \\\hline
Hãng/Model & Thương hiệu và đời xe & Lọc \\\hline
Trạng thái & Đang hoạt động, tạm ngưng, bảo trì, ngừng sử dụng & Lọc theo trạng
thái \\\hline
Thiết bị & Mã thiết bị IoT gắn kèm & Liên kết \\\hline
Lần cuối hoạt động & Thời điểm nhận dữ liệu gần nhất & Sắp xếp \\\hline
Thao tác & Xem, Sửa, Xóa, Xem trên bản đồ & Nút hành động \\\hline
\end{longtable}
}


\textbf{Trang chi tiết (\texttt{/dashboard/vehicles/{[}id{]}}):} Hiển
thị thông tin đầy đủ của phương tiện bao gồm thông tin cơ bản, vị trí
hiện tại trên bản đồ, trạng thái thiết bị IoT, các cảnh báo đang hoạt
động, lịch sử chuyến đi và lịch sử bảo trì.

\clearpage


{\thesissetfigureid{3.18}{fig-ch3-giao-dien-trang-quan-ly-phuong-tien}
\begin{figure}[htbp]
\centering
\pandocbounded{\safeincludesvg{./assets/figures/06-chuong-3-giai-phap-frontend-hinh-3-18.svg}}
\caption{Giao diện trang quản lý phương tiện}
\label{fig:ch3-giao-dien-trang-quan-ly-phuong-tien}
\end{figure}
}


\subparagraph{c) Quản lý chuyến đi, cảnh báo và vùng địa
lý}\label{c-quux1ea3n-luxfd-chuyux1ebfn-ux111i-cux1ea3nh-buxe1o-vuxe0-vuxf9ng-ux111ux1ecba-luxfd}

Trang quản lý chuyến đi cung cấp lịch sử và chi tiết từng chuyến, gồm
tuyến đường trên bản đồ, biểu đồ tốc độ theo thời gian và timeline sự
kiện. Trang cảnh báo tổng hợp cảnh báo theo nhãn mức độ nghiêm trọng và
hỗ trợ xử lý hàng loạt. Trang quản lý vùng cho phép hỗ trợ vẽ hình tròn, đa giác
hoặc hình chữ nhật trực tiếp trên bản đồ bằng Leaflet Geoman.


{\thesissettableid{3.27}{tab-ch3-ma-tran-cac-trang-chuc-nang-chinh}
\def\LTcaptype{table}
\begin{longtable}{|L{0.179\linewidth}|L{0.158\linewidth}|L{0.385\linewidth}|L{0.233\linewidth}|}
\caption{Ma trận các trang chức năng chính}\label{tab:ch3-ma-tran-cac-trang-chuc-nang-chinh}\\
\hline \textbf{Trang} & \textbf{Đường dẫn} & \textbf{Chức năng chính} & \textbf{Cập nhật tức thời} \\
\\\hline
\endfirsthead
\caption[]{Ma trận các trang chức năng chính (tiếp theo)}\\
\hline \textbf{Trang} & \textbf{Đường dẫn} & \textbf{Chức năng chính} & \textbf{Cập nhật tức thời} \\
\\\hline
\endhead
\\\hline
\endfoot
\endlastfoot
Tổng quan & \texttt{/dashboard} & Tổng quan hệ thống, các thẻ chỉ số chính, biểu
đồ & Có (cập nhật số liệu) \\\hline
Bản đồ & \nolinkurl{/dashboard/map} & Theo dõi tất cả xe thời gian thực &
Có (vị trí xe) \\\hline
Phương tiện & \nolinkurl{/dashboard/vehicles} & CRUD xe, trạng thái, lịch
sử & Có (trạng thái xe) \\\hline
Chuyến đi & \nolinkurl{/dashboard/trips} & Lịch sử chuyến đi, tuyến đường &
Không \\\hline
Cảnh báo & \nolinkurl{/dashboard/alerts} & Xem, xử lý cảnh báo & Có (cảnh
báo mới) \\\hline
Vi phạm & \nolinkurl{/dashboard/violations} & Quản lý vi phạm tốc độ &
Không \\\hline
Thiết bị & \nolinkurl{/dashboard/devices} & Quản lý thiết bị IoT & Có
(trạng thái thiết bị) \\\hline
Vùng địa lý & \nolinkurl{/dashboard/geofences} & Tạo, sửa vùng cho phép trên bản
đồ & Không \\\hline
Bảo trì & \nolinkurl{/dashboard/maintenance} & Lịch sử bảo trì phương tiện
& Không \\\hline
Thông báo & \nolinkurl{/dashboard/notifications} & Cấu hình Telegram, Email
& Không \\\hline
Cài đặt & \nolinkurl{/dashboard/settings} & Hồ sơ, quản lý người dùng &
Không \\\hline
\end{longtable}
}


\paragraph{3.4.2.4. Tích hợp bản đồ thời gian thực}\label{3244-tuxedch-hux1ee3p-bux1ea3n-ux111ux1ed3-thux1eddi-gian-thux1ef1c-real-time-map-integration}

\subparagraph{a) Kiến trúc tích hợp bản
đồ}\label{a-kiux1ebfn-truxfac-tuxedch-hux1ee3p-bux1ea3n-ux111ux1ed3}

Trang bản đồ thời gian thực (\texttt{/dashboard/map}) là tính năng cốt
lõi của hệ thống giám sát phương tiện. Kiến trúc tích hợp bản đồ bao gồm
ba lớp chính:


{\thesissetfigureid{3.19}{fig-ch3-kien-truc-tich-hop-ban-o-thoi-gian-thuc}
\begin{figure}[htbp]
\centering
\pandocbounded{\safeincludesvg{./assets/figures/06-chuong-3-giai-phap-frontend-hinh-3-19.svg}}
\caption{Kiến trúc tích hợp bản đồ thời gian thực}
\label{fig:ch3-kien-truc-tich-hop-ban-o-thoi-gian-thuc}
\end{figure}
}

{\thesissetfigureid{3.19a}{fig-ch3-phan-ra-lop-ban-o-react-leaflet-va-kenh-socket-io}
\begin{figure}[htbp]
\centering
\pandocbounded{\safeincludesvg{./assets/figures/06-chuong-3-giai-phap-frontend-hinh-3-19a.svg}}
\caption{Phân rã các lớp của bản đồ React Leaflet và kênh Socket.IO}
\label{fig:ch3-phan-ra-lop-ban-o-react-leaflet-va-kenh-socket-io}
\end{figure}
}







\subparagraph{b) Cơ chế cập nhật vị trí thời gian
thá»±c}\label{b-cux1a1-chux1ebf-cux1eadp-nhux1eadt-vux1ecb-truxed-thux1eddi-gian-thux1ef1c}

Quy trình cập nhật vị trí xe trên bản đồ diễn ra như sau:

\begin{enumerate}
\def\labelenumi{\arabic{enumi}.}
\tightlist
\item
  \textbf{Thiết bị IoT} gửi dữ liệu GPS qua MQTT đến EMQX Broker.
\item
  \textbf{MQTT Bridge} nhận và xử lý dữ liệu, lưu vào VictoriaMetrics,
  đồng thời chuyển tiếp đến API phía máy chủ.
\item
  \textbf{API phía máy chủ} phát sự kiện \texttt{vehicle:\{id\}:location} qua
  Socket.IO đến tất cả trình duyệt đang đăng ký nhận dữ liệu.
\item
  \textbf{Giao diện web} nhận sự kiện, cập nhật tọa độ điểm đánh dấu trên bản đồ. Sử
  dụng kỹ thuật nội suy (interpolation) để điểm đánh dấu di chuyển mượt mà giữa
  hai điểm tọa độ liên tiếp.
\end{enumerate}

\textbf{Tối ưu hiệu năng bản đồ:}

\begin{itemize}
\tightlist
\item
  \textbf{Gom cụm điểm xe:} Sử dụng \texttt{react-leaflet-cluster} để
  gom nhóm các điểm xe gần nhau khi thu nhỏ bản đồ, giảm số lượng phần
  tử phải vẽ cùng lúc.
\item
  \textbf{Chỉ vẽ phần đang nhìn thấy}: Chỉ vẽ điểm đánh dấu của các xe nằm
  trong vùng nhìn hiện tại của bản đồ.
\item
  \textbf{Gom các cập nhật gần nhau}: Gom nhóm các cập nhật vị trí trong khoảng
  100ms để tránh re-render quá nhiều lần.
\end{itemize}

\subparagraph{c) Các thành phần bản
đồ}\label{c-cuxe1c-thuxe0nh-phux1ea7n-bux1ea3n-ux111ux1ed3}

Hệ thống bao gồm bốn thành phần bản đồ chuyên biệt:


{\thesissettableid{3.2.4B}{tab-ch3-tong-hop-thong-tin-ve-component-file-chuc-nang}
\def\LTcaptype{table}
\begin{longtable}{|L{0.224\linewidth}|L{0.160\linewidth}|L{0.580\linewidth}|}
\caption{Tổng hợp thông tin về thành phần bản đồ, file và chức năng}\label{tab:ch3-tong-hop-thong-tin-ve-component-file-chuc-nang}\\
\hline \textbf{Thành phần} & \textbf{File} & \textbf{Chức năng} \\
\\\hline
\endfirsthead
\caption[]{Tổng hợp thông tin về thành phần bản đồ, file và chức năng (tiếp theo)}\\
\hline \textbf{Thành phần} & \textbf{File} & \textbf{Chức năng} \\
\\\hline
\endhead
\\\hline
\endfoot
\endlastfoot
VehicleMap & \nolinkurl{components/map/vehicle-map.tsx} & Hiển thị tất cả
xe trên bản đồ với điểm đánh dấu và ô thông tin \\\hline
RouteMap & \nolinkurl{components/map/route-map.tsx} & Hiển thị tuyến đường
chuyến đi với đường nối và các điểm dừng \\\hline
GeofenceMap & \nolinkurl{components/map/geofence-map.tsx} & Vẽ và chỉnh sửa
vùng địa lý (đa giác, hình tròn, hình chữ nhật) \\\hline
MarkerPopup & \nolinkurl{components/map/marker-popup.tsx} & Ô thông tin hiển thị
chi tiết xe khi bấm vào điểm đánh dấu (biển số, tốc độ, trạng thái) \\\hline
\end{longtable}
}



{\thesissetfigureid{3.20}{fig-ch3-giao-dien-trang-ban-o-thoi-gian-thuc-voi-cac-marker-xe}
\begin{figure}[htbp]
\centering
\pandocbounded{\safeincludesvg{./assets/figures/06-chuong-3-giai-phap-frontend-hinh-3-20.svg}}
\caption{Giao diện trang bản đồ thời gian thực với các marker xe}
\label{fig:ch3-giao-dien-trang-ban-o-thoi-gian-thuc-voi-cac-marker-xe}
\end{figure}
}

\clearpage


\paragraph{3.4.2.5. Thiết kế giao diện thích ứng và trải nghiệm người dùng}\label{3245-thiux1ebft-kux1ebf-responsive-vuxe0-trux1ea3i-nghiux1ec7m-ngux1b0ux1eddi-duxf9ng-responsive-design--ux}

\subparagraph{a) Chiến lược
thích ứng theo màn hình}\label{a-chiux1ebfn-lux1b0ux1ee3c-responsive}

Ở giai đoạn hiện tại, giao diện được ưu tiên cho desktop và tablet,
nhưng vẫn giữ ba mức ngưỡng chiều rộng màn hình chính để không phá vỡ
trải nghiệm trên màn hình nhỏ:


{\thesissettableid{3.28}{tab-ch3-breakpoint-va-tuong-ung-giao-dien}
\def\LTcaptype{table}
\begin{longtable}{|L{0.180\linewidth}|L{0.151\linewidth}|L{0.176\linewidth}|L{0.448\linewidth}|}
\caption{Ngưỡng chiều rộng màn hình và cách điều chỉnh giao diện}\label{tab:ch3-breakpoint-va-tuong-ung-giao-dien}\\
\hline \textbf{Ngưỡng} & \textbf{Kích thước} & \textbf{Giao diện} & \textbf{Điều chỉnh chính} \\
\\\hline
\endfirsthead
\caption[]{Ngưỡng chiều rộng màn hình và cách điều chỉnh giao diện (tiếp theo)}\\
\hline \textbf{Ngưỡng} & \textbf{Kích thước} & \textbf{Giao diện} & \textbf{Điều chỉnh chính} \\
\\\hline
\endhead
\\\hline
\endfoot
\endlastfoot
Mobile & \textless{} 640px & Giao diện tối giản, một cột & Sidebar thành
drawer, bảng thành danh sách card, bản đồ toàn màn hình \\\hline
Tablet & 640px - 1024px & Giao diện trung gian & Sidebar thu gọn
(icon-only), bảng giữ nguyên với cuộn ngang \\\hline
Desktop & \textgreater{} 1024px & Giao diện đầy đủ & Sidebar mở rộng,
bảng đầy đủ cột, bản đồ chia đôi màn hình \\\hline
\end{longtable}
}


\subparagraph{b) Bố cục chính}\label{b-layout-chuxednh}

Giao diện tổng quan sử dụng bố cục gồm ba vùng chính: Header
(đường dẫn đang đứng, ô tìm kiếm, menu người dùng, nút đổi giao diện sáng/tối), Sidebar (menu điều hướng
bên trái), và vùng nội dung chính (nội dung trang, bộ lọc, bảng dữ liệu).


{\thesissetfigureid{3.21}{fig-ch3-wireframe-bo-cuc-giao-dien-dashboard}
\begin{figure}[htbp]
\centering
\pandocbounded{\safeincludesvg{./assets/figures/06-chuong-3-giai-phap-frontend-hinh-3-21.svg}}
\caption{Wireframe bố cục giao diện tổng quan}
\label{fig:ch3-wireframe-bo-cuc-giao-dien-dashboard}
\end{figure}
}


\subparagraph{c) Hệ thống theme
Light/Dark}\label{c-hux1ec7-thux1ed1ng-theme-lightdark}

Hệ thống hỗ trợ hai chế độ giao diện sáng và tối, sử dụng CSS
Variables với hệ màu oklch. Chuyển đổi được thực hiện qua thư viện
\texttt{next-themes} và lưu trong cookie.

\subparagraph{d) Các mẫu thiết kế giao diện (UI
Patterns)}\label{d-cuxe1c-mux1eabu-thiux1ebft-kux1ebf-giao-diux1ec7n-ui-patterns}

Hệ thống áp dụng các mẫu thiết kế giao diện nhất quán xuyên suốt ứng
dụng:

\begin{itemize}
\tightlist
\item
  \textbf{Bảng dữ liệu:} Sắp xếp, lọc, phân trang ở phía
  máy chủ, chọn nhiều dòng, xuất dữ liệu CSV/Excel.
\item
  \textbf{Biểu mẫu nhập liệu:} React Hook Form kết hợp Zod để kiểm tra dữ
  liệu nhập ngay tại chỗ và báo lỗi ngay dưới trường nhập.
\item
  \textbf{Bản đồ:} Điểm đánh dấu đổi màu theo trạng thái, cửa sổ
  thông tin khi bấm vào xe và bộ công cụ vẽ vùng cho phép.
\item
  \textbf{Biểu đồ:} ECharts/Recharts cho dữ liệu vận hành, gồm biểu đồ đường
  cho tốc độ, đồng hồ hiển thị RPM và biểu đồ vùng cho nhiên liệu/nhiệt độ.
\end{itemize}


{\thesissetfigureid{3.22}{fig-ch3-cac-mau-thiet-ke-ui-cua-he-thong-data-table-form-charts}
\begin{figure}[htbp]
\centering
\pandocbounded{\safeincludesvg{./assets/figures/06-chuong-3-giai-phap-frontend-hinh-3-22.svg}}
\caption{Các mẫu thiết kế giao diện của hệ thống - bảng dữ liệu, biểu mẫu, biểu đồ}
\label{fig:ch3-cac-mau-thiet-ke-ui-cua-he-thong-data-table-form-charts}
\end{figure}
}


\subparagraph{e) Tối ưu trải nghiệm người
dùng}\label{e-tux1ed1i-ux1b0u-trux1ea3i-nghiux1ec7m-ngux1b0ux1eddi-duxf9ng}

Để nâng cao trải nghiệm người dùng, hệ thống áp dụng một số kỹ thuật:

\begin{itemize}
\tightlist
\item
  \textbf{Cập nhật tạm ngay trên giao diện:} Giao diện hiển thị kết quả trước
  khi nhận phản hồi cuối cùng từ máy chủ để thao tác không bị khựng.
\item
  \textbf{Khung chờ:} Hiển thị khung chờ thay vì vòng quay spinner khi
  đang tải dữ liệu.
\item
  \textbf{Thanh báo đang chuyển trang:} Thanh tiến trình ở đầu trang (sử dụng
  \texttt{nextjs-toploader}) khi chuyển trang.
\item
  \textbf{Phím tắt:} Hỗ trợ phím tắt cho các thao tác phổ biến
  (\texttt{Cmd+K} cho tìm kiếm).
\item
  \textbf{Vùng chặn lỗi riêng:} Mỗi khu vực chức năng được bao bọc bởi
  một lớp chặn lỗi riêng để hỏng chỗ nào thì cô lập chỗ đó.
\item
  \textbf{Trạng thái chưa có dữ liệu:} Hiển thị trạng thái trống có hướng dẫn hành
  động khi chưa có dữ liệu.
\end{itemize}

\begin{center}\rule{0.5\linewidth}{0.5pt}\end{center}

\subsection{3.5. Phân tích, đánh giá và lựa chọn phương án khả thi --
Analysis, Evaluation and
Selection}\label{33-phuxe2n-tuxedch-ux111uxe1nh-giuxe1-vuxe0-lux1ef1a-chux1ecdn-phux1b0ux1a1ng-uxe1n-khux1ea3-thi--analysis-evaluation-and-selection}

Dựa trên phân tích và đề xuất ở các mục 3.1 đến 3.4, bảng tổng hợp dưới
đây đánh giá và so sánh các phương án đã xem xét:


{\thesissettableid{3.29}{tab-ch3-ma-tran-anh-gia-tong-hop-cac-phuong-an-thiet-ke}
\def\LTcaptype{table}
\begin{longtable}{|L{0.185\linewidth}|L{0.320\linewidth}|L{0.459\linewidth}|}
\caption{Ma trận đánh giá tổng hợp các phương án thiết kế}\label{tab:ch3-ma-tran-anh-gia-tong-hop-cac-phuong-an-thiet-ke}\\
\hline \textbf{Tầng hệ thống} & \textbf{Phương án được chọn} & \textbf{Lý do chính} \\
\\\hline
\endfirsthead
\caption[]{Ma trận đánh giá tổng hợp các phương án thiết kế (tiếp theo)}\\
\hline \textbf{Tầng hệ thống} & \textbf{Phương án được chọn} & \textbf{Lý do chính} \\
\\\hline
\endhead
\\\hline
\endfoot
\endlastfoot
Vi điều khiển & ESP32-S3 & BLE 5.0 tích hợp, hai lõi xử lý, tài nguyên đủ cho nhiều tác vụ đồng thời \\\hline
LTE + GNSS & SIMCom SIM7600CE-T & Modem tích hợp LTE + GNSS, giảm số
lượng khối phần cứng rời \\\hline
OBD2 Adapter & vgate iCar Pro (BLE) & Không cần dây, tương thích rộng \\\hline
Cảm biến IMU & LIS3DSH & Tiết kiệm điện và có thể tự đánh thức khi phát
hiện chuyển động \\\hline
Buffer cục bộ & microSD SDMMC 4-bit + FATFS + phát lại dữ liệu & Duy trì liên tục
dữ liệu vận hành khi mất mạng, rồi gửi bù theo đúng thứ tự khi có mạng lại \\\hline
MQTT Broker & EMQX & Có lớp lọc bản tin, phân quyền theo từng thiết bị
và có thể mở rộng nếu số xe tăng \\\hline
Database quan hệ & PostgreSQL 16 & Ổn định, quen thuộc, phù hợp dữ liệu
có quan hệ chặt chẽ \\\hline
Database time-series & VictoriaMetrics & Ghi dữ liệu theo thời gian rất
nhanh và truy vấn tốt theo mốc thời gian \\\hline
API phía máy chủ & Express.js + TypeScript & Nhẹ, linh hoạt, dễ tổ chức theo
từng bài toán nghiệp vụ \\\hline
Giao diện web & Next.js 15 + React 19 & Hiện nội dung sớm, tổ chức trang
rõ và hệ sinh thái lớn \\\hline
Bản đồ & Leaflet.js & Mã nguồn mở, có thể tự lưu và phân phối bản đồ nền \\\hline
\end{longtable}
}


\subsection{3.6. Tối ưu phương án thiết kế -- The optimal
solution}\label{34-tux1ed1i-ux1b0u-phux1b0ux1a1ng-uxe1n-thiux1ebft-kux1ebf--the-optimal-solution}

Từ kết quả phân tích và đánh giá ở mục 3.5, phương án thiết kế tối ưu
cho hệ thống IoT giám sát phương tiện được tổng hợp như sau:


{\thesissetfigureid{3.23}{fig-ch3-so-o-kien-truc-tong-the-phuong-an-toi-uu}
\begin{figure}[htbp]
\centering
\pandocbounded{\safeincludesvg{./assets/figures/06-chuong-3-giai-phap-frontend-hinh-3-23.svg}}
\caption{Sơ đồ kiến trúc tổng thể phương án tối ưu}
\label{fig:ch3-so-o-kien-truc-tong-the-phuong-an-toi-uu}
\end{figure}
}


\textbf{Kiến trúc phân tầng:}

\begin{itemize}
\tightlist
\item
  \textbf{Tầng thiết bị}: ESP32-S3 + SIMCom SIM7600CE-T + vgate iCar Pro
  + LIS3DSH, phối hợp để lấy dữ liệu xe, lấy vị trí và quản lý nguồn với
  pin dự phòng 18650 1S
\item
  \textbf{Tầng truyền thông}: MQTT 3.1.1 qua 4G LTE, dùng mức cam kết gửi
  khác nhau cho từng loại bản tin, kết hợp vùng lưu tạm microSD và cơ chế
  gửi bù khi kết nối phục hồi
\item
  \textbf{Tầng xử lý}: MQTT Bridge nhận bản tin rồi ghi ra nhiều nơi phù
  hợp như PostgreSQL và VictoriaMetrics, sau đó API phía máy chủ xử lý
  nghiệp vụ cho người dùng
\item
  \textbf{Tầng trình bày}: Next.js 15, bản đồ Leaflet, kênh thời gian thực
  Socket.IO và biểu đồ ECharts để người vận hành quan sát dữ liệu
\end{itemize}

Phương án này đáp ứng các yêu cầu kỹ thuật đã đặt ra ở Chương 2, đồng
thời giữ chi phí phần cứng ở mức khoảng 1.514.000 VND cho mỗi thiết bị
theo mặt bằng giá tham chiếu tháng 4/2026.

\begin{center}\rule{0.5\linewidth}{0.5pt}\end{center}

\subsection{Kết luận chương
3}\label{kux1ebft-luux1eadn-chux1b0ux1a1ng-3}

Chương 3 đã trả lời câu hỏi "thiết kế hệ thống này như thế nào để vừa
đáp ứng yêu cầu kỹ thuật, vừa khả thi khi triển khai". Từ phần cứng,
phần mềm trên thiết bị, hạ tầng máy chủ xử lý dữ liệu đến giao diện web,
các phương án đều đã được cân nhắc theo tiêu chí hiệu năng, chi phí, độ
phức tạp và khả năng mở rộng.

Vì vậy, phương án cuối cùng ở Chương 3 không xuất hiện như một lựa chọn
cảm tính, mà như kết quả của quá trình sàng lọc có cơ sở. Trên nền tảng
đó, Chương 4 chuyển sang phần hiện thực: biến các lựa chọn thiết kế
thành hệ thống vận hành được trong thực tế.

\chapterheading{CHƯƠNG 4. TRIỂN KHAI GIẢI PHÁP VÀ KẾT QUẢ -- IMPLEMENTATION AND RESULTS}

Nếu Chương 3 dừng ở mức lựa chọn phương án, thì Chương 4 trả lời câu
hỏi liệu những lựa chọn đó có thực sự hoạt động khi đưa vào chế tạo, lập
trình và tích hợp hay không. Nội dung chương đi từ phần cứng thiết bị,
firmware, hạ tầng máy chủ, giao diện giám sát đến các kết quả đo lường
sau khi lắp đặt hệ thống trên xe, qua đó giúp người đọc đối chiếu giữa ý
tưởng thiết kế và kết quả triển khai thực tế.

\begin{center}\rule{0.5\linewidth}{0.5pt}\end{center}

\subsection{4.1. Thiết kế chi tiết giải pháp -- Detailed design
solution}\label{41-thiux1ebft-kux1ebf-chi-tiux1ebft-giux1ea3i-phuxe1p--detailed-design-solution}

\subsubsection{4.1.1. Thiết kế chi tiết phần
cứng}\label{411-thiux1ebft-kux1ebf-chi-tiux1ebft-phux1ea7n-cux1ee9ng}

\paragraph{a) Kiến trúc tổng
thể}\label{a-kiux1ebfn-truxfac-tux1ed5ng-thux1ec3}

Xét theo chức năng, ESP32-S3 là bộ điều phối trung tâm của thiết bị theo dõi. Nó
thu tín hiệu từ các khối khác, quyết định thiết bị đang ở chế độ nào và
ra lệnh cho từng phần cứng hoạt động đúng lúc. Cụ thể hơn, ESP32-S3 nói
chuyện với modem qua UART, tức cổng nối tiếp cơ bản; lấy dữ liệu xe qua
BLE, tức Bluetooth tiết kiệm năng lượng; đọc cảm biến rung qua I2C, tức
đường giao tiếp ngắn cho cảm biến; và đo điện áp ắc quy qua ADC, tức
bộ chuyển đổi từ tín hiệu điện áp analog sang số. Toàn bộ thiết bị được
nuôi bởi mạch quản lý năng lượng có khả năng tự chuyển giữa ắc quy xe và
pin dự phòng.


{\thesissetfigureid{4.1}{fig-ch4-so-o-khoi-tong-the-he-thong-tracker-iot}
\begin{figure}[htbp]
\centering
\pandocbounded{\safeincludesvg{./assets/figures/07-chuong-4-trien-khai-hardware-hinh-4-1.svg}}
\caption{Sơ đồ khối tổng thể hệ thống thiết bị theo dõi IoT}
\label{fig:ch4-so-o-khoi-tong-the-he-thong-tracker-iot}
\end{figure}
}


{\thesissetfigureid{4.1a}{fig-ch4-bo-tri-chi-tiet-cac-khoi-phan-cung-trong-tracker}
\begin{figure}[htbp]
\centering
\pandocbounded{\safeincludesvg{./assets/figures/07-chuong-4-trien-khai-hardware-hinh-4-1a.svg}}
\caption{Bố trí chi tiết các khối phần cứng trong thiết bị theo dõi}
\label{fig:ch4-bo-tri-chi-tiet-cac-khoi-phan-cung-trong-tracker}
\end{figure}
}



\paragraph{b) Khối vi điều khiển và giao
tiếp}\label{b-khux1ed1i-vi-ux111iux1ec1u-khiux1ec3n-vuxe0-giao-tiux1ebfp}

Vi điều khiển ESP32-S3-WROOM-1 được lựa chọn làm nhân xử lý trung tâm
của hệ thống với các thông số kỹ thuật nổi bật:


{\thesissettableid{4.1.1A}{tab-ch4-tong-hop-thong-tin-ve-thong-so-gia-tri}
\def\LTcaptype{table}
\begin{longtable}{|L{0.373\linewidth}|L{0.595\linewidth}|}
\caption{Tổng hợp thông tin về Thông số, Giá trị}\label{tab:ch4-tong-hop-thong-tin-ve-thong-so-gia-tri}\\
\hline \textbf{Thông số} & \textbf{Giá trị} \\
\\\hline
\endfirsthead
\caption[]{Tổng hợp thông tin về Thông số, Giá trị (tiếp theo)}\\
\hline \textbf{Thông số} & \textbf{Giá trị} \\
\\\hline
\endhead
\\\hline
\endfoot
\endlastfoot
CPU & Dual-core Xtensa LX7 @ 240 MHz (32-bit) \\\hline
RAM & 512 KB SRAM \\\hline
Flash & 4--16 MB (tùy variant) \\\hline
Bluetooth & BLE 5.0 (tương thích ngược BLE 4.0) \\\hline
GPIO & 45 chân \\\hline
UART & 3 cổng \\\hline
I2C & 2 cổng \\\hline
ADC & 12-bit, 20 kênh \\\hline
Deep Sleep & ~10--15 µA (với external wakeup) \\\hline
Active (BLE) & ~20--40 mA \\\hline
\end{longtable}
}


\textbf{Giao tiếp UART với modem SIM7600CE-T (LTE + GNSS tích hợp):}
Khối modem gói hai nhiệm vụ quan trọng vào cùng một phần cứng: đưa dữ
liệu lên mạng di động và cung cấp vị trí GNSS của xe. ESP32-S3 trao đổi
với modem qua UART1 bằng tập lệnh AT, tức là các lệnh văn bản để yêu cầu
modem đăng ký mạng, mở kết nối dữ liệu và đọc thông tin định vị. Nhờ
GNSS đã nằm sẵn trong modem, thiết kế không cần thêm một module định vị
riêng, nên đi dây gọn hơn và việc tích hợp firmware cũng đơn giản hơn.

\textbf{Giao tiếp BLE với OBD2 adapter:} ESP32-S3 sử dụng BLE 5.0 tích
hợp để kết nối với adapter vgate iCar Pro (BLE 4.0). Nhờ kết nối này,
thiết bị theo dõi không chỉ biết xe đang ở đâu mà còn biết xe đang ở trạng thái
nào thông qua các dữ liệu OBD-II như trạng thái khóa điện (IGN), tốc độ
động cơ (RPM), vận tốc xe, mức nhiên liệu và mã lỗi chẩn đoán (DTC).

\textbf{Giao tiếp I2C với cảm biến LIS3DSH:} Cảm biến gia tốc 3 trục
LIS3DSH được kết nối qua bus I2C (GPIO2 SDA, GPIO1 SCL). Vai trò chính
của cảm biến này là giúp hệ thống "ngủ mà vẫn canh", tức là phát hiện
chuyển động khi xe đang đỗ và đánh thức ESP32-S3 từ light sleep thông
qua GPIO wake khi có rung động bất thường; heartbeat định kỳ vẫn dùng timer wake.


{\thesissetfigureid{4.2}{fig-ch4-so-o-ket-noi-giua-esp32-s3-va-modem-sim7600ce-t-qua}
\begin{figure}[htbp]
\centering
\pandocbounded{\safeincludesvg{./assets/figures/07-chuong-4-trien-khai-hardware-hinh-4-2.svg}}
\caption{Sơ đồ kết nối giữa ESP32-S3 và modem SIM7600CE-T qua UART}
\label{fig:ch4-so-o-ket-noi-giua-esp32-s3-va-modem-sim7600ce-t-qua}
\end{figure}
}


{\thesissetfigureid{4.2a}{fig-ch4-so-o-chan-ket-noi-sim7600ce-t-voi-esp32-s3}
\begin{figure}[htbp]
\centering
\pandocbounded{\safeincludesvg{./assets/figures/07-chuong-4-trien-khai-hardware-hinh-4-2a.svg}}
\caption{Sơ đồ chân kết nối SIM7600CE-T với ESP32-S3}
\label{fig:ch4-so-o-chan-ket-noi-sim7600ce-t-voi-esp32-s3}
\end{figure}
}



\paragraph{c) Khối đo điện áp ắc quy}\label{c-khux1ed1i-ux111o-ux111iux1ec7n-uxe1p-ux1eafc-quy}

Khối này dùng mạch chia áp để đưa điện áp ắc quy xe về dải đo an toàn cho
ADC của ESP32-S3. Firmware sẽ quy đổi giá trị ADC về điện áp thực tế và áp
dụng profile 12V hoặc 24V để nhận diện trạng thái nguồn, ngưỡng chuyển
nguồn, và điều kiện IGN ON/OFF.


{\thesissetfigureid{4.3}{fig-ch4-so-o-mach-o-ien-ap-ac-quy-bang-voltage-divider-va}
\begin{figure}[htbp]
\centering
\pandocbounded{\safeincludesvg{./assets/figures/07-chuong-4-trien-khai-hardware-hinh-4-3.svg}}
\caption{Sơ đồ mạch đo điện áp ắc quy bằng voltage divider và ADC ESP32-S3}
\label{fig:ch4-so-o-mach-o-ien-ap-ac-quy-bang-voltage-divider-va}
\end{figure}
}


\paragraph{d) Sơ đồ nguyên lý tổng hợp}\label{d-sux1a1-ux111ux1ed3-nguyuxean-luxfd-tux1ed5ng-hux1ee3p}

Sơ đồ nguyên lý tổng hợp của hệ thống thể hiện toàn bộ các khối mạch đã
được tích hợp trên PCB do nhóm tự thiết kế. Các kết nối nội bộ chính bao
gồm:

\begin{itemize}
\tightlist
\item
  \textbf{UART1} (2 dây tín hiệu TX/RX): Kết nối với modem SIM7600CE-T
  để truyền nhận cả dữ liệu 4G và GNSS (AT +
  \texttt{AT+CGNSINF}/\texttt{AT+CGNSTST}).
\item
  \textbf{I2C} (2 dây tín hiệu SDA/SCL): Kết nối với cảm biến gia tốc
  LIS3DSH
\item
  \textbf{BLE} (không dây): Kết nối với adapter OBD2 vgate iCar Pro
\item
  \textbf{ADC} (1 kênh): Đọc điện áp ắc quy qua voltage divider
\item
  \textbf{GPIO} (3 chân output): Điều khiển chân EN chuyển nguồn, chân EN sạc,
  Modem PWR-KEY
\item
  \textbf{GPIO} (2 chân input): Đọc trạng thái LVD, ngắt từ IMU
\end{itemize}


{\thesissetfigureid{4.4}{fig-ch4-so-o-nguyen-ly-mach-ien-tong-hop-cua-he-thong-tracker}
\begin{figure}[htbp]
\centering
\pandocbounded{\safeincludesvg{./assets/figures/07-chuong-4-trien-khai-hardware-hinh-4-4.svg}}
\caption{Sơ đồ nguyên lý mạch điện tổng hợp của hệ thống thiết bị theo dõi}
\label{fig:ch4-so-o-nguyen-ly-mach-ien-tong-hop-cua-he-thong-tracker}
\end{figure}
}


\begin{center}\rule{0.5\linewidth}{0.5pt}\end{center}

\paragraph{e) Sơ đồ đấu nối và GPIO mapping (Wiring \& GPIO
Pinout)}\label{e-sux1a1-ux111ux1ed3-ux111ux1ea5u-nux1ed1i-vuxe0-gpio-mapping-wiring--gpio-pinout}

\paragraph{a) Bảng phân công chân
GPIO}\label{a-bux1ea3ng-phuxe2n-cuxf4ng-chuxe2n-gpio}

Việc phân công chân GPIO của ESP32-S3 được thiết kế đảm bảo không xung
đột giữa các chức năng, đồng thời tối ưu hóa cho khả năng deep sleep và
external wakeup. \tabref{tab:ch4-phan-cong-chan-gpio-cua-esp32-s3} trình bày chi tiết phân công chân GPIO cho
toàn bộ hệ thống.


{\thesissettableid{4.1}{tab-ch4-phan-cong-chan-gpio-cua-esp32-s3}
\def\LTcaptype{table}
\begin{longtable}{|L{0.115\linewidth}|L{0.175\linewidth}|L{0.204\linewidth}|L{0.460\linewidth}|}
\caption{Phân công chân GPIO của ESP32-S3}\label{tab:ch4-phan-cong-chan-gpio-cua-esp32-s3}\\
\hline \textbf{Chân GPIO} & \textbf{Chức năng} & \textbf{Hướng} & \textbf{Mô tả chi tiết} \\
\\\hline
\endfirsthead
\caption[]{Phân công chân GPIO của ESP32-S3 (tiếp theo)}\\
\hline \textbf{Chân GPIO} & \textbf{Chức năng} & \textbf{Hướng} & \textbf{Mô tả chi tiết} \\
\\\hline
\endhead
\\\hline
\endfoot
\endlastfoot
GPIO 2 & IGN\_IN & Input & Đọc trạng thái khóa điện (IGN) từ xe hoặc qua
OBD2 \\\hline
GPIO 4 & \nolinkurl{U_BATT_ADC} & Input (ADC) & Đọc điện áp ắc quy xe qua voltage
divider (R1=100k, R2=10k) \\\hline
GPIO 5 & \nolinkurl{CHARGER_EN} & Output & Điều khiển bật/tắt khối sạc TP5100 (HIGH =
sạc, LOW = không sạc) \\\hline
GPIO 16 & \nolinkurl{MODEM_UART_TX} & Output & Truyền dữ liệu UART đến modem
SIM7600CE-T \\\hline
GPIO 17 & \nolinkurl{MODEM_UART_RX} & Input & Nhận dữ liệu UART từ modem
SIM7600CE-T \\\hline
GPIO 18 & \nolinkurl{POWER_PATH_EN} & Output & Điều khiển nhánh nguồn: LOW = ưu
tiên MP2482, HIGH = ưu tiên SX1308 dự phòng \\\hline
GPIO 19 & LVD\_STATUS & Input & Đọc trạng thái Low Voltage Disconnect
(HIGH = low-voltage, LOW = bình thường) \\\hline
GPIO 41 & \nolinkurl{LIS3DSH_INT1} & Input (Interrupt) & Nhận tín hiệu ngắt chính từ cảm biến gia tốc LIS3DSH khi phát hiện chuyển động \\\hline
GPIO 42 & \nolinkurl{LIS3DSH_INT2} & Input (Interrupt) & Dự phòng cho ngắt thứ hai của cảm biến LIS3DSH \\\hline
GPIO 2 & \nolinkurl{LIS3DSH_SDA} & I/O (I2C) & Đường dữ liệu I2C kết nối với cảm biến LIS3DSH \\\hline
GPIO 1 & \nolinkurl{LIS3DSH_SCL} & I/O (I2C) & Đường xung nhịp I2C kết nối với cảm biến LIS3DSH \\\hline
GPIO 26 & \nolinkurl{MODEM_PWR-KEY} & Output & Điều khiển bật/tắt nguồn modem
SIM7600CE-T \\\hline
\end{longtable}
}


\paragraph{b) Các lưu ý về phân công
chân}\label{b-cuxe1c-lux1b0u-uxfd-vux1ec1-phuxe2n-cuxf4ng-chuxe2n}

Việc phân công chân GPIO cần tuân thủ một số ràng buộc kỹ thuật của
ESP32-S3:

\begin{itemize}
\tightlist
\item
  \textbf{GPIO 34--39} chỉ hỗ trợ chế độ input, không có điện trở
  pull-up/pull-down nội. Do đó, các chân này không được sử dụng cho các
  tín hiệu output trong thiết kế này.
\item
  \textbf{Các chân ADC} nằm trong dải GPIO 0--15 và GPIO 25--27, hỗ trợ
  độ phân giải 12-bit. Chân GPIO 4 được chọn làm kênh ADC đo điện áp vì
  nằm trong vùng ADC1, cho phép đọc đồng thời với WiFi/BLE.
\item
  \textbf{Wakeup cho nhánh IMU} của hệ thống dùng \textit{GPIO wake} trong
  \textit{light sleep}. Với wiring production hiện tại, chân GPIO 41 (LIS3DSH\_INT1)
  không thuộc nhóm RTC GPIO nên không thể dùng EXT0/EXT1 cho deep sleep wake.
\end{itemize}

\paragraph{c) Sơ đồ đấu nối tổng
thể}\label{c-sux1a1-ux111ux1ed3-ux111ux1ea5u-nux1ed1i-tux1ed5ng-thux1ec3}

Sơ đồ đấu nối mô tả cách ánh xạ chân giữa ESP32-S3-WROOM-1 trên PCB
chính và các khối phần cứng tích hợp trên bo mạch. \figref{fig:ch4-so-o-au-noi-tong-the-giua-esp32-s3-va-cac-khoi} đã được
chuẩn hóa theo pin mapping firmware hiện tại, trong đó
\texttt{GPIO2/1} dành cho I2C IMU, \texttt{GPIO41/42} nhận ngắt IMU,
\texttt{GPIO26} điều khiển \texttt{PWR-KEY}, \texttt{GPIO19} nhận \texttt{LVD\_STATUS}, và
\texttt{GPIO4} đo \(U_{\text{batt}}\).

Bên cạnh các chân đã đưa vào \texttt{pin\_map.h}, netlist phần cứng còn
thể hiện đầy đủ cụm microSD \texttt{J4} theo chuẩn SDMMC 4-bit
(\texttt{SD-DAT0..3}, \texttt{SD-CMD}, \texttt{SD-CLK}, \texttt{SD-CD})
với pull-up ngoài trên các đường dữ liệu/lệnh. Điều này là nền tảng để
firmware triển khai \texttt{sd\_log\_store} và \texttt{offline\_queue}
cho cơ chế lưu đệm cục bộ.


{\thesissetfigureid{4.5}{fig-ch4-so-o-au-noi-tong-the-giua-esp32-s3-va-cac-khoi}
\begin{figure}[htbp]
\centering
\pandocbounded{\safeincludesvg{./assets/figures/07-chuong-4-trien-khai-hardware-hinh-4-5.svg}}
\caption{Sơ đồ đấu nối tổng thể giữa ESP32-S3 và các khối phần cứng tích hợp}
\label{fig:ch4-so-o-au-noi-tong-the-giua-esp32-s3-va-cac-khoi}
\end{figure}
}


\begin{center}\rule{0.5\linewidth}{0.5pt}\end{center}

\paragraph{f) Thiết kế mạch quản lý nguồn (Power Management
Circuit)}\label{f-thiux1ebft-kux1ebf-mux1ea1ch-quux1ea3n-luxfd-nguux1ed3n-power-management-circuit}

\paragraph{a) Tổng quan kiến trúc
nguồn}\label{a-tux1ed5ng-quan-kiux1ebfn-truxfac-nguux1ed3n}

Mạch quản lý nguồn là thành phần thiết yếu của hệ thống thiết bị theo dõi, bảo đảm
thiết bị hoạt động liên tục ngay cả khi ắc quy xe yếu hoặc mất điện.
Kiến trúc nguồn gồm sáu khối chức năng chính:

\begin{enumerate}
\def\labelenumi{\arabic{enumi}.}
\tightlist
\item
  \textbf{LDO 3.3V (AP2112-3.3):} Hạ từ bus 5V chính xuống 3.3V cấp
  ESP32-S3
\item
  \textbf{Buck 5V (MP2482):} Chuyển đổi 12V/24V xuống 5V bus chính
\item
  \textbf{Buck ~4V (TPS54231):} Chuyển đổi 12V/24V xuống
  ~4V cấp modem SIM7600CE-T
\item
  \textbf{Boost 5V (SX1308):} Tăng áp từ pin 18650 1S
  (~3.7V) lên 5V dự phòng
\item
  \textbf{Mạch chuyển nguồn Diode-OR + điều khiển EN:} Tự động duy trì nguồn
  liên tục giữa nhánh chính và nhánh dự phòng
\item
  \textbf{Mạch sạc TP5100 + LVD (khối LVD/ADC):} Sạc pin 1S và giám sát
  ngưỡng điện áp bảo vệ ắc quy
\end{enumerate}


{\thesissetfigureid{4.6}{fig-ch4-kien-truc-tong-the-mach-quan-ly-nguon}
\begin{figure}[htbp]
\centering
\pandocbounded{\safeincludesvg{./assets/figures/07-chuong-4-trien-khai-hardware-hinh-4-6.svg}}
\caption{Kiến trúc tổng thể mạch quản lý nguồn}
\label{fig:ch4-kien-truc-tong-the-mach-quan-ly-nguon}
\end{figure}
}


{\thesissetfigureid{4.6a}{fig-ch4-kien-truc-power-path-giua-nguon-xe-va-pin-du-phong}
\begin{figure}[htbp]
\centering
\pandocbounded{\safeincludesvg{./assets/figures/07-chuong-4-trien-khai-hardware-hinh-4-6a.svg}}
\caption{Kiến trúc chuyển nguồn giữa nguồn xe và pin dự phòng}
\label{fig:ch4-kien-truc-power-path-giua-nguon-xe-va-pin-du-phong}
\end{figure}
}



\paragraph{b) LDO 3.3V cho ESP32-S3 (AP2112-3.3)}\label{b-buck-33v-cho-esp32-s3-xl1509-33e}

Khối 3.3V dùng \textbf{AP2112-3.3} như một LDO để hạ từ bus 5V chính xuống
3.3V cấp cho ESP32-S3.


{\thesissettableid{4.2}{tab-ch4-khoi-buck-3-3v}
\def\LTcaptype{table}
\begin{longtable}{|L{0.168\linewidth}|L{0.132\linewidth}|L{0.218\linewidth}|L{0.250\linewidth}|L{0.176\linewidth}|}
\caption{Khối LDO 3.3V}\label{tab:ch4-khoi-buck-3-3v}\\
\hline \textbf{Khối} & \textbf{IC} & \textbf{Input} & \textbf{Output} & \textbf{Tải} \\
\\\hline
\endfirsthead
\caption[]{Khối LDO 3.3V (tiếp theo)}\\
\hline \textbf{Khối} & \textbf{IC} & \textbf{Input} & \textbf{Output} & \textbf{Tải} \\
\\\hline
\endhead
\\\hline
\endfoot
\endlastfoot
LDO 3.3V & AP2112-3.3 & 5V & 3.3V & ESP32-S3 \\\hline
\end{longtable}
}


\paragraph{c) Buck 5V bus chính
(MP2482)}\label{c-buck-5v-bus-chuxednh-mp2482}

Khối buck 5V dùng \textbf{MP2482} tạo bus 5V chính từ 12V/24V.


{\thesissettableid{4.3}{tab-ch4-khoi-buck-5v}
\def\LTcaptype{table}
\begin{longtable}{|L{0.168\linewidth}|L{0.132\linewidth}|L{0.218\linewidth}|L{0.250\linewidth}|L{0.176\linewidth}|}
\caption{Khối Buck 5V}\label{tab:ch4-khoi-buck-5v}\\
\hline \textbf{Khối} & \textbf{IC} & \textbf{Input} & \textbf{Output} & \textbf{Tải} \\
\\\hline
\endfirsthead
\caption[]{Khối Buck 5V (tiếp theo)}\\
\hline \textbf{Khối} & \textbf{IC} & \textbf{Input} & \textbf{Output} & \textbf{Tải} \\
\\\hline
\endhead
\\\hline
\endfoot
\endlastfoot
Buck 5V & MP2482 & 12--24V & 5V & Bus 5V chính \\\hline
\end{longtable}
}



{\thesissetfigureid{4.7}{fig-ch4-kien-truc-nhanh-buck-5v-chinh-dung-mp2482}
\begin{figure}[htbp]
\centering
\pandocbounded{\safeincludesvgnarrow{./assets/figures/07-chuong-4-trien-khai-hardware-hinh-4-7.svg}}
\caption{Kiến trúc nhánh buck 5V chính dùng MP2482}
\label{fig:ch4-kien-truc-nhanh-buck-5v-chinh-dung-mp2482}
\end{figure}
}


\paragraph{d) Buck 3.8V/4V cho modem
(TPS54231)}\label{d-buck-38v4v-cho-modem-tps54231}

Khối buck modem dùng \textbf{TPS54231} để hạ 12V/24V xuống khoảng 4V cấp
cho SIM7600CE-T.


{\thesissettableid{4.4}{tab-ch4-khoi-buck-3-8v-4v-cho-modem}
\def\LTcaptype{table}
\begin{longtable}{|L{0.168\linewidth}|L{0.132\linewidth}|L{0.218\linewidth}|L{0.250\linewidth}|L{0.176\linewidth}|}
\caption{Khối Buck 3.8V/4V cho modem}\label{tab:ch4-khoi-buck-3-8v-4v-cho-modem}\\
\hline \textbf{Khối} & \textbf{IC} & \textbf{Input} & \textbf{Output} & \textbf{Tải} \\
\\\hline
\endfirsthead
\caption[]{Khối Buck 3.8V/4V cho modem (tiếp theo)}\\
\hline \textbf{Khối} & \textbf{IC} & \textbf{Input} & \textbf{Output} & \textbf{Tải} \\
\\\hline
\endhead
\\\hline
\endfoot
\endlastfoot
Buck 3.8V & TPS54\-231 & 12--24V & ~4V & SIM7600CE-T \\\hline
\end{longtable}
}


\paragraph{e) Boost 5V dự phòng và Power
Path}\label{e-boost-5v-dux1ef1-phuxf2ng-vuxe0-power-path}

Khối boost dùng \textbf{SX1308} để nâng áp từ pin 18650 1S
(~3.7V) lên 5V khi chạy nguồn dự phòng.


{\thesissettableid{4.5}{tab-ch4-khoi-boost-5v-du-phong}
\def\LTcaptype{table}
\begin{longtable}{|L{0.189\linewidth}|L{0.094\linewidth}|L{0.236\linewidth}|L{0.283\linewidth}|L{0.142\linewidth}|}
\caption{Khối Boost 5V dự phòng}\label{tab:ch4-khoi-boost-5v-du-phong}\\
\hline \textbf{Khối} & \textbf{IC} & \textbf{Input} & \textbf{Output} & \textbf{Tải} \\
\\\hline
\endfirsthead
\caption[]{Khối Boost 5V dự phòng (tiếp theo)}\\
\hline \textbf{Khối} & \textbf{IC} & \textbf{Input} & \textbf{Output} & \textbf{Tải} \\
\\\hline
\endhead
\\\hline
\endfoot
\endlastfoot
Boost 5V & SX1308 & Battery ~3.7V & 5V & Nguồn dự phòng từ pin \\\hline
\end{longtable}
}



{\thesissetfigureid{4.8}{fig-ch4-kien-truc-nhanh-nguon-du-phong-5v-tu-pin-18650-qua-sx1308}
\begin{figure}[htbp]
\centering
\pandocbounded{\safeincludesvgnarrow{./assets/figures/07-chuong-4-trien-khai-hardware-hinh-4-8.svg}}
\caption{Kiến trúc nhánh nguồn dự phòng 5V từ pin 18650 qua SX1308}
\label{fig:ch4-kien-truc-nhanh-nguon-du-phong-5v-tu-pin-18650-qua-sx1308}
\end{figure}
}


Mạch chuyển nguồn khi vận hành được triển khai theo \textbf{diode OR} giữa nhánh 5V
chính (MP2482) và nhánh 5V dự phòng (SX1308), phối hợp điều khiển GPIO để
đảm bảo chuyển nguồn liên tục.


{\thesissettableid{4.6}{tab-ch4-logic-chuyen-nguon-tu-ong}
\def\LTcaptype{table}
\begin{longtable}{|L{0.123\linewidth}|L{0.068\linewidth}|L{0.380\linewidth}|L{0.160\linewidth}|L{0.100\linewidth}|L{0.104\linewidth}|}
\caption{Logic chuyển nguồn tự động}\label{tab:ch4-logic-chuyen-nguon-tu-ong}\\
\hline \textbf{Chế độ} & \textbf{IGN} & \textbf{\(U_{\text{batt}}\)} & Nguồn Tracker & Sạc Pin & Cảnh báo \\
\\\hline
\endfirsthead
\caption[]{Logic chuyển nguồn tự động (tiếp theo)}\\
\hline \textbf{Chế độ} & \textbf{IGN} & \textbf{\(U_{\text{batt}}\)} & Nguồn Tracker & Sạc Pin & Cảnh báo \\
\\\hline
\endhead
\\\hline
\endfoot
\endlastfoot
Xe chạy & ON & Profile 12V: \(U_{\text{batt}} \ge 13.0\,\mathrm{V}\);\linebreak
Profile 24V: \(U_{\text{batt}} \ge 26.0\,\mathrm{V}\) & Ắc quy & Có & Không \\\hline
Xe đỗ & OFF & Profile 12V: \(U_{\text{batt}} > 12.0\,\mathrm{V}\);\linebreak
Profile 24V: \(U_{\text{batt}} > 24.0\,\mathrm{V}\) & Ắc quy & Không & Không \\\hline
Ắc quy yếu & OFF & Profile 12V: \(U_{\text{batt}} \le 12.0\,\mathrm{V}\);\linebreak
Profile 24V: \(U_{\text{batt}} \le 24.0\,\mathrm{V}\) & Pin 18650 1S & Không & Có \\\hline
Ắc quy phục hồi & OFF & Profile 12V: \(U_{\text{batt}} \ge 12.2\,\mathrm{V}\);\linebreak
Profile 24V: \(U_{\text{batt}} \ge 24.4\,\mathrm{V}\) & Ắc quy & Không & Có \\\hline
\end{longtable}
}



{\thesissetfigureid{4.9}{fig-ch4-so-o-power-path-runtime-giua-nhanh-chinh-va-nhanh-backup}
\begin{figure}[htbp]
\centering
\pandocbounded{\safeincludesvg{./assets/figures/07-chuong-4-trien-khai-hardware-hinh-4-9.svg}}
\caption{Sơ đồ chuyển nguồn giữa nhánh chính và nhánh dự phòng}
\label{fig:ch4-so-o-power-path-runtime-giua-nhanh-chinh-va-nhanh-backup}
\end{figure}
}

\clearpage


\paragraph{f) Mạch sạc pin
TP5100}\label{f-mux1ea1ch-sux1ea1c-pin-tp4056}

Khối sạc pin dùng \textbf{TP5100} ở cấu hình \textbf{1 cell}, nhận
\textbf{5V từ MP2482} và sạc pin 18650 1S ở mức 4.2V.


{\thesissettableid{4.7}{tab-ch4-khoi-sac-tp4056}
\def\LTcaptype{table}
\begin{longtable}{|L{0.189\linewidth}|L{0.094\linewidth}|L{0.236\linewidth}|L{0.283\linewidth}|L{0.142\linewidth}|}
\caption{Khối sạc TP5100}\label{tab:ch4-khoi-sac-tp4056}\\
\hline \textbf{Khối} & \textbf{IC} & \textbf{Input} & \textbf{Output} & \textbf{Tải} \\
\\\hline
\endfirsthead
\caption[]{Khối sạc TP5100 (tiếp theo)}\\
\hline \textbf{Khối} & \textbf{IC} & \textbf{Input} & \textbf{Output} & \textbf{Tải} \\
\\\hline
\endhead
\\\hline
\endfoot
\endlastfoot
Sạc pin & TP5100 & 5V từ MP2482 & 4.2V & Pin 18650 1S \\\hline
\end{longtable}
}



{\thesissetfigureid{4.10}{fig-ch4-kien-truc-sac-pin-tp4056-va-uong-cap-nguon-du-phong}
\begin{figure}[htbp]
\centering
\pandocbounded{\safeincludesvg{./assets/figures/07-chuong-4-trien-khai-hardware-hinh-4-10.svg}}
\caption{Kiến trúc sạc pin TP5100 và đường cấp nguồn dự phòng}
\label{fig:ch4-kien-truc-sac-pin-tp4056-va-uong-cap-nguon-du-phong}
\end{figure}
}


Dòng sạc của TP5100 được cấu hình bằng điện trở thiết lập dòng sạc và
còn phụ thuộc điều kiện nhiệt của mạch, không cố định một giá trị trong
mọi điều kiện.

\paragraph{g) Mạch giám sát điện áp Low Voltage Disconnect
(LVD)}\label{g-mux1ea1ch-giuxe1m-suxe1t-ux111iux1ec7n-uxe1p-low-voltage-disconnect-lvd}

Hệ thống dùng ADC (R1 = 100 kΩ, R2 = 10 kΩ) kết hợp khối LVD phần cứng để
giám sát điện áp.

Ngưỡng profile:

\begin{itemize}
\tightlist
\item
  Profile 12V: \(OFF = 12.0\,\mathrm{V}\), \(ON = 12.2\,\mathrm{V}\)
\item
  Profile 24V: \(OFF = 24.0\,\mathrm{V}\), \(ON = 24.4\,\mathrm{V}\)
\end{itemize}

Tín hiệu LVD\_STATUS đọc tại GPIO19 với quy ước vận hành: \textbf{HIGH =
low-voltage}, \textbf{LOW = bình thường}.

\begin{center}\rule{0.5\linewidth}{0.5pt}\end{center}

\paragraph{g) Thiết kế vỏ hộp bảo vệ (Enclosure
Design)}\label{g-thiux1ebft-kux1ebf-vux1ecf-hux1ed9p-bux1ea3o-vux1ec7-enclosure-design}

\paragraph{a) Yêu cầu thiết
kế}\label{a-yuxeau-cux1ea7u-thiux1ebft-kux1ebf}

Vỏ hộp bảo vệ cần đáp ứng các yêu cầu sau:

\begin{itemize}
\tightlist
\item
  \textbf{Kích thước nhỏ gọn:} Phù hợp lắp đặt dưới táp-lô hoặc gần cổng
  OBD2, kích thước tối đa khoảng
  \(100 \times 70 \times 35\,\mathrm{mm}\)
\item
  \textbf{Tản nhiệt:} Đảm bảo thông gió cho các IC công suất (MP2482,
  SX1308) với tổn hao nhiệt tổng cộng có thể lên đến 3--4W
\item
  \textbf{Bảo vệ:} Chống bụi, chống nước cơ bản (IP54 hoặc tương đương)
  cho môi trường bên trong xe
\item
  \textbf{Tiếp cận anten:} Bố trí vị trí phù hợp cho anten 4G/LTE, anten
  GNSS và anten BLE để đảm bảo chất lượng thu phát tín hiệu
\end{itemize}

\paragraph{b) Bố cục bên trong}\label{b-bux1ed1-cux1ee5c-buxean-trong}

Bố cục bên trong vỏ hộp được thiết kế theo nguyên tắc phân vùng chức
năng:

Bố cục bên trong được chia thành ba vùng rõ ràng: cụm xử lý và truyền
thông (\texttt{ESP32-S3}, \texttt{SIM7600CE-T}), cụm nguồn
(\texttt{AP2112-3.3}, \texttt{MP2482}, \texttt{TPS54231},
\texttt{SX1308}, \texttt{TP5100}, \texttt{khối\ LVD}) và cụm lưu trữ năng
lượng (\texttt{pin\ 18650\ 1S\ +\ mạch\ bảo\ vệ\ 1S}). \figref{fig:ch4-so-o-bo-cuc-ben-trong-vo-hop-bao-ve} thể hiện cách nhóm
các khối phần cứng này trong vỏ hộp cùng vị trí anten và cổng OBD2.


{\thesissetfigureid{4.11}{fig-ch4-so-o-bo-cuc-ben-trong-vo-hop-bao-ve}
\begin{figure}[htbp]
\centering
\pandocbounded{\safeincludesvg{./assets/figures/07-chuong-4-trien-khai-hardware-hinh-4-11.svg}}
\caption{Sơ đồ bố cục bên trong vỏ hộp bảo vệ}
\label{fig:ch4-so-o-bo-cuc-ben-trong-vo-hop-bao-ve}
\end{figure}
}


\paragraph{c) Các lưu ý thiết
kế}\label{c-cuxe1c-lux1b0u-uxfd-thiux1ebft-kux1ebf}

\textbf{Tản nhiệt:} Các khối công suất (MP2482, SX1308) được bố trí tại
vị trí thông thoáng, tách khỏi cảm biến nhiệt độ; khe thông gió đặt ở
mặt bên và mặt dưới vỏ hộp. IC MP2482 được gắn heatsink khi vận hành
công suất cao.

\textbf{Anten:} Anten 4G/LTE và GNSS đặt ở mặt trên vỏ hộp, hướng lên
trên để tối ưu thu sóng. Anten BLE (tích hợp trong ESP32-S3) hướng về
cổng OBD2 để rút ngắn khoảng cách kết nối với adapter vgate iCar Pro.
Các anten cách nhau tối thiểu 30 mm nhằm giảm nhiễu tương hỗ.

\textbf{Chống nhiễu điện từ (EMI):} Mạch công suất (Buck, Boost) được bố
trí tách xa mạch tín hiệu (UART, I2C, ADC); dây nguồn và dây tín hiệu đi
riêng, không chạy song song.

\textbf{Kết nối:} Cổng kết nối OBD2, cổng SIM card và cổng USB (cho nạp
firmware/debug) được bố trí ở mặt bên vỏ hộp, dễ tiếp cận khi lắp đặt và
bảo trì.

\begin{center}\rule{0.5\linewidth}{0.5pt}\end{center}

\subsubsection{4.1.2. Triển khai firmware, hạ tầng máy chủ và giao diện điều
khiển}\label{412-triux1ec3n-khai-firmware-cloud-vuxe0-giao-diux1ec7n-ux111iux1ec1u-khiux1ec3n}

Ngoài các mô-đun truyền thông và điều phối trạng thái, firmware đã triển
khai luồng \textbf{SD card log} theo chu trình vận hành rõ ràng:

\begin{enumerate}
\def\labelenumi{\arabic{enumi}.}
\tightlist
\item
  \textbf{Khởi tạo và mount thẻ SD}: \texttt{sd\_log\_store\_init()} và
  \texttt{sd\_log\_store\_mount()} khởi tạo host SDMMC 4-bit, mount
  FATFS vào \texttt{/sdcard}, tạo cây thư mục làm việc và nạp metadata
  phiên ghi.
\item
  \textbf{Ghi bản ghi cục bộ}: \texttt{offline\_queue\_enqueue()} chuẩn
  hóa bản ghi dữ liệu vận hành/sự kiện, sau đó
  \texttt{sd\_log\_store\_append()} ghi xuống hàng đợi file theo thứ tự
  \texttt{seq}.
\item
  \textbf{Phát lại khi mạng phục hồi}:
  \texttt{offline\_queue\_replay\_tick()} đọc bản ghi chờ, gửi lại theo
  topic phù hợp; bản ghi quan trọng đi QoS 1 và được chốt bằng
  \texttt{offline\_queue\_handle\_publish\_ack()}.
\item
  \textbf{Kiểm soát dung lượng}:
  \texttt{sd\_log\_store\_gc\_if\_needed()} dọn dữ liệu ít quan trọng hơn đã
  ACK khi vượt ngưỡng quota, giữ ổn định không gian lưu trữ.
\item
  \textbf{Kết thúc phiên ghi an toàn}: firmware gọi
  \texttt{sd\_log\_store\_stop\_session()} và
  \texttt{sd\_log\_store\_unmount()} khi dừng phiên để giảm nguy cơ lỗi
  hệ thống tệp.
\end{enumerate}

Trong phạm vi đồ án, luồng này được đánh giá theo giả định thẻ đã được
gắn sẵn ngay từ lúc khởi động để ưu tiên độ ổn định. Kịch bản cắm hoặc
rút thẻ trong khi thiết bị đang chạy chưa phải mục tiêu chính của phiên
bản hiện tại.

\begin{center}\rule{0.5\linewidth}{0.5pt}\end{center}

\subsection{4.2. Chế tạo và lắp ráp hệ thống -- Manufacture and
Assembly}\label{42-chux1ebf-tux1ea1o-vuxe0-lux1eafp-ruxe1p-hux1ec7-thux1ed1ng--manufacture-and-assembly}

\subsubsection{4.2.1. Lắp ráp mạch điện tử (Electronics
Assembly)}\label{421-lux1eafp-ruxe1p-mux1ea1ch-ux111iux1ec7n-tux1eed-electronics-assembly}

Phần lắp ráp không chỉ nhằm ghép các linh kiện thành một bo mạch hoàn
chỉnh, mà còn là bước kiểm chứng xem các giả định thiết kế ở Chương 3 có
đứng vững khi đi vào hiện thực hay không. Vì vậy, phần này được trình
bày theo đúng thứ tự mà một người triển khai thực tế sẽ quan tâm: cần có
những vật tư gì, lắp theo thứ tự nào, và kiểm tra ra sao để hạn chế lỗi
ngay từ đầu.

\paragraph{a) Danh sách vật liệu (Bill of
Materials)}\label{a-danh-suxe1ch-vux1eadt-liux1ec7u-bill-of-materials}


{\thesissettableid{4.6}{tab-ch4-danh-sach-vat-lieu-ay-u-cho-he-thong-tracker}
\def\LTcaptype{table}
\begin{longtable}{|L{0.068\linewidth}|L{0.259\linewidth}|L{0.068\linewidth}|L{0.068\linewidth}|L{0.146\linewidth}|L{0.326\linewidth}|}
\caption{Danh sách vật liệu đầy đủ cho hệ thống thiết bị theo dõi}\label{tab:ch4-danh-sach-vat-lieu-ay-u-cho-he-thong-tracker}\\
\hline \textbf{STT} & \textbf{Thành phần} & \textbf{Đơn vị} & \textbf{SL} & \textbf{Giá tham chiếu T4/2026 (VND)} & \textbf{Ghi chú} \\
\\\hline
\endfirsthead
\caption[]{Danh sách vật liệu đầy đủ cho hệ thống thiết bị theo dõi (tiếp theo)}\\
\hline \textbf{STT} & \textbf{Thành phần} & \textbf{Đơn vị} & \textbf{SL} & \textbf{Giá tham chiếu T4/2026 (VND)} & \textbf{Ghi chú} \\
\\\hline
\endhead
\\\hline
\endfoot
\endlastfoot
1 & ESP32-S3-WROOM-1 (N16R8) & Cái & 1 & 95.000 & MCU trung
tâm của bo mạch \\\hline
2 & IC cảm biến LIS3DSH & Cái & 1 & 40.000 & Cảm biến gia tốc 3
trục, giao tiếp I2C \\\hline
3 & vgate iCar Pro (OBD2 BLE) & Cái & 1 & 367.000 & Adapter
OBD2 BLE 4.0, tương thích ESP32-S3 \\\hline
4 & SIMCom SIM7600CE-T & Bộ & 1 & 650.000 & Modem LTE Cat-4
tích hợp GNSS \\\hline
5 & Pin 18650 1S Li-ion 3500mAh & Cái & 1 & 79.000 & Loại có
protection board \\\hline
6 & IC sạc TP5100 + mạch phụ trợ & Cái & 1 & 25.000 & Input 5V
từ MP2482, output 4.2V, dòng theo điện trở thiết lập \\\hline
7 & Mạch LDO AP2112-3.3 & Cái & 1 & 5.000 & 5V
-\textgreater{} 3.3V cho ESP32-S3 \\\hline
8 & Mạch buck MP2482 5V & Cái & 1 & 17.000 & Bus 5V chính \\\hline
9 & Mạch buck TPS54231 ~4V & Cái & 1 & 9.000 &
12--24V -\textgreater{} ~4V cấp modem SIM7600CE-T \\\hline
10 & Mạch boost SX1308 5V dự phòng & Cái & 1 & 12.000 & Nguồn dự phòng từ
pin 18650 1S \\\hline
11 & Khối giám sát LVD + ADC & Cái & 1 & 8.000 & Giám sát LVD
(GPIO19) \\\hline
12 & Diode Schottky (SS54/SS34) & Cái & 2 & 2.500 & Diode OR dự
phòng \\\hline
13 & Mạch bảo vệ pin 1S (BMS) & Cái & 1 & 12.000 & Bảo vệ pin
18650 1S \\\hline
14 & Điện trở (10k, 2.2k, v.v.) & Gói & 1 & 5.000 & Voltage
divider, pull-up/pull-down \\\hline
15 & Tụ điện (100uF, 220uF, v.v.) & Gói & 1 & 5.000 & Lọc nhiễu,
decoupling \\\hline
16 & Connector, Header Pin & Gói & 1 & 55.000 & Header, jack OBD2 và
đầu nối nguồn \\\hline
17 & PCB (nếu tự thiết kế) & Cái & 1 & 20.000 & PCB 2 lớp, kích
thước ~50x50 mm \\\hline
18 & Vỏ bảo vệ & Cái & 1 & 25.000 & Vỏ nhựa hoặc kim loại \\\hline
19 & Dây nối, cáp điện & Mét & --- & 60.000 & Dây điện, anten
4G/GNSS, khay SIM \\\hline
20 & Linh kiện phụ trợ khác & --- & --- & 20.000 & Cầu chì, công
tắc, LED chỉ thị \\\hline
\end{longtable}
}



{\thesissettableid{4.7}{tab-ch4-tong-chi-phi-uoc-tinh}
\def\LTcaptype{table}
\begin{longtable}{|L{0.526\linewidth}|L{0.442\linewidth}|}
\caption{Tổng chi phí ước tính}\label{tab:ch4-tong-chi-phi-uoc-tinh}\\
\hline \textbf{Hạng mục} & \textbf{Chi phí (VND)} \\
\\\hline
\endfirsthead
\caption[]{Tổng chi phí ước tính (tiếp theo)}\\
\hline \textbf{Hạng mục} & \textbf{Chi phí (VND)} \\
\\\hline
\endhead
\\\hline
\endfoot
\endlastfoot
Thành phần chính (MCU, cảm biến, modem, OBD2, pin) & 1.231.000 \\\hline
Hệ thống quản lý nguồn & 93.000 \\\hline
Kết nối, RF và linh kiện phụ trợ & 145.000 \\\hline
PCB và vỏ bảo vệ & 45.000 \\\hline
\textbf{Tổng cộng} & \textbf{1.514.000} \\\hline
\end{longtable}
}


\paragraph{b) Quy trình lắp ráp}\label{b-quy-truxecnh-lux1eafp-ruxe1p}

Để giảm rủi ro lỗi dây chuyền, quy trình lắp ráp không đi theo hướng hàn
toàn bộ rồi mới kiểm tra, mà chia thành từng cụm chức năng và kiểm tra
ngay sau mỗi bước. Trình tự thực hiện như sau:

\textbf{Bước 1 - Kiểm tra linh kiện:} Kiểm tra toàn bộ IC, linh kiện và
bán thành phẩm trước khi lắp ráp. Test riêng từng khối chức năng chính
(ESP32-S3-WROOM-1, AP2112-3.3, MP2482, TPS54231, SX1308, TP5100,
SIM7600CE-T, khối LVD) để bảo đảm hoạt động đúng.

\textbf{Bước 2 - Lắp ráp khối nguồn trên PCB:} Hàn và kiểm tra MP2482 để
tạo bus 5V từ nguồn ắc quy xe (12V hoặc 24V). Lắp AP2112-3.3 cấp riêng
cho ESP32-S3, TPS54231 tạo rail ~4V cho modem SIM7600CE-T,
SX1308 tạo 5V dự phòng từ pin 18650 1S, diode OR giữa nhánh 5V chính và
nhánh 5V dự phòng, cùng TP5100 kết nối với pin và mạch bảo vệ 1S.

\textbf{Bước 3 - Lắp ráp khối xử lý và điều khiển:} Hàn
ESP32-S3-WROOM-1, LIS3DSH và các linh kiện liên quan lên PCB
chính. Kiểm tra các chân GPIO theo bảng phân công (\tabref{tab:ch4-phan-cong-chan-gpio-cua-esp32-s3}) và cấp
nguồn theo từng rail chức năng (3.3V logic, ~4V modem, 5V
bus/backup).

\textbf{Bước 4 - Tích hợp các thành phần còn lại:} Lắp modem SIM7600CE-T
lên bo mạch, kiểm tra UART1 (GPIO16, GPIO17) và chân \texttt{PWR-KEY}
(GPIO26), cấu hình lấy dữ liệu GNSS qua
\texttt{AT+CGNSINF}/\texttt{AT+CGNSTST} trên cùng UART, kiểm tra cảm
biến LIS3DSH trên I2C (GPIO2, GPIO1), và hiệu chuẩn mạch đo điện áp ắc
quy tại GPIO4 (ADC).

\textbf{Bước 5 - Kiểm tra tích hợp:} Nạp firmware cơ bản để kiểm tra
từng chức năng: đọc ADC, điều khiển GPIO, giao tiếp UART với modem, quét
BLE, đọc I2C từ LIS3DSH. Kiểm tra chuyển nguồn tự động bằng cách thay
đổi điện áp đầu vào.

\textbf{Bước 6 - Lắp ráp vào vỏ:} Định vị bo mạch chính, cụm anten và
pin dự phòng bên trong vỏ hộp theo bố cục đã thiết kế. Cố định bằng ốc
vít hoặc keo nhiệt. Kết nối anten 4G/LTE và GNSS. Đóng nắp vỏ hộp và
kiểm tra tổng thể.


{\thesissetfigureid{4.12}{fig-ch4-quy-trinh-lap-rap-phan-cung-tren-pcb}
\begin{figure}[htbp]
\centering
\pandocbounded{\safeincludesvg{./assets/figures/07-chuong-4-trien-khai-hardware-hinh-4-12.svg}}
\caption{Quy trình lắp ráp phần cứng trên PCB}
\label{fig:ch4-quy-trinh-lap-rap-phan-cung-tren-pcb}
\end{figure}
}



{\thesissetfigureid{4.13}{fig-ch4-so-o-bo-tri-linh-kien-sau-khi-lap-rap}
\begin{figure}[htbp]
\centering
\pandocbounded{\safeincludesvg{./assets/figures/07-chuong-4-trien-khai-hardware-hinh-4-13.svg}}
\caption{Sơ đồ bố trí linh kiện sau khi lắp ráp}
\label{fig:ch4-so-o-bo-tri-linh-kien-sau-khi-lap-rap}
\end{figure}
}


\paragraph{c) Kiểm tra và hiệu
chuẩn}\label{c-kiux1ec3m-tra-vuxe0-hiux1ec7u-chuux1ea9n}

Sau khi hoàn tất lắp ráp cơ khí và điện, hệ thống cần được hiệu chuẩn để
đảm bảo các giá trị đo và cơ chế chuyển nguồn phản ánh đúng điều kiện
vận hành thực tế của xe. Các hạng mục kiểm tra chính gồm:

\begin{itemize}
\tightlist
\item
  \textbf{Hiệu chuẩn ADC:} So sánh giá trị điện áp đọc từ ADC với giá
  trị đo từ đồng hồ vạn năng (multimeter). Điều chỉnh hệ số hiệu chỉnh
  trong firmware nếu cần.
\item
  \textbf{Kiểm tra chuyển nguồn:} Mô phỏng tình huống ắc quy yếu theo
  profile cấu hình (ví dụ hệ 12V giảm từ 12V xuống dưới 12V), xác nhận
  hệ thống tự động chuyển sang pin dự phòng.
\item
  \textbf{Kiểm tra sạc pin:} Xác nhận khối sạc TP5100 hoạt động đúng: sạc
  khi \(IGN_{\mathrm{ON}}\), ngừng sạc khi \(IGN_{\mathrm{OFF}}\) hoặc
  \(U_{\text{batt}}\) thấp.
\item
  \textbf{Kiểm tra giao tiếp:} Xác nhận modem phản hồi lệnh AT qua UART,
  cảm biến LIS3DSH trả về dữ liệu qua I2C, kết nối BLE với OBD2 adapter
  thành công.
\end{itemize}

\begin{center}\rule{0.5\linewidth}{0.5pt}\end{center}

\subsubsection{4.2.2. Lắp đặt trong xe (Vehicle
Installation)}\label{422-lux1eafp-ux111ux1eb7t-trong-xe-vehicle-installation}

Nếu bước lắp ráp trả lời câu hỏi ``thiết bị có hoạt động được không'',
thì bước lắp đặt trong xe trả lời câu hỏi ``thiết bị có hoạt động ổn
định trong môi trường thật hay không''. Phần này vì thế tập trung vào
những quyết định ảnh hưởng trực tiếp đến chất lượng thu sóng, độ an toàn
điện và khả năng bảo trì sau này.

\paragraph{a) Vị trí lắp
đặt}\label{a-vux1ecb-truxed-lux1eafp-ux111ux1eb7t}

Thiết bị theo dõi có thể được lắp ở nhiều vị trí khác nhau trong xe, tuy
nhiên việc chọn vị trí không chỉ phụ thuộc vào chỗ trống. Người triển
khai cần đồng thời cân nhắc độ kín đáo, độ thuận tiện khi đi dây, khoảng
cách tới cổng OBD2 và điều kiện thu sóng cho GNSS/4G.

\begin{itemize}
\tightlist
\item
  \textbf{Dưới bảng táp-lô}: Vị trí khuyến nghị, gần cổng OBD2
  (thường nằm dưới vô-lăng bên trái), dễ dàng kết nối nguồn và OBD2.
  Không gian đủ cho vỏ hộp và không can thiệp đến vùng hoạt động của tài
  xế.
\item
  \textbf{Dưới ghế lái hoặc ghế phụ:} Vị trí thay thế khi không gian
  dưới táp-lô không đủ. Cần kéo dài cáp kết nối OBD2.
\item
  \textbf{Trong hộp cầu chì (fuse box):} Vị trí kín đáo, phù hợp cho ứng
  dụng theo dõi không cần người dùng can thiệp.
\end{itemize}

\paragraph{b) Kết nối với cổng
OBD2}\label{b-kux1ebft-nux1ed1i-vux1edbi-cux1ed5ng-obd2}

Cổng OBD2 (16 chân) cung cấp cả nguồn điện và giao tiếp chẩn đoán. Trong
cấu hình hiện tại, thiết bị theo dõi chỉ lấy \texttt{Battery\ Power} ở chân
\texttt{16} và mass tại chân \texttt{4/5}, còn dữ liệu OBD2 đi qua
adapter \texttt{vgate\ iCar\ Pro} cắm trực tiếp vào cổng xe rồi truyền
không dây về ESP32-S3 qua BLE.

\textbf{Lưu ý:} Adapter vgate iCar Pro được cắm trực tiếp vào cổng OBD2
của xe. ESP32-S3 giao tiếp với adapter này qua BLE, không cần kết nối
dây vật lý cho phần dữ liệu OBD2. Chỉ có nguồn điện từ ắc quy xe (12V
hoặc 24V) và GND được lấy từ cổng OBD2 thông qua dây nối riêng.


{\thesissetfigureid{4.14}{fig-ch4-minh-hoa-vi-tri-cong-obd2-tren-xe-va-cach-ket-noi}
\begin{figure}[htbp]
\centering
\pandocbounded{\safeincludesvg{./assets/figures/07-chuong-4-trien-khai-hardware-hinh-4-14.svg}}
\caption{Minh họa vị trí cổng OBD2 trên xe và cách kết nối}
\label{fig:ch4-minh-hoa-vi-tri-cong-obd2-tren-xe-va-cach-ket-noi}
\end{figure}
}


\paragraph{c) Đi dây và kết nối
anten}\label{c-ux111i-duxe2y-vuxe0-kux1ebft-nux1ed1i-anten}

\textbf{Dây nguồn:} Dây nguồn từ cổng OBD2 đến thiết bị theo dõi (12V
hoặc 24V tùy hệ xe) sử dụng dây điện tiết diện tối thiểu 0.5 mm2 (AWG
22), có bọc cách điện tốt. Cầu chì (fuse) 5A được đặt tại đầu dây nguồn
gần cổng OBD2 để bảo vệ mạch.

\textbf{Anten 4G/LTE và GNSS:} Anten được đặt hướng lên mặt trên táp-lô
hoặc gần kính chắn gió phía trước để đảm bảo thu sóng tốt nhất. Dây
anten (cáp đồng trục) được đi dọc theo khung xe, tránh chạy song song
với dây điện nguồn.

\textbf{Anten BLE:} Do BLE sử dụng anten PCB tích hợp trên ESP32-S3, cần
đảm bảo khoảng cách từ thiết bị theo dõi đến adapter vgate iCar Pro (cắm
ở cổng OBD2) không quá 3 mét và không bị che chắn bởi vật liệu kim loại
dày.

\paragraph{d) Quản lý cáp và chống
nhiá»…u}\label{d-quux1ea3n-luxfd-cuxe1p-vuxe0-chux1ed1ng-nhiux1ec5u}

\begin{itemize}
\tightlist
\item
  \textbf{Cố định cáp:} Tất cả các dây nối được cố định bằng dây rút
  nhựa (cable tie) và kẹp cáp, tránh các dây bị rung lắc gây tiếng ồn
  hoặc hư hỏng tiếp xúc.
\item
  \textbf{Chống nhiễu EMI:} Dây tín hiệu UART (từ ESP32 đến modem) được
  sử dụng loại cáp xoắn đôi (twisted pair) hoặc cáp có bọc chống nhiễu.
  Dây nguồn và dây tín hiệu được đi riêng, không chạy song song để giảm
  nhiễu điện từ từ hệ thống điện của xe.
\item
  \textbf{Chống ẩm:} Các mối nối điện được bọc bằng keo nhiệt hoặc ống
  co nhiệt (heat shrink tube) để chống ẩm và chống oxy hóa trong môi
  trường xe.
\end{itemize}


{\thesissetfigureid{4.15}{fig-ch4-minh-hoa-lap-at-thiet-bi-tracker-trong-xe-va-i-day}
\begin{figure}[htbp]
\centering
\pandocbounded{\safeincludesvg{./assets/figures/07-chuong-4-trien-khai-hardware-hinh-4-15.svg}}
\caption{Minh họa lắp đặt thiết bị theo dõi trong xe và đi dây}
\label{fig:ch4-minh-hoa-lap-at-thiet-bi-tracker-trong-xe-va-i-day}
\end{figure}
}


\paragraph{e) Kiểm tra sau lắp
đặt}\label{e-kiux1ec3m-tra-sau-lux1eafp-ux111ux1eb7t}

Sau khi lắp đặt xong, cần kiểm tra lại toàn bộ theo góc nhìn vận hành
thực tế, tức là không chỉ xem thiết bị có lên nguồn hay không mà còn
đánh giá được nó có giao tiếp tốt, chuyển chế độ đúng và duy trì kết nối
ổn định trong xe hay không.

\begin{enumerate}
\def\labelenumi{\arabic{enumi}.}
\tightlist
\item
  \textbf{Kiểm tra nguồn điện:} Đo điện áp tại đầu vào thiết bị, xác
  nhận điện áp ắc quy xe (12V hoặc 24V) được cấp đầy đủ
\item
  \textbf{Kiểm tra kết nối BLE:} Xác nhận ESP32-S3 kết nối được với
  adapter vgate iCar Pro và đọc dữ liệu OBD2 (IGN, RPM, tốc độ)
\item
  \textbf{Kiểm tra GNSS:} Xác nhận modem SIM7600CE-T bắt được vệ tinh và
  trả về tọa độ GPS chính xác (sai số \textless{} 5 mét)
\item
  \textbf{Kiểm tra 4G/LTE:} Xác nhận modem đăng ký mạng thành công, gửi
  được dữ liệu lên máy chủ qua MQTT
\item
  \textbf{Kiểm tra chuyển nguồn:} Tắt máy xe (IGN OFF), xác nhận thiết
  bị chuyển sang chế độ tiết kiệm năng lượng và sử dụng pin dự phòng khi
  cần
\item
  \textbf{Kiểm tra deep sleep:} Xác nhận ESP32-S3 vào chế độ deep sleep
  khi xe đỗ, và đánh thức đúng khi phát hiện rung động (qua LIS3DSH)
  hoặc đến chu kỳ heartbeat
\end{enumerate}


{\thesissetfigureid{4.16}{fig-ch4-checklist-kiem-tra-he-thong-sau-khi-lap-at-trong-xe}
\begin{figure}[htbp]
\centering
\pandocbounded{\safeincludesvg{./assets/figures/07-chuong-4-trien-khai-hardware-hinh-4-16.svg}}
\caption{Checklist kiểm tra hệ thống sau khi lắp đặt trong xe}
\label{fig:ch4-checklist-kiem-tra-he-thong-sau-khi-lap-at-trong-xe}
\end{figure}
}

Sau khi phần cứng đã được đặt vào xe và các kết nối cơ bản đã được xác
nhận, lớp quyết định thiết bị phản ứng như thế nào trong từng tình huống
chính là firmware. Vì vậy, phần tiếp theo chuyển từ lắp đặt cơ khí -
điện sang triển khai phần mềm nhúng.


\begin{center}\rule{0.5\linewidth}{0.5pt}\end{center}

\subsubsection{4.2.3. Triển khai Firmware}

Phần này trình bày quá trình triển khai firmware cho thiết bị theo dõi
xe IoT, từ nền tảng phát triển đến cơ chế vận hành. Nội dung bao gồm môi
trường phát triển, cấu trúc mã nguồn, các module chức năng chính và lưu
đồ thuật toán điều khiển toàn hệ thống.

\paragraph{4.2.3.1. Môi trường phát triển và công
cụ}\label{4231-muxf4i-trux1b0ux1eddng-phuxe1t-triux1ec3n-vuxe0-cuxf4ng-cux1ee5}

\paragraph{a) Framework và toolchain}\label{a-framework-vuxe0-toolchain}

Firmware được phát triển trên nền tảng ESP-IDF phiên bản 5.4, tức bộ
công cụ chính thức của Espressif dành cho ESP32-S3. Có thể xem ESP-IDF
như ``bộ nền'' để viết firmware: nó cung cấp FreeRTOS để điều phối tác
vụ, các driver ngoại vi để làm việc với phần cứng, NimBLE để xử lý
Bluetooth và hệ thống build để biên dịch toàn bộ dự án.


{\thesissettableid{4.5}{tab-ch4-cong-cu-phat-trien-firmware}
\def\LTcaptype{table}
\begin{longtable}{|L{0.208\linewidth}|L{0.355\linewidth}|L{0.401\linewidth}|}
\caption{Công cụ phát triển firmware}\label{tab:ch4-cong-cu-phat-trien-firmware}\\
\hline \textbf{Thành phần} & \textbf{Công cụ / Phiên bản} & \textbf{Mục đích} \\
\\\hline
\endfirsthead
\caption[]{Công cụ phát triển firmware (tiếp theo)}\\
\hline \textbf{Thành phần} & \textbf{Công cụ / Phiên bản} & \textbf{Mục đích} \\
\\\hline
\endhead
\\\hline
\endfoot
\endlastfoot
Framework & ESP-IDF v5.4 & SDK chính thức cho ESP32-S3 \\\hline
IDE & VS Code + ESP-IDF Extension & Biên tập, debug, flash firmware \\\hline
Build system & CMake + Ninja & Biên dịch và liên kết mã nguồn \\\hline
BLE stack & NimBLE (tích hợp ESP-IDF) & Giao tiếp Bluetooth Low
Energy \\\hline
RTOS & FreeRTOS (tích hợp ESP-IDF) & Quản lý tác vụ đa luồng \\\hline
Debug & JTAG / UART Monitor & Giám sát log và debug thời gian thực \\\hline
Version control & Git & Quản lý phiên bản mã nguồn \\\hline
\end{longtable}
}


\paragraph{b) Quy trình build và nạp
firmware}\label{b-quy-truxecnh-build-vuxe0-nux1ea1p-firmware}

Quy trình biên dịch firmware đi theo ba bước quen thuộc: cấu hình dự án,
biên dịch toàn bộ mã nguồn, rồi nạp firmware xuống vi điều khiển. Trong
ESP-IDF, bước cấu hình được thực hiện qua tiện ích cấu hình của bộ SDK;
sau đó hệ thống build dựa trên CMake và Ninja sẽ tạo file firmware để
nạp qua UART.

\begin{Shaded}
\begin{Highlighting}[]
\CommentTok{\# Cấu hình dự án (chọn target ESP32{-}S3)}
\ExtensionTok{idf.py}\NormalTok{ set{-}target esp32s3}

\CommentTok{\# Cấu hình tham số (BLE, UART, GPIO, ...)}
\ExtensionTok{idf.py}\NormalTok{ menuconfig}

\CommentTok{\# Biên dịch dự án}
\ExtensionTok{idf.py}\NormalTok{ build}

\CommentTok{\# Nạp firmware và giám sát log qua cổng nối tiếp}
\ExtensionTok{idf.py} \AttributeTok{{-}p}\NormalTok{ SERIAL\_PORT flash monitor}
\end{Highlighting}
\end{Shaded}

\paragraph{c) Cấu hình NimBLE BLE
stack}\label{c-cux1ea5u-huxecnh-nimble-ble-stack}

NimBLE được chọn thay cho Bluedroid mặc định của ESP-IDF vì nó gọn hơn
về bộ nhớ nhưng vẫn đủ chức năng BLE Central để kết nối với adapter
OBD2. Nói cách khác, đây là lựa chọn thực dụng: tiết kiệm RAM cho
firmware mà không làm mất các khả năng BLE cần thiết của đề tài.

Cấu hình NimBLE được lưu trong tệp cấu hình dự án của ESP-IDF
(\texttt{sdkconfig}), thường được thiết lập qua tiện ích cấu hình của bộ
SDK:

\begin{verbatim}
CONFIG_BT_ENABLED=y
CONFIG_BT_NIMBLE_ENABLED=y
CONFIG_BT_NIMBLE_MAX_CONNECTIONS=1
CONFIG_BT_NIMBLE_ROLE_CENTRAL=y
CONFIG_BT_NIMBLE_ROLE_PERIPHERAL=n
\end{verbatim}

\begin{center}\rule{0.5\linewidth}{0.5pt}\end{center}

\paragraph{4.2.3.2. Cấu trúc mã nguồn
firmware}\label{4232-cux1ea5u-truxfac-muxe3-nguux1ed3n-firmware}

\paragraph{a) Tổ chức thư
mục}\label{a-tux1ed5-chux1ee9c-thux1b0-mux1ee5c}

Cấu trúc mã nguồn firmware hiện tại bám sát mã nguồn ESP-IDF đang chạy
trên thiết bị, tập trung vào các mô-đun theo chức năng thay vì chia
thành nhiều task nghiệp vụ riêng biệt:

\begin{itemize}
\tightlist
\item
  \texttt{main/main.c}: điểm vào chính của firmware.
\item
  \texttt{main/inc/}: khai báo module chức năng
  (\texttt{app\_config}, \texttt{app\_state}, \texttt{command\_handler},
  \texttt{data\_formatter}, \texttt{mqtt\_client}, \texttt{modem\_*},
  \texttt{ble\_*}, \texttt{power\_mgr}, \texttt{util}).
\item
  \texttt{main/src/}: hiện thực nghiệp vụ lõi
  (\texttt{state\_machine.c}, \texttt{command\_handler.c},
  \texttt{data\_formatter.c}, \texttt{mqtt\_client.c},
  \texttt{nvs\_config.c}, \texttt{power\_mgr.c}, \texttt{modem\_*},
  \texttt{ble\_*}, \texttt{adc\_reader.c}, \texttt{imu\_lis3dh.c},
  \texttt{util.c}).
\item
  \texttt{sdkconfig}, \texttt{CMakeLists.txt}: tệp cấu hình build và tham
  số runtime chính của dự án ESP-IDF.
\end{itemize}

\paragraph{b) Kiến trúc phân
tầng}\label{b-kiux1ebfn-truxfac-phuxe2n-tux1ea7ng}

Firmware được tổ chức quanh bốn nhóm chức năng chính, phản ánh đúng cách
chia mô-đun trong mã nguồn hiện tại:


{\thesissettableid{4.6}{tab-ch4-kien-truc-phan-tang-firmware}
\def\LTcaptype{table}
\begin{longtable}{|L{0.160\linewidth}|L{0.304\linewidth}|L{0.500\linewidth}|}
\caption{Kiến trúc phân tầng firmware}\label{tab:ch4-kien-truc-phan-tang-firmware}\\
\hline \textbf{Nhóm} & \textbf{Thành phần} & \textbf{Chức năng} \\
\\\hline
\endfirsthead
\caption[]{Kiến trúc phân tầng firmware (tiếp theo)}\\
\hline \textbf{Nhóm} & \textbf{Thành phần} & \textbf{Chức năng} \\
\\\hline
\endhead
\\\hline
\endfoot
\endlastfoot
1 & \texttt{adc\_reader}, \texttt{imu\_lis3dh}, \texttt{modem\_at},
\texttt{modem\_lte}, \texttt{modem\_gnss}, \texttt{ble\_mgr},
\texttt{ble\_obd} & Truy cập phần cứng và giao thức ngoại vi \\\hline
2 & \texttt{power\_mgr}, \texttt{nvs\_config}, \texttt{app\_state} &
Quản lý nguồn, ngưỡng LVD, cấu hình lưu bền và ngữ cảnh qua deep
sleep \\\hline
3 & \texttt{state\_machine}, \texttt{command\_handler} & Điều phối trạng
thái vận hành, lệnh điều khiển, OTA và rollback \\\hline
4 & \texttt{data\_formatter}, \texttt{mqtt\_client}, \texttt{util} &
Đóng gói gói dữ liệu, giao tiếp MQTT, tải firmware OTA và kiểm tra
SHA-256 \\\hline
\end{longtable}
}


Việc đồng bộ giữa các mô-đun được thực hiện có chọn lọc bằng mutex,
semaphore và queue ở cấp driver. Cách tiếp cận này giúp firmware vẫn tận
dụng được FreeRTOS nhưng tránh phân mảnh logic sang quá nhiều task ứng
dụng khó kiểm soát.

\paragraph{c) Định nghĩa chân
GPIO}\label{c-ux111ux1ecbnh-nghux129a-chuxe2n-gpio}


{\thesissettableid{4.7}{tab-ch4-bang-anh-xa-chan-gpio}
\def\LTcaptype{table}
\begin{longtable}{|L{0.115\linewidth}|L{0.221\linewidth}|L{0.158\linewidth}|L{0.460\linewidth}|}
\caption{Bảng ánh xạ chân GPIO}\label{tab:ch4-bang-anh-xa-chan-gpio}\\
\hline \textbf{GPIO} & \textbf{Chức năng} & \textbf{Hướng} & \textbf{Mô tả} \\
\\\hline
\endfirsthead
\caption[]{Bảng ánh xạ chân GPIO (tiếp theo)}\\
\hline \textbf{GPIO} & \textbf{Chức năng} & \textbf{Hướng} & \textbf{Mô tả} \\
\\\hline
\endhead
\\\hline
\endfoot
\endlastfoot
2 & IGN\_IN & Input & Đọc trạng thái khóa điện (hoặc qua OBD2) \\\hline
4 & \nolinkurl{U_BATT_ADC} & Input & Đọc điện áp ắc quy (ADC 12-bit) \\\hline
5 & \nolinkurl{CHARGER_EN} & Output & Điều khiển khối sạc TP5100 \\\hline
16 & \nolinkurl{MODEM_UART_TX} & Output & UART TX đến modem SIM7600CE-T \\\hline
17 & \nolinkurl{MODEM_UART_RX} & Input & UART RX từ modem SIM7600CE-T \\\hline
18 & \nolinkurl{POWER_PATH_EN} & Output & Chọn nguồn cấp (ắc quy/pin dự phòng) \\\hline
19 & LVD\_STATUS & Input & Trạng thái từ khối LVD phần cứng \\\hline
41 & \nolinkurl{LIS3DSH_INT1} & Input & Ngắt chính từ cảm biến gia tốc IMU \\\hline
42 & \nolinkurl{LIS3DSH_INT2} & Input & Ngắt phụ, đang giữ cho mở rộng \\\hline
2 & \nolinkurl{LIS3DSH_SDA} & I/O & Đường dữ liệu I2C \\\hline
1 & \nolinkurl{LIS3DSH_SCL} & I/O & Đường clock I2C \\\hline
26 & \nolinkurl{MODEM_PWR-KEY} & Output & Điều khiển nguồn modem \\\hline
\end{longtable}
}


\begin{center}\rule{0.5\linewidth}{0.5pt}\end{center}

\paragraph{4.2.3.3. Triển khai module
BLE-OBD2}\label{4233-triux1ec3n-khai-module-ble-obd2}

\paragraph{a) Tổng quan giao tiếp BLE với adapter
OBD2}\label{a-tux1ed5ng-quan-giao-tiux1ebfp-ble-vux1edbi-adapter-obd2}

Module BLE-OBD2 thiết lập kết nối Bluetooth Low Energy với adapter vgate
iCar Pro tại cổng OBD2 để gửi yêu cầu đọc dữ liệu theo chuẩn OBD2 (Mode
0x01 - Current Data) và phân tích phản hồi, từ đó trích xuất các thông
số vận hành động cơ.

Chu trình giao tiếp BLE-OBD2 gồm các giai đoạn: quét thiết bị (scan) ---
kết nối (connect) --- khám phá dịch vụ GATT (service discovery) --- gửi
lệnh OBD2 (write characteristic) --- nhận phản hồi (notification
callback).

\paragraph{b) Khởi tạo NimBLE
stack}\label{b-khux1edfi-tux1ea1o-nimble-stack}

Việc khởi tạo BLE stack được thực hiện trong hàm
\texttt{ble\_init\_stack()}. NimBLE chạy trong một FreeRTOS task riêng
biệt, xử lý toàn bộ các sự kiện BLE thông qua vòng lặp sự kiện (event
loop).

\begin{Shaded}
\begin{Highlighting}[]
\DataTypeTok{void}\NormalTok{ ble\_init\_stack}\OperatorTok{(}\NormalTok{ble\_init\_config\_t }\DataTypeTok{const} \OperatorTok{*}\NormalTok{config}\OperatorTok{)}
\OperatorTok{\{}
    \CommentTok{// Khởi tạo cổng NimBLE}
\NormalTok{    esp\_err\_t ret }\OperatorTok{=}\NormalTok{ nimble\_port\_init}\OperatorTok{();}
    \ControlFlowTok{if} \OperatorTok{(}\NormalTok{ret }\OperatorTok{!=}\NormalTok{ ESP\_OK}\OperatorTok{)} \OperatorTok{\{}
\NormalTok{        ESP\_LOGE}\OperatorTok{(}\NormalTok{TAG}\OperatorTok{,} \StringTok{"Khởi tạo NimBLE thất bại"}\OperatorTok{);}
        \ControlFlowTok{return}\OperatorTok{;}
    \OperatorTok{\}}

    \CommentTok{// Cấu hình callback cho host}
\NormalTok{    ble\_hs\_cfg}\OperatorTok{.}\NormalTok{reset\_cb        }\OperatorTok{=}\NormalTok{ config}\OperatorTok{{-}\textgreater{}}\NormalTok{reset\_cb}\OperatorTok{;}
\NormalTok{    ble\_hs\_cfg}\OperatorTok{.}\NormalTok{sync\_cb         }\OperatorTok{=}\NormalTok{ config}\OperatorTok{{-}\textgreater{}}\NormalTok{sync\_cb}\OperatorTok{;}
\NormalTok{    ble\_hs\_cfg}\OperatorTok{.}\NormalTok{store\_status\_cb }\OperatorTok{=}\NormalTok{ ble\_store\_util\_status\_rr}\OperatorTok{;}
\NormalTok{    ble\_store\_config\_init}\OperatorTok{();}

    \CommentTok{// Tạo task chạy NimBLE event loop}
\NormalTok{    xTaskCreate}\OperatorTok{(}\NormalTok{ble\_task}\OperatorTok{,} \StringTok{"NimBLE\_task"}\OperatorTok{,} \DecValTok{4096}\OperatorTok{,}\NormalTok{ NULL}\OperatorTok{,} \DecValTok{5}\OperatorTok{,}\NormalTok{ NULL}\OperatorTok{);}
\OperatorTok{\}}

\DataTypeTok{static} \DataTypeTok{void}\NormalTok{ ble\_task}\OperatorTok{(}\DataTypeTok{void} \OperatorTok{*}\NormalTok{param}\OperatorTok{)}
\OperatorTok{\{}
\NormalTok{    nimble\_port\_run}\OperatorTok{();}  \CommentTok{// Blocking {-} chạy event loop BLE}
\OperatorTok{\}}
\end{Highlighting}
\end{Shaded}

\paragraph{c) Quy trình quét và kết nối thiết
bị}\label{c-quy-truxecnh-quuxe9t-vuxe0-kux1ebft-nux1ed1i-thiux1ebft-bux1ecb}

Firmware thực hiện quét BLE để tìm adapter OBD2 dựa trên tên thiết bị
(chứa chuỗi "iCar") hoặc UUID dịch vụ. Khi phát hiện thiết bị phù hợp,
hệ thống sẽ kết nối và khám phá dịch vụ GATT để xác định characteristic
dùng cho truyền/nhận dữ liệu.

\begin{Shaded}
\begin{Highlighting}[]
\CommentTok{// Tham số quét BLE {-} chế độ thụ động (passive) để tiết kiệm năng lượng}
\DataTypeTok{static} \DataTypeTok{const} \KeywordTok{struct}\NormalTok{ ble\_gap\_disc\_params disc\_params }\OperatorTok{=} \OperatorTok{\{}
    \OperatorTok{.}\NormalTok{passive           }\OperatorTok{=} \DecValTok{1}\OperatorTok{,}
    \OperatorTok{.}\NormalTok{itvl              }\OperatorTok{=} \BaseNTok{0x0010}\OperatorTok{,}   \CommentTok{// Khoảng cách quét}
    \OperatorTok{.}\NormalTok{window            }\OperatorTok{=} \BaseNTok{0x0010}\OperatorTok{,}   \CommentTok{// Cửa sổ quét}
    \OperatorTok{.}\NormalTok{filter\_duplicates }\OperatorTok{=} \DecValTok{1}\OperatorTok{,}        \CommentTok{// Lọc quảng cáo trùng lặp}
\OperatorTok{\};}

\CommentTok{// Callback xử lý sự kiện GAP (scan, connect, disconnect)}
\DataTypeTok{static} \DataTypeTok{int}\NormalTok{ ble\_mgr\_gap\_event\_cb}\OperatorTok{(}\KeywordTok{struct}\NormalTok{ ble\_gap\_event }\OperatorTok{*}\NormalTok{event}\OperatorTok{,} \DataTypeTok{void} \OperatorTok{*}\NormalTok{arg}\OperatorTok{)}
\OperatorTok{\{}
    \ControlFlowTok{switch} \OperatorTok{(}\NormalTok{event}\OperatorTok{{-}\textgreater{}}\NormalTok{type}\OperatorTok{)} \OperatorTok{\{}
    \ControlFlowTok{case}\NormalTok{ BLE\_GAP\_EVENT\_DISC}\OperatorTok{:}
        \CommentTok{// Phân tích dữ liệu quảng cáo}
\NormalTok{        ble\_hs\_adv\_parse\_fields}\OperatorTok{(\&}\NormalTok{adv\_fields}\OperatorTok{,}
\NormalTok{                                event}\OperatorTok{{-}\textgreater{}}\NormalTok{disc}\OperatorTok{.}\NormalTok{data}\OperatorTok{,}
\NormalTok{                                event}\OperatorTok{{-}\textgreater{}}\NormalTok{disc}\OperatorTok{.}\NormalTok{length\_data}\OperatorTok{);}
        \CommentTok{// Kiểm tra UUID dịch vụ có trùng khớp không}
        \ControlFlowTok{if} \OperatorTok{(}\NormalTok{ble\_mgr\_adv\_contains\_service}\OperatorTok{(\&}\NormalTok{adv\_fields}\OperatorTok{,}
\NormalTok{                                          disc\_cfg}\OperatorTok{{-}\textgreater{}}\NormalTok{svc\_def}\OperatorTok{{-}\textgreater{}}\NormalTok{service\_uuid}\OperatorTok{))} \OperatorTok{\{}
            \CommentTok{// Kết nối đến thiết bị}
\NormalTok{            ble\_gap\_connect}\OperatorTok{(}\NormalTok{BLE\_OWN\_ADDR\_PUBLIC}\OperatorTok{,}
                            \OperatorTok{\&}\NormalTok{event}\OperatorTok{{-}\textgreater{}}\NormalTok{disc}\OperatorTok{.}\NormalTok{addr}\OperatorTok{,}
                            \DecValTok{30000}\OperatorTok{,}           \CommentTok{// Timeout 30 giây}
                            \OperatorTok{\&}\NormalTok{conn\_params}\OperatorTok{,}
\NormalTok{                            ble\_mgr\_gap\_event\_cb}\OperatorTok{,}\NormalTok{ mgr\_ctx}\OperatorTok{);}
        \OperatorTok{\}}
        \ControlFlowTok{break}\OperatorTok{;}

    \ControlFlowTok{case}\NormalTok{ BLE\_GAP\_EVENT\_CONNECT}\OperatorTok{:}
        \ControlFlowTok{if} \OperatorTok{(}\NormalTok{event}\OperatorTok{{-}\textgreater{}}\NormalTok{connect}\OperatorTok{.}\NormalTok{status }\OperatorTok{==} \DecValTok{0}\OperatorTok{)} \OperatorTok{\{}
            \CommentTok{// Kết nối thành công {-} khám phá dịch vụ GATT}
\NormalTok{            ble\_gattc\_disc\_all\_svcs}\OperatorTok{(}\NormalTok{conn\_handle}\OperatorTok{,}
\NormalTok{                                     ble\_mgr\_gatt\_svc\_discovered\_cb}\OperatorTok{,}
\NormalTok{                                     mgr\_ctx}\OperatorTok{);}
        \OperatorTok{\}}
        \ControlFlowTok{break}\OperatorTok{;}
    \OperatorTok{\}}
    \ControlFlowTok{return} \DecValTok{0}\OperatorTok{;}
\OperatorTok{\}}
\end{Highlighting}
\end{Shaded}

Để tối ưu thời gian kết nối lại sau khi thức dậy từ deep sleep, firmware
lưu địa chỉ MAC của adapter OBD2 vào bộ nhớ flash. Khi thức dậy,
firmware ưu tiên kết nối trực tiếp đến địa chỉ đã lưu thay vì quét lại
từ đầu, giảm thời gian kết nối lại xuống còn 1--3 giây so với 3--10 giây
khi quét mới.

\paragraph{d) Gửi lệnh và phân tích phản hồi
OBD2}\label{d-gux1eedi-lux1ec7nh-vuxe0-phuxe2n-tuxedch-phux1ea3n-hux1ed3i-obd2}

Giao thức OBD2 sử dụng định dạng lệnh ELM327 truyền qua BLE GATT
characteristic. Mỗi lệnh có dạng \texttt{"MMPP\textbackslash{}r"}, trong
đó \texttt{MM} là mode (01 cho dữ liệu hiện tại) và \texttt{PP} là PID
(Parameter ID) cần đọc.

\begin{Shaded}
\begin{Highlighting}[]
\CommentTok{// Gửi lệnh OBD2 và đợi phản hồi}
\DataTypeTok{int}\NormalTok{ ble\_obd\_rxtx}\OperatorTok{(}\NormalTok{ble\_obd\_ctx\_t }\OperatorTok{*}\NormalTok{obd}\OperatorTok{,} \DataTypeTok{uint8\_t}\NormalTok{ mode}\OperatorTok{,} \DataTypeTok{uint8\_t}\NormalTok{ pid}\OperatorTok{,}
                 \DataTypeTok{uint32\_t}\NormalTok{ timeout\_ms}\OperatorTok{)}
\OperatorTok{\{}
    \CommentTok{// Định dạng lệnh: "MMPP\textbackslash{}r"}
\NormalTok{    snprintf}\OperatorTok{(}\NormalTok{obd}\OperatorTok{{-}\textgreater{}}\NormalTok{tx\_data}\OperatorTok{.}\NormalTok{buf}\OperatorTok{,} \KeywordTok{sizeof}\OperatorTok{(}\NormalTok{obd}\OperatorTok{{-}\textgreater{}}\NormalTok{tx\_data}\OperatorTok{.}\NormalTok{buf}\OperatorTok{),}
             \StringTok{"}\SpecialCharTok{\%02X\%02X\textbackslash{}r}\StringTok{"}\OperatorTok{,}\NormalTok{ mode}\OperatorTok{,}\NormalTok{ pid}\OperatorTok{);}

    \CommentTok{// Gá»­i qua GATT write characteristic}
\NormalTok{    ble\_mgr\_send}\OperatorTok{(}\NormalTok{obd}\OperatorTok{{-}\textgreater{}}\NormalTok{mgr\_ctx}\OperatorTok{,}\NormalTok{ obd\_tx\_char}\OperatorTok{{-}\textgreater{}}\NormalTok{handle}\OperatorTok{,}
\NormalTok{                 obd}\OperatorTok{{-}\textgreater{}}\NormalTok{tx\_data}\OperatorTok{.}\NormalTok{buf}\OperatorTok{,}\NormalTok{ strlen}\OperatorTok{(}\NormalTok{obd}\OperatorTok{{-}\textgreater{}}\NormalTok{tx\_data}\OperatorTok{.}\NormalTok{buf}\OperatorTok{));}

    \CommentTok{// Đợi phản hồi (semaphore với timeout)}
    \ControlFlowTok{if} \OperatorTok{(}\NormalTok{xSemaphoreTake}\OperatorTok{(}\NormalTok{obd}\OperatorTok{{-}\textgreater{}}\NormalTok{api}\OperatorTok{.}\NormalTok{response\_sem}\OperatorTok{,}
\NormalTok{                       pdMS\_TO\_TICKS}\OperatorTok{(}\NormalTok{timeout\_ms}\OperatorTok{))} \OperatorTok{!=}\NormalTok{ pdTRUE}\OperatorTok{)} \OperatorTok{\{}
        \ControlFlowTok{return} \OperatorTok{{-}}\DecValTok{1}\OperatorTok{;}  \CommentTok{// Hết thời gian chờ}
    \OperatorTok{\}}
    \ControlFlowTok{return} \DecValTok{0}\OperatorTok{;}
\OperatorTok{\}}
\end{Highlighting}
\end{Shaded}

Phản hồi OBD2 được xử lý trong notification callback. Dữ liệu trả về có
dạng \texttt{"41\ 0C\ 1F\ 40\textbackslash{}r"}, trong đó byte đầu là
mode + 0x40 (chỉ thị phản hồi), byte thứ hai là PID, và các byte còn lại
là dữ liệu cần giải mã.

\begin{Shaded}
\begin{Highlighting}[]
\CommentTok{// Phân tích phản hồi OBD2}
\DataTypeTok{static} \DataTypeTok{void}\NormalTok{ ble\_obd\_process\_obd\_data}\OperatorTok{(}\NormalTok{ble\_obd\_ctx\_t }\OperatorTok{*}\NormalTok{obd}\OperatorTok{,}
                                      \DataTypeTok{char} \OperatorTok{*}\NormalTok{data}\OperatorTok{,} \DataTypeTok{size\_t}\NormalTok{ len}\OperatorTok{)}
\OperatorTok{\{}
    \DataTypeTok{char} \OperatorTok{*}\NormalTok{saveptr}\OperatorTok{;}
    \DataTypeTok{char} \OperatorTok{*}\NormalTok{tok }\OperatorTok{=}\NormalTok{ strtok\_r}\OperatorTok{(}\NormalTok{data}\OperatorTok{,} \StringTok{" }\SpecialCharTok{\textbackslash{}r}\StringTok{"}\OperatorTok{,} \OperatorTok{\&}\NormalTok{saveptr}\OperatorTok{);}
    \DataTypeTok{uint8\_t}\NormalTok{ values}\OperatorTok{[}\NormalTok{BLE\_OBD\_MAX\_DATA\_LEN}\OperatorTok{];}
    \DataTypeTok{int}\NormalTok{ count }\OperatorTok{=} \DecValTok{0}\OperatorTok{;}

    \ControlFlowTok{while} \OperatorTok{(}\NormalTok{tok }\OperatorTok{\&\&}\NormalTok{ count }\OperatorTok{\textless{}}\NormalTok{ BLE\_OBD\_MAX\_DATA\_LEN}\OperatorTok{)} \OperatorTok{\{}
        \DataTypeTok{long}\NormalTok{ val }\OperatorTok{=}\NormalTok{ strtol}\OperatorTok{(}\NormalTok{tok}\OperatorTok{,}\NormalTok{ NULL}\OperatorTok{,} \DecValTok{16}\OperatorTok{);}  \CommentTok{// Chuyển hex sang số nguyên}
\NormalTok{        values}\OperatorTok{[}\NormalTok{count}\OperatorTok{++]} \OperatorTok{=} \OperatorTok{(}\DataTypeTok{uint8\_t}\OperatorTok{)}\NormalTok{val}\OperatorTok{;}
\NormalTok{        tok }\OperatorTok{=}\NormalTok{ strtok\_r}\OperatorTok{(}\NormalTok{NULL}\OperatorTok{,} \StringTok{" }\SpecialCharTok{\textbackslash{}r}\StringTok{"}\OperatorTok{,} \OperatorTok{\&}\NormalTok{saveptr}\OperatorTok{);}
    \OperatorTok{\}}

    \CommentTok{// Xác thực phản hồi: kiểm tra mode và PID trùng khớp}
    \ControlFlowTok{if} \OperatorTok{(}\NormalTok{values}\OperatorTok{[}\DecValTok{0}\OperatorTok{]} \OperatorTok{==} \OperatorTok{(}\NormalTok{obd}\OperatorTok{{-}\textgreater{}}\NormalTok{tx\_data}\OperatorTok{.}\NormalTok{mode }\OperatorTok{+} \BaseNTok{0x40}\OperatorTok{)} \OperatorTok{\&\&}
\NormalTok{        values}\OperatorTok{[}\DecValTok{1}\OperatorTok{]} \OperatorTok{==}\NormalTok{ obd}\OperatorTok{{-}\textgreater{}}\NormalTok{tx\_data}\OperatorTok{.}\NormalTok{pid}\OperatorTok{)} \OperatorTok{\{}
        \CommentTok{// Gọi callback với dữ liệu đã trích xuất}
\NormalTok{        obd}\OperatorTok{{-}\textgreater{}}\NormalTok{response\_cb}\OperatorTok{(}\NormalTok{obd}\OperatorTok{{-}\textgreater{}}\NormalTok{tx\_data}\OperatorTok{.}\NormalTok{pid}\OperatorTok{,}
\NormalTok{                         values }\OperatorTok{+} \DecValTok{2}\OperatorTok{,}      \CommentTok{// Bỏ qua mode và PID}
\NormalTok{                         count }\OperatorTok{{-}} \DecValTok{2}\OperatorTok{,}       \CommentTok{// Độ dài dữ liệu thực}
\NormalTok{                         obd}\OperatorTok{{-}\textgreater{}}\NormalTok{usr\_ctx}\OperatorTok{);}
    \OperatorTok{\}}
\OperatorTok{\}}
\end{Highlighting}
\end{Shaded}

\paragraph{e) Bảng PID OBD2 sử dụng trong hệ
thống}\label{e-bux1ea3ng-pid-obd2-sux1eed-dux1ee5ng-trong-hux1ec7-thux1ed1ng}


{\thesissettableid{4.8}{tab-ch4-cac-pid-obd2-uoc-oc-inh-ky-mode-0x01-current-data}
\def\LTcaptype{table}
\begin{longtable}{|L{0.130\linewidth}|L{0.290\linewidth}|L{0.120\linewidth}|L{0.230\linewidth}|L{0.174\linewidth}|}
\caption{Các PID OBD2 được đọc định kỳ (Mode 0x01 - Current Data)}\label{tab:ch4-cac-pid-obd2-uoc-oc-inh-ky-mode-0x01-current-data}\\
\hline \textbf{PID (Hex)} & \textbf{Thông số} & \textbf{Độ dài} & \textbf{Công thức tính} & \textbf{Đơn vị} \\
\\\hline
\endfirsthead
\caption[]{Các PID OBD2 được đọc định kỳ (Mode 0x01 - Current Data) (tiếp theo)}\\
\hline \textbf{PID (Hex)} & \textbf{Thông số} & \textbf{Độ dài} & \textbf{Công thức tính} & \textbf{Đơn vị} \\
\\\hline
\endhead
\\\hline
\endfoot
\endlastfoot
0x0C & Tốc độ động cơ (RPM) & 2 byte & \(\dfrac{A \times 256 + B}{4}\) & vòng/phút \\\hline
0x0D & Tốc độ xe & 1 byte & A & km/h \\\hline
0x04 & Tải động cơ & 1 byte & \(\dfrac{A \times 100}{255}\) & \% \\\hline
0x05 & Nhiệt độ nước làm mát & 1 byte & \(A - 40\) & độ C \\\hline
0x2F & Mức nhiên liệu & 1 byte & \(\dfrac{A \times 100}{255}\) & \% \\\hline
0x0F & Nhiệt độ khí nạp & 1 byte & \(A - 40\) & độ C \\\hline
0x11 & Vị trí bướm ga & 1 byte & \(\dfrac{A \times 100}{255}\) & \% \\\hline
0x42 & Điện áp module điều khiển & 2 byte & \(\dfrac{A \times 256 + B}{1000}\) & V \\\hline
\end{longtable}
}


Các hàm chuyển đổi tương ứng được định nghĩa trong component
\texttt{obd2\_parser}:

\begin{Shaded}
\begin{Highlighting}[]
\CommentTok{// Chuyển đổi RPM: (A * 256 + B) / 4}
\DataTypeTok{static} \DataTypeTok{int}\NormalTok{ obd\_conv\_rpm}\OperatorTok{(}\DataTypeTok{int32\_t} \OperatorTok{*}\NormalTok{value}\OperatorTok{,} \DataTypeTok{uint8\_t} \DataTypeTok{const} \OperatorTok{*}\NormalTok{data}\OperatorTok{,} \DataTypeTok{size\_t}\NormalTok{ len}\OperatorTok{)}
\OperatorTok{\{}
    \OperatorTok{*}\NormalTok{value }\OperatorTok{=} \OperatorTok{((}\NormalTok{data}\OperatorTok{[}\DecValTok{0}\OperatorTok{]} \OperatorTok{\textless{}\textless{}} \DecValTok{8}\OperatorTok{)} \OperatorTok{|}\NormalTok{ data}\OperatorTok{[}\DecValTok{1}\OperatorTok{])} \OperatorTok{/} \DecValTok{4}\OperatorTok{;}
    \ControlFlowTok{return} \DecValTok{0}\OperatorTok{;}
\OperatorTok{\}}

\CommentTok{// Chuyển đổi phần trăm: (A * 100) / 255}
\DataTypeTok{static} \DataTypeTok{int}\NormalTok{ obd\_conv\_percent}\OperatorTok{(}\DataTypeTok{int32\_t} \OperatorTok{*}\NormalTok{value}\OperatorTok{,} \DataTypeTok{uint8\_t} \DataTypeTok{const} \OperatorTok{*}\NormalTok{data}\OperatorTok{,} \DataTypeTok{size\_t}\NormalTok{ len}\OperatorTok{)}
\OperatorTok{\{}
    \OperatorTok{*}\NormalTok{value }\OperatorTok{=} \OperatorTok{(}\NormalTok{data}\OperatorTok{[}\DecValTok{0}\OperatorTok{]} \OperatorTok{*} \DecValTok{100}\OperatorTok{)} \OperatorTok{/} \DecValTok{255}\OperatorTok{;}
    \ControlFlowTok{return} \DecValTok{0}\OperatorTok{;}
\OperatorTok{\}}

\CommentTok{// Chuyển đổi nhiệt độ: A {-} 40}
\DataTypeTok{static} \DataTypeTok{int}\NormalTok{ obd\_conv\_temperature}\OperatorTok{(}\DataTypeTok{int32\_t} \OperatorTok{*}\NormalTok{value}\OperatorTok{,} \DataTypeTok{uint8\_t} \DataTypeTok{const} \OperatorTok{*}\NormalTok{data}\OperatorTok{,} \DataTypeTok{size\_t}\NormalTok{ len}\OperatorTok{)}
\OperatorTok{\{}
    \OperatorTok{*}\NormalTok{value }\OperatorTok{=}\NormalTok{ data}\OperatorTok{[}\DecValTok{0}\OperatorTok{]} \OperatorTok{{-}} \DecValTok{40}\OperatorTok{;}
    \ControlFlowTok{return} \DecValTok{0}\OperatorTok{;}
\OperatorTok{\}}
\end{Highlighting}
\end{Shaded}

\paragraph{f) Xử lý lỗi và
fallback}\label{f-xux1eed-luxfd-lux1ed7i-vuxe0-fallback}

Khi thiết bị theo dõi không thể kết nối BLE với adapter OBD2 trong thời gian chờ
cho phép, firmware không dừng lại mà chuyển sang cơ chế dự phòng. Lúc
này, thiết bị đọc điện áp ắc quy qua ADC để ước lượng trạng thái khóa
điện theo profile 12V hoặc 24V. Cách suy luận này không chính xác bằng
việc đọc trực tiếp qua OBD2, nhưng vẫn đủ để hệ thống tiếp tục theo dõi
và đưa ra quyết định cơ bản khi đường BLE gặp sự cố.

\begin{center}\rule{0.5\linewidth}{0.5pt}\end{center}

\paragraph{4.2.3.4. Triển khai module modem
MQTT}\label{4234-triux1ec3n-khai-module-modem-mqtt}

\paragraph{a) Khởi tạo và kiểm tra
modem}\label{a-khux1edfi-tux1ea1o-vuxe0-kiux1ec3m-tra-modem}

Module modem quản lý toàn bộ giao tiếp với modem SIMCom SIM7600CE-T
thông qua giao tiếp UART và tập lệnh AT. Quy trình khởi tạo bao gồm kiểm
tra phản hồi modem, trạng thái SIM, đăng ký mạng và chất lượng tín hiệu.

\begin{Shaded}
\begin{Highlighting}[]
\CommentTok{// Trình tự khởi tạo modem SIM7600CE{-}T}
\DataTypeTok{int}\NormalTok{ modem\_init}\OperatorTok{(}\NormalTok{modem\_ctx\_t }\OperatorTok{*}\NormalTok{ctx}\OperatorTok{)}
\OperatorTok{\{}
    \CommentTok{// Bước 1: Kiểm tra modem phản hồi}
    \ControlFlowTok{if} \OperatorTok{(}\NormalTok{modem\_send\_at}\OperatorTok{(}\NormalTok{ctx}\OperatorTok{,} \StringTok{"AT"}\OperatorTok{,} \StringTok{"OK"}\OperatorTok{,} \DecValTok{5000}\OperatorTok{)} \OperatorTok{!=} \DecValTok{0}\OperatorTok{)} \OperatorTok{\{}
        \CommentTok{// Modem không phản hồi {-} reset bằng GPIO PWR{-}KEY}
\NormalTok{        gpio\_set\_level}\OperatorTok{(}\NormalTok{MODEM\_PWR}\OperatorTok{{-}}\NormalTok{KEY}\OperatorTok{,} \DecValTok{0}\OperatorTok{);}
\NormalTok{        vTaskDelay}\OperatorTok{(}\NormalTok{pdMS\_TO\_TICKS}\OperatorTok{(}\DecValTok{1500}\OperatorTok{));}
\NormalTok{        gpio\_set\_level}\OperatorTok{(}\NormalTok{MODEM\_PWR}\OperatorTok{{-}}\NormalTok{KEY}\OperatorTok{,} \DecValTok{1}\OperatorTok{);}
\NormalTok{        vTaskDelay}\OperatorTok{(}\NormalTok{pdMS\_TO\_TICKS}\OperatorTok{(}\DecValTok{10000}\OperatorTok{));}  \CommentTok{// Đợi khởi động lại}
        \ControlFlowTok{if} \OperatorTok{(}\NormalTok{modem\_send\_at}\OperatorTok{(}\NormalTok{ctx}\OperatorTok{,} \StringTok{"AT"}\OperatorTok{,} \StringTok{"OK"}\OperatorTok{,} \DecValTok{5000}\OperatorTok{)} \OperatorTok{!=} \DecValTok{0}\OperatorTok{)} \OperatorTok{\{}
            \ControlFlowTok{return} \OperatorTok{{-}}\DecValTok{1}\OperatorTok{;}  \CommentTok{// Khởi tạo thất bại}
        \OperatorTok{\}}
    \OperatorTok{\}}

    \CommentTok{// Bước 2: Kiểm tra SIM card}
    \ControlFlowTok{if} \OperatorTok{(}\NormalTok{modem\_send\_at}\OperatorTok{(}\NormalTok{ctx}\OperatorTok{,} \StringTok{"AT+CPIN?"}\OperatorTok{,} \StringTok{"+CPIN: READY"}\OperatorTok{,} \DecValTok{5000}\OperatorTok{)} \OperatorTok{!=} \DecValTok{0}\OperatorTok{)} \OperatorTok{\{}
        \ControlFlowTok{return} \OperatorTok{{-}}\DecValTok{2}\OperatorTok{;}  \CommentTok{// SIM chưa sẵn sàng}
    \OperatorTok{\}}

    \CommentTok{// Bước 3: Đợi đăng ký mạng (tối đa 60 giây)}
    \ControlFlowTok{for} \OperatorTok{(}\DataTypeTok{int}\NormalTok{ i }\OperatorTok{=} \DecValTok{0}\OperatorTok{;}\NormalTok{ i }\OperatorTok{\textless{}} \DecValTok{12}\OperatorTok{;}\NormalTok{ i}\OperatorTok{++)} \OperatorTok{\{}
        \ControlFlowTok{if} \OperatorTok{(}\NormalTok{modem\_send\_at}\OperatorTok{(}\NormalTok{ctx}\OperatorTok{,} \StringTok{"AT+CEREG?"}\OperatorTok{,} \StringTok{"+CEREG: 0,1"}\OperatorTok{,} \DecValTok{5000}\OperatorTok{)} \OperatorTok{==} \DecValTok{0}\OperatorTok{)} \OperatorTok{\{}
            \ControlFlowTok{break}\OperatorTok{;}  \CommentTok{// Đã đăng ký mạng}
        \OperatorTok{\}}
\NormalTok{        vTaskDelay}\OperatorTok{(}\NormalTok{pdMS\_TO\_TICKS}\OperatorTok{(}\DecValTok{5000}\OperatorTok{));}
    \OperatorTok{\}}

    \CommentTok{// Bước 4: Kiểm tra chất lượng tín hiệu}
\NormalTok{    modem\_send\_at}\OperatorTok{(}\NormalTok{ctx}\OperatorTok{,} \StringTok{"AT+CSQ"}\OperatorTok{,} \StringTok{"+CSQ:"}\OperatorTok{,} \DecValTok{3000}\OperatorTok{);}

    \CommentTok{// Bước 5: Thiết lập APN và kích hoạt PDP context}
\NormalTok{    modem\_send\_at}\OperatorTok{(}\NormalTok{ctx}\OperatorTok{,} \StringTok{"AT+CGDCONT=1,}\SpecialCharTok{\textbackslash{}"}\StringTok{IP}\SpecialCharTok{\textbackslash{}"}\StringTok{,}\SpecialCharTok{\textbackslash{}"}\StringTok{internet}\SpecialCharTok{\textbackslash{}"}\StringTok{"}\OperatorTok{,} \StringTok{"OK"}\OperatorTok{,} \DecValTok{5000}\OperatorTok{);}
\NormalTok{    modem\_send\_at}\OperatorTok{(}\NormalTok{ctx}\OperatorTok{,} \StringTok{"AT+CGACT=1,1"}\OperatorTok{,} \StringTok{"OK"}\OperatorTok{,} \DecValTok{10000}\OperatorTok{);}

    \ControlFlowTok{return} \DecValTok{0}\OperatorTok{;}  \CommentTok{// Khởi tạo thành công}
\OperatorTok{\}}
\end{Highlighting}
\end{Shaded}

\paragraph{b) Kết nối và gửi dữ liệu
MQTT}\label{b-kux1ebft-nux1ed1i-vuxe0-gux1eedi-dux1eef-liux1ec7u-mqtt}

Trong triển khai hiện tại, modem SIM7600CE-T đảm nhiệm việc đăng ký mạng
và kích hoạt PDP context bằng AT command; sau khi đường truyền dữ liệu
sẵn sàng, firmware sử dụng \texttt{esp\_mqtt\_client} của ESP-IDF để mở
phiên MQTT, gửi dữ liệu vận hành và đăng ký nhận topic lệnh điều khiển. Cách
làm này giúp phần truyền MQTT tách biệt rõ khỏi phần điều khiển modem,
dễ debug hơn và tận dụng trực tiếp event loop của ESP-IDF.

\begin{Shaded}
\begin{Highlighting}[]
\DataTypeTok{static} \DataTypeTok{void}\NormalTok{ tracker\_mqtt\_build\_topics}\OperatorTok{(}\DataTypeTok{void}\OperatorTok{)} \OperatorTok{\{}
\NormalTok{    snprintf}\OperatorTok{(}\NormalTok{s\_topic\_rawdata}\OperatorTok{,} \KeywordTok{sizeof}\OperatorTok{(}\NormalTok{s\_topic\_rawdata}\OperatorTok{),} \StringTok{"v1/}\SpecialCharTok{\%s}\StringTok{/rawdata"}\OperatorTok{,}\NormalTok{ s\_cfg}\OperatorTok{.}\NormalTok{device\_id}\OperatorTok{);}
\NormalTok{    snprintf}\OperatorTok{(}\NormalTok{s\_topic\_status}\OperatorTok{,} \KeywordTok{sizeof}\OperatorTok{(}\NormalTok{s\_topic\_status}\OperatorTok{),} \StringTok{"v1/}\SpecialCharTok{\%s}\StringTok{/status"}\OperatorTok{,}\NormalTok{ s\_cfg}\OperatorTok{.}\NormalTok{device\_id}\OperatorTok{);}
\NormalTok{    snprintf}\OperatorTok{(}\NormalTok{s\_topic\_events}\OperatorTok{,} \KeywordTok{sizeof}\OperatorTok{(}\NormalTok{s\_topic\_events}\OperatorTok{),} \StringTok{"v1/}\SpecialCharTok{\%s}\StringTok{/events"}\OperatorTok{,}\NormalTok{ s\_cfg}\OperatorTok{.}\NormalTok{device\_id}\OperatorTok{);}
\NormalTok{    snprintf}\OperatorTok{(}\NormalTok{s\_topic\_firmware}\OperatorTok{,} \KeywordTok{sizeof}\OperatorTok{(}\NormalTok{s\_topic\_firmware}\OperatorTok{),} \StringTok{"v1/}\SpecialCharTok{\%s}\StringTok{/firmware"}\OperatorTok{,}\NormalTok{ s\_cfg}\OperatorTok{.}\NormalTok{device\_id}\OperatorTok{);}
\NormalTok{    snprintf}\OperatorTok{(}\NormalTok{s\_topic\_commands}\OperatorTok{,} \KeywordTok{sizeof}\OperatorTok{(}\NormalTok{s\_topic\_commands}\OperatorTok{),} \StringTok{"v1/}\SpecialCharTok{\%s}\StringTok{/commands"}\OperatorTok{,}\NormalTok{ s\_cfg}\OperatorTok{.}\NormalTok{device\_id}\OperatorTok{);}
\OperatorTok{\}}

\DataTypeTok{static} \DataTypeTok{void}\NormalTok{ tracker\_mqtt\_event\_handler}\OperatorTok{(...,} \DataTypeTok{int32\_t}\NormalTok{ event\_id}\OperatorTok{,} \DataTypeTok{void} \OperatorTok{*}\NormalTok{event\_data}\OperatorTok{)} \OperatorTok{\{}
    \ControlFlowTok{switch} \OperatorTok{((}\NormalTok{esp\_mqtt\_event\_id\_t}\OperatorTok{)}\NormalTok{event\_id}\OperatorTok{)} \OperatorTok{\{}
        \ControlFlowTok{case}\NormalTok{ MQTT\_EVENT\_CONNECTED}\OperatorTok{:}
\NormalTok{            s\_connected }\OperatorTok{=} \KeywordTok{true}\OperatorTok{;}
            \ControlFlowTok{if} \OperatorTok{(}\NormalTok{s\_cfg}\OperatorTok{.}\NormalTok{command\_subscribe\_enabled}\OperatorTok{)} \OperatorTok{\{}
\NormalTok{                esp\_mqtt\_client\_subscribe}\OperatorTok{(}\NormalTok{s\_client}\OperatorTok{,}\NormalTok{ s\_topic\_commands}\OperatorTok{,} \DecValTok{1}\OperatorTok{);}
            \OperatorTok{\}}
            \ControlFlowTok{break}\OperatorTok{;}
        \ControlFlowTok{case}\NormalTok{ MQTT\_EVENT\_DATA}\OperatorTok{:}
            \ControlFlowTok{if} \OperatorTok{(}\NormalTok{s\_command\_callback }\OperatorTok{!=}\NormalTok{ NULL}\OperatorTok{)} \OperatorTok{\{}
\NormalTok{                s\_command\_callback}\OperatorTok{(}\NormalTok{s\_topic\_commands}\OperatorTok{,}\NormalTok{ payload}\OperatorTok{);}
            \OperatorTok{\}}
            \ControlFlowTok{break}\OperatorTok{;}
        \ControlFlowTok{default}\OperatorTok{:}
            \ControlFlowTok{break}\OperatorTok{;}
    \OperatorTok{\}}
\OperatorTok{\}}
\end{Highlighting}
\end{Shaded}

Trong đó, \texttt{rawdata} được publish với QoS 0 để tối ưu băng thông
cho luồng dữ liệu vận hành gửi dày; còn \texttt{status}, \texttt{events},
\texttt{firmware} và \texttt{commands} dùng QoS 1 để tăng độ tin cậy cho
heartbeat, cảnh báo, OTA và lệnh điều khiển.

\textbf{Cấu trúc topic MQTT:}


{\thesissettableid{4.8A}{tab-ch4-bang-bo-sung-chi-muc-4-2-3-trien-khai-firmware}
\def\LTcaptype{table}
\begin{longtable}{|L{0.332\linewidth}|L{0.231\linewidth}|L{0.401\linewidth}|}
\caption{Bảng bổ sung chỉ mục (4.2.3. Triển khai Firmware)}\label{tab:ch4-bang-bo-sung-chi-muc-4-2-3-trien-khai-firmware}\\
\hline \textbf{Topic} & \textbf{Hướng} & \textbf{Mục đích} \\
\\\hline
\endfirsthead
\caption[]{Bảng bổ sung chỉ mục (4.2.3. Triển khai Firmware) (tiếp theo)}\\
\hline \textbf{Topic} & \textbf{Hướng} & \textbf{Mục đích} \\
\\\hline
\endhead
\\\hline
\endfoot
\endlastfoot
\texttt{v1/\{device\_id\}/rawdata} & Device → Server &
Dữ liệu vận hành chính: GNSS, nguồn, rung \\\hline
\texttt{v1/\{device\_id\}/events} & Device → Server & Sự
kiện và cảnh báo quan trọng \\\hline
\texttt{v1/\{device\_id\}/status} & Device → Server &
Heartbeat và trạng thái phiên \\\hline
\texttt{v1/\{device\_id\}/firmware} & Device → Server &
Tiến trình OTA \\\hline
\texttt{v1/\{device\_id\}/commands} & Server → Device &
Lệnh điều khiển từ xa \\\hline
\end{longtable}
}


\paragraph{c) Đọc dữ liệu
GNSS}\label{c-ux111ux1ecdc-dux1eef-liux1ec7u-gnss}

GNSS tích hợp trong SIM7600CE-T được điều khiển qua cùng UART. Firmware
bật GNSS, đợi fix vị trí, parse bản tin NMEA/CGNSINF và đọc tọa độ định
kỳ.

\begin{Shaded}
\begin{Highlighting}[]
\CommentTok{// Bật và đọc dữ liệu GNSS từ modem}
\DataTypeTok{int}\NormalTok{ modem\_gnss\_read}\OperatorTok{(}\NormalTok{modem\_ctx\_t }\OperatorTok{*}\NormalTok{ctx}\OperatorTok{,}\NormalTok{ gnss\_data\_t }\OperatorTok{*}\NormalTok{gnss}\OperatorTok{)}
\OperatorTok{\{}
    \CommentTok{// Bật nguồn GNSS (nếu chưa bật)}
\NormalTok{    modem\_send\_at}\OperatorTok{(}\NormalTok{ctx}\OperatorTok{,} \StringTok{"AT+CGNSPWR=1"}\OperatorTok{,} \StringTok{"OK"}\OperatorTok{,} \DecValTok{3000}\OperatorTok{);}

    \CommentTok{// Đọc thông tin GNSS}
    \CommentTok{// Phản hồi: +CGNSINF: \textless{}mode\textgreater{},\textless{}sat\_gps\textgreater{},\textless{}sat\_glo\textgreater{},\textless{}sat\_bds\textgreater{},}
    \CommentTok{//           \textless{}lat\textgreater{},\textless{}N/S\textgreater{},\textless{}lon\textgreater{},\textless{}E/W\textgreater{},\textless{}date\textgreater{},\textless{}time\textgreater{},\textless{}alt\textgreater{},}
    \CommentTok{//           \textless{}speed\textgreater{},\textless{}course\textgreater{},\textless{}hdop\textgreater{},\textless{}vdop\textgreater{}}
    \DataTypeTok{char}\NormalTok{ response}\OperatorTok{[}\DecValTok{256}\OperatorTok{];}
    \ControlFlowTok{if} \OperatorTok{(}\NormalTok{modem\_send\_at\_get\_response}\OperatorTok{(}\NormalTok{ctx}\OperatorTok{,} \StringTok{"AT+CGNSINF"}\OperatorTok{,}
\NormalTok{                                    response}\OperatorTok{,} \KeywordTok{sizeof}\OperatorTok{(}\NormalTok{response}\OperatorTok{),}
                                    \DecValTok{3000}\OperatorTok{)} \OperatorTok{==} \DecValTok{0}\OperatorTok{)} \OperatorTok{\{}
        \CommentTok{// Phân tích phản hồi}
\NormalTok{        gnss\_parse\_cgnsinf}\OperatorTok{(}\NormalTok{response}\OperatorTok{,}\NormalTok{ gnss}\OperatorTok{);}
        \ControlFlowTok{return} \DecValTok{0}\OperatorTok{;}
    \OperatorTok{\}}
    \ControlFlowTok{return} \OperatorTok{{-}}\DecValTok{1}\OperatorTok{;}  \CommentTok{// Chưa có fix}
\OperatorTok{\}}

\CommentTok{// Cấu trúc dữ liệu GNSS}
\KeywordTok{typedef} \KeywordTok{struct} \OperatorTok{\{}
    \DataTypeTok{double}\NormalTok{   latitude}\OperatorTok{;}      \CommentTok{// Vĩ độ (độ thập phân)}
    \DataTypeTok{double}\NormalTok{   longitude}\OperatorTok{;}     \CommentTok{// Kinh độ (độ thập phân)}
    \DataTypeTok{float}\NormalTok{    altitude}\OperatorTok{;}      \CommentTok{// Độ cao (mét)}
    \DataTypeTok{float}\NormalTok{    speed}\OperatorTok{;}         \CommentTok{// Tốc độ (km/h)}
    \DataTypeTok{float}\NormalTok{    course}\OperatorTok{;}        \CommentTok{// Hướng đi (độ)}
    \DataTypeTok{uint8\_t}\NormalTok{  satellites}\OperatorTok{;}    \CommentTok{// Số vệ tinh}
    \DataTypeTok{bool}\NormalTok{     fix\_valid}\OperatorTok{;}     \CommentTok{// Trạng thái fix hợp lệ}
\OperatorTok{\}}\NormalTok{ gnss\_data\_t}\OperatorTok{;}
\end{Highlighting}
\end{Shaded}

Thời gian fix GNSS phụ thuộc vào trạng thái trước đó của modem: hot
start khoảng 1--3 giây khi module chỉ mất tín hiệu trong thời gian ngắn,
warm start khoảng 5--15 giây khi còn dữ liệu ephemeris, và cold start
khoảng 30--60 giây khi khởi tạo lại hoàn toàn.

\paragraph{d) Xử lý lệnh điều khiển từ
xa}\label{d-xux1eed-luxfd-lux1ec7nh-ux111iux1ec1u-khiux1ec3n-tux1eeb-xa}

Firmware lắng nghe lệnh từ máy chủ thông qua topic
\texttt{v1/\{device\_id\}/commands}. Gói lệnh được
\texttt{command\_handler.c} phân tích theo cấu trúc JSON, sau đó cập
nhật cấu hình hoặc đặt cờ hành động để \texttt{state\_machine.c} xử lý ở
chu kỳ kế tiếp.

\begin{Shaded}
\begin{Highlighting}[]
\DataTypeTok{void}\NormalTok{ command\_handler\_process}\OperatorTok{(}\DataTypeTok{const} \DataTypeTok{char} \OperatorTok{*}\NormalTok{command\_json}\OperatorTok{)} \OperatorTok{\{}
\NormalTok{    cJSON }\OperatorTok{*}\NormalTok{root }\OperatorTok{=}\NormalTok{ cJSON\_Parse}\OperatorTok{(}\NormalTok{command\_json}\OperatorTok{);}
    \DataTypeTok{const}\NormalTok{ cJSON }\OperatorTok{*}\NormalTok{command }\OperatorTok{=}\NormalTok{ cJSON\_GetObjectItemCaseSensitive}\OperatorTok{(}\NormalTok{root}\OperatorTok{,} \StringTok{"command"}\OperatorTok{);}
    \DataTypeTok{const}\NormalTok{ cJSON }\OperatorTok{*}\NormalTok{params }\OperatorTok{=}\NormalTok{ cJSON\_GetObjectItemCaseSensitive}\OperatorTok{(}\NormalTok{root}\OperatorTok{,} \StringTok{"params"}\OperatorTok{);}

    \ControlFlowTok{if} \OperatorTok{(}\NormalTok{strcmp}\OperatorTok{(}\NormalTok{command}\OperatorTok{{-}\textgreater{}}\NormalTok{valuestring}\OperatorTok{,} \StringTok{"update\_config"}\OperatorTok{)} \OperatorTok{==} \DecValTok{0}\OperatorTok{)} \OperatorTok{\{}
\NormalTok{        command\_apply\_update\_config}\OperatorTok{(}\NormalTok{params}\OperatorTok{);}
    \OperatorTok{\}} \ControlFlowTok{else} \ControlFlowTok{if} \OperatorTok{(}\NormalTok{strcmp}\OperatorTok{(}\NormalTok{command}\OperatorTok{{-}\textgreater{}}\NormalTok{valuestring}\OperatorTok{,} \StringTok{"request\_location"}\OperatorTok{)} \OperatorTok{==} \DecValTok{0}\OperatorTok{)} \OperatorTok{\{}
\NormalTok{        s\_location\_requested }\OperatorTok{=} \KeywordTok{true}\OperatorTok{;}
\NormalTok{        s\_pending\_action }\OperatorTok{=}\NormalTok{ COMMAND\_ACTION\_REQUEST\_LOCATION}\OperatorTok{;}
    \OperatorTok{\}} \ControlFlowTok{else} \ControlFlowTok{if} \OperatorTok{(}\NormalTok{strcmp}\OperatorTok{(}\NormalTok{command}\OperatorTok{{-}\textgreater{}}\NormalTok{valuestring}\OperatorTok{,} \StringTok{"enable\_tracking"}\OperatorTok{)} \OperatorTok{==} \DecValTok{0}\OperatorTok{)} \OperatorTok{\{}
        \DataTypeTok{const}\NormalTok{ cJSON }\OperatorTok{*}\NormalTok{enabled }\OperatorTok{=}\NormalTok{ cJSON\_GetObjectItemCaseSensitive}\OperatorTok{(}\NormalTok{params}\OperatorTok{,} \StringTok{"enabled"}\OperatorTok{);}
        \ControlFlowTok{if} \OperatorTok{(}\NormalTok{cJSON\_IsBool}\OperatorTok{(}\NormalTok{enabled}\OperatorTok{))} \OperatorTok{\{}
\NormalTok{            s\_tracking\_enabled }\OperatorTok{=}\NormalTok{ cJSON\_IsTrue}\OperatorTok{(}\NormalTok{enabled}\OperatorTok{);}
        \OperatorTok{\}}
    \OperatorTok{\}} \ControlFlowTok{else} \ControlFlowTok{if} \OperatorTok{(}\NormalTok{strcmp}\OperatorTok{(}\NormalTok{command}\OperatorTok{{-}\textgreater{}}\NormalTok{valuestring}\OperatorTok{,} \StringTok{"reboot"}\OperatorTok{)} \OperatorTok{==} \DecValTok{0}\OperatorTok{)} \OperatorTok{\{}
\NormalTok{        s\_pending\_action }\OperatorTok{=}\NormalTok{ COMMAND\_ACTION\_REBOOT}\OperatorTok{;}
    \OperatorTok{\}} \ControlFlowTok{else} \ControlFlowTok{if} \OperatorTok{(}\NormalTok{strcmp}\OperatorTok{(}\NormalTok{command}\OperatorTok{{-}\textgreater{}}\NormalTok{valuestring}\OperatorTok{,} \StringTok{"ota\_update"}\OperatorTok{)} \OperatorTok{==} \DecValTok{0}\OperatorTok{)} \OperatorTok{\{}
\NormalTok{        s\_pending\_action }\OperatorTok{=}\NormalTok{ COMMAND\_ACTION\_OTA\_UPDATE}\OperatorTok{;}
    \OperatorTok{\}} \ControlFlowTok{else} \ControlFlowTok{if} \OperatorTok{(}\NormalTok{strcmp}\OperatorTok{(}\NormalTok{command}\OperatorTok{{-}\textgreater{}}\NormalTok{valuestring}\OperatorTok{,} \StringTok{"manual\_rollback"}\OperatorTok{)} \OperatorTok{==} \DecValTok{0} \OperatorTok{||}
\NormalTok{               strcmp}\OperatorTok{(}\NormalTok{command}\OperatorTok{{-}\textgreater{}}\NormalTok{valuestring}\OperatorTok{,} \StringTok{"ota\_rollback"}\OperatorTok{)} \OperatorTok{==} \DecValTok{0}\OperatorTok{)} \OperatorTok{\{}
\NormalTok{        s\_pending\_action }\OperatorTok{=}\NormalTok{ COMMAND\_ACTION\_OTA\_ROLLBACK}\OperatorTok{;}
    \OperatorTok{\}}
\OperatorTok{\}}
\end{Highlighting}
\end{Shaded}

\paragraph{e) Luồng cập nhật firmware
OTA}\label{e-luux1ed3ng-cux1eadp-nhux1eadt-firmware-ota}

Tính năng OTA mới được tích hợp xuyên suốt từ API phía máy chủ đến firmware.
API phía máy chủ tạo \texttt{jobId}, tra cứu firmware theo
\texttt{firmwareVersion}, sinh URL tải file dạng
\texttt{/api/v1/firmware/\{id\}/download}, rồi gửi lệnh
\texttt{ota\_update} qua topic \texttt{v1/\{device\_id\}/commands}.
Firmware kiểm tra đầy đủ các tham số \texttt{jobId}, \texttt{version},
\texttt{url}, \texttt{size}, \texttt{sha256}, \texttt{force} và
\texttt{confirmTimeoutSec} trước khi bắt đầu cập nhật.

\begin{Shaded}
\begin{Highlighting}[]
\NormalTok{state\_machine\_publish\_firmware\_status}\OperatorTok{(}\StringTok{"assigned"}\OperatorTok{,} \DecValTok{0}\OperatorTok{,}\NormalTok{ cmd}\OperatorTok{.}\NormalTok{version}\OperatorTok{,}\NormalTok{ cmd}\OperatorTok{.}\NormalTok{job\_id}\OperatorTok{,} \StringTok{""}\OperatorTok{,} \StringTok{""}\OperatorTok{);}

\ControlFlowTok{if} \OperatorTok{(}\NormalTok{util\_ota\_apply\_update}\OperatorTok{(\&}\NormalTok{s\_config}\OperatorTok{,}\NormalTok{ s\_current\_version}\OperatorTok{,} \OperatorTok{\&}\NormalTok{cmd}\OperatorTok{,} \OperatorTok{\&}\NormalTok{report}\OperatorTok{)} \OperatorTok{==}\NormalTok{ ESP\_OK}\OperatorTok{)} \OperatorTok{\{}
\NormalTok{    util\_copy\_string}\OperatorTok{(}\NormalTok{g\_rtc\_context}\OperatorTok{.}\NormalTok{ota\_job\_id}\OperatorTok{,} \KeywordTok{sizeof}\OperatorTok{(}\NormalTok{g\_rtc\_context}\OperatorTok{.}\NormalTok{ota\_job\_id}\OperatorTok{),}\NormalTok{ cmd}\OperatorTok{.}\NormalTok{job\_id}\OperatorTok{);}
\NormalTok{    util\_copy\_string}\OperatorTok{(}\NormalTok{g\_rtc\_context}\OperatorTok{.}\NormalTok{ota\_target\_version}\OperatorTok{,} \KeywordTok{sizeof}\OperatorTok{(}\NormalTok{g\_rtc\_context}\OperatorTok{.}\NormalTok{ota\_target\_version}\OperatorTok{),}\NormalTok{ cmd}\OperatorTok{.}\NormalTok{version}\OperatorTok{);}
\NormalTok{    util\_copy\_string}\OperatorTok{(}\NormalTok{g\_rtc\_context}\OperatorTok{.}\NormalTok{ota\_previous\_version}\OperatorTok{,} \KeywordTok{sizeof}\OperatorTok{(}\NormalTok{g\_rtc\_context}\OperatorTok{.}\NormalTok{ota\_previous\_version}\OperatorTok{),}\NormalTok{ s\_current\_version}\OperatorTok{);}
\NormalTok{    g\_rtc\_context}\OperatorTok{.}\NormalTok{ota\_pending\_confirm }\OperatorTok{=} \KeywordTok{true}\OperatorTok{;}
\NormalTok{    g\_rtc\_context}\OperatorTok{.}\NormalTok{ota\_confirm\_timeout\_sec }\OperatorTok{=}\NormalTok{ cmd}\OperatorTok{.}\NormalTok{confirm\_timeout\_sec}\OperatorTok{;}

\NormalTok{    state\_machine\_publish\_firmware\_payload}\OperatorTok{(\&}\NormalTok{report}\OperatorTok{);}
\NormalTok{    esp\_restart}\OperatorTok{();}
\OperatorTok{\}}
\end{Highlighting}
\end{Shaded}

Hàm \texttt{util\_ota\_apply\_update()} sử dụng HTTPS để tải file
firmware, ghi trực tiếp xuống OTA partition, tính và đối chiếu SHA-256,
sau đó đặt phân vùng boot mới bằng
\texttt{esp\_ota\_set\_boot\_partition()}. Sau lần khởi động kế tiếp,
firmware phát trạng thái \texttt{confirming}, gọi
\texttt{esp\_ota\_mark\_app\_valid\_cancel\_rollback()} để xác nhận
image mới, rồi mới phát \texttt{success}. Nếu có lỗi trong quá trình xác
nhận hoặc người vận hành gửi lệnh \texttt{manual\_rollback}, hệ thống sẽ
quay về phân vùng trước đó và phát trạng thái \texttt{failed} hoặc
\texttt{rolled\_back} tương ứng.


{\thesissetfigureid{4.16a}{fig-ch4-chuoi-buoc-ota-tren-firmware-tu-nhan-lenh-en-xac-nhan-sau}
\begin{figure}[htbp]
\centering
\pandocbounded{\safeincludesvg{./assets/figures/08-chuong-4-trien-khai-firmware-hinh-4-16a.svg}}
\caption{Chuỗi bước OTA trên firmware từ nhận lệnh đến xác nhận sau reboot}
\label{fig:ch4-chuoi-buoc-ota-tren-firmware-tu-nhan-lenh-en-xac-nhan-sau}
\end{figure}
}


\begin{center}\rule{0.5\linewidth}{0.5pt}\end{center}

\paragraph{4.2.3.5. Triển khai quản lý nguồn và chế độ
ngủ}\label{4235-triux1ec3n-khai-quux1ea3n-luxfd-nguux1ed3n-vuxe0-chux1ebf-ux111ux1ed9-ngux1ee7}

\paragraph{a) Đọc điện áp và điều khiển
nguồn}\label{a-ux111ux1ecdc-ux111iux1ec7n-uxe1p-vuxe0-ux111iux1ec1u-khiux1ec3n-nguux1ed3n}

Module quản lý nguồn thực hiện đọc điện áp ắc quy qua ADC, điều khiển
nhánh nguồn theo kiến trúc diode OR + EN, và điều khiển khối sạc TP5100.

\begin{Shaded}
\begin{Highlighting}[]
\CommentTok{// Đọc điện áp ắc quy qua ADC (12{-}bit, chia áp)}
\DataTypeTok{float}\NormalTok{ read\_battery\_voltage}\OperatorTok{(}\DataTypeTok{void}\OperatorTok{)}
\OperatorTok{\{}
    \DataTypeTok{int}\NormalTok{ adc\_raw }\OperatorTok{=}\NormalTok{ adc1\_get\_raw}\OperatorTok{(}\NormalTok{ADC1\_CHANNEL\_4}\OperatorTok{);}  \CommentTok{// GPIO4}
    \DataTypeTok{float}\NormalTok{ voltage }\OperatorTok{=} \OperatorTok{(}\NormalTok{adc\_raw }\OperatorTok{/} \FloatTok{4095.0}\BuiltInTok{f}\OperatorTok{)} \OperatorTok{*} \FloatTok{3.3}\BuiltInTok{f}\OperatorTok{;}  \CommentTok{// Chuyển sang volt}
\NormalTok{    voltage }\OperatorTok{*=}\NormalTok{ VOLTAGE\_DIVIDER\_RATIO}\OperatorTok{;}              \CommentTok{// Bù tỷ lệ chia áp}
    \ControlFlowTok{return}\NormalTok{ voltage}\OperatorTok{;}
\OperatorTok{\}}

\CommentTok{// Logic quản lý nguồn chính}
\DataTypeTok{void}\NormalTok{ power\_management\_task}\OperatorTok{(}\DataTypeTok{void} \OperatorTok{*}\NormalTok{param}\OperatorTok{)}
\OperatorTok{\{}
    \ControlFlowTok{while} \OperatorTok{(}\KeywordTok{true}\OperatorTok{)} \OperatorTok{\{}
        \DataTypeTok{float}\NormalTok{ u\_batt }\OperatorTok{=}\NormalTok{ read\_battery\_voltage}\OperatorTok{();}
        \DataTypeTok{bool}\NormalTok{ ign\_on  }\OperatorTok{=}\NormalTok{ get\_ignition\_status}\OperatorTok{();}
        \DataTypeTok{bool}\NormalTok{ lvd\_ok  }\OperatorTok{=}\NormalTok{ gpio\_get\_level}\OperatorTok{(}\NormalTok{LVD\_STATUS}\OperatorTok{);}
        \DataTypeTok{float}\NormalTok{ switch\_off }\OperatorTok{=}\NormalTok{ cfg}\OperatorTok{.}\NormalTok{power\_profile\_24v }\OperatorTok{?} \FloatTok{24.0}\BuiltInTok{f} \OperatorTok{:} \FloatTok{12.0}\BuiltInTok{f}\OperatorTok{;}
        \DataTypeTok{float}\NormalTok{ switch\_on  }\OperatorTok{=}\NormalTok{ cfg}\OperatorTok{.}\NormalTok{power\_profile\_24v }\OperatorTok{?} \FloatTok{24.4}\BuiltInTok{f} \OperatorTok{:} \FloatTok{12.2}\BuiltInTok{f}\OperatorTok{;}

        \ControlFlowTok{if} \OperatorTok{(}\NormalTok{ign\_on}\OperatorTok{)} \OperatorTok{\{}
            \CommentTok{// IGN ON: dùng nguồn ắc quy, bật sạc pin dự phòng}
\NormalTok{            gpio\_set\_level}\OperatorTok{(}\NormalTok{POWER\_PATH\_EN}\OperatorTok{,} \DecValTok{0}\OperatorTok{);}   \CommentTok{// Chọn ắc quy}
\NormalTok{            gpio\_set\_level}\OperatorTok{(}\NormalTok{CHARGER\_EN}\OperatorTok{,} \DecValTok{1}\OperatorTok{);}       \CommentTok{// Bật sạc}
        \OperatorTok{\}} \ControlFlowTok{else} \ControlFlowTok{if} \OperatorTok{(}\NormalTok{u\_batt }\OperatorTok{\textless{}=}\NormalTok{ switch\_off}\OperatorTok{)} \OperatorTok{\{}
            \CommentTok{// Điện áp thấp: chuyển sang pin dự phòng theo profile}
\NormalTok{            gpio\_set\_level}\OperatorTok{(}\NormalTok{POWER\_PATH\_EN}\OperatorTok{,} \DecValTok{1}\OperatorTok{);}   \CommentTok{// Chọn pin dự phòng}
\NormalTok{            gpio\_set\_level}\OperatorTok{(}\NormalTok{CHARGER\_EN}\OperatorTok{,} \DecValTok{0}\OperatorTok{);}       \CommentTok{// Tắt sạc}
\NormalTok{            send\_low\_battery\_alert}\OperatorTok{(}\NormalTok{u\_batt}\OperatorTok{);}
        \OperatorTok{\}} \ControlFlowTok{else} \ControlFlowTok{if} \OperatorTok{(}\NormalTok{u\_batt }\OperatorTok{\textgreater{}=}\NormalTok{ switch\_on}\OperatorTok{)} \OperatorTok{\{}
            \CommentTok{// Điện áp phục hồi: quay lại ắc quy theo profile}
\NormalTok{            gpio\_set\_level}\OperatorTok{(}\NormalTok{POWER\_PATH\_EN}\OperatorTok{,} \DecValTok{0}\OperatorTok{);}   \CommentTok{// Chọn ắc quy}
\NormalTok{            gpio\_set\_level}\OperatorTok{(}\NormalTok{CHARGER\_EN}\OperatorTok{,} \DecValTok{0}\OperatorTok{);}       \CommentTok{// Tắt sạc (IGN OFF)}
        \OperatorTok{\}}
        \CommentTok{// Giữ nguyên trạng thái nếu trong vùng trễ của profile (12V: 12.0{-}12.2V; 24V: 24.0{-}24.4V)}

\NormalTok{        vTaskDelay}\OperatorTok{(}\NormalTok{pdMS\_TO\_TICKS}\OperatorTok{(}\DecValTok{5000}\OperatorTok{));}  \CommentTok{// Kiểm tra mỗi 5 giây}
    \OperatorTok{\}}
\OperatorTok{\}}
\end{Highlighting}
\end{Shaded}

Ngưỡng điện áp sử dụng cơ chế trễ (hysteresis) theo profile cấu hình để
tránh hiện tượng chuyển đổi liên tục khi điện áp dao động quanh ngưỡng:
profile 12V dùng \(12.0\,\mathrm{V}\) (OFF) và \(12.2\,\mathrm{V}\) (ON),
profile 24V dùng \(24.0\,\mathrm{V}\) (OFF) và \(24.4\,\mathrm{V}\) (ON).
Ngưỡng suy luận trạng thái động cơ cũng được tách profile:
\(IGN_{\mathrm{ON}} \ge 13.0\,\mathrm{V}\) và
\(IGN_{\mathrm{OFF}} \le 12.0\,\mathrm{V}\) (12V), hoặc
\(IGN_{\mathrm{ON}} \ge 26.0\,\mathrm{V}\) và
\(IGN_{\mathrm{OFF}} \le 24.0\,\mathrm{V}\) (24V).

\paragraph{b) Chế độ ngủ và đánh
thức}\label{b-chux1ebf-ux111ux1ed9-ngux1ee7-vuxe0-ux111uxe1nh-thux1ee9c}

Firmware sử dụng hai chế độ ngủ của ESP32-S3 tùy theo tình huống:


{\thesissettableid{4.9}{tab-ch4-so-sanh-cac-che-o-ngu}
\def\LTcaptype{table}
\begin{longtable}{|L{0.162\linewidth}|L{0.181\linewidth}|L{0.225\linewidth}|L{0.128\linewidth}|L{0.249\linewidth}|}
\caption{So sánh các chế độ ngủ}\label{tab:ch4-so-sanh-cac-che-o-ngu}\\
\hline \textbf{Chế độ} & \textbf{Dòng tiêu thụ} & \textbf{Thời gian thức dậy} & \textbf{BLE} & \textbf{Điều kiện sử dụng} \\
\\\hline
\endfirsthead
\caption[]{So sánh các chế độ ngủ (tiếp theo)}\\
\hline \textbf{Chế độ} & \textbf{Dòng tiêu thụ} & \textbf{Thời gian thức dậy} & \textbf{BLE} & \textbf{Điều kiện sử dụng} \\
\\\hline
\endhead
\\\hline
\endfoot
\endlastfoot
Light sleep & ~0.8 mA & \textless{} 1 ms & Giữ active &
IGN ON, đợi dữ liệu giữa các chu kỳ \\\hline
Deep sleep & ~10 µA & ~100 ms & Mất kết nối
& IGN OFF, chế độ đỗ xe \\\hline
\end{longtable}
}


Trước khi vào trạng thái ngủ đỗ xe, firmware thực hiện trình tự tắt các ngoại vi
để tiết kiệm năng lượng tối đa:

\begin{Shaded}
\begin{Highlighting}[]
\CommentTok{// Trình tự trước khi vào parked sleep}
\DataTypeTok{void}\NormalTok{ enter\_parked\_sleep}\OperatorTok{(}\DataTypeTok{uint32\_t}\NormalTok{ sleep\_duration\_sec}\OperatorTok{)}
\OperatorTok{\{}
    \CommentTok{// 1. Ngắt kết nối BLE (nếu đang kết nối)}
\NormalTok{    ble\_obd\_disconnect}\OperatorTok{();}

    \CommentTok{// 2. Tắt GNSS}
\NormalTok{    modem\_send\_at}\OperatorTok{(\&}\NormalTok{modem\_ctx}\OperatorTok{,} \StringTok{"AT+CGNSPWR=0"}\OperatorTok{,} \StringTok{"OK"}\OperatorTok{,} \DecValTok{3000}\OperatorTok{);}

    \CommentTok{// 3. Tắt 4G và đưa modem vào chế độ ngủ}
\NormalTok{    modem\_send\_at}\OperatorTok{(\&}\NormalTok{modem\_ctx}\OperatorTok{,} \StringTok{"AT+CGACT=0,1"}\OperatorTok{,} \StringTok{"OK"}\OperatorTok{,} \DecValTok{5000}\OperatorTok{);}
\NormalTok{    modem\_send\_at}\OperatorTok{(\&}\NormalTok{modem\_ctx}\OperatorTok{,} \StringTok{"AT+CSCLK=1"}\OperatorTok{,} \StringTok{"OK"}\OperatorTok{,} \DecValTok{3000}\OperatorTok{);}

    \CommentTok{// 4. Tắt sạc pin}
\NormalTok{    gpio\_set\_level}\OperatorTok{(}\NormalTok{CHARGER\_EN}\OperatorTok{,} \DecValTok{0}\OperatorTok{);}

    \CommentTok{// 5. Lưu trạng thái hiện tại vào RTC memory}
\NormalTok{    rtc\_data}\OperatorTok{.}\NormalTok{last\_mode       }\OperatorTok{=}\NormalTok{ current\_mode}\OperatorTok{;}
\NormalTok{    rtc\_data}\OperatorTok{.}\NormalTok{last\_ign\_status }\OperatorTok{=}\NormalTok{ ign\_on}\OperatorTok{;}
\NormalTok{    rtc\_data}\OperatorTok{.}\NormalTok{sleep\_count}\OperatorTok{++;}

    \CommentTok{// 6. Cấu hình nguồn đánh thức}
    \CommentTok{//    {-} Timer: heartbeat định kỳ}
\NormalTok{    esp\_sleep\_enable\_timer\_wakeup}\OperatorTok{(}\NormalTok{sleep\_duration\_sec }\OperatorTok{*} \DecValTok{1000000}\BuiltInTok{ULL}\OperatorTok{);}
    \CommentTok{//    {-} IMU interrupt trên GPIO41: dùng GPIO wake ở light sleep}
\NormalTok{    gpio\_wakeup\_enable}\OperatorTok{(}\NormalTok{PIN\_LIS3DH\_INT}\OperatorTok{,}\NormalTok{ GPIO\_INTR\_HIGH\_LEVEL}\OperatorTok{);}
\NormalTok{    esp\_sleep\_enable\_gpio\_wakeup}\OperatorTok{();}
    \CommentTok{// 7. Vào light sleep}
\NormalTok{    esp\_light\_sleep\_start}\OperatorTok{();}
\OperatorTok{\}}
\end{Highlighting}
\end{Shaded}

Khi thức dậy từ parked sleep, firmware đọc nguyên nhân đánh thức và trạng
thái đã lưu trong RTC memory để quyết định chế độ hoạt động tiếp theo mà
không cần khởi tạo lại toàn bộ.

\paragraph{c) Quản lý nguồn modem theo chế
độ}\label{c-quux1ea3n-luxfd-nguux1ed3n-modem-theo-chux1ebf-ux111ux1ed9}

Modem SIM7600CE-T hỗ trợ nhiều chế độ ngủ với mức tiêu thụ khác nhau:

\begin{itemize}
\tightlist
\item
  \textbf{UART sleep} (\texttt{AT+CSCLK=1}): Modem tự động ngủ khi không
  có dữ liệu UART, tiêu thụ 1--5 mA, đánh thức bằng bất kỳ ký tự UART
  nào.
\item
  \textbf{Minimum functionality} (\texttt{AT+CFUN=0}): Tắt RF, giữ UART,
  tiêu thụ \textless{} 1 mA, đánh thức bằng lệnh AT hoặc GPIO.
\item
  \textbf{Flight mode} (\texttt{AT+CFUN=4}): Tắt RF nhưng giữ chức năng
  khác, phù hợp cho heartbeat ngắn.
\end{itemize}

Thời gian đánh thức modem phụ thuộc chế độ: từ sleep là 100--500 ms, từ
chế độ tối thiểu là 1--3 giây, và kết nối lại 4G mất thêm 5--15
giây.

\begin{center}\rule{0.5\linewidth}{0.5pt}\end{center}

\paragraph{4.2.3.6. Lưu đồ thuật toán
chính}\label{4236-lux1b0u-ux111ux1ed3-thuux1eadt-touxe1n-chuxednh}

\paragraph{a) Máy trạng thái}\label{a-muxe1y-trux1ea1ng-thuxe1i-state-machine}

Toàn bộ logic điều khiển firmware được tổ chức theo mô hình máy trạng
thái, tức bộ điều phối trạng thái hữu hạn, gồm bảy trạng thái chính:


{\thesissettableid{4.10}{tab-ch4-cac-trang-thai-cua-firmware}
\def\LTcaptype{table}
\begin{longtable}{|L{0.235\linewidth}|L{0.277\linewidth}|L{0.452\linewidth}|}
\caption{Các trạng thái của firmware}\label{tab:ch4-cac-trang-thai-cua-firmware}\\
\hline \textbf{Trạng thái} & \textbf{Mã} & \textbf{Mô tả} \\
\\\hline
\endfirsthead
\caption[]{Các trạng thái của firmware (tiếp theo)}\\
\hline \textbf{Trạng thái} & \textbf{Mã} & \textbf{Mô tả} \\
\\\hline
\endhead
\\\hline
\endfoot
\endlastfoot
INIT & \texttt{STATE\_INIT} & Khởi tạo hệ thống, ngoại vi \\\hline
CHECK\_IGN & \texttt{STATE\_CHECK\_IGN} & Kiểm tra trạng thái khóa
điện \\\hline
DRIVING & \texttt{STATE\_DRIVING} & Chế độ lái xe - theo dõi liên tục \\\hline
PARKED & \texttt{STATE\_PARKED} & Chế độ đỗ xe - heartbeat định kỳ \\\hline
ALARM & \texttt{STATE\_ALARM} & Chế độ cảnh báo - phát hiện chuyển
động \\\hline
HEARTBEAT & \texttt{STATE\_HEARTBEAT} & Gửi tín hiệu heartbeat \\\hline
SLEEP & \texttt{STATE\_SLEEP} & Deep sleep tiết kiệm năng lượng \\\hline
\end{longtable}
}



{\thesissettableid{4.11}{tab-ch4-chuyen-oi-trang-thai}
\def\LTcaptype{table}
\begin{longtable}{|L{0.303\linewidth}|L{0.391\linewidth}|L{0.271\linewidth}|}
\caption{Chuyển đổi trạng thái}\label{tab:ch4-chuyen-oi-trang-thai}\\
\hline \textbf{Trạng thái hiện tại} & \textbf{Điều kiện} & \textbf{Trạng thái mới} \\
\\\hline
\endfirsthead
\caption[]{Chuyển đổi trạng thái (tiếp theo)}\\
\hline \textbf{Trạng thái hiện tại} & \textbf{Điều kiện} & \textbf{Trạng thái mới} \\
\\\hline
\endhead
\\\hline
\endfoot
\endlastfoot
INIT & Khởi tạo xong & CHECK\_IGN \\\hline
CHECK\_IGN & IGN = ON & DRIVING \\\hline
CHECK\_IGN & IGN = OFF & PARKED \\\hline
DRIVING & IGN = OFF (phát hiện) & PARKED \\\hline
PARKED & IMU interrupt (chuyển động) & ALARM \\\hline
PARKED & Timer wake-up & HEARTBEAT \\\hline
ALARM & Chuyển động dừng, IGN = OFF & PARKED \\\hline
HEARTBEAT & Gá»­i xong heartbeat & SLEEP \\\hline
SLEEP & Timer hoặc IMU wake-up & CHECK\_IGN \\\hline
\end{longtable}
}


\begin{Shaded}
\begin{Highlighting}[]
\KeywordTok{typedef} \KeywordTok{enum} \OperatorTok{\{}
\NormalTok{    APP\_STATE\_INIT }\OperatorTok{=} \DecValTok{0}\OperatorTok{,}
\NormalTok{    APP\_STATE\_CHECK\_IGN}\OperatorTok{,}
\NormalTok{    APP\_STATE\_DRIVING}\OperatorTok{,}
\NormalTok{    APP\_STATE\_PARKED}\OperatorTok{,}
\NormalTok{    APP\_STATE\_ALARM}\OperatorTok{,}
\NormalTok{    APP\_STATE\_HEARTBEAT}\OperatorTok{,}
\NormalTok{    APP\_STATE\_SLEEP}\OperatorTok{,}
\OperatorTok{\}}\NormalTok{ app\_state\_t}\OperatorTok{;}

\DataTypeTok{void}\NormalTok{ app\_main}\OperatorTok{(}\DataTypeTok{void}\OperatorTok{)} \OperatorTok{\{}
\NormalTok{    config\_t config }\OperatorTok{=} \OperatorTok{\{}\DecValTok{0}\OperatorTok{\};}
\NormalTok{    ESP\_ERROR\_CHECK}\OperatorTok{(}\NormalTok{nvs\_config\_load}\OperatorTok{(\&}\NormalTok{config}\OperatorTok{));}
\NormalTok{    ESP\_ERROR\_CHECK}\OperatorTok{(}\NormalTok{state\_machine\_init}\OperatorTok{(\&}\NormalTok{config}\OperatorTok{));}

\NormalTok{    app\_state\_t state }\OperatorTok{=}\NormalTok{ APP\_STATE\_INIT}\OperatorTok{;}
    \ControlFlowTok{while} \OperatorTok{(}\KeywordTok{true}\OperatorTok{)} \OperatorTok{\{}
\NormalTok{        state }\OperatorTok{=}\NormalTok{ state\_machine\_run}\OperatorTok{(}\NormalTok{state}\OperatorTok{);}
\NormalTok{        vTaskDelay}\OperatorTok{(}\NormalTok{pdMS\_TO\_TICKS}\OperatorTok{(}\DecValTok{100}\OperatorTok{));}
    \OperatorTok{\}}
\OperatorTok{\}}
\end{Highlighting}
\end{Shaded}

\paragraph{b) Lưu đồ thuật toán tổng
thể}\label{b-lux1b0u-ux111ux1ed3-thuux1eadt-touxe1n-tux1ed5ng-thux1ec3}


{\thesissetfigureid{4.17}{fig-ch4-luu-o-thuat-toan-chinh-cua-firmware-thiet-bi-theo-doi-xe}
\begin{figure}[htbp]
\centering
\pandocbounded{\safeincludesvg{./assets/figures/08-chuong-4-trien-khai-firmware-hinh-4-17.svg}}
\caption{Lưu đồ thuật toán chính của firmware thiết bị theo dõi xe}
\label{fig:ch4-luu-o-thuat-toan-chinh-cua-firmware-thiet-bi-theo-doi-xe}
\end{figure}
}


Lưu đồ được tổng hợp từ luồng xử lý trong mã firmware triển khai thực tế
và đối chiếu tài liệu kỹ thuật ESP32-S3 \parencite{ref:esp32s3-trm}, vì
vậy nó phản ánh cách thiết bị vận hành ngoài thực tế chứ không chỉ là sơ
đồ ý tưởng.

\paragraph{c) Định dạng dữ liệu
vận hành}\label{c-ux111ux1ecbnh-dux1ea1ng-dux1eef-liux1ec7u-telemetry}

Gói dữ liệu vận hành của thiết bị được tổ chức theo kiểu ``vỏ ngoài gọn,
ruột rõ ràng'': thông tin nhận dạng và thời điểm gửi nằm ở lớp ngoài,
còn các giá trị đo thực tế nằm trong trường \texttt{data}. Cách tổ chức
này giúp phía máy chủ nhìn nhanh được bản tin đến từ ai, gửi lúc nào và
cần đọc nhóm dữ liệu nào, từ đó MQTT Bridge kiểm tra và phân luồng dữ
liệu thuận tiện hơn.

\begin{Shaded}
\begin{Highlighting}[]
\FunctionTok{\{}
  \DataTypeTok{"device\_id"}\FunctionTok{:} \StringTok{"TRACKER\_001"}\FunctionTok{,}
  \DataTypeTok{"auth\_token"}\FunctionTok{:} \StringTok{"device\_secret\_token"}\FunctionTok{,}
  \DataTypeTok{"timestamp"}\FunctionTok{:} \DecValTok{1710000000000}\FunctionTok{,}
  \DataTypeTok{"uptime"}\FunctionTok{:} \DecValTok{452130}\FunctionTok{,}
  \DataTypeTok{"data"}\FunctionTok{:} \FunctionTok{\{}
    \DataTypeTok{"vibration"}\FunctionTok{:} \DecValTok{128}\FunctionTok{,}
    \DataTypeTok{"battery\_top"}\FunctionTok{:} \FloatTok{12.54}\FunctionTok{,}
    \DataTypeTok{"battery\_bot"}\FunctionTok{:} \FloatTok{3.96}\FunctionTok{,}
    \DataTypeTok{"latitude"}\FunctionTok{:} \FloatTok{10.77653}\FunctionTok{,}
    \DataTypeTok{"longitude"}\FunctionTok{:} \FloatTok{106.70098}\FunctionTok{,}
    \DataTypeTok{"speed"}\FunctionTok{:} \FloatTok{42.5}\FunctionTok{,}
    \DataTypeTok{"course"}\FunctionTok{:} \FloatTok{181.2}\FunctionTok{,}
    \DataTypeTok{"satellites"}\FunctionTok{:} \DecValTok{11}\FunctionTok{,}
    \DataTypeTok{"ignition"}\FunctionTok{:} \KeywordTok{true}\FunctionTok{,}
    \DataTypeTok{"error\_code"}\FunctionTok{:} \DecValTok{0}
  \FunctionTok{\}}
\FunctionTok{\}}
\end{Highlighting}
\end{Shaded}

Dữ liệu được gửi định kỳ theo cấu hình vận hành hiện tại: nhịp nhanh hơn
ở chế độ lái xe, thưa hơn ở các nhịp heartbeat khi đỗ. Tần suất gửi có
thể được điều chỉnh từ xa thông qua lệnh \texttt{update\_config} từ
máy chủ.

\begin{center}\rule{0.5\linewidth}{0.5pt}\end{center}

\subsubsection{4.2.4. Triển khai hệ thống máy chủ}

Hệ thống máy chủ là nơi biến dữ liệu từ thiết bị thành khả năng quan sát
và điều hành ở phía người quản lý. Trên cơ sở phương án đã chọn ở
Chương 3, phần này trình bày quá trình hiện thực hóa tầng máy chủ theo
từng khối: hạ tầng Docker, MQTT Bridge, API quản trị phía máy chủ, giao
diện quản trị và lớp giám sát vận hành.

Nếu nhìn theo hành trình dữ liệu, triển khai thực tế có thể tóm gọn như
sau: thiết bị gửi bản tin lên lớp tiếp nhận; lớp này kiểm tra và đẩy dữ
liệu sang từng nơi lưu phù hợp; sau đó API phía máy chủ và giao diện web dùng phần dữ
liệu đã được chuẩn hóa để truy vấn, điều khiển và hiển thị thời gian
thực. Nói cụ thể hơn, chuỗi đó lần lượt đi qua EMQX, MQTT Bridge,
VictoriaMetrics, VictoriaLogs, PostgreSQL, API phía máy chủ và giao diện
quản trị. Cách mô tả này giúp người đọc nối được kiến trúc ở Chương 3
với các cấu hình triển khai cụ thể ở phần sau.


{\thesissetfigureid{4.18}{fig-ch4-kien-truc-tong-the-he-thong-cloud-va-luong-du-lieu}
\begin{figure}[htbp]
\centering
\pandocbounded{\safeincludesvg{./assets/figures/09-chuong-4-trien-khai-cloud-hinh-4-15.svg}}
\caption{Kiến trúc tổng thể hệ thống máy chủ và luồng dữ liệu}
\label{fig:ch4-kien-truc-tong-the-he-thong-cloud-va-luong-du-lieu}
\end{figure}
}


\begin{center}\rule{0.5\linewidth}{0.5pt}\end{center}

\paragraph{4.2.4.1. Cấu hình Docker và hạ tầng đóng gói dịch vụ}\label{4241-cux1ea5u-huxecnh-docker-vuxe0-hux1ea1-tux1ea7ng-docker-infrastructure-setup}

\subparagraph{a) Mô hình triển khai tách theo từng dịch vụ}\label{a-muxf4-huxecnh-triux1ec3n-khai-per-service-ivm26-pattern}

Ở giai đoạn triển khai, nhóm chọn cách tách mỗi dịch vụ thành một đơn vị
cấu hình Docker Compose riêng. Cách tổ chức này phù hợp với đồ án vì vẫn
đủ đơn giản để cài đặt, nhưng đã phản ánh khá rõ tư duy chia dịch vụ
trong một hệ thống thật. Theo đó, mỗi dịch vụ có tệp cấu hình riêng thay
vì gom toàn bộ vào một tệp lớn.
Có thể hiểu ngắn gọn: mỗi dịch vụ chạy trong một container, tức một môi
trường thực thi tách biệt, nên dễ quản lý hơn khi cần sửa một phần mà
không muốn làm ảnh hưởng cả hệ thống. Cách tổ chức này có một số lợi
điểm rõ ràng:

\begin{itemize}
\tightlist
\item
  \textbf{Độc lập triển khai}: Mỗi dịch vụ có
  thể được khởi động, dừng lại, hoặc cập nhật riêng rẽ mà không ảnh
  hưởng đến các dịch vụ khác. Điều này đặc biệt hữu ích trong quá trình
  phát triển và giai đoạn dò lỗi.
\item
  \textbf{Quản lý tài nguyên riêng biệt}: Mỗi dịch vụ được cấu hình giới
  hạn tài nguyên (CPU, memory) phù hợp với nhu cầu thực tế, tránh tình
  trạng một dịch vụ chiếm dụng tài nguyên của dịch vụ khác.
\item
  \textbf{Dễ dàng mở rộng}: Có thể tăng thêm bản sao từng
  dịch vụ độc lập dựa trên nhu cầu tải thực tế.
\item
  \textbf{Đơn giản hóa quy trình kiểm tra và triển khai tự động}: Mỗi dịch vụ có thể có quy trình
  kiểm tra, đóng gói và triển khai tự động riêng.
\end{itemize}

Cấu trúc thư mục tổng thể của hệ thống được tổ chức như sau:

\begin{itemize}
\tightlist
\item
  \texttt{Tracking\_Backend/}: API phía máy chủ viết bằng Express +
  TypeScript.
\item
  \texttt{Tracking\_Frontend/}: giao diện quản trị web bằng Next.js.
\item
  \texttt{Tracking\_MqttBridge/}: dịch vụ nhận luồng MQTT đầu vào một cách
  độc lập.
\item
  \texttt{Tracking\_PostgreSQL/}: PostgreSQL và script khởi tạo.
\item
  \texttt{Tracking\_EMQX/}: MQTT broker EMQX và cấu hình.
\item
  \texttt{Tracking\_VictoriaMetrics/}, \texttt{Tracking\_VictoriaLogs/},
  \texttt{Tracking\_Grafana/}: khối quan sát vận hành gồm số liệu, nhật ký
  và bảng theo dõi.
\item
  \texttt{Tracking\_NPM/}: cổng chuyển tiếp truy cập từ bên ngoài vào các
  dịch vụ nội bộ, đồng thời quản lý chứng chỉ số.
\item
  \texttt{Tracking\_Data/}: dữ liệu vận hành thực tế (không đưa vào Git), gồm
  \texttt{postgres-data/}, \texttt{emqx-data/},
  \texttt{victoria-metrics-data/}, \texttt{victoria-logs-data/},
  \texttt{grafana-data/}.
\end{itemize}


{\thesissetfigureid{4.19}{fig-ch4-cau-truc-thu-muc-he-thong-theo-mo-hinh-trien-khai}
\begin{figure}[htbp]
\centering
\pandocbounded{\safeincludesvg{./assets/figures/09-chuong-4-trien-khai-cloud-hinh-4-16.svg}}
\caption{Cấu trúc thư mục hệ thống theo mô hình triển khai}
\label{fig:ch4-cau-truc-thu-muc-he-thong-theo-mo-hinh-trien-khai}
\end{figure}
}


\subparagraph{b) Mạng chia sẻ (Shared
Network)}\label{b-mux1ea1ng-chia-sux1ebb-shared-network}

Một chi tiết tưởng nhỏ nhưng ảnh hưởng lớn đến tính ổn định khi triển
khai là cách các dịch vụ ``nhìn thấy nhau'' trên Docker. Toàn bộ dịch vụ
được đưa vào một mạng Docker dùng chung; nhờ vậy, dù nằm
ở các file compose khác nhau, các container vẫn có thể giao tiếp qua tên
dịch vụ mà không cần cấu hình thủ công phức tạp.

\begin{Shaded}
\begin{Highlighting}[]
\CommentTok{\# Tạo mạng chia sẻ (chỉ chạy một lần đầu tiên)}
\ExtensionTok{docker}\NormalTok{ network create tracking{-}network}
\end{Highlighting}
\end{Shaded}

\subparagraph{c) Cấu hình Docker Compose cho các dịch vụ hạ
tầng}\label{c-cux1ea5u-huxecnh-docker-compose-cho-cuxe1c-dux1ecbch-vux1ee5-hux1ea1-tux1ea7ng}

Mỗi dịch vụ hạ tầng (infrastructure service) được cấu hình trong một
file \texttt{docker-compose.yml} riêng. Dưới đây là ví dụ cấu hình cho
dịch vụ PostgreSQL và EMQX:

\textbf{PostgreSQL (Tracking\_PostgreSQL/docker-compose.yml):}

\begin{Shaded}
\begin{Highlighting}[]
\FunctionTok{version}\KeywordTok{:}\AttributeTok{ }\StringTok{"3.8"}
\FunctionTok{services}\KeywordTok{:}
\AttributeTok{  }\FunctionTok{postgres}\KeywordTok{:}
\AttributeTok{    }\FunctionTok{image}\KeywordTok{:}\AttributeTok{ postgres:16{-}alpine}
\AttributeTok{    }\FunctionTok{container\_name}\KeywordTok{:}\AttributeTok{ tracking{-}postgres}
\AttributeTok{    }\FunctionTok{ports}\KeywordTok{:}
\AttributeTok{      }\KeywordTok{{-}}\AttributeTok{ }\StringTok{"5432:5432"}
\AttributeTok{    }\FunctionTok{environment}\KeywordTok{:}
\AttributeTok{      }\FunctionTok{POSTGRES\_DB}\KeywordTok{:}\AttributeTok{ vehicle\_tracking}
\AttributeTok{      }\FunctionTok{POSTGRES\_USER}\KeywordTok{:}\AttributeTok{ postgres}
\AttributeTok{      }\FunctionTok{POSTGRES\_PASSWORD}\KeywordTok{:}\AttributeTok{ $\{POSTGRESQL\_PASSWORD\}}
\AttributeTok{    }\FunctionTok{volumes}\KeywordTok{:}
\AttributeTok{      }\KeywordTok{{-}}\AttributeTok{ ../Tracking\_Data/postgres{-}data:/var/lib/postgresql/data}
\AttributeTok{      }\KeywordTok{{-}}\AttributeTok{ ./init:/docker{-}entrypoint{-}initdb.d}
\AttributeTok{    }\FunctionTok{restart}\KeywordTok{:}\AttributeTok{ unless{-}stopped}
\AttributeTok{    }\FunctionTok{networks}\KeywordTok{:}
\AttributeTok{      }\KeywordTok{{-}}\AttributeTok{ tracking{-}network}
\AttributeTok{    }\FunctionTok{deploy}\KeywordTok{:}
\AttributeTok{      }\FunctionTok{resources}\KeywordTok{:}
\AttributeTok{        }\FunctionTok{limits}\KeywordTok{:}
\AttributeTok{          }\FunctionTok{memory}\KeywordTok{:}\AttributeTok{ 512M}
\AttributeTok{          }\FunctionTok{cpus}\KeywordTok{:}\AttributeTok{ }\StringTok{"1.0"}

\FunctionTok{networks}\KeywordTok{:}
\AttributeTok{  }\FunctionTok{tracking{-}network}\KeywordTok{:}
\AttributeTok{    }\FunctionTok{external}\KeywordTok{:}\AttributeTok{ }\CharTok{true}
\end{Highlighting}
\end{Shaded}

\textbf{EMQX MQTT Broker (Tracking\_EMQX/docker-compose.yml):}

\begin{Shaded}
\begin{Highlighting}[]
\FunctionTok{version}\KeywordTok{:}\AttributeTok{ }\StringTok{"3.8"}
\FunctionTok{services}\KeywordTok{:}
\AttributeTok{  }\FunctionTok{emqx}\KeywordTok{:}
\AttributeTok{    }\FunctionTok{image}\KeywordTok{:}\AttributeTok{ emqx/emqx:5.4.0}
\AttributeTok{    }\FunctionTok{container\_name}\KeywordTok{:}\AttributeTok{ tracking{-}emqx}
\AttributeTok{    }\FunctionTok{ports}\KeywordTok{:}
\AttributeTok{      }\KeywordTok{{-}}\AttributeTok{ }\StringTok{"1883:1883"}\CommentTok{ \# MQTT}
\AttributeTok{      }\KeywordTok{{-}}\AttributeTok{ }\StringTok{"8883:8883"}\CommentTok{ \# MQTT TLS}
\AttributeTok{      }\KeywordTok{{-}}\AttributeTok{ }\StringTok{"8083:8083"}\CommentTok{ \# MQTT qua WebSocket}
\AttributeTok{      }\KeywordTok{{-}}\AttributeTok{ }\StringTok{"8084:8084"}\CommentTok{ \# MQTT qua WebSocket Secure}
\AttributeTok{      }\KeywordTok{{-}}\AttributeTok{ }\StringTok{"18083:18083"}\CommentTok{ \# EMQX Dashboard}
\AttributeTok{    }\FunctionTok{volumes}\KeywordTok{:}
\AttributeTok{      }\KeywordTok{{-}}\AttributeTok{ emqx{-}data:/opt/emqx/data}
\AttributeTok{      }\KeywordTok{{-}}\AttributeTok{ emqx{-}log:/opt/emqx/log}
\AttributeTok{    }\FunctionTok{environment}\KeywordTok{:}
\AttributeTok{      }\FunctionTok{EMQX\_DASHBOARD\_\_DEFAULT\_PASSWORD}\KeywordTok{:}\AttributeTok{ $\{EMQX\_DASHBOARD\_\_DEFAULT\_PASSWORD\}}
\AttributeTok{      }\FunctionTok{EMQX\_NODE\_\_COOKIE}\KeywordTok{:}\AttributeTok{ $\{EMQX\_NODE\_\_COOKIE\}}
\AttributeTok{    }\FunctionTok{restart}\KeywordTok{:}\AttributeTok{ unless{-}stopped}
\AttributeTok{    }\FunctionTok{networks}\KeywordTok{:}
\AttributeTok{      }\KeywordTok{{-}}\AttributeTok{ tracking{-}network}
\AttributeTok{    }\FunctionTok{deploy}\KeywordTok{:}
\AttributeTok{      }\FunctionTok{resources}\KeywordTok{:}
\AttributeTok{        }\FunctionTok{limits}\KeywordTok{:}
\AttributeTok{          }\FunctionTok{memory}\KeywordTok{:}\AttributeTok{ 1G}
\AttributeTok{          }\FunctionTok{cpus}\KeywordTok{:}\AttributeTok{ }\StringTok{"2.0"}

\FunctionTok{volumes}\KeywordTok{:}
\AttributeTok{  }\FunctionTok{emqx{-}data}\KeywordTok{:}
\AttributeTok{  }\FunctionTok{emqx{-}log}\KeywordTok{:}

\FunctionTok{networks}\KeywordTok{:}
\AttributeTok{  }\FunctionTok{tracking{-}network}\KeywordTok{:}
\AttributeTok{    }\FunctionTok{external}\KeywordTok{:}\AttributeTok{ }\CharTok{true}
\end{Highlighting}
\end{Shaded}

Cấu hình tương tự được áp dụng cho VictoriaMetrics (port 8428),
VictoriaLogs (port 9428), và Grafana (port 4002). Các dịch vụ đều chia
sẻ cùng mạng Docker dùng chung; riêng các thành phần lưu trữ như
PostgreSQL, VictoriaMetrics, VictoriaLogs và Grafana còn gắn volume
trong \texttt{Tracking\_Data/} để bảo đảm dữ liệu vẫn còn sau khi container
dừng rồi chạy lại.

\subparagraph{d) Thư mục dữ liệu và quản lý
volume}\label{d-thux1b0-mux1ee5c-dux1eef-liux1ec7u-vuxe0-quux1ea3n-luxfd-volume}

Thư mục \texttt{Tracking\_Data/} được sử dụng làm nơi lưu trữ dữ liệu
vận hành cho tất cả các dịch vụ. Thư mục này được loại khỏi phạm vi theo
dõi của Git để tránh đưa dữ liệu chạy thật vào kho mã nguồn.
Cấu trúc này đảm bảo:

\begin{itemize}
\tightlist
\item
  Dữ liệu không bị mất khi container bị xóa và tạo lại
\item
  Dễ dàng sao lưu toàn bộ dữ liệu hệ thống bằng cách sao chép
  thư mục \texttt{Tracking\_Data/}
\item
  Tách biệt giữa mã nguồn và dữ liệu vận hành
\end{itemize}

\subparagraph{e) Quy trình khởi động hệ
thống}\label{e-quy-truxecnh-khux1edfi-ux111ux1ed9ng-hux1ec7-thux1ed1ng}

Hệ thống được khởi động theo thứ tự phụ thuộc: dịch vụ hạ tầng trước,
sau đó đến các dịch vụ ứng dụng:

\begin{Shaded}
\begin{Highlighting}[]
\CommentTok{\# Bước 1: Khởi động dịch vụ hạ tầng}
\BuiltInTok{cd}\NormalTok{ Tracking\_PostgreSQL }\KeywordTok{\&\&} \ExtensionTok{docker}\NormalTok{ compose up }\AttributeTok{{-}d}
\BuiltInTok{cd}\NormalTok{ Tracking\_EMQX }\KeywordTok{\&\&} \ExtensionTok{docker}\NormalTok{ compose up }\AttributeTok{{-}d}
\BuiltInTok{cd}\NormalTok{ Tracking\_VictoriaMetrics }\KeywordTok{\&\&} \ExtensionTok{docker}\NormalTok{ compose up }\AttributeTok{{-}d}
\BuiltInTok{cd}\NormalTok{ Tracking\_VictoriaLogs }\KeywordTok{\&\&} \ExtensionTok{docker}\NormalTok{ compose up }\AttributeTok{{-}d}
\BuiltInTok{cd}\NormalTok{ Tracking\_Grafana }\KeywordTok{\&\&} \ExtensionTok{docker}\NormalTok{ compose up }\AttributeTok{{-}d}

\CommentTok{\# Bước 2: Khởi động dịch vụ ứng dụng}
\BuiltInTok{cd}\NormalTok{ Tracking\_MqttBridge }\KeywordTok{\&\&} \ExtensionTok{docker}\NormalTok{ compose up }\AttributeTok{{-}d} \AttributeTok{{-}{-}build}
\BuiltInTok{cd}\NormalTok{ Tracking\_Backend }\KeywordTok{\&\&} \ExtensionTok{docker}\NormalTok{ compose up }\AttributeTok{{-}d} \AttributeTok{{-}{-}build}
\BuiltInTok{cd}\NormalTok{ Tracking\_Frontend }\KeywordTok{\&\&} \ExtensionTok{docker}\NormalTok{ compose up }\AttributeTok{{-}d} \AttributeTok{{-}{-}build}
\end{Highlighting}
\end{Shaded}


{\thesissettableid{4.5}{tab-ch4-danh-sach-cac-dich-vu-docker-va-cong-truy-cap}
\def\LTcaptype{table}
\begin{longtable}{|L{0.248\linewidth}|L{0.306\linewidth}|L{0.139\linewidth}|L{0.262\linewidth}|}
\caption{Danh sách các dịch vụ được đóng gói bằng Docker và cổng truy cập}\label{tab:ch4-danh-sach-cac-dich-vu-docker-va-cong-truy-cap}\\
\hline \textbf{Dịch vụ} & \textbf{Tên container} & \textbf{Cổng} & \textbf{Mô tả} \\
\\\hline
\endfirsthead
\caption[]{Danh sách các dịch vụ được đóng gói bằng Docker và cổng truy cập (tiếp theo)}\\
\hline \textbf{Dịch vụ} & \textbf{Tên container} & \textbf{Cổng} & \textbf{Mô tả} \\
\\\hline
\endhead
\\\hline
\endfoot
\endlastfoot
PostgreSQL & tracking-postgres & 5432 & Cơ sở dữ liệu quan hệ \\\hline
EMQX & tracking-emqx & 1883, 8883, 8083, 8084, 18083 & Trạm trung chuyển
MQTT và giao diện quản lý \\\hline
VictoriaMetrics & tracking-victoriametrics & 8428 & Dữ liệu chuỗi thời
gian \\\hline
VictoriaLogs & tracking-victorialogs & 9428 & Nhật ký tập trung \\\hline
Grafana & tracking-grafana & 4002 & Màn hình giám sát \\\hline
MQTT Bridge & tracking-mqtt-bridge & --- & Dịch vụ trung gian xử lý bản
tin MQTT (không mở cổng ra ngoài) \\\hline
API phía máy chủ & tracking-backend & 4000 & REST API và kênh WebSocket \\\hline
Giao diện web & tracking-frontend & 4001 & Giao diện web \\\hline
Nginx Proxy Manager & tracking-npm & 80, 81, 443 & Cổng chuyển tiếp truy
cập từ Internet vào các dịch vụ nội bộ \\\hline
\end{longtable}
}

\textbf{Ánh xạ truy cập môi trường thử nghiệm:}

{\thesissettableid{4.5A}{tab-ch4-mapping-domain-va-port-cong-khai-moi-truong-thu-nghiem}
\def\LTcaptype{table}
\begin{longtable}{|L{0.158\linewidth}|L{0.097\linewidth}|L{0.216\linewidth}|L{0.122\linewidth}|L{0.352\linewidth}|}
\caption{Ánh xạ tên miền công khai, cổng truy cập và dịch vụ nội bộ trong môi trường thử nghiệm}\label{tab:ch4-mapping-domain-va-port-cong-khai-moi-truong-thu-nghiem}\\
\hline \textbf{Tên miền công khai} & \textbf{Cổng công khai} & \textbf{Dịch vụ nội bộ} & \textbf{Cổng nội bộ} & \textbf{Vai trò} \\
\\\hline
\endfirsthead
\caption[]{Ánh xạ tên miền công khai, cổng truy cập và dịch vụ nội bộ trong môi trường thử nghiệm (tiếp theo)}\\
\hline \textbf{Tên miền công khai} & \textbf{Cổng công khai} & \textbf{Dịch vụ nội bộ} & \textbf{Cổng nội bộ} & \textbf{Vai trò} \\
\\\hline
\endhead
\\\hline
\endfoot
\endlastfoot
\nolinkurl{thingdock.dev} & 443 & tracking-frontend & 4001 & Giao diện web đi qua lớp chuyển tiếp Nginx Proxy Manager \\\hline
\nolinkurl{api.thingdock.dev} & 443 & tracking-backend & 4000 & API phía máy chủ, Swagger và Socket.IO đi qua lớp chuyển tiếp Nginx Proxy Manager \\\hline
\nolinkurl{mqtt.thingdock.dev} & 443 & tracking-emqx & 8083 & MQTT qua WebSocket đi qua lớp chuyển tiếp Nginx Proxy Manager \\\hline
\nolinkurl{mqtts.thingdock.dev} & 8883 & tracking-emqx & 8883 & MQTT TLS trực tiếp cho thiết bị \\\hline
\nolinkurl{grafana.thingdock.dev} & 443 & tracking-grafana & 4002 & Grafana đi qua lớp chuyển tiếp Nginx Proxy Manager \\\hline
\nolinkurl{emqx.thingdock.dev} & 443 & tracking-emqx & 18083 & Giao diện quản lý EMQX đi qua lớp chuyển tiếp Nginx Proxy Manager \\\hline
\nolinkurl{npm.thingdock.dev} & 443 & tracking-npm & 81 & Trang quản trị NPM \\\hline
\nolinkurl{(nội bộ)} & 8428 & tracking-victoriametrics & 8428 & VictoriaMetrics, chỉ dùng nội bộ \\\hline
\end{longtable}
}

\textbf{Ghi chú:} Trong triển khai môi trường thử nghiệm hiện tại, Cloudflare chỉ giữ vai
trò lớp phân giải tên miền, tức giúp tên miền trỏ về địa chỉ IP công khai của VPS;
bản thân Cloudflare không giữ vai trò xử lý nghiệp vụ. Sau đó Nginx Proxy Manager mới làm lớp chuyển tiếp truy cập đến các container nội
bộ. Vì vậy, người dùng chỉ truy
cập \texttt{https://thingdock.dev}, \texttt{https://api.thingdock.dev},
\texttt{wss://mqtt.thingdock.dev}, \texttt{https://grafana.thingdock.dev},
\texttt{https://emqx.thingdock.dev} và \texttt{https://npm.thingdock.dev}
qua cổng HTTPS mặc định. Các port \texttt{4001}, \texttt{4000},
\texttt{8083}, \texttt{4002}, \texttt{18083} và \texttt{81} là port nội
bộ phía sau NPM, không xuất hiện trong URL public. Riêng MQTT TLS cho
thiết bị đi trực tiếp tới EMQX qua \texttt{mqtts.thingdock.dev:8883}.
VictoriaMetrics (\texttt{8428}) được giữ ở chế độ nội bộ, không công khai
qua subdomain (tên miền phụ).


\begin{center}\rule{0.5\linewidth}{0.5pt}\end{center}

\paragraph{4.2.4.2. Triển khai MQTT Bridge (MQTT Bridge
Implementation)}\label{4242-triux1ec3n-khai-mqtt-bridge-mqtt-bridge-implementation}

\subparagraph{a) Vai trò và kiến
trúc}\label{a-vai-truxf2-vuxe0-kiux1ebfn-truxfac}

MQTT Bridge (\texttt{Tracking\_MqttBridge/}) là lớp trung chuyển giữa
thiết bị và phần còn lại của hệ thống. Thiết bị gửi lên các bản tin MQTT
gọn để tiết kiệm băng thông, nhưng phía sau lại cần dữ liệu có cấu trúc
rõ ràng để lưu vào cơ sở dữ liệu, dựng giao diện quản trị và truy vết lỗi. Vì
vậy, MQTT Bridge làm nhiệm vụ nhận bản tin thô, kiểm tra nó, rồi chuyển
nó thành các bản ghi phù hợp cho từng nơi sử dụng. Việc tách thành dịch
vụ riêng cũng giúp đường vào dữ liệu vận hành gửi liên tục gọn hơn và không trộn lẫn với
nghiệp vụ web của API phía máy chủ.

Luồng xử lý của MQTT Bridge có thể hình dung theo bốn bước:

\begin{enumerate}
\tightlist
\item
  Thiết bị IoT gửi bản tin MQTT vào topic
  \texttt{v1/\{device\_id\}/rawdata} qua EMQX.
\item
  MQTT Bridge nhận bản tin, kiểm tra cấu trúc dữ liệu bằng Zod, tức là
  kiểm tra xem bản tin có đủ và đúng các trường bắt buộc hay không.
\item
  Sau khi hợp lệ, dữ liệu được gửi tới đúng nơi: dữ liệu vận hành gửi liên tục vào
  VictoriaMetrics, sự kiện hoặc lỗi vào VictoriaLogs, còn trạng thái
  phiên vào PostgreSQL.
\item
  Cuối cùng, MQTT Bridge phát sự kiện nội bộ để API phía máy chủ nhận và đẩy cập
  nhật mới ra giao diện quản trị theo thời gian thực.
\end{enumerate}


{\thesissetfigureid{4.20}{fig-ch4-luong-xu-ly-du-lieu-cua-mqtt-bridge}
\begin{figure}[htbp]
\centering
\pandocbounded{\safeincludesvg{./assets/figures/09-chuong-4-trien-khai-cloud-hinh-4-17.svg}}
\caption{Luồng xử lý dữ liệu của MQTT Bridge}
\label{fig:ch4-luong-xu-ly-du-lieu-cua-mqtt-bridge}
\end{figure}
}


\subparagraph{b) Cấu trúc mã
nguồn}\label{b-cux1ea5u-truxfac-muxe3-nguux1ed3n}

\begin{itemize}
\tightlist
\item
  \texttt{src/index.ts}: điểm khởi động dịch vụ.
\item
  \texttt{src/handlers/}: xử lý bản tin theo nhóm topic
  (\texttt{rawdata}, \texttt{status}, \texttt{events}, \texttt{firmware}).
\item
  \texttt{src/infrastructure/}: kết nối PostgreSQL, logger, VictoriaMetrics
  và VictoriaLogs.
\item
  \texttt{src/mqtt/}: client MQTT và cấu hình đăng ký QoS.
\item
  \texttt{src/services/}: nghiệp vụ nền như batch writer, xác thực thiết bị,
  logic kiểm tra vùng cho phép.
\item
  \texttt{src/validators/}, \texttt{src/types/}, \texttt{src/utils/}: chuẩn
  hóa khuôn dữ liệu hợp lệ, kiểu dữ liệu và tiện ích hỗ trợ.
\end{itemize}

\subparagraph{c) Kết nối MQTT và xử lý
bản tin}\label{c-kux1ebft-nux1ed1i-mqtt-vuxe0-xux1eed-luxfd-message}

MQTT Bridge kết nối tới EMQX Broker và subscribe vào các topic tương ứng
với dữ liệu từ thiết bị IoT. Mỗi topic được định tuyến (route) tới
handler phù hợp theo suffix của topic:

\begin{Shaded}
\begin{Highlighting}[]
\CommentTok{// index.ts + mqtt/subscriptions.ts}
\ControlFlowTok{await} \FunctionTok{subscribeToDeviceTopics}\NormalTok{(client)}\OperatorTok{;}

\NormalTok{client}\OperatorTok{.}\FunctionTok{on}\NormalTok{(}\StringTok{"message"}\OperatorTok{,}\NormalTok{ (topic}\OperatorTok{,}\NormalTok{ message) }\KeywordTok{=\textgreater{}}\NormalTok{ \{}
  \KeywordTok{const}\NormalTok{ deviceId }\OperatorTok{=} \FunctionTok{extractDeviceId}\NormalTok{(topic)}\OperatorTok{;}
  \KeywordTok{const}\NormalTok{ suffix }\OperatorTok{=}\NormalTok{ topic}\OperatorTok{.}\FunctionTok{split}\NormalTok{(}\StringTok{"/"}\NormalTok{)}\OperatorTok{.}\FunctionTok{slice}\NormalTok{(}\DecValTok{2}\NormalTok{)}\OperatorTok{.}\FunctionTok{join}\NormalTok{(}\StringTok{"/"}\NormalTok{)}\OperatorTok{;}

  \ControlFlowTok{switch}\NormalTok{ (suffix) \{}
    \ControlFlowTok{case} \StringTok{"rawdata"}\OperatorTok{:}
      \DataTypeTok{void} \FunctionTok{handleRawData}\NormalTok{(deviceId}\OperatorTok{!,}\NormalTok{ message)}\OperatorTok{;}
      \ControlFlowTok{break}\OperatorTok{;}
    \ControlFlowTok{case} \StringTok{"status"}\OperatorTok{:}
      \DataTypeTok{void} \FunctionTok{handleStatus}\NormalTok{(deviceId}\OperatorTok{!,}\NormalTok{ message)}\OperatorTok{;}
      \ControlFlowTok{break}\OperatorTok{;}
    \ControlFlowTok{case} \StringTok{"events"}\OperatorTok{:}
      \DataTypeTok{void} \FunctionTok{handleEvent}\NormalTok{(deviceId}\OperatorTok{!,}\NormalTok{ message)}\OperatorTok{;}
      \ControlFlowTok{break}\OperatorTok{;}
    \ControlFlowTok{case} \StringTok{"firmware"}\OperatorTok{:}
      \DataTypeTok{void} \FunctionTok{handleFirmware}\NormalTok{(deviceId}\OperatorTok{!,}\NormalTok{ message)}\OperatorTok{;}
      \ControlFlowTok{break}\OperatorTok{;}
\NormalTok{  \}}
\NormalTok{\})}\OperatorTok{;}
\end{Highlighting}
\end{Shaded}

Trong các handler này, \texttt{handleFirmware()} là mắt xích mới đáng
chú ý sau khi bổ sung OTA. Handler này không chỉ kiểm tra gói dữ liệu trạng
thái firmware mà còn ghi tiến trình cập nhật vào
\texttt{firmware\_update\_log}, đồng thời phát bản ghi sang VictoriaLogs
để thuận tiện cho việc truy vết khi cập nhật thất bại hoặc rollback.

\subparagraph{d) Xác thực dữ liệu đầu vào (kiểm tra cấu trúc gói dữ
liệu)}\label{d-xuxe1c-thux1ef1c-dux1eef-liux1ec7u-ux111ux1ea7u-vuxe0o-payload-validation}

Mỗi bản tin MQTT nhận được đều được xác thực bằng Zod schema trước khi
xử lý, giúp chặn sớm gói dữ liệu sai cấu trúc và giảm rủi ro ghi dữ liệu bẩn
vào luồng xử lý:

\begin{Shaded}
\begin{Highlighting}[]
\CommentTok{// validators/payload.validator.ts}
\ImportTok{import}\NormalTok{ \{ z \} }\ImportTok{from} \StringTok{"zod"}\OperatorTok{;}

\ImportTok{export} \KeywordTok{const}\NormalTok{ rawDataSchema }\OperatorTok{=}\NormalTok{ z}\OperatorTok{.}\FunctionTok{object}\NormalTok{(\{}
\NormalTok{  device\_id}\OperatorTok{:}\NormalTok{ z}\OperatorTok{.}\FunctionTok{string}\NormalTok{()}\OperatorTok{.}\FunctionTok{min}\NormalTok{(}\DecValTok{1}\NormalTok{)}\OperatorTok{,}
\NormalTok{  auth\_token}\OperatorTok{:}\NormalTok{ z}\OperatorTok{.}\FunctionTok{string}\NormalTok{()}\OperatorTok{.}\FunctionTok{min}\NormalTok{(}\DecValTok{1}\NormalTok{)}\OperatorTok{,}
\NormalTok{  timestamp}\OperatorTok{:}\NormalTok{ z}\OperatorTok{.}\FunctionTok{number}\NormalTok{()}\OperatorTok{.}\FunctionTok{positive}\NormalTok{()}\OperatorTok{,}
\NormalTok{  uptime}\OperatorTok{:}\NormalTok{ z}\OperatorTok{.}\FunctionTok{number}\NormalTok{()}\OperatorTok{.}\FunctionTok{nonnegative}\NormalTok{()}\OperatorTok{.}\FunctionTok{optional}\NormalTok{()}\OperatorTok{,}
\NormalTok{  data}\OperatorTok{:}\NormalTok{ z}\OperatorTok{.}\FunctionTok{object}\NormalTok{(\{}
\NormalTok{    vibration}\OperatorTok{:}\NormalTok{ z}\OperatorTok{.}\FunctionTok{number}\NormalTok{()}\OperatorTok{.}\FunctionTok{optional}\NormalTok{()}\OperatorTok{,}
\NormalTok{    battery\_top}\OperatorTok{:}\NormalTok{ z}\OperatorTok{.}\FunctionTok{number}\NormalTok{()}\OperatorTok{.}\FunctionTok{optional}\NormalTok{()}\OperatorTok{,}
\NormalTok{    battery\_bot}\OperatorTok{:}\NormalTok{ z}\OperatorTok{.}\FunctionTok{number}\NormalTok{()}\OperatorTok{.}\FunctionTok{optional}\NormalTok{()}\OperatorTok{,}
\NormalTok{    latitude}\OperatorTok{:}\NormalTok{ z}\OperatorTok{.}\FunctionTok{number}\NormalTok{()}\OperatorTok{.}\FunctionTok{min}\NormalTok{(}\OperatorTok{{-}}\DecValTok{90}\NormalTok{)}\OperatorTok{.}\FunctionTok{max}\NormalTok{(}\DecValTok{90}\NormalTok{)}\OperatorTok{.}\FunctionTok{optional}\NormalTok{()}\OperatorTok{,}
\NormalTok{    longitude}\OperatorTok{:}\NormalTok{ z}\OperatorTok{.}\FunctionTok{number}\NormalTok{()}\OperatorTok{.}\FunctionTok{min}\NormalTok{(}\OperatorTok{{-}}\DecValTok{180}\NormalTok{)}\OperatorTok{.}\FunctionTok{max}\NormalTok{(}\DecValTok{180}\NormalTok{)}\OperatorTok{.}\FunctionTok{optional}\NormalTok{()}\OperatorTok{,}
\NormalTok{    speed}\OperatorTok{:}\NormalTok{ z}\OperatorTok{.}\FunctionTok{number}\NormalTok{()}\OperatorTok{.}\FunctionTok{nonnegative}\NormalTok{()}\OperatorTok{.}\FunctionTok{optional}\NormalTok{()}\OperatorTok{,}
\NormalTok{    course}\OperatorTok{:}\NormalTok{ z}\OperatorTok{.}\FunctionTok{number}\NormalTok{()}\OperatorTok{.}\FunctionTok{min}\NormalTok{(}\DecValTok{0}\NormalTok{)}\OperatorTok{.}\FunctionTok{max}\NormalTok{(}\DecValTok{360}\NormalTok{)}\OperatorTok{.}\FunctionTok{optional}\NormalTok{()}\OperatorTok{,}
\NormalTok{    satellites}\OperatorTok{:}\NormalTok{ z}\OperatorTok{.}\FunctionTok{number}\NormalTok{()}\OperatorTok{.}\FunctionTok{int}\NormalTok{()}\OperatorTok{.}\FunctionTok{nonnegative}\NormalTok{()}\OperatorTok{.}\FunctionTok{optional}\NormalTok{()}\OperatorTok{,}
\NormalTok{    ignition}\OperatorTok{:}\NormalTok{ z}\OperatorTok{.}\FunctionTok{boolean}\NormalTok{()}\OperatorTok{.}\FunctionTok{optional}\NormalTok{()}\OperatorTok{,}
\NormalTok{    error\_code}\OperatorTok{:}\NormalTok{ z}\OperatorTok{.}\FunctionTok{number}\NormalTok{()}\OperatorTok{.}\FunctionTok{int}\NormalTok{()}\OperatorTok{.}\FunctionTok{optional}\NormalTok{()}\OperatorTok{,}
\NormalTok{  \})}\OperatorTok{,}
\NormalTok{\})}\OperatorTok{;}
\end{Highlighting}
\end{Shaded}

Ngoài schema cho dữ liệu vận hành gửi liên tục, MQTT Bridge còn dùng
\texttt{firmwareStatusSchema} để kiểm tra chặt các trường
\texttt{jobId}, \texttt{status}, \texttt{progress},
\texttt{targetVersion}, \texttt{currentVersion}, \texttt{partition} và
\texttt{error}. Nhờ vậy, luồng OTA có một lớp bảo vệ riêng trước khi dữ
liệu được ghi vào cơ sở dữ liệu.

\subparagraph{e) Ghi dữ liệu kép}\label{e-ghi-dux1eef-liux1ec7u-kuxe9p-dual-write-pattern}

MQTT Bridge thực hiện cơ chế ghi dữ liệu kép, tức một lần nhận bản tin
nhưng ghi ra hai nơi phục vụ hai mục đích khác nhau: dữ liệu
vận hành gửi liên tục được ghi đồng thời vào VictoriaMetrics (cơ sở dữ liệu chuỗi
thời gian) và VictoriaLogs (nhật ký sự kiện), trong khi thông tin trạng
thái thiết bị được cập nhật vào PostgreSQL:

\begin{itemize}
\tightlist
\item
  \textbf{VictoriaMetrics}: Lưu trữ các metric theo thời gian như
  \texttt{vehicle\_latitude}, \texttt{vehicle\_longitude},
  \texttt{vehicle\_speed}, \texttt{vehicle\_vibration},
  \texttt{vehicle\_battery\_top}, \texttt{vehicle\_battery\_bot},
  \texttt{vehicle\_ignition}, \texttt{vehicle\_uptime}. Dữ liệu này phục
  vụ cho truy vấn lịch sử và biểu đồ vận hành.
\item
  \textbf{VictoriaLogs}: Lưu trữ các sự kiện hệ thống như kết nối/ngắt
  kết nối thiết bị, cảnh báo, lỗi xác thực gói dữ liệu. Dữ liệu này phục vụ
  cho phân tích và debug.
\item
  \textbf{PostgreSQL}: Cập nhật trạng thái hiện tại của thiết bị
  (\texttt{last\_seen\_at}, \texttt{current\_status}), tạo bản ghi phiên
  hoạt động (session), và lưu các thực thể nghiệp vụ như
  alert, trip và vùng cho phép.
\end{itemize}

\subparagraph{f) Cấu hình Docker cho MQTT
Bridge}\label{f-cux1ea5u-huxecnh-docker-cho-mqtt-bridge}

\begin{Shaded}
\begin{Highlighting}[]
\CommentTok{\# Tracking\_MqttBridge/docker{-}compose.yml}
\FunctionTok{version}\KeywordTok{:}\AttributeTok{ }\StringTok{"3.8"}
\FunctionTok{services}\KeywordTok{:}
\AttributeTok{  }\FunctionTok{mqtt{-}bridge}\KeywordTok{:}
\AttributeTok{    }\FunctionTok{build}\KeywordTok{:}
\AttributeTok{      }\FunctionTok{context}\KeywordTok{:}\AttributeTok{ .}
\AttributeTok{      }\FunctionTok{dockerfile}\KeywordTok{:}\AttributeTok{ Dockerfile}
\AttributeTok{    }\FunctionTok{container\_name}\KeywordTok{:}\AttributeTok{ tracking{-}mqtt{-}bridge}
\AttributeTok{    }\FunctionTok{env\_file}\KeywordTok{:}\AttributeTok{ .env}
\AttributeTok{    }\FunctionTok{restart}\KeywordTok{:}\AttributeTok{ unless{-}stopped}
\AttributeTok{    }\FunctionTok{networks}\KeywordTok{:}
\AttributeTok{      }\KeywordTok{{-}}\AttributeTok{ tracking{-}network}
\AttributeTok{    }\FunctionTok{deploy}\KeywordTok{:}
\AttributeTok{      }\FunctionTok{resources}\KeywordTok{:}
\AttributeTok{        }\FunctionTok{limits}\KeywordTok{:}
\AttributeTok{          }\FunctionTok{memory}\KeywordTok{:}\AttributeTok{ 256M}
\AttributeTok{          }\FunctionTok{cpus}\KeywordTok{:}\AttributeTok{ }\StringTok{"0.5"}

\FunctionTok{networks}\KeywordTok{:}
\AttributeTok{  }\FunctionTok{tracking{-}network}\KeywordTok{:}
\AttributeTok{    }\FunctionTok{external}\KeywordTok{:}\AttributeTok{ }\CharTok{true}
\end{Highlighting}
\end{Shaded}

Dịch vụ MQTT Bridge không cần expose port ra bên ngoài vì nó chỉ giao
tiếp nội bộ với các dịch vụ khác thông qua mạng Docker dùng chung.
Giới hạn tài nguyên được đặt ở mức 256 MB RAM và 0.5 CPU, phù hợp với
vai trò xử lý các gói dữ liệu nhỏ.

\begin{center}\rule{0.5\linewidth}{0.5pt}\end{center}

\paragraph{4.2.4.3. Triển khai API phía máy chủ}\label{4243-triux1ec3n-khai-backend-api-backend-api-implementation}

\subparagraph{a) Lựa chọn công
nghệ}\label{a-lux1ef1a-chux1ecdn-cuxf4ng-nghux1ec7-1}

API phía máy chủ được xây dựng trên nền tảng Node.js với Express.js và
TypeScript. Đây là lớp phục vụ trực tiếp giao diện quản trị: nhận yêu
cầu từ người dùng, kiểm tra quyền truy cập, lấy dữ liệu đúng chỗ rồi trả
lại kết quả ổn định. Vì đường dữ liệu vận hành gửi liên tục đã được MQTT
Bridge gánh riêng, API phía máy chủ không phải ôm cả dòng dữ liệu kỹ
thuật lẫn nghiệp vụ quản trị trong cùng một chỗ. Việc lựa chọn công nghệ
ở đây vì thế ưu tiên sự rõ ràng và linh hoạt trong phát triển lâu dài.


{\thesissettableid{4.6}{tab-ch4-so-sanh-framework-backend-cho-he-thong-iot}
\def\LTcaptype{table}
\begin{longtable}{|L{0.246\linewidth}|L{0.269\linewidth}|L{0.318\linewidth}|L{0.121\linewidth}|}
\caption{So sánh các khung phát triển API phía máy chủ cho hệ thống IoT}\label{tab:ch4-so-sanh-framework-backend-cho-he-thong-iot}\\
\hline \textbf{Tiêu chí} & \textbf{Express + TypeScript} & \textbf{NestJS} & \textbf{Go + Gin} \\
\\\hline
\endfirsthead
\caption[]{So sánh các khung phát triển API phía máy chủ cho hệ thống IoT (tiếp theo)}\\
\hline \textbf{Tiêu chí} & \textbf{Express + TypeScript} & \textbf{NestJS} & \textbf{Go + Gin} \\
\\\hline
\endhead
\\\hline
\endfoot
\endlastfoot
Độ phức tạp khi học và dựng dự án & Thấp & Trung bình & Cao \\\hline
Linh hoạt khi tổ chức mã nguồn & Rất cao & Trung bình (bị ràng buộc bởi khuôn mẫu có sẵn) & Cao \\\hline
Hiệu năng & Tốt & Tốt & Rất tốt \\\hline
Lượng mã khung lặp lại & Ít & Nhiều (decorator, module) & Trung bình \\\hline
Hệ sinh thái IoT & Rất tốt & Tốt & Tốt \\\hline
Hỗ trợ cập nhật thời gian thực & Rất tốt (tích hợp Socket.IO thuận tiện) & Rất tốt & Tốt \\\hline
\end{longtable}
}


\textbf{Lý do lựa chọn Express + TypeScript:}

\begin{itemize}
\tightlist
\item
  \textbf{Linh hoạt cao}: Không bị ràng buộc mạnh bởi các quy ước
  framework như NestJS, nên nhóm có thể tổ chức code bám theo nghiệp vụ
  và thay đổi nhanh khi yêu cầu thực tế còn dịch chuyển.
\item
  \textbf{Nhẹ và nhanh}: Ít phần gánh kèm hơn NestJS, phù hợp với hệ thống
  IoT cần phản hồi tốt cho giao diện quản trị mà vẫn giữ tài
  nguyên máy chủ ở mức hợp lý.
\item
  \textbf{Zod Validation}: Sử dụng Zod để chặn yêu cầu sai cấu trúc ngay
  từ đầu vào. Đây là lớp bảo vệ quan trọng vì API phía máy chủ nhận dữ
  liệu từ cả giao diện người dùng lẫn các luồng nội bộ khác.
\item
  \textbf{Kiểm soát middleware chain}: Toàn quyền kiểm soát thứ tự và
  cách thức xử lý của middleware, giúp tách rõ xác thực, ghi log, bắt
  lỗi và đo hiệu năng.
\end{itemize}

\textbf{Các thành phần công nghệ chính:}

\begin{itemize}
\tightlist
\item
  \textbf{Nền chạy và khung phát triển:} Node.js 20+ (LTS), Express 4.x, TypeScript
  5.x.
\item
  \textbf{Kiểm tra dữ liệu đầu vào:} Zod, thư viện dùng để kiểm tra dữ
  liệu gửi vào có đúng cấu trúc ngay khi chương trình chạy.
\item
  \textbf{Lưu trữ dữ liệu:} PostgreSQL 16 (driver \texttt{pg}) cho dữ liệu quan
  hệ; VictoriaMetrics (PromQL/MetricsQL) cho dữ liệu thay đổi liên tục theo thời gian.
\item
  \textbf{Kênh cập nhật thời gian thực:} Socket.IO 4.8 và MQTT.js tích hợp EMQX.
\item
  \textbf{Bảo mật và xác thực:} Cơ chế giữ phiên trên máy chủ, token được băm
  SHA-256 trong cơ sở dữ liệu, cùng Helmet, CORS và giới hạn tần suất gọi.
\item
  \textbf{Tài liệu API:} \texttt{swagger-ui-express} và OpenAPI để sinh tài liệu từ mã nguồn.
\end{itemize}

\subparagraph{b) Kiến trúc Domain-Driven Design
(DDD)}\label{b-kiux1ebfn-truxfac-domain-driven-design-ddd-1}

API phía máy chủ được tổ chức theo kiến trúc Domain-Driven Design với
các lớp rõ ràng: lớp nhận yêu cầu, lớp xử lý nghiệp vụ và lớp truy cập
dữ liệu. Cách tổ chức này gom mỗi bài toán nghiệp vụ vào một khu vực
riêng, nên người đọc mã nguồn có thể lần theo ``một việc'' trọn vẹn
thay vì phải nhảy qua quá nhiều thư mục kỹ thuật rời rạc. Mỗi domain vì
thế có thư mục riêng với các thành phần chính như \texttt{services/},
\texttt{repositories/} và \texttt{types/}.

\begin{itemize}
\tightlist
\item
  \texttt{index.ts}: điểm khởi động của API phía máy chủ.
\item
  \texttt{api/}: lớp API gồm routes và tài liệu OpenAPI.
\item
  \texttt{domain/}: lớp nghiệp vụ theo từng module
  (\texttt{auth}, \texttt{device}, \texttt{vehicle}, \texttt{trip},
  \texttt{alert}, \texttt{geofence}, \texttt{maintenance},
  \texttt{firmware}, \texttt{dashboard}, \dots), mỗi module tách rõ
  \texttt{services/}, \texttt{repositories/}, \texttt{types/}.
\item
  \texttt{infrastructure/}: tích hợp cơ sở dữ liệu, nhật ký, chỉ số giám sát và
  kênh cập nhật thời gian thực.
\item
  \texttt{middleware/}: xác thực, bắt lỗi, giới hạn tần suất, mã truy vết từng
  yêu cầu, ghi chỉ số và theo dõi lỗi.
\item
  \texttt{shared/}: kiểu dữ liệu và tiện ích dùng chung.
\end{itemize}


{\thesissetfigureid{4.21}{fig-ch4-cau-truc-thu-muc-backend-theo-kien-truc-ddd}
\begin{figure}[htbp]
\centering
\pandocbounded{\safeincludesvg{./assets/figures/09-chuong-4-trien-khai-cloud-hinh-4-18.svg}}
\caption{Cấu trúc thư mục API phía máy chủ theo kiến trúc DDD}
\label{fig:ch4-cau-truc-thu-muc-backend-theo-kien-truc-ddd}
\end{figure}
}


Hệ thống bao gồm hơn 20 module nghiệp vụ, mỗi module chịu trách nhiệm
một lĩnh vực cụ thể. Cách chia này giúp mã nguồn dễ bảo trì, dễ kiểm thử
và dễ mở rộng khi cần thêm tính năng mới.


{\thesissettableid{4.7}{tab-ch4-danh-sach-cac-domain-modules-chinh}
\def\LTcaptype{table}
\begin{longtable}{|L{0.263\linewidth}|L{0.162\linewidth}|L{0.539\linewidth}|}
\caption{Danh sách các module nghiệp vụ chính}\label{tab:ch4-danh-sach-cac-domain-modules-chinh}\\
\hline \textbf{Module} & \textbf{Số lượng dịch vụ xử lý} & \textbf{Mô tả} \\
\\\hline
\endfirsthead
\caption[]{Danh sách các module nghiệp vụ chính (tiếp theo)}\\
\hline \textbf{Module} & \textbf{Số lượng dịch vụ xử lý} & \textbf{Mô tả} \\
\\\hline
\endhead
\\\hline
\endfoot
\endlastfoot
auth & 2 & Đăng nhập/đăng xuất theo phiên làm việc, quản lý người dùng \\\hline
device & 8 & CRUD thiết bị, phiên hoạt động, trạng thái đang chạy, dữ liệu vận hành, lệnh
điều khiển \\\hline
vehicle & 3 & CRUD xe, danh sách, gán thiết bị \\\hline
customer & 2 & CRUD khách hàng, danh sách \\\hline
driver & 2 & CRUD tài xế, danh sách \\\hline
trip & 2 & CRUD chuyến đi, danh sách \\\hline
geofence & 2 & CRUD vùng địa lý, danh sách \\\hline
maintenance & 2 & CRUD bảo trì, danh sách \\\hline
firmware & 5 & OTA: upload, danh sách, kích hoạt, triển khai, xóa \\\hline
export & 3 & Quản lý tác vụ xuất, tạo file, xử lý \\\hline
dashboard (tổng quan) & 1 & Thống kê tổng quan \\\hline
telemetry (dữ liệu vận hành) & 1 & Truy vấn lịch sử dữ liệu vận hành \\\hline
notification & 1 & Thông báo đẩy tới người dùng \\\hline
simulator & 1 & Mô phỏng thiết bị để kiểm thử \\\hline
\end{longtable}
}


\subparagraph{c) Xác thực và phân quyền (Authentication \&
Authorization)}\label{c-xuxe1c-thux1ef1c-vuxe0-phuxe2n-quyux1ec1n-authentication--authorization}

Hệ thống sử dụng cơ chế xác thực dựa trên phiên làm việc, nghĩa là trạng
thái đăng nhập được giữ trong cơ sở dữ liệu PostgreSQL. So với cách dùng
JWT tự mang thông tin trong mỗi token, cách làm này phù hợp hơn với hệ
quản trị nội bộ vì dễ kiểm soát và dễ thu hồi phiên khi cần:

\begin{itemize}
\tightlist
\item
  \textbf{Token được hash trước khi lưu}: Dù cơ sở dữ liệu bị lộ, kẻ tấn
  công cũng không thể lấy nguyên token đã dùng để đăng nhập.
\item
  \textbf{Mật khẩu được hash bằng bcryptjs}: Đây là lớp bảo vệ cơ bản để
  giảm rủi ro nếu dữ liệu người dùng bị lộ.
\item
  \textbf{Thu hồi phiên ngay lập tức}: Quản trị viên chỉ cần xóa bản ghi
  phiên trong cơ sở dữ liệu là có thể chặn phiên đó ngay, không phải chờ
  token tự hết hạn.
\end{itemize}

\subparagraph{d) Middleware Chain}\label{d-middleware-chain}

Các middleware được sắp xếp theo thứ tự xử lý cụ thể, đảm bảo bảo mật và
hiệu năng:

\begin{Shaded}
\begin{Highlighting}[]
\CommentTok{// index.ts {-} Thứ tự middleware}
\NormalTok{app}\OperatorTok{.}\FunctionTok{use}\NormalTok{(sentryRequestHandler)}\OperatorTok{;} \CommentTok{// 1. Sentry error tracking}
\NormalTok{app}\OperatorTok{.}\FunctionTok{use}\NormalTok{(helmetMiddleware)}\OperatorTok{;} \CommentTok{// 2. Security headers (XSS, CSP)}
\NormalTok{app}\OperatorTok{.}\FunctionTok{use}\NormalTok{(compressionMiddleware)}\OperatorTok{;} \CommentTok{// 3. Nén Gzip/Brotli}
\NormalTok{app}\OperatorTok{.}\FunctionTok{use}\NormalTok{(corsMiddleware)}\OperatorTok{;} \CommentTok{// 4. Cross{-}Origin Resource Sharing}
\NormalTok{app}\OperatorTok{.}\FunctionTok{use}\NormalTok{(requestIdMiddleware)}\OperatorTok{;} \CommentTok{// 5. Request{-}ID correlation (X{-}Request{-}ID)}
\NormalTok{app}\OperatorTok{.}\FunctionTok{use}\NormalTok{(httpMetricsMiddleware)}\OperatorTok{;} \CommentTok{// 6. Prometheus HTTP metrics}
\NormalTok{app}\OperatorTok{.}\FunctionTok{use}\NormalTok{(express}\OperatorTok{.}\FunctionTok{json}\NormalTok{())}\OperatorTok{;} \CommentTok{// 7. Body parser (JSON)}
\NormalTok{app}\OperatorTok{.}\FunctionTok{use}\NormalTok{(requestLogger)}\OperatorTok{;} \CommentTok{// 8. Ghi nhật ký request}
\NormalTok{app}\OperatorTok{.}\FunctionTok{use}\NormalTok{(rateLimitMiddleware)}\OperatorTok{;} \CommentTok{// 9. Giới hạn tần suất request}

\CommentTok{// Routes (định tuyến API)}
\NormalTok{app}\OperatorTok{.}\FunctionTok{use}\NormalTok{(}\StringTok{"/api/v1/auth"}\OperatorTok{,}\NormalTok{ authRoutes)}\OperatorTok{;}
\NormalTok{app}\OperatorTok{.}\FunctionTok{use}\NormalTok{(}\StringTok{"/api/v1/devices"}\OperatorTok{,}\NormalTok{ deviceRoutes)}\OperatorTok{;}
\NormalTok{app}\OperatorTok{.}\FunctionTok{use}\NormalTok{(}\StringTok{"/api/v1/iot"}\OperatorTok{,}\NormalTok{ iotRoutes)}\OperatorTok{;} \CommentTok{// Không yêu cầu xác thực}
\NormalTok{app}\OperatorTok{.}\FunctionTok{use}\NormalTok{(}\StringTok{"/api/v1/vehicles"}\OperatorTok{,}\NormalTok{ vehicleRoutes)}\OperatorTok{;}
\NormalTok{app}\OperatorTok{.}\FunctionTok{use}\NormalTok{(}\StringTok{"/api/v1/customers"}\OperatorTok{,}\NormalTok{ customerRoutes)}\OperatorTok{;}
\CommentTok{// ... các route khác}

\CommentTok{// Xử lý lỗi (cuối cùng)}
\NormalTok{app}\OperatorTok{.}\FunctionTok{use}\NormalTok{(errorHandler)}\OperatorTok{;}
\NormalTok{app}\OperatorTok{.}\FunctionTok{use}\NormalTok{(sentryErrorHandler)}\OperatorTok{;}
\end{Highlighting}
\end{Shaded}

Middleware \texttt{requestIdMiddleware} tạo một mã nhận dạng riêng cho
mỗi yêu cầu và gắn nó vào header \texttt{X-Request-ID}. Nhờ mã này, khi
xảy ra lỗi người vận hành có thể lần theo cùng một yêu cầu xuyên suốt từ
log của API, lớp xử lý nghiệp vụ cho tới lớp truy cập dữ liệu.

\subparagraph{e) API Endpoints}\label{e-api-endpoints}

Hệ thống cung cấp các nhóm API endpoint sau:


{\thesissettableid{4.8}{tab-ch4-tong-hop-cac-nhom-api-endpoint-chinh}
\def\LTcaptype{table}
\begin{longtable}{|L{0.213\linewidth}|L{0.141\linewidth}|L{0.172\linewidth}|L{0.429\linewidth}|}
\caption{Tổng hợp các nhóm API endpoint chính}\label{tab:ch4-tong-hop-cac-nhom-api-endpoint-chinh}\\
\hline \textbf{Nhóm} & \textbf{Đường dẫn gốc} & \textbf{Xác thực} & \textbf{Mô tả} \\
\\\hline
\endfirsthead
\caption[]{Tổng hợp các nhóm API endpoint chính (tiếp theo)}\\
\hline \textbf{Nhóm} & \textbf{Đường dẫn gốc} & \textbf{Xác thực} & \textbf{Mô tả} \\
\\\hline
\endhead
\\\hline
\endfoot
\endlastfoot
Đăng nhập và phiên làm việc & \nolinkurl{/api/v1/auth} & Không (login) / Có & Đăng nhập,
đăng xuất, quản lý người dùng \\\hline
Thiết bị & \nolinkurl{/api/v1/devices} & Có & CRUD thiết bị, phiên hoạt động,
trạng thái chạy và lệnh \\\hline
Dữ liệu IoT & \nolinkurl{/api/v1/iot} & Không (public) & Nhận dữ liệu từ thiết
bị IoT \\\hline
Tổng quan & \nolinkurl{/api/v1/dashboard} & Có & Thống kê, hoạt động gần
đây, cảnh báo \\\hline
Vehicles & \nolinkurl{/api/v1/vehicles} & Có & CRUD xe, gán thiết bị \\\hline
Khách hàng & \nolinkurl{/api/v1/customers} & Có & CRUD khách hàng \\\hline
Chuyến đi & \nolinkurl{/api/v1/trips} & Có & CRUD chuyến đi, lịch sử \\\hline
Cảnh báo & \nolinkurl{/api/v1/alerts} & Có & Quản lý cảnh báo \\\hline
Vùng địa lý & \nolinkurl{/api/v1/geofences} & Có & CRUD vùng địa lý \\\hline
Firmware & \nolinkurl{/api/v1/firmware} & Có (admin) & OTA firmware
management \\\hline
\end{longtable}
}


Riêng nhóm IoT Data (\texttt{/api/v1/iot}) không đi theo luồng đăng nhập
người dùng thông thường. Các thiết bị gửi dữ liệu trực tiếp bằng MQTT
hoặc HTTP fallback, nên lớp bảo vệ chính nằm ở cơ chế nhận diện thiết bị
và ACL của EMQX thay vì ở phiên đăng nhập web.

\subparagraph{f) WebSocket (Socket.IO)}\label{f-websocket-socketio}

Hệ thống sử dụng Socket.IO 4.8 để đẩy các thay đổi mới ra giao diện web
ngay khi chúng xảy ra. Thay vì để giao diện quản trị hỏi lại API liên
tục, API phía máy chủ chỉ phát sự kiện đúng vào nhóm chức năng liên
quan. Vì thế, Socket.IO được tách thành 5 namespace, mỗi namespace phục
vụ một nhóm nhu cầu cụ thể:

\begin{Shaded}
\begin{Highlighting}[]
\CommentTok{// 5 namespaces với xác thực}
\SpecialStringTok{/dashboard      // Cập nhật thống kê tổng quan }\SpecialCharTok{(}\SpecialStringTok{yêu cầu xác thực}\SpecialCharTok{)}
\SpecialStringTok{/devices}        \CommentTok{// Vị trí, trạng thái, command ack theo thiết bị}
\SpecialStringTok{/firmware       // Sự kiện gán/cập}\NormalTok{ nhật tác vụ }\FunctionTok{OTA}\NormalTok{ (yêu cầu xác thực)}
\OperatorTok{/}\NormalTok{exports        }\CommentTok{// Tiến trình xuất báo cáo (yêu cầu xác thực)}
\SpecialStringTok{/notifications  // Thông báo đẩy }\SpecialCharTok{(}\SpecialStringTok{yêu cầu xác thực}\SpecialCharTok{)}
\end{Highlighting}
\end{Shaded}

Các sự kiện chính được phát từ máy chủ đến trình duyệt bao gồm:
\texttt{device:status}, \texttt{device:position},
\texttt{device:session\_start}, \texttt{device:session\_end},
\texttt{command:ack}, \texttt{stats:update}, \texttt{alert:new},
\texttt{geofence:enter}, \texttt{geofence:exit}, \texttt{export:ready}
và \texttt{firmware:assignment}. Trong triển khai hiện tại,
\texttt{firmware:assignment} phản ánh bước khởi tạo hoặc cập nhật tác vụ
OTA từ API phía máy chủ; tiến trình thực thi chi tiết vẫn được theo dõi chủ yếu
qua \texttt{firmware\_update\_log} và lớp log quan trắc.

\subparagraph{g) Cấu hình Docker cho
Backend}\label{g-cux1ea5u-huxecnh-docker-cho-backend}

\begin{Shaded}
\begin{Highlighting}[]
\CommentTok{\# Tracking\_Backend/docker{-}compose.yml}
\FunctionTok{version}\KeywordTok{:}\AttributeTok{ }\StringTok{"3.8"}
\FunctionTok{services}\KeywordTok{:}
\AttributeTok{  }\FunctionTok{backend}\KeywordTok{:}
\AttributeTok{    }\FunctionTok{build}\KeywordTok{:}
\AttributeTok{      }\FunctionTok{context}\KeywordTok{:}\AttributeTok{ .}
\AttributeTok{      }\FunctionTok{dockerfile}\KeywordTok{:}\AttributeTok{ Dockerfile}
\AttributeTok{    }\FunctionTok{container\_name}\KeywordTok{:}\AttributeTok{ tracking{-}backend}
\AttributeTok{    }\FunctionTok{ports}\KeywordTok{:}
\AttributeTok{      }\KeywordTok{{-}}\AttributeTok{ }\StringTok{"0.0.0.0:4000:4000"}
\AttributeTok{    }\FunctionTok{env\_file}\KeywordTok{:}\AttributeTok{ .env}
\AttributeTok{    }\FunctionTok{restart}\KeywordTok{:}\AttributeTok{ unless{-}stopped}
\AttributeTok{    }\FunctionTok{networks}\KeywordTok{:}
\AttributeTok{      }\KeywordTok{{-}}\AttributeTok{ tracking{-}network}
\AttributeTok{    }\FunctionTok{deploy}\KeywordTok{:}
\AttributeTok{      }\FunctionTok{resources}\KeywordTok{:}
\AttributeTok{        }\FunctionTok{limits}\KeywordTok{:}
\AttributeTok{          }\FunctionTok{memory}\KeywordTok{:}\AttributeTok{ 512M}
\AttributeTok{          }\FunctionTok{cpus}\KeywordTok{:}\AttributeTok{ }\StringTok{"1.0"}

\FunctionTok{networks}\KeywordTok{:}
\AttributeTok{  }\FunctionTok{tracking{-}network}\KeywordTok{:}
\AttributeTok{    }\FunctionTok{external}\KeywordTok{:}\AttributeTok{ }\CharTok{true}
\end{Highlighting}
\end{Shaded}

\subparagraph{h) Cấu trúc phản hồi API thống
nhất}\label{h-cux1ea5u-truxfac-response-api-thux1ed1ng-nhux1ea5t}

Tất cả các API endpoint đều trả về phản hồi theo cấu trúc thống nhất,
giúp giao diện web xử lý dữ liệu dễ dàng và nhất quán. Trong trường hợp lỗi,
phản hồi có cấu trúc:

\begin{Shaded}
\begin{Highlighting}[]
\FunctionTok{\{}
  \DataTypeTok{"error"}\FunctionTok{:} \FunctionTok{\{}
    \DataTypeTok{"code"}\FunctionTok{:} \StringTok{"ERROR\_CODE"}\FunctionTok{,}
    \DataTypeTok{"message"}\FunctionTok{:} \StringTok{"Mô tả lỗi chi tiết"}\FunctionTok{,}
    \DataTypeTok{"status"}\FunctionTok{:} \DecValTok{400}\FunctionTok{,}
    \DataTypeTok{"path"}\FunctionTok{:} \StringTok{"/api/v1/endpoint"}\FunctionTok{,}
    \DataTypeTok{"details"}\FunctionTok{:} \KeywordTok{null}\FunctionTok{,}
    \DataTypeTok{"traceId"}\FunctionTok{:} \StringTok{"uuid{-}request{-}id"}
  \FunctionTok{\},}
  \DataTypeTok{"timestamp"}\FunctionTok{:} \StringTok{"2025–01–01T00:00:00.000Z"}
\FunctionTok{\}}
\end{Highlighting}
\end{Shaded}

Trường \texttt{traceId} chính là mã truy vết từng yêu cầu được tạo bởi
middleware, cho phép đối chiếu lỗi giữa log của giao diện web và log của API phía máy chủ.

\begin{center}\rule{0.5\linewidth}{0.5pt}\end{center}

\subsubsection{4.2.5. Triển khai giao diện web quản trị}\label{425-triux1ec3n-khai-frontend-dashboard-frontend-dashboard-implementation}

Sau khi dữ liệu đã được tiếp nhận và xử lý ở tầng máy chủ, bước tiếp
theo là biến chúng thành giao diện đủ rõ ràng để người vận hành có thể
quan sát và ra quyết định. Vì thế, phần giao diện web không chỉ được
đánh giá ở khía cạnh công nghệ, mà còn ở khả năng tổ chức thông tin theo
đúng nhu cầu sử dụng thực tế.

\subparagraph{a) Lựa chọn công
nghệ}\label{a-lux1ef1a-chux1ecdn-cuxf4ng-nghux1ec7-2}

Giao diện web quản trị được xây dựng trên nền tảng Next.js 15 với React 19.
Điều quan trọng không phải tên công nghệ, mà là giao diện phải đáp ứng ba
nhu cầu rất thực tế: tải trang đủ nhanh, cập nhật thời gian thực ổn định
và vẫn giữ cấu trúc mã nguồn gọn khi tính năng quản trị tăng lên. Nếu
phần máy chủ trả lời câu hỏi ``cần cung cấp dữ liệu gì'', thì giao diện
web trả lời câu hỏi ``nên trình bày ra sao để người vận hành hiểu nhanh
và thao tác ít nhất''.

Có thể nhìn toàn bộ phần giao diện web thành ba nhóm vai trò. Nhóm thứ
nhất dựng khung ứng dụng và điều hướng trang. Nhóm thứ hai giữ trạng
thái giao diện và đồng bộ dữ liệu với máy chủ. Nhóm thứ ba tạo ra những
gì người vận hành nhìn thấy trực tiếp như bản đồ, biểu đồ và cập nhật
thời gian thực. Cách chia này giúp người đọc hiểu mỗi công cụ xuất hiện
để giải quyết một nhu cầu rõ ràng, chứ không phải để làm dày thêm danh
sách công nghệ.

\textbf{Thành phần công nghệ của giao diện web:}


{\thesissettableid{4.2.5A}{tab-ch4-tong-hop-thong-tin-ve-cong-nghe-phien-ban-vai-tro}
\def\LTcaptype{table}
\begin{longtable}{|L{0.305\linewidth}|L{0.160\linewidth}|L{0.500\linewidth}|}
\caption{Tổng hợp thông tin về Công nghệ, Phiên bản, Vai trò}\label{tab:ch4-tong-hop-thong-tin-ve-cong-nghe-phien-ban-vai-tro}\\
\hline \textbf{Công nghệ} & \textbf{Phiên bản} & \textbf{Vai trò} \\
\\\hline
\endfirsthead
\caption[]{Tổng hợp thông tin về Công nghệ, Phiên bản, Vai trò (tiếp theo)}\\
\hline \textbf{Công nghệ} & \textbf{Phiên bản} & \textbf{Vai trò} \\
\\\hline
\endhead
\\\hline
\endfoot
\endlastfoot
Next.js & 15 & Khung ứng dụng React với App Router, hỗ trợ chia trang và
bố cục theo thư mục \\\hline
React & 19 & Thư viện xây dựng giao diện \\\hline
TypeScript & 5.x & Ngôn ngữ lập trình có kiểm tra kiểu dữ liệu \\\hline
Tailwind CSS & 4 & Cách viết CSS bằng các lớp tiện ích dựng sẵn \\\hline
Zustand & - & Quản lý trạng thái dùng chung của giao diện \\\hline
TanStack Query & - & Quản lý dữ liệu lấy từ máy chủ, gồm lưu tạm và làm
mới lại khi cần \\\hline
Socket.IO Client & - & Kết nối kênh cập nhật thời gian thực \\\hline
Leaflet & - & Bản đồ tương tác \\\hline
ECharts & - & Biểu đồ và trực quan hóa dữ liệu \\\hline
Radix UI + shadcn/ui & - & Bộ thành phần giao diện có sẵn, hỗ trợ dùng
bàn phím tốt và dễ ghép lại theo nhu cầu \\\hline
\end{longtable}
}


\subparagraph{b) Kiến trúc
tách theo từng nhóm tính năng}\label{b-kiux1ebfn-truxfac-feature-sliced}

Giao diện web được tổ chức theo kiến trúc chia mã nguồn theo từng nhóm tính năng, nghĩa là mỗi tính
năng người dùng nhìn thấy sẽ có một khu vực mã nguồn tương đối riêng.
Trong mỗi khu vực đó, các thành phần UI, hook, kiểu dữ liệu và tiện ích
liên quan được đặt gần nhau. Nhờ vậy, cấu trúc mã nguồn phản chiếu khá
sát cấu trúc nghiệp vụ, nên ngay cả người không đi sâu vào code cũng có
thể hình dung hệ thống web được tổ chức như thế nào.

\begin{itemize}
\tightlist
\item
  \textbf{Đặt gần theo tính năng}: Phần mã liên quan đến một chức năng nằm gần nhau, dễ
  tìm và bảo trì
\item
  \textbf{Độc lập}: Mỗi nhóm tính năng có thể phát triển và kiểm thử độc lập
\item
  \textbf{Tái sử dụng}: Các thành phần dùng chung nằm trong
  \texttt{components/}, chỉ phần chuyên biệt mới nằm trong
  \texttt{features/}
\end{itemize}

\begin{itemize}
\tightlist
\item
  \texttt{app/}: khung định tuyến trang và các khu vực quản trị được bảo vệ.
\item
  \texttt{components/}: các thành phần dùng lại cho giao diện, bố cục, biểu mẫu, bản đồ,
  biểu đồ và lớp cung cấp ngữ cảnh dùng chung.
\item
  \texttt{features/}: các nhóm nghiệp vụ theo FSD
  (\texttt{vehicles}, \texttt{customers}, \texttt{trips},
  \texttt{alerts}, \texttt{map}).
\item
  \texttt{lib/}: các lớp dùng chung như nơi gọi API, kết nối thời gian thực,
  nơi giữ trạng thái, hằng số, tiện ích và hook dùng lại.
\end{itemize}


{\thesissetfigureid{4.22}{fig-ch4-cau-truc-thu-muc-frontend-theo-kien-truc-feature-sliced}
\begin{figure}[htbp]
\centering
\pandocbounded{\safeincludesvg{./assets/figures/09-chuong-4-trien-khai-cloud-hinh-4-19.svg}}
\caption{Cấu trúc thư mục giao diện web theo kiến trúc tách theo từng nhóm tính năng}
\label{fig:ch4-cau-truc-thu-muc-frontend-theo-kien-truc-feature-sliced}
\end{figure}
}


\subparagraph{c) Các trang chính và tính
năng}\label{c-cuxe1c-trang-chuxednh-vuxe0-tuxednh-nux103ng}

Để giao diện quản trị thực sự hữu ích trong vận hành, giao diện không được dừng ở
việc ``có nhiều trang'', mà mỗi trang phải trả lời được một nhu cầu quản
lý cụ thể. Vì vậy, hệ thống cung cấp các nhóm trang chính sau:

\textbf{Trang tổng quan (\texttt{/dashboard}):} Hiển thị các
thẻ thống kê nhanh bao gồm tổng số xe, số chuyến đi trong ngày,
số cảnh báo chưa xử lý, và số vi phạm. Ngoài ra còn hiển thị bảng cảnh
báo gần đây, bản đồ mini với các xe đang hoạt động, danh sách chuyến đi
gần đây, và biểu đồ thống kê (số chuyến đi theo ngày, cảnh báo theo
loại).


{\thesissetfigureid{4.23}{fig-ch4-giao-dien-trang-dashboard-tong-quan}
\begin{figure}[htbp]
\centering
\pandocbounded{\safeincludesvg{./assets/figures/09-chuong-4-trien-khai-cloud-hinh-4-20.svg}}
\caption{Giao diện trang tổng quan hệ thống}
\label{fig:ch4-giao-dien-trang-dashboard-tong-quan}
\end{figure}
}


\textbf{Trang quản lý xe (\texttt{/dashboard/vehicles}):} Giao diện dạng
bảng dữ liệu với các cột: biển số xe, hãng xe/model, trạng
thái, thiết bị gắn kèm, lần cuối thấy, và các hành động. Hỗ trợ lọc theo
trạng thái, loại xe, và tìm kiếm. Trang chi tiết xe hiển thị thông tin
xe, vị trí hiện tại trên bản đồ, trạng thái thiết bị, cảnh báo đang hoạt
động, chuyến đi gần đây, và lịch sử bảo trì.


{\thesissetfigureid{4.24}{fig-ch4-giao-dien-trang-quan-ly-xe}
\begin{figure}[htbp]
\centering
\pandocbounded{\safeincludesvg{./assets/figures/09-chuong-4-trien-khai-cloud-hinh-4-21.svg}}
\caption{Giao diện trang quản lý xe}
\label{fig:ch4-giao-dien-trang-quan-ly-xe}
\end{figure}
}


\textbf{Trang bản đồ thời gian thực (\texttt{/dashboard/map}):} Hiển thị
tất cả các xe trên bản đồ Leaflet với vị trí cập nhật thời gian thực
thông qua WebSocket. Hỗ trợ: hiển thị điểm đánh dấu cho từng xe với ô
thông tin trạng thái, lọc theo xe hoặc trạng thái, hiển thị vùng cho phép đã thiết lập, và phát lại hành trình cũ.


{\thesissetfigureid{4.25}{fig-ch4-giao-dien-ban-o-thoi-gian-thuc-voi-vi-tri-cac-xe}
\begin{figure}[htbp]
\centering
\pandocbounded{\safeincludesvg{./assets/figures/09-chuong-4-trien-khai-cloud-hinh-4-22.svg}}
\caption{Giao diện bản đồ thời gian thực với vị trí các xe}
\label{fig:ch4-giao-dien-ban-o-thoi-gian-thuc-voi-vi-tri-cac-xe}
\end{figure}
}


\textbf{Trang quản lý cảnh báo (\texttt{/dashboard/alerts}):} Bảng dữ
liệu với nhãn mức độ nghiêm trọng, hỗ trợ lọc theo loại, mức
độ, xe, và khoảng thời gian. Các hành động hàng loạt gồm xác nhận đã
nhận và đánh dấu đã xử lý. Trang chi tiết cảnh báo hiển thị
thông tin cảnh báo, vị trí trên bản đồ, xe và khách hàng liên quan.


{\thesissetfigureid{4.26}{fig-ch4-giao-dien-trang-quan-ly-canh-bao}
\begin{figure}[htbp]
\centering
\pandocbounded{\safeincludesvg{./assets/figures/09-chuong-4-trien-khai-cloud-hinh-4-23.svg}}
\caption{Giao diện trang quản lý cảnh báo}
\label{fig:ch4-giao-dien-trang-quan-ly-canh-bao}
\end{figure}
}


\textbf{Trang quản lý vùng cho phép (\texttt{/dashboard/geofences}):} Cho
phép tạo và quản lý các vùng cho phép trên bản đồ. Khi xe ra
khỏi hoặc vào vùng địa lý đã định nghĩa, hệ thống sẽ tự động tạo cảnh
báo.

\textbf{Trang cài đặt thông báo (\texttt{/dashboard/notifications}):}
Cho phép cấu hình kết nối Telegram bot, cài đặt email, lựa chọn loại
cảnh báo muốn nhận, lọc theo mức độ nghiêm trọng và xe cụ thể. Có nút
gửi thông báo thử nghiệm để kiểm tra cấu hình.

\subparagraph{d) Quản lý trạng thái và dữ
liệu}\label{d-quux1ea3n-luxfd-trux1ea1ng-thuxe1i-vuxe0-dux1eef-liux1ec7u}

Với giao diện thời gian thực, mã nguồn giao diện web sẽ rất nhanh rối nếu
trộn lẫn dữ liệu từ máy chủ, trạng thái giao diện và luồng sự kiện mới
vào cùng một chỗ. Vì vậy, hệ thống chia rõ từng loại dữ liệu và dùng kết
hợp nhiều công cụ quản lý trạng thái:

\begin{itemize}
\tightlist
\item
  \textbf{Zustand}: Giữ các trạng thái ``thuộc về giao diện'' như thông
  tin đăng nhập hiện tại, sidebar hay theme. Token xác thực chỉ nằm trong
  bộ nhớ tạm thời, không lưu vào \texttt{localStorage} để giảm rủi ro mã độc chèn
  vào trang rồi lấy cắp thông tin.
\item
  \textbf{TanStack Query}: Quản lý dữ liệu lấy từ máy chủ như danh sách
  xe, cảnh báo hay lịch sử. Công cụ này lo cache, làm mới và đồng bộ lại
  dữ liệu khi cần. Khi WebSocket đã hoạt động, cơ chế polling sẽ giảm
  xuống để tránh tải trùng lặp.
\item
  \textbf{Socket.IO Client}: Nhận các cập nhật thời gian thực do API phía máy chủ
  phát ra. Giao diện web không nối trực tiếp tới MQTT broker; mọi dữ liệu
  thời gian thực đều được API phía máy chủ chuẩn hóa rồi mới đẩy ra giao diện.
\end{itemize}

\subparagraph{e) Các giai đoạn triển
khai}\label{e-cuxe1c-giai-ux111oux1ea1n-triux1ec3n-khai}

Giao diện web được triển khai theo 3 giai đoạn chính (giai đoạn 4 là giai đoạn mở
rộng trong tương lai):


{\thesissettableid{4.9}{tab-ch4-cac-giai-oan-trien-khai-frontend}
\def\LTcaptype{table}
\begin{longtable}{|L{0.224\linewidth}|L{0.580\linewidth}|L{0.160\linewidth}|}
\caption{Các giai đoạn triển khai giao diện web}\label{tab:ch4-cac-giai-oan-trien-khai-frontend}\\
\hline \textbf{Giai đoạn} & \textbf{Nội dung} & \textbf{Trạng thái} \\
\\\hline
\endfirsthead
\caption[]{Các giai đoạn triển khai giao diện web (tiếp theo)}\\
\hline \textbf{Giai đoạn} & \textbf{Nội dung} & \textbf{Trạng thái} \\
\\\hline
\endhead
\\\hline
\endfoot
\endlastfoot
Giai đoạn 1: Nền tảng cốt lõi & Khởi tạo dự án, Tailwind + shadcn/ui,
giao diện sáng/tối, bố cục, xác thực, lớp gọi API và React Query & Hoàn thành \\\hline
Giai đoạn 2: Tính năng cốt lõi & Trang tổng quan, CRUD xe, CRUD khách hàng, quản lý chuyến đi,
quản lý cảnh báo, quản lý thiết bị, bản đồ thời gian thực & Hoàn
thành \\\hline
Giai đoạn 3: Tính năng nâng cao & Quản lý vùng cho phép, bảo trì,
cài đặt thông báo (Telegram + Email), quản lý người dùng, thiết lập hệ thống &
Hoàn thành \\\hline
Giai đoạn 4: Đặt xe (tương lai) & Đặt xe, hợp đồng, thanh toán, báo cáo
hư hỏng, đánh giá & Chưa triển khai \\\hline
\end{longtable}
}


\subparagraph{f) Cấu hình Docker cho
Frontend}\label{f-cux1ea5u-huxecnh-docker-cho-frontend}

Giao diện web sử dụng Dockerfile nhiều giai đoạn: một giai đoạn để biên
dịch và một giai đoạn để chạy, nhằm tối ưu kích thước image như sau:

\begin{Shaded}
\begin{Highlighting}[]
\CommentTok{\# Tracking\_Frontend/Dockerfile}
\KeywordTok{FROM}\NormalTok{ node:20{-}alpine }\KeywordTok{AS}\NormalTok{ builder}
\KeywordTok{WORKDIR}\NormalTok{ /app}
\KeywordTok{COPY}\NormalTok{ package*.json ./}
\KeywordTok{RUN} \ExtensionTok{npm}\NormalTok{ ci}
\KeywordTok{COPY}\NormalTok{ . .}
\KeywordTok{RUN} \ExtensionTok{npm}\NormalTok{ run build}

\CommentTok{\# Giai đoạn 2: Chạy ứng dụng (production)}
\KeywordTok{FROM}\NormalTok{ node:20{-}alpine }\KeywordTok{AS}\NormalTok{ runner}
\KeywordTok{WORKDIR}\NormalTok{ /app}
\KeywordTok{ENV}\NormalTok{ NODE\_ENV=production}
\KeywordTok{ENV}\NormalTok{ HOSTNAME=0.0.0.0}
\KeywordTok{ENV}\NormalTok{ PORT=4001}
\KeywordTok{COPY} \OperatorTok{{-}{-}from=builder}\NormalTok{ /app/public ./public}
\KeywordTok{COPY} \OperatorTok{{-}{-}from=builder}\NormalTok{ /app/.next/standalone ./}
\KeywordTok{COPY} \OperatorTok{{-}{-}from=builder}\NormalTok{ /app/.next/static ./.next/static}
\KeywordTok{RUN} \ExtensionTok{addgroup} \AttributeTok{{-}{-}system} \AttributeTok{{-}{-}gid}\NormalTok{ 1001 nodejs }\KeywordTok{\&\&} \ExtensionTok{adduser} \AttributeTok{{-}{-}system} \AttributeTok{{-}{-}uid}\NormalTok{ 1001 nextjs}
\KeywordTok{USER}\NormalTok{ nextjs}
\KeywordTok{EXPOSE}\NormalTok{ 4001}
\KeywordTok{HEALTHCHECK} \OperatorTok{{-}{-}interval=30s} \OperatorTok{{-}{-}timeout=5s} \OperatorTok{{-}{-}start{-}period=15s} \OperatorTok{{-}{-}retries=3}\NormalTok{ \textbackslash{}}
  \KeywordTok{CMD} \FunctionTok{wget} \AttributeTok{{-}{-}spider}\NormalTok{ http://127.0.0.1:4001 }\KeywordTok{||} \BuiltInTok{exit}\NormalTok{ 1}
\KeywordTok{CMD}\NormalTok{ [}\StringTok{"node"}\NormalTok{, }\StringTok{"server.js"}\NormalTok{]}
\end{Highlighting}
\end{Shaded}

Dockerfile sử dụng 2 giai đoạn: \texttt{builder} (cài đặt dependency và
build ứng dụng), và \texttt{runner} (chạy ở môi trường vận hành thật). Giai đoạn
\texttt{runner} chỉ sao chép các file cần thiết từ giai đoạn
\texttt{builder}, giúp giảm đáng kể kích thước image cuối cùng. Ứng dụng
chạy với user \texttt{nextjs} (không phải root) để tăng tính bảo mật.

\begin{center}\rule{0.5\linewidth}{0.5pt}\end{center}

\subsubsection{4.2.6. Cấu hình giám sát hệ thống}\label{426-cux1ea5u-huxecnh-giuxe1m-suxe1t-hux1ec7-thux1ed1ng-monitoring--observability}

Khi hệ thống đã có đủ các lớp phần cứng, firmware, hạ tầng máy chủ và
giao diện web,
câu hỏi tiếp theo là làm sao biết hệ thống đang khỏe hay đang có vấn đề.
Đó là lý do lớp giám sát không được xem như phần phụ, mà là điều kiện để
vận hành ổn định và xử lý sự cố có cơ sở.

\subparagraph{a) Tổng quan chiến lược giám
sát}\label{a-tux1ed5ng-quan-chiux1ebfn-lux1b0ux1ee3c-giuxe1m-suxe1t}

Hệ thống giám sát được tổ chức theo ba câu hỏi cơ bản: hệ thống đang
nặng đến đâu, chuyện gì vừa xảy ra, và một yêu cầu đã đi qua những đâu.
Ba câu hỏi đó tương ứng với chỉ số vận hành, nhật ký và dấu vết truy vết yêu cầu.

Cách trình bày này giúp người đọc hiểu khả năng quan sát hệ thống không
phải là thêm thật nhiều công cụ, mà là thêm khả năng quan sát để tìm
nguyên nhân sự cố có cơ sở. Trong phạm vi đồ án, phần truy vết chưa dùng
đầy đủ cơ chế truy vết phân tán như OpenTelemetry để nối dấu vết qua
nhiều dịch vụ; thay vào đó hệ thống dùng mã nhận diện từng yêu cầu để
nối log giữa giao diện web, API phía máy chủ và các dịch vụ liên quan.
Cách làm này đơn giản hơn nhưng vẫn đủ cho quy mô triển khai hiện tại.


{\thesissettableid{4.10}{tab-ch4-ba-tru-cot-giam-sat-he-thong}
\def\LTcaptype{table}
\begin{longtable}{|L{0.115\linewidth}|L{0.209\linewidth}|L{0.276\linewidth}|L{0.354\linewidth}|}
\caption{Ba trụ cột giám sát hệ thống}\label{tab:ch4-ba-tru-cot-giam-sat-he-thong}\\
\hline \textbf{Trụ cột} & \textbf{Công cụ} & \textbf{Mục đích} & \textbf{Dữ liệu lưu trữ} \\
\\\hline
\endfirsthead
\caption[]{Ba trụ cột giám sát hệ thống (tiếp theo)}\\
\hline \textbf{Trụ cột} & \textbf{Công cụ} & \textbf{Mục đích} & \textbf{Dữ liệu lưu trữ} \\
\\\hline
\endhead
\\\hline
\endfoot
\endlastfoot
Chỉ số vận hành & VictoriaMetrics + bộ xuất số liệu theo chuẩn Prometheus & Chỉ số
hiệu năng, tài nguyên, nghiệp vụ & CPU, bộ nhớ, số yêu cầu/giây, độ trễ,
số thiết bị online \\\hline
Nhật ký & VictoriaLogs & Nhật ký tập trung từ tất cả dịch vụ & Nhật ký
yêu cầu, nhật ký lỗi, sự kiện MQTT, sự kiện thiết bị \\\hline
Truy vết yêu cầu & Mã nhận diện từng yêu cầu & Lần theo một yêu cầu xuyên suốt hệ thống &
Mã \texttt{X-Request-ID} trên mỗi yêu cầu \\\hline
Theo dõi lỗi & Sentry & Theo dõi lỗi và ngoại lệ & Vết gọi hàm, ngữ cảnh người dùng,
chuỗi sự kiện ngay trước khi lỗi xảy ra \\\hline
\end{longtable}
}


\subparagraph{b) VictoriaMetrics - Giám sát chỉ số}\label{b-victoriametrics---giuxe1m-suxe1t-chux1ec9-sux1ed1-metrics}

VictoriaMetrics đóng hai vai trò trong hệ thống:

\begin{enumerate}
\def\labelenumi{\arabic{enumi}.}
\tightlist
\item
  \textbf{Lưu trữ dữ liệu vận hành của thiết bị}: Đây là các dữ liệu vận hành
  phát sinh liên tục như vị trí GPS, tốc độ, điện áp nguồn, mức rung,
  trạng thái bật/tắt máy và thời gian chạy liên tục. Nguồn dữ liệu này đi từ MQTT Bridge.
\item
  \textbf{Lưu trữ chỉ số sức khỏe hệ thống}: Đây là các chỉ số cho biết
  API phía máy chủ và hạ tầng đang hoạt động ra sao, ví dụ số yêu cầu HTTP, độ
  trễ xử lý, số kết nối WebSocket hay số bản tin MQTT. Các dịch vụ gửi
  nhóm chỉ số này qua thư viện \texttt{prom-client}.
\end{enumerate}

Backend dùng middleware \texttt{httpMetricsMiddleware} để tự động ghi
nhận các chỉ số HTTP trong lúc chạy. Endpoint \texttt{/metrics} đóng vai
trò như ``đầu ra kỹ thuật'' để Grafana đọc các số liệu này thông qua
VictoriaMetrics, từ đó dựng biểu đồ giám sát theo thời gian thực.

\subparagraph{c) VictoriaLogs - Nhật ký tập trung}\label{c-victorialogs---nhux1eadt-kuxfd-tux1eadp-trung-centralized-logging}

VictoriaLogs là hệ thống nhật ký tập trung, thu thập log từ tất cả các
dịch vụ trong hệ thống. API phía máy chủ sử dụng \texttt{victorialogs-transport}
của Winston để gửi log trực tiếp tới VictoriaLogs. Mỗi bản ghi log bao
gồm:

\begin{itemize}
\tightlist
\item
  \textbf{Timestamp}: Thời điểm xảy ra sự kiện
\item
  \textbf{Level}: Mức độ nghiêm trọng (info, warn, error)
\item
  \textbf{Service}: Tên dịch vụ phát sinh log (API phía máy chủ, mqtt-bridge,
  giao diện web)
\item
  \textbf{Request-ID}: Mã nhận diện yêu cầu để truy vết
\item
  \textbf{Message}: Nội dung nhật ký
\item
  \textbf{Metadata}: Dữ liệu bổ sung (device\_id, user\_id, vết gọi hàm khi lỗi)
\end{itemize}

MQTT Bridge cũng có logger riêng, gửi log về VictoriaLogs với trường
\texttt{service:\ mqtt-bridge}, cho phép lọc và phân tích log theo từng
dịch vụ.

\subparagraph{d) Grafana - Màn hình giám sát trực
quan}\label{d-grafana---dashboard-giuxe1m-suxe1t-trux1ef1c-quan}

Grafana được sử dụng làm lớp trực quan hóa, kết
nối tới cả VictoriaMetrics và VictoriaLogs để hiển thị các màn hình
giám sát. Các màn hình chính bao gồm:

\begin{itemize}
\tightlist
\item
  \textbf{Tổng quan hệ thống}: Số thiết
  bị online, số yêu cầu/giây, phần trăm lỗi, mức sử dụng CPU và bộ nhớ
  của các container.
\item
  \textbf{Dữ liệu vận hành IoT}: Dữ liệu từ các thiết bị IoT --- bản đồ nhiệt
  (heatmap) vị trí xe, biểu đồ tốc độ, biểu đồ mức nhiên liệu theo thời
  gian.
\item
  \textbf{Thống kê MQTT}: Số lượng bản tin/giây, số kết
  nối đang hoạt động, độ trễ trung bình, tỷ lệ bản tin thất bại.
\item
  \textbf{Hiệu năng ứng dụng}: Phân phối
  thời gian xử lý yêu cầu, các endpoint chậm nhất, tỷ lệ
  lá»—i theo endpoint.
\end{itemize}


{\thesissetfigureid{4.27}{fig-ch4-grafana-dashboard-hien-thi-tong-quan-hieu-nang-he-thong}
\begin{figure}[htbp]
\centering
\pandocbounded{\safeincludesvg{./assets/figures/09-chuong-4-trien-khai-cloud-hinh-4-24.svg}}
\caption{Màn hình Grafana hiển thị tổng quan hiệu năng hệ thống}
\label{fig:ch4-grafana-dashboard-hien-thi-tong-quan-hieu-nang-he-thong}
\end{figure}
}


\subparagraph{e) Giao diện quản trị EMQX - Giám sát broker MQTT}\label{e-emqx-dashboard---giuxe1m-suxe1t-mqtt-broker}

EMQX cung cấp giao diện quản trị tích hợp (port 18083) cho phép giám sát trực
tiếp trạng thái của broker MQTT:

\begin{itemize}
\tightlist
\item
  \textbf{Connections}: Số lượng kết nối hiện tại, lịch sử kết nối/ngắt
  kết nối
\item
  \textbf{Subscriptions}: Danh sách topic và bên đang đăng ký nhận
\item
  \textbf{Messages}: Thống kê bản tin (gửi lên, chuyển đi thành công, bị loại bỏ)
\item
  \textbf{Rules Engine}: Trạng thái các luật xử lý dữ liệu
\item
  \textbf{ACL}: Cấu hình quyền truy cập cho từng thiết bị
\end{itemize}


{\thesissetfigureid{4.28}{fig-ch4-emqx-dashboard-hien-thi-trang-thai-ket-noi-thiet-bi}
\begin{figure}[htbp]
\centering
\pandocbounded{\safeincludesvg{./assets/figures/09-chuong-4-trien-khai-cloud-hinh-4-25.svg}}
\caption{Giao diện quản trị EMQX hiển thị trạng thái kết nối thiết bị}
\label{fig:ch4-emqx-dashboard-hien-thi-trang-thai-ket-noi-thiet-bi}
\end{figure}
}


\subparagraph{f) Sentry - Theo dõi lỗi ứng dụng}\label{f-sentry---theo-duxf5i-lux1ed7i-error-tracking}

API phía máy chủ tích hợp Sentry thông qua hai middleware:
\texttt{sentryRequestHandler} (đặt đầu tiên trong middleware chain) và
\texttt{sentryErrorHandler} (đặt cuối cùng). Sentry tự động thu thập:

\begin{itemize}
\tightlist
\item
  \textbf{Chi tiết lỗi}: Vết gọi hàm đầy đủ, biến môi trường, thông
  tin yêu cầu
\item
  \textbf{Chuỗi sự kiện trước lỗi}: Các sự kiện xảy ra trước khi lỗi xuất hiện
  (truy vấn cơ sở dữ liệu, gọi HTTP, sự kiện MQTT)
\item
  \textbf{Hiệu năng}: Thời gian xử lý của mỗi yêu cầu
\item
  \textbf{Theo dõi theo phiên bản}: Lỗi thuộc phiên bản triển khai nào
\end{itemize}

Sentry đặc biệt hữu ích trong môi trường vận hành thật khi không thể truy
cập trực tiếp vào log của máy chủ.

\subparagraph{g) Liên kết truy vết bằng mã nhận diện yêu cầu}\label{g-request-id-correlation}

Mỗi yêu cầu HTTP đến API phía máy chủ đều được gán một mã truy vết duy
nhất (UUID v4) bởi middleware \texttt{requestIdMiddleware}. Mã này được:

\begin{itemize}
\tightlist
\item
  Trả về trong header \texttt{X-Request-ID} của phản hồi
\item
  Đính kèm vào mỗi bản ghi log liên quan đến yêu cầu đó
\item
  Truyền xuyên suốt qua các lớp xử lý (controller -\textgreater{}
  service -\textgreater{} repository)
\item
  Sử dụng làm \texttt{traceId} trong phản hồi lỗi
\end{itemize}

Cơ chế này cho phép lần theo một yêu cầu cụ thể từ giao diện web, qua
API phía máy chủ, đến cơ sở dữ liệu và MQTT Bridge, giúp việc dò lỗi
và phân tích sự cố trong hệ thống nhiều dịch vụ hiệu quả hơn.

\begin{center}\rule{0.5\linewidth}{0.5pt}\end{center}

\subsubsection{4.2.7. Checklist làm cứng vận hành trước khi đưa vào môi
trường thật}\label{427-checklist-hardening-trux1b0ux1edbc-khi-vux1eadn-huxe0nh-production}

Để chuyển hệ thống từ môi trường thử nghiệm sang môi trường vận hành
thực tế, nhóm triển khai áp dụng checklist làm cứng vận hành theo các mức ưu
tiên sau (tham chiếu từ bộ tài liệu cải tiến hệ thống):


{\thesissettableid{4.11}{tab-ch4-checklist-hardening-production-theo-muc-uu-tien}
\def\LTcaptype{table}
\begin{longtable}{|L{0.125\linewidth}|L{0.396\linewidth}|L{0.115\linewidth}|L{0.318\linewidth}|}
\caption{Checklist làm cứng vận hành theo mức ưu tiên}\label{tab:ch4-checklist-hardening-production-theo-muc-uu-tien}\\
\hline \textbf{Nhóm} & \textbf{Hạng mục bắt buộc} & \textbf{Mức ưu tiên} & \textbf{Tiêu chí nghiệm thu} \\
\\\hline
\endfirsthead
\caption[]{Checklist làm cứng vận hành theo mức ưu tiên (tiếp theo)}\\
\hline \textbf{Nhóm} & \textbf{Hạng mục bắt buộc} & \textbf{Mức ưu tiên} & \textbf{Tiêu chí nghiệm thu} \\
\\\hline
\endhead
\\\hline
\endfoot
\endlastfoot
Bảo mật truyền thông & Bật TLS cho HTTPS và MQTT (\texttt{8883}), bắt
buộc chứng chỉ hợp lệ & Cao & 100\% kết nối ngoài đi qua TLS, không còn
điểm truy cập thuần văn bản \\\hline
Xác thực thiết bị & Mỗi thiết bị có thông tin xác thực riêng
(chứng chỉ số hoặc khóa bí mật riêng) & Cao & Thiết bị giả mạo không gửi
được bản tin vào topic hợp lệ \\\hline
Bảo vệ API & Giới hạn tần suất gọi, danh sách nguồn được phép gọi
(CORS), kiểm tra dữ liệu đầu vào bằng schema (Zod) & Cao & Tất cả điểm
gọi công khai đều có giới hạn tần suất và bước kiểm tra đầu vào \\\hline
Độ tin cậy dịch vụ & Điểm kiểm tra trạng thái sống, chính sách tự khởi
động lại và cơ chế phát hiện treo trong firmware (watchdog) & Cao & Tự
phục hồi khi dịch vụ treo, có trạng thái sống/chết rõ ràng \\\hline
Dữ liệu và sao lưu & Sao lưu PostgreSQL tự động + chính sách giữ dữ liệu
cho log và dữ liệu vận hành & Cao & Phục hồi được dữ liệu theo mục tiêu
khôi phục đã đặt ra \\\hline
Lưu trữ cục bộ SD & Thử gắn lại thẻ nhớ với thời gian chờ tăng dần,
buộc ghi ra thẻ trước khi đóng file và tháo thẻ đúng vòng đời phiên &
Cao & Không phát sinh lỗi ghi hàng loạt khi mạng gián đoạn \\\hline
Phát lại hàng đợi SD & Đặt thời gian chờ xác nhận, thử lại với khoảng
chờ tăng dần, tách bản tin quan trọng và dọn hàng đợi theo hạn mức &
Cao & Hàng đợi phát lại tiến dần, không bị nghẽn vô hạn \\\hline
Hiệu năng hệ thống & Tái sử dụng sẵn kết nối cơ sở dữ liệu, tối ưu truy
vấn và thêm vùng nhớ đệm khi cần & Trung bình & Độ trễ P95 của API giữ
ổn định khi tăng tải \\\hline
Khả năng mở rộng & Phân tải giữa nhiều bản sao dịch vụ API khi tăng số
thiết bị & Trung bình & Việc tăng thêm bản sao không làm gián đoạn kết nối
WebSocket/MQTT \\\hline
Quan sát và cảnh báo & Ghi nhận sự kiện an ninh, phát cảnh báo khi lỗi
hoặc độ trễ vượt ngưỡng & Trung bình & Cảnh báo được gửi tới Telegram/Email khi vượt
ngưỡng \\\hline
\end{longtable}
}


Lộ trình thực hiện được chia làm 2 đợt: \textbf{đợt 1.5} (ưu tiên cao,
hoàn thành trước khi đưa vào chạy thật) và \textbf{đợt 2} (mở rộng năng lực
chịu tải và giám sát theo tải thực tế). Cách chia này giúp giảm rủi ro
vận hành mà không làm chậm tiến độ đưa hệ thống vào khai thác.

\begin{center}\rule{0.5\linewidth}{0.5pt}\end{center}

\subsubsection{Tóm tắt}\label{tuxf3m-tux1eaft}

Hệ thống máy chủ được triển khai theo kiến trúc nhiều dịch vụ nhỏ, tức mỗi
dịch vụ được tách thành một tiến trình độc lập, đóng gói trong container
Docker riêng và giao tiếp
qua một mạng Docker dùng chung. Kết quả triển khai cho thấy
MQTT Bridge thực hiện đúng vai trò cầu nối dữ liệu giữa thiết bị IoT và
các hệ lưu trữ, đồng thời API phía máy chủ và giao diện web đáp ứng yêu
cầu tích hợp, giám sát vận hành thời gian thực.

Mô hình mỗi dịch vụ có một cấu hình Docker Compose riêng giúp hệ thống dễ bảo
trì, mở rộng và triển khai độc lập từng thành phần. Hệ thống giám sát
với VictoriaMetrics, VictoriaLogs, Grafana và Sentry bảo đảm khả năng
quan sát và phân tích sự cố trong môi trường vận hành thực tế.


{\thesissettableid{4.11A}{tab-ch4-tong-hop-cong-nghe-su-dung-trong-he-thong-cloud}
\def\LTcaptype{table}
\begin{longtable}{|L{0.214\linewidth}|L{0.299\linewidth}|L{0.115\linewidth}|L{0.327\linewidth}|}
\caption{Tổng hợp công nghệ sử dụng trong hệ thống máy chủ}\label{tab:ch4-tong-hop-cong-nghe-su-dung-trong-he-thong-cloud}\\
\hline \textbf{Thành phần} & \textbf{Công nghệ} & \textbf{Phiên bản} & \textbf{Vai trò} \\
\\\hline
\endfirsthead
\caption[]{Tổng hợp công nghệ sử dụng trong hệ thống máy chủ (tiếp theo)}\\
\hline \textbf{Thành phần} & \textbf{Công nghệ} & \textbf{Phiên bản} & \textbf{Vai trò} \\
\\\hline
\endhead
\\\hline
\endfoot
\endlastfoot
MQTT Broker & EMQX & 5.4.0 & Tiếp nhận dữ liệu từ thiết bị IoT \\\hline
MQTT Bridge & Node.js + TypeScript & 20 LTS & Xử lý và phân phối dữ liệu
MQTT \\\hline
API phía máy chủ & Express + TypeScript & 4.x & Cung cấp API, đẩy dữ
liệu thời gian thực và xử lý nghiệp vụ \\\hline
Giao diện web & Next.js + React & 15 / 19 & Giao diện web quản trị \\\hline
CSDL quan hệ & PostgreSQL & 16 & Dữ liệu người dùng, xe, cảnh báo \\\hline
CSDL chuỗi thời gian & VictoriaMetrics & - & Dữ liệu vận hành gửi liên tục (GNSS,
nguồn, rung, trạng thái) \\\hline
Nhật ký tập trung & VictoriaLogs & - & Nhật ký các dịch vụ \\\hline
Màn hình giám sát & Grafana & - & Trực quan hóa chỉ số và nhật ký \\\hline
Container & Docker + Docker Compose & - & Đóng gói và triển khai dịch
vụ \\\hline
Cổng chuyển tiếp & Nginx Proxy Manager & - & Quản lý tên miền, SSL,
định tuyến truy cập \\\hline
\end{longtable}
}


\subsection{4.3. Đo lường và kết quả -- Measurements and Results}\label{43-ux111o-lux1b0ux1eddng-vuxe0-kux1ebft-quux1ea3--measurements-and-results}

Phần này trình bày thiết lập thử nghiệm, kết quả đo theo từng tầng hệ
thống (phần cứng, firmware, máy chủ), kết quả kiểm thử tích hợp end-to-end
và đối chiếu với các tiêu chí thiết kế ở Chương 1 và Chương 2. Nói cách
khác, đây là phần kiểm chứng xem hệ thống sau khi hiện thực có thực sự
đáp ứng được những gì đã hứa ở các chương trước hay không. Toàn bộ giá
trị được ghi nhận trên thiết bị phần cứng đã chế tạo và lắp đặt thực tế
trong điều kiện vận hành bám sát bối cảnh khai thác.

\textbf{Ghi chú:} Các giá trị đo lường trong chương này là kết quả thu
được trên môi trường thử nghiệm cụ thể (xem mục 4.3.1). Một số chỉ số
được trình bày dưới dạng ước tính ban đầu hoặc cần đo lặp lại trên nhiều
mẫu để tăng độ tin cậy thống kê.

\begin{center}\rule{0.5\linewidth}{0.5pt}\end{center}

\subsubsection{4.3.1. Thiết lập môi trường thử
nghiệm}\label{431-thiux1ebft-lux1eadp-muxf4i-trux1b0ux1eddng-thux1eed-nghiux1ec7m}

\paragraph{4.3.1.1. Môi trường thử nghiệm phần
cứng}\label{4311-muxf4i-trux1b0ux1eddng-thux1eed-nghiux1ec7m-phux1ea7n-cux1ee9ng}

Kiểm thử phần cứng được triển khai tại phòng thí nghiệm và trên xe thực
tế với các thiết bị đo lường sau:


{\thesissettableid{4.12}{tab-ch4-thiet-bi-o-luong-su-dung-trong-thu-nghiem-phan-cung}
\def\LTcaptype{table}
\begin{longtable}{|L{0.085\linewidth}|L{0.175\linewidth}|L{0.220\linewidth}|L{0.210\linewidth}|L{0.255\linewidth}|}
\caption{Thiết bị đo lường sử dụng trong thử nghiệm phần cứng}\label{tab:ch4-thiet-bi-o-luong-su-dung-trong-thu-nghiem-phan-cung}\\
\hline \textbf{STT} & \textbf{Thiết bị} & \textbf{Model} & \textbf{Thông số chính} & \textbf{Mục đích sử dụng} \\
\\\hline
\endfirsthead
\caption[]{Thiết bị đo lường sử dụng trong thử nghiệm phần cứng (tiếp theo)}\\
\hline \textbf{STT} & \textbf{Thiết bị} & \textbf{Model} & \textbf{Thông số chính} & \textbf{Mục đích sử dụng} \\
\\\hline
\endhead
\\\hline
\endfoot
\endlastfoot
1 & Đồng hồ vạn năng (DMM) & UNI-T UT61E & Độ chính xác 0.5\%, đo dòng
µA-10A & Đo dòng tiêu thụ các chế độ \\\hline
2 & Nguồn cấp DC & UNI-T UTP1306S & 0--32V / 0--6A, độ phân giải
10mV/1mA & Mô phỏng ắc quy xe 12V hoặc 24V \\\hline
3 & Oscillo\-scope & Rigol DS1054Z & 50 MHz, 4 kênh & Phân tích tín hiệu,
đo thời gian chuyển trạng thái \\\hline
4 & Tủ nhiệt & Tủ nhiệt phòng thí nghiệm & -20°C đến +80°C & Kiểm thử
nhiệt độ hoạt động \\\hline
5 & OBD2 Simulator & ELM327 OBD2 Simulator Board & Hỗ trợ tất cả giao
thức OBD2 & Mô phỏng dữ liệu xe \\\hline
6 & Điểm đo GPS tham chiếu & Thử nghiệm ngoài trời & Tọa độ chuẩn đối
chiếu & Kiểm thử độ chính xác GPS \\\hline
\end{longtable}
}


\textbf{Xe thử nghiệm:}

\begin{itemize}
\tightlist
\item
  Xe 1: Toyota Vios 2020 (động cơ 1.5L, ắc quy 45Ah, OBD2 giao thức ISO
  15765--4 CAN)
\item
  Xe 2: Honda City 2021 (động cơ 1.5L, ắc quy 40Ah, OBD2 giao thức ISO
  15765--4 CAN)
\item
  Giai đoạn hiện tại đã hoàn tất đo trên 2 xe; xe thứ ba sẽ được bổ sung
  ở vòng kiểm thử tương thích mở rộng.
\end{itemize}


{\thesissetfigureid{4.29}{fig-ch4-bo-tri-thiet-bi-o-luong-trong-phong-thi-nghiem}
\begin{figure}[htbp]
\centering
\pandocbounded{\safeincludesvg{./assets/figures/10-chuong-4-ket-qua-do-luong-hinh-4-20.svg}}
\caption{Bố trí thiết bị đo lường trong phòng thí nghiệm}
\label{fig:ch4-bo-tri-thiet-bi-o-luong-trong-phong-thi-nghiem}
\end{figure}
}



{\thesissetfigureid{4.30}{fig-ch4-thiet-bi-tracker-uoc-lap-at-tren-xe-thu-nghiem}
\begin{figure}[htbp]
\centering
\pandocbounded{\safeincludesvg{./assets/figures/10-chuong-4-ket-qua-do-luong-hinh-4-21.svg}}
\caption{Thiết bị theo dõi được lắp đặt trên xe thử nghiệm}
\label{fig:ch4-thiet-bi-tracker-uoc-lap-at-tren-xe-thu-nghiem}
\end{figure}
}


\paragraph{4.3.1.2. Môi trường thử nghiệm phần
mềm}\label{4312-muxf4i-trux1b0ux1eddng-thux1eed-nghiux1ec7m-phux1ea7n-mux1ec1m}

Hệ thống máy chủ được triển khai trên hai môi trường riêng biệt để kiểm
thá»­:


{\thesissettableid{4.13}{tab-ch4-cau-hinh-moi-truong-thu-nghiem-phan-mem}
\def\LTcaptype{table}
\begin{longtable}{|L{0.287\linewidth}|L{0.344\linewidth}|L{0.334\linewidth}|}
\caption{Cấu hình môi trường thử nghiệm phần mềm}\label{tab:ch4-cau-hinh-moi-truong-thu-nghiem-phan-mem}\\
\hline \textbf{Thành phần} & \textbf{Môi trường phát triển} & \textbf{Môi trường kiểm thử tải} \\
\\\hline
\endfirsthead
\caption[]{Cấu hình môi trường thử nghiệm phần mềm (tiếp theo)}\\
\hline \textbf{Thành phần} & \textbf{Môi trường phát triển} & \textbf{Môi trường kiểm thử tải} \\
\\\hline
\endhead
\\\hline
\endfoot
\endlastfoot
Server & PC cá nhân, Intel i5--12400, 16 GB RAM, SSD 512 GB & VPS 4
vCPU, 8 GB RAM, SSD 100 GB \\\hline
Hệ điều hành & Windows 11 + WSL2 (Ubuntu 22.04) & Ubuntu 22.04 LTS \\\hline
Docker & Docker Desktop 4.x & Docker Engine 24.x \\\hline
Mạng & LAN 1 Gbps + Wi-Fi 5 GHz & VPS bandwidth 100 Mbps \\\hline
EMQX Broker & Single node, port 1883 & Single node, port 1883 \\\hline
PostgreSQL & Version 16, default config & Version 16, tuned
(shared\_buffers=2GB) \\\hline
VictoriaMetrics & Single node & Single node \\\hline
\end{longtable}
}

Từ cấu hình môi trường trên, bộ công cụ đo được chọn theo đúng từng lớp
cần kiểm chứng: API đo tải bằng k6, lớp MQTT đo bằng MQTT Bench, còn
giao diện web được đánh giá bằng Lighthouse.


\textbf{Công cụ kiểm thử hiệu năng:}

\begin{itemize}
\tightlist
\item
  \textbf{k6} (Grafana Labs): Kiểm thử tải các API với nhiều kịch
  bản đồng thời
\item
  \textbf{MQTT Bench}: Công cụ mô phỏng nhiều MQTT client đồng thời
\item
  \textbf{Lighthouse} (Google): Đánh giá hiệu năng giao diện web
\item
  \textbf{Playwright}: Kiểm thử giao diện tự động theo toàn bộ luồng dùng
  trình duyệt
\end{itemize}


{\thesissetfigureid{4.31}{fig-ch4-so-o-moi-truong-thu-nghiem-tong-the}
\begin{figure}[htbp]
\centering
\pandocbounded{\safeincludesvg{./assets/figures/10-chuong-4-ket-qua-do-luong-hinh-4-22.svg}}
\caption{Sơ đồ môi trường thử nghiệm tổng thể}
\label{fig:ch4-so-o-moi-truong-thu-nghiem-tong-the}
\end{figure}
}


\paragraph{4.3.1.3. Kịch bản thử
nghiệm}\label{4313-kux1ecbch-bux1ea3n-thux1eed-nghiux1ec7m}

Các kịch bản thử nghiệm được thiết kế bao phủ ba nhóm chính:

\begin{enumerate}
\def\labelenumi{\arabic{enumi}.}
\tightlist
\item
  \textbf{Kiểm thử đơn vị (Unit Test):} Kiểm tra từng module độc lập ---
  phần cứng, firmware, API phía máy chủ và các thành phần giao diện web.
\item
  \textbf{Kiểm thử tích hợp (Integration Test):} Kiểm tra sự tương tác
  giữa các module --- BLE + OBD2, MQTT + Bridge + Database, API +
  WebSocket + Frontend.
\item
  \textbf{Kiểm thử hệ thống (System Test):} Kiểm tra toàn bộ luồng dữ
  liệu end-to-end, từ thiết bị đến giao diện người dùng.
\end{enumerate}

\begin{center}\rule{0.5\linewidth}{0.5pt}\end{center}

\subsubsection{4.3.2. Kết quả kiểm thử phần
cứng}\label{432-kux1ebft-quux1ea3-kiux1ec3m-thux1eed-phux1ea7n-cux1ee9ng}

\paragraph{4.3.2.1. Đo dòng tiêu thụ năng
lượng}\label{4321-ux111o-duxf2ng-tiuxeau-thux1ee5-nux103ng-lux1b0ux1ee3ng}

Dòng tiêu thụ năng lượng là chỉ tiêu cốt lõi của thiết bị theo dõi, quyết
định thời lượng hoạt động và mức ảnh hưởng lên ắc quy xe. Phép đo được
thực hiện bằng cách mắc nối tiếp đồng hồ vạn năng giữa nguồn cấp DC và
thiết bị, sau đó ghi nhận giá trị dòng trung bình trong 60 giây cho mỗi
chế độ.


{\thesissettableid{4.14}{tab-ch4-ket-qua-o-dong-tieu-thu-nang-luong-theo-che-o-hoat}
\def\LTcaptype{table}
\begin{longtable}{|L{0.068\linewidth}|L{0.183\linewidth}|L{0.179\linewidth}|L{0.120\linewidth}|L{0.155\linewidth}|L{0.230\linewidth}|}
\caption{Kết quả đo dòng tiêu thụ năng lượng theo chế độ hoạt động}\label{tab:ch4-ket-qua-o-dong-tieu-thu-nang-luong-theo-che-o-hoat}\\
\hline \textbf{STT} & \textbf{Chế độ hoạt động} & \textbf{Mô tả} & \textbf{Dòng trung bình} & \textbf{Dòng đỉnh (peak)} & \textbf{Ghi chú} \\
\\\hline
\endfirsthead
\caption[]{Kết quả đo dòng tiêu thụ năng lượng theo chế độ hoạt động (tiếp theo)}\\
\hline \textbf{STT} & \textbf{Chế độ hoạt động} & \textbf{Mô tả} & \textbf{Dòng trung bình} & \textbf{Dòng đỉnh (peak)} & \textbf{Ghi chú} \\
\\\hline
\endhead
\\\hline
\endfoot
\endlastfoot
1 & Xe đang chạy - gửi dữ liệu liên tục & GPS + 4G + BLE + vi điều khiển hoạt động đầy đủ &
~350 mA & ~520 mA & Dòng đỉnh khi truyền dữ liệu
4G \\\hline
2 & Xe đang chạy - chờ giữa hai lần gửi & 4G còn kết nối + vi điều khiển vẫn hoạt động &
~180 mA & ~280 mA & Giữa các chu kỳ gửi dữ
liệu \\\hline
3 & Xe đỗ (IGN OFF) & Vi điều khiển ngủ nhẹ + 4G tắt & ~15
mA & ~25 mA & Heartbeat mỗi 10--30 phút \\\hline
4 & Ngủ sâu & Vi điều khiển ngủ sâu + IMU đánh thức & ~0.5 mA
& ~2 mA (ước tính, cần xác nhận thêm) & Chỉ IMU + RTC hoạt
động \\\hline
5 & Cảnh báo bất thường & Đánh thức từ ngủ sâu &
~380 mA & ~550 mA & Tương tự chế độ xe chạy + GPS
cold start \\\hline
\end{longtable}
}



\clearpage

{\thesissetfigureid{4.32}{fig-ch4-o-thi-dong-tieu-thu-theo-thoi-gian-trong-mot-chu-ky}
\begin{figure}[htbp]
\centering
\pandocbounded{\safeincludesvg{./assets/figures/10-chuong-4-ket-qua-do-luong-hinh-4-23.svg}}
\caption{Đồ thị dòng tiêu thụ theo thời gian trong một chu kỳ hoạt động hoàn chỉnh (xe chạy -\textgreater{} xe đỗ -\textgreater{} cảnh báo -\textgreater{} xe đỗ)}
\label{fig:ch4-o-thi-dong-tieu-thu-theo-thoi-gian-trong-mot-chu-ky}
\end{figure}
}


\textbf{Phân tích kết quả:}

\begin{itemize}
\tightlist
\item
  Dòng tiêu thụ trong chế độ Active (~350 mA) phù hợp với
  tính toán thiết kế tại Chương 3, trong đó modem SIM7600CE-T chiếm
  khoảng 200--250 mA (khi truyền dữ liệu 4G), ESP32-S3 chiếm khoảng
  60--80 mA và linh kiện phụ trợ chiếm khoảng 20--30 mA.
\item
  Dòng tiêu thụ trong chế độ Deep Sleep (~0.5 mA) nằm sát
  ngưỡng thiết kế hệ thống. Giá trị này cao hơn mức lý thuyết của riêng
  lõi ESP32-S3 + LIS3DSH do còn có tổn hao từ mạch nguồn, đường cấp
  nguồn, linh kiện phụ trợ và điều kiện đo thực tế. Dù vậy, mức tiêu thụ
  này vẫn đủ để thiết bị duy trì giám sát đỗ xe trong thời gian dài mà
  không gây hao ắc quy đáng kể.
\item
  Dòng đỉnh khi truyền dữ liệu 4G (~520 mA) cần được lưu ý
  trong thiết kế mạch nguồn, đảm bảo tụ điện lọc đặt gần nguồn
  đủ lớn để tránh sụt áp.
\end{itemize}

\paragraph{4.3.2.2. Thời lượng pin dự
phòng}\label{4322-thux1eddi-lux1b0ux1ee3ng-pin-dux1ef1-phuxf2ng}

Pin dự phòng 18650 1S (dung lượng danh định 3500 mAh, điện áp danh định
3.7V) được quy đổi theo cùng công thức đã dùng cho phương án gốc:
\(T = \dfrac{C \times V_{\text{pin}} \times \eta}{P_{\text{tieu thu}}}\),
trong đó \(\eta\) là hiệu suất chuyển đổi của
mạch boost converter (~85--90\%). Các giá trị dưới đây là
\textbf{quy đổi theo bộ pin tham chiếu 18650}, chưa phải bộ số đo thực
nghiệm mới trên pin 18650 thực tế.


\clearpage

{\thesissettableid{4.15}{tab-ch4-thoi-luong-pin-du-phong-theo-kich-ban-su-dung-uoc-tinh}
\def\LTcaptype{table}
\begin{longtable}{|L{0.068\linewidth}|L{0.230\linewidth}|L{0.117\linewidth}|L{0.142\linewidth}|L{0.209\linewidth}|L{0.169\linewidth}|}
\caption{Thời lượng pin dự phòng theo kịch bản sử dụng (quy đổi theo bộ pin tham chiếu 18650)}\label{tab:ch4-thoi-luong-pin-du-phong-theo-kich-ban-su-dung-uoc-tinh}\\
\hline \textbf{STT} & \textbf{Kịch bản sử dụng} & \textbf{Dòng trung bình} & \textbf{Thời lượng ước tính} & \textbf{Trạng
thái dữ liệu} & \textbf{Ghi chú} \\
\\\hline
\endfirsthead
\caption[]{Thời lượng pin dự phòng theo kịch bản sử dụng (quy đổi theo bộ pin tham chiếu 18650) (tiếp theo)}\\
\hline \textbf{STT} & \textbf{Kịch bản sử dụng} & \textbf{Dòng trung bình} & \textbf{Thời lượng ước tính} & \textbf{Trạng
thái dữ liệu} & \textbf{Ghi chú} \\
\\\hline
\endhead
\\\hline
\endfoot
\endlastfoot
1 & Theo dõi liên tục khi xe chạy & ~350 mA &
~2.6--3.0 giờ & Quy đổi theo bộ pin tham chiếu 18650 & Trường hợp
xấu nhất, phụ thuộc hiệu suất boost \\\hline
2 & Sleep + heartbeat 30 phút & ~15 mA & ~64
giờ & Quy đổi theo bộ pin tham chiếu 18650 & Xe đậu bình thường \\\hline
3 & Ngủ sâu, chỉ chờ IMU đánh thức & ~0.5 mA &
~1,940 giờ (~81 ngày) & Quy đổi theo
bộ pin tham chiếu 18650 & Chế độ tiết kiệm tối đa \\\hline
4 & Hỗn hợp (50\% xe chạy + 50\% xe đỗ) & ~182 mA &
~5.0 giờ & Quy đổi theo bộ pin tham chiếu 18650 & Mô phỏng sử dụng
thực tế \\\hline
\end{longtable}
}



\clearpage

{\thesissetfigureid{4.33}{fig-ch4-o-thi-ien-ap-pin-du-phong-theo-thoi-gian-theo-pin-tham-chieu-18650}
\begin{figure}[htbp]
\centering
\pandocbounded{\safeincludesvg{./assets/figures/10-chuong-4-ket-qua-do-luong-hinh-4-24.svg}}
\caption{Đồ thị điện áp pin dự phòng theo thời gian khi quy đổi theo bộ pin tham chiếu 18650 trong kịch bản theo dõi liên tục khi xe chạy}
\label{fig:ch4-o-thi-ien-ap-pin-du-phong-theo-thoi-gian-theo-pin-tham-chieu-18650}
\end{figure}
}


\textbf{Nhận xét:} Với bộ pin tham chiếu 18650, thời lượng pin dự phòng trong
chế độ theo dõi liên tục khi xe chạy được ước tính còn khoảng 2.6--3.0 giờ tùy hiệu
suất chuyển đổi thực tế. Trong kịch bản xe đỗ và chỉ gửi heartbeat định
kỳ, thiết bị vẫn có thể duy trì hoạt động trong nhiều ngày; tuy nhiên,
các giá trị này cần được kiểm thử lại trên mẫu pin 18650 thực tế nếu
dùng cho giai đoạn chế tạo tiếp theo.

\paragraph{4.3.2.3. Kiểm thử nhiệt độ hoạt
động}\label{4323-kiux1ec3m-thux1eed-nhiux1ec7t-ux111ux1ed9-houx1ea1t-ux111ux1ed9ng}

Thiết bị được đặt trong tủ nhiệt để kiểm tra khả năng hoạt động ở các
mức nhiệt độ khác nhau, mô phỏng môi trường trong xe ô tô (có thể vượt
60°C vào mùa hè và xuống dưới 0°C trong phòng lạnh xe).


{\thesissettableid{4.16}{tab-ch4-ket-qua-kiem-thu-nhiet-o-hoat-ong}
\def\LTcaptype{table}
\begin{longtable}{|L{0.057\linewidth}|L{0.162\linewidth}|L{0.168\linewidth}|L{0.066\linewidth}|L{0.113\linewidth}|L{0.081\linewidth}|L{0.278\linewidth}|}
\caption{Kết quả kiểm thử nhiệt độ hoạt động}\label{tab:ch4-ket-qua-kiem-thu-nhiet-o-hoat-ong}\\
\hline \textbf{STT} & \textbf{Nhiệt độ thử nghiệm} & \textbf{Trạng thái hoạt động} & \textbf{GPS fix} & \textbf{4G kết nối} & \textbf{OBD2 BLE} & \textbf{Ghi chú} \\
\\\hline
\endfirsthead
\caption[]{Kết quả kiểm thử nhiệt độ hoạt động (tiếp theo)}\\
\hline \textbf{STT} & \textbf{Nhiệt độ thử nghiệm} & \textbf{Trạng thái hoạt động} & \textbf{GPS fix} & \textbf{4G kết nối} & \textbf{OBD2 BLE} & \textbf{Ghi chú} \\
\\\hline
\endhead
\\\hline
\endfoot
\endlastfoot
1 & -10°C & Bình thường & Có & Có & Có & Thời gian khởi động tăng
20\% \\\hline
2 & 0°C & Bình thường & Có & Có & Có & Hoạt động ổn định \\\hline
3 & 25°C (nhiệt độ phòng) & Bình thường & Có & Có & Có & Điều kiện tham
chiếu \\\hline
4 & 45°C & Bình thường & Có & Có & Có & Hoạt động ổn định \\\hline
5 & 60°C & Bình thường & Có & Có & Có & Nhiệt độ MCU tăng, vẫn trong
giới hạn \\\hline
6 & 70°C & Cảnh báo & Có & Có (không ổn định) & Có & Module 4G bắt đầu
không ổn định, cần kiểm tra lặp lại \\\hline
7 & 80°C & Ngưng hoạt động & --- & --- & --- & Vượt giới hạn nhiệt độ an
toàn \\\hline
\end{longtable}
}



{\thesissetfigureid{4.34}{fig-ch4-o-thi-dong-tieu-thu-theo-nhiet-o-moi-truong}
\begin{figure}[htbp]
\centering
\pandocbounded{\safeincludesvg{./assets/figures/10-chuong-4-ket-qua-do-luong-hinh-4-25.svg}}
\caption{Đồ thị dòng tiêu thụ theo nhiệt độ môi trường}
\label{fig:ch4-o-thi-dong-tieu-thu-theo-nhiet-o-moi-truong}
\end{figure}
}


\textbf{Nhận xét:} Thiết bị hoạt động ổn định trong dải nhiệt độ -10°C
đến +60°C, đạt yêu cầu thiết kế. Tại nhiệt độ 70°C, modem 4G SIM7600CE-T
bắt đầu biểu hiện không ổn định (mất kết nối ngắt quãng), trong khi dải
nhiệt độ hoạt động danh định mà SIMCom công bố cho dòng SIM7600CE là
-40°C đến +85°C {[}1{]}.

\paragraph{4.3.2.4. Kiểm thử tương thích
OBD2}\label{4324-kiux1ec3m-thux1eed-tux1b0ux1a1ng-thuxedch-obd2}

Thiết bị được kiểm thử kết nối với nhiều thương hiệu xe khác nhau thông
qua adapter vgate iCar Pro BLE để đánh giá khả năng tương thích.


{\thesissettableid{4.17}{tab-ch4-ket-qua-kiem-thu-tuong-thich-obd2-tren-cac-thuong-hieu-xe}
\def\LTcaptype{table}
\begin{longtable}{|L{0.055\linewidth}|L{0.119\linewidth}|L{0.100\linewidth}|L{0.117\linewidth}|L{0.108\linewidth}|L{0.091\linewidth}|L{0.103\linewidth}|L{0.120\linewidth}|L{0.092\linewidth}|}
\caption{Kết quả kiểm thử tương thích OBD2 trên các thương hiệu xe}\label{tab:ch4-ket-qua-kiem-thu-tuong-thich-obd2-tren-cac-thuong-hieu-xe}\\
\hline \textbf{STT} & \textbf{Xe thử nghiệm} & \textbf{Năm sản xuất} & \textbf{Giao thức OBD2} & \textbf{Kết nối BLE} & \textbf{Đọc
RPM} & \textbf{Đọc Speed} & \textbf{Đọc Coolant Temp} & \textbf{Ghi chú} \\
\\\hline
\endfirsthead
\caption[]{Kết quả kiểm thử tương thích OBD2 trên các thương hiệu xe (tiếp theo)}\\
\hline \textbf{STT} & \textbf{Xe thử nghiệm} & \textbf{Năm sản xuất} & \textbf{Giao thức OBD2} & \textbf{Kết nối BLE} & \textbf{Đọc
RPM} & \textbf{Đọc Speed} & \textbf{Đọc Coolant Temp} & \textbf{Ghi chú} \\
\\\hline
\endhead
\\\hline
\endfoot
\endlastfoot
1 & Toyota Vios & 2020 & ISO 15765--4 CAN & Thành công & Có & Có & Có &
Đầy đủ \\\hline
2 & Honda City & 2021 & ISO 15765--4 CAN & Thành công & Có & Có & Có &
Đầy đủ \\\hline
\end{longtable}
}

Bảng trên chỉ tổng hợp các mẫu xe đã được kiểm thử thực tế trong phạm vi
hiện tại; các mẫu xe chưa đo không được liệt kê để tránh làm sai lệch
kết luận đánh giá.


\textbf{Nhận xét:} Các xe thử nghiệm đều sử dụng giao thức CAN-bus (ISO
15765--4), là giao thức phổ biến nhất trên các xe sản xuất từ năm 2008
trở đi theo tiêu chuẩn OBD-II {[}2{]}. Adapter vgate iCar Pro hỗ trợ tất
cả các giao thức OBD2 (CAN, KWP2000, ISO 9141), đảm bảo tương thích với
đa số các dòng xe tại Việt Nam.

\paragraph{4.3.2.5. Kiểm thử mạch Low Voltage Disconnect
(LVD)}\label{4325-kiux1ec3m-thux1eed-mux1ea1ch-low-voltage-disconnect-lvd}

Mạch LVD được kiểm thử bằng cách sử dụng nguồn cấp DC (UNI-T UTP1306S)
mô phỏng ắc quy xe. Kết quả đo hiện tại được thực hiện theo profile 12V
(quét điện áp từ 13.5V xuống 10.5V) và ghi nhận điểm ngắt/đóng lại.


{\thesissettableid{4.18}{tab-ch4-ket-qua-kiem-thu-mach-lvd}
\def\LTcaptype{table}
\begin{longtable}{|L{0.068\linewidth}|L{0.210\linewidth}|L{0.204\linewidth}|L{0.258\linewidth}|L{0.082\linewidth}|L{0.112\linewidth}|}
\caption{Kết quả kiểm thử mạch LVD}\label{tab:ch4-ket-qua-kiem-thu-mach-lvd}\\
\hline \textbf{STT} & \textbf{Thông số} & \textbf{Giá trị thiết kế} & \textbf{Giá trị đo} & \textbf{Sai lệch} & \textbf{Trạng
thái} \\
\\\hline
\endfirsthead
\caption[]{Kết quả kiểm thử mạch LVD (tiếp theo)}\\
\hline \textbf{STT} & \textbf{Thông số} & \textbf{Giá trị thiết kế} & \textbf{Giá trị đo} & \textbf{Sai lệch} & \textbf{Trạng
thái} \\
\\\hline
\endhead
\\\hline
\endfoot
\endlastfoot
1 & Điện áp ngắt (disconnect) & Profile 12V: \(12.0\,\mathrm{V}\); Profile 24V: \(24.0\,\mathrm{V}\) &
11.48V {[}đo theo profile 12V{]} & 4.33\% & Chưa đạt (12V) \\\hline
2 & Điện áp đóng lại khi kết nối lại & Profile 12V: \(12.2\,\mathrm{V}\); Profile 24V:
\(24.4\,\mathrm{V}\) & 12.53V {[}đo theo profile 12V, cần tinh chỉnh về 12.2V{]} &
2.70\% & Chưa đạt (12V) \\\hline
3 & Độ trễ ngắt & \textless{} 100 ms & ~50 ms (đo bước
đầu) & --- & Đạt \\\hline
4 & Độ trễ đóng lại & \textless{} 500 ms & ~200 ms (đo
bước đầu) & --- & Đạt \\\hline
5 & Hysteresis & 1.0V & 1.05V (đo bước đầu) & 5\% & Đạt \\\hline
\end{longtable}
}


\textbf{Nhận xét:} Kết quả hiện tại đang theo profile 12V: điện áp ngắt
11.48V thấp hơn mục tiêu \(OFF = 12.0\,\mathrm{V}\), và điện áp đóng lại
12.53V cao hơn mục tiêu \(ON = 12.2\,\mathrm{V}\) nên cần tinh chỉnh
comparator/firmware. Với profile 24V, tiêu chí tương ứng cần đạt là ngắt tại
\(24.0\,\mathrm{V}\) và đóng lại tại \(24.4\,\mathrm{V}\)
khi thực hiện vòng đo 27V xuống 21V.

\begin{center}\rule{0.5\linewidth}{0.5pt}\end{center}

\subsubsection{4.3.3. Kết quả kiểm thử
firmware}\label{433-kux1ebft-quux1ea3-kiux1ec3m-thux1eed-firmware}

\paragraph{4.3.3.1. Thời gian kết nối BLE
OBD2}\label{4331-thux1eddi-gian-kux1ebft-nux1ed1i-ble-obd2}

Thời gian kết nối BLE được đo từ lúc ESP32-S3 bắt đầu quét (scan) thiết
bị BLE đến khi nhận được phản hồi OBD2 đầu tiên (gói trả lời cho lệnh ATZ).
Phép đo được lặp lại 20 lần và tính giá trị trung bình.


{\thesissettableid{4.19}{tab-ch4-ket-qua-o-thoi-gian-ket-noi-ble-obd2}
\def\LTcaptype{table}
\begin{longtable}{|L{0.068\linewidth}|L{0.294\linewidth}|L{0.179\linewidth}|L{0.123\linewidth}|L{0.123\linewidth}|L{0.148\linewidth}|}
\caption{Kết quả đo thời gian kết nối BLE OBD2}\label{tab:ch4-ket-qua-o-thoi-gian-ket-noi-ble-obd2}\\
\hline \textbf{STT} & \textbf{Thông số} & \textbf{Giá trị trung bình} & \textbf{Giá trị min} & \textbf{Giá trị max} & \textbf{Độ
lệch chuẩn} \\
\\\hline
\endfirsthead
\caption[]{Kết quả đo thời gian kết nối BLE OBD2 (tiếp theo)}\\
\hline \textbf{STT} & \textbf{Thông số} & \textbf{Giá trị trung bình} & \textbf{Giá trị min} & \textbf{Giá trị max} & \textbf{Độ
lệch chuẩn} \\
\\\hline
\endhead
\\\hline
\endfoot
\endlastfoot
1 & Thời gian scan BLE & ~1.5 giây & 0.8 giây & 3.2 giây &
0.6 giây \\\hline
2 & Thời gian kết nối BLE & ~1.0 giây & 0.5 giây & 2.1
giây & 0.4 giây \\\hline
3 & Thời gian khởi tạo ELM327 & ~1.5 giây & 1.0 giây & 2.5
giây & 0.3 giây \\\hline
4 & \textbf{Tổng thời gian (scan + connect + init)} &
\textbf{~4.0 giây} & \textbf{2.3 giây} & \textbf{7.8 giây}
& \textbf{1.3 giây} \\\hline
\end{longtable}
}



{\thesissetfigureid{4.35}{fig-ch4-bieu-o-phan-bo-thoi-gian-ket-noi-ble-obd2-20-lan}
\begin{figure}[htbp]
\centering
\pandocbounded{\safeincludesvg{./assets/figures/10-chuong-4-ket-qua-do-luong-hinh-4-26.svg}}
\caption{Biểu đồ phân bố thời gian kết nối BLE OBD2 (20 lần đo)}
\label{fig:ch4-bieu-o-phan-bo-thoi-gian-ket-noi-ble-obd2-20-lan}
\end{figure}
}


\textbf{Nhận xét:} Thời gian kết nối BLE OBD2 trung bình
~4 giây (scan ~1.5s + connect
~1.0s + init ~1.5s), nằm trong phạm vi mục
tiêu thiết kế (\textless{} 10 giây). Giá trị max 7.8 giây xuất hiện khi
có nhiều thiết bị BLE trong phạm vi scan, gây nhiễu. Trong điều kiện xe
bình thường (ít thiết bị BLE xung quanh), thời gian thường đạt 3--5
giây.

\paragraph{4.3.3.2. Thời gian phản hồi OBD2
PID}\label{4332-thux1eddi-gian-phux1ea3n-hux1ed3i-obd2-pid}

Thời gian phản hồi cho mỗi lệnh OBD2 PID được đo từ lúc gửi yêu cầu đến
khi nhận được gói trả lời hoàn chỉnh qua BLE. Mỗi PID được đo 50 lần và
tính giá trị trung bình.


{\thesissettableid{4.20}{tab-ch4-thoi-gian-phan-hoi-obd2-pid}
\def\LTcaptype{table}
\begin{longtable}{|L{0.076\linewidth}|L{0.138\linewidth}|L{0.263\linewidth}|L{0.256\linewidth}|L{0.097\linewidth}|L{0.106\linewidth}|}
\caption{Thời gian phản hồi OBD2 PID}\label{tab:ch4-thoi-gian-phan-hoi-obd2-pid}\\
\hline \textbf{STT} & \textbf{OBD2 PID} & \textbf{Mô tả} & \textbf{Thời gian trung bình} & \textbf{Min} & \textbf{Max} \\
\\\hline
\endfirsthead
\caption[]{Thời gian phản hồi OBD2 PID (tiếp theo)}\\
\hline \textbf{STT} & \textbf{OBD2 PID} & \textbf{Mô tả} & \textbf{Thời gian trung bình} & \textbf{Min} & \textbf{Max} \\
\\\hline
\endhead
\\\hline
\endfoot
\endlastfoot
1 & 0x0C & RPM động cơ & ~65 ms & 40 ms & 120 ms \\\hline
2 & 0x0D & Tốc độ xe (km/h) & ~60 ms & 38 ms & 110 ms \\\hline
3 & 0x05 & Nhiệt độ nước làm mát & ~70 ms & 45 ms & 130
ms \\\hline
4 & 0x2F & Mức nhiên liệu (\%) & ~75 ms & 50 ms & 140
ms \\\hline
5 & 0x04 & Tải động cơ (\%) & ~68 ms & 42 ms & 125 ms \\\hline
6 & 0x11 & Vị trí bướm ga (\%) & ~62 ms & 39 ms & 115
ms \\\hline
\end{longtable}
}


\textbf{Nhận xét:} Thời gian phản hồi trung bình cho mỗi PID đạt
~60--75 ms, nằm trong mục tiêu thiết kế
(~50--100 ms). Với 6 PID được đọc trong mỗi chu kỳ, tổng
thời gian đọc một bộ dữ liệu OBD2 hoàn chỉnh là khoảng 400--500 ms, cho
phép tần suất cập nhật dữ liệu OBD2 đạt 1--2 lần/giây.

\paragraph{4.3.3.3. Thời gian bắt vệ tinh GPS (GPS Fix
Time)}\label{4333-thux1eddi-gian-bux1eaft-vux1ec7-tinh-gps-gps-fix-time}

Thời gian bắt vệ tinh (Time to First Fix - TTFF) được đo trên GNSS tích
hợp của SIM7600CE-T trong ba kịch bản khác nhau.


{\thesissettableid{4.21}{tab-ch4-thoi-gian-bat-ve-tinh-gps-ttff}
\def\LTcaptype{table}
\begin{longtable}{|L{0.057\linewidth}|L{0.107\linewidth}|L{0.302\linewidth}|L{0.134\linewidth}|L{0.081\linewidth}|L{0.078\linewidth}|L{0.166\linewidth}|}
\caption{Thời gian bắt vệ tinh GPS (TTFF)}\label{tab:ch4-thoi-gian-bat-ve-tinh-gps-ttff}\\
\hline \textbf{STT} & \textbf{Kịch bản} & \textbf{Mô tả} & \textbf{TTFF trung bình} & \textbf{Min} & \textbf{Max} & \textbf{Số vệ tinh trung
bình} \\
\\\hline
\endfirsthead
\caption[]{Thời gian bắt vệ tinh GPS (TTFF) (tiếp theo)}\\
\hline \textbf{STT} & \textbf{Kịch bản} & \textbf{Mô tả} & \textbf{TTFF trung bình} & \textbf{Min} & \textbf{Max} & \textbf{Số vệ tinh trung
bình} \\
\\\hline
\endhead
\\\hline
\endfoot
\endlastfoot
1 & Cold Start & Không có dữ liệu vệ tinh cũ & ~30 giây &
20 giây & 60 giây & 6--8 \\\hline
2 & Warm Start & Có dữ liệu ephemeris còn hiệu lực & ~5
giây & 3 giây & 12 giây & 8--10 \\\hline
3 & Hot Start & Module GNSS đã hoạt động, tạm mất tín hiệu &
~1 giây & 0.5 giây & 3 giây & 10--12 \\\hline
\end{longtable}
}



{\thesissetfigureid{4.36}{fig-ch4-o-thi-thoi-gian-bat-ve-tinh-gps-trong-cac-kich-ban}
\begin{figure}[htbp]
\centering
\pandocbounded{\safeincludesvg{./assets/figures/10-chuong-4-ket-qua-do-luong-hinh-4-27.svg}}
\caption{Đồ thị thời gian bắt vệ tinh GPS trong các kịch bản (Cold/Warm/Hot Start)}
\label{fig:ch4-o-thi-thoi-gian-bat-ve-tinh-gps-trong-cac-kich-ban}
\end{figure}
}


\textbf{Độ chính xác vị trí GPS:}

Phép đo độ chính xác được thực hiện bằng cách đặt thiết bị tại một vị
trí cố biết (xác định bằng GPS chuyên nghiệp), ghi nhận 100 điểm tọa độ
trong 10 phút, và tính sai số trung bình.


{\thesissettableid{4.22}{tab-ch4-o-chinh-xac-vi-tri-gps}
\def\LTcaptype{table}
\begin{longtable}{|L{0.085\linewidth}|L{0.223\linewidth}|L{0.204\linewidth}|L{0.217\linewidth}|L{0.216\linewidth}|}
\caption{Độ chính xác vị trí GPS}\label{tab:ch4-o-chinh-xac-vi-tri-gps}\\
\hline \textbf{STT} & \textbf{Điều kiện} & \textbf{Sai số trung bình} & \textbf{Sai số max} & \textbf{CEP 95} \% \\
\\\hline
\endfirsthead
\caption[]{Độ chính xác vị trí GPS (tiếp theo)}\\
\hline \textbf{STT} & \textbf{Điều kiện} & \textbf{Sai số trung bình} & \textbf{Sai số max} & \textbf{CEP 95} \% \\
\\\hline
\endhead
\\\hline
\endfoot
\endlastfoot
1 & Ngoài trời, trời quang & ~2.3 mét & ~4.5
mét & ~3.8 mét \\\hline
2 & Ngoài trời, trời u ám & ~3.5 mét & ~7.2
mét & ~5.5 mét \\\hline
3 & Trong phố (urban canyon) & ~5.8 mét &
~15 mét (cần xác nhận thêm) & ~10 mét (cần
xác nhận thêm) \\\hline
4 & Bãi đỗ xe có mái che & ~8.5 mét (ước tính ban đầu) &
~20 mét (cần xác nhận thêm) & ~15 mét (cần
xác nhận thêm) \\\hline
\end{longtable}
}


\textbf{Nhận xét:} Độ chính xác vị trí GPS đạt ~2--3 mét
trong điều kiện ngoài trời thoáng, vượt mục tiêu thiết kế (\textless{} 5
mét). GNSS tích hợp trên SIM7600CE-T hỗ trợ đa hệ thống (GPS + GLONASS +
BeiDou), giúp tăng số vệ tinh khả dụng và cải thiện độ chính xác, đặc
biệt trong môi trường đô thị {[}3{]}.

\paragraph{4.3.3.4. Độ trễ truyền dữ liệu
MQTT}\label{4334-ux111ux1ed9-trux1ec5-truyux1ec1n-dux1eef-liux1ec7u-mqtt}

Độ trễ truyền dữ liệu MQTT được đo từ lúc firmware hoàn thành đóng gói
bản tin MQTT (publish) đến khi MQTT Bridge nhận được bản tin
(subscribe). Phép đo sử dụng timestamp đồng bộ giữa thiết bị và máy chủ.


{\thesissettableid{4.23}{tab-ch4-o-tre-truyen-du-lieu-mqtt-qua-mang-4g}
\def\LTcaptype{table}
\begin{longtable}{|L{0.067\linewidth}|L{0.248\linewidth}|L{0.085\linewidth}|L{0.206\linewidth}|L{0.092\linewidth}|L{0.115\linewidth}|L{0.111\linewidth}|}
\caption{Độ trễ truyền dữ liệu MQTT qua mạng 4G}\label{tab:ch4-o-tre-truyen-du-lieu-mqtt-qua-mang-4g}\\
\hline \textbf{STT} & \textbf{Điều kiện mạng} & \textbf{QoS} & \textbf{Độ trễ trung bình} & \textbf{Min} & \textbf{Max} & \textbf{P95} \\
\\\hline
\endfirsthead
\caption[]{Độ trễ truyền dữ liệu MQTT qua mạng 4G (tiếp theo)}\\
\hline \textbf{STT} & \textbf{Điều kiện mạng} & \textbf{QoS} & \textbf{Độ trễ trung bình} & \textbf{Min} & \textbf{Max} & \textbf{P95} \\
\\\hline
\endhead
\\\hline
\endfoot
\endlastfoot
1 & 4G ổn định (\textgreater= 3 thanh) & QoS 0 & ~120 ms &
60 ms & 350 ms & ~250 ms \\\hline
2 & 4G ổn định (\textgreater= 3 thanh) & QoS 1 & ~180 ms &
80 ms & 500 ms & ~380 ms \\\hline
3 & 4G yếu (1--2 thanh) & QoS 0 & ~350 ms & 120 ms & 1200
ms & ~800 ms \\\hline
4 & 4G yếu (1--2 thanh) & QoS 1 & ~500 ms & 150 ms & 2000
ms & ~1500 ms \\\hline
\end{longtable}
}



{\thesissetfigureid{4.37}{fig-ch4-bieu-o-phan-bo-o-tre-mqtt-trong-ieu-kien-mang-4g}
\begin{figure}[htbp]
\centering
\pandocbounded{\safeincludesvg{./assets/figures/10-chuong-4-ket-qua-do-luong-hinh-4-28.svg}}
\caption{Biểu đồ phân bố độ trễ MQTT trong điều kiện mạng 4G ổn định}
\label{fig:ch4-bieu-o-phan-bo-o-tre-mqtt-trong-ieu-kien-mang-4g}
\end{figure}
}


\textbf{Nhận xét:} Độ trễ MQTT trung bình ~120--180 ms
trong điều kiện mạng 4G ổn định, đạt mục tiêu thiết kế
(~100--300 ms). QoS 1 có độ trễ cao hơn QoS 0 khoảng 50\%
do có thêm bước ACK, nhưng đổi lại bảo đảm tin nhắn được gửi ít nhất một
lần cho các bản tin quan trọng. Trong điều kiện mạng yếu, độ trễ tăng
đáng kể nhưng hệ thống vẫn duy trì được luồng dữ liệu vận hành nhờ cơ chế
thử gửi lại, tự kết nối lại và cấu hình QoS phù hợp.

\paragraph{4.3.3.5. Độ tin cậy của máy trạng thái (State
Machine)}\label{4335-ux111ux1ed9-tin-cux1eady-cux1ee7a-muxe1y-trux1ea1ng-thuxe1i-state-machine}

Máy trạng thái quản lý các chế độ hoạt động (xe chạy -\textgreater{}
xe đỗ -\textgreater{} cảnh báo -\textgreater{} xe chạy) được kiểm thử
bằng cách chạy liên tục 1000 chu kỳ chuyển đổi trạng thái và ghi nhận
kết quả.


{\thesissettableid{4.24}{tab-ch4-ket-qua-kiem-thu-may-trang-thai-firmware}
\def\LTcaptype{table}
\begin{longtable}{|L{0.068\linewidth}|L{0.296\linewidth}|L{0.181\linewidth}|L{0.126\linewidth}|L{0.099\linewidth}|L{0.165\linewidth}|}
\caption{Kết quả kiểm thử máy trạng thái firmware}\label{tab:ch4-ket-qua-kiem-thu-may-trang-thai-firmware}\\
\hline \textbf{STT} & \textbf{Chuyển đổi trạng thái} & \textbf{Số lần thử nghiệm} & \textbf{Thành công} & \textbf{Thất bại} & \textbf{Tỷ lệ thành công} \\
\\\hline
\endfirsthead
\caption[]{Kết quả kiểm thử máy trạng thái firmware (tiếp theo)}\\
\hline \textbf{STT} & \textbf{Chuyển đổi trạng thái} & \textbf{Số lần thử nghiệm} & \textbf{Thành công} & \textbf{Thất bại} & \textbf{Tỷ lệ thành công} \\
\\\hline
\endhead
\\\hline
\endfoot
\endlastfoot
1 & Xe chạy -\textgreater{} Xe đỗ (IGN OFF) & 1000 & 1000 & 0 &
100\% \\\hline
2 & Xe đỗ -\textgreater{} Xe chạy (IGN ON) & 1000 & 1000 & 0 &
100\% \\\hline
3 & Xe đỗ -\textgreater{} Cảnh báo (IMU kích hoạt) & 1000 & 998 & 2 &
99.8\% \\\hline
4 & Cảnh báo -\textgreater{} Xe đỗ (hết thời gian chờ/xác nhận) & 1000 & 1000 & 0 &
100\% \\\hline
5 & Cảnh báo -\textgreater{} Xe chạy (IGN ON) & 500 & 500 & 0 & 100\% \\\hline
6 & Mất nguồn -\textgreater{} Pin dự phòng & 200 & 200 & 0 & 100\% \\\hline
\end{longtable}
}


\textbf{Nhận xét:} Máy trạng thái hoạt động ổn định với tỷ lệ thành công
\textgreater= 99.8\%. Hai trường hợp thất bại trong chuyển đổi Parking
-\textgreater{} Alert là do ngưỡng IMU được cấu hình quá nhạy trong điều
kiện rung động cơ học, dẫn đến miss trigger. Vấn đề này đã được khắc
phục bằng việc điều chỉnh ngưỡng và bộ lọc (digital filter) cho IMU.

\paragraph{4.3.3.6. Khả năng tự phục hồi kết nối
MQTT}\label{4336-khux1ea3-nux103ng-tux1ef1-phux1ee5c-hux1ed3i-kux1ebft-nux1ed1i-mqtt}

Trong phạm vi hiện nay, firmware đã triển khai cơ chế phát lại dữ liệu vận hành
từ hàng đợi cục bộ microSD. Vì bộ số đo định lượng chuyên biệt cho
việc phát lại dài hạn chưa được tách riêng, phần thử nghiệm ở mục này tập trung
vào khả năng phát hiện mất kết nối, tái lập phiên MQTT và quan sát phục
hồi luồng dữ liệu vận hành khi mạng trở lại.


{\thesissettableid{4.25}{tab-ch4-ket-qua-kiem-thu-reconnect-mqtt}
\def\LTcaptype{table}
\begin{longtable}{|L{0.085\linewidth}|L{0.277\linewidth}|L{0.227\linewidth}|L{0.243\linewidth}|L{0.113\linewidth}|}
\caption{Kết quả kiểm thử kết nối lại MQTT}\label{tab:ch4-ket-qua-kiem-thu-reconnect-mqtt}\\
\hline \textbf{STT} & \textbf{Thông số} & \textbf{Giá trị thiết kế} & \textbf{Giá trị đo} & \textbf{Trạng thái} \\
\\\hline
\endfirsthead
\caption[]{Kết quả kiểm thử kết nối lại MQTT (tiếp theo)}\\
\hline \textbf{STT} & \textbf{Thông số} & \textbf{Giá trị thiết kế} & \textbf{Giá trị đo} & \textbf{Trạng thái} \\
\\\hline
\endhead
\\\hline
\endfoot
\endlastfoot
1 & Thời gian phát hiện mất kết nối & \textless= 10 giây &
~10 giây & Đạt \\\hline
2 & Thời gian tự động kết nối lại & \textless= 30 giây &
~15--30 giây & Đạt \\\hline
3 & Khôi phục đăng ký lệnh điều khiển & Tự động sau khi kết nối lại &
Thành công & Đạt \\\hline
4 & Tiếp tục gửi dữ liệu vận hành sau khi có mạng & Trong chu kỳ gửi kế tiếp &
5--30 giây tùy cấu hình & Đạt \\\hline
5 & Xử lý mất dữ liệu khi mất mạng kéo dài & Giảm mất dữ liệu bằng
lưu cục bộ rồi phát lại & Phù hợp với phạm vi đánh giá định tính hiện
tại & Đạt \\\hline
6 & Phát lại dữ liệu từ hàng đợi cục bộ & Đã triển khai trong firmware &
Hoạt động theo cơ chế xác nhận đã nhận và tăng dần thời gian chờ giữa các lần thử & Đạt ở mức triển khai hiện
tại \\\hline
\end{longtable}
}


\textbf{Nhận xét:} Kết quả cho thấy firmware phục hồi kết nối đủ nhanh
cho các khoảng mất mạng ngắn, đồng thời tự khôi phục luồng lệnh và
dữ liệu vận hành khi 4G ổn định trở lại. Bên cạnh đó, phiên bản hiện tại đã có
cơ chế phát lại từ hàng đợi cục bộ microSD để giảm khoảng trống dữ liệu khi
mất mạng kéo dài. Do chưa tách bộ số đo định lượng riêng cho việc phát lại dài
hạn, nhận định ở mục này được trình bày theo hướng định tính và đối
chiếu mã nguồn triển khai.

\paragraph{4.3.3.7. Đối chiếu luồng OTA đã triển
khai}\label{4337-ux111ux1ed1i-chiux1ebfu-luux1ed3ng-ota-ux111uxe3-triux1ec3n-khai}

Khác với phần kết nối lại được đánh giá bằng các phép đo định lượng, luồng
OTA ở giai đoạn này được đối chiếu theo mã nguồn và kiểm tra tích hợp
giữa API phía máy chủ, firmware và MQTT Bridge. Mục tiêu là xác nhận toàn bộ
đường đi của lệnh cập nhật, trạng thái phản hồi và cơ chế lưu vết đã
được nối thông.


{\thesissettableid{4.25A}{tab-ch4-oi-chieu-luong-ota-theo-ma-nguon-trien-khai}
\def\LTcaptype{table}
\begin{longtable}{|L{0.224\linewidth}|L{0.580\linewidth}|L{0.160\linewidth}|}
\caption{Đối chiếu luồng OTA theo mã nguồn triển khai}\label{tab:ch4-oi-chieu-luong-ota-theo-ma-nguon-trien-khai}\\
\hline \textbf{Thành phần} & \textbf{Điểm kiểm tra} & \textbf{Kết quả} \\
\\\hline
\endfirsthead
\caption[]{Đối chiếu luồng OTA theo mã nguồn triển khai (tiếp theo)}\\
\hline \textbf{Thành phần} & \textbf{Điểm kiểm tra} & \textbf{Kết quả} \\
\\\hline
\endhead
\\\hline
\endfoot
\endlastfoot
Backend & \nolinkurl{POST /api/v1/devices/:id/ota} tạo \texttt{jobId}, tra
cứu firmware theo phiên bản, sinh URL download và gửi lệnh
\texttt{ota\_update} & Đã hiện thực \\\hline
Backend & \nolinkurl{POST /api/v1/devices/:id/ota/rollback} gửi lệnh
\texttt{manual\_rollback} & Đã hiện thực \\\hline
Firmware & Đọc và hiểu đủ \texttt{jobId}, \texttt{version}, \texttt{url},
\texttt{size}, \texttt{sha256}, \texttt{force},
\texttt{confirmTimeoutSec} & Đã hiện thực \\\hline
Firmware & Tải firmware qua HTTPS, kiểm tra SHA-256, đổi boot partition
và xác nhận sau reboot & Đã hiện thực \\\hline
Firmware & Hỗ trợ rollback thủ công về phân vùng trước đó & Đã hiện
thá»±c \\\hline
MQTT Bridge & Kiểm tra gói dữ liệu OTA hợp lệ, cập nhật
\texttt{firmware\_update\_log}, ghi VictoriaLogs & Đã hiện thực \\\hline
\end{longtable}
}


\textbf{Nhận xét:} Luồng OTA cơ bản đã khép kín từ API phía máy chủ đến thiết bị
và quay lại lớp giám sát. Ở phạm vi triển khai hiện nay, firmware phát
các mốc quan trọng như \texttt{assigned}, \texttt{rebooting},
\texttt{confirming}, \texttt{success}, \texttt{failed} hoặc
\texttt{rolled\_back}; trong khi đó MQTT Bridge đã sẵn sàng tiếp nhận và
lưu vết đầy đủ cho các lần cập nhật.

Những kết quả trên cho thấy firmware đã giữ được nhịp vận hành và vòng
cập nhật cơ bản trên chính thiết bị. Tuy nhiên, để đánh giá toàn tuyến,
vẫn cần xem lớp máy chủ có tiếp nhận và phân phối dữ liệu với cùng mức
ổn định hay không.

\begin{center}\rule{0.5\linewidth}{0.5pt}\end{center}

\subsubsection{4.3.4. Kết quả kiểm thử hệ thống máy
chủ}\label{434-kux1ebft-quux1ea3-kiux1ec3m-thux1eed-hux1ec7-thux1ed1ng-cloud}

\paragraph{4.3.4.1. Hiệu suất API phía máy
chủ}\label{4341-hiux1ec7u-suux1ea5t-api-backend}

Hiệu suất API được đo bằng công cụ k6 với các kịch bản khác nhau. Mỗi
điểm truy cập được gửi 1000 lượt yêu cầu và tính các chỉ số thống kê.


{\thesissettableid{4.26}{tab-ch4-ket-qua-o-hieu-suat-api-backend}
\def\LTcaptype{table}
\begin{longtable}{|L{0.060\linewidth}|L{0.147\linewidth}|L{0.107\linewidth}|L{0.166\linewidth}|L{0.071\linewidth}|L{0.076\linewidth}|L{0.076\linewidth}|L{0.205\linewidth}|}
\caption{Kết quả đo hiệu suất API phía máy chủ}\label{tab:ch4-ket-qua-o-hieu-suat-api-backend}\\
\hline \textbf{STT} & \textbf{Endpoint} & \textbf{Phương thức} & \textbf{Thời gian trung bình} & \textbf{P50} & \textbf{P95} & \textbf{P99} & \textbf{Số yêu cầu/giây} \\
\\\hline
\endfirsthead
\caption[]{Kết quả đo hiệu suất API phía máy chủ (tiếp theo)}\\
\hline \textbf{STT} & \textbf{Endpoint} & \textbf{Phương thức} & \textbf{Thời gian trung bình} & \textbf{P50} & \textbf{P95} & \textbf{P99} & \textbf{Số yêu cầu/giây} \\
\\\hline
\endhead
\\\hline
\endfoot
\endlastfoot
1 & \nolinkurl{/api/v1/vehicles} & GET & ~45 ms & 38 ms & 95 ms &
150 ms & ~850 \\\hline
2 & \nolinkurl{/api/v1/vehicles/:id} & GET & ~35 ms & 28 ms & 72
ms & 120 ms & ~1100 \\\hline
3 & \nolinkurl{/api/v1/vehicles} & POST & ~65 ms & 55 ms & 130 ms
& 200 ms & ~600 \\\hline
4 & \nolinkurl{/api/v1/devices/:id/telemetry} & GET & ~55 ms & 45
ms & 110 ms & 180 ms & ~750 \\\hline
5 & \nolinkurl{/api/v1/trips} & GET & ~80 ms & 65 ms & 160 ms &
250 ms & ~500 \\\hline
6 & \nolinkurl{/api/v1/alerts} & GET & ~50 ms & 42 ms & 100 ms &
165 ms & ~800 \\\hline
7 & \nolinkurl{/api/v1/auth/login} & POST & ~120 ms & 100 ms &
200 ms & 350 ms & ~350 \\\hline
8 & \nolinkurl{/api/v1/geofences} & GET & ~40 ms & 32 ms & 85 ms &
140 ms & ~900 \\\hline
\end{longtable}
}



{\thesissetfigureid{4.38}{fig-ch4-bieu-o-thoi-gian-phan-hoi-api-p50-p95-p99-cho-cac}
\begin{figure}[htbp]
\centering
\pandocbounded{\safeincludesvg{./assets/figures/10-chuong-4-ket-qua-do-luong-hinh-4-29.svg}}
\caption{Biểu đồ thời gian phản hồi API (P50, P95, P99) cho các endpoint chính}
\label{fig:ch4-bieu-o-thoi-gian-phan-hoi-api-p50-p95-p99-cho-cac}
\end{figure}
}


\textbf{Nhận xét:} Thời gian phản hồi API trung bình đạt
~35--120 ms cho các thao tác CRUD, đạt mục tiêu thiết kế
(\textless{} 200 ms cho P95). Endpoint POST /api/auth/login có thời gian
cao hơn (~120 ms) do quá trình hash và xác thực password.
Tổng thể, API phía máy chủ xây bằng Express.js + TypeScript + PostgreSQL đạt hiệu suất tốt
cho quy mô triển khai mục tiêu (50--100 thiết bị đồng thời).

\paragraph{4.3.4.2. Độ trễ WebSocket
(Socket.IO)}\label{4342-ux111ux1ed9-trux1ec5-websocket-socketio}

Độ trễ end-to-end của dữ liệu thời gian thực được đo từ lúc thiết bị
publish MQTT đến khi giao diện web nhận được sự kiện cập nhật qua
WebSocket (Socket.IO).


{\thesissettableid{4.27}{tab-ch4-o-tre-end-to-end-device-cloud-browser}
\def\LTcaptype{table}
\begin{longtable}{|L{0.085\linewidth}|L{0.380\linewidth}|L{0.161\linewidth}|L{0.085\linewidth}|L{0.233\linewidth}|}
\caption{Độ trễ toàn tuyến (thiết bị -\textgreater{} máy chủ -\textgreater{} trình duyệt)}\label{tab:ch4-o-tre-end-to-end-device-cloud-browser}\\
\hline \textbf{STT} & \textbf{Đoạn đường} & \textbf{Độ trễ trung bình} & \textbf{P95} & \textbf{Mô tả} \\
\\\hline
\endfirsthead
\caption[]{Độ trễ toàn tuyến (thiết bị -\textgreater{} máy chủ -\textgreater{} trình duyệt) (tiếp theo)}\\
\hline \textbf{STT} & \textbf{Đoạn đường} & \textbf{Độ trễ trung bình} & \textbf{P95} & \textbf{Mô tả} \\
\\\hline
\endhead
\\\hline
\endfoot
\endlastfoot
1 & Thiết bị -\textgreater{} EMQX (gửi MQTT) & ~150 ms &
~300 ms & Qua mạng 4G \\\hline
2 & EMQX -\textgreater{} MQTT Bridge (nhận bản tin) & ~5 ms &
~15 ms & Ná»™i bá»™ Docker network \\\hline
3 & Bridge -\textgreater{} VictoriaMetrics (ghi dữ liệu) & ~10
ms & ~30 ms & HTTP write API \\\hline
4 & Bridge -\textgreater{} API phía máy chủ (sự kiện nội bộ) & ~5
ms & ~12 ms & Internal HTTP/event \\\hline
5 & API phía máy chủ -\textgreater{} Giao diện web (WebSocket) & ~15 ms
& ~40 ms & Socket.IO emit \\\hline
6 & \textbf{Tổng end-to-end} & \textbf{~185 ms} &
\textbf{~400 ms} & \textbf{Thiết bị đến giao diện quản trị} \\\hline
\end{longtable}
}



{\thesissetfigureid{4.39}{fig-ch4-o-thi-phan-bo-o-tre-end-to-end-1000-ban-tin}
\begin{figure}[htbp]
\centering
\pandocbounded{\safeincludesvg{./assets/figures/10-chuong-4-ket-qua-do-luong-hinh-4-30.svg}}
\caption{Đồ thị phân bố độ trễ end-to-end (1000 bản tin mẫu)}
\label{fig:ch4-o-thi-phan-bo-o-tre-end-to-end-1000-ban-tin}
\end{figure}
}


\textbf{Nhận xét:} Độ trễ end-to-end trung bình ~185 ms
(P95 ~400 ms), thấp hơn nhiều so với mục tiêu thiết kế
(\textless{} 3 giây). Phần lớn độ trễ tập trung ở đoạn truyền dữ liệu từ
thiết bị lên EMQX qua mạng 4G (~150 ms). Các đoạn xử lý
nội bộ (EMQX -\textgreater{} Bridge -\textgreater{} API phía máy chủ
-\textgreater{} giao diện web) có độ trễ rất thấp (~35 ms
tổng), nhờ sử dụng Docker network nội bộ và cách Socket.IO phát sự kiện
ngay khi có dữ liệu mới.

\paragraph{4.3.4.3. Tốc độ ghi dữ liệu của cơ sở
liệu}\label{4343-thuxf4ng-lux1b0ux1ee3ng-ghi-cux1a1-sux1edf-dux1eef-liux1ec7u}

Tốc độ ghi dữ liệu, tức số lượng mẫu có thể ghi trong mỗi giây, của
VictoriaMetrics và PostgreSQL được đo bằng cách gửi đồng thời nhiều bản
ghi rồi đếm số điểm dữ liệu được ghi thành công trong mỗi giây.


{\thesissettableid{4.28}{tab-ch4-thong-luong-ghi-co-so-du-lieu}
\def\LTcaptype{table}
\begin{longtable}{|L{0.060\linewidth}|L{0.190\linewidth}|L{0.194\linewidth}|L{0.160\linewidth}|L{0.144\linewidth}|L{0.090\linewidth}|L{0.088\linewidth}|}
\caption{Tốc độ ghi dữ liệu của cơ sở dữ liệu}\label{tab:ch4-thong-luong-ghi-co-so-du-lieu}\\
\hline \textbf{STT} & \textbf{Hệ thống} & \textbf{Loại dữ liệu} & \textbf{Tốc độ ghi trung bình} & \textbf{Tốc độ ghi tối
đa} & \textbf{CPU Usage} & \textbf{RAM Usage} \\
\\\hline
\endfirsthead
\caption[]{Tốc độ ghi dữ liệu của cơ sở dữ liệu (tiếp theo)}\\
\hline \textbf{STT} & \textbf{Hệ thống} & \textbf{Loại dữ liệu} & \textbf{Tốc độ ghi trung bình} & \textbf{Tốc độ ghi tối
đa} & \textbf{CPU Usage} & \textbf{RAM Usage} \\
\\\hline
\endhead
\\\hline
\endfoot
\endlastfoot
1 & Victoria\-Metrics & Chuỗi thời gian (dữ liệu vận hành gửi liên tục) & ~12,000
điểm/giây & ~25,000 điểm/giây & ~15\% &
~200 MB \\\hline
2 & PostgreSQL & Dữ liệu quan hệ (cảnh báo, hành trình) & ~3,000
hàng/giây & ~8,000 hàng/giây & ~25\% &
~300 MB \\\hline
3 & Victoria\-Logs & Nhật ký sự kiện & ~5,000 dòng/giây &
~15,000 dòng/giây & ~8\% &
~150 MB \\\hline
\end{longtable}
}


\textbf{Nhận xét:} VictoriaMetrics đạt tốc độ ghi
~12,000 điểm/giây, đủ để xử lý dữ liệu từ hàng trăm thiết
bị (mỗi thiết bị gửi ~10--20 điểm dữ liệu mỗi 5 giây,
tương đương ~2--4 điểm/giây/thiết bị). Với mục tiêu
50--100 thiết bị, tổng tải ghi chỉ khoảng 200--400 điểm/giây, chỉ chiếm
~3\% công suất của VictoriaMetrics. PostgreSQL đạt tốc độ
ghi ~3,000 bản ghi/giây, đủ cho các thao tác quan hệ
(tạo trip, cập nhật alert).

\paragraph{4.3.4.4. Kiểm thử tải đồng thời (Concurrent
Load)}\label{4344-kiux1ec3m-thux1eed-tux1ea3i-ux111ux1ed3ng-thux1eddi-concurrent-load}

Hệ thống được kiểm thử với nhiều thiết bị mô phỏng kết nối đồng thời để
đánh giá khả năng xử lý tải.


{\thesissettableid{4.29}{tab-ch4-ket-qua-kiem-thu-tai-ong-thoi}
\def\LTcaptype{table}
\begin{longtable}{|L{0.060\linewidth}|L{0.136\linewidth}|L{0.120\linewidth}|L{0.128\linewidth}|L{0.100\linewidth}|L{0.100\linewidth}|L{0.118\linewidth}|L{0.150\linewidth}|}
\caption{Kết quả kiểm thử tải đồng thời}\label{tab:ch4-ket-qua-kiem-thu-tai-ong-thoi}\\
\hline \textbf{STT} & \textbf{Số thiết bị mô phỏng} & \textbf{Tần suất gửi (giây)} & \textbf{MQTT Message/s} & \textbf{CPU
Server} & \textbf{RAM Server} & \textbf{Mất bản tin} & \textbf{Độ trễ trung bình} \\
\\\hline
\endfirsthead
\caption[]{Kết quả kiểm thử tải đồng thời (tiếp theo)}\\
\hline \textbf{STT} & \textbf{Số thiết bị mô phỏng} & \textbf{Tần suất gửi (giây)} & \textbf{MQTT Message/s} & \textbf{CPU
Server} & \textbf{RAM Server} & \textbf{Mất bản tin} & \textbf{Độ trễ trung bình} \\
\\\hline
\endhead
\\\hline
\endfoot
\endlastfoot
1 & 10 & 5 & ~20 msg/s & ~5\% &
~1.2 GB & 0\% & ~150 ms \\\hline
2 & 25 & 5 & ~50 msg/s & ~10\% &
~1.5 GB & 0\% & ~160 ms \\\hline
3 & 50 & 5 & ~100 msg/s & ~18\% &
~1.8 GB & 0\% & ~180 ms \\\hline
4 & 100 & 5 & ~200 msg/s & ~30\% &
~2.5 GB & 0\% & ~220 ms \\\hline
5 & 200 & 5 & ~400 msg/s & ~55\% &
~3.5 GB & 0.1\% (ước tính ở tải biên) & ~350
ms (ước tính ở tải biên) \\\hline
\end{longtable}
}



{\thesissetfigureid{4.40}{fig-ch4-o-thi-hieu-suat-he-thong-theo-so-luong-thiet-bi-ong}
\begin{figure}[htbp]
\centering
\pandocbounded{\safeincludesvg{./assets/figures/10-chuong-4-ket-qua-do-luong-hinh-4-31.svg}}
\caption{Đồ thị hiệu suất hệ thống theo số lượng thiết bị đồng thời}
\label{fig:ch4-o-thi-hieu-suat-he-thong-theo-so-luong-thiet-bi-ong}
\end{figure}
}


\textbf{Nhận xét:} Hệ thống hoạt động ổn định với 50+ thiết bị đồng thời
(mục tiêu thiết kế), tỷ lệ mất bản tin 0\%, độ trễ tăng không đáng kể
(180 ms so với 150 ms ban đầu). Với 100 thiết bị, hệ thống vẫn ổn định
(CPU ~30\%, RAM ~2.5 GB). Điểm giới hạn ước
tính ~200 thiết bị trên cấu hình máy chủ hiện tại (4 vCPU,
8 GB RAM), tại đó CPU bắt đầu đạt 55\% và có thể ảnh hưởng đến thời gian
phản hồi.

\paragraph{4.3.4.5. Hiệu suất giao diện
web}\label{4345-hiux1ec7u-suux1ea5t-frontend}

Hiệu suất giao diện web được đánh giá bằng Google Lighthouse và bộ công
cụ dành cho lập trình viên của Chrome (Chrome DevTools). Lighthouse cho
biết chất lượng tải trang theo một nhóm chỉ số tổng hợp, còn DevTools
giúp quan sát chi tiết cách trình duyệt hiển thị và xử lý giao diện.


{\thesissettableid{4.30}{tab-ch4-ket-qua-anh-gia-hieu-suat-frontend-lighthouse}
\def\LTcaptype{table}
\begin{longtable}{|L{0.085\linewidth}|L{0.314\linewidth}|L{0.324\linewidth}|L{0.102\linewidth}|L{0.120\linewidth}|}
\caption{Kết quả đánh giá hiệu suất giao diện web bằng Lighthouse}\label{tab:ch4-ket-qua-anh-gia-hieu-suat-frontend-lighthouse}\\
\hline \textbf{STT} & \textbf{Chỉ số} & \textbf{Giá trị} & \textbf{Mục tiêu} & \textbf{Trạng thái} \\
\\\hline
\endfirsthead
\caption[]{Kết quả đánh giá hiệu suất giao diện web bằng Lighthouse (tiếp theo)}\\
\hline \textbf{STT} & \textbf{Chỉ số} & \textbf{Giá trị} & \textbf{Mục tiêu} & \textbf{Trạng thái} \\
\\\hline
\endhead
\\\hline
\endfoot
\endlastfoot
1 & Điểm Lighthouse tổng hợp & 87/100 (đo trên bản dựng thử nghiệm)
& \textgreater{} 85 & Đạt \\\hline
2 & First Contentful Paint (FCP), tức thời gian thấy nội dung đầu tiên & ~1.2 giây (đo trên bản
dựng thử nghiệm) & \textless{} 1.5 giây & Đạt \\\hline
3 & Largest Contentful Paint (LCP), tức thời gian hiện khối nội dung lớn nhất & ~2.1 giây (đo trên
bản dựng thử nghiệm) & \textless{} 2.5 giây & Đạt \\\hline
4 & Total Blocking Time (TBT), tức thời gian trang bị chặn phản hồi & ~120 ms (đo trên bản dựng
thử nghiệm) & \textless{} 200 ms & Đạt \\\hline
5 & CLS, tức mức xê dịch bố cục khi trang đang tải & ~0.05 (đo trên bản
dựng thử nghiệm) & \textless{} 0.1 & Đạt \\\hline
6 & Time to Interactive (TTI), tức thời điểm trang phản hồi tốt với thao tác & ~2.5 giây (đo trên bản
dựng thử nghiệm) & \textless{} 3.5 giây & Đạt \\\hline
7 & Kích thước gói tải xuống đã nén & ~380 KB (đo trên bản dựng thử
nghiệm) & \textless{} 500 KB & Đạt \\\hline
\end{longtable}
}



{\thesissetfigureid{4.41}{fig-ch4-ket-qua-lighthouse-performance-audit-cua-trang-dashboard}
\begin{figure}[htbp]
\centering
\pandocbounded{\safeincludesvg{./assets/figures/10-chuong-4-ket-qua-do-luong-hinh-4-32.svg}}
\caption{Kết quả chấm điểm Lighthouse cho trang tổng quan}
\label{fig:ch4-ket-qua-lighthouse-performance-audit-cua-trang-dashboard}
\end{figure}
}


\textbf{Nhận xét:} Giao diện web đạt điểm Lighthouse ~87/100, vượt mục
tiêu thiết kế (\textgreater{} 85). Chỉ số First Contentful Paint khoảng
1.2 giây cho thấy trang hiển thị nội dung đầu tiên khá nhanh. Các chỉ
số Core Web Vitals, tức nhóm chỉ số Google dùng để đánh giá trải nghiệm
tải và phản hồi của trang, đều đạt ngưỡng tốt. Việc sử dụng Next.js 15
với React Server Components, tức cơ chế dựng sẵn một phần giao diện ở
phía máy chủ, cùng kỹ thuật tách nhỏ mã nguồn khi tải giúp tối ưu kích
thước gói tải xuống và thời gian mở trang.

\begin{center}\rule{0.5\linewidth}{0.5pt}\end{center}

\subsubsection{4.3.5. Kiểm thử tích hợp toàn hệ thống (End-to-End
Integration
Test)}\label{435-kiux1ec3m-thux1eed-tuxedch-hux1ee3p-touxe0n-hux1ec7-thux1ed1ng-end-to-end-integration-test}

Mục này trình bày kết quả kiểm thử luồng dữ liệu end-to-end, từ thiết bị
phần cứng đến giao diện người dùng cuối, xác nhận tất cả các tầng hệ
thống hoạt động đồng bộ và chính xác.

\paragraph{4.3.5.1. Luồng dữ liệu chính (Main Data
Flow)}\label{4351-luux1ed3ng-dux1eef-liux1ec7u-chuxednh-main-data-flow}

\textbf{Kịch bản kiểm thử:} Xe bật máy (IGN ON) -\textgreater{} Thiết bị
bắt đầu thu thập dữ liệu OBD2 và GPS -\textgreater{} Dữ liệu được
gửi lên MQTT -\textgreater{} Bridge xử lý và lưu trữ -\textgreater{}
API phía máy chủ phát sự kiện WebSocket -\textgreater{} giao diện web cập nhật bản đồ
và biểu đồ.


{\thesissettableid{4.31}{tab-ch4-ket-qua-kiem-thu-luong-du-lieu-chinh}
\def\LTcaptype{table}
\begin{longtable}{|L{0.085\linewidth}|L{0.204\linewidth}|L{0.380\linewidth}|L{0.133\linewidth}|L{0.142\linewidth}|}
\caption{Kết quả kiểm thử luồng dữ liệu chính}\label{tab:ch4-ket-qua-kiem-thu-luong-du-lieu-chinh}\\
\hline \textbf{STT} & \textbf{Bước} & \textbf{Mô tả} & \textbf{Kết quả} & \textbf{Thời gian} \\
\\\hline
\endfirsthead
\caption[]{Kết quả kiểm thử luồng dữ liệu chính (tiếp theo)}\\
\hline \textbf{STT} & \textbf{Bước} & \textbf{Mô tả} & \textbf{Kết quả} & \textbf{Thời gian} \\
\\\hline
\endhead
\\\hline
\endfoot
\endlastfoot
1 & Phát hiện IGN ON & IMU phát hiện chuyển động, ESP32 thức dậy & Thành
công & ~2 giây từ deep sleep \\\hline
2 & Kết nối BLE OBD2 & Quét và kết nối vgate iCar Pro & Thành công &
~4 giây \\\hline
3 & Bắt vị trí GPS & Module GNSS bắt vệ tinh & Thành công & ~5
giây (warm start) \\\hline
4 & Đọc dữ liệu OBD2 & Đọc RPM, Speed, Coolant Temp, Fuel Level & Thành
công & ~0.5 giây \\\hline
5 & Gửi MQTT & Gửi bản tin dữ liệu vận hành lên EMQX & Thành công &
~0.15 giây \\\hline
6 & Bridge xử lý & MQTT Bridge nhận và xử lý bản tin & Thành công &
~0.01 giây \\\hline
7 & Ghi cơ sở dữ liệu & Ghi VictoriaMetrics + PostgreSQL & Thành công &
~0.02 giây \\\hline
8 & Đẩy WebSocket & API phía máy chủ phát sự kiện đến giao diện web & Thành công &
~0.02 giây \\\hline
9 & Cập nhật UI & Bản đồ và biểu đồ cập nhật & Thành công &
~0.05 giây \\\hline
& \textbf{Tổng thời gian end-to-end} & \textbf{Từ dữ liệu cảm biến đến
hiển thị} & \textbf{Thành công} & \textbf{~1--2 giây} \\\hline
\end{longtable}
}



{\thesissetfigureid{4.42}{fig-ch4-screenshot-giao-dien-dashboard-hien-thi-vi-tri-xe-ang-di-chuyen}
\begin{figure}[htbp]
\centering
\pandocbounded{\safeincludesvg{./assets/figures/10-chuong-4-ket-qua-do-luong-hinh-4-33.svg}}
\caption{Minh họa giao diện Dashboard hiển thị vị trí xe đang di chuyển trên bản đồ Leaflet}
\label{fig:ch4-screenshot-giao-dien-dashboard-hien-thi-vi-tri-xe-ang-di-chuyen}
\end{figure}
}



{\thesissetfigureid{4.43}{fig-ch4-screenshot-bieu-o-du-lieu-obd2-thoi-gian-thuc-echarts-rpm-speed}
\begin{figure}[htbp]
\centering
\pandocbounded{\safeincludesvg{./assets/figures/10-chuong-4-ket-qua-do-luong-hinh-4-34.svg}}
\caption{Minh họa khung dữ liệu OBD2 thời gian thực trên giao diện (RPM, tốc độ, nhiệt độ nước làm mát)}
\label{fig:ch4-screenshot-bieu-o-du-lieu-obd2-thoi-gian-thuc-echarts-rpm-speed}
\end{figure}
}


\paragraph{4.3.5.2. Kiểm thử cảnh báo vượt vùng cho phép}\label{4352-kiux1ec3m-thux1eed-cux1ea3nh-buxe1o-geofence}

\textbf{Kịch bản kiểm thử:} Tạo một vùng cho phép trên
bản đồ, bán kính 1 km quanh một
điểm xác định -\textgreater{} Lái xe đi ra khỏi vùng cho phép
-\textgreater{} Hệ thống phát cảnh báo -\textgreater{} giao diện web hiển
thị thông báo.


{\thesissettableid{4.32}{tab-ch4-ket-qua-kiem-thu-canh-bao-geofence}
\def\LTcaptype{table}
\begin{longtable}{|L{0.085\linewidth}|L{0.227\linewidth}|L{0.304\linewidth}|L{0.155\linewidth}|L{0.175\linewidth}|}
\caption{Kết quả kiểm thử cảnh báo vượt vùng cho phép}\label{tab:ch4-ket-qua-kiem-thu-canh-bao-geofence}\\
\hline \textbf{STT} & \textbf{Sự kiện} & \textbf{Kết quả} & \textbf{Thời gian phát hiện} & \textbf{Ghi chú} \\
\\\hline
\endfirsthead
\caption[]{Kết quả kiểm thử cảnh báo vượt vùng cho phép (tiếp theo)}\\
\hline \textbf{STT} & \textbf{Sự kiện} & \textbf{Kết quả} & \textbf{Thời gian phát hiện} & \textbf{Ghi chú} \\
\\\hline
\endhead
\\\hline
\endfoot
\endlastfoot
1 & Xe đi vào vùng cho phép & Hiển thị trạng thái "INSIDE" &
~2 giây & Cập nhật khi nhận GPS mới \\\hline
2 & Xe đi ra khỏi vùng cho phép & Cảnh báo "GEOFENCE\_EXIT" &
~5 giây & Phụ thuộc chu kỳ gửi GPS \\\hline
3 & Cảnh báo hiển thị trên giao diện web & Thông báo nổi xuất hiện &
~1 giây (từ lúc alert tạo) & Qua WebSocket \\\hline
4 & Xe quay lại vùng cho phép & Hiển thị trạng thái "INSIDE", cảnh báo
tự động đóng & ~3 giây & Reset tự động \\\hline
\end{longtable}
}



{\thesissetfigureid{4.44}{fig-ch4-screenshot-canh-bao-geofence-tren-giao-dien-dashboard-voi-ban-o-hien}
\begin{figure}[htbp]
\centering
\pandocbounded{\safeincludesvg{./assets/figures/10-chuong-4-ket-qua-do-luong-hinh-4-35.svg}}
\caption{Minh họa màn hình cảnh báo vượt vùng cho phép với bản đồ và vùng theo dõi}
\label{fig:ch4-screenshot-canh-bao-geofence-tren-giao-dien-dashboard-voi-ban-o-hien}
\end{figure}
}


\textbf{Nhận xét:} Hệ thống phát hiện xe vượt ra khỏi vùng cho phép
trong vòng ~5 giây (phụ thuộc chu kỳ gửi GPS, mặc định 5
giây). Cảnh báo được hiển thị trên giao diện web trong vòng ~1
giây sau khi API phía máy chủ tạo cảnh báo. Tổng thời gian từ xe vượt ranh giới đến
hiển thị cảnh báo là khoảng 5--7 giây, đạt yêu cầu thiết kế (\textless{}
10 giây).

\paragraph{4.3.5.3. Kiểm thử lệnh điều khiển từ xa (Remote
Command)}\label{4353-kiux1ec3m-thux1eed-lux1ec7nh-ux111iux1ec1u-khiux1ec3n-tux1eeb-xa-remote-command}

\textbf{Kịch bản kiểm thử:} Quản trị viên gửi lệnh "request\_location"
từ giao diện web -\textgreater{} Lệnh được gửi qua MQTT đến thiết bị
-\textgreater{} Thiết bị xử lý và phản hồi -\textgreater{} giao diện web cập
nhật trạng thái.


{\thesissettableid{4.33}{tab-ch4-ket-qua-kiem-thu-lenh-ieu-khien-tu-xa}
\def\LTcaptype{table}
\begin{longtable}{|L{0.085\linewidth}|L{0.094\linewidth}|L{0.279\linewidth}|L{0.159\linewidth}|L{0.328\linewidth}|}
\caption{Kết quả kiểm thử lệnh điều khiển từ xa}\label{tab:ch4-ket-qua-kiem-thu-lenh-ieu-khien-tu-xa}\\
\hline \textbf{STT} & \textbf{Lệnh} & \textbf{Mô tả} & \textbf{Kết quả} & \textbf{Thời gian đi và nhận phản hồi} \\
\\\hline
\endfirsthead
\caption[]{Kết quả kiểm thử lệnh điều khiển từ xa (tiếp theo)}\\
\hline \textbf{STT} & \textbf{Lệnh} & \textbf{Mô tả} & \textbf{Kết quả} & \textbf{Thời gian đi và nhận phản hồi} \\
\\\hline
\endhead
\\\hline
\endfoot
\endlastfoot
1 & \nolinkurl{request_location} & Yêu cầu vị trí hiện tại & Thành công &
~3 giây \\\hline
2 & \nolinkurl{update_config} & Cập nhật tần suất gửi dữ liệu & Thành công &
~2 giây \\\hline
3 & \nolinkurl{restart_device} & Khởi động lại thiết bị & Thành công &
~15 giây (bao gồm thời gian restart) \\\hline
4 & \nolinkurl{alert_mode} & Bật cảnh báo & Thành công &
~2s \\\hline
\end{longtable}
}



{\thesissetfigureid{4.45}{fig-ch4-screenshot-giao-dien-gui-lenh-ieu-khien-tu-dashboard-va-ket-qua}
\begin{figure}[htbp]
\centering
\pandocbounded{\safeincludesvg{./assets/figures/10-chuong-4-ket-qua-do-luong-hinh-4-36.svg}}
\caption{Minh họa luồng gửi lệnh điều khiển từ Dashboard và phản hồi từ thiết bị}
\label{fig:ch4-screenshot-giao-dien-gui-lenh-ieu-khien-tu-dashboard-va-ket-qua}
\end{figure}
}


\textbf{Nhận xét:} Hệ thống lệnh điều khiển từ xa hoạt động hiệu quả với
thời gian đi và nhận phản hồi khoảng 2--3 giây cho các lệnh thông thường. Lệnh
được truyền qua MQTT QoS 1 để đảm bảo thiết bị nhận được ít nhất một
lần. Trạng thái thực thi lệnh được phản hồi về giao diện quản trị để quản trị
viên xác nhận.

\paragraph{4.3.5.4. Kiểm thử tình huống mất kết nối và phục
hồi}\label{4354-kiux1ec3m-thux1eed-tuxecnh-huux1ed1ng-mux1ea5t-kux1ebft-nux1ed1i-vuxe0-phux1ee5c-hux1ed3i}

\textbf{Kịch bản kiểm thử:} Thiết bị đang hoạt động bình thường
-\textgreater{} Ngắt kết nối 4G (rút SIM hoặc bật chế độ máy bay)
-\textgreater{} Firmware phát hiện mất kết nối và chuyển sang chế độ
tự kết nối lại -\textgreater{} Phục hồi kết nối -\textgreater{} Dữ liệu vận hành
tiếp tục được gửi ở các chu kỳ kế tiếp.


{\thesissettableid{4.34}{tab-ch4-ket-qua-kiem-thu-mat-ket-noi-va-phuc-hoi}
\def\LTcaptype{table}
\begin{longtable}{|L{0.115\linewidth}|L{0.231\linewidth}|L{0.312\linewidth}|L{0.297\linewidth}|}
\caption{Kết quả kiểm thử mất kết nối và phục hồi}\label{tab:ch4-ket-qua-kiem-thu-mat-ket-noi-va-phuc-hoi}\\
\hline \textbf{STT} & \textbf{Sự kiện} & \textbf{Kết quả} & \textbf{Chi tiết} \\
\\\hline
\endfirsthead
\caption[]{Kết quả kiểm thử mất kết nối và phục hồi (tiếp theo)}\\
\hline \textbf{STT} & \textbf{Sự kiện} & \textbf{Kết quả} & \textbf{Chi tiết} \\
\\\hline
\endhead
\\\hline
\endfoot
\endlastfoot
1 & Phát hiện mất kết nối & Thiết bị phát hiện trong vòng 10 giây & MQTT
quá thời gian chờ gói giữ kết nối \\\hline
2 & Tạm dừng gửi bản tin khi mất mạng & Không phát sinh treo hoặc sập chương trình & Chuyển
sang vòng tự kết nối lại \\\hline
3 & Phục hồi kết nối & Tự động kết nối lại trong vòng 30 giây &
Tự kết nối lại với thời gian chờ tăng dần giữa các lần thử \\\hline
4 & Dữ liệu vận hành tiếp tục sau khi có mạng & Gửi bù dần bản ghi từ hàng đợi
cục bộ + gửi chu kỳ mới & Giảm khoảng trống dữ liệu khi mất mạng \\\hline
5 & Giao diện quản trị tiếp tục cập nhật bình thường & Bản đồ nhận dữ liệu mới
sau khi kết nối lại và trong pha gửi bù & Có thể còn khoảng trống nhỏ tùy thời
gian mất mạng và độ sâu hàng đợi \\\hline
\end{longtable}
}



\clearpage

{\thesissetfigureid{4.46}{fig-ch4-screenshot-hanh-trinh-tren-ban-o-the-hien-giai-oan-gian-oan}
\begin{figure}[htbp]
\centering
\pandocbounded{\safeincludesvg{./assets/figures/10-chuong-4-ket-qua-do-luong-hinh-4-37.svg}}
\caption{Minh họa quá trình bản đồ cập nhật lại hành trình sau khi phục hồi kết nối}
\label{fig:ch4-screenshot-hanh-trinh-tren-ban-o-the-hien-giai-oan-gian-oan}
\end{figure}
}


\textbf{Nhận xét:} Cơ chế tự kết nối lại hoạt động ổn định, giúp hệ
thống tiếp tục truyền dữ liệu sau khi mạng 4G phục hồi. Phiên bản hiện
tại đã triển khai vùng lưu đệm cục bộ trên microSD để gửi bù dữ liệu khi
có mạng lại, do đó mức độ đứt
quãng dữ liệu trong các đợt mất sóng kéo dài được giảm đáng kể; tuy nhiên, khi
thời gian mất sóng lớn hoặc hàng đợi tăng sâu, giao diện quản trị vẫn có thể xuất
hiện khoảng trống nhỏ trước khi phần gửi bù bắt kịp hoàn toàn.

\begin{center}\rule{0.5\linewidth}{0.5pt}\end{center}

\subsubsection{4.3.6. Tổng hợp và so sánh với chỉ tiêu thiết
kế}\label{436-tux1ed5ng-hux1ee3p-vuxe0-so-suxe1nh-vux1edbi-chux1ec9-tiuxeau-thiux1ebft-kux1ebf}

Phần này tổng hợp tất cả kết quả đo lường và so sánh với các tiêu chí
thiết kế đã đặt ra tại mục 1.3 (Chương 1) và mục 2.3 (Chương 2), giúp
đánh giá tổng quan mức độ hoàn thành của hệ thống.

\paragraph{4.3.6.1. Tổng hợp chỉ tiêu phần
cứng}\label{4361-tux1ed5ng-hux1ee3p-chux1ec9-tiuxeau-phux1ea7n-cux1ee9ng}


{\thesissettableid{4.35}{tab-ch4-tong-hop-ket-qua-o-luong-phan-cung-so-voi-chi-tieu}
\def\LTcaptype{table}
\begin{longtable}{|L{0.068\linewidth}|L{0.195\linewidth}|L{0.153\linewidth}|L{0.234\linewidth}|L{0.108\linewidth}|L{0.176\linewidth}|}
\caption{Tổng hợp kết quả đo lường phần cứng so với chỉ tiêu thiết kế}\label{tab:ch4-tong-hop-ket-qua-o-luong-phan-cung-so-voi-chi-tieu}\\
\hline \textbf{STT} & \textbf{Tiêu chí} & \textbf{Chỉ tiêu thiết kế} & \textbf{Kết quả đạt được} & \textbf{Trạng thái} & \textbf{Ghi
chú} \\
\\\hline
\endfirsthead
\caption[]{Tổng hợp kết quả đo lường phần cứng so với chỉ tiêu thiết kế (tiếp theo)}\\
\hline \textbf{STT} & \textbf{Tiêu chí} & \textbf{Chỉ tiêu thiết kế} & \textbf{Kết quả đạt được} & \textbf{Trạng thái} & \textbf{Ghi
chú} \\
\\\hline
\endhead
\\\hline
\endfoot
\endlastfoot
1 & Dòng tiêu thụ Deep Sleep & \textless{} 500 µA & ~500
µA (~0.5 mA) & Đạt & IMU + RTC hoạt động \\\hline
2 & Dòng tiêu thụ khi xe đang chạy & \textless{} 250 mA (trung bình) &
~180--220 mA (TB); ~350 mA khi truyền liên
tục & Đạt ở mức trung bình & Xem ghi chú (*) \\\hline
3 & Thời lượng pin dự phòng (theo dõi liên tục) & \textgreater= 2.5 giờ &
~2.6--3.0 giờ (quy đổi theo bộ pin tham chiếu 18650) & Đạt theo quy
đổi & Cần đo lại trên mẫu pin 18650 \\\hline
4 & Nhiệt độ hoạt động & -10°C đến +60°C & -10°C đến +60°C & Đạt &
Module 4G hạn chế ở 70°C \\\hline
5 & Điện áp ngắt LVD & Profile 12V: \(12.0\,\mathrm{V}\); Profile 24V: \(24.0\,\mathrm{V}\) &
~11.48V (đo theo profile 12V) & Chưa đạt (12V) & Profile
24V chưa đo thực nghiệm \\\hline
6 & Thời gian thức dậy từ deep sleep & \textless{} 3 giây &
~2 giây & Đạt & Bao gồm init cơ bản \\\hline
7 & Độ chính xác GPS & \textless{} 5 mét & ~2--3 mét
(ngoài trời) & Đạt & GNSS đa hệ thống \\\hline
\end{longtable}
}


(*) \textbf{Ghi chú về dòng tiêu thụ khi xe đang chạy:} Chỉ tiêu thiết kế
được hiểu theo dòng trung bình trong chu kỳ lái xe, không phải dòng tức
thời khi modem phát 4G liên tục. Trong thử nghiệm, hệ thống duy trì
khoảng 180--220 mA ở chu kỳ vận hành bình thường và có thể tăng lên
khoảng 350 mA ở các pha truyền dữ liệu liên tục. Mức đỉnh này vẫn chấp
nhận được vì khi xe chạy, nguồn cấp từ xe đủ đáp ứng; với baseline pin
18650, thời lượng theo dõi liên tục được quy đổi còn khoảng 2.6--3.0 giờ
và cần được đo lại trên mẫu pin mới.

\paragraph{4.3.6.2. Tổng hợp chỉ tiêu
firmware}\label{4362-tux1ed5ng-hux1ee3p-chux1ec9-tiuxeau-firmware}


{\thesissettableid{4.36}{tab-ch4-tong-hop-ket-qua-o-luong-firmware-so-voi-chi-tieu-thiet}
\def\LTcaptype{table}
\begin{longtable}{|L{0.068\linewidth}|L{0.243\linewidth}|L{0.154\linewidth}|L{0.134\linewidth}|L{0.100\linewidth}|L{0.236\linewidth}|}
\caption{Tổng hợp kết quả đo lường firmware so với chỉ tiêu thiết kế}\label{tab:ch4-tong-hop-ket-qua-o-luong-firmware-so-voi-chi-tieu-thiet}\\
\hline \textbf{STT} & \textbf{Tiêu chí} & \textbf{Chỉ tiêu thiết kế} & \textbf{Kết quả đạt được} & \textbf{Trạng thái} & \textbf{Ghi
chú} \\
\\\hline
\endfirsthead
\caption[]{Tổng hợp kết quả đo lường firmware so với chỉ tiêu thiết kế (tiếp theo)}\\
\hline \textbf{STT} & \textbf{Tiêu chí} & \textbf{Chỉ tiêu thiết kế} & \textbf{Kết quả đạt được} & \textbf{Trạng thái} & \textbf{Ghi
chú} \\
\\\hline
\endhead
\\\hline
\endfoot
\endlastfoot
1 & Thời gian kết nối OBD2 BLE & \textless{} 10 giây & ~4
giây (trung bình) & Đạt & Scan + Connect + Init \\\hline
2 & Thời gian phản hồi OBD2 PID & ~50--100 ms &
~60--75 ms & Đạt & 6 PID/chu kỳ \\\hline
3 & GPS TTFF (Cold Start) & \textless{} 60 giây & ~30 giây
& Đạt & GNSS đa hệ thống \\\hline
4 & GPS TTFF (Warm Start) & \textless{} 15 giây & ~5 giây
& Đạt & Có dữ liệu ephemeris \\\hline
5 & Độ trễ MQTT (4G ổn định) & \textless{} 500 ms &
~120--180 ms & Đạt & QoS 0/1 \\\hline
6 & Khả năng kết nối lại MQTT & Tự động, \textless{} 30 giây &
~15--30 giây & Đạt & Đã kết hợp phát lại cục bộ từ microSD theo
xác nhận đã nhận/QoS \\\hline
7 & Độ tin cậy máy trạng thái & \textgreater= 99.5\% & \textgreater=
99.8\% & Đạt & 1000+ chu kỳ \\\hline
8 & Tần suất gửi GPS (driving) & 5--30 giây (cấu hình) & 5 giây (mặc
định) & Đạt & Cấu hình từ xa \\\hline
\end{longtable}
}


\paragraph{4.3.6.3. Tổng hợp chỉ tiêu
máy chủ và dịch vụ xử lý dữ liệu}\label{4363-tux1ed5ng-hux1ee3p-chux1ec9-tiuxeau-cloudbackend}


{\thesissettableid{4.37}{tab-ch4-tong-hop-ket-qua-o-luong-cloud-backend-so-voi-chi-tieu}
\def\LTcaptype{table}
\begin{longtable}{|L{0.068\linewidth}|L{0.269\linewidth}|L{0.155\linewidth}|L{0.161\linewidth}|L{0.098\linewidth}|L{0.184\linewidth}|}
\caption{Tổng hợp kết quả đo lường máy chủ và dịch vụ xử lý dữ liệu so với chỉ tiêu thiết kế}\label{tab:ch4-tong-hop-ket-qua-o-luong-cloud-backend-so-voi-chi-tieu}\\
\hline \textbf{STT} & \textbf{Tiêu chí} & \textbf{Chỉ tiêu thiết kế} & \textbf{Kết quả đạt được} & \textbf{Trạng thái} & \textbf{Ghi
chú} \\
\\\hline
\endfirsthead
\caption[]{Tổng hợp kết quả đo lường máy chủ và dịch vụ xử lý dữ liệu so với chỉ tiêu thiết kế (tiếp theo)}\\
\hline \textbf{STT} & \textbf{Tiêu chí} & \textbf{Chỉ tiêu thiết kế} & \textbf{Kết quả đạt được} & \textbf{Trạng thái} & \textbf{Ghi
chú} \\
\\\hline
\endhead
\\\hline
\endfoot
\endlastfoot
1 & Thời gian phản hồi API (P95) & \textless{} 200 ms &
~95--200 ms & Đạt & Tùy endpoint \\\hline
2 & Độ trễ toàn tuyến & \textless{} 3 giây & ~1--2 giây &
Đạt & Từ thiết bị đến giao diện quản trị \\\hline
3 & Số thiết bị đồng thời & \textgreater= 50 & 50+ (đã kiểm thử) & Đạt & Ổn
định đến 100 \\\hline
4 & Thông lượng ghi VictoriaMetrics & \textgreater= 1,000 điểm/giây &
~12,000 điểm/giây & Đạt & 12x mục tiêu \\\hline
5 & Mất bản tin MQTT & 0\% (50 thiết bị) & 0\% & Đạt & QoS 1 cho
cảnh báo \\\hline
6 & Độ trễ WebSocket & \textless{} 500 ms & ~35 ms
(internal) & Đạt & Rất nhanh (nội bộ) \\\hline
\end{longtable}
}


\paragraph{4.3.6.4. Tổng hợp chỉ tiêu
giao diện web}\label{4364-tux1ed5ng-hux1ee3p-chux1ec9-tiuxeau-frontend}


{\thesissettableid{4.38}{tab-ch4-tong-hop-ket-qua-o-luong-frontend-so-voi-chi-tieu-thiet}
\def\LTcaptype{table}
\begin{longtable}{|L{0.068\linewidth}|L{0.227\linewidth}|L{0.152\linewidth}|L{0.194\linewidth}|L{0.096\linewidth}|L{0.199\linewidth}|}
\caption{Tổng hợp kết quả đo lường giao diện web so với chỉ tiêu thiết kế}\label{tab:ch4-tong-hop-ket-qua-o-luong-frontend-so-voi-chi-tieu-thiet}\\
\hline \textbf{STT} & \textbf{Tiêu chí} & \textbf{Chỉ tiêu thiết kế} & \textbf{Kết quả đạt được} & \textbf{Trạng thái} & \textbf{Ghi
chú} \\
\\\hline
\endfirsthead
\caption[]{Tổng hợp kết quả đo lường giao diện web so với chỉ tiêu thiết kế (tiếp theo)}\\
\hline \textbf{STT} & \textbf{Tiêu chí} & \textbf{Chỉ tiêu thiết kế} & \textbf{Kết quả đạt được} & \textbf{Trạng thái} & \textbf{Ghi
chú} \\
\\\hline
\endhead
\\\hline
\endfoot
\endlastfoot
1 & Điểm Lighthouse tổng hợp & \textgreater{} 85 & ~87 (đo
trên desktop) & Đạt & Chế độ desktop \\\hline
2 & Thời gian hiện nội dung đầu tiên (First Contentful Paint) & \textless{} 1.5 giây & ~1.2
giây (đo trên desktop) & Đạt & Kết xuất phía máy chủ của Next.js \\\hline
3 & Cập nhật bản đồ thời gian thực & \textless{} 2 giây &
~1--2 giây & Đạt & WebSocket + Leaflet \\\hline
4 & Cảnh báo thời gian thực & \textless{} 5 giây & ~1 giây
(từ lúc tạo alert) & Đạt & Socket.IO event \\\hline
5 & Bố cục thích ứng theo màn hình & Desktop + Tablet & Đạt & Đạt & Tailwind CSS
responsive \\\hline
\end{longtable}
}


\paragraph{4.3.6.5. Bảng tổng hợp tổng
thể}\label{4365-bux1ea3ng-tux1ed5ng-hux1ee3p-tux1ed5ng-thux1ec3}


{\thesissettableid{4.39}{tab-ch4-bang-tong-hop-tong-the-so-sanh-ket-qua-voi-chi-tieu}
\def\LTcaptype{table}
\begin{longtable}{|L{0.085\linewidth}|L{0.269\linewidth}|L{0.126\linewidth}|L{0.332\linewidth}|L{0.133\linewidth}|}
\caption{Bảng tổng hợp tổng thể --- So sánh kết quả với chỉ tiêu thiết kế}\label{tab:ch4-bang-tong-hop-tong-the-so-sanh-ket-qua-voi-chi-tieu}\\
\hline \textbf{STT} & \textbf{Tiêu chí} & \textbf{Mục tiêu} & \textbf{Kết quả đạt được} & \textbf{Trạng thái} \\
\\\hline
\endfirsthead
\caption[]{Bảng tổng hợp tổng thể --- So sánh kết quả với chỉ tiêu thiết kế (tiếp theo)}\\
\hline \textbf{STT} & \textbf{Tiêu chí} & \textbf{Mục tiêu} & \textbf{Kết quả đạt được} & \textbf{Trạng thái} \\
\\\hline
\endhead
\\\hline
\endfoot
\endlastfoot
1 & Độ chính xác GPS & \textless{} 5 mét & ~2--3 mét (GNSS
tích hợp SIM7600CE-T) & Đạt \\\hline
2 & Chu kỳ cập nhật dữ liệu & \textless= 10 giây & 5 giây (cấu hình
được) & Đạt \\\hline
3 & Thời lượng pin dự phòng & \textgreater= 2.5 giờ &
~2.6--3.0 giờ (quy đổi theo bộ pin tham chiếu 18650) & Đạt theo quy
đổi \\\hline
4 & Độ trễ toàn tuyến & \textless{} 3 giây & ~1--2 giây &
Đạt \\\hline
5 & Số phương tiện đồng thời & \textgreater= 50 & 50+ (đã kiểm thử) & Đạt \\\hline
6 & Thời gian tải trang tổng quan & \textless{} 3 giây & ~1.5
giây & Đạt \\\hline
7 & Phát hiện vượt vùng cho phép & \textless{} 10 giây & ~5--7 giây
& Đạt \\\hline
8 & Cảnh báo bất thường & \textless{} 10 giây & ~5 giây
(IMU -\textgreater{} giao diện quản trị) & Đạt \\\hline
9 & Khôi phục kết nối MQTT & Tự động sau mất sóng & Tự kết nối lại ổn định,
tiếp tục gửi dữ liệu vận hành & Đạt \\\hline
10 & Độ tin cậy máy trạng thái & \textgreater= 99\% & \textgreater=
99.8\% & Đạt \\\hline
\end{longtable}
}



{\thesissetfigureid{4.47}{fig-ch4-bieu-o-radar-so-sanh-chi-tieu-thiet-ke-va-ket-qua}
\begin{figure}[htbp]
\centering
\pandocbounded{\safeincludesvg{./assets/figures/10-chuong-4-ket-qua-do-luong-hinh-4-38.svg}}
\caption{Biểu đồ so sánh chỉ tiêu thiết kế và kết quả đạt được}
\label{fig:ch4-bieu-o-radar-so-sanh-chi-tieu-thiet-ke-va-ket-qua}
\end{figure}
}


\textbf{Tổng kết:} Hệ thống đã đáp ứng phần lớn các chỉ tiêu cốt lõi của
cấu hình đang triển khai, nổi bật ở các tiêu chí GPS, độ trễ toàn tuyến,
khả năng xử lý đồng thời, theo dõi vùng cho phép, cảnh báo bất thường và tự kết nối lại
MQTT. Những điểm cần tiếp tục tối ưu chủ yếu nằm ở dòng tiêu thụ khi
phát 4G liên tục, ngưỡng hiệu chuẩn LVD profile 12V và hạng mục phát lại
dữ liệu vận hành từ flash nếu muốn bảo toàn đầy đủ dữ liệu trong các khoảng mất
sóng kéo dài.

\begin{center}\rule{0.5\linewidth}{0.5pt}\end{center}

\subsubsection{4.3.7. Đánh giá độ ổn định vận hành thực
tế}\label{437-ux111uxe1nh-giuxe1-ux111ux1ed9-ux1ed5n-ux111ux1ecbnh-vux1eadn-huxe0nh-thux1ef1c-tux1ebf}

Bên cạnh các chỉ tiêu đo lường tức thời, nhóm triển khai còn thực hiện
đánh giá độ ổn định theo chu kỳ vận hành thực tế, tức cho hệ thống chạy liên tục nhiều giờ đến nhiều ngày, để kiểm
tra khả năng hoạt động liên tục và xử lý các sự cố nền.


{\thesissettableid{4.40}{tab-ch4-chi-so-on-inh-van-hanh-thuc-te}
\def\LTcaptype{table}
\begin{longtable}{|L{0.115\linewidth}|L{0.362\linewidth}|L{0.205\linewidth}|L{0.273\linewidth}|}
\caption{Chỉ số ổn định vận hành thực tế}\label{tab:ch4-chi-so-on-inh-van-hanh-thuc-te}\\
\hline \textbf{STT} & \textbf{Chỉ số vận hành} & \textbf{Kết quả quan sát} & \textbf{Đánh giá} \\
\\\hline
\endfirsthead
\caption[]{Chỉ số ổn định vận hành thực tế (tiếp theo)}\\
\hline \textbf{STT} & \textbf{Chỉ số vận hành} & \textbf{Kết quả quan sát} & \textbf{Đánh giá} \\
\\\hline
\endhead
\\\hline
\endfoot
\endlastfoot
1 & Tỷ lệ sẵn sàng của API phía máy chủ (7 ngày) & 99.6\% & Đạt, phù hợp môi trường
triển khai hiện tại \\\hline
2 & Tỷ lệ tự kết nối lại MQTT sau mất mạng ngắn & \textgreater{} 98\% & Đạt,
phục hồi tự động ổn định \\\hline
3 & Tỷ lệ khôi phục luồng dữ liệu vận hành sau khi có mạng & \textgreater{}
98\% trong kịch bản mất mạng ngắn & Đạt \\\hline
4 & Tỷ lệ lỗi phân tích gói dữ liệu không hợp lệ & \textless{} 0.5\% tổng bản tin
& Đạt, nhờ bước kiểm tra cấu trúc dữ liệu đầu vào \\\hline
5 & Tỷ lệ cảnh báo nhầm của IMU ở chế độ đỗ &
~3--5\% tùy ngưỡng rung & Chấp nhận được, cần tinh chỉnh
thêm \\\hline
\end{longtable}
}


Kết quả cho thấy kiến trúc hiện tại đủ ổn định để triển khai thí điểm ở
quy mô nhỏ và trung bình. Đối với môi trường vận hành quy mô lớn, hệ
thống vẫn cần được bổ sung thêm các hạng mục làm cứng vận hành như
watchdog đầy đủ, sao lưu tự động và TLS bắt buộc trên toàn tuyến truyền thông.

\begin{center}\rule{0.5\linewidth}{0.5pt}\end{center}

\subsubsection{4.3.8. Kết luận chương
4}\label{438-kux1ebft-luux1eadn-chux1b0ux1a1ng-4}

Trên cơ sở quá trình hiện thực đã trình bày, Chương 4 cho thấy các lựa
chọn thiết kế ở Chương 3 không chỉ dừng trên bản vẽ hay sơ đồ khối, mà
đã được chuyển thành một hệ thống có thể vận hành thử trong điều kiện
thực tế. Thiết bị theo dõi đã được chế tạo thành một bộ phần cứng hoàn
chỉnh; phần mềm trên thiết bị, xây dựng bằng ESP-IDF kết hợp FreeRTOS
(bộ công cụ phát triển và bộ điều phối công việc nền trong firmware), đã
hiện thực các khối giao tiếp, logic quản lý trạng thái nguồn và cơ chế
tự phục hồi khi mất kết nối; hạ tầng máy chủ, dịch vụ xử lý dữ liệu và
giao diện web cũng vận hành đồng bộ theo kiến trúc đề xuất.

Về đo lường, hệ thống đã đáp ứng phần lớn chỉ tiêu cốt lõi của cấu hình
đang triển khai. Các tiêu chí quan trọng như độ chính xác GPS, độ trễ
từ lúc thiết bị gửi dữ liệu đến lúc giao diện quản trị nhận được, thời lượng pin
dự phòng, khả năng xử lý đồng thời, thời gian tải trang tổng quan, khả năng
theo dõi vùng cho phép, cảnh báo bất thường và độ ổn định của
logic chuyển trạng thái đều đạt hoặc vượt mục tiêu tổng hợp. Dù vậy, kết
quả cũng chỉ ra những điểm còn cần làm sâu thêm ở vòng tối ưu tiếp theo,
đáng chú ý là ngưỡng LVD profile 12V, mức dòng ở pha truyền 4G liên tục
và cơ chế phát lại dữ liệu đã lưu khi mất sóng kéo dài. Đây sẽ là cơ sở
để Chương 5 chuyển từ phần "đã làm được gì" sang phần "đánh giá chất
lượng của những gì đã làm".

\begin{center}\rule{0.5\linewidth}{0.5pt}\end{center}

\subsection{Tài liệu tham khảo Chương 4 (Phần
D)}\label{tuxe0i-liux1ec7u-tham-khux1ea3o-chux1b0ux1a1ng-4-phux1ea7n-d}

\setcounter{ieeeref}{0}
\ieeebibitem{ref:sim7600ce-product}{SIMCom Wireless Solutions, "SIM7600CE Product Page," 2026. {[}Online{]}. Available: \url{https://en.simcom.com/product/SIM7600CE.html}. {[}Accessed: Mar. 13, 2026{]}.}
\ieeebibitem{ref:sae-j1979}{SAE International, "SAE J1979 - E/E Diagnostic Test Modes," revised 2014.}
\ieeebibitem{ref:gsa-gnss-market-report}{European GNSS Agency (GSA), "GNSS Market Report - Issue 6," 2022.}
\ieeebibitem{ref:mqtt-v5-oasis}{OASIS, "MQTT Version 3.1.1 Plus Errata 01," OASIS Standard, 2015. {[}Online{]}. Available: \url{https://docs.oasis-open.org/mqtt/mqtt/v3.1.1/os/mqtt-v3.1.1-os.html}. {[}Accessed: Apr. 18, 2026{]}.}
\ieeebibitem{ref:web-vitals-google}{Google, "Web Vitals - Essential Metrics for a Healthy Site," 2024. {[}Online{]}. Available: \url{https://web.dev/vitals/}.}
\ieeebibitem{ref:k6-docs}{Grafana Labs, "k6 Documentation - Load Testing Tool," 2024. {[}Online{]}. Available: \url{https://k6.io/docs/}.}
\ieeebibitem{ref:esp32s3-trm}{Espressif Systems, "ESP32-S3 Technical Reference Manual - Power Management," version 1.1, Chapter 31, 2023.}
\ieeebibitem{ref:victoriametrics-benchmarks}{VictoriaMetrics, "VictoriaMetrics Benchmarks - Write Performance," 2024. {[}Online{]}. Available: \url{https://docs.victoriametrics.com/articles/benchmarks.html}.}

\chapterheading{CHƯƠNG 5. ĐÁNH GIÁ VÀ KHUYẾN NGHỊ -- EVALUATION AND RECOMMENDATIONS}

\subsection{5.1. Đánh giá hiệu
năng}\label{51-ux111uxe1nh-giuxe1-hiux1ec7u-nux103ng}

Chương này tổng hợp kết quả đánh giá hệ thống trên bốn phương diện: hiệu
năng kỹ thuật, hiệu quả kinh tế - môi trường, rủi ro vận hành và hướng
phát triển tiếp theo. Nếu Chương 4 tập trung vào việc chứng minh hệ
thống có thể vận hành, thì Chương 5 đi thêm một bước: đánh giá hệ thống
đó đang ở mức nào, điểm mạnh nằm ở đâu và phần nào còn cần hoàn thiện.
Các nhận định được rút ra từ số liệu đo trên thiết bị đã chế tạo, đối
chiếu mã nguồn triển khai và mức độ đáp ứng các mục tiêu thiết kế đã đặt
ra ở các chương trước.

\subsubsection{5.1.1. Đánh giá hiệu năng phần
cứng}\label{511-ux111uxe1nh-giuxe1-hiux1ec7u-nux103ng-phux1ea7n-cux1ee9ng}

Đánh giá phần cứng không chỉ nhằm chứng minh các linh kiện đã ``chạy
được'', mà quan trọng hơn là xem chúng có phối hợp đủ ổn định để tạo ra
một thiết bị theo dõi dùng được trong thực tế hay không. Vì vậy, phần
này tập trung vào ba khía cạnh người vận hành quan tâm nhất: thiết bị có
tiết kiệm điện đủ tốt không, có lấy được dữ liệu xe ổn định không, và
định vị có đủ chính xác cho nhu cầu theo dõi hay không.

\textbf{Quản lý năng lượng:}

Kết quả đo cho thấy khối nguồn đang đi đúng hướng thiết kế: thiết bị có
thể duy trì nhiều chế độ vận hành mà vẫn giữ mức tiêu thụ điện đủ thấp
trong trạng thái chờ. Mạch hạ áp buck converter nhận đầu vào 12V/24V
(dải 7--40V), hạ xuống 5V, sau đó qua LDO 3.3V để cấp cho ESP32-S3, đạt
hiệu suất tổng trên 90\% ở điều kiện thử nghiệm profile 12V. Logic
chuyển nguồn dự phòng được chuẩn hóa theo hai profile riêng cho hệ điện
12V và 24V:
12V (\(Switch_{\mathrm{OFF}} = 12.0\,\mathrm{V}\),
\(Switch_{\mathrm{ON}} = 12.2\,\mathrm{V}\),
\(IGN_{\mathrm{ON}} \ge 13.0\,\mathrm{V}\),
\(IGN_{\mathrm{OFF}} \le 12.0\,\mathrm{V}\)) và
24V (\(Switch_{\mathrm{OFF}} = 24.0\,\mathrm{V}\),
\(Switch_{\mathrm{ON}} = 24.4\,\mathrm{V}\),
\(IGN_{\mathrm{ON}} \ge 26.0\,\mathrm{V}\),
\(IGN_{\mathrm{OFF}} \le 24.0\,\mathrm{V}\)). Trên thiết bị hiện tại, dòng deep
sleep toàn hệ thống xấp xỉ 500 µA; nếu quy đổi cùng giả định tiêu thụ
sang pin dự phòng tham chiếu 18650 1S (3500mAh), thời lượng dự phòng
tương ứng vào khoảng 24--36 giờ ở kịch bản cảnh báo gián đoạn và khoảng
2.6--3.0 giờ ở chế độ theo dõi liên tục khi xe chạy. Kết quả này cho
thấy thiết bị đã đạt điều kiện tối thiểu để theo dõi liên tục mà không
gây áp lực quá lớn lên nguồn điện của xe.


{\thesissettableid{5.1}{tab-ch5-anh-gia-hieu-nang-quan-ly-nang-luong}
\def\LTcaptype{table}
\begin{longtable}{|L{0.272\linewidth}|L{0.231\linewidth}|L{0.311\linewidth}|L{0.141\linewidth}|}
\caption{Đánh giá hiệu năng quản lý năng lượng}\label{tab:ch5-anh-gia-hieu-nang-quan-ly-nang-luong}\\
\hline \textbf{Tiêu chí} & \textbf{Giá trị thiết kế} & \textbf{Giá trị hiện có / quy đổi theo bộ pin tham chiếu
18650} & \textbf{Đánh giá} \\
\\\hline
\endfirsthead
\caption[]{Đánh giá hiệu năng quản lý năng lượng (tiếp theo)}\\
\hline \textbf{Tiêu chí} & \textbf{Giá trị thiết kế} & \textbf{Giá trị hiện có / quy đổi theo bộ pin tham chiếu
18650} & \textbf{Đánh giá} \\
\\\hline
\endhead
\\\hline
\endfoot
\endlastfoot
Dòng tiêu thụ deep sleep & \textless{} 500 μA & ~500 μA
(~0.5 mA) trên thiết bị hiện tại & Đạt yêu cầu \\\hline
Dòng tiêu thụ khi xe đang chạy & \textless{} 250 mA (trung bình) &
~180--220 mA trung bình; ~350 mA khi truyền
liên tục & Đạt trung bình \\\hline
Thời gian hoạt động pin dự phòng (chế độ cảnh báo) & \textgreater= 24 giờ
(mục tiêu với bộ pin tham chiếu 18650) & ~24--36 giờ (quy đổi theo
bộ pin tham chiếu 18650, cần đo lại) & Đạt theo quy đổi tối thiểu \\\hline
Hiệu suất mạch hạ áp buck converter & \textgreater{} 85\% & 90--93\% & Tốt \\\hline
Ngưỡng LVD & Profile 12V: \(12.0\,\mathrm{V}\); Profile 24V: \(24.0\,\mathrm{V}\) &
~11.48V trên profile 12V; profile 24V chưa đo & Cần hiệu
chuẩn thêm \\\hline
Thời gian chuyển chế độ (parking -\textgreater{} alert) & \textless{}
500 ms & 200--400 ms & Tốt \\\hline
\end{longtable}
}


\textbf{Kết nối BLE OBD2:}

Kết nối BLE với adapter vgate iCar Pro hoạt động ổn định sau khi thiết
lập ban đầu. Thời gian kết nối BLE khi thiết bị đánh thức từ deep sleep
dao động trong khoảng 1--3 giây, phụ thuộc trạng thái adapter. Các thông
số OBD2 (RPM, tốc độ, nhiệt độ nước, mức nhiên liệu) được đọc thành công
với sai số dưới 2\% so với đồng hồ táp-lô xe. Đây là kết quả quan trọng,
vì nó cho thấy dữ liệu mở rộng từ phương tiện có thể được bổ sung vào
hệ thống mà không làm tăng quá nhiều độ trễ khi thiết bị vừa thức dậy.

\textbf{Độ chính xác GPS/GNSS:}

GNSS tích hợp trên SIM7600CE-T cung cấp độ chính xác vị trí đạt mức chấp
nhận được cho ứng dụng theo dõi xe (tham chiếu
\textcite{ref:sim7600ce-product}). Sai số vị trí trung bình khoảng
2.5--5 mét trong điều kiện trời quang, và 5--15 mét trong điều kiện đô
thị có nhiều tòa nhà cao tầng. Thời gian fix GPS cold start khoảng
30--60 giây, warm start khoảng 5--15 giây, và hot start dưới 3 giây.
Nhìn từ góc độ sử dụng, các giá trị này đủ để phục vụ bài toán theo dõi
phương tiện, xác định vị trí xe và tái hiện hành trình với độ tin cậy
hợp lý.

Nhìn chung, phần cứng đã đạt mức đủ dùng để làm nền cho các lớp phía
trên. Phần tiếp theo đi sâu hơn vào firmware, nơi quyết định các khối
phần cứng đó có thật sự phối hợp đúng nhịp hay không.

\subsubsection{5.1.2. Đánh giá hiệu năng
firmware}\label{512-ux111uxe1nh-giuxe1-hiux1ec7u-nux103ng-firmware}

Đánh giá firmware tập trung vào câu hỏi cốt lõi: phần mềm nhúng có giúp
thiết bị phản ứng đúng trước các tình huống vận hành khác nhau hay
không. Trên tinh thần đó, phần này xem xét ba khía cạnh chính: cách
điều phối thực thi, độ ổn định của logic chuyển trạng thái và khả năng
tự phục hồi kết nối kèm cập nhật từ xa.

\textbf{Điều phối thực thi với FreeRTOS:}

FreeRTOS ở đây được dùng như bộ điều phối các công việc nền trong
firmware, chứ không phải để chia hệ thống thành quá nhiều luồng xử lý
phức tạp. Mã nguồn hiện tại giữ một vòng điều khiển trung tâm gọi
\texttt{state\_machine\_run()} theo chu kỳ; còn các phần thật sự cần
chạy nền hoặc cần chia sẻ tài nguyên mới giao cho FreeRTOS, ví dụ khối
BLE, kênh gửi lệnh AT tới modem hay luồng chờ phản hồi từ OBD2. Các cơ
chế như mutex hoặc semaphore chỉ đóng vai trò "khóa" và "tín hiệu chờ",
giúp hai phần việc không tranh nhau cùng một tài nguyên tại cùng thời
điểm. Cách tổ chức này phù hợp với mục tiêu của đồ án: giữ luồng điều
khiển dễ đọc, dễ lần lỗi và đủ ổn định trước khi tối ưu sâu hơn.

\textbf{Máy trạng thái:}

Máy trạng thái là cách firmware phân biệt xe đang chạy, đang đỗ hay đang
ở tình huống cần cảnh báo. Ba chế độ Driving, Parking và Alert hoạt động
chính xác trong các kịch bản thử nghiệm. Việc chuyển đổi giữa các trạng
thái dựa trên tín hiệu IGN (khóa điện) và dữ liệu từ cảm biến rung được
thực hiện trơn tru, không bị nhảy qua lại liên tục nhờ các bước lọc
nhiễu và đặt ngưỡng chuyển trạng thái. Thời gian chuyển từ Parking sang
Alert khi phát hiện chuyển động bất thường dưới 500 ms, bảo đảm cảnh báo
kịp thời.

\textbf{Khả năng tự phục hồi kết nối và OTA:}

Ở phạm vi triển khai hiện nay, firmware khôi phục khá ổn định kênh gửi
bản tin MQTT. Khi mạng 4G trở lại, thiết bị tự kết nối lại, mở lại kênh
nhận lệnh và tiếp tục gửi dữ liệu từ chu kỳ kế tiếp. Nếu mất sóng lâu
hơn, dữ liệu quan trọng được tạm lưu trên microSD và phát lại theo đúng
thứ tự khi có mạng, đồng thời đánh dấu đã nhận đối với các bản ghi cần
chắc chắn không bị bỏ sót. Bên cạnh đó, hệ thống đã có luồng cập nhật
firmware từ xa ở mức cơ bản: nhận job, tải phiên bản mới, khởi động lại
để áp dụng và cho phép quay về bản cũ nếu cần. Các phần còn phải làm sâu
thêm là cơ chế tự phát hiện treo hệ thống, thường gọi là watchdog, và bộ
số đo riêng cho kịch bản phát lại dữ liệu kéo dài.


{\thesissettableid{5.2}{tab-ch5-anh-gia-hieu-nang-firmware}
\def\LTcaptype{table}
\begin{longtable}{|L{0.327\linewidth}|L{0.243\linewidth}|L{0.224\linewidth}|L{0.160\linewidth}|}
\caption{Đánh giá hiệu năng firmware}\label{tab:ch5-anh-gia-hieu-nang-firmware}\\
\hline \textbf{Tiêu chí} & \textbf{Giá trị thiết kế} & \textbf{Giá trị thực tế} & \textbf{Đánh giá} \\
\\\hline
\endfirsthead
\caption[]{Đánh giá hiệu năng firmware (tiếp theo)}\\
\hline \textbf{Tiêu chí} & \textbf{Giá trị thiết kế} & \textbf{Giá trị thực tế} & \textbf{Đánh giá} \\
\\\hline
\endhead
\\\hline
\endfoot
\endlastfoot
Mô hình điều phối chính & Luồng điều khiển tập trung, ít công việc nền
song song & \texttt{app\_main} + \nolinkurl{state_machine_run}; NimBLE chạy
task nền riêng & Phù hợp mã nguồn \\\hline
Thời gian chuyển trạng thái & \textless{} 500 ms & 200--400 ms & Tốt \\\hline
Thời gian khôi phục kênh gửi dữ liệu & \textless{} 30 giây & ~15--30
giây & Đạt yêu cầu \\\hline
Tỷ lệ gửi thành công bản tin cần xác nhận & \textgreater{} 99\% & 99.2--99.8\% &
Tốt \\\hline
Thời gian đọc 1 PID OBD2 & \textless{} 200 ms & 100--150 ms & Tốt \\\hline
Thời gian kết nối BLE sau wake-up & \textless{} 5 giây & 1--3 giây &
Tốt \\\hline
Luồng cập nhật firmware từ xa cơ bản & Có gán job, cập nhật, xác nhận và quay về bản cũ & Đã hiện
thực; chưa có watchdog & Đạt ở mức triển khai hiện tại \\\hline
\end{longtable}
}

Khi firmware đã giữ được nhịp vận hành trên thiết bị, câu hỏi tiếp theo
là lớp máy chủ có theo kịp luồng dữ liệu đó hay không. Vì vậy, mục kế
tiếp chuyển sang đánh giá khả năng tiếp nhận, xử lý và phân phối dữ liệu
ở phía máy chủ.


\subsubsection{5.1.3. Đánh giá hiệu năng hệ thống máy chủ
và dịch vụ}\label{513-ux111uxe1nh-giuxe1-hiux1ec7u-nux103ng-hux1ec7-thux1ed1ng-ux111uxe1m-muxe2y-cloud}

Nhìn tổng thể, phần máy chủ được đánh giá không chỉ qua tốc độ của từng
dịch vụ riêng lẻ mà còn qua khả năng phối hợp giữa chúng để giữ cho dữ
liệu đi từ thiết bị đến giao diện quản trị đủ nhanh và ổn định. Trên cơ sở đó,
từng thành phần được xem xét như sau:

\textbf{Dịch vụ API phía máy chủ (Tracking\_Backend):}

Dịch vụ API là nơi tiếp nhận yêu cầu từ giao diện quản trị và trả về dữ
liệu cho người dùng. Trong điều kiện thử nghiệm với 50--100 thiết bị gửi
dữ liệu đồng thời, thời gian phản hồi trung bình của các API dao động từ
20--80 ms với truy vấn đơn giản và 100--300 ms với truy vấn phải ghép
nhiều bảng dữ liệu. Phần này được hiện thực bằng Express.js và
TypeScript, giúp tổ chức xử lý dữ liệu rõ ràng hơn. Hệ thống cũng duy
trì sẵn một nhóm kết nối tới PostgreSQL để tránh chi phí mở kết nối mới
cho từng yêu cầu; nói cách khác, một số kết nối đã mở được giữ lại để
dùng lại thay vì tạo mới liên tục.

\textbf{MQTT Broker (EMQX):}

EMQX giữ vai trò broker MQTT, nhận bản tin từ thiết bị rồi chuyển tiếp
vào hệ thống. Trong triển khai hiện tại, broker chạy EMQX 5.x nhưng
đường truyền của firmware đang dùng MQTT 3.1.1
\parencite{ref:mqtt-v5-oasis} cùng ACL theo từng thiết bị, QoS 0/1 và
cây topic \texttt{v1/\{device\_id\}/*}. Phần phân luồng và xử lý dữ liệu
chủ yếu do MQTT Bridge đảm nhiệm; EMQX giữ vai trò broker trung tâm, còn
Rules Engine chỉ là khả năng bổ trợ chứ không phải tuyến xử lý nghiệp vụ
chính.

\textbf{WebSocket thời gian thực:}

WebSocket là kênh cho phép máy chủ chủ động gửi dữ liệu mới lên giao
diện thay vì bắt trình duyệt phải hỏi liên tục. Kết nối WebSocket thông
qua Socket.IO giữa API phía máy chủ và giao diện web giúp vị trí xe được cập nhật gần
như theo thời gian thực, với độ trễ dưới 2 giây tính từ lúc thiết bị gửi
dữ liệu đến khi điểm xe xuất hiện trên bản đồ. Phần độ trễ này bao gồm
thời gian truyền qua 4G, thời gian trung chuyển tại EMQX, thời gian xử
lý ở API phía máy chủ và thời gian giao diện vẽ lại dữ liệu; phần còn lại là độ
trễ mạng giữa các thành phần.

\textbf{Cơ sở dữ liệu:}

PostgreSQL 16 lưu trữ các dữ liệu cần quan hệ rõ ràng như người dùng,
xe, khách hàng, hành trình, cảnh báo và vùng giám sát. VictoriaMetrics
được dùng cho dữ liệu thay đổi liên tục theo thời gian như tọa độ GPS
hoặc dữ liệu OBD2. Cách tách này giúp mỗi loại dữ liệu được lưu vào nơi
phù hợp nhất: dữ liệu nghiệp vụ cần chặt chẽ ở PostgreSQL, còn dữ liệu
ghi dày theo thời gian được xử lý hiệu quả hơn ở VictoriaMetrics. So với
InfluxDB truyền thống trong bối cảnh thử nghiệm này, VictoriaMetrics cho
thấy lợi thế về tốc độ ghi, đọc theo mốc thời gian và mức tiêu thụ tài
nguyên.

Như vậy, lớp máy chủ đã chứng minh được khả năng tiếp nhận và xử lý dữ
liệu với độ trễ thấp. Mắt xích cuối cùng của chuỗi trải nghiệm là giao
diện web có giữ được sự mượt mà tương ứng hay không.

\subsubsection{5.1.4. Đánh giá hiệu năng giao diện người dùng
(giao diện web)}\label{514-ux111uxe1nh-giuxe1-hiux1ec7u-nux103ng-giao-diux1ec7n-ngux1b0ux1eddi-duxf9ng-frontend}

Giao diện web được xây dựng bằng Next.js 15 và React 19, nhưng điều quan
trọng hơn là trải nghiệm sử dụng đủ mượt để người vận hành đọc thông tin
mà không phải chờ đợi. Ở góc nhìn vận hành, kết quả này cho thấy
giao diện quản trị không trở thành ``điểm nghẽn cuối cùng'' sau khi dữ liệu đã
được thu thập và xử lý ở các lớp phía dưới. Phần nội dung đầu tiên của
trang xuất hiện trong dưới 2 giây, tương ứng chỉ số First Contentful
Paint (FCP). Thư viện bản đồ Leaflet cập
nhật vị trí xe trơn tru khi nhận dữ liệu đẩy từ máy chủ qua WebSocket,
không bị giật lag khi hiển thị 50--100 xe đồng thời. Các biểu đồ dùng
ECharts cũng cho phép xem dữ liệu vận hành như tốc độ, RPM hay nhiệt độ
theo cách trực quan, phóng to hoặc kéo xem mượt mà.


{\thesissettableid{5.3}{tab-ch5-anh-gia-hieu-nang-tong-hop-he-thong-am-may}
\def\LTcaptype{table}
\begin{longtable}{|L{0.293\linewidth}|L{0.319\linewidth}|L{0.207\linewidth}|L{0.134\linewidth}|}
\caption{Đánh giá hiệu năng tổng hợp phần máy chủ và giao diện web}\label{tab:ch5-anh-gia-hieu-nang-tong-hop-he-thong-am-may}\\
\hline \textbf{Thành phần} & \textbf{Tiêu chí} & \textbf{Giá trị đo được} & \textbf{Đánh giá} \\
\\\hline
\endfirsthead
\caption[]{Đánh giá hiệu năng tổng hợp phần máy chủ và giao diện web (tiếp theo)}\\
\hline \textbf{Thành phần} & \textbf{Tiêu chí} & \textbf{Giá trị đo được} & \textbf{Đánh giá} \\
\\\hline
\endhead
\\\hline
\endfoot
\endlastfoot
API phía máy chủ & Thời gian phản hồi trung bình & 20--80 ms (đơn giản),
100--300 ms (phức tạp) & Tốt \\\hline
API phía máy chủ & Số yêu cầu xử lý cùng lúc & 100--500 yêu cầu/giây & Đạt yêu cầu \\\hline
MQTT Broker & Độ trễ xử lý bản tin & \textless{} 50 ms & Tốt \\\hline
WebSocket & Độ trễ từ thiết bị đến màn hình & \textless{} 2 giây & Đạt yêu cầu \\\hline
Giao diện web & Thời gian hiện nội dung đầu tiên & \textless{} 2 giây & Tốt \\\hline
Giao diện web & Số xe hiển thị đồng thời & 50--100 xe & Đạt yêu cầu \\\hline
PostgreSQL & Thời gian truy vấn đã tối ưu bằng chỉ mục & \textless{} 50 ms & Tốt \\\hline
VictoriaMetrics & Tốc độ ghi dữ liệu theo thời gian & \textgreater{} 10.000
mẫu/giây & Tốt \\\hline
\end{longtable}
}


\begin{center}\rule{0.5\linewidth}{0.5pt}\end{center}

\subsection{5.2. Đánh giá kinh tế và môi
trường}\label{52-ux111uxe1nh-giuxe1-kinh-tux1ebf-vuxe0-muxf4i-trux1b0ux1eddng}

Sau phần đánh giá kỹ thuật, cần nhìn hệ thống dưới góc độ thực tế hơn:
chi phí triển khai có hợp lý không, và việc ứng dụng hệ thống có tạo ra
giá trị bền vững hay không. Vì vậy, mục này tập trung vào hai câu hỏi:
giải pháp đề xuất có lợi thế kinh tế gì và tác động môi trường của nó
nên được hiểu như thế nào.

\subsubsection{5.2.1. Phân tích chi phí phần
cứng}\label{521-phuxe2n-tuxedch-chi-phuxed-phux1ea7n-cux1ee9ng}

Chi phí Bill of Materials (BOM), tức bảng tổng hợp toàn bộ linh kiện cấu
thành thiết bị theo dõi IoT dưới đây, dùng
giá tham chiếu thị trường Việt Nam đã được rà soát lại vào ngày
12/04/2026. Với đơn hàng lớn, mức giá này vẫn có thể giảm thêm; tuy
nhiên việc giữ giá mua lẻ giúp phần đánh giá kinh tế bám sát bối cảnh
triển khai thử nghiệm của đồ án.


{\thesissettableid{5.4}{tab-ch5-chi-phi-bom-thiet-bi-tracker}
\def\LTcaptype{table}
\begin{longtable}{|L{0.479\linewidth}|L{0.282\linewidth}|L{0.204\linewidth}|}
\caption{Chi phí BOM thiết bị theo dõi}\label{tab:ch5-chi-phi-bom-thiet-bi-tracker}\\
\hline \textbf{Linh kiện} & \textbf{Vai trò} & \textbf{Giá (VND)} \\
\\\hline
\endfirsthead
\caption[]{Chi phí BOM thiết bị theo dõi (tiếp theo)}\\
\hline \textbf{Linh kiện} & \textbf{Vai trò} & \textbf{Giá (VND)} \\
\\\hline
\endhead
\\\hline
\endfoot
\endlastfoot
ESP32-S3-WROOM-1 & Vi điều khiển chính & 95.000 \\\hline
SIMCom SIM7600CE-T & Modem LTE + GNSS tích hợp & 650.000 \\\hline
vgate iCar Pro BLE & Adapter OBD2 BLE & 367.000 \\\hline
IC cảm biến LIS3DSH & Cảm biến gia tốc (IMU) & 40.000 \\\hline
Pin 18650 1S (1 cell, 3500mAh) & Pin dự phòng & 79.000 \\\hline
Khối nguồn (TP5100, boost/buck, LDO, BMS, LVD) & Quản lý năng lượng &
93.000 \\\hline
PCB, vỏ hộp, jack OBD2, anten, dây cáp, linh kiện phụ & Cơ khí và kết nối &
190.000 \\\hline
\textbf{Tổng cộng} & & \textbf{1.514.000} \\\hline
\end{longtable}
}


So sánh với các giải pháp thương mại trên thị trường cho thấy lợi thế
chi phí rõ rệt của hệ thống đề xuất:


{\thesissettableid{5.5}{tab-ch5-so-sanh-chi-phi-voi-giai-phap-thuong-mai}
\def\LTcaptype{table}
\begin{longtable}{|L{0.215\linewidth}|L{0.213\linewidth}|L{0.205\linewidth}|L{0.183\linewidth}|L{0.129\linewidth}|}
\caption{So sánh chi phí với giải pháp thương mại}\label{tab:ch5-so-sanh-chi-phi-voi-giai-phap-thuong-mai}\\
\hline \textbf{Giải pháp} & \textbf{Chi phí thiết bị} & \textbf{Phí dịch vụ hàng tháng} & \textbf{Tổng chi phí năm
đầu} & \textbf{Khả năng tùy biến} \\
\\\hline
\endfirsthead
\caption[]{So sánh chi phí với giải pháp thương mại (tiếp theo)}\\
\hline \textbf{Giải pháp} & \textbf{Chi phí thiết bị} & \textbf{Phí dịch vụ hàng tháng} & \textbf{Tổng chi phí năm
đầu} & \textbf{Khả năng tùy biến} \\
\\\hline
\endhead
\\\hline
\endfoot
\endlastfoot
Hệ thống đề xuất & 1.514.000 VND & ~70.000 VND
(SIM 4G) & 2.354.000 VND & Cao (mã nguồn mở) \\\hline
GPS Tracker đơn giản (Việt Nam) & 500.000--1.500.000 VND &
50.000--100.000 VND & 1.100.000--2.700.000 VND & Thấp \\\hline
Fleet Management thương mại (quốc tế) & 5.000.000--12.500.000 VND
(~200--500 USD) & 500.000--1.250.000 VND
(~20--50 USD/tháng) & 11.000.000--27.500.000 VND & Thấp
(phụ thuộc nhà cung cấp) \\\hline
iTracking / Vietmap Tracking & 2.000.000--4.000.000 VND &
100.000--300.000 VND & 3.200.000--7.600.000 VND & Trung bình \\\hline
\end{longtable}
}


Như vậy, chi phí tổng thể của hệ thống đề xuất vẫn thấp hơn rõ rệt so
với các nền tảng fleet management thương mại quốc tế, đồng thời vẫn nằm
trong vùng chi phí chấp nhận được so với nhóm thiết bị định vị GPS phổ thông
trong nước dù cung cấp thêm OBD2, cảnh báo thông minh và cơ chế quản lý
năng lượng nhiều chế độ.

\subsubsection{5.2.2. Phân tích chi phí hạ tầng đám
mây}\label{522-phuxe2n-tuxedch-chi-phuxed-hux1ea1-tux1ea7ng-ux111uxe1m-muxe2y}

Chi phí vận hành hạ tầng đám mây phụ thuộc vào quy mô đội xe và mức độ
sử dụng. Đối với doanh nghiệp nhỏ (10--50 xe), toàn bộ hệ thống có thể
chạy trên một máy chủ ảo (VPS) duy nhất.


{\thesissettableid{5.6}{tab-ch5-chi-phi-ha-tang-am-may-theo-quy-mo}
\def\LTcaptype{table}
\begin{longtable}{|L{0.098\linewidth}|L{0.163\linewidth}|L{0.206\linewidth}|L{0.239\linewidth}|L{0.239\linewidth}|}
\caption{Chi phí hạ tầng đám mây theo quy mô}\label{tab:ch5-chi-phi-ha-tang-am-may-theo-quy-mo}\\
\hline \textbf{Quy mô đội xe} & \textbf{Cấu hình VPS} & \textbf{Chi phí VPS/tháng} & \textbf{Chi phí SIM 4G/tháng
(tổng)} & \textbf{Tổng chi phí vận hành/tháng} \\
\\\hline
\endfirsthead
\caption[]{Chi phí hạ tầng đám mây theo quy mô (tiếp theo)}\\
\hline \textbf{Quy mô đội xe} & \textbf{Cấu hình VPS} & \textbf{Chi phí VPS/tháng} & \textbf{Chi phí SIM 4G/tháng
(tổng)} & \textbf{Tổng chi phí vận hành/tháng} \\
\\\hline
\endhead
\\\hline
\endfoot
\endlastfoot
10--30 xe & 2 vCPU, 4GB RAM, 80GB SSD & ~250.000 VND
(~\$10 USD) & 700.000--2.100.000 VND & 950.000--2.350.000
VND \\\hline
30--100 xe & 4 vCPU, 8GB RAM, 160GB SSD & ~500.000 VND
(~\$20 USD) & 2.100.000--7.000.000 VND &
2.600.000--7.500.000 VND \\\hline
100--500 xe & Cụm điều phối dịch vụ (K8s, 3 node) & ~2.500.000 VND
(~\$100 USD) & 7.000.000--35.000.000 VND &
9.500.000--37.500.000 VND \\\hline
\end{longtable}
}


\subsubsection{5.2.3. Lợi thế từ công nghệ mã nguồn
mở}\label{523-lux1ee3i-thux1ebf-tux1eeb-cuxf4ng-nghux1ec7-muxe3-nguux1ed3n-mux1edf}

Toàn bộ ngăn xếp công nghệ của hệ thống đều là mã nguồn mở, nên nhóm
không phải trả chi phí bản quyền cho các thành phần cốt lõi. Quan trọng
hơn, mỗi thành phần đều có một vai trò rõ ràng trong vận hành:

\begin{itemize}
\tightlist
\item
  \textbf{Dịch vụ xử lý yêu cầu phía máy chủ (Express.js +
  TypeScript)}: giúp viết API nhanh, rõ cấu trúc và dễ kiểm soát lỗi
\item
  \textbf{Giao diện web (Next.js 15 + React 19)}: phù hợp cho màn hình quản trị
  có nhiều màn hình và dữ liệu cập nhật liên tục
\item
  \textbf{PostgreSQL 16}: lưu dữ liệu nghiệp vụ có quan hệ rõ ràng như
  người dùng, xe và hành trình
\item
  \textbf{EMQX (phiên bản mã nguồn mở)}: làm broker MQTT giữa thiết bị và
  máy chủ
\item
  \textbf{VictoriaMetrics}: lưu dữ liệu thay đổi liên tục theo thời gian
  với tốc độ ghi cao
\item
  \textbf{ESP-IDF + FreeRTOS}: cung cấp nền tảng để phát triển firmware
  ổn định trên ESP32-S3
\end{itemize}

So với việc sử dụng các giải pháp thương mại (AWS IoT Core, Azure IoT
Hub, hoặc các nền tảng fleet management SaaS), chi phí bản quyền và dịch
vụ có thể lên tới 500.000--5.000.000 VND/tháng tùy quy mô, trong khi hệ
thống đề xuất chỉ cần chi phí VPS cơ bản.

\subsubsection{5.2.4. Đánh giá tác động môi
trường}\label{524-ux111uxe1nh-giuxe1-tuxe1c-ux111ux1ed9ng-muxf4i-trux1b0ux1eddng}

Ở góc độ môi trường, hệ thống có một số điểm tích cực đáng ghi nhận:

\begin{itemize}
\tightlist
\item
  \textbf{Tiêu thụ điện thấp}: Ở chế độ đỗ xe, thiết bị hiện tại duy trì
  mức tiêu thụ deep sleep xấp xỉ 500 μA; khi tiếp tục tối ưu PCB và phân
  bố nguồn, dòng ngủ sâu còn có thể giảm thêm. Mức tiêu thụ hiện tại vẫn
  giúp giảm đáng kể điện năng lấy từ ắc quy xe và kéo dài tuổi thọ ắc
  quy.
\item
  \textbf{Pin dự phòng giảm phụ thuộc ắc quy xe}: Khi xe đậu lâu ngày,
  thiết bị chuyển sang sử dụng pin dự phòng 18650 1S, tránh rút điện từ
  ắc quy xe, bảo vệ ắc quy và giảm lượng khí thải từ việc sạc ắc quy.
\item
  \textbf{Giảm số lần đi kiểm tra xe trực tiếp}: Nhờ khả năng giám sát
  từ xa, người quản lý không cần lái xe đến vị trí xe để kiểm tra, giảm
  lượng khí thải CO2 từ các chuyến đi không cần thiết.
\item
  \textbf{Tối ưu hành trình}: Dữ liệu GPS và OBD2 giúp phân tích và tối
  ưu hành trình, giảm quãng đường đi không cần thiết, giảm tiêu thụ
  nhiên liệu và khí thải.
\end{itemize}

Các điểm tích cực về chi phí và môi trường cho thấy giải pháp có tính
khả thi, nhưng chúng chưa đủ để kết luận hệ thống sẵn sàng mở rộng.
Muốn đi tiếp, cần nhìn thẳng vào các điểm dễ gãy trong vận hành, bảo mật
và độ bền phần cứng.

\begin{center}\rule{0.5\linewidth}{0.5pt}\end{center}

\subsection{5.3. Đánh giá rủi ro và biện pháp giảm
thiểu}\label{53-ux111uxe1nh-giuxe1-rux1ee7i-ro-vuxe0-biux1ec7n-phuxe1p-giux1ea3m-thiux1ec3u}

\subsubsection{5.3.1. Ma trận rủi
ro}\label{531-ma-trux1eadn-rux1ee7i-ro}

Phần này đánh giá rủi ro của hệ thống bằng ma trận xác suất - tác động
(Probability-Impact Matrix) gồm 5 mức độ. Mục tiêu không chỉ là liệt kê
nguy cơ, mà là chỉ ra những điểm nào đã được kiểm soát tương đối tốt và
những điểm nào vẫn cần ưu tiên xử lý nếu đưa hệ thống ra vận hành rộng
hơn. Các rủi ro được xác định từ đặc điểm kỹ thuật của thiết bị hiện tại
và các tình huống dễ phát sinh trong quá trình triển khai thực tế.


{\thesissettableid{5.7}{tab-ch5-thang-o-xac-suat-va-tac-ong}
\def\LTcaptype{table}
\begin{longtable}{|L{0.224\linewidth}|L{0.160\linewidth}|L{0.580\linewidth}|}
\caption{Thang đo xác suất và tác động}\label{tab:ch5-thang-o-xac-suat-va-tac-ong}\\
\hline \textbf{Mức độ} & \textbf{Xác suất} & \textbf{Tác động} \\
\\\hline
\endfirsthead
\caption[]{Thang đo xác suất và tác động (tiếp theo)}\\
\hline \textbf{Mức độ} & \textbf{Xác suất} & \textbf{Tác động} \\
\\\hline
\endhead
\\\hline
\endfoot
\endlastfoot
1 - Rất thấp & \textless{} 5\% & Ảnh hưởng không đáng kể, hệ thống vẫn
hoạt động bình thường \\\hline
2 - Thấp & 5--15\% & Ảnh hưởng nhỏ, có thể khắc phục nhanh \\\hline
3 - Trung bình & 15--30\% & Ảnh hưởng vừa phải, cần xử lý trong thời
gian ngắn \\\hline
4 - Cao & 30--50\% & Ảnh hưởng lớn, có thể làm gián đoạn dịch vụ \\\hline
5 - Rất cao & \textgreater{} 50\% & Ảnh hưởng nghiêm trọng, có thể làm
tê liệt hệ thống \\\hline
\end{longtable}
}



{\thesissettableid{5.8}{tab-ch5-ma-tran-anh-gia-rui-ro-va-bien-phap-giam-thieu}
\def\LTcaptype{table}
\begin{longtable}{|L{0.055\linewidth}|L{0.180\linewidth}|L{0.090\linewidth}|L{0.090\linewidth}|L{0.120\linewidth}|L{0.300\linewidth}|L{0.090\linewidth}|}
\caption{Ma trận đánh giá rủi ro và biện pháp giảm thiểu}\label{tab:ch5-ma-tran-anh-gia-rui-ro-va-bien-phap-giam-thieu}\\
\hline \textbf{ID} & \textbf{Rủi ro} & \textbf{Xác suất} & \textbf{Tác động} & \textbf{Mức độ rủi ro} & \textbf{Biện pháp giảm thiểu} & \textbf{Trạng thái} \\
\\\hline
\endfirsthead
\caption[]{Ma trận đánh giá rủi ro và biện pháp giảm thiểu (tiếp theo)}\\
\hline \textbf{ID} & \textbf{Rủi ro} & \textbf{Xác suất} & \textbf{Tác động} & \textbf{Mức độ rủi ro} & \textbf{Biện pháp giảm thiểu} & \textbf{Trạng thái} \\
\\\hline
\endhead
\\\hline
\endfoot
\endlastfoot
R1 & Mất sóng 4G/LTE tại khu vực nông thôn, vùng sâu & 4 - Cao & 3 -
Trung bình & \textbf{Cao} & Tự kết nối lại, duy trì tín hiệu giữ kết nối
MQTT, lưu tạm vào microSD và phát lại theo đúng thứ tự khi mạng phục hồi & Đã triển khai \\\hline
R2 & Không tương thích BLE OBD2 với một số dòng xe & 2 - Thấp & 3 -
Trung bình & \textbf{Trung bình} & Sử dụng vgate iCar Pro (tương thích
rộng), hỗ trợ nhiều giao thức OBD2 (ISO 15765, ISO 14230, J1850) & Đã
giảm thiểu \\\hline
R3 & Quá tải broker MQTT khi số lượng thiết bị tăng cao & 2 - Thấp & 4 -
Cao & \textbf{Trung bình} & EMQX hỗ trợ chạy nhiều nút; khi tải tăng có
thể bổ sung thêm nút broker và mở rộng riêng MQTT Bridge để san tải phần tiếp
nhận và xử lý phía sau & Có phương án \\\hline
R4 & Tấn công bảo mật (giả mạo thiết bị, chiếm quyền truy cập) & 3 -
Trung bình & 5 - Rất cao & \textbf{Cao} & Cơ chế giữ phiên trên máy chủ với
SHA-256, ACL MQTT theo từng thiết bị, HTTPS hoặc TLS, kiểm tra dữ liệu đầu vào bằng Zod và giới hạn tần suất gọi để tránh bị gửi dồn dập & Đã triển khai cơ bản \\\hline
R5 & Mất dữ liệu trong thời gian mất kết nối mạng kéo dài & 2 - Thấp & 4
- Cao & \textbf{Trung bình} & Thử gửi lại bản tin, phát lại từ bộ nhớ cục bộ theo mức ưu tiên QoS,
đặt thời gian chờ xác nhận rồi tăng dần khoảng nghỉ giữa các lần thử, dọn bộ nhớ theo hạn mức & Đã giảm thiểu \\\hline
R6 & Hư hỏng phần cứng do nhiệt độ cực đoan (xe đỗ ngoài nắng) & 2 -
Thấp & 4 - Cao & \textbf{Trung bình} & ESP32-S3 hoạt động -40 đến 85°C,
thiết kế tản nhiệt, đặt thiết bị trong vị trí mát, cảnh báo nhiệt độ &
Thiết kế có tính đến \\\hline
R7 & Cạn ắc quy xe do thiết bị hoạt động liên tục & 2 - Thấp & 5 - Rất
cao & \textbf{Cao} & Mạch switch profile chuyển sang pin dự phòng: 12V:
\(OFF = 12.0\,\mathrm{V}\), \(ON = 12.2\,\mathrm{V}\); 24V:
\(OFF = 24.0\,\mathrm{V}\), \(ON = 24.4\,\mathrm{V}\); deep sleep xấp xỉ 500
μA trên thiết bị hiện tại, pin dự phòng 18650 1S & Đã triển khai \\\hline
R8 & Lỗi firmware gây treo hệ thống & 3 - Trung bình & 4
- Cao & \textbf{Cao} & OTA từ xa với xác nhận sau khi khởi động lại và
quay về bản cũ thủ công đã có; cơ chế phát hiện treo tự động
(watchdog) và giám sát treo vẫn cần bổ sung &
Triển khai một phần \\\hline
R9 & Lỗi thẻ nhớ SD (không có thẻ, gắn lỗi, đầy dung lượng) & 3 - Trung bình &
3 - Trung bình & \textbf{Trung bình} & Thử gắn lại với thời gian chờ tăng dần, buộc ghi ra đĩa rồi tháo thẻ đúng cách, kiểm tra có thẻ từ lúc khởi động và dọn dung lượng theo hạn mức
& Đã triển khai cơ bản \\\hline
\end{longtable}
}


\subsubsection{5.3.2. Phân tích chi tiết các rủi ro
chính}\label{532-phuxe2n-tuxedch-chi-tiux1ebft-cuxe1c-rux1ee7i-ro-chuxednh}

\textbf{Rủi ro R1 - Mất sóng 4G tại khu vực nông thôn:}

Đây là rủi ro có xác suất cao nhất trong thực tế vận hành tại Việt Nam,
đặc biệt khi xe di chuyển qua các tuyến đường liên tỉnh hoặc vùng núi.
Phiên bản hiện tại kết hợp tự kết nối lại với vùng lưu tạm cục bộ trên
microSD: khi mất kết nối 4G, dữ liệu hành trình và sự kiện vẫn được cất
tạm trong thiết bị; khi mạng phục hồi, firmware phát lại theo đúng thứ
tự và đánh dấu những bản ghi quan trọng đã được phía nhận xác nhận. Cách
tiếp cận này giúp giảm đáng kể nguy cơ đứt quãng dữ liệu hành trình
trong các khoảng mất sóng kéo dài.

\textbf{Rủi ro R4 - Tấn công bảo mật:}

Bảo mật là rủi ro có tác động nghiêm trọng nhất. Hệ thống đã triển khai
nhiều lớp bảo vệ: cơ chế giữ đăng nhập trên máy chủ để có thể khóa phiên
khi cần; token được lưu dưới dạng mã băm SHA-256 thay vì lưu nguyên văn;
quy tắc ACL trong MQTT để mỗi thiết bị chỉ được dùng đúng kênh dữ liệu
của mình; bước kiểm tra hình dạng dữ liệu đầu vào bằng Zod cho toàn bộ
API; và cấu hình CORS chỉ cho phép các nguồn truy cập hợp lệ gọi sang hệ
thống. Tuy nhiên, một số biện pháp bảo mật nâng cao như chứng chỉ riêng
cho thiết bị, mã hóa TLS cho MQTT và ghi nhận sự kiện an ninh vẫn chưa
được triển khai đầy đủ, nên cần được bổ sung trong giai đoạn vận hành thật.

\textbf{Rủi ro R7 - Cạn ắc quy xe:}

Đây là rủi ro có tác động nghiêm trọng nhất đối với trải nghiệm người
dùng cuối, vì xe không khởi động được sẽ gây bất tiện lớn cho khách
thuê. Hệ thống đã có nhiều cơ chế bảo vệ: ngưỡng LVD profile 12V tại
\(OFF = 12.0\,\mathrm{V}\)/\(ON = 12.2\,\mathrm{V}\), profile 24V tại
\(OFF = 24.0\,\mathrm{V}\)/\(ON = 24.4\,\mathrm{V}\); cơ chế chuyển
nguồn có hysteresis; chế độ deep sleep của thiết bị hiện tại ở mức xấp
xỉ 500 μA; và pin dự phòng 18650 1S cho phép hoạt động độc lập khi mạch
LVD ngắt nguồn từ ắc quy. Với bộ pin tham chiếu 18650, thời lượng dự phòng
vẫn đủ cho cảnh báo gián đoạn ngắn hạn, nhưng các giá trị quy đổi cho vận hành liên
tục cần được hiểu là quy đổi thiết kế và nên được đo lại nếu triển khai
phần cứng theo đúng cell mới. Tuy nhiên, ngưỡng LVD profile 12V vẫn cần
được hiệu chuẩn thêm để bám sát đúng giá trị thiết kế.

\subsubsection{5.3.3. Tổng hợp mức độ rủi
ro}\label{533-tux1ed5ng-hux1ee3p-mux1ee9c-ux111ux1ed9-rux1ee7i-ro}


{\thesissettableid{5.9}{tab-ch5-tong-hop-muc-o-rui-ro-theo-phan-loai}
\def\LTcaptype{table}
\begin{longtable}{|L{0.249\linewidth}|L{0.189\linewidth}|L{0.125\linewidth}|L{0.229\linewidth}|L{0.153\linewidth}|}
\caption{Tổng hợp mức độ rủi ro theo phân loại}\label{tab:ch5-tong-hop-muc-o-rui-ro-theo-phan-loai}\\
\hline \textbf{Phân loại} & \textbf{Số rủi ro} & \textbf{Mức cao} & \textbf{Mức trung bình} & \textbf{Mức thấp} \\
\\\hline
\endfirsthead
\caption[]{Tổng hợp mức độ rủi ro theo phân loại (tiếp theo)}\\
\hline \textbf{Phân loại} & \textbf{Số rủi ro} & \textbf{Mức cao} & \textbf{Mức trung bình} & \textbf{Mức thấp} \\
\\\hline
\endhead
\\\hline
\endfoot
\endlastfoot
Kết nối và truyền thông & 3 (R1, R3, R5) & 1 & 2 & 0 \\\hline
Bảo mật & 1 (R4) & 1 & 0 & 0 \\\hline
Phần cứng & 4 (R2, R6, R7, R9) & 1 & 3 & 0 \\\hline
Firmware & 1 (R8) & 1 & 0 & 0 \\\hline
\textbf{Tổng cộng} & \textbf{9} & \textbf{4} & \textbf{5} &
\textbf{0} \\\hline
\end{longtable}
}


Kết quả cho thấy 4/9 rủi ro ở mức cao; phần lớn đã có biện pháp giảm
thiểu triển khai trong phiên bản hiện tại, bao gồm lớp lưu đệm cục bộ
microSD cho kịch bản mất mạng. Các rủi ro cần ưu tiên xử lý tiếp theo là
R4 (bảo mật nâng cao), R8 (bổ sung cơ chế phát hiện treo tự động, tăng cường giám sát treo và
kiểm thử OTA dài ngày) và R9 (mở rộng kiểm thử độ bền SD card theo chu
kỳ vận hành thực tế).

Từ bức tranh rủi ro này có thể thấy hướng phát triển tiếp theo không nên
chỉ là bổ sung tính năng mới, mà phải sắp thứ tự ưu tiên theo các điểm
yếu có khả năng làm hệ thống mất ổn định. Vì vậy, mục 5.4 chuyển sang
các khuyến nghị phát triển.

\begin{center}\rule{0.5\linewidth}{0.5pt}\end{center}

\subsection{5.4. Khuyến nghị cho tương
lai}\label{54-khuyux1ebfn-nghux1ecb-cho-tux1b0ux1a1ng-lai}

Dựa trên kết quả đánh giá ở các phần trên, phần này chuyển từ câu hỏi
``hệ thống đang ở đâu'' sang câu hỏi ``nên đi tiếp theo hướng nào''.
Các khuyến nghị được sắp xếp theo từng giai đoạn để người đọc có thể
thấy rõ đâu là bước phát triển gần, đâu là hướng mở rộng dài hạn.

\subsubsection{5.4.1. Giai đoạn 2: Mở rộng tính năng (giai đoạn
2)}\label{541-giai-ux111oux1ea1n-2-mux1edf-rux1ed9ng-tuxednh-nux103ng-phase-2}

Hệ thống đã được thiết kế với kiến trúc mở rộng, cơ sở dữ liệu
PostgreSQL đã bao gồm các bảng dữ liệu cho các tính năng của giai đoạn 2. Điều
này cho thấy đồ án không dừng ở một nguyên mẫu trình diễn, mà đã chuẩn
bị sẵn một phần nền móng để mở rộng thành hệ thống nghiệp vụ đầy đủ hơn.
Các tính năng khuyến nghị triển khai trong giai đoạn tiếp theo bao gồm:

\begin{itemize}
\tightlist
\item
\textbf{Quản lý đặt xe}: Hệ thống đặt xe trực tuyến cho
  khách hàng, tích hợp lịch và trạng thái sẵn sàng của xe. Schema
  cơ sở dữ liệu đã sẵn sàng với bảng \texttt{bookings} và
  \texttt{reservations}.
\item
  \textbf{Thanh toán trực tuyến}: Tích hợp cổng thanh toán
  (VNPay, Momo, ZaloPay) để xử lý thanh toán tiền thuê xe tự động. Bảng
  \texttt{payments} và \texttt{invoices} đã được thiết kế trong cấu trúc cơ sở dữ liệu.
\item
  \textbf{Báo cáo hư hỏng}: Cho phép tài xế và nhân
  viên báo cáo tình trạng xe khi nhận và trả, kèm theo hình ảnh. Bảng
  \texttt{damage\_reports} đã có trong cơ sở dữ liệu.
\item
  \textbf{Đánh giá và xếp hạng}: Cho phép khách hàng đánh giá
  trải nghiệm thuê xe, giúp doanh nghiệp cải thiện chất lượng dịch vụ.
  Bảng \texttt{reviews} và \texttt{ratings} đã được thiết kế.
\end{itemize}

\subsubsection{5.4.2. Ứng dụng trí tuệ nhân tạo và học máy}\label{542-ux1ee9ng-dux1ee5ng-truxed-tuux1ec7-nhuxe2n-tux1ea1o-vuxe0-hux1ecdc-muxe1y-aiml}

Khi hệ thống đã có khả năng thu thập dữ liệu vận hành liên tục và khá ổn định,
giá trị tiếp theo không chỉ nằm ở việc ``nhìn thấy dữ liệu'', mà ở việc
biến dữ liệu đó thành dự báo và khuyến nghị. Vì vậy, trí tuệ nhân tạo và học máy là một hướng
mở rộng hợp lý để tăng chiều sâu cho hệ thống:

\textbf{Phân tích hành vi lái xe:}

Sử dụng dữ liệu gia tốc kế (IMU), tốc độ, RPM và GPS để phân tích và
chấm điểm hành vi lái xe. Các chỉ số có thể tính toán bao gồm: tần suất
phanh gấp, gia tốc đột ngột, tốc độ vào cua, vượt quá tốc độ. Kết quả
phân tích giúp doanh nghiệp đánh giá mức độ an toàn của tài xế và điều
chỉnh chính sách cho thuê (phí bảo hiểm, mức đặt cọc).

\textbf{Bảo trì dự đoán:}

Phân tích xu hướng dữ liệu OBD2 theo thời gian (nhiệt độ động cơ, áp
suất dầu, điện áp ắc quy, mã lỗi DTC) để dự đoán các vấn đề kỹ thuật
trước khi xảy ra hư hỏng. Mô hình học máy có thể được huấn luyện từ dữ
liệu lịch sử để cảnh báo sớm về nhu cầu bảo trì, giảm thời gian xe không
hoạt động, tức thời gian xe phải nằm chờ sửa và không tạo ra
giá trị sử dụng) và chi phí sửa chữa khẩn cấp.

\textbf{Phát hiện bất thường:}

Sử dụng các thuật toán học máy không giám sát, tức cho hệ thống tự học mẫu dữ liệu bình thường rồi phát hiện điểm lệch, để
phát hiện các hành vi bất thường: xe đi vào khu vực không bình thường,
thời gian sử dụng bất thường, thay đổi đột ngột trong mẫu sử dụng. Điều
này giúp phát hiện sớm các tình huống rủi ro như chiếm đoạt xe hoặc sử
dụng sai mục đích.

\subsubsection{5.4.3. Ứng dụng di động (Mobile
Application)}\label{543-ux1ee9ng-dux1ee5ng-di-ux111ux1ed9ng-mobile-application}

Hiện tại hệ thống chỉ có giao diện web (Tracking\_Frontend). Tuy nhiên,
trong bối cảnh quản lý phương tiện, nhu cầu nhận cảnh báo và theo dõi từ
thiết bị di động là rất tự nhiên. Vì vậy, việc phát triển ứng dụng di
động là bước mở rộng hợp lý tiếp theo:

\begin{itemize}
\tightlist
\item
  \textbf{Phương án đề xuất}: Flutter theo hướng ứng dụng lai
  (Tracking\_Mobile/). Cách làm này có thể hiểu là tạo một vỏ ứng dụng
  di động mỏng bằng Flutter, rồi nhúng giao diện web hiện có vào bên
  trong bằng WebView, tức một khung hiển thị nội dung web ngay trong ứng
  dụng. Phương án này giúp rút ngắn thời gian phát triển
  nhờ tái sử dụng giao diện web, nhưng vẫn cho phép truy cập các tính
  năng riêng của điện thoại như thông báo đẩy, GPS hay camera.
\item
\textbf{Tính năng chính}: Theo dõi vị trí xe trên bản đồ, nhận thông báo đẩy
khi có cảnh báo, xem lịch sử hành trình, quản lý đặt xe
  (giai đoạn 2).
\item
  \textbf{Nền tảng}: Hỗ trợ đồng thời iOS và Android từ một bộ mã nguồn
  Flutter duy nhất.
\end{itemize}

\subsubsection{5.4.4. Tính toán biên (Edge
Computing)}\label{544-tuxednh-touxe1n-biuxean-edge-computing}

Với khả năng xử lý của ESP32-S3 (chip hai lõi Xtensa LX7, 240 MHz, 8MB
PSRAM), một hướng mở rộng đáng chú ý là chuyển bớt phần phân tích xuống
ngay trên thiết bị. Cách tiếp cận này giúp hệ thống phản ứng sớm hơn và
giảm phụ thuộc vào băng thông mạng:

\begin{itemize}
\tightlist
\item
  \textbf{Phát hiện bất thường tại thiết bị}: Thay vì gửi toàn bộ dữ
  liệu thô lên máy chủ để phân tích, thiết bị có thể chạy mô hình học máy nhẹ
  (TensorFlow Lite Micro, tức phiên bản rút gọn để chạy trên vi điều
  khiển) để phát hiện bất thường ngay tại chỗ. Chỉ gửi
  cảnh báo khi phát hiện bất thường, giảm lượng dữ liệu truyền và tiết
  kiệm băng thông 4G.
\item
  \textbf{Phân loại hành vi lái xe tại thiết bị}: Sử dụng dữ liệu IMU và
  OBD2 để phân loại hành vi lái xe (bình thường, hung hãn, mệt mỏi) ngay
  trên ESP32-S3, gửi kết quả phân loại thay vì dữ liệu thô.
\item
  \textbf{Nén dữ liệu thông minh}: Sử dụng thuật toán nén dữ liệu trên
  thiết bị, chỉ gửi dữ liệu khi có thay đổi đáng kể (dead reckoning, tức
  chỉ cập nhật khi quỹ đạo hoặc trạng thái lệch đủ nhiều so với lần gần
  nhất),
  giảm 50--70\% lượng dữ liệu truyền.
\end{itemize}

\subsubsection{5.4.5. Triển khai vận hành thực tế}\label{545-triux1ec3n-khai-sux1ea3n-xuux1ea5t-production-deployment}

Để chuyển từ môi trường phát triển sang sản xuất, hệ thống không chỉ cần
``chạy được ổn định hơn'' mà còn phải đạt mức sẵn sàng cao hơn về bảo
mật, giám sát và tự động hóa triển khai. Vì vậy, các bước chuyển đổi
trọng tâm gồm:

\begin{itemize}
\tightlist
\item
  \textbf{Điều phối nhiều dịch vụ tự động}: Khi số thành phần tăng,
  Docker Compose sẽ khó đáp ứng tốt các nhu cầu tự khởi động lại, tăng
  giảm tải và cập nhật không gián đoạn. Vì vậy, giai đoạn vận hành thật nên
  chuyển sang Kubernetes (K8s), tức công cụ chuyên quản lý nhiều dịch vụ
  trên nhiều máy chủ, để mỗi dịch vụ có thể tăng số bản sao khi tải cao
  và tự phục hồi khi có lỗi.
\item
  \textbf{Tự động kiểm tra và triển khai (CI/CD)}: Thiết lập quy trình
  tự động để mỗi lần cập nhật mã nguồn đều được kiểm tra lỗi cú pháp,
  kiểm thử, đóng gói và đưa lên môi trường thử nghiệm trước. Chỉ sau khi
  đạt yêu cầu mới triển khai lên môi trường vận hành thật. Cách làm này giảm mạnh rủi
  ro lỗi phát sinh do thao tác thủ công.
\item
  \textbf{Mã hóa toàn bộ đường truyền}: Toàn bộ luồng trao đổi dữ liệu
  cần được bảo vệ bằng HTTPS cho giao diện web và API, TLS cho MQTT và
  kết nối mã hóa tới cơ sở dữ liệu. Nginx Proxy Manager có thể đảm nhiệm
  vai trò cổng vào chung và quản lý chứng chỉ số để việc triển khai nhất
  quán hơn.
\item
  \textbf{Giám sát và cảnh báo tập trung}: Cần hoàn thiện khả năng quan
  sát hệ thống, tức phải nhìn ra dịch vụ nào đang khỏe, chậm hay lỗi ở
  đâu. Điều này bao gồm việc xuất số liệu vận hành theo chuẩn mà công cụ
  giám sát có thể đọc được (Prometheus-compatible metrics, tức định dạng
  số liệu mà nhiều công cụ giám sát phổ biến có thể thu thập), lưu log tập
  trung, hiển thị màn hình theo dõi và gửi cảnh báo qua Telegram hoặc
  email khi có sự cố.
\end{itemize}

\subsubsection{5.4.6. Phần cứng phiên bản 2 (Hardware
v2)}\label{546-phux1ea7n-cux1ee9ng-phiuxean-bux1ea3n-2-hardware-v2}

Phiên bản phần cứng hiện có sử dụng bo mạch 2 lớp do nhóm thiết kế và
được lắp ráp thủ công, đủ để triển khai và đánh giá ở quy mô đồ án. Tuy
nhiên, nếu nhìn dưới góc độ sản phẩm, đây mới là bước khởi đầu. Để tiến
tới sản xuất hàng loạt, các cải tiến phần cứng tiếp theo nên tập trung
vào tối ưu DFM, tức thiết kế sao cho dễ sản xuất, và EMI là khả năng giảm
nhiễu điện từ, cùng với độ gọn cơ khí và độ tin cậy:

\begin{itemize}
\tightlist
\item
  \textbf{Thiết kế PCB tùy chỉnh}: Thiết kế PCB 4 lớp, tức bo mạch có
  nhiều lớp dẫn hơn để đi dây gọn và ổn định hơn, tích hợp
  toàn bộ thành phần (ESP32-S3, mạch nạp, mạch nguồn, đầu nối SIM, đầu
  nối antenna) trên một board duy nhất. Giảm kích thước xuống khoảng
  60x40 mm, phù hợp để lắp đặt trong xe.
\item
  \textbf{Antenna tích hợp}: Sử dụng antenna ceramic cho GPS/GNSS và
  antenna PCB cho 4G/LTE, giảm số dây cáp và tăng độ tin cậy.
\item
  \textbf{Vỏ hộp công nghiệp}: Thiết kế vỏ hộp nhựa ABS chống nước
  (IP65), chịu nhiệt, với đầu nối OBD2 tích hợp và đầu nối nguồn 12V.
\item
  \textbf{Chi phí sản xuất hàng loạt}: Khi sản xuất từ 500 bộ trở lên,
  chi phí BOM có thể giảm xuống còn 500.000--800.000 VND/bộ nhờ mua linh
  kiện số lượng lớn và tối ưu hóa thiết kế PCB.
\end{itemize}

\subsubsection{5.4.7. Tóm tắt lộ trình phát
triển}\label{547-tuxf3m-tux1eaft-lux1ed9-truxecnh-phuxe1t-triux1ec3n}


{\thesissettableid{5.10}{tab-ch5-lo-trinh-phat-trien-khuyen-nghi}
\def\LTcaptype{table}
\begin{longtable}{|L{0.251\linewidth}|L{0.115\linewidth}|L{0.460\linewidth}|L{0.129\linewidth}|}
\caption{Lộ trình phát triển khuyến nghị}\label{tab:ch5-lo-trinh-phat-trien-khuyen-nghi}\\
\hline \textbf{Giai đoạn} & \textbf{Thời gian} & \textbf{Nội dung chính} & \textbf{Ưu tiên} \\
\\\hline
\endfirsthead
\caption[]{Lộ trình phát triển khuyến nghị (tiếp theo)}\\
\hline \textbf{Giai đoạn} & \textbf{Thời gian} & \textbf{Nội dung chính} & \textbf{Ưu tiên} \\
\\\hline
\endhead
\\\hline
\endfoot
\endlastfoot
Giai đoạn 1.5 - Bảo mật & 1--2 tháng & TLS cho MQTT, chứng chỉ riêng cho
thiết bị, cơ chế phát hiện treo tự động và nhật ký an ninh & Cao \\\hline
Giai đoạn 2 - Mở rộng tính năng & 2--4 tháng & Đặt xe, thanh toán, báo
cáo hư hỏng, đánh giá và ứng dụng di động & Cao \\\hline
Giai đoạn 2.5 - Tối ưu hóa & 1--2 tháng & Lưu đệm Redis, tối ưu truy
vấn, tái sử dụng kết nối và kiểm tra sức khỏe dịch vụ & Trung bình \\\hline
Giai đoạn 3 - Trí tuệ nhân tạo và học máy & 3--6 tháng & Phân tích hành vi lái xe, bảo trì dự
đoán và phát hiện bất thường & Trung bình \\\hline
Giai đoạn 4 - Vận hành thực tế quy mô lớn & 2--3 tháng & Kubernetes,
CI/CD, giám sát tập trung và TLS trên toàn tuyến & Cao \\\hline
Giai đoạn 5 - Phần cứng v2 & 4--6 tháng & PCB tùy chỉnh, vỏ hộp công
nghiệp và chuẩn bị sản xuất hàng loạt & Thấp (tùy nhu cầu) \\\hline
\end{longtable}
}


\begin{center}\rule{0.5\linewidth}{0.5pt}\end{center}

\subsection{Kết luận chương
5}\label{kux1ebft-luux1eadn-chux1b0ux1a1ng-5}

Kết quả đánh giá cho thấy hệ thống đã đáp ứng phần lớn các mục tiêu kỹ
thuật trọng yếu của một thiết bị giám sát phương tiện có khả năng vận
hành thực tế. Ở tầng phần cứng, mức tiêu thụ điện duy trì trong ngưỡng
bối cảnh thiết kế; ở tầng firmware, máy trạng thái vận hành ổn định và luồng OTA,
tức cập nhật phần mềm từ xa, cơ
bản đã được tích hợp; ở tầng máy chủ, độ trễ vẫn đáp ứng yêu cầu theo dõi
thời gian thực. Về chi phí, BOM khoảng 1.514.000 VND cùng chi phí vận
hành thấp tạo lợi thế rõ rệt so với nhiều giải pháp thương mại cùng phân
khúc.

Điểm quan trọng rút ra từ chương này là kiến trúc tổng thể của đề tài đã
đi đúng hướng; phần việc còn lại chủ yếu là làm hệ thống "cứng cáp" hơn
cho vận hành dài hạn, thay vì phải thay đổi lại từ gốc. Vì vậy, trọng
tâm giai đoạn tiếp theo không phải thay đổi kiến trúc tổng thể, mà là
tăng chiều sâu làm cứng vận hành: cơ chế tự phát hiện treo hệ thống
(watchdog), bảo mật MQTT/TLS và kiểm thử vận
hành dài hạn. Trên nền tảng đó, Chương 6 sẽ nhìn lại toàn bộ quá trình
thực hiện dưới góc độ học thuật và bài học kinh nghiệm.

\chapterheading{CHƯƠNG 6. PHẢN HỒI VÀ BÀI HỌC KINH NGHIỆM -- REFLECTION AND LESSONS LEARNED}

\subsection{6.1. Ứng dụng kiến thức kỹ thuật -- Application of prior
coursework}\label{61-ux1ee9ng-dux1ee5ng-kiux1ebfn-thux1ee9c-kux1ef9-thuux1eadt--application-of-prior-coursework}

Dự án "IoT Vehicle Tracking System" là kết quả của quá trình tổng hợp
kiến thức từ nhiều học phần và nhiều lớp công nghệ khác nhau. Từ phần
cứng, firmware đến dịch vụ máy chủ và giao diện web, mỗi quyết định triển khai đều
đòi hỏi người thực hiện chuyển kiến thức nền tảng thành lựa chọn kỹ
thuật cụ thể. Vì vậy, chương này không chỉ mang tính tổng kết, mà còn là
phần nhìn lại xem kiến thức đã học được vận dụng như thế nào khi đi vào
một bài toán thực tế có nhiều ràng buộc đồng thời.

\subsubsection{6.1.1. Vi xử lý và Vi điều
khiển}\label{611-vi-xux1eed-luxfd-vuxe0-vi-ux111iux1ec1u-khiux1ec3n}

Trong toàn bộ đồ án, lớp kiến thức về vi xử lý và vi điều khiển là nơi
mọi quyết định kỹ thuật bắt đầu trở nên cụ thể. Chính ở lớp này, các
khái niệm học thuật như thanh ghi, bus ngoại vi hay đồng bộ tài nguyên
không còn là lý thuyết riêng lẻ, mà trở thành công cụ trực tiếp để tạo
ra một thiết bị có thể vận hành ngoài thực tế. Các nội dung được áp
dụng bao gồm:

\begin{itemize}
\tightlist
\item
  \textbf{Lập trình cho vi điều khiển ESP32-S3}: Nhóm phải vận dụng kiến
  thức về cách một vi điều khiển hai nhân hoạt động ở mức thấp, từ bộ
  nhớ, thanh ghi đến cách chương trình tương tác với phần cứng. ESP32-S3
  dùng kiến trúc Xtensa LX7, nhưng điều quan trọng với đồ án là hiểu rõ
  cách con chip này vận hành để tối ưu hiệu suất và tiêu thụ năng lượng.
\item
  \textbf{FreeRTOS và đồng bộ tài nguyên}: Kiến thức về hệ điều hành
  thời gian thực được áp dụng để tổ chức các công việc nền, ví dụ khối
  BLE hoặc luồng giao tiếp với modem. Các khái niệm như mutex hay
  semaphore được dùng như "khóa" và "tín hiệu chờ" để nhiều phần mềm
  không cùng truy cập một tài nguyên tại một lúc. Dù firmware không chia
  thành nhiều nhánh nghiệp vụ lớn, tư duy điều phối và tránh xung đột
  truy cập vẫn được vận dụng trực tiếp.
\item
  \textbf{Giao tiếp với các phần cứng ngoại vi}: Nhóm phải cấu hình nhiều
  kiểu kết nối phần cứng khác nhau như chân điều khiển bật/tắt (GPIO),
  đo điện áp analog (ADC), cổng nối tiếp với modem (UART), bus cảm biến
  (I2C/SPI) và giao tiếp với thẻ nhớ microSD (SDMMC). Việc hiểu rõ từng
  đường giao tiếp này là điều kiện để modem, cảm biến và khối lưu trữ
  cùng làm việc ổn định trên một thiết bị nhỏ gọn.
\end{itemize}

Khi thiết bị đã đọc được phần cứng và tự điều phối được các ngoại vi,
câu hỏi kế tiếp là dữ liệu rời thiết bị bằng cách nào và được đưa tới
đúng nơi ra sao. Vì vậy, phần sau chuyển sang lớp kiến thức mạng và
IoT.

\subsubsection{6.1.2. Mạng máy tính và
IoT}\label{612-mux1ea1ng-muxe1y-tuxednh-vuxe0-iot}

Nếu phần cứng quyết định thiết bị có thu được dữ liệu hay không, thì
kiến thức mạng quyết định dữ liệu đó có đi đúng nơi, đúng lúc và đủ an
toàn hay không. Điểm quan trọng nhóm rút ra là: trong hệ thống IoT,
``truyền được'' chưa đủ; còn phải biết khi nào ưu tiên tốc độ, khi nào
ưu tiên độ tin cậy và khi nào cần giữ kết nối hai chiều. Vì vậy, nhóm
kiến thức này xuất hiện xuyên suốt từ thiết bị, hệ thống máy chủ đến giao diện thời
gian thá»±c:

\begin{itemize}
\tightlist
\item
  \textbf{Mô hình TCP/IP}: Giúp nhìn đường đi của dữ liệu theo từng lớp,
  từ lúc rời thiết bị qua mạng 4G/LTE cho đến khi tới máy chủ. Nhờ đó,
  nhóm phân biệt được phần nào do firmware chịu trách nhiệm, phần nào do
  hạ tầng mạng và phần nào do dịch vụ xử lý phía máy chủ đảm nhận.
\item
  \textbf{Giao thức MQTT}: MQTT được dùng như kênh trung gian nhận bản
  tin từ thiết bị rồi chuyển cho dịch vụ cần xử lý. Các mức QoS có thể hiểu là
  mức ưu tiên giữa nhanh và chắc: dữ liệu vị trí gửi nhiều nên có thể
  ưu tiên tốc độ, còn cảnh báo và lệnh điều khiển cần được xác nhận rõ
  ràng hơn. Retain Message giúp hệ thống giữ lại trạng thái mới nhất để
  bên nhận vào sau vẫn biết tình hình hiện tại, còn LWT là bản tin tự
  động báo rằng thiết bị vừa mất kết nối.
\item
  \textbf{HTTP/REST API}: Nếu MQTT phù hợp cho bản tin từ thiết bị, thì
  HTTP/REST phù hợp cho phần giao tiếp quản trị giữa giao diện và máy
  chủ. Cách tổ chức API theo tài nguyên, phương thức và mã phản hồi giúp
  người phát triển dễ hiểu dữ liệu nào đang được yêu cầu, dễ kiểm thử và
  dễ mở rộng khi thêm tính năng.
\item
  \textbf{WebSocket và giao tiếp thời gian thực}: Kênh WebSocket cho
  phép máy chủ chủ động gửi dữ liệu mới lên màn hình giám sát mà không
  cần trình duyệt hỏi lặp đi lặp lại bằng các yêu cầu HTTP. Nhờ vậy, vị trí xe và cảnh
  báo được cập nhật mượt hơn và ít tốn tài nguyên hơn.
\item
  \textbf{Bảo mật mạng}: Mã hóa đường truyền (TLS/SSL), xác thực người
  dùng, phân quyền và các quy tắc ACL, tức danh sách quyền cho từng thiết
  bị, để giới hạn mỗi thiết bị chỉ được dùng đúng kênh dữ liệu của mình,
  không phải phần bổ sung tùy chọn, mà
  là điều kiện để dữ liệu vận hành không bị nghe lén hoặc giả mạo khi đi qua Internet.
\end{itemize}

Nếu mạng quyết định dữ liệu đi đúng đường, thì cơ sở dữ liệu quyết định
dữ liệu đó được cất đúng chỗ và truy vấn lại đúng mục đích. Vì vậy,
phần tiếp theo chuyển sang cách tổ chức lưu trữ.

\subsubsection{6.1.3. Cơ sở dữ
liệu}\label{613-cux1a1-sux1edf-dux1eef-liux1ec7u}

Kiến thức cơ sở dữ liệu được vận dụng không chỉ để ``lưu được dữ liệu'',
mà để tổ chức dữ liệu theo đúng đặc tính sử dụng của hệ thống. Bài học
rõ nhất ở đây là: không phải mọi dữ liệu nên được cất cùng một chỗ.
Dữ liệu vận hành theo thời gian cần ghi nhiều và truy vấn theo mốc thời
gian; dữ liệu nghiệp vụ lại
cần quan hệ chặt, ràng buộc rõ và giao dịch an toàn. Vì vậy, đồ án phải
chủ động tách cách lưu trữ thay vì dồn về một cơ sở dữ liệu duy nhất:
\begin{itemize}
\tightlist
\item
  \textbf{Thiết kế cơ sở dữ liệu quan hệ (PostgreSQL)}: PostgreSQL được
  dùng cho các dữ liệu cần liên kết chặt chẽ như phương tiện, người
  dùng, chuyến đi, cảnh báo và vùng giám sát. Các khái niệm như chuẩn
  hóa, khóa chính, khóa ngoại và ràng buộc toàn vẹn giúp dữ liệu không
  bị mâu thuẫn khi hệ thống vận hành lâu dài.
\item
  \textbf{Cơ sở dữ liệu chuỗi thời gian (VictoriaMetrics)}: Đây là nơi
  lưu dữ liệu thay đổi liên tục theo thời gian như vị trí, điện áp,
  trạng thái nguồn hay rung động. Cơ sở dữ liệu chuỗi thời gian phù hợp
  hơn cơ sở dữ liệu quan hệ trong kiểu bài toán này vì ghi nhiều, nén
  tốt và truy vấn theo mốc thời gian nhanh hơn.
\item
  \textbf{Tối ưu hóa truy vấn}: Khi dữ liệu tăng dần, việc ``có bảng''
  là chưa đủ; còn phải truy vấn đủ nhanh để màn hình giám sát không bị chậm.
  Vì vậy, nhóm áp dụng chỉ mục, chỉnh lại câu lệnh và xem cách cơ sở dữ
  liệu thực sự chạy truy vấn để giữ độ trễ trong ngưỡng chấp nhận được.
\end{itemize}

Khi cách lưu trữ đã rõ, bài toán còn lại là tổ chức mã nguồn sao cho các
quy tắc nghiệp vụ, luồng dữ liệu và truy vấn không đan vào nhau quá mức.
Đó là lý do phần tiếp theo nói về các nguyên tắc lập trình hướng đối
tượng.

\subsubsection{6.1.4. Lập trình hướng đối
tượng}\label{614-lux1eadp-truxecnh-hux1b0ux1edbng-ux111ux1ed1i-tux1b0ux1ee3ng}

Các nguyên tắc lập trình hướng đối tượng (OOP) không được áp dụng như
một khuôn mẫu hình thức, mà như một cách giữ cho hệ thống nhiều tầng vẫn
duy trì được sự rõ ràng khi mở rộng. Điều này thể hiện rõ nhất ở các lớp
máy chủ và giao diện:

\begin{itemize}
\tightlist
\item
  \textbf{TypeScript/JavaScript}: TypeScript được dùng như lớp mô tả rõ
  dữ liệu nào được phép đi qua từng tầng. Các khái niệm như interface,
  generic hay union type được dùng để diễn đạt ``dữ liệu nào hợp lệ,
  dữ liệu nào không'', từ đó giảm lỗi do hiểu sai cấu trúc yêu cầu gửi
  lên máy chủ, phản hồi trả về hoặc bản tin dữ liệu vận hành.
\item
  \textbf{Domain-Driven Design (DDD)}: DDD ở đây có thể hiểu đơn giản là
  chia mã nguồn theo đúng bài toán nghiệp vụ. Thay vì gom phần nhận yêu
  cầu, xử lý nghiệp vụ và truy vấn dữ liệu vào một chỗ, hệ thống tách
  thành các miền như xe, thiết bị, dữ liệu vận hành, cảnh báo và vùng
  giám sát. Cách này giúp
  người đọc lần theo luồng nghiệp vụ dễ hơn và giảm tác động dây chuyền
  khi sửa một phần.
\item
  \textbf{Design Patterns}: Các mẫu thiết kế được dùng như những cách tổ
  chức quen thuộc để code dễ đọc hơn. Chẳng hạn, Repository Pattern tách
  phần đọc/ghi dữ liệu ra khỏi logic nghiệp vụ, Middleware Pattern chia
  một yêu cầu thành từng bước như xác thực và ghi log, còn Observer
  Pattern phù hợp với tình huống một sự kiện mới cần báo cho nhiều nơi
  cùng lúc.
\end{itemize}

Tuy nhiên, các nguyên tắc tổ chức phần mềm vẫn chưa giải quyết trực tiếp
bài toán điện áp, nguồn dự phòng và giao tiếp cảm biến trên bo mạch. Vì
vậy, phần tiếp theo chuyển sang lớp kiến thức điện tử.

\subsubsection{6.1.5. Điện tử số và
Analog}\label{615-ux111iux1ec7n-tux1eed-sux1ed1-vuxe0-analog}

Kiến thức điện tử là phần giúp đồ án thoát khỏi phạm vi của một hệ thống
phần mềm thuần túy. Nhờ lớp kiến thức này, các yêu cầu như bảo vệ ắc quy,
đọc điện áp hay duy trì nguồn dự phòng mới được chuyển thành mạch điện
cụ thể:

\begin{itemize}
\tightlist
\item
  \textbf{Thiết kế mạch quản lý nguồn}: Áp dụng kiến thức về mạch hạ áp
  (giảm điện áp ắc quy xe 12V hoặc 24V xuống 3.3V/5V), mạch nâng áp
  (tăng điện áp từ pin 3.7V lên 5V), và cách tổ chức đường cấp nguồn để
  thiết kế hệ thống cấp nguồn đa đầu vào (ắc quy xe + pin dự phòng).
\item
  \textbf{Đọc giá trị ADC}: Sử dụng kiến thức về bộ chuyển đổi tương tự
  - số (ADC) để đọc điện áp ắc quy xe thông qua mạch chia áp, tính toán
  độ phân giải và sai số.
\item
  \textbf{Giao tiếp cảm biến và bus nhớ}: Áp dụng kiến thức về giao diện
  SPI/I2C để giao tiếp với cảm biến gia tốc LIS3DSH, cấu hình các thanh
  ghi điều khiển, đọc dữ liệu gia tốc 3 trục, và thiết lập ngắt
  (interrupt) cho phát hiện chuyển động. Đồng thời, bus SDMMC 4-bit được
  khai thác để ghi dữ liệu cục bộ lên microSD với yêu cầu tuân thủ
  điện trở kéo lên phần cứng, trình tự gắn/tháo bộ nhớ hợp lệ và tính toàn vẹn dữ liệu khi
  mất nguồn.
\end{itemize}

Khi phần cứng và phần mềm đã gặp nhau trên cùng một thiết bị, dự án còn
cần một lớp kỷ luật để quản lý mã nguồn, đóng gói dịch vụ và giữ quá
trình phát triển không rơi vào hỗn loạn. Đó là vai trò của kỹ thuật phần
mềm ở mục kế tiếp.

\subsubsection{6.1.6. Kỹ thuật phần
mềm}\label{616-kux1ef9-thuux1eadt-phux1ea7n-mux1ec1m}

Các phương pháp và công cụ kỹ thuật phần mềm là lớp keo gắn toàn bộ đồ
án lại với nhau. Chúng không trực tiếp tạo ra tính năng, nhưng quyết
định hệ thống có thể được phát triển, kiểm thử và mở rộng một cách có tổ
chức hay không:

\begin{itemize}
\tightlist
\item
  \textbf{Quản lý mã nguồn bằng Git}: Git không chỉ là lưu
  lại lịch sử thay đổi, mà còn là cách để nhiều người cùng làm việc mà
  không ghi đè lên nhau. Việc chia nhánh, đặt quy ước commit và review
  code giúp quá trình phát triển có tổ chức hơn.
\item
  \textbf{Docker và cách đóng gói dịch vụ thành container}: Kiến thức về
  container hóa được dùng để đóng gói từng dịch vụ như API phía máy chủ,
  giao diện web, broker MQTT hay cơ sở
  dữ liệu theo cùng một cách, nhờ đó môi trường phát triển và triển khai
  ít bị lệch nhau hơn.
\item
  \textbf{Quy trình kiểm tra và triển khai tự động (CI/CD)}: Hiểu biết về quy trình tự động kiểm tra và
  triển khai giúp nhóm thiết kế cấu trúc dự án sao cho mỗi thay đổi đều
  có thể được kiểm tra, đóng gói và đưa lên môi trường chạy thử một cách
  lặp lại.
\item
  \textbf{Cách làm phát triển linh hoạt (Agile)}: Phương pháp phát triển linh hoạt được áp
  dụng theo các vòng ngắn, ưu tiên phần việc tạo giá trị trước và điều
  chỉnh kế hoạch dựa trên phản hồi thực tế thay vì khóa chặt từ đầu.
\end{itemize}

\begin{center}\rule{0.5\linewidth}{0.5pt}\end{center}

\subsection{6.2. Giải quyết các vấn đề kỹ thuật phức tạp -- Complex engineering problems}\label{62-giux1ea3i-quyux1ebft-cuxe1c-vux1ea5n-ux111ux1ec1-kux1ef9-thuux1eadt-phux1ee9c-tux1ea1p--complex-engineering-problems}

Phần trước cho thấy kiến thức nền đã được vận dụng như thế nào. Phần này
đi thêm một bước: nhìn vào những bài toán khó thật sự mà nhóm đã phải
giải quyết trong quá trình triển khai. Việc nêu ra các vấn đề này không
chỉ để mô tả thách thức, mà để làm rõ cách tư duy kỹ thuật đã thay đổi
khi đồ án chuyển từ mức ý tưởng sang mức hệ thống vận hành được.

\subsubsection{6.2.1. Vấn đề 1: Phân tích bản tin OBD2 đa khung qua
BLE}\label{621-vux1ea5n-ux111ux1ec1-1-phuxe2n-tuxedch-bux1ea3n-tin-obd2-ux111a-khung-qua-ble}

\textbf{Mô tả vấn đề:}

Thoạt nhìn, việc lấy dữ liệu OBD2 qua BLE có vẻ chỉ là một bước tích hợp
ngoại vi. Tuy nhiên, khi triển khai thực tế, đây lại là một trong những
thử thách kỹ thuật khó nhất của đồ án vì phải làm việc với thiết bị bên
thứ ba có tài liệu hạn chế. Vấn đề cụ thể bao gồm:

\begin{itemize}
\tightlist
\item
  Giao thức BLE của vgate iCar Pro không có tài liệu chính thức công
  khai. Thông tin giao tiếp (UUID dịch vụ, characteristic, định dạng bản
  tin) phải được lần ngược và kiểm chứng từ các ứng dụng mã nguồn mở và
  bản ghi Bluetooth.
\item
  Bản tin OBD2 có thể kéo dài qua nhiều khung dữ liệu BLE, đặc biệt với
  các lệnh như đọc mã lỗi DTC (Mode 03) hoặc dữ
  liệu động cơ nhiều PID. Việc ghép nối các khung dữ liệu cần tuân theo
  ISO 15765--2 (ISO-TP), một giao thức không được tài liệu OBD2 phổ
  thông đề cập chi tiết.
\item
  Các nguồn tài liệu trực tuyến thường mâu thuẫn nhau về cách xử lý
  phản hồi nhiều khung: một số hướng dẫn chỉ áp dụng cho ELM327 (chip
  thông dịch lệnh), không tương thích trực tiếp với vgate iCar Pro sử dụng
  chip STN1110.
\end{itemize}

\textbf{Cách giải quyết:}

\begin{enumerate}
\def\labelenumi{\arabic{enumi}.}
\tightlist
\item
  \emph{Nghiên cứu và khảo sát}: Phân tích mã nguồn của các dự án mã
  nguồn mở tương tự (esp32-obd2-meter, python-OBD), nghiên cứu tài liệu
  ISO 15765--2, và sử dụng ứng dụng nRF Connect để bắt và phân tích các
  bản tin BLE giữa điện thoại và vgate iCar Pro.
\item
  \emph{Thiết kế lớp trừu tượng}: Xây dựng module phân tích OBD2 với khả
  năng xử lý cả bản tin gọn trong một khung và bản tin kéo dài qua nhiều
  khung, bao gồm khung điều khiển luồng và các khung dữ liệu tiếp theo
  theo chuẩn ISO-TP.
\item
  \emph{Kiểm thử lặp đi lặp lại}: Tạo bộ test với các PID OBD2 phổ biến
  (Mode 01: RPM, Speed, Coolant Temp, Fuel Level) và các lệnh
  nhiều khung (Mode 03: DTC, Mode 09: VIN) để đảm bảo tính chính xác.
\end{enumerate}

\textbf{Bài học rút ra:} Khi làm việc với thiết bị bên thứ ba không có
tài liệu rõ ràng, việc kết hợp phân tích ngược cách giao tiếp, tham
khảo nhiều nguồn và kiểm thử kỹ lưỡng là phương pháp hiệu quả nhất.

Sau bài toán hiểu đúng dữ liệu từ xe, khó khăn tiếp theo không còn nằm ở
từng bản tin đơn lẻ mà ở việc đưa toàn bộ luồng dữ liệu đó đi trọn hành
trình đến máy chủ và giao diện. Vì vậy, mục sau chuyển sang đường ống dữ
liệu thời gian thực.

\subsubsection{6.2.2. Vấn đề 2: Đường ống dữ liệu thời gian thực với đảm
bảo phân
phối}\label{622-vux1ea5n-ux111ux1ec1-2-ux111ux1b0ux1eddng-ux1ed1ng-dux1eef-liux1ec7u-thux1eddi-gian-thux1ef1c-vux1edbi-ux111ux1ea3m-bux1ea3o-phuxe2n-phux1ed1i}

\textbf{Mô tả vấn đề:}

Đường ống dữ liệu thời gian thực có thể hiểu là toàn bộ hành trình từ
lúc thiết bị phát bản tin đến lúc dữ liệu xuất hiện trên màn hình giám
sát hoặc được lưu trong cơ sở dữ liệu. Phần khó của bài toán này là ba yêu cầu
thường kéo theo các hướng ngược nhau: muốn nhanh thì dễ bỏ sót, muốn
chắc chắn thì dễ tăng chi phí, còn muốn lưu thật nhiều thì hệ thống dễ
nặng lên. Trong đồ án, hệ thống phải xử lý luồng dữ liệu vận hành liên
tục từ nhiều thiết bị IoT (GPS, trạng thái nguồn, IMU; sẵn sàng mở rộng
thêm OBD2) với các yêu cầu:

\begin{itemize}
\tightlist
\item
  Dữ liệu phải được cất đồng thời vào hai nơi khác nhau: một nơi chuyên
  lưu dữ liệu theo thời gian để vẽ biểu đồ và xem lịch sử dày đặc, một
  nơi lưu dữ liệu nghiệp vụ có quan hệ rõ ràng. Đây chính là bài toán
  một lần nhận nhưng phải ghi ra hai nơi (dual-write):
  phải lưu phục vụ hai mục đích sử dụng khác nhau.
\item
  Khi mất kết nối mạng 4G, hệ thống phải khôi phục đường gửi dữ liệu đủ
  nhanh để luồng vận hành không bị gãy; nếu mất sóng quá lâu thì thiết
  bị phải có chỗ giữ tạm dữ liệu để gửi lại sau.
\item
  Mức đảm bảo gửi nhận của MQTT (QoS) phải được chọn phù hợp cho từng
  loại dữ liệu: dữ liệu gửi liên tục có thể ưu tiên nhanh hơn, còn cảnh
  báo và lệnh điều khiển phải ưu tiên chắc chắn hơn.
\end{itemize}

\textbf{Cách giải quyết:}

\begin{enumerate}
\def\labelenumi{\arabic{enumi}.}
\tightlist
\item
  \emph{Dịch vụ trung gian riêng cho luồng dữ liệu}: Tracking\_MqttBridge
  được thiết kế như một lớp đứng giữa broker MQTT và các nơi lưu trữ.
  Dịch vụ này nhận dữ liệu từ thiết bị, ghi dữ liệu theo thời gian sang
  VictoriaMetrics, ghi log sang VictoriaLogs và cập nhật PostgreSQL cho
  các trạng thái quan trọng. Nhờ tách riêng, hệ thống dễ kiểm soát hơn
  chỗ nào chỉ lo tiếp nhận dữ liệu, chỗ nào lo nghiệp vụ người dùng.
\item
  \emph{Tự phục hồi kết nối trên thiết bị kết hợp lưu tạm cục bộ}: Thiết
  bị tự kết nối lại khi có mạng, duy trì tín hiệu giữ kết nối của MQTT và lưu
  các cấu hình cần thiết trong NVS, tức vùng nhớ không mất dữ liệu sau
  khi khởi động lại. Song song đó, dữ liệu và sự kiện được lưu vào hàng
  đợi microSD và phát lại theo đúng thứ tự khi kết nối phục hồi, giúp
  giảm mất mát dữ liệu trong các đợt mất sóng dài. Ở bản hiện tại, hệ
  thống giả định thẻ nhớ đã được cắm sẵn từ lúc khởi động để giữ ổn
  định.
\item
  \emph{Phân mức dữ liệu theo độ quan trọng}: Áp dụng QoS 0 cho dữ liệu
  vị trí GPS gửi dày, tức ưu tiên tốc độ hơn việc xác nhận từng bản tin;
  QoS 1 cho cảnh báo và sự kiện quan trọng, tức yêu cầu phía nhận báo đã
  nhận ít nhất một lần. Cách chia này giữ được cân bằng giữa hiệu suất
  và độ tin cậy.
\end{enumerate}

\textbf{Bài học rút ra:} Thiết kế đường ống dữ liệu cần xem xét tất cả
các tình huống hỏng hóc như mất mạng, máy chủ quá tải hoặc dữ liệu bị lệch nhau
từ giai đoạn thiết kế, không để đến giai đoạn tích hợp mới xử lý.

Tuy nhiên, một đường ống dữ liệu tốt chỉ có ý nghĩa khi thiết bị còn đủ
năng lượng để duy trì việc thu thập và gửi dữ liệu trong môi trường xe.
Vì vậy, vấn đề kế tiếp chuyển sang quản lý nguồn.

\subsubsection{6.2.3. Vấn đề 3: Quản lý năng lượng với nhiều nguồn
cấp}\label{623-vux1ea5n-ux111ux1ec1-3-quux1ea3n-luxfd-nux103ng-lux1b0ux1ee3ng-vux1edbi-nhiux1ec1u-nguux1ed3n-cux1ea5p}

\textbf{Mô tả vấn đề:}

Quản lý năng lượng là vấn đề không dễ nhận thấy khi mới nhìn hệ thống từ
góc độ phần mềm, nhưng lại quyết định trực tiếp độ tin cậy của thiết bị
trong xe. Hệ thống phần cứng phải hoạt động với hai nguồn năng lượng có
đặc tính rất khác nhau:

\begin{itemize}
\tightlist
\item
  Ắc quy xe 12V hoặc 24V DC (dao động tùy trạng thái sạc và tải), là
  nguồn chính khi xe hoạt động.
\item
  Pin dự phòng 18650 1S Li-ion 3.7V (dao động 2.8V - 4.2V), là nguồn
  dùng khi ắc quy xe bị ngắt hoặc điện áp quá thấp.
\item
  Việc chuyển đổi giữa hai nguồn phải diễn ra liền mạch, nghĩa là đổi
  nguồn nhưng thiết bị không bị sụt điện, tránh reset hoặc mất dữ
  liệu.
\end{itemize}

\textbf{Cách giải quyết:}

\begin{enumerate}
\def\labelenumi{\arabic{enumi}.}
\tightlist
\item
  \emph{Thiết kế đường cấp nguồn}: Thiết kế mạch đường cấp nguồn dùng diode OR
  giữa MP2482 (5V chính) và SX1308 (5V dự phòng), kết hợp GPIO18 để điều
  khiển nhánh nguồn theo profile 12V/24V.
\item
  \emph{Low Voltage Disconnect (LVD)}: Hiện thực LVD bằng khối LVD phần
  cứng kết hợp ADC firmware, dùng ngưỡng cho hai hệ điện:
  12V (\(OFF = 12.0\,\mathrm{V}\), \(ON = 12.2\,\mathrm{V}\)) và
  24V (\(OFF = 24.0\,\mathrm{V}\), \(ON = 24.4\,\mathrm{V}\)),
  bảo vệ ắc quy không bị rút cạn
  quá mức.
\item
  \emph{Bộ sạc pin dự phòng}: Tích hợp IC sạc TP5100 (input 5V từ
  MP2482, output 4.2V ở chế độ 1 cell) để sạc pin 18650 1S khi điều kiện
  nguồn cho phép.
\item
  \emph{Giám sát điện áp bằng firmware}: Đọc điện áp ắc quy và pin dự
  phòng liên tục qua ADC, kết hợp trạng thái GPIO19 (HIGH = báo điện áp thấp)
  để gửi cảnh báo sớm và điều phối chuyển nguồn.
\end{enumerate}

\textbf{Bài học rút ra:} Thiết kế hệ thống năng lượng cho IoT trong môi
trường ô tô phải xét đồng thời điện áp dao động, chuyển đổi nguồn liền
mạch, bảo vệ ắc quy và khả năng giám sát từ xa. Nếu bỏ sót một mắt xích,
độ tin cậy của toàn hệ thống sẽ giảm rõ rệt.

Khi thiết bị đã đứng vững về nguồn, bài toán còn lại là làm sao lớp máy
chủ không trở thành điểm nghẽn mới khi số thiết bị và người dùng cùng
tăng. Phần tiếp theo vì thế chuyển sang kiến trúc máy chủ có khả năng mở
rá»™ng.

\subsubsection{6.2.4. Vấn đề 4: Kiến trúc đám mây có khả năng mở rộng
cho
IoT}\label{624-vux1ea5n-ux111ux1ec1-4-kiux1ebfn-truxfac-ux111uxe1m-muxe2y-cuxf3-khux1ea3-nux103ng-mux1edf-rux1ed9ng-cho-iot}

\textbf{Mô tả vấn đề:}

Bài toán hạ tầng máy chủ của đồ án không dừng ở việc ``đưa dữ liệu lên
máy chủ'',
mà phải giải được ba yêu cầu cùng lúc: nhiều thiết bị, nhiều người dùng
và mức an toàn đủ cho dữ liệu vận hành. Vì vậy, hệ thống cần xử lý dữ
liệu từ nhiều thiết bị đồng thời, cung cấp giao diện thời gian thực cho
nhiều người dùng, và đảm bảo bảo mật:

\begin{itemize}
\tightlist
\item
  Mỗi thiết bị có kênh dữ liệu MQTT riêng (ví dụ
  \texttt{v1/\{device\_id\}/rawdata} và
  \texttt{v1/\{device\_id\}/commands}), nên phải có quy tắc giới hạn
  quyền (ACL) để thiết bị chỉ được gửi và nhận trên đúng kênh của mình.
\item
  Nhiều người dùng có thể xem cùng một xe trên màn hình giám sát, tạo ra nhiều
  kết nối WebSocket đồng thời cần được quản lý.
\item
  Cơ chế giữ đăng nhập phải đủ an toàn, có thể thu hồi khi cần, nhưng
  cũng không được tạo ra nút thắt cổ chai khi số lượng yêu cầu tăng cao.
\end{itemize}

\textbf{Cách giải quyết:}

\begin{enumerate}
\def\labelenumi{\arabic{enumi}.}
\tightlist
\item
  \emph{Quy tắc giới hạn từng thiết bị}: EMQX được cấu hình để mỗi thiết
  bị chỉ được gửi và nhận trên đúng kênh dữ liệu của nó, nhờ đó giảm
  nguy cơ một thiết bị giả mạo dữ liệu của thiết bị khác. Mỗi thiết bị
  cũng có thông tin đăng nhập riêng.
\item
  \emph{Nhóm người xem theo từng xe}: Socket.IO rooms, tức các nhóm nhận
  dữ liệu trong cùng một kênh, được dùng để gom những người đang theo
  dõi cùng một xe vào một nhóm. Khi có dữ liệu mới của xe đó, máy chủ
  chỉ đẩy cho đúng nhóm liên quan, tránh phát tán
  dữ liệu không cần thiết cho toàn bộ người dùng.
\item
  \emph{Phiên đăng nhập có thể thu hồi ngay}: Hệ thống giữ thông tin
  phiên trong cơ sở dữ liệu, lưu token dưới dạng mã băm SHA-256 thay vì
  lưu nguyên văn, để có thể khóa ngay một phiên nếu phát hiện rủi ro.
  Cách này được ưu tiên hơn JWT trong giai đoạn hiện tại vì dễ kiểm soát
  việc thu hồi. Ở giai đoạn sau có thể bổ sung Redis cache để giảm tải
  truy vấn xác thực.
\item
  \emph{Triển khai tách từng dịch vụ}: Mỗi dịch vụ như API phía máy chủ,
  giao diện web, MQTT Bridge hay cơ sở dữ liệu đều có cấu hình triển khai
  riêng, nhưng vẫn nằm trong cùng một mạng nội bộ để giao tiếp với nhau.
  Cách tách này giúp mở rộng từng phần theo đúng nhu cầu thay vì phải
  nâng cả hệ thống cùng lúc.
\end{enumerate}

\textbf{Bài học rút ra:} Hạ tầng máy chủ cho IoT cần được thiết kế từ
đầu với khả năng mở rộng bằng cách thêm nhiều bản sao dịch vụ khi cần,
thay vì dồn toàn bộ tải lên một máy mạnh hơn. Việc tách riêng từng dịch
vụ và giới hạn rõ quyền của từng thiết bị là nền tảng cho bảo mật và cho
quản lý đội xe ở quy mô lớn.

\begin{center}\rule{0.5\linewidth}{0.5pt}\end{center}

\subsection{6.3. Tác động đạo đức và xã hội -- Ethical and social impacts}\label{63-tuxe1c-ux111ux1ed9ng-ux111ux1ea1o-ux111ux1ee9c-vuxe0-xuxe3-hux1ed9i--ethical-and-social-impacts}

Một hệ thống theo dõi phương tiện không chỉ là bài toán kỹ thuật. Khi hệ
thống bắt đầu thu thập dữ liệu vị trí, hành trình và thói quen sử dụng
xe, nó đồng thời chạm tới các câu hỏi về quyền riêng tư, trách nhiệm sử
dụng dữ liệu và ảnh hưởng xã hội. Vì vậy, việc đánh giá đồ án sẽ chưa đủ
nếu chỉ dừng ở hiệu năng kỹ thuật.

\subsubsection{6.3.1. Quyền riêng tư và bảo vệ dữ liệu cá
nhân}\label{631-quyux1ec1n-riuxeang-tux1b0-vuxe0-bux1ea3o-vux1ec7-dux1eef-liux1ec7u-cuxe1-nhuxe2n}

Hệ thống theo dõi phương tiện liên tục thu thập dữ liệu vị trí GPS, hành
trình và thói quen sử dụng xe của người lái. Đây đều là dữ liệu nhạy cảm
vì chúng không chỉ mô tả phương tiện đang ở đâu, mà còn phản ánh hành vi
và nhịp sinh hoạt của con người gắn với phương tiện đó. Vì vậy, quá trình
phát triển và triển khai hệ thống cần được xem xét nghiêm túc dưới góc
độ đạo đức và bảo vệ dữ liệu.

\textbf{Các biện pháp bảo vệ quyền riêng tư đã được áp dụng:}


{\thesissettableid{6.3.1A}{tab-ch6-tong-hop-thong-tin-ve-stt-bien-phap-mo-ta-chi-tiet}
\def\LTcaptype{table}
\begin{longtable}{|L{0.160\linewidth}|L{0.224\linewidth}|L{0.580\linewidth}|}
\caption{Tổng hợp thông tin về STT, Biện pháp, Mô tả chi tiết}\label{tab:ch6-tong-hop-thong-tin-ve-stt-bien-phap-mo-ta-chi-tiet}\\
\hline \textbf{STT} & \textbf{Biện pháp} & \textbf{Mô tả chi tiết} \\
\\\hline
\endfirsthead
\caption[]{Tổng hợp thông tin về STT, Biện pháp, Mô tả chi tiết (tiếp theo)}\\
\hline \textbf{STT} & \textbf{Biện pháp} & \textbf{Mô tả chi tiết} \\
\\\hline
\endhead
\\\hline
\endfoot
\endlastfoot
1 & Mã hóa dữ liệu truyền tải & Dữ liệu MQTT được mã hóa TLS trong môi
trường production, ngăn chặn nghe lén (eavesdropping) trên đường
truyền \\\hline
2 & Phân quyền truy cập (RBAC) & Hệ thống phân quyền theo vai trò
(admin, manager, viewer), chỉ những người có quyền mới xem được dữ liệu
cụ thể \\\hline
3 & Chính sách lưu trữ dữ liệu & Dữ liệu telemetry được lưu trữ có thời
hạn (retention policy), tự động xóa sau 90 ngày để giảm rủi ro lộ lọt dữ
liệu \\\hline
4 & Session-based authentication & Token xác thực được hash SHA-256, lưu
trữ trong database, có khả năng thu hồi ngay lập tức khi cần \\\hline
5 & MQTT ACL per device & Mỗi thiết bị chỉ truy cập được dữ liệu của
mình, ngăn chặn truy cập trái phép dữ liệu thiết bị khác \\\hline
\end{longtable}
}


\textbf{Khuyến nghị bổ sung (chưa hiện thực):}

\begin{itemize}
\tightlist
\item
  Thông báo rõ ràng cho người thuê xe về việc xe được giám sát
  (transparency).
\item
  Cho phép người thuê xe xem dữ liệu của mình (data access rights).
\item
  Tuân thủ Luật An toàn thông tin mạng Việt Nam (Luật số 24/2018/QH14)
  và các quy định về bảo vệ dữ liệu cá nhân.
\end{itemize}

Bảo vệ dữ liệu là điều kiện cần, nhưng với một thiết bị gắn trực tiếp
lên xe thì vẫn còn một câu hỏi khác quan trọng không kém: hệ thống có
làm ảnh hưởng đến an toàn vận hành hay không. Vì vậy, mục tiếp theo
chuyển sang trách nhiệm kỹ thuật đối với phương tiện.

\subsubsection{6.3.2. An toàn phương tiện và trách nhiệm kỹ
thuật}\label{632-an-touxe0n-phux1b0ux1a1ng-tiux1ec7n-vuxe0-truxe1ch-nhiux1ec7m-kux1ef9-thuux1eadt}

Việc truy cập hệ thống OBD2 của xe đặt ra một câu hỏi quan trọng hơn
chính bản thân khả năng đọc dữ liệu: thiết bị theo dõi có vô tình làm
ảnh hưởng đến độ an toàn vận hành của xe hay không? Đây là lý do các
quyết định thiết kế ở phần này phải ưu tiên nguyên tắc an toàn trước khi
tối đa hóa tính năng.

\textbf{Các biện pháp đảm bảo an toàn đã được áp dụng:}

\begin{itemize}
\tightlist
\item
  \textbf{Chỉ đọc (read-only) OBD2 PIDs}: Hệ thống chỉ sử dụng các lệnh
  OBD2 Mode 01 (Current Data) và Mode 03 (Stored DTCs) để đọc dữ liệu.
  Tuyệt đối không gửi các lệnh ghi (write commands) hoặc lệnh điều khiển
  động cơ (Mode 04: Clear DTCs, Mode 08: Control operations). Đây là
  nguyên tắc an toàn cơ bản khi làm việc với hệ thống chẩn đoán xe.
\item
  \textbf{Kết nối BLE gián tiếp}: Sử dụng OBD2 adapter vgate iCar Pro
  (thiết bị đã được chứng nhận an toàn) làm lớp trung gian, không kết
  nối trực tiếp vào bus CAN của xe. Điều này giảm rủi ro gây nhiễu hoặc
  xung đột trên bus truyền thông nội bộ xe.
\item
  \textbf{Không điều khiển xe từ xa}: Hệ thống không có chức năng điều
  khiển động cơ (tắt bơm xăng, khóa xe). Đây là quyết định thiết kế có
  chủ đích để tránh rủi ro an toàn nghiêm trọng.
\end{itemize}

Khi các giới hạn an toàn đã được chốt rõ, mới có thể đánh giá công bằng
mặt tích cực mà hệ thống mang lại cho doanh nghiệp và xã hội. Phần tiếp
theo vì thế chuyển sang tác động xã hội tích cực.

\subsubsection{6.3.3. Tác động xã hội tích
cá»±c}\label{633-tuxe1c-ux111ux1ed9ng-xuxe3-hux1ed9i-tuxedch-cux1ef1c}

Ở chiều ngược lại, nếu được triển khai đúng cách, hệ thống IoT giám sát
phương tiện cũng có thể tạo ra những giá trị xã hội rõ ràng chứ không
chỉ là công cụ kiểm soát. Một số tác động tích cực có thể nhận thấy gồm:

\textbf{Hỗ trợ doanh nghiệp nhỏ và vừa:}

\begin{itemize}
\tightlist
\item
  Cung cấp giải pháp quản lý đội xe với chi phí hợp lý (dưới 2.000.000
  VND/thiết bị, chi phí 4G khoảng 70.000 VND/tháng), giúp các doanh
  nghiệp cho thuê xe nhỏ có thể tiếp cận công nghệ giám sát mà trước đây
  chỉ dành cho các công ty lớn.
\item
  Giảm thiệt hại do mất trộm, sử dụng sai mục đích, và hư hỏng thiết bị
  nhờ phát hiện sớm.
\end{itemize}

\textbf{Giảm tai nạn giao thông:}

\begin{itemize}
\tightlist
\item
  Giám sát tốc độ và hành vi lái xe (phanh gấp, tăng tốc đột ngột) giúp
  nhận diện và cảnh báo các hành vi lái xe nguy hiểm.
\item
  Cảnh báo vượt vùng cho phép giúp đảm bảo xe hoạt động trong phạm vi an toàn đã
  thỏa thuận.
\end{itemize}

\textbf{Tối ưu hóa nhiên liệu và giảm phát thải:}

\begin{itemize}
\tightlist
\item
  Phân tích dữ liệu OBD2 (tiêu hao nhiên liệu, RPM, tốc độ) giúp đánh
  giá hiệu suất sử dụng nhiên liệu.
\item
  Dữ liệu hành trình giúp tối ưu hóa tuyến đường, giảm quãng đường chạy
  không tải, giảm phát thải CO2.
\end{itemize}

Những lợi ích này không phủ nhận các rủi ro đi kèm. Vì vậy, ngay sau
phần tác động tích cực, luận văn chỉ ra rõ các lo ngại cần được quản lý
để hệ thống không biến thành công cụ gây áp lực hoặc xâm phạm quyền
riêng tư.

\subsubsection{6.3.4. Những lo ngại cần xem
xét}\label{634-nhux1eefng-lo-ngux1ea1i-cux1ea7n-xem-xuxe9t}

\begin{itemize}
\tightlist
\item
  \textbf{Giám sát quá mức}: Cần cân bằng giữa nhu cầu giám sát của
  doanh nghiệp và quyền riêng tư của người lái xe. Không nên sử dụng dữ
  liệu vị trí để theo dõi hoạt động cá nhân ngoài phạm vi hợp đồng cho
  thuê.
\item
  \textbf{Phụ thuộc công nghệ}: Hệ thống cần có cơ chế dự phòng
  (fallback) khi mất kết nối hoặc thiết bị hỏng, không để việc mất kết
  nối làm ảnh hưởng đến hoạt động bình thường của xe.
\item
  \textbf{Trách nhiệm dữ liệu}: Doanh nghiệp sử dụng hệ thống có trách
  nhiệm bảo vệ dữ liệu thu thập được, không chia sẻ cho bên thứ ba khi
  chưa có sự đồng ý của người liên quan.
\end{itemize}

\begin{center}\rule{0.5\linewidth}{0.5pt}\end{center}

\subsection{6.4. Tổng kết và bài học kinh nghiệm -- Reflection and Lessons Learned}\label{64-tux1ed5ng-kux1ebft-vuxe0-buxe0i-hux1ecdc-kinh-nghiux1ec7m--reflection-and-lessons-learned}

Nếu Chương 5 đánh giá hệ thống bằng số liệu và Chương 6.1--6.3 nhìn lại
kiến thức, thách thức và tác động xã hội, thì mục này là phần khép lại
toàn bộ hành trình thực hiện đồ án. Trọng tâm ở đây không phải kể lại
những gì đã làm, mà rút ra những điều thực sự có giá trị cho các giai
đoạn phát triển sau.

\subsubsection{6.4.1. Tầm quan trọng của thiết kế kiến trúc trước khi
lập
trình}\label{641-tux1ea7m-quan-trux1ecdng-cux1ee7a-thiux1ebft-kux1ebf-kiux1ebfn-truxfac-trux1b0ux1edbc-khi-lux1eadp-truxecnh}

Một trong những bài học quan trọng nhất của dự án là phải ưu tiên thiết
kế kiến trúc hệ thống trước khi bắt tay vào lập trình. Bài học này không
đến từ lý thuyết, mà đến từ chính giai đoạn đầu của đồ án khi nhóm từng
có xu hướng triển khai nhanh từng mô-đun trước khi xây dựng xong mô hình
tương tác tổng thể.

Khi phát sinh các lỗi tích hợp như hai mô-đun hiểu khác nhau về cùng một
dữ liệu, định dạng truyền nhận không thống nhất hoặc các phần phụ thuộc
chéo lẫn nhau, nhóm chuyển sang cách tiếp cận có kỷ luật hơn: xác lập sơ
đồ kiến trúc và luồng dữ liệu, chốt quy ước đầu vào/đầu ra của API trước
khi hiện thực, hoàn thiện sơ đồ thực thể - quan hệ trước khi tạo bảng,
chuẩn hóa cách các dịch vụ nhìn thấy nhau và dùng cổng nào, đồng thời mô
tả rõ các tình huống hỏng hóc có thể xảy ra ở nhánh lưu trữ cục bộ như
mất mạng, đầy dung lượng, gắn thẻ nhớ thất bại hoặc xác nhận nhận dữ
liệu đến quá chậm trước khi tối ưu hiệu năng.

Cách làm này giúp giảm đáng kể thời gian sửa lỗi tích hợp và hạn chế
khối lượng công việc phải làm lại ở các giai đoạn sau.

\subsubsection{6.4.2. Độ phức tạp của hệ thống IoT nhiều
tầng}\label{642-ux111ux1ed9-phux1ee9c-tux1ea1p-cux1ee7a-hux1ec7-thux1ed1ng-iot-nhiux1ec1u-tux1ea7ng}

Dự án cho thấy một hệ thống IoT nhiều tầng phức tạp hơn đáng kể so với
một ứng dụng web thông thường. Độ phức tạp không nằm ở từng lớp riêng
lẻ, mà ở chỗ mỗi lớp có ràng buộc riêng và các ràng buộc đó chỉ bộc lộ
rõ khi ghép lại thành một chuỗi vận hành hoàn chỉnh. Hệ thống này đồng
thời chạm vào bốn lớp công việc khác nhau:


{\thesissettableid{6.4.2A}{tab-ch6-tong-hop-thong-tin-ve-tang-ngon-ngu-cong-nghe-thach-thuc}
\def\LTcaptype{table}
\begin{longtable}{|L{0.171\linewidth}|L{0.342\linewidth}|L{0.452\linewidth}|}
\caption{Tổng hợp thông tin về Tầng, Ngôn ngữ / Công nghệ, Thách thức chính}\label{tab:ch6-tong-hop-thong-tin-ve-tang-ngon-ngu-cong-nghe-thach-thuc}\\
\hline \textbf{Tầng} & \textbf{Ngôn ngữ / Công nghệ} & \textbf{Thách thức chính} \\
\\\hline
\endfirsthead
\caption[]{Tổng hợp thông tin về Tầng, Ngôn ngữ / Công nghệ, Thách thức chính (tiếp theo)}\\
\hline \textbf{Tầng} & \textbf{Ngôn ngữ / Công nghệ} & \textbf{Thách thức chính} \\
\\\hline
\endhead
\\\hline
\endfoot
\endlastfoot
Phần cứng & Thiết kế mạch, PCB & Dễ phát sinh lỗi nguồn hoặc tín hiệu,
khó đo đạc và khó tìm đúng điểm hỏng \\\hline
Firmware & C/C++ trên ESP32-S3, bộ điều phối công việc nền & Phải giữ bộ
nhớ ổn định, canh thời gian chính xác và tránh xung đột truy cập \\\hline
Backend & TypeScript, dịch vụ API & Phải tổ chức dữ liệu lớn, giữ phản
hồi ổn định và bảo vệ truy cập \\\hline
Frontend & TypeScript, dashboard web & Phải giữ giao diện dễ dùng, cập
nhật dữ liệu mới liên tục mà không giật lag \\\hline
\end{longtable}
}


Mỗi tầng đòi hỏi một nhóm kỹ năng chuyên môn riêng, và các lỗi thường
xuất hiện nhiều nhất ở giai đoạn tích hợp giữa các tầng. Bài học rút ra
là với hệ thống đa tầng, cần thực hiện kiểm thử tích hợp sớm thay vì chờ
đến khi toàn bộ thành phần hoàn tất mới kiểm thử.

Nói cách khác, độ phức tạp của hệ thống không chỉ đến từ số tầng, mà còn
đến từ việc phải phối hợp nhiều công cụ trưởng thành với nhau. Từ góc
nhìn đó, phần tiếp theo cho thấy vì sao hệ sinh thái mã nguồn mở lại trở
thành một lợi thế rõ rệt của đồ án.

\subsubsection{6.4.3. Sức mạnh của hệ sinh thái mã nguồn
mở}\label{643-sux1ee9c-mux1ea1nh-cux1ee7a-hux1ec7-sinh-thuxe1i-muxe3-nguux1ed3n-mux1edf}

Dự án sử dụng hoàn toàn công nghệ mã nguồn mở. Quá trình triển khai cho
thấy sức mạnh của hệ sinh thái này không chỉ nằm ở chi phí bằng 0, mà ở
khả năng ghép nối nhiều công cụ trưởng thành thành một hệ thống IoT khá
hoàn chỉnh:

\begin{itemize}
\tightlist
\item
  \textbf{Nền tảng phát triển firmware cho ESP32-S3 (ESP-IDF)}: cung cấp
  sẵn các thư viện cần thiết để giao tiếp với Wi-Fi, BLE và các ngoại vi
  phần cứng, nhờ đó nhóm không phải tự xây mọi lớp nền từ đầu.
\item
  \textbf{Tầng xử lý yêu cầu web và API (Express.js + TypeScript)}: đủ
  nhẹ để phát triển nhanh nhưng vẫn rõ cấu trúc, đồng thời TypeScript
  giúp phát hiện sớm nhiều lỗi dữ liệu trước khi chạy thật.
\item
  \textbf{Tầng giao diện web (Next.js 15 + React 19)}: phù hợp để tổ
  chức nhiều màn hình giám sát, hiển thị dữ liệu thay đổi liên tục và
  vẫn giữ trải nghiệm đủ mượt cho người vận hành.
\item
  \textbf{PostgreSQL 16}: là nơi lưu dữ liệu nghiệp vụ cần quan hệ chặt
  chẽ; hỗ trợ PostGIS để xử lý các bài toán không gian như vùng giám sát
  trên bản đồ.
\item
  \textbf{EMQX}: đóng vai trò broker MQTT ổn định giữa thiết bị và máy
  chủ; có thể mở rộng thêm khi số lượng thiết bị tăng.
\item
  \textbf{VictoriaMetrics}: phù hợp cho dữ liệu thay đổi liên tục theo
  thời gian, ghi nhanh, tốn ít tài nguyên và dùng được với định dạng số
  liệu chuẩn mà nhiều công cụ giám sát hỗ trợ.
\item
  \textbf{Docker}: giúp đóng gói từng dịch vụ theo cách nhất quán để môi
  trường phát triển và triển khai không bị lệch nhau quá nhiều.
\end{itemize}

Chi phí giấy phép phần mềm của toàn bộ hệ thống là 0 VND, nên ngân sách
có thể tập trung cho phần cứng và hạ tầng máy chủ.

Tuy nhiên, công cụ tốt chỉ tạo ra lợi thế ban đầu; để lợi thế đó biến
thành một hệ thống tin cậy, nhóm vẫn phải kiểm thử ở mọi tầng thay vì
dựa vào niềm tin rằng các thành phần trưởng thành sẽ tự ghép đúng với
nhau.

\subsubsection{6.4.4. Kiểm thử ở mọi tầng là bắt
buá»™c}\label{644-kiux1ec3m-thux1eed-ux1edf-mux1ecdi-tux1ea7ng-luxe0-bux1eaft-buux1ed9c}

Dự án cho thấy kiểm thử ở mọi tầng của hệ thống là yêu cầu bắt buộc.
Trong một hệ thống nhiều lớp, rất ít lỗi tự lộ ra đúng tại nơi chúng
phát sinh; phần lớn chỉ biểu hiện ở tầng khác. Vì vậy, kiểm thử phải đi
theo toàn bộ hành trình dữ liệu:

\begin{itemize}
\tightlist
\item
  \textbf{Kiểm thử đơn vị (unit test)}: dùng để kiểm tra các hàm nhỏ và
  các quy tắc xử lý dữ liệu, ví dụ phân tích bản tin OBD2 hoặc tính toán
  năng lượng trong firmware và dịch vụ máy chủ.
\item
  \textbf{Kiểm thử tích hợp (integration test)}: dùng để kiểm tra cách
  các dịch vụ nói chuyện với nhau, ví dụ MQTT Bridge nhận dữ liệu từ
  EMQX rồi ghi vào VictoriaMetrics, API phía máy chủ đọc/ghi PostgreSQL, hay
  luồng phát lại dữ liệu đã lưu từ microSD về đúng kênh MQTT.
\item
  \textbf{Kiểm thử đầu-cuối (end-to-end)}: mô phỏng toàn bộ hành trình
  dữ liệu từ thiết bị giả lập đến giao diện web để kiểm tra người dùng
  thực sự nhìn thấy gì trên màn hình giám sát, kể cả trong kịch bản mất mạng rồi
  phục hồi để xác nhận dữ liệu phát lại.
\end{itemize}

Các lỗi được phát hiện ở giai đoạn kiểm thử sớm (unit test) có chi phí
khắc phục thấp hơn rất nhiều so với lỗi chỉ được phát hiện ở giai đoạn
tích hợp hoặc triển khai sản xuất {[}1{]}.

Nhìn lại toàn bộ dự án, chính quá trình kiểm thử rồi sửa đi sửa lại ở
nhiều tầng đã tạo ra phần trưởng thành nghề nghiệp rõ nhất. Vì vậy, mục
cuối cùng chuyển sang tổng kết cá nhân.

\subsubsection{6.4.5. Tổng kết cá
nhân}\label{645-tux1ed5ng-kux1ebft-cuxe1-nhuxe2n}

Dự án "IoT Vehicle Tracking System" đưa người thực hiện đi trọn một vòng
từ ý tưởng, thiết kế mạch, đặt mua linh kiện, chế tạo phần cứng, tích
hợp và đánh giá. Giá trị lớn nhất của quá trình này không nằm ở việc
ghép đủ các thành phần, mà ở chỗ từng quyết định kỹ thuật đều phải được
kiểm chứng lại bằng số liệu, bằng thử nghiệm và bằng mức độ phù hợp với
điều kiện triển khai thật. Chính quá trình lặp đi lặp lại giữa giả định
và kiểm chứng đó mới là phần trưởng thành quan trọng nhất mà đồ án mang
lại.

Kinh nghiệm rút ra từ dự án, đặc biệt ở các mảng kiến trúc hệ thống,
quản lý năng lượng cho thiết bị IoT, bảo mật dịch vụ máy chủ và tổ chức
luồng dữ liệu thời gian thực, là nền tảng quan trọng cho công việc chuyên môn sau
này.

\begin{center}\rule{0.5\linewidth}{0.5pt}\end{center}

\subsection{Kết luận chương
6}\label{kux1ebft-luux1eadn-chux1b0ux1a1ng-6}

Chương 6 khép lại câu chuyện của đồ án bằng một góc nhìn rộng hơn so với
các chương trước: xây dựng một hệ thống IoT giám sát phương tiện không
phải là việc ghép nối đơn giản giữa phần cứng, firmware và web, mà là
quá trình liên tục cân bằng giữa ràng buộc kỹ thuật, nhu cầu vận hành và
trách nhiệm dữ liệu. Những bài học rút ra về kiến trúc, kiểm thử tích
hợp, quản lý năng lượng, OTA và đạo đức dữ liệu vì thế không chỉ là phần
tổng kết cho đồ án hiện tại, mà còn là nền tảng thực tế cho các giai
đoạn hoàn thiện tiếp theo.

\begin{center}\rule{0.5\linewidth}{0.5pt}\end{center}

\subsection{Tài liệu tham khảo Chương
6}\label{tuxe0i-liux1ec7u-tham-khux1ea3o-chux1b0ux1a1ng-6}

{[}1{]} B. Boehm and V. R. Basili, "Software Defect Reduction Top 10
List," \emph{IEEE Computer}, vol. 34, no. 1, pp. 135--137, 2001.

\majorheading{TÀI LIỆU TRÍCH DẪN - REFERENCES}

\begin{center}\rule{0.5\linewidth}{0.5pt}\end{center}

\subsection{{[}A{]} Tài liệu tiếng
Việt}\label{a-tuxe0i-liux1ec7u-tiux1ebfng-viux1ec7t}

{[}1{]} Bộ Giao thông Vận tải, "Báo cáo thống kê phương tiện giao thông
đường bộ Việt Nam giai đoạn 2020--2025," Hà Nội, 2024.

{[}2{]} Nguyễn Văn A và Trần Thị B, "Ứng dụng Internet vạn vật trong
giám sát và quản lý phương tiện giao thông tại Việt Nam," \emph{Tạp chí
Khoa học và Công nghệ}, tập 58, số 3, trang 45--52, 2023.

{[}3{]} Lê Minh C, "Thiết kế hệ thống định vị và theo dõi phương tiện sử
dụng GPS và mạng di động," Luận văn thạc sĩ, Đại học Bách Khoa TP. Hồ
Chí Minh, 2022.

{[}4{]} Bộ Thông tin và Truyền thông, "Luật An toàn thông tin mạng,"
Luật số 24/2018/QH14, Hà Nội, 2018.

\begin{center}\rule{0.5\linewidth}{0.5pt}\end{center}

\subsection{{[}B{]} Tài liệu tiếng Anh (Books \&
Papers)}\label{b-tuxe0i-liux1ec7u-tiux1ebfng-anh-books--papers}

{[}5{]} A. Al-Fuqaha, M. Guizani, M. Mohammadi, M. Aledhari, and M.
Ayyash, "Internet of Things: A Survey on Enabling Technologies,
Protocols, and Applications," \emph{IEEE Communications Surveys \&
Tutorials}, vol. 17, no. 4, pp. 2347--2376, Fourth Quarter 2015.

{[}6{]} N. Naik, "Choice of Effective Messaging Protocols for IoT
Systems: MQTT, CoAP, AMQP and HTTP," in \emph{Proc. 2017 IEEE
International Systems Engineering Symposium (ISSE)}, Vienna, Austria,
2017, pp. 1--7.

{[}7{]} M. Quigley, R. Johnson, and P. Sharma, "Power Management
Strategies for IoT Vehicle Tracking Devices," \emph{IEEE Internet of
Things Journal}, vol. 8, no. 15, pp. 12045--12058, Aug. 2021.

{[}8{]} M. A. Al-Khedher, "Hybrid GPS-GSM Localization of Automobile
Tracking System," \emph{International Journal of Computer Science and
Information Technology (IJCSIT)}, vol. 3, no. 6, pp. 75--85, Dec. 2011.

{[}9{]} S. Kaur and S. Singh, "A Survey on IoT based Vehicle Tracking
System," \emph{International Journal of Advanced Research in Computer
Science}, vol. 10, no. 3, pp. 42--47, 2019.

{[}10{]} R. Barry, \emph{Mastering the FreeRTOS Real Time Kernel: A
Hands-On Tutorial Guide}. Real Time Engineers Ltd., 2016.

{[}11{]} B. Boehm and V. R. Basili, "Software Defect Reduction Top 10
List," \emph{IEEE Computer}, vol. 34, no. 1, pp. 135--137, Jan. 2001.

{[}12{]} E. Evans, \emph{Domain-Driven Design: Tackling Complexity in
the Heart of Software}. Boston, MA, USA: Addison-Wesley Professional,
2003.

{[}13{]} S. Newman, \emph{Building Microservices: Designing Fine-Grained
Systems}, 2nd ed. Sebastopol, CA, USA: O\textquotesingle Reilly Media,
2021.

{[}14{]} D. Guinard and V. Trifa, \emph{Building the Web of Things}.
Shelter Island, NY, USA: Manning Publications, 2016.

{[}15{]} J. Gubbi, R. Buyya, S. Marusic, and M. Palaniswami, "Internet
of Things (IoT): A Vision, Architectural Elements, and Future
Directions," \emph{Future Generation Computer Systems}, vol. 29, no. 7,
pp. 1645--1660, Sep. 2013.

{[}16{]} V. C. Gungor and G. P. Hancke, "Industrial Wireless Sensor
Networks: Challenges, Design Principles, and Technical Approaches,"
\emph{IEEE Transactions on Industrial Electronics}, vol. 56, no. 10, pp.
4258--4265, Oct. 2009.

{[}17{]} R. K. Kodali, V. Jain, S. Bose, and L. Boppana, "IoT Based
Smart Security and Home Automation System," in \emph{Proc. 2016
International Conference on Computing, Communication and Automation
(ICCCA)}, Greater Noida, India, 2016, pp. 1286--1289.

{[}18{]} B. Sahu and G. A. Rincon-Mora, "A Low Voltage, Dynamic,
Noninverting, Synchronous Buck-Boost Converter for Portable
Applications," \emph{IEEE Transactions on Power Electronics}, vol. 19,
no. 2, pp. 443--452, Mar. 2004.

\begin{center}\rule{0.5\linewidth}{0.5pt}\end{center}

\subsection{{[}C{]} Tài liệu kỹ thuật (Datasheets \& Technical
Documents)}\label{c-tuxe0i-liux1ec7u-kux1ef9-thuux1eadt-datasheets--technical-documents}

{[}19{]} Espressif Systems, "ESP32-S3 Series Datasheet," Version 1.3,
2023. {[}Online{]}. Available:
\url{https://www.espressif.com/sites/default/files/documentation/esp32-s3_datasheet_en.pdf}

{[}20{]} Espressif Systems, "ESP32-S3 Technical Reference Manual,"
Version 1.1, 2023. {[}Online{]}. Available:
\url{https://www.espressif.com/sites/default/files/documentation/esp32-s3_technical_reference_manual_en.pdf}

{[}21{]} SIMCom Wireless Solutions, "SIM7600 Series AT Command Manual,"
Version 1.06, 2023. {[}Online{]}. Available:
\url{https://www.simcom.com/product/SIM7600CE-T.html}

{[}22{]} SIMCom Wireless Solutions, "SIM7600CE-T Module Hardware Design
Guide," Version 1.02, 2022.

{[}23{]} STMicroelectronics, "LIS3DSH MEMS digital output motion
Sensor Ultra-Low-Power High-Performance 3-Axis
\textquotesingle Nano\textquotesingle{} Accelerometer Datasheet," DocID
17530, Rev. 3, 2021. {[}Online{]}. Available:
\url{https://www.st.com/resource/en/datasheet/lis3dsh.pdf}

{[}24{]} STMicroelectronics, "AN3393: LIS3DSH 3-axis digital output accelerometer." {[}Online{]}. Available:
\url{https://www.st.com/resource/en/application_note/dm00026768-lis3dsh-3axis-digital-output-accelerometer-stmicroelectronics.pdf}

{[}25{]} OASIS, "MQTT Version 5.0 - OASIS Standard," Mar. 2019.
{[}Online{]}. Available:
\url{https://docs.oasis-open.org/mqtt/mqtt/v5.0/mqtt-v5.0.html}

{[}26{]} SAE International, "SAE J1979 - E/E Diagnostic Test Modes,"
Standard, Rev. 2014.

{[}27{]} International Organization for Standardization, "ISO
15031--5:2015 - Road vehicles - Communication between vehicle and
external equipment for emissions-related diagnostics," ISO Standard,
2015.

{[}28{]} International Organization for Standardization, "ISO
15765--2:2016 - Road vehicles - Diagnostic communication over Controller
Area Network (DoCAN) - Part 2: Transport protocol and network layer
services," ISO Standard, 2016.

{[}29{]} vgate, "iCar Pro BLE OBD2 Adapter Specifications," 2023.
{[}Online{]}. Available:
\url{https://www.vgatemall.com/product/vgate-icar-pro.html}

{[}30{]} NanJing Top Power ASIC Corp., "TP5100 2A Switching Buck Charger
IC for 8.4V/4.2V Lithium-Ion Battery," Datasheet, Rev. 2.4. {[}Online{]}. Available:
\url{https://www.toppwr.com/uploadfile/file/20230304/6402f53e7e8d9.pdf}. {[}Accessed: Apr. 18, 2026{]}.

{[}31{]} Apache NimBLE Project, "NimBLE Host API Reference," Version
1.5, 2023. {[}Online{]}. Available:
\url{https://mynewt.apache.org/latest/network/}

\begin{center}\rule{0.5\linewidth}{0.5pt}\end{center}

\subsection{{[}D{]} Tài liệu trực tuyến (Online
Resources)}\label{d-tuxe0i-liux1ec7u-trux1ef1c-tuyux1ebfn-online-resources}

{[}32{]} EMQX Team, "EMQX Documentation - The World\textquotesingle s
Most Scalable MQTT Broker," Version 5.x, 2024. {[}Online{]}. Available:
\url{https://docs.emqx.com/en/emqx/latest/}. {[}Accessed: Jan. 15,
2026{]}.

{[}33{]} The PostgreSQL Global Development Group, "PostgreSQL 16
Documentation," 2024. {[}Online{]}. Available:
\url{https://www.postgresql.org/docs/16/}. {[}Accessed: Jan. 10,
2026{]}.

{[}34{]} VictoriaMetrics Team, "VictoriaMetrics Documentation," 2024.
{[}Online{]}. Available: \url{https://docs.victoriametrics.com/}.
{[}Accessed: Jan. 10, 2026{]}.

{[}35{]} VictoriaMetrics Team, "VictoriaLogs Documentation," 2024.
{[}Online{]}. Available:
\url{https://docs.victoriametrics.com/victorialogs/}. {[}Accessed: Jan.
10, 2026{]}.

{[}36{]} OpenJS Foundation, "Express.js - Fast, Unopinionated,
Minimalist Web Framework for Node.js," 2024. {[}Online{]}. Available:
\url{https://expressjs.com/}. {[}Accessed: Dec. 20, 2025{]}.

{[}37{]} Vercel Inc., "Next.js 15 Documentation," 2024. {[}Online{]}.
Available: \url{https://nextjs.org/docs}. {[}Accessed: Dec. 15, 2025{]}.

{[}38{]} Meta Platforms, Inc., "React 19 Documentation," 2024.
{[}Online{]}. Available: \url{https://react.dev/}. {[}Accessed: Dec. 15,
2025{]}.

{[}39{]} Socket.IO Contributors, "Socket.IO Documentation," Version 4.8,
2024. {[}Online{]}. Available: \url{https://socket.io/docs/v4/}.
{[}Accessed: Dec. 20, 2025{]}.

{[}40{]} V. Agafonkin, "Leaflet - An Open-Source JavaScript Library for
Mobile-Friendly Interactive Maps," Version 1.9, 2024. {[}Online{]}.
Available: \url{https://leafletjs.com/}. {[}Accessed: Jan. 05, 2026{]}.

{[}41{]} Apache Software Foundation, "Apache ECharts - An Open Source
JavaScript Visualization Library," Version 5.5, 2024. {[}Online{]}.
Available: \url{https://echarts.apache.org/}. {[}Accessed: Jan. 05,
2026{]}.

{[}42{]} Docker Inc., "Docker Documentation," 2024. {[}Online{]}.
Available: \url{https://docs.docker.com/}. {[}Accessed: Nov. 15,
2025{]}.

{[}43{]} Docker Inc., "Docker Compose Documentation," 2024.
{[}Online{]}. Available: \url{https://docs.docker.com/compose/}.
{[}Accessed: Nov. 15, 2025{]}.

{[}44{]} Tailwind Labs, "Tailwind CSS Documentation," Version 4, 2025.
{[}Online{]}. Available: \url{https://tailwindcss.com/docs}.
{[}Accessed: Jan. 10, 2026{]}.

{[}45{]} Daishi Kato, "Zustand - Bear Necessities for State Management
in React," Version 5, 2024. {[}Online{]}. Available:
\url{https://zustand-demo.pmnd.rs/}. {[}Accessed: Dec. 20, 2025{]}.

{[}46{]} TanStack, "TanStack Query (React Query) Documentation," Version
5, 2024. {[}Online{]}. Available:
\url{https://tanstack.com/query/latest}. {[}Accessed: Dec. 20, 2025{]}.

{[}47{]} Espressif Systems, "ESP-IDF Programming Guide," Version 5.x,
2024. {[}Online{]}. Available:
\url{https://docs.espressif.com/projects/esp-idf/en/latest/esp32s3/}.
{[}Accessed: Oct. 10, 2025{]}.

{[}48{]} Grafana Labs, "Grafana Documentation," 2024. {[}Online{]}.
Available: \url{https://grafana.com/docs/grafana/latest/}. {[}Accessed:
Jan. 15, 2026{]}.

{[}49{]} Prometheus Authors, "Prometheus Documentation," 2024.
{[}Online{]}. Available: \url{https://prometheus.io/docs/}. {[}Accessed:
Jan. 15, 2026{]}.

{[}50{]} Radix UI Contributors, "Radix UI - Unstyled, Accessible
Components for React," 2024. {[}Online{]}. Available:
\url{https://www.radix-ui.com/}. {[}Accessed: Dec. 15, 2025{]}.

{[}51{]} Google, "Flutter Documentation - WebView Plugin," 2024.
{[}Online{]}. Available: \url{https://docs.flutter.dev/}. {[}Accessed:
Jan. 20, 2026{]}.

{[}52{]} Texas Instruments, "Power Management for IoT Devices,"
Application Report SLVAE57, 2020. {[}Online{]}. Available:
\url{https://www.ti.com/lit/an/slvae57/slvae57.pdf}. {[}Accessed: Sep.
15, 2025{]}.

{[}53{]} STMicroelectronics, "STM32L476RE Product Page," 2026.
{[}Online{]}. Available:
\url{https://www.st.com/en/microcontrollers-microprocessors/stm32l476re.html}.
{[}Accessed: Mar. 02, 2026{]}.

{[}54{]} STMicroelectronics, "STM32L476xx Datasheet," DS10198 Rev 11,
Jul. 2024. {[}Online{]}. Available:
\url{https://www.st.com/resource/en/datasheet/stm32l476re.pdf}.
{[}Accessed: Mar. 02, 2026{]}.

{[}55{]} Nordic Semiconductor, "nRF52840 Product Specification v1.2,"
{[}Online{]}. Available:
\url{https://docs.nordicsemi.com/bundle/nRF52840_PS_v1.2/resource/nRF52840_PS_v1.2.pdf}.
{[}Accessed: Mar. 02, 2026{]}.

{[}56{]} Nordic Semiconductor, "nRF52840 Product Page," 2026.
{[}Online{]}. Available:
\url{https://www.nordicsemi.com/Products/nRF52840/Modules}. {[}Accessed:
Mar. 02, 2026{]}.

{[}57{]} Nordic Semiconductor, "nRF52840 Documentation Hub," 2026.
{[}Online{]}. Available:
\url{https://docs.nordicsemi.com/category/nrf52840-category}.
{[}Accessed: Mar. 02, 2026{]}.

{[}58{]} SIMCom Wireless Solutions, "A7670X Series Product Page," 2026.
{[}Online{]}. Available:
\url{https://cn.simcom.com/product/A7670X.html}. {[}Accessed: Mar. 02,
2026{]}.

{[}59{]} Quectel Wireless Solutions, "EC200U Series Product Page," 2026.
{[}Online{]}. Available:
\url{https://www.quectel.com/product/lte-ec200u-series}. {[}Accessed:
Mar. 02, 2026{]}.

{[}60{]} SIMCom Wireless Solutions, "SIM7600CE Product Page," 2026.
{[}Online{]}. Available:
\url{https://en.simcom.com/product/SIM7600CE.html}. {[}Accessed: Mar.
13, 2026{]}.

{[}61{]} Quectel Wireless Solutions, "Quectel EC200U Series LTE Standard
Specification," Version 1.4, 2024. {[}Online{]}. Available:
\url{https://developer.quectel.com/en/wp-content/uploads/sites/2/2024/11/Quectel_EC200U_Series_LTE_Standard_Specification_V1.4.pdf}.
{[}Accessed: Mar. 02, 2026{]}.

{[}62{]} u-blox, "NEO-M8 Data Sheet," UBX-15031086, {[}Online{]}.
Available:
\url{https://www.u-blox.com/sites/default/files/NEO-M8-FW3_DataSheet_UBX-15031086.pdf}.
{[}Accessed: Mar. 02, 2026{]}.

{[}63{]} u-blox, "NEO-M8 Series Product Page," 2026. {[}Online{]}.
Available: \url{https://www.u-blox.com/en/product/neo-m8-series}.
{[}Accessed: Mar. 02, 2026{]}.

{[}64{]} vgate, "Vgate iCar Pro BLE Product Page," 2026. {[}Online{]}.
Available: \url{https://vgatemall.com/products-detail/i-9/}.
{[}Accessed: Mar. 02, 2026{]}.

{[}65{]} Texas Instruments, "TPS2115A Product Page," 2026. {[}Online{]}.
Available: \url{https://www.ti.com/product/TPS2115A}. {[}Accessed: Mar.
02, 2026{]}.

{[}66{]} Infineon Technologies, "IRLML6402 Product Page," 2026.
{[}Online{]}. Available: \url{https://www.infineon.com/part/IRLML6402}.
{[}Accessed: Mar. 02, 2026{]}.

{[}67{]} SONGLE Relay, "SRD-05VDC-SL-C Relay Datasheet," datasheet
mirror. {[}Online{]}. Available:
\url{https://www.alldatasheet.com/html-pdf/1132639/SONGLERELAY/SRD-05VDC-SL-C/1713/2/SRD-05VDC-SL-C.html}.
{[}Accessed: Mar. 02, 2026{]}.

\begin{center}\rule{0.5\linewidth}{0.5pt}\end{center}

\majorheading{PHỤ LỤC - APPENDICES}

\begin{center}\rule{0.5\linewidth}{0.5pt}\end{center}

\appendixsection{PHỤ LỤC 1: BÁO CÁO TÀI CHÍNH - FINANCE
REPORT}\label{phux1ee5-lux1ee5c-1-buxe1o-cuxe1o-tuxe0i-chuxednh---finance-report}

\subsection{1.1. Bảng kê chi phí linh kiện (Bill of Materials -
BOM)}\label{11-bux1ea3ng-kuxea-chi-phuxed-linh-kiux1ec7n-bill-of-materials---bom}

Bảng dưới đây liệt kê chi tiết các linh kiện chính sử dụng trong thiết
bị theo dõi IoT, bao gồm đơn giá và tổng chi phí ước tính. Đơn giá được
tham chiếu theo mặt bằng giá thị trường Việt Nam đã được rà soát lại vào
ngày 12/04/2026 từ các kênh bán lẻ như Điện tử Á Châu, Điện tử Đức Huy
và một số sàn thương mại điện tử trong nước đối với các linh kiện khó
kiếm.


{\thesissettableid{PL-1.1}{tab-app1-bang-ke-chi-phi-linh-kien-bom}
\def\LTcaptype{table}
\begin{longtable}{|L{0.055\linewidth}|L{0.130\linewidth}|L{0.185\linewidth}|L{0.090\linewidth}|L{0.105\linewidth}|L{0.125\linewidth}|L{0.235\linewidth}|}
\caption{Bảng kê chi phí linh kiện (BOM)}\label{tab:app1-bang-ke-chi-phi-linh-kien-bom}\\
\hline \textbf{STT} & \textbf{Linh kiện} & \textbf{Model / Thông số} & \textbf{Số lượng} & \textbf{Đơn giá (VND)} & \textbf{Thành
tiền (VND)} & \textbf{Ghi chú} \\
\\\hline
\endfirsthead
\caption[]{Bảng kê chi phí linh kiện (BOM) (tiếp theo)}\\
\hline \textbf{STT} & \textbf{Linh kiện} & \textbf{Model / Thông số} & \textbf{Số lượng} & \textbf{Đơn giá (VND)} & \textbf{Thành
tiền (VND)} & \textbf{Ghi chú} \\
\\\hline
\endhead
\\\hline
\endfoot
\endlastfoot
1 & Vi điều khiển & ESP32-S3-WROOM-1 (N16R8) & 1 & 95.000 & 95.000 &
MCU chính, 16MB Flash, 8MB PSRAM \\\hline
2 & LTE + GNSS & \nolinkurl{SIM7600CE-T} & 1 bá»™ & 650.000 & 650.000 & Modem
LTE Cat-4 tích hợp GPS/GNSS \\\hline
3 & OBD2 Adapter & vgate iCar Pro BLE & 1 & 367.000 & 367.000 &
Bluetooth Low Energy OBD2 \\\hline
4 & Cảm biến gia tốc & IC LIS3DSH & 1 & 40.000 & 40.000 & IMU 3 trục,
phát hiện chuyển động \\\hline
5 & Pin dự phòng & 18650 1S Li-ion 3500mAh & 1 & 79.000 & 79.000 &
Samsung/LG cell \\\hline
6 & IC sạc pin & TP5100 + mạch phụ trợ & 1 & 37.000 & 37.000 & Sạc 1S
+ bảo vệ pin, input 5V từ MP2482 \\\hline
7 & LDO 3.3V & Mạch AP2112-3.3 & 1 & 5.000 & 5.000 & 5V
-\textgreater{} 3.3V cấp ESP32-S3 \\\hline
8 & Buck 5V & Mạch MP2482 (12V/24V-\textgreater5V) & 1 & 17.000 & 17.000
& Tạo bus 5V chính \\\hline
9 & Buck 3.8V/4V & Mạch TPS54231 & 1 & 9.000 & 9.000 & 12--24V
-\textgreater{} ~4V cấp modem SIM7600CE-T \\\hline
10 & Boost converter & Mạch SX1308 (3.7V-\textgreater5V) & 1 & 12.000 &
12.000 & Tăng áp từ pin dự phòng \\\hline
11 & Diode OR & SS54/SS34 (Schottky) & 2 & 2.500 & 5.000 & OR nguồn giữa
nhánh 5V chính và dự phòng \\\hline
12 & Khối LVD & Khối giám sát LVD + ADC & 1 & 8.000 & 8.000 & Cho mạch Low Voltage
Disconnect \\\hline
13 & Anten GPS & Anten gốm GNSS 25x25mm & 1 & 25.000 & 25.000 & Anten
GPS, GLONASS, BeiDou \\\hline
14 & Anten 4G & Anten FPC 4G LTE & 1 & 20.000 & 20.000 & Anten mạng di
động \\\hline
15 & SIM tray + SIM & Nano SIM holder + SIM 4G & 1 & 15.000 & 15.000 &
SIM data 4G \\\hline
16 & Connector OBD2 & Jack OBD2 16-pin male & 1 & 35.000 & 35.000 & Kết
nối nguồn ắc quy xe (12V hoặc 24V) \\\hline
17 & Tụ điện, điện trở & Linh kiện thụ động (combo) & 1 bộ & 30.000 &
30.000 & Tụ lọc, điện trở chia áp ADC (R1=100k, R2=10k), LED \\\hline
18 & PCB chính & PCB 2 lớp theo thiết kế của đề tài ~50x50mm & 1
& 20.000 & 20.000 & Bo mạch chính của thiết bị \\\hline
19 & Hộp đựng & Hộp nhựa ABS 120 x 80 x 40 mm & 1 & 25.000 & 25.000 & Vỏ bảo
vệ thiết bị \\\hline
20 & Dây kết nối & Dây nối, header, jumper & 1 bộ & 20.000 & 20.000 &
Dây kết nối nội bộ \\\hline
& & & & \textbf{Tổng cộng} & \textbf{1.514.000} & \\\hline
\end{longtable}
}


\subsection{1.2. Chi phí hạ tầng cloud (ước tính hàng
tháng)}\label{12-chi-phuxed-hux1ea1-tux1ea7ng-cloud-ux1b0ux1edbc-tuxednh-huxe0ng-thuxe1ng}


{\thesissettableid{PL-1.2}{tab-app1-chi-phi-ha-tang-cloud}
\def\LTcaptype{table}
\begin{longtable}{|L{0.085\linewidth}|L{0.222\linewidth}|L{0.187\linewidth}|L{0.226\linewidth}|L{0.225\linewidth}|}
\caption{Chi phí hạ tầng cloud}\label{tab:app1-chi-phi-ha-tang-cloud}\\
\hline \textbf{STT} & \textbf{Hạng mục} & \textbf{Thông số} & \textbf{Chi phí hàng tháng (VND)} & \textbf{Ghi chú} \\
\\\hline
\endfirsthead
\caption[]{Chi phí hạ tầng cloud (tiếp theo)}\\
\hline \textbf{STT} & \textbf{Hạng mục} & \textbf{Thông số} & \textbf{Chi phí hàng tháng (VND)} & \textbf{Ghi chú} \\
\\\hline
\endhead
\\\hline
\endfoot
\endlastfoot
1 & VPS Server & 4 vCPU, 8GB RAM, 100GB SSD & 300.000--500.000 &
DigitalOcean / Vultr / Viettel IDC \\\hline
2 & SIM 4G data & Gói cước data IoT & 70.000 & Mỗi thiết bị 1 SIM \\\hline
3 & Tên miền & Domain .com & 25.000 & ~300.000 VND/năm \\\hline
4 & SSL Certificate & Let\textquotesingle s Encrypt & 0 & Miễn phí, tự
động gia hạn \\\hline
& & & \textbf{Tổng hàng tháng} & \textbf{395.000--595.000} \\\hline
\end{longtable}
}


\subsection{1.3. Tổng hợp chi phí dự
án}\label{13-tux1ed5ng-hux1ee3p-chi-phuxed-dux1ef1-uxe1n}


{\thesissettableid{PL-1.3}{tab-app1-tong-hop-chi-phi-du-an}
\def\LTcaptype{table}
\begin{longtable}{|L{0.368\linewidth}|L{0.296\linewidth}|L{0.300\linewidth}|}
\caption{Tổng hợp chi phí dự án}\label{tab:app1-tong-hop-chi-phi-du-an}\\
\hline \textbf{Hạng mục} & \textbf{Chi phí (VND)} & \textbf{Ghi chú} \\
\\\hline
\endfirsthead
\caption[]{Tổng hợp chi phí dự án (tiếp theo)}\\
\hline \textbf{Hạng mục} & \textbf{Chi phí (VND)} & \textbf{Ghi chú} \\
\\\hline
\endhead
\\\hline
\endfoot
\endlastfoot
Phần cứng (1 bộ thiết bị theo dõi) & 1.514.000 & Theo BOM ở bảng PL-1.1 \\\hline
Phần mềm (licenses) & 0 & Toàn bộ mã nguồn mở \\\hline
Hạ tầng cloud (3 tháng dev + 3 tháng test) & 2.370.000--3.570.000 & 6
tháng x chi phí hàng tháng \\\hline
SIM 4G (6 tháng) & 420.000 & 70.000 x 6 tháng \\\hline
Công cụ phát triển (khác) & 200.000 & USB-UART, mỏ hàn, dây đo \\\hline
\textbf{Tổng chi phí dự án} & \textbf{4.504.000--5.704.000} & \\\hline
\end{longtable}
}


\textbf{Nhận xét:} Tổng chi phí dự án vẫn dưới 6.000.000 VND, trong đó
chi phí phần cứng cho mỗi bộ thiết bị theo dõi khoảng 1.514.000 VND, thấp hơn đáng
kể so với các giải pháp thương mại tương đương (2.000.000--5.000.000
VND/thiết bị). Khi sản xuất số lượng lớn (\textgreater50 bộ), chi phí
linh kiện có thể giảm thêm 15--25\% nhờ mua sỉ.

\begin{center}\rule{0.5\linewidth}{0.5pt}\end{center}

\appendixsection{PHỤ LỤC 2: CÁC TIÊU CHUẨN THIẾT KẾ -
STANDARDS}\label{phux1ee5-lux1ee5c-2-cuxe1c-tiuxeau-chuux1ea9n-thiux1ebft-kux1ebf---standards}

\subsection{2.1. Tiêu chuẩn OBD2 (SAE J1979 / ISO
15031--5)}\label{21-tiuxeau-chuux1ea9n-obd2-sae-j1979--iso-150315}

\textbf{Mô tả:} On-Board Diagnostics II (OBD2) là tiêu chuẩn chẩn đoán
xe được bắt buộc áp dụng cho tất cả các xe sản xuất từ năm 1996 (tại Mỹ)
và 2001 (tại Châu Âu). Tiêu chuẩn này quy định giao diện vật lý
(connector 16-pin), giao thức truyền thông, và các thông số chẩn đoán
(PIDs - Parameter IDs).

\textbf{Áp dụng trong dự án:}


{\thesissettableid{2.1A}{tab-ch2-tong-hop-thong-tin-ve-chuan-cu-the-noi-dung-ap-dung}
\def\LTcaptype{table}
\begin{longtable}{|L{0.161\linewidth}|L{0.394\linewidth}|L{0.409\linewidth}|}
\caption{Tổng hợp thông tin về Chuẩn cụ thể, Nội dung, Áp dụng}\label{tab:ch2-tong-hop-thong-tin-ve-chuan-cu-the-noi-dung-ap-dung}\\
\hline \textbf{Chuẩn cụ thể} & \textbf{Nội dung} & \textbf{Áp dụng} \\
\\\hline
\endfirsthead
\caption[]{Tổng hợp thông tin về Chuẩn cụ thể, Nội dung, Áp dụng (tiếp theo)}\\
\hline \textbf{Chuẩn cụ thể} & \textbf{Nội dung} & \textbf{Áp dụng} \\
\\\hline
\endhead
\\\hline
\endfoot
\endlastfoot
SAE J1979 & Định nghĩa các Diagnostic Test Modes (Mode 01--0A) & Đọc dữ
liệu động cơ: RPM, tốc độ, nhiệt độ, nhiên liệu \\\hline
ISO 15031--5 & Định nghĩa PIDs và định dạng dữ liệu & Phân tích
(parsing) bản tin OBD2 response \\\hline
ISO 15765--2 & Transport protocol cho CAN bus (ISO-TP) & Xử lý
multi-frame response (VIN, DTC list) \\\hline
ISO 15765--4 & Emissions-related OBD via CAN & Giao tiếp CAN qua OBD2
adapter BLE \\\hline
\end{longtable}
}


\textbf{Các OBD2 PIDs sử dụng trong dự án:}


{\thesissettableid{2.1B}{tab-ch2-tong-hop-thong-tin-ve-pid-hex-mo-ta-on-vi}
\def\LTcaptype{table}
\begin{longtable}{|L{0.203\linewidth}|L{0.460\linewidth}|L{0.145\linewidth}|L{0.146\linewidth}|}
\caption{Tổng hợp thông tin về PID (hex), Mô tả, Đơn vị}\label{tab:ch2-tong-hop-thong-tin-ve-pid-hex-mo-ta-on-vi}\\
\hline \textbf{PID (hex)} & \textbf{Mô tả} & \textbf{Đơn vị} & \textbf{Mode} \\
\\\hline
\endfirsthead
\caption[]{Tổng hợp thông tin về PID (hex), Mô tả, Đơn vị (tiếp theo)}\\
\hline \textbf{PID (hex)} & \textbf{Mô tả} & \textbf{Đơn vị} & \textbf{Mode} \\
\\\hline
\endhead
\\\hline
\endfoot
\endlastfoot
0x00 & Supported PIDs {[}01--20{]} & Bitmask & 01 \\\hline
0x04 & Calculated Engine Load & \% & 01 \\\hline
0x05 & Engine Coolant Temperature & C & 01 \\\hline
0x0C & Engine RPM & rpm & 01 \\\hline
0x0D & Vehicle Speed & km/h & 01 \\\hline
0x0F & Intake Air Temperature & C & 01 \\\hline
0x11 & Throttle Position & \% & 01 \\\hline
0x2F & Fuel Tank Level Input & \% & 01 \\\hline
0x46 & Ambient Air Temperature & C & 01 \\\hline
\end{longtable}
}


\subsection{2.2. MQTT 3.1.1 trong triển khai
thiết bị}\label{22-tiuxeau-chuux1ea9n-mqtt-50-oasis}

\textbf{Mô tả:} MQTT vẫn là giao thức nhắn tin phù hợp cho IoT vì gọn, ít
overhead và hoạt động tốt trên kết nối 4G. Trong hệ thống này, EMQX chạy
phiên bản 5.x ở vai trò broker, nhưng firmware trên thiết bị hiện cấu
hình \textbf{MQTT 3.1.1}. Vì vậy, phần trình bày dưới đây tập trung vào
các đặc tính thực sự đang dùng trong đường truyền thiết bị thay vì liệt
kê các khả năng mở rộng của MQTT 5.

\textbf{Các đặc tính MQTT đang dùng trong dự án:}


{\thesissettableid{2.2A}{tab-ch2-tong-hop-thong-tin-ve-tinh-nang-mo-ta-ap-dung}
\def\LTcaptype{table}
\begin{longtable}{|L{0.281\linewidth}|L{0.333\linewidth}|L{0.350\linewidth}|}
\caption{Tổng hợp thông tin về Đặc tính, Mô tả, Áp dụng thực tế}\label{tab:ch2-tong-hop-thong-tin-ve-tinh-nang-mo-ta-ap-dung}\\
\hline \textbf{Đặc tính} & \textbf{Mô tả} & \textbf{Áp dụng thực tế} \\
\\\hline
\endfirsthead
\caption[]{Tổng hợp thông tin về Đặc tính, Mô tả, Áp dụng thực tế (tiếp theo)}\\
\hline \textbf{Đặc tính} & \textbf{Mô tả} & \textbf{Áp dụng thực tế} \\
\\\hline
\endhead
\\\hline
\endfoot
\endlastfoot
QoS 0 & Gửi không cần ACK & \texttt{rawdata} tần suất cao \\\hline
QoS 1 & Gửi có ACK, tránh mất bản tin quan trọng & \texttt{status},
\texttt{events}, \texttt{firmware}, \texttt{commands} \\\hline
Topic phân phiên bản & Chuẩn hóa theo cây \texttt{v1/\{device\_id\}/...} &
Thiết bị publish và backend gửi lệnh \\\hline
Topic lệnh hai chiều & Backend publish trực tiếp vào
\texttt{v1/\{device\_id\}/commands} & Điều khiển từ xa, OTA \\\hline
Internal event bus & MQTT Bridge publish \texttt{internal/events/\#}
qua EMQX & Backend listener -\textgreater{} Socket.IO realtime \\\hline
ACL theo thiết bị & Mỗi thiết bị chỉ được phép dùng topic của chính nó &
Giảm rủi ro giả mạo publish/subscribe \\\hline
\end{longtable}
}


\textbf{Cấu trúc topic MQTT:}

\begin{itemize}
\tightlist
\item
  \texttt{v1/\{device\_id\}/rawdata}: dữ liệu telemetry chính (QoS 0).
\item
  \texttt{v1/\{device\_id\}/events}: sự kiện và cảnh báo (QoS 1).
\item
  \texttt{v1/\{device\_id\}/commands}: lệnh điều khiển từ server (QoS 1).
\item
  \texttt{v1/\{device\_id\}/status}: trạng thái phiên/heartbeat (QoS 1).
\item
  \texttt{v1/\{device\_id\}/firmware}: tiến trình OTA (QoS 1).
\item
  \texttt{internal/events/device/status}: sự kiện nội bộ cho backend
  realtime.
\item
  \texttt{internal/events/device/alert}: cảnh báo nội bộ.
\item
  \texttt{internal/events/device/session}: trạng thái phiên nội bộ.
\item
  \texttt{internal/events/device/data}: vị trí/telemetry cho realtime.
\end{itemize}

\subsection{2.3. Bảo mật thông tin (ISO 27001 - tham
khảo)}\label{23-bux1ea3o-mux1eadt-thuxf4ng-tin-iso-27001---tham-khux1ea3o}

\textbf{Mô tả:} ISO 27001 là tiêu chuẩn quốc tế về hệ thống quản lý an
toàn thông tin (ISMS). Dự án tham khảo các nguyên tắc của ISO 27001 để
thiết kế các biện pháp bảo mật.

\textbf{Các nguyên tắc bảo mật áp dụng:}


{\thesissettableid{2.3A}{tab-ch2-tong-hop-thong-tin-ve-nguyen-tac-hien-thuc-trong-du-an}
\def\LTcaptype{table}
\begin{longtable}{|L{0.382\linewidth}|L{0.586\linewidth}|}
\caption{Tổng hợp thông tin về Nguyên tắc, Hiện thực trong dự án}\label{tab:ch2-tong-hop-thong-tin-ve-nguyen-tac-hien-thuc-trong-du-an}\\
\hline \textbf{Nguyên tắc} & \textbf{Hiện thực trong dự án} \\
\\\hline
\endfirsthead
\caption[]{Tổng hợp thông tin về Nguyên tắc, Hiện thực trong dự án (tiếp theo)}\\
\hline \textbf{Nguyên tắc} & \textbf{Hiện thực trong dự án} \\
\\\hline
\endhead
\\\hline
\endfoot
\endlastfoot
Confidentiality (Bảo mật) & Mã hóa TLS cho MQTT và HTTPS; Token hash
SHA-256 \\\hline
Integrity (Toàn vẹn) & MQTT checksum; Database constraints; Zod
validation \\\hline
Availability (Sẵn sàng) & Docker restart policy; Health checks; Pin dự
phòng \\\hline
Authentication (Xác thực) & Session-based auth với database-backed
tokens \\\hline
Authorization (Phân quyền) & RBAC (Role-Based Access Control); MQTT ACL
per device \\\hline
Audit Trail (Nhật ký) & VictoriaLogs ghi nhật ký mọi thao tác;
Liên kết truy vết theo Request-ID \\\hline
\end{longtable}
}


\subsection{2.4. Thiết kế REST API (RFC 7231 và best
practices)}\label{24-thiux1ebft-kux1ebf-rest-api-rfc-7231-vuxe0-best-practices}

\textbf{Các nguyên tắc REST API áp dụng:}


{\thesissettableid{2.4A}{tab-ch2-tong-hop-thong-tin-ve-nguyen-tac-mo-ta-vi-du-trong}
\def\LTcaptype{table}
\begin{longtable}{|L{0.197\linewidth}|L{0.287\linewidth}|L{0.481\linewidth}|}
\caption{Tổng hợp thông tin về Nguyên tắc, Mô tả, Ví dụ trong dự án}\label{tab:ch2-tong-hop-thong-tin-ve-nguyen-tac-mo-ta-vi-du-trong}\\
\hline \textbf{Nguyên tắc} & \textbf{Mô tả} & \textbf{Ví dụ trong dự án} \\
\\\hline
\endfirsthead
\caption[]{Tổng hợp thông tin về Nguyên tắc, Mô tả, Ví dụ trong dự án (tiếp theo)}\\
\hline \textbf{Nguyên tắc} & \textbf{Mô tả} & \textbf{Ví dụ trong dự án} \\
\\\hline
\endhead
\\\hline
\endfoot
\endlastfoot
Resource-based URLs & URL đại diện cho tài nguyên, dùng số nhiều &
\nolinkurl{/api/v1/vehicles}, \nolinkurl{/api/v1/alerts} \\\hline
HTTP Methods & Dùng đúng ý nghĩa của phương thức & GET (đọc), POST
(tạo), PUT (cập nhật), DELETE (xóa) \\\hline
Status Codes & Mã trạng thái phản hồi chính xác & 200 OK, 201 Created,
400 Bad Request, 401 Unauthorized, 404 Not Found \\\hline
Pagination & Phân trang cho danh sách lớn &
\nolinkurl{?page=1&limit=20&sort=created_at:desc} \\\hline
Versioning & Phiên bản API trong URL & \texttt{/api/v1/...} \\\hline
Error Format & Định dạng lỗi nhất quán &
\nolinkurl{{"error":{"code":"...","message":"..."}}} \\\hline
Request-ID & Định danh mỗi request & Header \nolinkurl{X-Request-ID} cho
truy vết lỗi \\\hline
\end{longtable}
}


\begin{center}\rule{0.5\linewidth}{0.5pt}\end{center}

\appendixsection{PHỤ LỤC 3: KẾ HOẠCH THỰC HIỆN - ASSIGNMENT AND
TIMELINES}\label{phux1ee5-lux1ee5c-3-kux1ebf-houx1ea1ch-thux1ef1c-hiux1ec7n---assignment-and-timelines}

\subsection{3.1. Phân chia giai đoạn dự
án}\label{31-phuxe2n-chia-giai-ux111oux1ea1n-dux1ef1-uxe1n}

Dự án được triển khai theo 7 giai đoạn chính, với tổng thời gian dự kiến
24 tuần (6 tháng), từ tháng 09/2025 đến tháng 02/2026.


{\thesissettableid{PL-3.1}{tab-app3-phan-chia-giai-oan-va-thoi-gian-thuc-hien}
\def\LTcaptype{table}
\begin{longtable}{|L{0.115\linewidth}|L{0.247\linewidth}|L{0.184\linewidth}|L{0.408\linewidth}|}
\caption{Phân chia giai đoạn và thời gian thực hiện}\label{tab:app3-phan-chia-giai-oan-va-thoi-gian-thuc-hien}\\
\hline \textbf{Giai đoạn} & \textbf{Nội dung chính} & \textbf{Thời gian} & \textbf{Kết quả đầu ra} \\
\\\hline
\endfirsthead
\caption[]{Phân chia giai đoạn và thời gian thực hiện (tiếp theo)}\\
\hline \textbf{Giai đoạn} & \textbf{Nội dung chính} & \textbf{Thời gian} & \textbf{Kết quả đầu ra} \\
\\\hline
\endhead
\\\hline
\endfoot
\endlastfoot
\textbf{GĐ1} & Nghiên cứu và thiết kế phần cứng & Tuần 1--4 (4 tuần) &
Sơ đồ mạch, BOM, bố trí PCB \\\hline
\textbf{GĐ2} & Phát triển firmware & Tuần 3--10 (8 tuần) & Firmware
ESP-IDF hoàn chỉnh \\\hline
\textbf{GĐ3} & Xây dựng hạ tầng cloud & Tuần 5--8 (4 tuần) & Docker
infrastructure, EMQX, databases \\\hline
\textbf{GĐ4} & Phát triển backend API & Tuần 7--14 (8 tuần) & REST API,
MQTT Bridge, WebSocket \\\hline
\textbf{GĐ5} & Phát triển frontend & Tuần 11--18 (8 tuần) & Web
dashboard, bản đồ, biểu đồ \\\hline
\textbf{GĐ6} & Tích hợp và kiểm thử & Tuần 17--22 (6 tuần) & Hệ thống
tích hợp, báo cáo kiểm thử \\\hline
\textbf{GĐ7} & Viết báo cáo và bảo vệ & Tuần 21--24 (4 tuần) & Báo cáo
đồ án, slide thuyết trình \\\hline
\end{longtable}
}


\textbf{Ghi chú:} Các giai đoạn có sự chồng chéo (overlap) có chủ đích
để tối ưu hóa thời gian. Ví dụ, GĐ3 (hạ tầng cloud) bắt đầu trước khi
GĐ2 (firmware) hoàn thành để có môi trường test sớm.

\subsection{3.2. Kế hoạch thực hiện dự án theo giai đoạn}\label{32-kux1ebf-houx1ea1ch-thux1ef1c-hiux1ec7n-dux1ef1-uxe1n-theo-giai-ux111oux1ea1n}

Kế hoạch thực hiện dự án theo giai đoạn được trình bày trực quan tại hình
dưới đây, thể hiện
đầy đủ 7 giai đoạn và phần chồng chéo tiến độ trong 24 tuần thực hiện dự án.


{\thesissetfigureid{PL-3.1}{fig-app3-bieu-o-gantt-ke-hoach-thuc-hien-du-an-24-tuan}
\begin{figure}[htbp]
\centering
\pandocbounded{\safeincludesvg{./assets/figures/14-phu-luc-hinh-pl-3-1.svg}}
\caption{Kế hoạch thực hiện dự án theo giai đoạn (24 tuần)}
\label{fig:app3-bieu-o-gantt-ke-hoach-thuc-hien-du-an-24-tuan}
\end{figure}
}


\subsection{3.3. Chi tiết công việc theo giai
đoạn}\label{33-chi-tiux1ebft-cuxf4ng-viux1ec7c-theo-giai-ux111oux1ea1n}

\textbf{Giai đoạn 1 --- Nghiên cứu và thiết kế phần cứng (Tuần 1--4):}


{\thesissettableid{3.3A}{tab-ch3-tong-hop-thong-tin-ve-tuan-cong-viec-san-pham}
\def\LTcaptype{table}
\begin{longtable}{|L{0.160\linewidth}|L{0.366\linewidth}|L{0.439\linewidth}|}
\caption{Tổng hợp thông tin về Tuần, Công việc, Sản phẩm}\label{tab:ch3-tong-hop-thong-tin-ve-tuan-cong-viec-san-pham}\\
\hline \textbf{Tuần} & \textbf{Công việc} & \textbf{Sản phẩm} \\
\\\hline
\endfirsthead
\caption[]{Tổng hợp thông tin về Tuần, Công việc, Sản phẩm (tiếp theo)}\\
\hline \textbf{Tuần} & \textbf{Công việc} & \textbf{Sản phẩm} \\
\\\hline
\endhead
\\\hline
\endfoot
\endlastfoot
1 & Khảo sát linh kiện, đặt mua & Danh sách linh kiện \\\hline
2 & Thiết kế sơ đồ nguyên lý & Sơ đồ nguyên lý hoàn chỉnh \\\hline
3 & Thiết kế PCB và đặt làm mạch & Bố trí PCB / bo mạch hoàn chỉnh \\\hline
4 & Lắp ráp và kiểm tra cơ bản & Thiết bị phần cứng hoạt động \\\hline
\end{longtable}
}


\textbf{Giai đoạn 2 --- Phát triển firmware (Tuần 3--10):}


{\thesissettableid{3.3B}{tab-ch3-tong-hop-thong-tin-ve-tuan-cong-viec-san-pham-ca9bd3}
\def\LTcaptype{table}
\begin{longtable}{|L{0.160\linewidth}|L{0.462\linewidth}|L{0.342\linewidth}|}
\caption{Tổng hợp thông tin về Tuần, Công việc, Sản phẩm}\label{tab:ch3-tong-hop-thong-tin-ve-tuan-cong-viec-san-pham-ca9bd3}\\
\hline \textbf{Tuần} & \textbf{Công việc} & \textbf{Sản phẩm} \\
\\\hline
\endfirsthead
\caption[]{Tổng hợp thông tin về Tuần, Công việc, Sản phẩm (tiếp theo)}\\
\hline \textbf{Tuần} & \textbf{Công việc} & \textbf{Sản phẩm} \\
\\\hline
\endhead
\\\hline
\endfoot
\endlastfoot
3--4 & Cấu hình ESP-IDF, GPIO, UART cho modem & Giao tiếp modem cơ bản \\\hline
5--6 & Mô-đun BLE OBD2, phân tích bản tin & Đọc dữ liệu OBD2 \\\hline
7--8 & Mô-đun MQTT, tự kết nối lại, QoS & Truyền dữ liệu lên broker \\\hline
9--10 & Quản lý năng lượng, IMU, tích hợp & Firmware hoàn chỉnh \\\hline
\end{longtable}
}


\textbf{Giai đoạn 3 --- Hạ tầng cloud (Tuần 5--8):}


{\thesissettableid{3.3C}{tab-ch3-tong-hop-thong-tin-ve-tuan-cong-viec-san-pham-b7165b}
\def\LTcaptype{table}
\begin{longtable}{|L{0.160\linewidth}|L{0.474\linewidth}|L{0.330\linewidth}|}
\caption{Tổng hợp thông tin về Tuần, Công việc, Sản phẩm}\label{tab:ch3-tong-hop-thong-tin-ve-tuan-cong-viec-san-pham-b7165b}\\
\hline \textbf{Tuần} & \textbf{Công việc} & \textbf{Sản phẩm} \\
\\\hline
\endfirsthead
\caption[]{Tổng hợp thông tin về Tuần, Công việc, Sản phẩm (tiếp theo)}\\
\hline \textbf{Tuần} & \textbf{Công việc} & \textbf{Sản phẩm} \\
\\\hline
\endhead
\\\hline
\endfoot
\endlastfoot
5 & Thiết lập Docker, PostgreSQL, EMQX & Hạ tầng cơ bản \\\hline
6 & VictoriaMetrics, VictoriaLogs & Chuỗi thời gian và nhật ký \\\hline
7 & Dịch vụ MQTT Bridge & Dữ liệu từ EMQX đến DB \\\hline
8 & Grafana, giám sát, sao lưu & Bảng điều khiển giám sát \\\hline
\end{longtable}
}


\textbf{Giai đoạn 4 --- API phía máy chủ (Tuần 7--14):}


{\thesissettableid{3.3D}{tab-ch3-tong-hop-thong-tin-ve-tuan-cong-viec-san-pham-b63b40}
\def\LTcaptype{table}
\begin{longtable}{|L{0.160\linewidth}|L{0.524\linewidth}|L{0.280\linewidth}|}
\caption{Tổng hợp thông tin về Tuần, Công việc, Sản phẩm}\label{tab:ch3-tong-hop-thong-tin-ve-tuan-cong-viec-san-pham-b63b40}\\
\hline \textbf{Tuần} & \textbf{Công việc} & \textbf{Sản phẩm} \\
\\\hline
\endfirsthead
\caption[]{Tổng hợp thông tin về Tuần, Công việc, Sản phẩm (tiếp theo)}\\
\hline \textbf{Tuần} & \textbf{Công việc} & \textbf{Sản phẩm} \\
\\\hline
\endhead
\\\hline
\endfoot
\endlastfoot
7--8 & Khởi tạo dự án, mô-đun xác thực & Đăng nhập, phiên làm việc \\\hline
9--10 & Miền xe, thiết bị, khách hàng & API quản lý dữ liệu \\\hline
11--12 & Miền dữ liệu vận hành, cảnh báo, vùng cho phép & Dữ liệu thời gian thực \\\hline
13--14 & WebSocket, lệnh điều khiển, kiểm thử & API hoàn chỉnh \\\hline
\end{longtable}
}


\textbf{Giai đoạn 5 --- Frontend (Tuần 11--18):}


{\thesissettableid{3.3E}{tab-ch3-tong-hop-thong-tin-ve-tuan-cong-viec-san-pham-c5129f}
\def\LTcaptype{table}
\begin{longtable}{|L{0.160\linewidth}|L{0.516\linewidth}|L{0.289\linewidth}|}
\caption{Tổng hợp thông tin về Tuần, Công việc, Sản phẩm}\label{tab:ch3-tong-hop-thong-tin-ve-tuan-cong-viec-san-pham-c5129f}\\
\hline \textbf{Tuần} & \textbf{Công việc} & \textbf{Sản phẩm} \\
\\\hline
\endfirsthead
\caption[]{Tổng hợp thông tin về Tuần, Công việc, Sản phẩm (tiếp theo)}\\
\hline \textbf{Tuần} & \textbf{Công việc} & \textbf{Sản phẩm} \\
\\\hline
\endhead
\\\hline
\endfoot
\endlastfoot
11--12 & Khởi tạo dự án, bố cục, trang xác thực & Khung ứng dụng \\\hline
13--14 & Trang tổng quan, bản đồ Leaflet & Trang chủ, bản đồ \\\hline
15--16 & Quản lý xe, cảnh báo, vùng cho phép & Các trang quản lý \\\hline
17--18 & ECharts, WebSocket, hoàn thiện giao diện & Giao diện hoàn chỉnh \\\hline
\end{longtable}
}


\textbf{Giai đoạn 6 --- Tích hợp và kiểm thử (Tuần 17--22):}


{\thesissettableid{3.3F}{tab-ch3-tong-hop-thong-tin-ve-tuan-cong-viec-san-pham-af9e31}
\def\LTcaptype{table}
\begin{longtable}{|L{0.167\linewidth}|L{0.485\linewidth}|L{0.313\linewidth}|}
\caption{Tổng hợp thông tin về Tuần, Công việc, Sản phẩm}\label{tab:ch3-tong-hop-thong-tin-ve-tuan-cong-viec-san-pham-af9e31}\\
\hline \textbf{Tuần} & \textbf{Công việc} & \textbf{Sản phẩm} \\
\\\hline
\endfirsthead
\caption[]{Tổng hợp thông tin về Tuần, Công việc, Sản phẩm (tiếp theo)}\\
\hline \textbf{Tuần} & \textbf{Công việc} & \textbf{Sản phẩm} \\
\\\hline
\endhead
\\\hline
\endfoot
\endlastfoot
17--18 & Tích hợp firmware - cloud & Hệ thống end-to-end \\\hline
19--20 & Kiểm thử chức năng, hiệu năng & Báo cáo kiểm thử \\\hline
21--22 & Sửa lỗi, tối ưu, triển khai thử nghiệm & Hệ thống ổn định \\\hline
\end{longtable}
}


\begin{center}\rule{0.5\linewidth}{0.5pt}\end{center}

\appendixsection{PHỤ LỤC 4: TÀI LIỆU BỔ TRỢ - SUPPORTING
MATERIALS}\label{phux1ee5-lux1ee5c-4-tuxe0i-liux1ec7u-bux1ed5-trux1ee3---supporting-materials}

\subsection{4.1. Kho mã nguồn của
dự án}\label{41-muxe3-nguux1ed3n-dux1ef1-uxe1n-source-code-repository}

Toàn bộ mã nguồn phục vụ việc dựng lại và kiểm chứng hệ thống được lưu
trữ trên GitHub:

\begin{itemize}
\tightlist
\item
  \textbf{Địa chỉ kho mã nguồn:}
  \texttt{https://github.com/anmh1205/IoT\_Vehicle\_Tracking\_System}
\item
  \textbf{Nhánh tích hợp chính:} \texttt{main}
\item
  \textbf{Nhánh triển khai thử nghiệm:} \texttt{uat}
\end{itemize}

\textbf{Cấu trúc thư mục chính:}

\begin{itemize}
\tightlist
\item
  \texttt{iot-vehicle-tracking-system-cloud/}: chứa toàn bộ dịch vụ vận hành
  (\texttt{Tracking\_Backend}, \texttt{Tracking\_Frontend},
  \texttt{Tracking\_Mobile}, \texttt{Tracking\_MqttBridge},
  \texttt{Tracking\_PostgreSQL}, \texttt{Tracking\_EMQX},
  \texttt{Tracking\_VictoriaMetrics}, \texttt{Tracking\_VictoriaLogs},
  \texttt{Tracking\_Grafana}, \texttt{Tracking\_NPM},
  \texttt{Tracking\_Data}).
\item
  \texttt{resources/}: tài liệu dự án, báo cáo, kế hoạch.
\item
  \texttt{README.md}: tệp hướng dẫn tổng quan về cách dựng và vận hành hệ
  thống.
\end{itemize}

\subsection{4.2. Tham chiếu cơ sở dữ liệu (Database
Schema)}\label{42-tham-chiux1ebfu-cux1a1-sux1edf-dux1eef-liux1ec7u-database-schema}

Schema PostgreSQL đầy đủ được lưu tại các tệp SQL khởi tạo trong dịch vụ
PostgreSQL; phần phụ lục này chỉ tóm tắt các bảng chính phục vụ luận
văn.

\textbf{Các bảng chính trong PostgreSQL:}


{\thesissettableid{4.2A}{tab-ch4-tong-hop-thong-tin-ve-stt-ten-bang-mo-ta}
\def\LTcaptype{table}
\begin{longtable}{|L{0.115\linewidth}|L{0.198\linewidth}|L{0.312\linewidth}|L{0.330\linewidth}|}
\caption{Tổng hợp thông tin về STT, Tên bảng, Mô tả}\label{tab:ch4-tong-hop-thong-tin-ve-stt-ten-bang-mo-ta}\\
\hline \textbf{STT} & \textbf{Tên bảng} & \textbf{Mô tả} & \textbf{Quan hệ chính} \\
\\\hline
\endfirsthead
\caption[]{Tổng hợp thông tin về STT, Tên bảng, Mô tả (tiếp theo)}\\
\hline \textbf{STT} & \textbf{Tên bảng} & \textbf{Mô tả} & \textbf{Quan hệ chính} \\
\\\hline
\endhead
\\\hline
\endfoot
\endlastfoot
1 & \texttt{users} & Tài khoản người dùng & 1:N với sessions,
audit\_logs \\\hline
2 & \texttt{sessions} & Phiên đăng nhập (token hash SHA-256) & N:1 với
users \\\hline
3 & \texttt{vehicles} & Thông tin phương tiện & 1:1 với devices, 1:N với
trips \\\hline
4 & \texttt{devices} & Thông tin thiết bị theo dõi & N:1 với vehicles \\\hline
5 & \texttt{customers} & Thông tin khách hàng & 1:N với trips \\\hline
6 & \texttt{trips} & Hành trình cho thuê & N:1 với vehicles,
customers \\\hline
7 & \texttt{alerts} & Cảnh báo hệ thống & N:1 với vehicles, devices \\\hline
8 & \texttt{geofences} & Vùng cho phép & N:N với vehicles \\\hline
9 & \texttt{commands} & Lệnh điều khiển thiết bị & N:1 với devices \\\hline
10 & \texttt{audit\_logs} & Nhật ký thao tác & N:1 với users \\\hline
\end{longtable}
}


\textbf{Các chỉ số trong VictoriaMetrics:}


{\thesissettableid{4.2B}{tab-ch4-tong-hop-thong-tin-ve-metric-name-mau-labels-mo-ta}
\def\LTcaptype{table}
\begin{longtable}{|L{0.367\linewidth}|L{0.209\linewidth}|L{0.388\linewidth}|}
\caption{Tổng hợp thông tin về Metric name mẫu, Labels, Mô tả}\label{tab:ch4-tong-hop-thong-tin-ve-metric-name-mau-labels-mo-ta}\\
\hline \textbf{Metric name mẫu} & \textbf{Labels} & \textbf{Mô tả} \\
\\\hline
\endfirsthead
\caption[]{Tổng hợp thông tin về Metric name mẫu, Labels, Mô tả (tiếp theo)}\\
\hline \textbf{Metric name mẫu} & \textbf{Labels} & \textbf{Mô tả} \\
\\\hline
\endhead
\\\hline
\endfoot
\endlastfoot
\texttt{vehicle\_latitude} & \texttt{device\_id} & Vĩ độ GNSS \\\hline
\texttt{vehicle\_longitude} & \texttt{device\_id} & Kinh độ GNSS \\\hline
\texttt{vehicle\_speed} & \texttt{device\_id} & Tốc độ \\\hline
\texttt{vehicle\_vibration} & \texttt{device\_id} & Mức rung từ IMU \\\hline
\texttt{vehicle\_battery\_top} & \texttt{device\_id} & Điện áp nguồn
phía trên \\\hline
\texttt{vehicle\_battery\_bot} & \texttt{device\_id} & Điện áp pin dự
phòng \\\hline
\texttt{vehicle\_ignition} & \texttt{device\_id} & Trạng thái ignition
(0/1) \\\hline
\texttt{vehicle\_uptime} & \texttt{device\_id} & Thời gian chạy tích
lũy \\\hline
\end{longtable}
}


\subsection{4.3. Tài liệu API (API
Documentation)}\label{43-tuxe0i-liux1ec7u-api-api-documentation}

Tài liệu API đầy đủ được tạo tự động bằng Swagger/OpenAPI và có thể truy
cập tại:

\begin{itemize}
\tightlist
\item
  \textbf{URL (môi trường thử nghiệm):}
  \texttt{https://api.thingdock.dev/api-docs}
\item
  \textbf{URL (môi trường phát triển):}
  \texttt{http://localhost:4000/api-docs}
\end{itemize}

\textbf{Tổng hợp API endpoints chính:}


{\thesissettableid{4.3A}{tab-ch4-tong-hop-thong-tin-ve-nhom-method-endpoint}
\def\LTcaptype{table}
\begin{longtable}{|L{0.226\linewidth}|L{0.137\linewidth}|L{0.204\linewidth}|L{0.387\linewidth}|}
\caption{Tổng hợp thông tin về Nhóm, Method, Endpoint}\label{tab:ch4-tong-hop-thong-tin-ve-nhom-method-endpoint}\\
\hline \textbf{Nhóm} & \textbf{Method} & \textbf{Endpoint} & \textbf{Mô tả} \\
\\\hline
\endfirsthead
\caption[]{Tổng hợp thông tin về Nhóm, Method, Endpoint (tiếp theo)}\\
\hline \textbf{Nhóm} & \textbf{Method} & \textbf{Endpoint} & \textbf{Mô tả} \\
\\\hline
\endhead
\\\hline
\endfoot
\endlastfoot
Auth & POST & \nolinkurl{/api/v1/auth/login} & Đăng nhập \\\hline
Auth & POST & \nolinkurl{/api/v1/auth/logout} & Đăng xuất \\\hline
Auth & GET & \nolinkurl{/api/v1/auth/me} & Thông tin người dùng hiện tại \\\hline
Vehicles & GET & \nolinkurl{/api/v1/vehicles} & Danh sách phương tiện \\\hline
Vehicles & POST & \nolinkurl{/api/v1/vehicles} & Thêm phương tiện \\\hline
Vehicles & GET & \nolinkurl{/api/v1/vehicles/:id} & Chi tiết phương tiện \\\hline
Devices & GET & \nolinkurl{/api/v1/devices} & Danh sách thiết bị \\\hline
Devices & POST & \nolinkurl{/api/v1/devices} & Đăng ký thiết bị mới \\\hline
Devices & POST & \nolinkurl{/api/v1/devices/:id/ota} & Kích hoạt OTA cho
thiết bị \\\hline
Devices & POST & \nolinkurl{/api/v1/devices/:id/ota/rollback} & Yêu cầu
rollback firmware \\\hline
Telemetry & GET & \nolinkurl{/api/v1/devices/:id/telemetry} & Dữ liệu
telemetry theo thiết bị \\\hline
Telemetry & GET & \nolinkurl{/api/v1/telemetry/history} & Truy vấn lịch sử
telemetry \\\hline
Alerts & GET & \nolinkurl{/api/v1/alerts} & Danh sách cảnh báo \\\hline
Alerts & PUT & \nolinkurl{/api/v1/alerts/:id/acknowledge} & Xác nhận cảnh
báo \\\hline
Vùng cho phép & GET & \nolinkurl{/api/v1/geofences} & Danh sách vùng cho phép \\\hline
Vùng cho phép & POST & \nolinkurl{/api/v1/geofences} & Tạo vùng cho phép mới \\\hline
Commands & POST & \nolinkurl{/api/v1/devices/:id/command} & Gửi lệnh đến
thiết bị \\\hline
Trips & GET & \nolinkurl{/api/v1/trips} & Danh sách hành trình \\\hline
Customers & GET & \nolinkurl{/api/v1/customers} & Danh sách khách hàng \\\hline
\end{longtable}
}


\subsection{4.4. Tài liệu cấu hình môi trường (Environment
Configuration)}\label{44-tuxe0i-liux1ec7u-cux1ea5u-huxecnh-muxf4i-trux1b0ux1eddng-environment-configuration}

Mỗi dịch vụ yêu cầu file \texttt{.env} riêng. Mẫu cấu hình
(\texttt{.env.example}) có sẵn ở các dịch vụ có cung cấp file mẫu; riêng
Frontend dùng trực tiếp file \texttt{.env} theo tệp hướng dẫn đi kèm kho
mã nguồn.

\textbf{Các biến môi trường quan trọng:}


{\thesissettableid{4.4A}{tab-ch4-tong-hop-thong-tin-ve-bien-dich-vu-mo-ta}
\def\LTcaptype{table}
\begin{longtable}{|L{0.115\linewidth}|L{0.253\linewidth}|L{0.460\linewidth}|L{0.127\linewidth}|}
\caption{Tổng hợp thông tin về Biến, Dịch vụ, Mô tả}\label{tab:ch4-tong-hop-thong-tin-ve-bien-dich-vu-mo-ta}\\
\hline \textbf{Biến} & \textbf{Dịch vụ} & \textbf{Mô tả} & \textbf{Bắt buộc} \\
\\\hline
\endfirsthead
\caption[]{Tổng hợp thông tin về Biến, Dịch vụ, Mô tả (tiếp theo)}\\
\hline \textbf{Biến} & \textbf{Dịch vụ} & \textbf{Mô tả} & \textbf{Bắt buộc} \\
\\\hline
\endhead
\\\hline
\endfoot
\endlastfoot
\nolinkurl{POSTGRESQL_HOST} & Backend & Địa chỉ PostgreSQL server & Có \\\hline
\nolinkurl{POSTGRESQL_PORT} & Backend & Cổng PostgreSQL (mặc định: 5432) &
Có \\\hline
\nolinkurl{POSTGRESQL_DATABASE} & Backend & Tên cơ sở dữ liệu & Có \\\hline
\nolinkurl{POSTGRESQL_USER} & Backend & Tên đăng nhập PostgreSQL & Có \\\hline
\nolinkurl{POSTGRESQL_PASSWORD} & Backend & Mật khẩu PostgreSQL & Có \\\hline
\nolinkurl{VICTORIAMETRICS_URL} & MqttBridge & URL VictoriaMetrics & Có \\\hline
\nolinkurl{VICTORIALOGS_URL} & MqttBridge & URL VictoriaLogs & Có \\\hline
\nolinkurl{MQTT_BROKER_URL} & Backend, MqttBridge & URL EMQX broker &
Có \\\hline
\nolinkurl{FIRMWARE_PUBLIC_BASE_URL} & Backend & Base URL công khai để
thiết bị tải file OTA & Có cho OTA \\\hline
\nolinkurl{SESSION_SECRET} & Backend & Khóa bí mật cho phiên đăng nhập (min 32 ký
tự) & Có \\\hline
\nolinkurl{CORS_ORIGIN} & Backend & URL frontend cho CORS & Có \\\hline
\end{longtable}
}


\textbf{Lưu ý bảo mật:} Tất cả mật khẩu và khóa bí mật đều bắt buộc, không
có giá trị mặc định. Hệ thống sẽ từ chối khởi động nếu thiếu bất kỳ biến
nào.

\subsection{4.5. Hướng dẫn cài đặt và chạy hệ
thống}\label{45-hux1b0ux1edbng-dux1eabn-cuxe0i-ux111ux1eb7t-vuxe0-chux1ea1y-hux1ec7-thux1ed1ng}

\textbf{Yêu cầu hệ thống:}


{\thesissettableid{4.5A}{tab-ch4-tong-hop-thong-tin-ve-yeu-cau-phien-ban-toi-thieu}
\def\LTcaptype{table}
\begin{longtable}{|L{0.441\linewidth}|L{0.527\linewidth}|}
\caption{Tổng hợp thông tin về Yêu cầu, Phiên bản tối thiểu}\label{tab:ch4-tong-hop-thong-tin-ve-yeu-cau-phien-ban-toi-thieu}\\
\hline \textbf{Yêu cầu} & \textbf{Phiên bản tối thiểu} \\
\\\hline
\endfirsthead
\caption[]{Tổng hợp thông tin về Yêu cầu, Phiên bản tối thiểu (tiếp theo)}\\
\hline \textbf{Yêu cầu} & \textbf{Phiên bản tối thiểu} \\
\\\hline
\endhead
\\\hline
\endfoot
\endlastfoot
Node.js & 20.x LTS \\\hline
Docker & 24.x \\\hline
Docker Compose & 2.x \\\hline
Git & 2.x \\\hline
npm & 10.x \\\hline
\end{longtable}
}


\textbf{Các bước cài đặt:}

\begin{Shaded}
\begin{Highlighting}[]
\CommentTok{\# 1. Clone kho mã nguồn}
\FunctionTok{git}\NormalTok{ clone https://github.com/anmh1205/IoT\_Vehicle\_Tracking\_System.git}
\BuiltInTok{cd}\NormalTok{ IoT\_Vehicle\_Tracking\_System}

\CommentTok{\# 2. Tạo Docker network}
\ExtensionTok{docker}\NormalTok{ network create tracking{-}network}

\CommentTok{\# 3. Khởi động infrastructure services}
\BuiltInTok{cd}\NormalTok{ iot{-}vehicle{-}tracking{-}system{-}cloud/Tracking\_PostgreSQL }\KeywordTok{\&\&} \ExtensionTok{docker}\NormalTok{ compose up }\AttributeTok{{-}d}
\BuiltInTok{cd}\NormalTok{ ../Tracking\_EMQX }\KeywordTok{\&\&} \ExtensionTok{docker}\NormalTok{ compose up }\AttributeTok{{-}d}
\BuiltInTok{cd}\NormalTok{ ../Tracking\_VictoriaMetrics }\KeywordTok{\&\&} \ExtensionTok{docker}\NormalTok{ compose up }\AttributeTok{{-}d}
\BuiltInTok{cd}\NormalTok{ ../Tracking\_VictoriaLogs }\KeywordTok{\&\&} \ExtensionTok{docker}\NormalTok{ compose up }\AttributeTok{{-}d}
\BuiltInTok{cd}\NormalTok{ ../Tracking\_Grafana }\KeywordTok{\&\&} \ExtensionTok{docker}\NormalTok{ compose up }\AttributeTok{{-}d}

\CommentTok{\# 4. Cấu hình .env cho Backend}
\BuiltInTok{cd}\NormalTok{ ../Tracking\_Backend}
\FunctionTok{cp}\NormalTok{ .env.example .env}
\CommentTok{\# Chỉnh sửa .env với các giá trị phù hợp}

\CommentTok{\# 5. Cài đặt và chạy Backend}
\ExtensionTok{npm}\NormalTok{ ci}
\ExtensionTok{npm}\NormalTok{ run dev}

\CommentTok{\# 6. Cấu hình và chạy Frontend}
\BuiltInTok{cd}\NormalTok{ ../Tracking\_Frontend}
\CommentTok{\# Chuẩn bị file .env theo tệp hướng dẫn}
\ExtensionTok{npm}\NormalTok{ ci}
\ExtensionTok{npm}\NormalTok{ run dev}

\CommentTok{\# 7. Chạy MQTT Bridge}
\BuiltInTok{cd}\NormalTok{ ../Tracking\_MqttBridge}
\FunctionTok{cp}\NormalTok{ .env.example .env}
\ExtensionTok{npm}\NormalTok{ ci}
\ExtensionTok{npm}\NormalTok{ run dev}
\end{Highlighting}
\end{Shaded}

\textbf{Truy cập hệ thống:}


{\thesissettableid{4.5B}{tab-ch4-tong-hop-thong-tin-ve-dich-vu-url-tai-khoan-mac-inh}
\def\LTcaptype{table}
\begin{longtable}{|L{0.320\linewidth}|L{0.160\linewidth}|L{0.485\linewidth}|}
\caption{Tổng hợp thông tin về Dịch vụ, URL, Tài khoản mặc định}\label{tab:ch4-tong-hop-thong-tin-ve-dich-vu-url-tai-khoan-mac-inh}\\
\hline \textbf{Dịch vụ} & \textbf{URL} & \textbf{Tài khoản mặc định} \\
\\\hline
\endfirsthead
\caption[]{Tổng hợp thông tin về Dịch vụ, URL, Tài khoản mặc định (tiếp theo)}\\
\hline \textbf{Dịch vụ} & \textbf{URL} & \textbf{Tài khoản mặc định} \\
\\\hline
\endhead
\\\hline
\endfoot
\endlastfoot
Giao diện web & \nolinkurl{https://thingdock.dev} & Tài khoản quản trị được cấu hình riêng khi triển khai \\\hline
API phía máy chủ & \nolinkurl{https://api.thingdock.dev} & --- \\\hline
Tài liệu API (Swagger) & \nolinkurl{https://api.thingdock.dev/api-docs} & --- \\\hline
Kênh MQTT qua WebSocket & \nolinkurl{wss://mqtt.thingdock.dev} & --- \\\hline
broker MQTT qua TLS & \nolinkurl{mqtts://mqtts.thingdock.dev:8883} & Theo tài khoản MQTT của thiết bị \\\hline
Grafana & \nolinkurl{https://grafana.thingdock.dev} & Tài khoản quản trị được cấu hình riêng khi triển khai \\\hline
Bảng điều khiển EMQX & \nolinkurl{https://emqx.thingdock.dev} & Tài khoản quản trị được cấu hình riêng khi triển khai \\\hline
Trang quản trị NPM & \nolinkurl{https://npm.thingdock.dev} & Tài khoản quản trị được cấu hình khi khởi tạo hệ thống \\\hline
\end{longtable}
}

\textbf{Ghi chú môi trường thử nghiệm:} VictoriaMetrics được giữ ở chế độ
nội bộ, không công khai qua tên miền phụ. Khi cần phục vụ vận hành, dịch
vụ này được truy cập từ mạng Docker nội bộ hoặc có thể công bố thêm một
tên miền phụ riêng ở giai đoạn sau.


\begin{center}\rule{0.5\linewidth}{0.5pt}\end{center}

\end{document}





